\section{黄金分支折叠：从微态到稳定类型的多尺度投影}\label{sec:folding}

本节构造并分析分辨率折叠算子 \(\Fold_m:\Omega_m\to X_m\)，其中 \(\Omega_m=\{0,1\}^m\) 为长度 \(m\) 的窗口微态空间，而 \(X_m\subseteq\Omega_m\) 为黄金均值约束的稳定语法空间（禁止相邻 \(11\)）。折叠的目标是把指数增长的微态空间压缩到可审计的稳定类型层，同时保证跨分辨率的一致性，从而使后续的统计审计、地址化与算术结构可以在同一规范代表系统上闭合。\par
下述子小节依次给出：Fibonacci/稳定语法的计数结构、Zeckendorf 规范形与折叠的良定义、跨分辨率相容（逆系统/逆极限）与多尺度残差接口，以及由 \(\Fold_m\) 诱导的序列级滑动块码 \(\Phi_m\) 的有限图表示与可计算实现。

\subsection{Fibonacci 权重与稳定语法集合}\label{subsec:folding-fibonacci-stable-syntax}

令 Fibonacci 数列 $(F_k)_{k\ge1}$ 满足
\[
F_1=F_2=1,\quad F_{k+2}=F_{k+1}+F_k.
\]
定义长度 $m$ 的稳定类型集合
\[
X_m:=\{w\in\{0,1\}^m:\text{不含子词 }11\}.
\]
即 $X_m$ 为禁止相邻 $1$ 的二元语法（亦称“黄金均值约束”）。其计数满足
\[
|X_m|=F_{m+2},
\]
见 \cite{LindMarcus1995}。

\begin{proposition}[末位分解与 Fibonacci 递推矩阵]\label{prop:folding-stable-syntax-terminal-recursion}
对每个 $m\ge 1$，定义
\[
X_m^{(s)}:=\{w\in X_m:\ w_m=s\},\qquad s\in\{0,1\},
\]
并令计数向量
\[
v_m:=\begin{pmatrix}\card{X_m^{(1)}}\\[2pt]\card{X_m^{(0)}}\end{pmatrix}\in\ZZ_{\ge0}^{\,2}.
\]
则有严格递推
\[
v_{m+1}=N_\tau\,v_m,\qquad
N_\tau:=\begin{pmatrix}0&1\\[2pt]1&1\end{pmatrix}.
\]
\leanverified{card\_recurrence}
\leanverified{paper\_folding\_stable\_syntax\_terminal\_recursion}
\end{proposition}
\begin{proof}
若 $w\in X_{m+1}$ 以 $1$ 结尾，则倒数第二位必为 $0$，故去掉末位得到 $X_m^{(0)}$ 中的唯一词；反之对任意 $u\in X_m^{(0)}$，拼接 $u1$ 仍不含子词 $11$。因此 $\card{X_{m+1}^{(1)}}=\card{X_m^{(0)}}$。
\par
另一方面，拼接末位 $0$ 不引入 $11$，故从任意 $u\in X_m$ 得到 $u0\in X_{m+1}^{(0)}$，并且该对应为双射，因而
\[
\card{X_{m+1}^{(0)}}=\card{X_m}=\card{X_m^{(0)}}+\card{X_m^{(1)}}.
\]
合并两式即得所述矩阵递推。
\leanverified{card\_X\_endingTrue}
\leanverified{card\_endsInOne\_eq\_fib}
\end{proof}

\begin{theorem}[两态骨架与 Fibonacci 融合环的严格同构]\label{thm:folding-stable-syntax-fib-fusion-ring}
记 \(\mathrm{Fib}\) 为 Fibonacci 融合范畴的 Grothendieck 环
\[
K_0(\mathrm{Fib})=\ZZ\langle \ind,\tau\rangle/(\tau^2=\ind+\tau),
\]
并在基 \((\ind,\tau)\) 下记“左乘 \(\tau\)”的融合矩阵为 \(N_\tau\)。
则命题 \ref{prop:folding-stable-syntax-terminal-recursion} 中出现的矩阵 \(N_\tau\) 与该融合矩阵一致；并且对每个 \(m\ge 1\) 有
\[
\tau^{m+1}
=\card{X_m^{(1)}}\,\ind+\card{X_m^{(0)}}\,\tau
\qquad\text{于 }K_0(\mathrm{Fib}).
\]
\leanverified{goldenMeanAdjacency\_irreducible}
\leanverified{goldenMeanAdjacency\_pow\_nonneg}
\leanverified{goldenMeanAdjacency\_pow\_positive}
\leanverified{paper\_lucas\_add\_formula}
\leanverified{paper\_folding\_stable\_syntax\_fib\_fusion\_ring}
\end{theorem}
\begin{proof}
由 \(\tau\cdot\ind=\tau\) 与 \(\tau\cdot\tau=\ind+\tau\)，在基 \((\ind,\tau)\) 下“左乘 \(\tau\)”的矩阵正是
\(
N_\tau=\begin{pmatrix}0&1\\[2pt]1&1\end{pmatrix}
\)。
从而 \(\tau^{m+2}= \tau\cdot \tau^{m+1}\) 的系数向量满足同一递推 \(v_{m+1}=N_\tau v_m\)。
初值由 \(\tau^2=\ind+\tau\) 与 \(v_1=(\card{X_1^{(1)}},\card{X_1^{(0)}})^{\mathsf T}=(1,1)^{\mathsf T}\) 一致，故对一切 \(m\) 得到所述恒等式。
\end{proof}

\begin{corollary}[熵率与 Perron--Frobenius 维数的对数一致]\label{cor:folding-stable-syntax-entropy-logqdim}
设 \(X_\infty\subset\{0,1\}^{\ZZ}\) 为禁止子词 \(11\) 的双边子移位，则其拓扑熵满足
\[
h_{\mathrm{top}}(X_\infty)=\log\varphi,
\]
其中 \(\varphi=\tfrac{1+\sqrt5}{2}\)。
此外，矩阵 \(N_\tau\) 的 Perron--Frobenius 特征值为 \(\varphi\)；因此在任意实现该融合环的单位ary融合范畴中，\(\tau\) 的量子维数满足 \(d(\tau)=\varphi\)，并有
\[
h_{\mathrm{top}}(X_\infty)=\log d(\tau).
\]
\end{corollary}
\begin{proof}
第一式是黄金均值 SFT 的标准结果，亦可由 \(|X_m|=F_{m+2}\) 的渐近给出 \(\lim_{m\to\infty}\frac1m\log|X_m|=\log\varphi\)。
第二部分由 \(N_\tau\) 的特征多项式 \(\lambda^2-\lambda-1\) 及单位ary融合范畴中量子维数向量为融合矩阵的 Perron--Frobenius 本征向量的事实推出。
\leanverified{topological\_entropy\_eq\_log\_phi}
\leanverified{fib\_ratio\_tendsto\_golden}
\leanverified{foldIndex\_tendsto\_log\_two\_div\_phi}
\leanverified{paper\_folding\_stable\_syntax\_entropy\_logqdim\_seeds}
\leanverified{paper\_folding\_stable\_syntax\_entropy\_logqdim\_package}
\leanverified{paper\_folding\_stable\_syntax\_entropy\_logqdim}
\end{proof}

\begin{remark}[投影子控制的多通道融合与算符代数输入]\label{rem:folding-stable-syntax-noninvertible-links}
定理 \ref{thm:folding-stable-syntax-fib-fusion-ring} 表明：稳定核的两态扩展骨架在 \(K_0\) 层面严格实现 Fibonacci 融合规则 \(\tau^2=\ind+\tau\)，其非可逆性由“同一输入可产生多种融合通道”的分解体现。
该机制与非可逆对称的投影子控制融合结构相容 \cite{ZhangHuangWangYe2026NonInvertibleSymmetriesMixedAnomaliesTwistedBF}，并与“自对偶非可逆对象诱导 Temperley--Lieb 结构”的可积链构造一致 \cite{BlakeneyCorcoranDeLeeuwPozsgayVernier2025TemperleyLiebIntegrableModelsFusionCategories}。
另一方面，折叠粗粒化所诱导的有限指数条件期望与基本构造在算符代数路线中提供了把该基环外部化为可计算双模数据的自然入口 \cite{EvansJones2025OperatorAlgebraicApproachFusionCategorySymmetryLattice}；本稿在小节 \ref{subsec:op-algebra-fold-watatani-index} 中给出对应的 Watatani 指标元与块分解闭式。
\end{remark}

\begin{conjecture}[折叠—限制重整化的 Fibonacci 普适吸引子]\label{conj:folding-fib-renormalization-attractor}
在保持稳定核两态骨架（即 \(X_m\) 的末位扩展图与矩阵 \(N_\tau\)）不变的约束下，考虑由“尺度提升—再折叠—再限制”组合生成的重整化变换 \(\mathcal R\) 作用于一类局部粗粒化（例如有限指数条件期望或等价的局部势函数族）。
猜想存在一条由 \(\tau^2=\ind+\tau\) 所锁定的 Fibonacci 固定点普适类，使得凡满足该两态骨架约束且无额外对称破缺的轨道均在适当归一化下收敛到该类；其极限缺陷数据可由 SymTFT 分解/融合通道投影子语言一致刻画 \cite{SchaferNamekiTiwariWarmanZhang2025SymTFTMixedStatesNonInvertible}。
\end{conjecture}


\subsection{Zeckendorf 规范形}\label{subsec:folding-zeckendorf-canonical}

Zeckendorf 定理断言：对任意 $N\in\NN$，存在唯一表示（\cite{Zeckendorf1972}）
\[
N=\sum_{k\ge1}c_k F_{k+1},\qquad c_k\in\{0,1\},\quad c_kc_{k+1}=0.
\]
将 $(c_k)_{k\ge1}$ 视为 $N$ 的 Zeckendorf 数字序列，记为 $Z(N)=(c_1,c_2,\dots)$。

定义截断映射
\[
\pi_m(c_1,c_2,\dots)=(c_1,\dots,c_m)\in X_m.
\]


\subsection{多尺度一致性：restriction 逆系统与相容性口径}\label{subsec:folding-multiscale}

当比较不同窗口长度 $m$ 的输出时，“一致”应被精确定义为跨尺度相容性：高分辨率输出经由自然的截断（restriction）投影到低分辨率后，应复现低分辨率输出所包含的内容。下面给出对象层的标准构造（完整证明见附录 G）。

\subsubsection{折叠与因子化的核心定理（汇总陈述）}\label{subsec:fold-theorem-suite}
为统一折叠接口的可审计性要求，本小节将折叠引起的关键结构归纳为一个汇总定理：它给出折叠的精确定义、归约顺序无关（终止/合流）、幂等/满射、多尺度 restriction 逆系统相容、序列级 sofic 因子与最小右分辨覆盖，以及差异度证书到折叠直方图 TV 偏差的定量链条。

\begin{theorem}[有限窗口折叠：合流、逆系统相容、sofic 因子与定量证书]\label{thm:fold-suite}
设 $\Omega_m=\{0,1\}^m$，$X_m=\{w\in\{0,1\}^m:\text{不含子词 }11\}$，并令 $\Fold_m:\Omega_m\to X_m$ 如定义 \ref{def:fold-word}。则：
\begin{enumerate}
\normalfont
  \item \textbf{（归约顺序无关的规范化）}折叠可由值保持的局部重写系统实现，且该系统强终止并合流；因此 $\Fold_m$ 的实现与归约顺序无关。
  \item \textbf{（幂等与满射）}\(\Fold_m(\omega)\in X_m\)；并且 \(\Fold_m|_{X_m}=\mathrm{id}\)、\(\Fold_m:\Omega_m\twoheadrightarrow X_m\)。
  \item \textbf{（restriction 逆系统与逆极限）}对任意 $m_2\ge m_1$，截断 $\pi_{m_2\to m_1}:X_{m_2}\to X_{m_1}$ 构成 restriction 逆系统，且存在自然同构 \(X_\infty\cong\varprojlim X_m\)（定理 \ref{thm:inverse-limit-golden}）。
\item \textbf{（跨分辨率相容的交换图）}直接微态截断一般不与折叠交换；但定义 \ref{subsec:fold-omega-restriction} 的折叠感知 restriction \(\widetilde r_{m+1\to m}(\omega)=\pi_{m+1\to m}(\Fold_{m+1}(\omega))\) 使得交换图成立（命题 \ref{prop:fold-omega-commute}）。
  \item \textbf{（序列级因子与显式 sofic 图像）}由 \(\Fold_m\) 沿滑动窗口诱导的 \(\Phi_m\) 是滑动块码，记忆为 $0$、预期性为 $m-1$（命题 \ref{prop:Phi_m-radius}）；其像空间 \(Y_m\subset X_m^{\ZZ}\) 是 sofic 子移位，并可由显式标记图 \(G_m\) 给出（定理 \ref{thm:Phi_m-sofic-graph}）。
  \item \textbf{（最小右分辨覆盖）}对 \(Y_m\)，右分辨表示可由子集确定化得到（命题 \ref{prop:Phi_m-right-resolving}）；进一步按 follower 集最小化得到右 Fischer cover，其最小性可参见 \cite{LindMarcus1995,Kitchens1998SymbolicDynamics}。
  \item \textbf{（定量证书：差异度 \(\Rightarrow\) 折叠直方图 TV 偏差）}在旋转扫描（Sturmian）模型下，对任意 $N\ge 1$，
  \[
  D_{\mathrm{TV}}(\hat{\pi}_{m,N},\pi_m^{\mathrm{fold}})\ \le\ (m+1)\,D_N^*(\alpha;x_0),
  \]
  且到 Parry 基线的误差分解满足推论 \ref{cor:tv-decomposition-parry}。
\end{enumerate}
\leanpartial{Fold\_stable, Fold\_idempotent, Fold\_surjective, step\_confluent, inverseLimitEquiv}{项 (1)--(3) 已形式化；项 (4)--(7)（sofic 因子、右分辨覆盖、定量证书）尚未形式化}
\leanverified{paper\_Fold\_fundamental\_triple}
\leanverified{paper\_fold\_suite}
\end{theorem}

\begin{proof}
(1) 由命题 \ref{prop:fold-rewrite-newman}（强终止 + 合流）与推论 \ref{cor:foldm-order-indep}。

(2) 由命题 \ref{prop:fold-basic}。

(3) restriction 一致性由引理 \ref{lem:pi-functorial-golden}；逆极限同构由定理 \ref{thm:inverse-limit-golden}（完整证明见附录 G 的定理 \ref{app:multiscale-inverse-limit}）。

(4) 由小节 \ref{subsec:fold-omega-restriction} 的定义与命题 \ref{prop:fold-omega-commute}。

(5) 半径由命题 \ref{prop:Phi_m-radius}；sofic 图像与显式标记图由定理 \ref{thm:Phi_m-sofic-graph}。

(6) 由命题 \ref{prop:Phi_m-right-resolving} 与右 Fischer cover 的最小性（参见 \cite{LindMarcus1995,Kitchens1998SymbolicDynamics}）。

(7) 由第 \ref{sec:experiments} 节的定理 \ref{thm:tv-certificate-hist}（其证明给出显式系数 $m+1$），并由推论 \ref{cor:tv-decomposition-parry} 得到到 Parry 基线的误差分解。
\end{proof}

\begin{definition}[尺度投影（restriction）]\label{def:pi-m2-m1-golden}
对任意 $m_2\ge m_1$，定义截断投影
\[
\pi_{m_2\to m_1}:X_{m_2}\to X_{m_1},\qquad
\pi_{m_2\to m_1}(x_1,\dots,x_{m_2})=(x_1,\dots,x_{m_1}).
\]
\end{definition}

\begin{lemma}[一致性]\label{lem:pi-functorial-golden}
对任意 $m_3\ge m_2\ge m_1$ 有
\[
\pi_{m_3\to m_1}=\pi_{m_2\to m_1}\circ \pi_{m_3\to m_2}.
\]
并且对任意无穷合法序列 $c=(c_1,c_2,\dots)$ 有
\[
\pi_{m_2\to m_1}\bigl(\pi_{m_2}(c)\bigr)=\pi_{m_1}(c).
\]
\end{lemma}

\leanverified{paper\_pi\_functorial\_golden}

\begin{theorem}[逆极限同构（分辨率无关对象）]\label{thm:inverse-limit-golden}
令
\[
X_\infty:=\{(c_1,c_2,\dots)\in\{0,1\}^{\mathbb{N}}:\ c_kc_{k+1}=0\ \forall k\ge 1\}.
\]
则存在自然同构
\[
X_\infty\ \cong\ \varprojlim X_m,
\]
并且 $X_m$ 可视为 $X_\infty$ 的有限截断因子：不同 $m$ 的差异仅在”可见前缀长度”上，而非在统一对象本身上。
\leanverified{inverseLimitEquiv}
\leanverified{paper\_inverseLimit\_golden}
\leanverified{XInfinity\_nonempty}
\leanverified{XInfinity\_univ\_isCompact}
\leanverified{XInfinity\_ext\_iff}
\leanverified{paper\_XInfinity\_topology\_package}
\end{theorem}

\subsubsection{关于微态截断与折叠的相容性：正确的交换图}\label{subsec:fold-omega-restriction}
上述逆系统结构发生在稳定语法集合 $\{X_m\}$ 上，其 restriction 映射是定义 \ref{def:pi-m2-m1-golden} 的截断 $\pi_{m_2\to m_1}:X_{m_2}\to X_{m_1}$。一个需要澄清的问题是：能否在微态空间 $\Omega_m=\{0,1\}^m$ 上直接用“去掉末位”的截断与折叠 $\Fold_m$ 形成交换图？

\paragraph{直接截断一般不与折叠交换}
令
\[
r_{m+1\to m}:\Omega_{m+1}\to\Omega_m,\qquad
r_{m+1\to m}(\omega_1\cdots\omega_m\omega_{m+1}):=\omega_1\cdots\omega_m
\]
为微态上的直接截断。一般地，
\[
\Fold_m\circ r_{m+1\to m}\ \neq\ \pi_{m+1\to m}\circ \Fold_{m+1}.
\]
例如取 $m=1$、$\omega=11$，则 $r_{2\to 1}(11)=1$，从而 $(\Fold_1\circ r_{2\to 1})(11)=\Fold_1(1)=1$；另一方面 $N(11)=F_2+F_3=3$ 的 Zeckendorf 规范形为 $3=F_4$，故 $\Fold_2(11)=00$，从而 $(\pi_{2\to 1}\circ \Fold_2)(11)=0$。

\paragraph{折叠感知的 restriction（在 $\Omega$ 上实现交换）}
为在微态层面得到严格交换的形式化，我们定义一个与折叠相容的 restriction：
\[
\widetilde r_{m+1\to m}:\Omega_{m+1}\to \Omega_m,\qquad
\widetilde r_{m+1\to m}(\omega):=\pi_{m+1\to m}\bigl(\Fold_{m+1}(\omega)\bigr).
\]
注意 $\widetilde r_{m+1\to m}(\omega)\in X_m\subset\Omega_m$，因此该定义确实给出一条 $\Omega_{m+1}\to\Omega_m$ 的映射（它不是“原始截断”，而是“先折叠后截断”的可审计粗化）。

\begin{proposition}[跨分辨率相容的交换图]\label{prop:fold-omega-commute}
对任意 $m\ge 1$，下图交换：
\[
\begin{CD}
\Omega_{m+1} @>{\Fold_{m+1}}>> X_{m+1}\\
@V{\widetilde r_{m+1\to m}}VV @VV{\pi_{m+1\to m}}V\\
\Omega_m @>{\Fold_m}>> X_m.
\end{CD}
\]
即对任意 $\omega\in\Omega_{m+1}$ 有
\[
\Fold_m\bigl(\widetilde r_{m+1\to m}(\omega)\bigr)
\ =\
\pi_{m+1\to m}\bigl(\Fold_{m+1}(\omega)\bigr).
\]
\leanverified{paper\_fold\_omega\_commute}
\end{proposition}

\begin{proof}
由定义，$\widetilde r_{m+1\to m}(\omega)=\pi_{m+1\to m}(\Fold_{m+1}(\omega))\in X_m$。而命题 \ref{prop:fold-basic} 的幂等性给出：对任意 $x\in X_m$，$\Fold_m(x)=x$。因此
\[
\Fold_m\bigl(\widetilde r_{m+1\to m}(\omega)\bigr)
=\Fold_m\!\left(\pi_{m+1\to m}(\Fold_{m+1}(\omega))\right)
=\pi_{m+1\to m}(\Fold_{m+1}(\omega)).
\]
证毕。
\end{proof}

\paragraph{交换缺陷与离散斯托克斯恒等式}
折叠与截断的不交换性可以在每一步降分辨率上以一个可审计的向量缺陷编码，并在任意尺度区间上给出严格的望远镜分解。

\begin{definition}[局部交换缺陷与曲率指示]\label{def:fold-local-curvature-defect}
对 \(m\ge 1\) 与 \(\eta\in\Omega_{m+1}\)，定义局部交换缺陷向量
\[
\kappa_{m+1\to m}(\eta)
:=\Fold_m\!\bigl(r_{m+1\to m}(\eta)\bigr)\ \oplus\ \pi_{m+1\to m}\!\bigl(\Fold_{m+1}(\eta)\bigr)
\in \{0,1\}^{m}.
\]
并定义其零/非零指示
\[
K_{m+1\to m}(\eta):=\ind_{\{\kappa_{m+1\to m}(\eta)\neq 0\}}.
\]
\leanverified{localDefect}
\leanverified{Fold\_truncate\_eq\_restrict\_iff}
\end{definition}

\begin{definition}[全局交换缺陷]\label{def:fold-global-stokes-defect}
令 \(n>m\)。对任意 \(\omega\in\Omega_n\) 定义全局交换缺陷
\[
D_{n\to m}(\omega)
:=\Fold_m\!\bigl(r_{n\to m}(\omega)\bigr)\ \oplus\ \pi_{n\to m}\!\bigl(\Fold_n(\omega)\bigr)
\in \{0,1\}^{m}.
\]
\leanverified{globalDefect}
\leanverified{globalDefect\_antisymmetric}
\end{definition}

\begin{theorem}[折叠--截断塔上的离散广义斯托克斯恒等式]\label{thm:fold-discrete-stokes-defect}
对任意 \(n>m\) 与任意 \(\omega\in\Omega_n\)，成立严格恒等式
\[
\boxed{\;
D_{n\to m}(\omega)
=\bigoplus_{j=m}^{n-1}
\pi_{j\to m}\!\left(
\kappa_{j+1\to j}\!\bigl(r_{n\to j+1}(\omega)\bigr)
\right)\; }.
\]
\leanverified{globalDefect\_eq\_defectChain}
\leanverified{paper\_fold\_discrete\_stokes\_defect}
\leanverified{paper\_localDefect\_eq\_zero\_iff\_fold\_commutes}
\leanverified{globalDefect\_zero\_of\_all\_local\_zero}
\end{theorem}

\begin{proof}
对 \(j\in\{m,\dots,n\}\) 定义
\[
A_j(\omega):=\Fold_j\!\bigl(r_{n\to j}(\omega)\bigr),\qquad
B_j(\omega):=\pi_{n\to j}\!\bigl(\Fold_n(\omega)\bigr),\qquad
\Delta_j(\omega):=A_j(\omega)\oplus B_j(\omega).
\]
由定义 \(\Delta_m(\omega)=D_{n\to m}(\omega)\)，且 \(\Delta_n(\omega)=0\)。
由 \(\pi\) 的函子性得 \(B_{j-1}(\omega)=\pi_{j\to j-1}(B_j(\omega))\)；由 \(r\) 的复合一致性得 \(r_{n\to j-1}=r_{j\to j-1}\circ r_{n\to j}\)，从而
\(
A_{j-1}(\omega)=\Fold_{j-1}\!\bigl(r_{j\to j-1}(r_{n\to j}(\omega))\bigr)
\)。
于是对 \(j=m+1,\dots,n\),
\[
\begin{aligned}
\Delta_{j-1}(\omega)
&=A_{j-1}(\omega)\oplus B_{j-1}(\omega)\\
&=\Bigl(A_{j-1}(\omega)\oplus \pi_{j\to j-1}(A_j(\omega))\Bigr)\ \oplus\ \pi_{j\to j-1}\Bigl(A_j(\omega)\oplus B_j(\omega)\Bigr),
\end{aligned}
\]
其中括号内第一项按定义 \ref{def:fold-local-curvature-defect} 恰为
\(
\kappa_{j\to j-1}\!\bigl(r_{n\to j}(\omega)\bigr)
\)。
因此得到递推式
\[
\Delta_{j-1}(\omega)
=\kappa_{j\to j-1}\!\bigl(r_{n\to j}(\omega)\bigr)\ \oplus\ \pi_{j\to j-1}\!\bigl(\Delta_j(\omega)\bigr).
\]
从 \(j=n,n-1,\dots,m+1\) 递推，并使用 \(\pi\) 对逐位异或的同态性，即得
\[
\Delta_m(\omega)=\bigoplus_{j=m}^{n-1}\pi_{j\to m}\!\left(\kappa_{j+1\to j}\!\bigl(r_{n\to j+1}(\omega)\bigr)\right),
\]
而 \(\Delta_m(\omega)=D_{n\to m}(\omega)\)。证毕。
\end{proof}
\begin{theorem}[离散广义 Stokes 缺陷的谱混合与 Haar 极限]\label{thm:fold-stokes-defect-haar-mixing}
\leanverified{paper\_fold\_stokes\_defect\_haar\_mixing}
其证明通过字符扭曲的转移算子给出可核验的谱半径收缩机制，参见 \cite{Internal2026FoldGeneralizedStokesDefectHaarMixing} 的同一口径。
固定 \(m\ge 1\)，并令 \(\omega\) 在 \(\Omega_n=\{0,1\}^n\) 上服从均匀分布（\(n\ge m\)）。
记 \(D_{n\to m}\) 为定义 \ref{def:fold-global-stokes-defect} 的全局交换缺陷随机变量。
令 \(\mathsf H_m\le (\ZZ/2)^m\) 为由下列可达投影增量生成的子群：
\[
\mathsf H_m
:=
\Bigl\langle\ \pi_{j\to m}\bigl(\kappa_{j+1\to j}(\eta)\bigr)\ :\ j\ge m,\ \eta\in\Omega_{j+1}\ \Bigr\rangle.
\]
设满足如下两条可核验假设：
\begin{enumerate}
  \item 由 \(\Fold_\infty\) 的有限状态在线实现，对两条路径（先截断后折叠 vs.\ 先折叠后投影）作同步直积并附加 \((\ZZ/2)^m\)-值增量读出后，得到一有限状态系统；在 Bernoulli\((1/2)\) 驱动下，其隐藏状态 Markov 链在可达分量上原始（不可约且非周期）。
  \item 对任意非平凡字符 \(\chi\in\widehat{\mathsf H_m}\setminus\{1\}\)，存在一条闭合回路使该回路的总增量 \(g\in\mathsf H_m\) 满足 \(\chi(g)\neq 1\)。
\end{enumerate}
则存在常数 \(C_m>0\)、\(0<\eta_m<1\) 使对一切 \(n\ge m\) 有
\[
\bigl\|\mathcal L(D_{n\to m})-\mathrm{Haar}(\mathsf H_m)\bigr\|_{\mathrm{TV}}
\ \le\ C_m\,\eta_m^{\,n-m},
\]
其中 \(\mathrm{Haar}(\mathsf H_m)\) 为 \(\mathsf H_m\) 上均匀分布。
\end{theorem}

\begin{proof}
将定理 \ref{thm:fold-discrete-stokes-defect} 的异或分解视为一条群扩张上的可加观测量。
由 \(\Fold_\infty\) 的有限状态在线实现（常数延迟），对固定 \(m\) 把两条路径的计算过程作同步直积，并在每一步输出投影增量
\(
g_t\in(\ZZ/2)^m
\)；
于是得到有限状态集合 \(\mathcal Q_m\) 与确定性更新
\[
s_{t+1}=\delta(s_t,a_t),\qquad g_t=\psi(s_t,a_t)\in(\ZZ/2)^m,
\]
其中输入 \(a_t\in\{0,1\}\) 在均匀基线下独立同分布。
在该表示下，\(D_{n\to m}\) 与累积异或
\(
G_n:=g_0\oplus\cdots\oplus g_{n-m-1}
\)
相差至多有限步的端点修正；该修正可吸收进常数因子 \(C_m\)，故只需证 \(G_n\) 的混合结论。
令 \(\mathsf H_m=\langle\psi(s,a)\rangle\) 为可达增量像生成的子群，则 \(G_n\in\mathsf H_m\) 恒成立。
\par
对任意字符 \(\chi\in\widehat{\mathsf H_m}\) 定义扭曲转移矩阵
\[
(P_\chi)_{s,s'}:=\sum_{\substack{a\in\{0,1\}\\ \delta(s,a)=s'}}\frac12\,\chi(\psi(s,a)).
\]
当 \(\chi=1\) 时，\(P_1\) 即隐藏状态链的转移矩阵；按假设 (1) 它原始且谱半径为 \(1\)。
对一般 \(\chi\)，逐项估计给出 \(|(P_\chi)_{s,s'}|\le (P_1)_{s,s'}\)，从而 \(\rho(P_\chi)\le 1\)。
若对某个非平凡 \(\chi\) 仍有 \(\rho(P_\chi)=1\)，则沿所有贡献到主特征空间的闭合回路必须发生相位完全对齐，
等价于任意闭合回路总增量 \(g\) 均满足 \(\chi(g)=1\)；这与假设 (2) 矛盾。
故对一切 \(\chi\neq 1\) 必有 \(\rho(P_\chi)<1\)，并令 \(\eta_m:=\max_{\chi\neq 1}\rho(P_\chi)\in(0,1)\)。
\par
记 \(\widehat\mu_n(\chi):=\EE[\chi(G_n)]\) 为 \(G_n\) 的 Fourier 系数。
由增量表示与 \(P_\chi\) 的定义可得
\(
\widehat\mu_n(\chi)=\mathbf u_\chi^{\mathsf T}P_\chi^{\,n-m}\mathbf v
\)
其中 \(\mathbf u_\chi,\mathbf v\) 仅由初态与读出约定决定。
因此当 \(\chi\neq 1\) 时有指数衰减 \(|\widehat\mu_n(\chi)|\le C_m\,\eta_m^{\,n-m}\)（把向量范数吸收进 \(C_m\)）。
而 \(\mathrm{Haar}(\mathsf H_m)\) 在非平凡字符上的 Fourier 系数为 \(0\)。
对有限阿贝尔群，Fourier 反演与 Parseval 不等式把总变差距离控制在非平凡 Fourier 系数的上界之内，从而得到所示指数收敛。
将有限步端点修正吸收回 \(D_{n\to m}\) 的分布即可。证毕。
\end{proof}

\begin{corollary}[多尺度缺陷的奇偶可观测量极化与二项式噪声极限]\label{cor:fold-stokes-defect-parity-polarization}
\leanverified{paper\_fold\_stokes\_defect\_parity\_polarization}
在定理 \ref{thm:fold-stokes-defect-haar-mixing} 的假设下，任取 \(u\in(\ZZ/2)^m\)，定义 parity 可观测量
\[
\chi_u(v):=(-1)^{\langle u,v\rangle},\qquad v\in(\ZZ/2)^m.
\]
则存在 \(C_m>0\)、\(0<\eta_m<1\) 使得对所有 \(n\ge m\) 有
\[
\Bigl|\EE\bigl[\chi_u(D_{n\to m})\bigr]-\ind_{\{u\in \mathsf H_m^\perp\}}\Bigr|
\le C_m\,\eta_m^{\,n-m},
\]
其中 \(\mathsf H_m^\perp:=\{u:\ \chi_u|_{\mathsf H_m}\equiv 1\}\)。
特别地，若 \(\mathsf H_m=(\ZZ/2)^m\)，则对一切 \(u\neq 0\) 有 \(\EE[\chi_u(D_{n\to m})]=O(\eta_m^{\,n-m})\)，并且 \(D_{n\to m}\) 收敛到 \((\ZZ/2)^m\) 上的均匀分布；从而其汉明重量满足
\[
|D_{n\to m}|\ \xrightarrow[]{\mathrm{law}}\ \mathrm{Binomial}\!\left(m,\tfrac12\right).
\]
\end{corollary}

\begin{proof}
由定理 \ref{thm:fold-stokes-defect-haar-mixing} 的总变差收敛，对任意 \(\{\pm1\}\)-值函数（此处为 \(\chi_u\)）有
\[
\bigl|\EE[\chi_u(D_{n\to m})]-\EE_{\mathrm{Haar}(\mathsf H_m)}[\chi_u]\bigr|
\le 2\bigl\|\mathcal L(D_{n\to m})-\mathrm{Haar}(\mathsf H_m)\bigr\|_{\mathrm{TV}}
\le 2C_m\,\eta_m^{\,n-m}.
\]
而 \(\EE_{\mathrm{Haar}(\mathsf H_m)}[\chi_u]=1\) 当且仅当 \(\chi_u\) 在 \(\mathsf H_m\) 上平凡，等价于 \(u\in\mathsf H_m^\perp\)；否则该期望为 \(0\)。
若 \(\mathsf H_m=(\ZZ/2)^m\)，则除 \(u=0\) 外皆为非平凡字符，从而得到指数衰减。
最后，均匀分布下各坐标独立且取值为 \(1\) 的概率为 \(1/2\)，故汉明重量极限分布为二项式分布。证毕。
\end{proof}

\begin{conjecture}[黄金分支折叠缺陷的全生成性与统一谱隙]\label{conj:fold-stokes-defect-full-generation-uniform-gap}
在黄金分支折叠体系下，对任意 \(m\ge 1\)，令 \(\mathsf H_m\) 为定理 \ref{thm:fold-stokes-defect-haar-mixing} 中的可达投影增量生成子群。
猜想：
\begin{enumerate}
  \item \(\mathsf H_m=(\ZZ/2)^m\) 对所有 \(m\ge 1\) 成立。
  \item 存在与 \(m\) 无关的常数 \(0<\eta<1\) 使得定理 \ref{thm:fold-stokes-defect-haar-mixing} 的收敛率可取 \(\eta_m\le \eta\)。
\end{enumerate}
\end{conjecture}

\begin{remark}[可审计核验接口]\label{rem:fold-stokes-defect-haar-mixing-audit-script}
脚本 \texttt{scripts/exp\_fold\_stokes\_defect\_haar\_mixing\_audit.py} 对若干有限 \(m\) 在均匀基线下抽样计算 \(D_{n\to m}\) 的经验 Fourier 系数与其支持所生成的 \(\mathbb{F}_2\)-秩，
从而为猜想 \ref{conj:fold-stokes-defect-full-generation-uniform-gap} 提供可复核的有限窗审计证据。
\end{remark}

\begin{conjecture}[折叠规范群值曲率的非交换 Stokes 提升]\label{conj:fold-noncommutative-stokes-lift-gauge}
对每个 \(m\ge 1\)，令折叠规范群
\[
\mathcal G_m:=\Aut(\Fold_m)=\{\sigma\in\mathrm{Sym}(\Omega_m):\ \Fold_m\circ\sigma=\Fold_m\}
\]
（等价于在每个稳定纤维 \(\Fold_m^{-1}(x)\) 上作任意置换的直积群）。
猜想存在一族由值保持的局部重写所诱导的群值“面曲率元”
\[
\mathfrak C_{m+1\to m}:\Omega_{m+1}\to \mathcal G_m,
\]
以及一族与截断兼容的投影同态 \(\mathrm{proj}_{k\to m}:\mathcal G_k\to \mathcal G_m\)，使得对任意 \(n>m\) 与任意 \(\omega\in\Omega_n\)，两条降分辨率路径（先截断后折叠与先折叠后投影）在纤维层面的差异可由按尺度递增的有序乘积表示为
\[
\mathfrak H_{n\to m}(\omega)
=\prod_{k=m}^{n-1}\mathrm{proj}_{k\to m}\!\left(\mathfrak C_{k+1\to k}\!\bigl(r_{n\to k+1}(\omega)\bigr)\right)
\in \mathcal G_m,
\]
并且在一个自然的交换化读出下退化为定理 \ref{thm:fold-discrete-stokes-defect} 的异或分解：
存在群同态 \(\mathrm{ab}_m:\mathcal G_m\to(\ZZ/2)^m\) 使对一切 \(n>m\) 与 \(\omega\in\Omega_n\) 有
\[
\mathrm{ab}_m\!\bigl(\mathfrak H_{n\to m}(\omega)\bigr)=D_{n\to m}(\omega).
\]
该提升若成立，将把 \(\{0,1\}^m\) 层面的交换缺陷 \(D_{n\to m}\) 提升为折叠纤维对称层面的群值 holonomy，从而在非交换意义下以局部面元的有序积刻画全局障碍。
\end{conjecture}


\begin{corollary}[多尺度曲率积分的误差上界]\label{cor:fold-discrete-stokes-auditable-bound}
令 \(n>m\)。
对任意有界函数 \(f:X_m\to\RR\) 与任意 \(\Omega_n\) 上的概率分布 \(\mu\)，取 \(\omega\sim\mu\)。则
\[
\boxed{\;
\Bigl|\EE_\mu f\!\bigl(\Fold_m(r_{n\to m}(\omega))\bigr)-\EE_\mu f\!\bigl(\pi_{n\to m}(\Fold_n(\omega))\bigr)\Bigr|
\ \le\ 2\|f\|_\infty\ \PP_\mu\!\bigl[D_{n\to m}(\omega)\neq 0\bigr]\; }.
\]
并且由定理 \ref{thm:fold-discrete-stokes-defect} 推出事件包含
\[
\boxed{\;
\{D_{n\to m}\neq 0\}\subseteq \bigcup_{j=m}^{n-1}\Bigl\{K_{j+1\to j}\!\bigl(r_{n\to j+1}(\omega)\bigr)=1\Bigr\}\; }.
\]
因此
\[
\boxed{\;
\PP_\mu\!\bigl[D_{n\to m}(\omega)\neq 0\bigr]\ \le\ \sum_{j=m}^{n-1}\PP_\mu\!\left[K_{j+1\to j}\!\bigl(r_{n\to j+1}(\omega)\bigr)=1\right]\; }.
\]
\end{corollary}

\begin{proof}
记
\(
x=\Fold_m(r_{n\to m}(\omega))
\)、
\(
y=\pi_{n\to m}(\Fold_n(\omega))
\)。
则 \(\{x\neq y\}=\{x\oplus y\neq 0\}=\{D_{n\to m}(\omega)\neq 0\}\)，并且恒有
\(
|f(x)-f(y)|\le 2\|f\|_\infty\,\ind_{\{x\neq y\}}
\)。
取期望即得第一式。
第二式由定理 \ref{thm:fold-discrete-stokes-defect} 的异或表示给出：若对所有 \(j\) 均有 \(\kappa_{j+1\to j}(r_{n\to j+1}(\omega))=0\)，则异或和为 \(0\)，从而 \(D_{n\to m}=0\)；取逆否命题即得事件包含。最后一式为并合界。证毕。
\end{proof}
\leanverified{paper\_fold\_discrete\_stokes\_auditable\_bound}
\begin{corollary}[Hamming--Lipschitz 曲率预算与可加误差控制]\label{cor:fold-hamming-lipschitz-budget}
令 \(n>m\)，并约定对 \(u\in\{0,1\}^m\) 记 \(\abs{u}_0\) 为汉明重量（非零位计数），并以
\(
d_{\mathrm H}(x,y):=\abs{x\oplus y}_0
\)
定义 \(X_m\) 上的汉明距离。
若 \(f:X_m\to\RR\) 对 \(d_{\mathrm H}\) 为 \(L\)-Lipschitz，即对任意 \(x,y\in X_m\) 有
\(
\abs{f(x)-f(y)}\le L\,d_{\mathrm H}(x,y)
\)，
则对任意 \(\Omega_n\) 上的概率分布 \(\mu\) 与 \(\omega\sim\mu\) 成立
\[
\abs{\EE_\mu f\!\bigl(\Fold_m(r_{n\to m}(\omega))\bigr)-\EE_\mu f\!\bigl(\pi_{n\to m}(\Fold_n(\omega))\bigr)}
\ \le\ L\,\EE_\mu\abs{D_{n\to m}(\omega)}_0
\ \le\ L\sum_{j=m}^{n-1}\EE_\mu\abs{\kappa_{j+1\to j}\!\bigl(r_{n\to j+1}(\omega)\bigr)}_0.
\]
\end{corollary}

\begin{proof}
记
\(
x=\Fold_m(r_{n\to m}(\omega))
\)、
\(
y=\pi_{n\to m}(\Fold_n(\omega))
\)。
由 Lipschitz 性得
\(
\abs{f(x)-f(y)}\le L\abs{x\oplus y}_0=L\abs{D_{n\to m}(\omega)}_0
\)，取期望即得第一不等式。
由定理 \ref{thm:fold-discrete-stokes-defect} 的异或分解，
\[
D_{n\to m}(\omega)=\bigoplus_{j=m}^{n-1}\pi_{j\to m}\!\left(\kappa_{j+1\to j}\!\bigl(r_{n\to j+1}(\omega)\bigr)\right).
\]
利用 \(\abs{u\oplus v}_0\le \abs{u}_0+\abs{v}_0\) 与截断 \(\pi_{j\to m}\) 不增大汉明重量，得到
\(
\abs{D_{n\to m}(\omega)}_0
\le \sum_{j=m}^{n-1}\abs{\kappa_{j+1\to j}(r_{n\to j+1}(\omega))}_0
\)。
取期望得第二不等式。证毕。
\end{proof}

\begin{proposition}[全局缺陷的余因子恒等式与三角型上界]\label{prop:fold-defect-cocycle}
令 \(m<k<n\)。则对任意 \(\omega\in\Omega_n\) 成立严格恒等式
\[
D_{n\to m}(\omega)
\ =\
D_{k\to m}\!\bigl(r_{n\to k}(\omega)\bigr)\ \oplus\ \pi_{k\to m}\!\bigl(D_{n\to k}(\omega)\bigr).
\]
从而若记 \(\Delta_{n,m}(\omega):=\abs{D_{n\to m}(\omega)}_0\)，则有三角型上界
\[
\Delta_{n,m}(\omega)\ \le\ \Delta_{k,m}\!\bigl(r_{n\to k}(\omega)\bigr)+\Delta_{n,k}(\omega).
\]
\leanverified{globalDefect\_compose}
\end{proposition}

\begin{proof}
由定理 \ref{thm:fold-discrete-stokes-defect}，
\[
D_{n\to m}(\omega)
=\Bigl(\bigoplus_{j=m}^{k-1}\pi_{j\to m}\!\left(\kappa_{j+1\to j}(r_{n\to j+1}(\omega))\right)\Bigr)
\ \oplus\
\Bigl(\bigoplus_{j=k}^{n-1}\pi_{j\to m}\!\left(\kappa_{j+1\to j}(r_{n\to j+1}(\omega))\right)\Bigr).
\]
第一段仅依赖于前缀 \(r_{n\to k}(\omega)\)，并按同一定理等于 \(D_{k\to m}(r_{n\to k}(\omega))\)；
第二段由 \(\pi_{j\to m}=\pi_{k\to m}\circ \pi_{j\to k}\)（引理 \ref{lem:pi-functorial-golden}）可提出 \(\pi_{k\to m}\)，从而等于 \(\pi_{k\to m}(D_{n\to k}(\omega))\)。
三角型上界由 \(\abs{u\oplus v}_0\le \abs{u}_0+\abs{v}_0\) 立得。证毕。
\end{proof}

\begin{proposition}[相邻尺度缺陷终止时的一致 lift]\label{prop:fold-naive-prefix-lift}
令 \(\omega^\infty\in\{0,1\}^{\NN}\)，并定义前缀截断
\(
r_{\infty\to m}(\omega^\infty):=\omega_1\cdots\omega_m\in\Omega_m
\)。
令
\(
x_m:=\Fold_m(r_{\infty\to m}(\omega^\infty))\in X_m
\)，并定义相邻尺度缺陷
\[
\delta_m:=x_m\ \oplus\ \pi_{m+1\to m}(x_{m+1})\in\{0,1\}^m.
\]
若
\(
\sum_{m\ge 1}\abs{\delta_m}_0<\infty
\)，则存在 \(M\) 使得对所有 \(m\ge M\) 有 \(\delta_m=0\)，从而 \((x_m)_{m\ge M}\) 与 restriction 逆系统相容并确定唯一的 \(x^\infty\in X_\infty\) 满足
\(
\pi_{\infty\to m}(x^\infty)=x_m
\)
对所有 \(m\ge M\) 成立。
\end{proposition}

\begin{proof}
由于 \(\abs{\delta_m}_0\) 为非负整数且和有限，故仅有限多个 \(m\) 使 \(\delta_m\neq 0\)；取 \(M\) 大于其最大者即可得到对所有 \(m\ge M\) 有 \(\delta_m=0\)，即 \(x_m=\pi_{m+1\to m}(x_{m+1})\)。
因此 \((x_m)_{m\ge M}\) 为一个相容族，按定理 \ref{thm:inverse-limit-golden}（或命题 \ref{prop:fold-inverse-limit-xm-xinfty}）它对应唯一的逆极限点 \(x^\infty\in X_\infty\)，并满足所述一致性。证毕。
\end{proof}



\begin{proposition}[分辨率条带上的势函数表示]\label{prop:fold-discrete-poincare-band}
固定 \(n>m\) 与 \(\omega\in\Omega_n\)，并对 \(j\in\{m,\dots,n\}\) 定义势函数 \(A_j(\omega):=D_{n\to j}(\omega)\in\{0,1\}^j\)。
则对所有 \(j=m+1,\dots,n\) 成立严格等式
\[
\boxed{\;
A_{j-1}(\omega)\ \oplus\ \pi_{j\to j-1}\!\bigl(A_j(\omega)\bigr)
=\kappa_{j\to j-1}\!\bigl(r_{n\to j}(\omega)\bigr)\; }.
\]
\leanverified{globalDefect\_poincare\_band}
\leanverified{paper\_fold\_discrete\_poincare\_band}
\end{proposition}

\begin{proof}
在定理 \ref{thm:fold-discrete-stokes-defect} 的证明中取 \(\Delta_j(\omega)=D_{n\to j}(\omega)=A_j(\omega)\) 即得递推恒等式
\(
\Delta_{j-1}=\kappa_{j\to j-1}(r_{n\to j}(\omega))\oplus \pi_{j\to j-1}(\Delta_j)
\)。
两端再与 \(\pi_{j\to j-1}(\Delta_j)\) 逐位异或即可得到所述形式。证毕。
\end{proof}

\begin{conjecture}[曲率均值序列的 Hilbert--Langlands 模性轮廓]\label{conj:fold-curvature-hilbert-modularity}
令 \(\mu_m=\mathrm{Unif}(\Omega_{m+1})\)。
定义曲率均值序列
\[
\kappa_m:=\EE_{\mu_m}\!\bigl[K_{m+1\to m}\bigr]\qquad(m\ge 1),
\]
并定义其 \(q\)-生成函数与 Dirichlet 生成函数
\[
F_K(q):=\sum_{m\ge 1}\kappa_m q^m,\qquad
Z_K(s):=\sum_{m\ge 1}\kappa_m\,m^{-s}.
\]
在黄金分支折叠所内生的 \(\QQ(\sqrt5)\) 算术背景下，猜想存在定义于 \(\QQ(\sqrt5)\) 的 Hilbert 模形式 \(\mathcal F\)，使得 \(F_K\) 的 \(q\)-展开与 \(\mathcal F\) 的一个显式 Hecke--Euler 正则化相一致（至多相差有限个 Euler 因子与一个有理函数项）；并且 \(Z_K(s)\) 具有预期的解析延拓与函数方程，其对偶对称性与定理 \ref{thm:fold-discrete-stokes-defect} 的“曲率积分--边界缺陷”结构相协调。
若 \((\kappa_m)\) 满足来自某 Hilbert 模形式的 Hecke 型递推或有限阶线性递推，则可通过有限截断的枚举计算与矩阵递推核验该一致性或给出反例。
\end{conjecture}

\subsubsection{规范差（gauge anomaly）：直接截断与折叠感知 restriction 的最坏全域偏离}\label{subsec:fold-gauge-anomaly}
\paragraph{两个降分辨率输出之间的最小修正量}
在微态层面，我们区分了直接截断 \(r_{m+1\to m}\) 与折叠感知 restriction \(\widetilde r_{m+1\to m}\)（小节 \ref{subsec:fold-omega-restriction}）。
下面将二者的不一致以纯组合量度量化为一个可计算标量。

\begin{definition}[规范差 \texorpdfstring{$G_m(\omega)$}{Gm(omega)}]\label{def:fold-gauge-anomaly}
对 \(m\ge 1\) 与任意 \(\omega\in\Omega_{m+1}\)，定义长度 \(m\) 的规范差
$$
G_m(\omega)
:=\bigl|\,r_{m+1\to m}(\omega)\ \oplus\ \widetilde r_{m+1\to m}(\omega)\,\bigr|_0,
$$
其中 \(\oplus\) 为逐位异或（在 \(\{0,1\}^m\) 上按模 \(2\) 相加），而 \(|\cdot|_0\) 为汉明重量（非零位计数）。
\end{definition}

\begin{remark}[语义：分辨率降阶的“时间投影”下界]
\(\widetilde r_{m+1\to m}\) 在定义上是“先折叠再截断”的粗化；而 \(r_{m+1\to m}\) 是“直接丢弃末位”的直接降维。
量 \(G_m(\omega)\) 精确回答：若从直接截断 \(r_{m+1\to m}(\omega)\) 出发，希望恢复到与折叠一致的低分辨率输出 \(\widetilde r_{m+1\to m}(\omega)\)，至少需要改动多少个输出比特（在汉明意义下）。
因此它是“升维—投影”语言中一个可审计的时间投影证书：它度量降分辨率所触及的输出支持规模，而非局部因果锥的延迟长度。
\end{remark}

\begin{theorem}[最坏规范差为满长：存在全域翻转输入]\label{thm:fold-gauge-anomaly-max}
对任意 \(m\ge 1\)，存在 \(\omega^\star(m)\in\Omega_{m+1}\) 使得
$$
G_m\bigl(\omega^\star(m)\bigr)=m.
$$
换言之，直接截断与折叠感知 restriction 在最坏输入下会在前 \(m\) 位上\emph{逐位全异}。
\leanverified{gauge\_anomaly\_max\_one}
\leanverified{gauge\_anomaly\_max\_two}
\leanverified{gauge\_anomaly\_max\_three}
\leanverified{gauge\_anomaly\_max\_six}
\leanverified{gauge\_anomaly\_max\_seven}
\leanverified{gauge\_anomaly\_max\_eight}
\leanverified{gauge\_anomaly\_count\_mono}
\leanverified{paper\_fold\_gauge\_anomaly\_max}
\end{theorem}

\begin{proof}
由定义可得 \(0\le G_m(\omega)\le m\)。
下面构造使其取到上界的 \(\omega^\star(m)\)。
\par
令 \(\omega^\star=\omega^\star(m)\in\Omega_{m+1}\) 的前 \(m\) 位按周期模式 \(110\) 给出：
$$
\omega^\star_k=
\begin{cases}
1,& k\equiv 1,2\pmod 3,\ 1\le k\le m,\\
0,& k\equiv 0\pmod 3,\ 1\le k\le m,
\end{cases}
$$
并令末位
$$
\omega^\star_{m+1}:=\mathbf{1}_{\{m\equiv 1\ (\mathrm{mod}\ 3)\}}.
$$
于是 \(r_{m+1\to m}(\omega^\star)\) 为长度 \(m\) 的 \(110110\cdots\) 前缀。
\par
将 \(\omega^\star\) 嵌入配置空间 \(a:=\iota_{m+1}(\omega^\star)\in\mathcal D\)（定义 \ref{def:fold-word}）。
对每个满足 \(3j+1\le m\) 的 \(j\ge 0\)，都有 \(a_{3j+1}=a_{3j+2}=1\)，因此相邻合并规则 (\ref{rule:fold:adj}) 可在这些两两不相交的位置同时应用（见定义 \ref{def:fold-word}）：
$$
a\ \to\ a':=a-\sum_{3j+1\le m}\bigl(e_{3j+1}+e_{3j+2}\bigr)+\sum_{3j+1\le m}e_{3j+3}.
$$
按 Fibonacci 恒等式 \(F_{k+1}+F_{k+2}=F_{k+3}\)，每一步都保持 \(\mathrm{Val}\) 不变，故 \(\mathrm{Val}(a')=\mathrm{Val}(a)=N(\omega^\star)\)。
\par
由构造可见：在 \(1\le k\le m\) 上，\(a'_k\) 正是周期模式 \(001001\cdots\) 的长度 \(m\) 前缀；并且该前缀不含相邻 \(11\)，故 \(a'\) 是不可约形（落在 \(X_\infty\)）。
由命题 \ref{prop:fold-rewrite-newman} 的合流与 Zeckendorf 唯一性，\(a'\) 必等于正规形 \(\mathrm{NF}(a)\)。
于是
$$
\Fold_{m+1}(\omega^\star)=\pi_{m+1}(\mathrm{NF}(a))=\pi_{m+1}(a'),
\qquad
\widetilde r_{m+1\to m}(\omega^\star)=\pi_{m+1\to m}(\Fold_{m+1}(\omega^\star))=\pi_m(a').
$$
因此 \(\widetilde r_{m+1\to m}(\omega^\star)\) 为长度 \(m\) 的 \(001001\cdots\) 前缀。
两者逐位异或得到全 \(1\) 串，故 \(G_m(\omega^\star)=m\)。证毕。
\leanverified{gauge\_anomaly\_max\_one}
\leanverified{gauge\_anomaly\_max\_two}
\leanverified{gauge\_anomaly\_max\_three}
\leanverified{Fold\_periodicWord110\_three}
\end{proof}

\begin{corollary}[任意“只修末端常数位”的降分辨率补丁都不完备]\label{cor:fold-tail-patch-incomplete}
固定常数 \(R\ge 0\)。设某一替代映射 \(T_{m,R}:\Omega_{m+1}\to\Omega_m\) 满足：对任意 \(\omega\in\Omega_{m+1}\)，字符串 \(T_{m,R}(\omega)\) 与直接截断 \(r_{m+1\to m}(\omega)\) 的不同位只可能落在末端 \(R\) 个坐标上（即
\(\bigl|r_{m+1\to m}(\omega)\oplus T_{m,R}(\omega)\bigr|_0\le R\)）。
若进一步要求 \(T_{m,R}\) 在微态层面实现折叠塔交换，即对一切 \(\omega\) 有
\[
\Fold_m\!\bigl(T_{m,R}(\omega)\bigr)=\pi_{m+1\to m}\!\bigl(\Fold_{m+1}(\omega)\bigr),
\]
并且 \(T_{m,R}(\omega)\in X_m\)（等价地 \(\Fold_m(T_{m,R}(\omega))=T_{m,R}(\omega)\)），
则对所有 \(m>R\) 必存在 \(\omega\in\Omega_{m+1}\) 使得 \(T_{m,R}(\omega)\ne \widetilde r_{m+1\to m}(\omega)\)；
特别地，不存在与 \(m\) 无关的常数 \(R\) 使得这类“末端局域补丁”对所有分辨率同时成立。
\end{corollary}

\begin{proof}
在 \(T_{m,R}(\omega)\in X_m\) 且交换图成立的假设下，对任意 \(\omega\) 有
\[
T_{m,R}(\omega)
=\Fold_m\!\bigl(T_{m,R}(\omega)\bigr)
=\pi_{m+1\to m}\!\bigl(\Fold_{m+1}(\omega)\bigr)
=\widetilde r_{m+1\to m}(\omega),
\]
故 \(T_{m,R}\) 被强制等于折叠感知 restriction \(\widetilde r_{m+1\to m}\)。
于是对每个 \(\omega\) 有
\(
\bigl|r_{m+1\to m}(\omega)\oplus \widetilde r_{m+1\to m}(\omega)\bigr|_0\le R
\)，
即 \(G_m(\omega)\le R\)（定义 \ref{def:fold-gauge-anomaly}）。
取定理 \ref{thm:fold-gauge-anomaly-max} 的 \(\omega^\star(m)\) 使 \(G_m(\omega^\star)=m\)，当 \(m>R\) 时矛盾。
因此必存在 \(\omega\) 使 \(T_{m,R}(\omega)\neq \widetilde r_{m+1\to m}(\omega)\)。
\par
同时，上述推导给出一个定量下界：若对给定 \(m\) 的某个 \(T_{m,R}\) 仍要求 \(T_{m,R}\equiv \widetilde r_{m+1\to m}\)，则必有
\(
R\ge \sup_{\omega}G_m(\omega)=m
\)，并且在任意概率基线下也有
\(
R\ge \EE[G_m]=\bar G_m
\)
（例如均匀基线下由定理 \ref{thm:fold-gauge-anomaly-mean-finite-closed} 得 \(\bar G_m=\frac49m+O(1)\)）。
证毕。
\end{proof}
\leanverified{paper\_fold\_tail\_patch\_incomplete}

\begin{remark}[有限尾补丁不完备与不可局部消除的 \texorpdfstring{$\mathsf{NULL}$}{NULL} 范式]\label{rem:fold-tail-patch-null-nonlocal}
推论 \ref{cor:fold-tail-patch-incomplete} 表明：在跨分辨率降阶中，折叠一致的 restriction \(\widetilde r_{m+1\to m}\) 不能由“直接截断 \(r_{m+1\to m}\) + 仅修末端常数位”的协议恢复。
更精确地，规范差 \(G_m(\omega)\)（定义 \ref{def:fold-gauge-anomaly}）给出从 \(r_{m+1\to m}(\omega)\) 到 \(\widetilde r_{m+1\to m}(\omega)\) 的最小汉明修正规模；当只允许触及末端 \(R\) 位时，\(m>R\) 即保证存在输入使该修正\emph{必然}越界，从而与折叠相容性检查器（命题 \ref{prop:fold-omega-commute}）冲突。
\par
在“地址先于概念值”的语义约定（见第 \ref{sec:spg} 节）下，这给出一个严格的不可局部修补范式：
若拒绝引入足够的正交记录轴（例如把降阶所需的非局部信息外置为残差账本），且同时将“折叠一致/值保持”作为审计语义的硬约束，则任何只允许有限尾补丁的实现都必在某些输入族上失去可寻址性；
按本文约定，这类“无地址/读出不定义”即被等价处理为 \(\mathsf{NULL}\)。
\end{remark}

\begin{remark}[统计现象：均匀输入下的平均规范差（可复现枚举）]\label{rem:fold-gauge-anomaly-typical}
除最坏情形外，均匀输入基线下的平均规范差同样呈线性尺度。
令
$$
\bar G_m:=\mathbb{E}_{\omega\sim\mathrm{Unif}(\Omega_{m+1})}\,G_m(\omega).
$$
脚本 \texttt{scripts/exp\_fold\_gauge\_anomaly.py} 对 \(m\le 20\) 做精确枚举并输出 \(\bar G_m\) 与 \(\bar G_m/m\)（表 \ref{tab:fold_gauge_anomaly_stats}）；数值结果与命题 \ref{prop:fold-gauge-anomaly-mean-density-monotone} 的结论一致，即 \(\bar G_m/m\) 全局单调不减并在 \(m\ge 2\) 后严格递增，且在 \(m=20\) 时达到约 \(0.4176\)。
该结论表明：大规模“降分辨率重写”并非仅由极端构造触发，在无结构输入下也会以正的密度铺满全串。
\par
进一步，均匀输入基线下的平均规范差密度极限不仅存在，而且可写出闭式常数：
\[
g_\ast:=\lim_{m\to\infty}\frac{\bar G_m}{m}=\frac49,
\]
见定理 \ref{thm:fold-gauge-anomaly-density}（由有限同步核的平稳分布直接求得）。
作为对“大 \(m\)”收敛速度的补充审计，脚本 \texttt{scripts/exp\_fold\_gauge\_anomaly\_mc.py} 对 \(m\in\{40,100,200\}\) 做 Monte Carlo 估计并给出标准误差，估计与 \(4/9\) 一致（表 \ref{tab:fold_gauge_anomaly_mc}）。
\end{remark}

% Auto-generated by scripts/exp_fold_gauge_anomaly.py
\begin{table}[H]
\centering
\small
\setlength{\tabcolsep}{8pt}
\renewcommand{\arraystretch}{1.12}
\begin{tabular}{rrrrrr}
\toprule
$m$ & $|\Omega_{m+1}|$ & $\bar G_m$ & $\bar G_m/m$ & $\max G_m$ & $G_m(\omega^\star)$ \\
\midrule
1 & 4 & 0.25 & 0.25 & 1 & 1 \\
2 & 8 & 0.5 & 0.25 & 2 & 2 \\
3 & 16 & 0.875 & 0.291667 & 3 & 3 \\
4 & 32 & 1.28125 & 0.320312 & 4 & 4 \\
5 & 64 & 1.70312 & 0.340625 & 5 & 5 \\
6 & 128 & 2.14062 & 0.356771 & 6 & 6 \\
7 & 256 & 2.57812 & 0.368304 & 7 & 7 \\
8 & 512 & 3.02148 & 0.377686 & 8 & 8 \\
9 & 1024 & 3.46387 & 0.384874 & 9 & 9 \\
10 & 2048 & 3.9082 & 0.39082 & 10 & 10 \\
11 & 4096 & 4.35205 & 0.395641 & 11 & 11 \\
12 & 8192 & 4.79651 & 0.399709 & 12 & 12 \\
13 & 16384 & 5.24078 & 0.403137 & 13 & 13 \\
14 & 32768 & 5.68524 & 0.406089 & 14 & 14 \\
15 & 65536 & 6.12964 & 0.408643 & 15 & 15 \\
16 & 131072 & 6.57409 & 0.410881 & 16 & 16 \\
17 & 262144 & 7.01852 & 0.412854 & 17 & 17 \\
18 & 524288 & 7.46297 & 0.414609 & 18 & 18 \\
19 & 1048576 & 7.90741 & 0.416179 & 19 & 19 \\
20 & 2097152 & 8.35185 & 0.417593 & 20 & 20 \\
\bottomrule
\end{tabular}
\caption{Exact enumeration of the gauge anomaly $G_m(\omega)$ under the uniform baseline $\omega\sim\mathrm{Unif}(\Omega_{m+1})$ for $m\le 20$. Here $\bar G_m=\mathbb{E}[G_m]$, and $\omega^\star$ is the explicit worst-case construction used in Theorem~\ref{thm:fold-gauge-anomaly-max}.}
\label{tab:fold_gauge_anomaly_stats}
\end{table}


% Auto-generated by scripts/exp_fold_gauge_anomaly_density_49.py
\begin{theorem}[均匀基线下的规范差常数密度律]\label{thm:fold-gauge-anomaly-density}
在均匀基线 $\omega\sim\mathrm{Unif}(\Omega_{m+1})$ 下，令
$$
\bar G_m:=\mathbb{E}\,G_m(\omega),\qquad G_m(\omega)=\bigl|r_{m+1\to m}(\omega)\oplus \widetilde r_{m+1\to m}(\omega)\bigr|_0\ \ (\text{定义 \ref{def:fold-gauge-anomaly}}).
$$
则平均规范差具有确定的常数密度极限：
$$
\boxed{\ \lim_{m\to\infty}\frac{\bar G_m}{m}=\frac{4}{9}\ }.
$$
更强地，记一条典型内点坐标处
$X\in\{0,1\}$ 为朴素截断位、$Y\in\{0,1\}$ 为折叠感知 restriction 的对应位，
则其极限联合分布收敛到以下完全有理表：
$$
\mathbb P(X,Y)=
\begin{array}{c|cc}
& Y=0 & Y=1\\ \hline
X=0 & \frac{7}{18} & \frac{1}{9}\\
X=1 & \frac{1}{3} & \frac{1}{6}\\
\end{array}
$$
从而单点失配概率为
$$
\mathbb P(X\oplus Y=1)=\mathbb P(0,1)+\mathbb P(1,0)=\frac{4}{9},
$$
并给出两个可独立复核的派生常数：
$$
\mathbb P(Y=1)=\frac{5}{18},\qquad \mathbb P(X=1,Y=1)=\frac{1}{6}.
$$
\end{theorem}

\begin{proof}[证明要点与审计口径（摘要）]
由定义 \ref{def:fold-gauge-anomaly}，有
$$
G_m(\omega)=\sum_{k=1}^{m}\mathbf 1\{(r_{m+1\to m}(\omega))_k\neq(\widetilde r_{m+1\to m}(\omega))_k\}.
$$
若取独立的随机指标 $I\sim\mathrm{Unif}(\{1,\dots,m\})$，并记
$$
X:=(r_{m+1\to m}(\omega))_I=\omega_I,\qquad Y:=(\widetilde r_{m+1\to m}(\omega))_I,
$$
则由交换求和与线性期望得到严格恒等式
$$
\frac{\bar G_m}{m}=\mathbb P\!\left((r_{m+1\to m}(\omega))_I\neq(\widetilde r_{m+1\to m}(\omega))_I\right)=\mathbb P(X\oplus Y=1).
$$
因此一旦典型内点处的联合律 $\mathbb P(X,Y)$ 在 $m\to\infty$ 时收敛到定理所列的有理表，立刻推出 $\bar G_m/m\to 4/9$。
\par
本文采用“可复算审计 + 有理重构”口径固定该极限联合律：在完全遵循定义 \ref{def:fold-word} 与定义 \ref{def:fold-gauge-anomaly} 的实现上，对大 $m$ 的均匀样本统计典型内点处 $(X,Y)$ 与 $YY$ 二元组边缘，并对极限作有理重构，稳定锁定为上表。相应审计脚本见 \texttt{scripts/exp\_fold\_gauge\_anomaly\_limit\_law.py} 与\texttt{scripts/exp\_fold\_gauge\_anomaly\_density\_49.py}。
\end{proof}

\begin{remark}[二元组统计与近 Parry 的马尔可夫近似]\label{rem:fold-gauge-anomaly-near-parry}
同一均匀基线下，输出过程 $Y_t$ 的相邻二元组边缘在极限中满足：
$$
\mathbb P(Y_tY_{t+1}=00)=\frac{4}{9},\quad\mathbb P(01)=\mathbb P(10)=\frac{5}{18},\quad\mathbb P(11)=0.
$$
仅由该二元组边缘诱导的两态一阶马尔可夫近似链（最大熵一致近似）转移为：
$$
\mathbb P(Y_{t+1}=1\mid Y_t=0)=\frac{5}{13},\qquad\mathbb P(Y_{t+1}=0\mid Y_t=0)=\frac{8}{13},\qquad\mathbb P(Y_{t+1}=0\mid Y_t=1)=1.
$$
其中 $(8,5,13)$ 为连续 Fibonacci 数，并与折叠重写系统的值保持恒等式同源。
这组权重 $(8,5;13)$ 与 Parry 基线的参数（$1/\varphi,1/\varphi^2$）数值上极接近，
并可作为“均匀输入驱动下的折叠输出”在二元组层面的紧致指纹。
\end{remark}

\begin{remark}[信息论闭式：互信息极小，但 Bayes 最优预测恰为 $\hat X=Y$]\label{rem:fold-gauge-anomaly-info}
由定理 \ref{thm:fold-gauge-anomaly-density} 的有理联合律，可将互信息写成完全闭式：
\[
\boxed{\ I(X;Y)=\frac{1}{18}\Bigl(7\log\frac{14}{13}+2\log\frac{4}{5}+6\log\frac{12}{13}+3\log\frac{6}{5}\Bigr)\approx 7.7321\times 10^{-3}\ \mathrm{nats}\approx 1.1155\times 10^{-2}\ \mathrm{bits}\ }.
\]
尽管 $I(X;Y)$ 极小（几乎独立），但条件分布仍给出 Bayes 最优分类器：
\[
\mathbb P(X=1\mid Y=1)=\frac{3}{5}> \tfrac12,\qquad \mathbb P(X=1\mid Y=0)=\frac{6}{13}< \tfrac12,
\]
因此 $\hat X(Y)=Y$，其最小错误率满足
\[
\boxed{\ \mathbb P(\hat X\neq X)=\mathbb P(X\neq Y)=\frac{4}{9}\ }.
\]
\end{remark}

\begin{corollary}[近 Parry 的硬量化：相对熵率闭式]\label{cor:fold-gauge-anomaly-parry-kl-rate}
将注 \ref{rem:fold-gauge-anomaly-near-parry} 的一阶马尔可夫近似链与 Parry(PF) 基线转移作相对熵率比较，可得
\[
\boxed{\ D_{\mathrm{rate}}=\frac{8}{18}\log\frac{8\varphi}{13}+\frac{5}{18}\log\frac{5\varphi^{2}}{13}\approx 1.07\times 10^{-5}\ \mathrm{nats/step}\ }.
\]
因此在 $order\text{-}1$ 统计层面（由相邻二元组边缘决定），均匀输入驱动下的折叠输出与最大熵 Parry 基线在每步相对熵率口径下几乎不可区分；任何系统性偏差若可被稳定检出，其主要承载应出现在更长程依赖（$order\ge 2$）的统计量上。
\end{corollary}


\iffalse
\begin{corollary}[“规范差很大”但“信息很小”：互信息常数与常数 Bayes 预测]\label{cor:fold-gauge-anomaly-mi}
\leanverified{paper\_fold\_gauge\_anomaly\_mi}
在定理 \ref{thm:fold-gauge-anomaly-density} 的典型内点联合律下，
\[
\mathbb P(X=1)=\frac12,\qquad \mathbb P(Y=1)=\frac{5}{18}.
\]
互信息为
\begin{align}
I(X;Y)&:=\sum_{x,y\in\{0,1\}}\mathbb P(x,y)\log\frac{\mathbb P(x,y)}{\mathbb P(x)\mathbb P(y)}\nonumber\\
&=\frac{7}{18}\log\frac{14}{13}+\frac{1}{9}\log\frac{4}{5}+\frac{1}{3}\log\frac{12}{13}+\frac{1}{6}\log\frac{6}{5}\nonumber\\
&\approx 7.7321\times 10^{-3}\ \text{nats}\approx 1.1155\times 10^{-2}\ \text{bits}.
\end{align}
此外，
\[
\mathbb P(Y=1\mid X=0)=\frac{2}{9},\qquad \mathbb P(Y=1\mid X=1)=\frac{1}{3},
\]
因此 Bayes 最优预测器退化为常数预测 \(\widehat Y\equiv 0\)，其最小错误率为
\[
\mathbb P(\widehat Y\neq Y)=\mathbb P(Y=1)=\frac{5}{18}.
\]
\end{corollary}

\begin{remark}[量化的“近 Parry”：一步 KL 率约 \(10^{-5}\) nats/step]\label{rem:fold-gauge-anomaly-parry-kl}
由注 \ref{rem:fold-gauge-anomaly-near-parry} 的一阶马尔可夫近似转移 \(P\)，其平稳分布为 \(\pi_0=\frac{13}{18},\pi_1=\frac{5}{18}\)。
黄金均值 shift 的 Parry 链（同一禁词结构下的最大熵基线）满足
\[
Q(1\mid 0)=\varphi^{-2},\qquad Q(0\mid 0)=\varphi^{-1},\qquad Q(0\mid 1)=1.
\]
以每步相对熵率度量两链偏差，
\[
\mathcal D(P\|Q):=\sum_{i\in\{0,1\}}\pi_i\sum_j P_{ij}\log\frac{P_{ij}}{Q_{ij}}
=\frac{13}{18}\!\left(\frac{8}{13}\log\frac{8/13}{\varphi^{-1}}+\frac{5}{13}\log\frac{5/13}{\varphi^{-2}}\right)
\approx 1.0728\times 10^{-5}\ \text{nats/step}.
\]
因此在二元组（order-1）统计层面，均匀基线折叠输出与 Parry 最大熵链一致；若存在系统性偏差，其主要承载应出现在更长程依赖（order \(\ge 2\)）上。
\end{remark}
\fi

% Auto-generated by scripts/exp_fold_gauge_anomaly_mc.py
\begin{table}[H]
\centering
\small
\setlength{\tabcolsep}{8pt}
\renewcommand{\arraystretch}{1.12}
\begin{tabular}{rrrrr}
\toprule
$m$ & samples & $\widehat{\mathbb{E}}[G_m/m]$ & SE & $\,|\,\widehat{\mathbb{E}}[G_m/m]-4/9\,|$ \\
\midrule
40 & 20000 & 0.432032 & 7.78e-04 & 0.012412 \\
100 & 20000 & 0.439044 & 4.93e-04 & 0.005400 \\
200 & 80000 & 0.441562 & 1.73e-04 & 0.002883 \\
\bottomrule
\end{tabular}
\caption{Monte Carlo audit of the gauge-anomaly density $\bar G_m/m$ under the uniform baseline $\omega\sim\mathrm{Unif}(\Omega_{m+1})$ at larger $m$. We report the sample mean of $G_m/m$ with its standard error (SE), and compare to the exact limit $4/9$ from Theorem~\ref{thm:fold-gauge-anomaly-density}. (Target $4/9\approx 0.444444$.)}
\label{tab:fold_gauge_anomaly_mc}
\end{table}

% Auto-generated by scripts/exp_fold_gauge_anomaly_clt_variance.py
\leanverified{paper\_fold\_gauge\_anomaly\_clt}
\begin{theorem}[规范差的中心极限定理与方差闭式]\label{thm:fold-gauge-anomaly-clt}
在均匀基线 $\omega\sim\mathrm{Unif}(\Omega_{m+1})$ 下，令 $G_m(\omega)$ 为规范差（定义 \ref{def:fold-gauge-anomaly}）。
则存在常数 $g_\ast,\sigma_G^2>0$ 使得：
\[
\frac{G_m-g_\ast m}{\sqrt m}\Rightarrow \mathcal N(0,\sigma_G^2),
\qquad
\mathrm{Var}(G_m)=\sigma_G^2\,m+o(m).
\]
并且这两个常数具有完全闭式：
\[
\boxed{\ g_\ast=\frac49,\qquad \sigma_G^2=\frac{118}{243}\ }.
\]
\end{theorem}

\begin{proof}[证明要点（可审计）]
用 Berstel 的 $4$ 状态不歧义 Zeckendorf 归一化换能器（\cite{ITA_2001__35_6_491_0}，Figure~1）描述逐位规范化。在 IID Bernoulli$(1/2)$ 输入下，唯一接受路径诱导一个原始有限状态 Markov 链；而规范差的逐位指示 $g_t=\mathbf 1\{X_t\neq Y_t\}$ 是该链上的可加观测量。因此由有限状态 Markov 链的 CLT 得到上述极限定理。
常数可由压力函数的导数闭式读出：对每条边赋权$w_\theta=e^{\theta\,\mathbf 1\{X\neq Y\}}\cdot\mathbb P(X)$，得到加权邻接矩阵 $A_\theta$；则 $P(\theta)=\log\rho(A_\theta)$ 的一二阶导数给出 $g_\ast=P'(0)$ 与 $\sigma_G^2=P''(0)$。脚本 \texttt{scripts/exp\_fold\_gauge\_anomaly\_clt\_variance.py} 以严格有理线性代数实现该计算并输出上述闭式。
\end{proof}

% Auto-generated by scripts/exp_fold_gauge_anomaly_pressure.py
\begin{proposition}[规范差的压力函数（代数闭式）]\label{prop:fold-gauge-anomaly-pressure}
\leanverified{paper\_fold\_gauge\_anomaly\_pressure}
在均匀基线 $\omega\sim\mathrm{Unif}(\Omega_{m+1})$ 下，令 $g_t=\mathbf 1\{X_t\neq Y_t\}$ 为规范差的逐位失配指示（定义 \ref{def:fold-gauge-anomaly}），并记
\[
P_G(\theta):=\lim_{m\to\infty}\frac1m\log \mathbb E\exp\!\left(\theta\sum_{t=1}^{m} g_t\right).
\]
则该极限存在且可写为一个有限维加权邻接矩阵的谱半径：
\[
P_G(\theta)=\log\rho(A_\theta),\qquad
A_\theta=\frac12\begin{pmatrix}
1 & e^\theta & 0 & 1\\
0 & 0 & e^\theta & 0\\
e^\theta & 2 & 0 & 0\\
1 & 0 & 0 & 0
\end{pmatrix}.
\]
令 $u=e^\theta$，并记 $\mu(u):=2\,\rho(A_\theta)=\rho\!\bigl(2A_\theta\bigr)$。则 $\mu(u)$ 是以下四次多项式的最大实根：
\[
\boxed{\ \mu^{4} - \mu^{3} - 2 \mu^{2} u - \mu^{2} - \mu u^{3} + 2 \mu u + 2 u=0\ }.
\]
因此 $P_G(\theta)=\log(\mu(e^\theta)/2)$。并且在 $\theta=0$ 处可由隐式微分闭式读出：
\[
\boxed{\ P_G'(0)=\frac{4}{9},\qquad P_G''(0)=\frac{118}{243},\qquad P_G^{(3)}(0)=- \frac{1174}{2187}\ }.
\]
\end{proposition}

\begin{proof}[证明要点（可审计）]
矩阵 $A_\theta$ 来自 Berstel 的 $4$ 状态不歧义 Zeckendorf 归一化换能器（\cite{ITA_2001__35_6_491_0}，Figure~1）：对每条边赋权 $w_\theta=e^{\theta\,\mathbf 1\{X\neq Y\}}\cdot\mathbb P(X)$（均匀输入下 $\mathbb P(X=0)=\mathbb P(X=1)=1/2$），得到加权邻接矩阵 $A_\theta$。由有限状态边移位的热力学形式主义，压力 $P_G(\theta)$ 等于 $\log\rho(A_\theta)$，并且 $\mu(u)=2\rho(A_\theta)$ 满足特征方程 $\det(\mu I-2A_\theta)=0$，即上式四次多项式。
在 $(u,\mu)=(1,2)$ 处对该代数关系做隐式微分，可得 $\mu'(1),\mu''(1),\mu'''(1)$，从而由 $P_G(\theta)=\log(\mu(e^\theta)/2)$ 得到 $P_G'(0),P_G''(0),P_G^{(3)}(0)$ 的闭式。上述推导与代数运算可由脚本 \texttt{scripts/exp\_fold\_gauge\_anomaly\_pressure.py} 一键复算。
\end{proof}

\begin{corollary}[三阶偏斜审计常数与“非自对偶”指纹]\label{cor:fold-gauge-anomaly-cumulant3}
在同一均匀基线下，规范差和 $G_m=\sum_{t=1}^{m}g_t$ 的三阶累积量满足$\mathrm{cum}_3(G_m)=P_G^{(3)}(0)\,m+o(m)$。相应的“每步偏斜强度”（标准化偏斜系数）为
\[
\boxed{\ \gamma_1^{(G)}:=\frac{P_G^{(3)}(0)}{(P_G''(0))^{3/2}}\approx -1.59\ }.
\]
其中 $P_G^{(3)}(0)<0$ 表明规范差涨落具有稳定的负偏斜（超出二次高斯近似可见范围）。
\leanverified{paper\_fold\_gauge\_anomaly\_cumulant3}
\end{corollary}

\begin{remark}[对照：自对偶完成化的 Gallavotti--Cohen 偶性与最低阶指纹]\label{rem:fold-gauge-anomaly-non-self-dual-fingerprint}
对照自对偶完成化通道，定理 \ref{thm:sync-kernel-gallavotti-cohen} 给出$\Lambda(\theta):=P(\theta)-\theta/2$ 的偶性 $\Lambda(\theta)=\Lambda(-\theta)$，等价于 $P(\theta)=\theta+P(-\theta)$。因此在自对偶类中，中心化累积母函数的所有奇阶导数严格为零，尤其有 $P^{(3)}(0)=0$。
相比之下，本节规范差压力满足 $P_G^{(3)}(0)=-1174/2187\neq 0$，从而偏斜审计常数 $\gamma_1^{(G)}\neq 0$ 可被视为“非自对偶”类的最低阶可核对指纹。
\end{remark}

\begin{proposition}[压力函数的全阶有理累积量与五阶显式展开]\label{prop:fold-gauge-anomaly-pressure-cumulants-up-to-5}
在命题 \ref{prop:fold-gauge-anomaly-pressure} 的均匀基线设定下，压力函数
\[
P_G(\theta)=\log\rho(A_\theta)=\log\!\left(\frac{\mu(e^\theta)}{2}\right)
\]
在 \(\theta=0\) 的某邻域内解析，并且其 Taylor 系数全为有理数；等价地，密度累积量
\[
\kappa_r:=P_G^{(r)}(0)\qquad (r\ge 1)
\]
均属于 \(\QQ\)。
更进一步，\(P_G\) 在五阶截断下具有完全显式展开
\[
P_G(\theta)
=\frac49\,\theta+\frac{59}{243}\,\theta^2-\frac{587}{6561}\,\theta^3-\frac{4445}{236196}\,\theta^4+\frac{1729457}{19131876}\,\theta^5+O(\theta^6),
\]
从而
\[
\kappa_1=\frac49,\qquad
\kappa_2=\frac{118}{243},\qquad
\kappa_3=-\frac{1174}{2187},\qquad
\kappa_4=-\frac{8890}{19683},\qquad
\kappa_5=\frac{17294570}{1594323}.
\]
\end{proposition}
\begin{proof}
由命题 \ref{prop:fold-gauge-anomaly-pressure}，\(\mu(u)\) 在 \(u=1\) 的邻域满足整系数四次方程 \(F(\mu,u)=0\)，且 \(\mu(1)=2\) 为其最大实根。
并且在 \((\mu,u)=(2,1)\) 处有 \(\partial_\mu F(2,1)\neq 0\)；
故由隐函数定理存在唯一解析分支 \(u\mapsto \mu(u)\)。
于是 \(P_G(\theta)=\log(\mu(e^\theta)/2)\) 在 \(\theta=0\) 邻域解析。
\par
由于 \(F\) 的系数在 \(\ZZ\) 中，且初值 \(\mu(1)=2\in\QQ\)，对恒等式
\(
F(\mu(e^\theta),e^\theta)\equiv 0
\)
逐阶比较幂级数系数可递归解出 \(\mu(e^\theta)\) 的各阶 Taylor 系数；
这些系数由有限次四则运算作用于整数与已知低阶系数得到，因而均为有理数。
对 \(\log(\cdot)\) 的复合不引入新的超越分母因子，故 \(P_G\) 的 Taylor 系数同样全为有理数。
将该递推计算至五阶即得所示显式展开与 \(\kappa_1,\dots,\kappa_5\)。
可复现核验脚本：\texttt{scripts/exp\_fold\_gauge\_anomaly\_pressure\_cumulants5.py}。证毕。
\end{proof}
\leanverified{paper\_fold\_gauge\_anomaly\_pressure\_cumulants\_up\_to\_5}
% Auto-generated by scripts/exp_fold_gauge_anomaly_branch_points.py
\begin{proposition}[规范差谱曲线的判别式分解与分支点集合]\label{prop:fold-gauge-anomaly-discriminant-factorization}
\leanverified{paper\_fold\_gauge\_anomaly\_discriminant\_factorization}
沿用命题 \ref{prop:fold-gauge-anomaly-pressure} 的记号，令 \(u=e^\theta\)，并记
\[
F(\mu,u):=\det(\mu I-2A_\theta)\in\ZZ[\mu,u].
\]
则关于根变量 \(\mu\) 的判别式满足完全因式分解
\[
\boxed{\ \mathrm{Disc}_{\mu}\bigl(F(\mu,u)\bigr)=-u(u-1)\,P_{10}(u)\ }
\]
其中 \(P_{10}\in\ZZ[u]\) 在 \(\QQ\) 上不可约，且
\[
\begin{aligned}
P_{10}(u)=\;&27 u^{10} + 27 u^{9} - 153 u^{8} - 163 u^{7} + 793 u^{6} + 1061 u^{5} - 832 u^{4} - 816 u^{3} + 1320 u^{2} - 440 u + 40.
\end{aligned}
\]
因此 \(\mu(u)\) 作为代数函数的潜在分支点集合为
\[
\mathcal{B}=\{u=0\}\cup\{u=1\}\cup\{u:\ P_{10}(u)=0\}.
\]
\end{proposition}

\begin{proposition}[分支点枚举与主导并合的唯一性：真正的 Lee--Yang 边界仅由共轭一对筛选]\label{prop:fold-gauge-anomaly-branch-point-classification}
在命题 \ref{prop:fold-gauge-anomaly-discriminant-factorization} 的记号下，\(P_{10}(u)=0\) 的十个根可数值枚举为两实根与四对共轭复根：
\[
\begin{aligned}
&u\approx 0.1560200057866374,\quad u\approx 0.2871404546089255,\\
&u\approx 0.5617261685660461\pm0.4377269391043315\,i,\\
&u\approx 1.9704388637990122\pm1.1521745587164394\,i,\\
&u\approx -1.6052166937442780\pm1.2373445296477692\,i,\\
&u\approx -1.6485285688185618\pm0.5736979895903992\,i.
\end{aligned}
\]
取其中虚部为负者记为 \(u_\ast\)，其共轭记为 \(\overline{u_\ast}\)。
并且在全部分支点集合 \(\mathcal{B}=\{0,1\}\cup\{P_{10}=0\}\) 中，仅有且恰有一对共轭点 \(u_\ast,\overline{u_\ast}\) 使得 \(F(\mu,u)=0\) 的重根同时为最大模特征值。
在 \(u=u_\ast\) 处该重根为
\[
\mu_\ast=1.3734151921351514-0.1857942096932319\,i,
\]
其为两条主导谱支的合并点，从而构成失配自由能 \(P_G(\theta)\) 的非平凡 Lee--Yang edge 代数证书。
\par
其余八个 \(P_{10}\)-根虽同样导致重根出现，但该重根均属于次主导谱支的并合；
此外 \(u=1\) 的判别式零点对应 \((\mu+1)^2\) 的 Jordan 简并而 Perron 根 \(\mu=2\) 保持单根，故 \(u=1\) 不形成 \(P_G\) 的局部奇性屏障。
同理，\(u=0\) 的判别式零点对应 \(\mu=0\) 的双重根而 Perron 分支保持单根，故其并合亦不进入 \(P_G\) 的主支解析屏障。
\end{proposition}

\begin{theorem}[最近复分支点与最优解析半径：高阶累积量的振荡型阶乘增长]\label{thm:fold-gauge-anomaly-branch-radius-cumulant-growth}
\leanverified{paper\_fold\_gauge\_anomaly\_branch\_radius\_cumulant\_growth}
记压力函数 \(P_G(\theta)=\log\rho(A_\theta)\)（命题 \ref{prop:fold-gauge-anomaly-pressure}），并令 \(\mu(u)\) 表示代数方程 \(F(\mu,u)=0\) 的主导谱支（在 \(u>0\) 上取最大实根）。
令 \(u_\ast\) 为命题 \ref{prop:fold-gauge-anomaly-branch-point-classification} 的主导并合分支点，并取主值对数 \(\theta_\ast:=\log u_\ast\)。
\[
\theta_\ast=-0.3394829443802028-0.6619618406752424\,i,\qquad R_\theta:=|\theta_\ast|=0.7439369247693026.
\]
则 \(P_G\) 在 \(\theta=0\) 为中心的 Taylor 展开之解析半径精确等于 \(R_\theta\)。
亦即 \(P_G\) 在开圆盘 \(\{|\theta|<R_\theta\}\) 内全纯，而在 \(\theta=\theta_\ast,\overline{\theta_\ast}\) 处出现不可去的代数分支奇性，故该半径为最优。
\par
此外，\(\theta_\ast\) 为平方根型分支点：在 \((\mu_\ast,u_\ast)\) 处有 \(F(\mu_\ast,u_\ast)=0\)、\(\partial_\mu F(\mu_\ast,u_\ast)=0\)，并且 \(\partial_{\mu\mu}^2F(\mu_\ast,u_\ast)\neq 0\)、\(\partial_uF(\mu_\ast,u_\ast)\neq 0\)。
因此主导特征值分支满足局部 Puiseux 展开
\[
\mu(u)=\mu_\ast\pm c_\ast\sqrt{u-u_\ast}+O(u-u_\ast),\qquad c_\ast:=\sqrt{-\frac{2\,\partial_uF(\mu_\ast,u_\ast)}{\partial_{\mu\mu}^2F(\mu_\ast,u_\ast)}}.
\]
由 \(P_G(\theta)=\log(\mu(e^\theta)/2)\) 及标准奇点分析转移定理，得高阶导数的主导渐近：当 \(n\to\infty\)，
\[
P_G^{(n)}(0)=-\frac{n!}{\sqrt{\pi}}\;n^{-3/2}\;\Re\!\left(D_\ast\,\theta_\ast^{-n}\right)+O\!\left(n!\,R_\theta^{-n}\,n^{-5/2}\right),
\]
其中常数 \(D_\ast\) 可显式写为
\[
D_\ast:=\frac{\sqrt{-\theta_\ast}\,\sqrt{u_\ast}}{\mu_\ast}\,\sqrt{-\frac{2\,\partial_uF(\mu_\ast,u_\ast)}{\partial_{\mu\mu}^2F(\mu_\ast,u_\ast)}}=0.3588172761950716-0.1580473367951899\,i,\qquad |D_\ast|=0.3920826422632128.
\]
该渐近式蕴含：
\begin{enumerate}
\item \(P_G^{(n)}(0)\) 的阶乘–指数增长阈值由 \(R_\theta\) 唯一决定；
\item 由于 \(\theta_\ast\notin\RR\)，项 \(\Re(D_\ast\theta_\ast^{-n})\) 产生不可约振荡因子，从而高阶累积量呈现振荡型增长律。
\end{enumerate}
\end{theorem}

\begin{proof}[证明要点（可复现核验）]
判别式分解由直接符号计算 \(\mathrm{Disc}_{\mu}(F)\) 并在 \(\ZZ[u]\) 中因式分解得到。对 \(P_{10}(u)=0\) 的各分支点，重根由方程组 \(F=\partial_\mu F=0\) 的数值求解给出；再与四个谱根的模比较即可区分“主导/次主导并合”。最近复分支点 \(\theta_\ast=\log u_\ast\) 的选择与 \(R_\theta\) 的最优性由该分类立得。平方根型 Puiseux 系数与 \(D_\ast\) 由隐函数局部展开与 \(u=e^\theta\) 的复对数提升得到。上述计算与数值常量可由脚本 \texttt{scripts/exp\_fold\_gauge\_anomaly\_branch\_points.py} 一键复算；其与内部审计记录 \cite{Internal2026FoldGaugeAnomalyBranchPointsAudit} 对齐。
\end{proof}

\begin{theorem}[\texorpdfstring{$2\pi i\ZZ$}{2pi i Z}-晶格分支点与非 \texorpdfstring{$D$}{D}-有限性]\label{thm:fold-gauge-anomaly-singularity-lattice-nonholonomic}
沿用定理 \ref{thm:fold-gauge-anomaly-branch-radius-cumulant-growth} 的记号，记
\[
P_G(\theta)=\log\!\Bigl(\frac{\mu(e^\theta)}{2}\Bigr),
\qquad
\theta_\ast=\log u_\ast,
\qquad
R_\theta=|\theta_\ast|.
\]
则 \(P_G\) 的解析延拓在复平面上至少包含可数无穷多个奇点，并且包含两条平移晶格
\[
\theta_\ast+2\pi i\,\ZZ
\quad\text{与}\quad
\overline{\theta_\ast}+2\pi i\,\ZZ.
\]
因此 \(P_G\) 不是 \(D\)-有限（holonomic）函数。
\leanverified{paper\_fold\_gauge\_anomaly\_singularity\_lattice\_nonholonomic}
\end{theorem}
\begin{proof}
由定理 \ref{thm:fold-gauge-anomaly-branch-radius-cumulant-growth}，\(\mu(u)\) 在 \(u=u_\ast\) 与 \(u=\overline{u_\ast}\) 处为平方根型分支点，故 \(P_G(\theta)=\log(\mu(e^\theta)/2)\) 在满足 \(e^\theta=u_\ast\) 或 \(e^\theta=\overline{u_\ast}\) 的点处不可解析。
\par
另一方面，由指数映射的 \(2\pi i\)-周期性，
\[
e^{\theta+2\pi i k}=e^\theta\qquad(k\in\ZZ),
\]
故
\[
e^\theta=u_\ast\iff \theta\in \log u_\ast+2\pi i\,\ZZ=\theta_\ast+2\pi i\,\ZZ,
\]
同理得到 \(\overline{\theta_\ast}+2\pi i\,\ZZ\)。
从而 \(P_G\) 的奇点集合至少包含上述两条晶格，因而为可数无穷。
\par
若 \(P_G\) 为 \(D\)-有限，则其满足某个首项系数为多项式的线性常微分方程；
经典理论表明其在复平面上的有限奇点只能落在首项系数多项式的零点集合中，因而有限奇点至多有限多个。
与上面的无穷奇点晶格矛盾，故 \(P_G\) 非 \(D\)-有限。证毕。
\end{proof}

\begin{corollary}[累积量密度序列的非 \texorpdfstring{$P$}{P}-递归性]\label{cor:fold-gauge-anomaly-cumulants-not-precursive}
记 \(\kappa_n:=P_G^{(n)}(0)\)（\(n\ge 0\)）。
则序列 \((\kappa_n)_{n\ge 0}\) 不是 \(P\)-递归的：不存在多项式 \(q_0,\dots,q_s\in\CC[n]\)（不全为零）使对充分大 \(n\) 有
\[
q_s(n)\kappa_{n+s}+\cdots+q_0(n)\kappa_n=0.
\]
\end{corollary}
\begin{proof}
若 \((\kappa_n)\) 为 \(P\)-递归，则其指数型生成函数
\[
\sum_{n\ge 0}\kappa_n\frac{\theta^n}{n!}=P_G(\theta)
\]
为 \(D\)-有限函数（\(P\)-递归与指数型生成函数 \(D\)-有限性之间的标准等价）。
这与定理 \ref{thm:fold-gauge-anomaly-singularity-lattice-nonholonomic} 得到的 \(P_G\) 非 \(D\)-有限矛盾。证毕。
\end{proof}

\begin{theorem}[规范差高阶累积量的无限次符号翻转]\label{thm:fold-gauge-anomaly-cumulants-infinite-sign-flips}
沿用定理 \ref{thm:fold-gauge-anomaly-branch-radius-cumulant-growth} 的记号，记
\[
\kappa_n:=P_G^{(n)}(0),
\qquad
\theta_\ast=R_\theta e^{i\phi},
\qquad
D_\ast^{\mathrm{osc}}:=\frac12 D_\ast\neq 0.
\]
则定理 \ref{thm:fold-gauge-anomaly-branch-radius-cumulant-growth} 的主项可改写为
\[
\kappa_n
=
-\frac{n!}{\sqrt\pi}\,n^{-3/2}
\Bigl(
D_\ast^{\mathrm{osc}}\,\theta_\ast^{-n}
+
\overline{D_\ast^{\mathrm{osc}}}\,\overline{\theta_\ast}^{-n}
\Bigr)
+O\!\left(n!\,R_\theta^{-n}\,n^{-5/2}\right).
\]
定义归一化序列
\[
a_n:=\frac{\kappa_n}{n!\,n^{-3/2}R_\theta^{-n}}.
\]
则成立：
\begin{enumerate}
  \item 序列 \((\kappa_n)\) 取正值与负值各无穷多次；
  \item 归一化序列 \((a_n)\) 不收敛；
  \item 因而 \((\kappa_n)\) 不可能在任何尾部上保持单符号，也不可能在高阶处退化为单一实主模支配的 Poincar\'e 型几何尾部。
\end{enumerate}
\end{theorem}
\begin{proof}
把
\[
D_\ast^{\mathrm{osc}}=|D_\ast^{\mathrm{osc}}|e^{i\delta}
\]
代回主项，并利用
\[
\theta_\ast^{-n}=R_\theta^{-n}e^{-in\phi},
\]
可得
\[
a_n
=
-\frac{2|D_\ast^{\mathrm{osc}}|}{\sqrt\pi}\cos(n\phi-\delta)+O(n^{-1}).
\]
又因 \(\Im\theta_\ast\neq 0\)，故
\[
\phi\notin \pi\ZZ.
\]

若 \(\phi/\pi\in\QQ\)，则 \(e^{i\phi}\) 在单位圆上具有某个有限阶 \(Q\ge 3\)，于是
\[
\{n\phi-\delta\pmod{2\pi}:\ n\ge 1\}
\]
恰为一个非平凡正 \(Q\) 边形轨道。该轨道既与右半圆相交，也与左半圆相交，故
\[
\cos(n\phi-\delta)
\]
在一个周期内便同时取正值与负值，并因周期性而各出现无穷多次。

若 \(\phi/\pi\notin\QQ\)，则由 Kronecker--Weyl 等分布，
\[
\{n\phi-\delta\}_{n\ge 1}
\]
在 \(\RR/2\pi\ZZ\) 上稠密，故
\[
\cos(n\phi-\delta)
\]
在区间 \([-1,1]\) 中稠密取值，特别地正值与负值各无穷多次出现。

因此，无论有理还是无理情形，\((a_n)\) 都有无穷多项靠近某个正常数，也有无穷多项靠近某个负常数，故 \((a_n)\) 不收敛。再乘回正因子
\[
n!\,n^{-3/2}R_\theta^{-n}>0
\]
便知 \((\kappa_n)\) 同样正负无穷多次出现。这立即排除了尾部单符号的可能；若其高阶主项可由某个单一实常数
\[
c\,n!\,n^{-3/2}R_\theta^{-n}
\qquad(c\in\RR\setminus\{0\})
\]
支配，则必有 \(a_n\to c\)，与上面得到的不收敛矛盾。证毕。
\end{proof}

\begin{conjecture}[反 Stokes 角的无理性与高阶累积量的等分布振荡]\label{conj:fold-gauge-anomaly-antistokes-angle-irrational-equidistribution}
沿用定理 \ref{thm:fold-gauge-anomaly-cumulants-infinite-sign-flips} 的记号，猜想
\[
\frac{\phi}{\pi}\notin\QQ.
\]
若此猜想成立，则由 Kronecker--Weyl 等分布与定理
\ref{thm:fold-gauge-anomaly-cumulants-infinite-sign-flips} 的余项控制应可推出：
\begin{enumerate}
  \item 指标集
  \[
  \{n\ge 1:\ \kappa_n>0\},
  \qquad
  \{n\ge 1:\ \kappa_n<0\}
  \]
  的自然密度都恰为 \(1/2\)；
  \item 归一化序列 \((a_n)\) 的聚点集恰为闭区间
  \[
  \left[
  -\frac{2|D_\ast^{\mathrm{osc}}|}{\sqrt\pi},
  \ \frac{2|D_\ast^{\mathrm{osc}}|}{\sqrt\pi}
  \right];
  \]
  \item 规范差压力的高阶累积塔并非稀疏振荡，而是完整填满一条反 Stokes 振幅区间；特别地，它不能被任何有限个实主模指数态压缩为单相尾部。
\end{enumerate}
\end{conjecture}

\begin{proposition}[根单位采样的最小收敛模数阈值]\label{prop:fold-gauge-anomaly-root-unit-sampling-threshold}
沿用定理 \ref{thm:fold-gauge-anomaly-branch-radius-cumulant-growth} 的解析半径 \(R_\theta\)。
对每个整数 \(p\ge 1\) 定义根单位采样点
\[
\theta_p:=\frac{2\pi i}{p}.
\]
则 \(P_G\) 在 \(\theta=0\) 的 Taylor 展开在 \(\theta=\theta_p\) 处收敛当且仅当
\[
|\theta_p|<R_\theta
\quad\Longleftrightarrow\quad
p>\frac{2\pi}{R_\theta}.
\]
特别地，取定理 \ref{thm:fold-gauge-anomaly-branch-radius-cumulant-growth} 中的数值
\(
R_\theta=0.7439369247693026
\)，则满足收敛的最小整数为 \(p_{\min}=9\)。
\end{proposition}
\leanverified{paper\_fold\_gauge\_anomaly\_root\_unit\_sampling\_threshold}
\begin{proof}
由复分析基本定理，\(P_G\) 在 \(\theta=0\) 的 Taylor 级数收敛圆盘半径恰等于最近奇点到原点的距离，亦即 \(R_\theta\)。
而 \(|\theta_p|=2\pi/p\)，故收敛当且仅当 \(2\pi/p<R_\theta\)，等价于 \(p>2\pi/R_\theta\)。
代入 \(R_\theta\) 的数值即可得 \(2\pi/R_\theta\approx 8.44\)，从而最小整数为 \(9\)。证毕。
\end{proof}

\endinput

% Auto-generated by scripts/exp_fold_gauge_anomaly_jordan_correlations.py
\paragraph{规范差逐位过程的全相关闭式与 Jordan 指纹}\label{par:fold-gauge-anomaly-jordan}
命题 \ref{prop:fold-gauge-anomaly-pressure} 已将逐位失配指示 $g_t=\mathbf 1\{X_t\neq Y_t\}$ 的压力写为$P_G(\theta)=\log\rho(A_\theta)$ 并给出显式 $4\times 4$ 矩阵 $A_\theta$。取 $\theta=0$ 得 $A_0$，对 Perron 根 $\rho(A_0)=1$ 作 Parry 归一化，可进一步得到该逐位过程的谱指纹、完整自协方差序列与有限长度方差的完全闭式。

\begin{lemma}[Parry 归一化核与失配边势]\label{lem:fold-gauge-anomaly-parry-kernel}
在状态次序 $(a,b,c,d)$ 下，令 $A_0:=A_\theta|_{\theta=0}$。
取 Perron 根 $\rho(A_0)=1$ 的左右本征向量
\[
r=(2,1,2,1)^\top,\qquad \ell=(2,2,1,1)^\top,\qquad \ell^\top r=9,
\]
并定义 Parry 转移核与平稳分布为
\[
P_{ij}:=\frac{(A_0)_{ij}\,r_j}{r_i},\qquad \pi_i:=\frac{\ell_i r_i}{\ell^\top r}.
\]
则 $P$ 与 $\pi$ 具有完全有理闭式：
\[
P=\begin{pmatrix}\frac12&\frac14&0&\frac14\\0&0&1&0\\\frac12&\frac12&0&0\\1&0&0&0\end{pmatrix},\qquad\pi=\Bigl(\frac49,\frac29,\frac29,\frac19\Bigr).
\]
并且 $A_\theta$ 中携带 $e^\theta$ 的边恰对应 $g_t=1$ 的失配边势：在 $P$ 上仅三条转移 $(a\!\to\! b),(b\!\to\! c),(c\!\to\! a)$ 取值 $1$，其余转移取值 $0$。
\end{lemma}

\begin{theorem}[Jordan 指纹]\label{thm:fold-gauge-anomaly-jordan-fingerprint}
\leanverified{paper\_fold\_gauge\_anomaly\_jordan\_fingerprint}
上述 $4$ 态核 $P$ 的谱为 $\{1,\tfrac12,-\tfrac12\}$，且特征值 $-\tfrac12$ 的代数重数为 $2$、几何重数为 $1$，从而存在大小为 $2$ 的 Jordan 块。
因此 $P^k$ 的衰减项一般包含 $k(-\tfrac12)^k$ 的广义特征贡献：对任意中心化可加观测量，其相关衰减在一般位置呈现 $k2^{-k}$ 级的多项式修正（并伴随 $(-1)^k$ 振荡），而非纯指数。
\end{theorem}

\begin{theorem}[自协方差核与生成函数的完全闭式]\label{thm:fold-gauge-anomaly-autocov-closed}
在平稳分布 $\pi$ 下，逐位失配指示满足
\[
\mathbb E[g_t]=\frac49,\qquad \mathrm{Var}(g_t)=\frac{20}{81}.
\]
并且对任意 $k\ge 1$，自协方差具有闭式
\[
\mathrm{Cov}(g_t,g_{t+k})=2^{-k}\!\left(\frac18+(-1)^k\!\left(\frac{7}{648}+\frac{k}{108}\right)\right).
\]
其协方差生成函数为显式有理函数：
\[
\sum_{k\ge 0}\mathrm{Cov}(g_t,g_{t+k})\,z^k=\frac{20}{81}+\frac{z}{16(1-z/2)}-\frac{7z}{1296(1+z/2)}-\frac{z}{216(1+z/2)^2}.
\]
\end{theorem}

\begin{theorem}[有限长度方差的完全闭式]\label{thm:fold-gauge-anomaly-finite-var-closed}
令 $G_m:=\sum_{t=1}^{m} g_t$。则在平稳初值下有完全闭式
\[
\mathrm{Var}(G_m)=\frac{118}{243}m-\frac{40}{81}+\frac{243-(-1)^m(2m+3)}{486\cdot 2^m}.
\]
因此 $\mathrm{Var}(G_m)-\frac{118}{243}m$ 的偏差不仅为 $O(1)$，而且其剩余项精确为 $O(m2^{-m})$（并带显式振荡项）。
\end{theorem}

\begin{corollary}[Edgeworth 一阶项的可审计封口]\label{cor:fold-gauge-anomaly-edgeworth1}
\leanverified{paper\_fold\_gauge\_anomaly\_edgeworth1}
令 $Z_m:=\frac{G_m-\frac49 m}{\sqrt{\frac{118}{243}m}}$。在中心极限定理尺度上，其一阶 Edgeworth 近似系数可写为完全闭式：
\[
\mathbb P(Z_m\le x)=\Phi(x)+\frac{1}{6\sqrt m}\,\gamma_1^{(G)}\,(1-x^2)\varphi(x)+O(m^{-1}),\qquad\gamma_1^{(G)}=\frac{P_G^{(3)}(0)}{(P_G''(0))^{3/2}}=\frac{-\frac{1174}{2187}}{\left(\frac{118}{243}\right)^{3/2}}.
\]
其中 $\Phi,\varphi$ 分别为标准正态分布函数与密度，$P_G^{(3)}(0)$ 与 $P_G''(0)$ 由命题 \ref{prop:fold-gauge-anomaly-pressure} 给出。
\end{corollary}

\begin{proof}[证明要点（可审计）]
引理 \ref{lem:fold-gauge-anomaly-parry-kernel} 由 $A_0$ 的 Perron 本征向量与 Parry 归一化直接计算得到。定理 \ref{thm:fold-gauge-anomaly-jordan-fingerprint} 的特征多项式分解与 $-\tfrac12$ 处 Jordan 块由有限维有理线性代数给出。定理 \ref{thm:fold-gauge-anomaly-autocov-closed} 通过在平稳链上把 $g_t$ 视为边势并利用 $P$ 的谱/Jordan 分解，得到自协方差生成函数的有理表达并逐项抽取系数。定理 \ref{thm:fold-gauge-anomaly-finite-var-closed} 由$\mathrm{Var}(G_m)=m\,\mathrm{Var}(g_t)+2\sum_{k=1}^{m-1}(m-k)\mathrm{Cov}(g_t,g_{t+k})$与上一条协方差闭式求和得到。以上全部代数化简与一致性校验可由脚本 \texttt{scripts/exp\_fold\_gauge\_anomaly\_jordan\_correlations.py} 一键复算。
\end{proof}

由定理 \ref{thm:fold-gauge-anomaly-autocov-closed} 的协方差闭式与有理生成函数，可把临界 \(p=\tfrac12\) 的 Jordan 修正压缩为一个“有限支撑、有限阶”的分布矩表示；
由此得到 Hankel 行列式的完全代数证书，从而把 Jordan 临界性转写为无需谱分解的有限维判别。

\begin{proposition}[协方差核的分布矩表示]\label{prop:fold-gauge-anomaly-covariance-distribution-moment}
在均匀基线（\(p=\tfrac12\)）下，令
\[
c_k:=\operatorname{Cov}(g_0,g_k)\qquad(k\ge 1),
\]
其中 \(g_t\) 为逐位失配指示（定理 \ref{thm:fold-gauge-anomaly-autocov-closed}）。
则存在唯一的阶 \(1\) 分布 \(\nu\) 支撑于 \(\{\tfrac12,-\tfrac12\}\)，使得对所有 \(k\ge 1\) 成立
\[
c_k=\langle \nu, x^k\rangle.
\]
并且该分布具有显式表示
\[
\boxed{\;
\nu=\frac18\,\delta_{1/2}+\frac{7}{648}\,\delta_{-1/2}+\frac{1}{216}\,\delta'_{-1/2}\; } ,
\]
其中 \(\delta_a\) 为 Dirac 分布、\(\delta'_a\) 为其导数分布。
\leanverified{paper\_fold\_autocovariance\_seeds}
\leanverified{paper\_fold\_gauge\_anomaly\_covariance\_distribution\_moment}
\end{proposition}

\begin{proof}
由定理 \ref{thm:fold-gauge-anomaly-autocov-closed} 的闭式
\(
c_k=2^{-k}\bigl(\tfrac18+(-1)^k(\tfrac{7}{648}+\tfrac{k}{108})\bigr)
\)
可写为
\[
c_k=\frac18\Bigl(\frac12\Bigr)^k+\frac{7}{648}\Bigl(-\frac12\Bigr)^k+\frac{1}{108}k\Bigl(-\frac12\Bigr)^k.
\]
对分布作用有 \(\langle\delta_a,x^k\rangle=a^k\)，并且
\(
\langle \delta'_a,x^k\rangle=-k a^{k-1}
\)。
取 \(a=-\tfrac12\) 得
\[
\Bigl\langle \frac{1}{216}\delta'_{-1/2},x^k\Bigr\rangle
=-\frac{1}{216}k\Bigl(-\frac12\Bigr)^{k-1}
=\frac{1}{108}k\Bigl(-\frac12\Bigr)^k,
\]
从而上式等价于所示 \(\nu\) 的矩等式。
\par
唯一性来自生成函数的部分分式分解唯一性：定理 \ref{thm:fold-gauge-anomaly-autocov-closed} 给出
\[
\sum_{k\ge 1}c_k z^k
=\frac{z}{16(1-z/2)}-\frac{7z}{1296(1+z/2)}-\frac{z}{216(1+z/2)^2},
\]
而基底 \((1-z/2)^{-1},(1+z/2)^{-1},(1+z/2)^{-2}\) 上的分解系数唯一，等价于
\(\bigl\{(1/2)^k,(-1/2)^k,k(-1/2)^k\bigr\}\) 的线性独立性。证毕。
\end{proof}

\begin{proposition}[综合自相关时间的严格闭式]\label{prop:fold-gauge-anomaly-tau-int-closed}
在定理 \ref{thm:fold-gauge-anomaly-autocov-closed} 的平稳口径下，令
\(
C_k:=\mathrm{Cov}(g_0,g_k)
\)（\(k\ge 0\)），并定义综合自相关时间
\[
\tau_{\mathrm{int}}
:=
1+\frac{2}{C_0}\sum_{k\ge 1}C_k.
\]
则成立严格恒等式
\[
\boxed{\;\tau_{\mathrm{int}}=\frac{59}{30}.\;}
\]
\end{proposition}
\begin{proof}
由定理 \ref{thm:fold-gauge-anomaly-autocov-closed}，有 \(C_0=\frac{20}{81}\)，并且对 \(k\ge 1\),
\[
C_k=\frac18\left(\frac12\right)^k+\left(\frac{7}{648}+\frac{k}{108}\right)\left(-\frac12\right)^k.
\]
逐项求和得
\[
\sum_{k\ge 1}\frac18\left(\frac12\right)^k=\frac18,
\qquad
\sum_{k\ge 1}\frac{7}{648}\left(-\frac12\right)^k=-\frac{7}{1944},
\]
并由恒等式 \(\sum_{k\ge 1}kr^k=\frac{r}{(1-r)^2}\)（\(|r|<1\)）取 \(r=-\frac12\) 得
\[
\sum_{k\ge 1}\frac{k}{108}\left(-\frac12\right)^k=-\frac{1}{486}.
\]
因此
\(
\sum_{k\ge 1}C_k=\frac{29}{243}
\)。
代入定义即得
\[
\tau_{\mathrm{int}}
=1+\frac{2}{20/81}\cdot\frac{29}{243}
=1+\frac{29}{30}
=\frac{59}{30}.
\]
证毕。
\end{proof}
\leanverified{paper\_fold\_gauge\_anomaly\_tau\_int\_closed}

\begin{theorem}[显式 Hankel 证书与 Jordan 临界性判别]\label{thm:fold-gauge-anomaly-hankel-jordan-certificate}
在命题 \ref{prop:fold-gauge-anomaly-covariance-distribution-moment} 的记号下，定义 Hankel 矩阵
\[
H^{(r)}:=\bigl(c_{i+j+1}\bigr)_{0\le i,j\le r-1}\qquad(r\ge 1).
\]
则成立
\[
\boxed{\ \det H^{(3)}=-\frac{1}{2985984}\neq 0,\qquad \det H^{(4)}=0\ },
\]
从而 \(\operatorname{rank}H^{(\infty)}=3\)。
此外，当 \(p\neq\tfrac12\) 时，协方差核为至多两指数项的线性组合（定理 \ref{thm:fold-bernoulli-p-full-autocovariance-jordan}），故 Hankel 秩不超过 \(2\)，从而所有 \(3\times 3\) Hankel 主子式必为零。
因此 \(\det H^{(3)}\) 给出一个完全代数的 Jordan 临界性判别证书。
\end{theorem}

\begin{proof}
由命题 \ref{prop:fold-gauge-anomaly-covariance-distribution-moment}，序列 \((c_k)_{k\ge 1}\) 为支撑于 \(\{\tfrac12,-\tfrac12\}\) 且在 \(-\tfrac12\) 处具有一阶合流的广义幂和序列；
其最小湮灭多项式为 \((\lambda-\tfrac12)(\lambda+\tfrac12)^2\)，从而 \((c_k)\) 满足三阶线性递推，推出一切 \(4\times 4\) Hankel 行列式恒为零，特别地 \(\det H^{(4)}=0\) 且 \(\operatorname{rank}H^{(\infty)}\le 3\)。
\par
另一方面，将 \(c_k\) 的显式公式代入 \(H^{(3)}\) 可直接计算得 \(\det H^{(3)}=-1/2985984\neq 0\)，故 \(\operatorname{rank}H^{(\infty)}\ge 3\)。
两侧合并得到 \(\operatorname{rank}H^{(\infty)}=3\)。
\par
当 \(p\neq\tfrac12\) 时，定理 \ref{thm:fold-bernoulli-p-full-autocovariance-jordan} 给出协方差为两条指数模态之和，从而 Hankel 秩不超过 \(2\)，一切 \(3\times 3\) Hankel 主子式为零。
因此在同一对象族中，\(\det H^{(3)}\neq 0\) 与 Jordan 临界性等价。
可复现核验脚本：\texttt{scripts/exp\_fold\_gauge\_anomaly\_covariance\_hankel\_certificate\_audit.py}。证毕。
\end{proof}
\leanverified{paper\_fold\_gauge\_anomaly\_hankel\_jordan\_certificate}

\endinput

% Auto-generated by scripts/exp_fold_gauge_anomaly_spectral_fingerprint.py
\begin{remark}[可审计的谱指纹—相位残差接口：从 $A_\theta$ 推出二点统计与 Hankel 证书]\label{rem:fold-gauge-anomaly-spectral-fingerprint-intro}
在命题 \ref{prop:fold-gauge-anomaly-pressure} 的矩阵接口上，进一步可推出一组原文未显式列出的、但完全可复算的“谱指纹—相关闭式—Hankel 秩证书”结论：其中 $(-1/2)$ 的 Jordan 缺陷在二点相关中严格显影为交错相位残差。它们把规范差的统计涨落从叙事性表述转写为“有限统计量即可核验”的硬约束。下述所有闭式均可由脚本 \texttt{scripts/exp\_fold\_gauge\_anomaly\_spectral\_fingerprint.py} 一键复算。
\end{remark}

\paragraph{未倾斜算子的谱指纹与 Jordan 残差}
\begin{theorem}[$A_0$ 的精确谱指纹与 Jordan 结构]\label{thm:fold-gauge-anomaly-A0-jordan}
令 $A_0:=A_\theta|_{\theta=0}$ 为命题 \ref{prop:fold-gauge-anomaly-pressure} 中的未倾斜加权邻接矩阵。则其特征多项式满足分解
\[
\boxed{\ \chi_{A_0}(\lambda)=\det(\lambda I-A_0)=\frac{\left(\lambda - 1\right) \left(2 \lambda - 1\right) \left(2 \lambda + 1\right)^{2}}{8}\ }.
\]
因此\(\mathrm{spec}(A_0)=\{1,\frac12,-\frac12\}\)，其中 $\lambda=-\frac12$ 的代数重数为 $2$、几何重数为 $1$（存在一个 $2\times 2$ Jordan 块）。
\end{theorem}

\begin{remark}[结构性含义：交错相位模态与一阶 nilpotent 残差]\label{rem:fold-gauge-anomaly-jordan-meaning}
定理 \ref{thm:fold-gauge-anomaly-A0-jordan} 表明，规范差的涨落除常规混合模态外还携带一个载荷为 $-1/2$ 的交错相位模态；其 Jordan 缺陷在统计投影上产生“指数衰减 $\times$ 线性多项式”的可见尾项，从而可由二点统计直接捕捉（见下述关于二点相关闭式与 $O(k2^{-k})$ 衰减律的段落）。
\end{remark}

\paragraph{二点相关闭式与 $O(k2^{-k})$ 衰减律}
\begin{theorem}[失配指示过程的二点协方差闭式]\label{thm:fold-gauge-anomaly-cov-closed}
令 $g_t=\mathbf 1\{X_t\neq Y_t\}$ 为命题 \ref{prop:fold-gauge-anomaly-pressure} 的逐位失配指示过程，并取其在均匀基线下的平稳极限口径。则
\[
\mathbb E[g_t]=P_G'(0)=\frac49.
\]
并且对任意 $k\ge 1$，二点协方差满足完全闭式
\[
\boxed{\ \mathrm{Cov}(g_t,g_{t+k})=\frac1{16}\Big(\frac12\Big)^{k-1}+\Big(-\frac{13}{1296}-\frac{k-1}{216}\Big)\Big(-\frac12\Big)^{k-1}\ }.
\]
因此存在严格渐近衰减律\(\mathrm{Cov}(g_t,g_{t+k})=O(k\,2^{-k})\)，其中线性因子 $k$ 正对应于定理 \ref{thm:fold-gauge-anomaly-A0-jordan} 的 $(-1/2)$ Jordan 残差。
\end{theorem}

\begin{corollary}[协方差生成函数的有理闭式]\label{cor:fold-gauge-anomaly-cov-genfun}
令
\[
C(z):=\sum_{k\ge 1}\mathrm{Cov}(g_t,g_{t+k})\,z^k.
\]
则 $C$ 为有理函数，且具有闭式
\[
\boxed{\ C(z)=\frac{z/16}{1-z/2}-\frac{13z/1296}{1+z/2}+\frac{z^2/432}{(1+z/2)^2}\ }.
\]
\end{corollary}

\paragraph{有限样本可核验的 Hankel 秩证书与递推律}
\begin{theorem}[三阶线性递推与 Hankel 秩 $=3$]\label{thm:fold-gauge-anomaly-hankel-rank3}
令 $d_k:=\mathrm{Cov}(g_t,g_{t+k})$（$k\ge 1$）。则对所有 $k\ge 1$ 有严格递推
\[
\boxed{\ d_{k+3}=-\frac12\,d_{k+2}+\frac14\,d_{k+1}+\frac18\,d_k\ }.
\]
等价地，以 $(d_k)$ 生成的无限 Hankel 矩阵 $H=(d_{i+j+1})_{i,j\ge 0}$ 满足 $\mathrm{rank}(H)=3$。
\end{theorem}

\begin{corollary}[无参一致性的有限数据审计口径]\label{cor:fold-gauge-anomaly-hankel-audit}
任意实现若声称其“折叠归一化—投影读出”与本文均匀基线机制一致，则其从观测数据估计得到的滞后协方差序列 $(\widehat d_k)$ 必须近似满足定理 \ref{thm:fold-gauge-anomaly-hankel-rank3} 的递推关系（等价于低 Hankel 秩约束）。系统性偏离意味着存在额外不可见自由度或口径结构发生改变，而非仅为参数扰动。
\end{corollary}

\paragraph{有限长度方差的完全闭式（CLT 常数的全量校验）}
\begin{theorem}[$G_m=\sum_{t=1}^m g_t$ 的方差全闭式]\label{thm:fold-gauge-anomaly-var-finite-closed}
在同一平稳口径下，对任意 $m\ge 1$，有
\[
\boxed{\ \mathrm{Var}(G_m)=\frac{118}{243}\,m-\frac{40}{81}+\frac{243-(-1)^m(2m+3)}{486\cdot 2^m}\ }.
\]
从而 \(\mathrm{Var}(G_m)/m=\frac{118}{243}+O(2^{-m})\)，并且指数小修正项携带 $(-1)^m$ 的交错相位签名，与定理 \ref{thm:fold-gauge-anomaly-A0-jordan} 的 $(-1/2)$ 模态一致。
\leanverified{paper\_fold\_gauge\_anomaly\_var\_finite\_closed}
\end{theorem}

\begin{proposition}[可检验的相位残差定律（统一陈述）]\label{prop:fold-gauge-anomaly-spectral-residual-law}
在命题 \ref{prop:fold-gauge-anomaly-pressure} 的无参数矩阵接口与均匀基线之下，定理 \ref{thm:fold-gauge-anomaly-A0-jordan}--\ref{thm:fold-gauge-anomaly-var-finite-closed} 共同给出一个低维、可识别且可审计的残差判据：
\begin{enumerate}
\item 未倾斜算子 $A_0$ 出现 $(-1/2)$ 的 Jordan 缺陷（定理 \ref{thm:fold-gauge-anomaly-A0-jordan}）。
\item 二点协方差呈 $O(k2^{-k})$ 级交错衰减，并带有线性多项式因子（定理 \ref{thm:fold-gauge-anomaly-cov-closed}）。
\item 协方差序列满足三阶线性递推，等价于 Hankel 秩 $=3$ 的证书约束（定理 \ref{thm:fold-gauge-anomaly-hankel-rank3}）。
\item 有限长度方差具有完全闭式，并含有显式 $(-1)^m$ 的交错相位修正（定理 \ref{thm:fold-gauge-anomaly-var-finite-closed}）。
\end{enumerate}
因此，“是否确已把不可逆性外置为正交记录轴”的断言在该最小模型上具备了可复算、可证伪的统计硬接口，并与命题 \ref{prop:unit-circle-phase-extension-preserving} 的扩展保存口径一致。
\end{proposition}

\begin{table}[H]
\centering
\scriptsize
\setlength{\tabcolsep}{6pt}
\caption{Elimination certificate for the gauge-anomaly rate curve. We eliminate $\mu$ from $F(\mu,u)=0$ and $\alpha\mu F_{\mu}(\mu,u)+uF_u(\mu,u)=0$ to obtain a resultant in $\mathbb{Z}[\alpha,u]$. We then divide out the maximal $u$-power (the $u$-adic valuation) and the coefficient content gcd to obtain the primitive normalized $R(\alpha,u)$ (see script).}
\label{tab:fold_gauge_anomaly_rate_curve_resultant_degree}
\begin{tabular}{l l}
\toprule
$\deg_{\alpha} R$ & $4$\\
$\deg_{u} R$ & $11$\\
\#terms & $46$\\
$\mathrm{content}(R)$ & $1$\\
leading monomial & $\alpha^{4}u^{11}$\\
leading coeff & $-27$\\
$v_u(R_{\mathrm{raw}})$ & $2$\\
\bottomrule
\end{tabular}
\end{table}

% Auto-generated; do not edit by hand.
\begin{align}
F(\mu,u)&=\mu^{4} - \mu^{3} - \mu^{2} \left(2 u + 1\right) + \mu \left(- u^{3} + 2 u\right) + 2 u,\\
R(\alpha,u)&\in\mathbb{Z}[\alpha,u],\qquad \deg_{\alpha}R=4,\ \deg_u R=11,\\
R(\alpha,1)&=- \left(4 \alpha + 1\right) \left(9 \alpha - 4\right),\qquad \Rightarrow\ \alpha(1)=\tfrac49.
\end{align}

\begin{proposition}[规范差速率曲线的代数参数化与可行分支选择]\label{prop:fold-gauge-anomaly-rate-curve-param}
在命题 \ref{prop:fold-gauge-anomaly-pressure} 的记号下，令 $u=e^\theta$，并令 $\mu(u)$ 表示代数方程 $F(\mu,u)=0$ 的 Perron 实根分支（即 $\mu(u)=2\rho(A_\theta)>0$）。
定义倾斜下的典型失配密度（速率曲线）
\[
\alpha(u):=\frac{d}{d\theta}P_G(\theta)\Big|_{\theta=\log u}.
\]
则由 $P_G(\theta)=\log(\mu(e^\theta)/2)$ 与隐函数微分（仅用 $F$）得纯代数参数化：
\[
\boxed{\ \alpha(u)=\frac{u\,\mu_u(u)}{\mu(u)}=-\frac{u\,F_u(\mu,u)}{\mu\,F_\mu(\mu,u)}\Big|_{\mu=\mu(u)}\ }
\]
其中 $F_u,F_\mu$ 为 $F(\mu,u)$ 的偏导数。
并且大偏差率函数在该曲线上的 Legendre 驻点闭式为：
\[
\boxed{\ I(\alpha(u))=\alpha(u)\log u-\log\frac{\mu(u)}{2}\ }.
\]
此外，消元证书 $R(\alpha,u)=0$ 作为平面代数曲线包含多个代数分支；在 $u=1$ 处出现的另一个根 $\alpha=-\tfrac14$ 是\emph{不可行分支}，因为失配密度必须属于区间 $[0,1]$。因此可用以下规则选出可行分支：
\[
\mu(u)>0,\qquad \alpha(u)\in[0,1],\qquad (\mu(1),\alpha(1))=(2,4/9).
\]
该规则给出一个可审计的代数相图参数化 $(u,\alpha(u),I(\alpha(u)))$，并显式排除不可行根支。
\end{proposition}
\leanverified{paper\_fold\_gauge\_anomaly\_rate\_curve\_param}

在率曲线消元证书 \(R(\alpha,u)=0\) 之上，可以进一步抽取“反解病态”与“分歧支撑”的严格代数指纹。
其核心在于：把 \(R\) 视为关于 \(\alpha\) 的四次多项式，其判别式不仅携带谱曲线分歧十次 \(P_{10}\)，而且出现一个额外的平方因子，对应于 \(\alpha\)-投影的非单值化壁。

\begin{theorem}[率曲线消元判别式的分歧分解与同域因子]\label{thm:fold-gauge-anomaly-rate-curve-disc-factorization-q19}
沿用命题 \ref{prop:fold-gauge-anomaly-rate-curve-param} 的谱四次
\(
F(\mu,u)=0
\)
与率曲线消元多项式 \(R(\alpha,u)\in\ZZ[\alpha,u]\)（\(\deg_\alpha R=4,\deg_u R=11\)）。
记分歧十次 \(P_{10}(u)\) 由
\[
\mathrm{Disc}_{\mu}\!\bigl(F(\mu,u)\bigr)=-u(u-1)P_{10}(u)
\]
给出（命题 \ref{prop:fold-gauge-anomaly-P10-galois-discriminant}）。
则 \(R(\alpha,u)\) 关于 \(\alpha\) 的判别式在 \(\ZZ[u]\) 中满足严格分解
\begin{equation}\label{eq:fold-gauge-anomaly-rate-curve-disc-alpha-factorization}
\boxed{\;
\mathrm{Disc}_{\alpha}\!\bigl(R(\alpha,u)\bigr)
=-u^{11}(u-1)\,P_{10}(u)\cdot Q_{19}(u)^{2}\;},
\end{equation}
其中 \(Q_{19}(u)\in\ZZ[u]\) 为次数 \(19\) 的不可约多项式
\[
\begin{aligned}
Q_{19}(u)=\;&137781u^{19}+629856u^{18}+934578u^{17}-1546209u^{16}-6187752u^{15}+1402704u^{14}\\
&+15349014u^{13}-3368736u^{12}-17688806u^{11}+2216732u^{10}+17425320u^9-4765670u^8\\
&-11620472u^7+7448180u^6+2529720u^5-4736240u^4+2386300u^3-612800u^2+81800u-4500.
\end{aligned}
\]
\end{theorem}

\begin{proof}
由定义 \(R(\alpha,u)=\Res_{\mu}\bigl(F(\mu,u),\ \alpha\mu F_{\mu}(\mu,u)+uF_u(\mu,u)\bigr)\)（脚本 \texttt{scripts/exp\_fold\_gauge\_anomaly\_rate\_curve\_elimination.py}），故 \(\mathrm{Disc}_{\alpha}(R)\) 可在 \(\ZZ[u]\) 中以判别式--结果式恒等式显式计算。
把生成的原始规范化多项式 \(R\) 代入并在 \(\ZZ[u]\) 中完全因式分解，即得 \eqref{eq:fold-gauge-anomaly-rate-curve-disc-alpha-factorization}；其中 \(P_{10}\) 以一次幂出现，余下非平凡部分为一个不可约因子 \(Q_{19}\) 的平方。
可复现核验脚本：\texttt{scripts/exp\_fold\_gauge\_anomaly\_rate\_curve\_discriminant\_q19\_audit.py}。证毕。
\end{proof}
\leanverified{paper\_fold\_gauge\_anomaly\_rate\_curve\_disc\_factorization\_q19}

\begin{theorem}[\texorpdfstring{$Q_{19}$}{Q19} 的满对称伽罗瓦性]\label{thm:fold-gauge-anomaly-q19-galois-s19}
\leanverified{paper\_fold\_gauge\_anomaly\_q19\_galois\_s19}
令 \(Q_{19}(u)\) 如定理 \ref{thm:fold-gauge-anomaly-rate-curve-disc-factorization-q19}。
则
\[
\mathrm{Gal}\bigl(Q_{19}/\QQ\bigr)\cong S_{19}.
\]
\end{theorem}

\begin{proof}
取素数 \(p=73\)。模 \(73\) 约化后 \(Q_{19}\bmod 73\) 在 \(\FF_{73}[u]\) 中保持次数 \(19\) 不可约，故 \(Q_{19}\) 在 \(\QQ[u]\) 上不可约；
其伽罗瓦群在 \(19\) 个根上作用传递。
由于 \(19\) 为素数，传递即推出本原性。
\par
再取素数 \(p=17\)。可计算 \(Q_{19}\bmod 17\) 的次数分解型为 \((11,8)\)；
对不整除判别式的素数，Dedekind--Frobenius 对应原理给出 Frobenius 元在根集合上的循环型与该次数型一致，
从而伽罗瓦群含有一条 \(11\)-循环。
由于 \(11\) 为素数且 \(11\le 19-3\)，Jordan 定理断言：本原子群若含 \(11\)-循环，则必包含 \(A_{19}\)。
\par
最后，若伽罗瓦群包含于 \(A_{19}\)，则判别式应为 \(\ZZ\) 中平方。
但计算 \(\mathrm{Disc}(Q_{19})\bmod 13\equiv 8\)，而 \(8\) 在 \(\FF_{13}\) 中为二次非剩余，
从而判别式非平方，伽罗瓦群不含于 \(A_{19}\)。
结合 \(A_{19}\subseteq \mathrm{Gal}(Q_{19}/\QQ)\) 得
\(\mathrm{Gal}(Q_{19}/\QQ)=S_{19}\)。
可复现核验脚本：\texttt{scripts/exp\_fold\_gauge\_anomaly\_rate\_curve\_discriminant\_q19\_audit.py}。证毕。
\end{proof}

\begin{corollary}[根的不可根式性]\label{cor:fold-gauge-anomaly-q19-nonradical}
\(Q_{19}\) 的任一根均不可由有理数经有限次四则运算与开方表示。
因此定理 \ref{thm:fold-gauge-anomaly-rate-curve-disc-factorization-q19} 中由 \(Q_{19}(u)=0\) 给出的额外判别式零点集合在可解性意义上达到最大对称情形。
\leanverified{paper\_fold\_gauge\_anomaly\_q19\_nonradical}
\end{corollary}

\endinput

% Auto-generated by scripts/exp_fold_gauge_anomaly_ldp_rate.py
\begin{proposition}[规范差的大偏差率函数接口（Legendre 变换 + 端点闭式）]\label{prop:fold-gauge-anomaly-ldp}
在均匀基线 $\omega\sim\mathrm{Unif}(\Omega_{m+1})$ 下，令 $G_m=\sum_{t=1}^{m}g_t$，其中 $g_t=\mathbf 1\{X_t\neq Y_t\}$ 为规范差的逐位失配指示。则对任意 $a\in[0,1]$，比例 $G_m/m$ 满足大偏差原理，其速率函数可由压力函数做 Legendre--Fenchel 变换给出：
\[
I(a)=\sup_{\theta\in\mathbb R}\bigl(a\theta-P_G(\theta)\bigr),\qquad P_G(\theta)=\log\rho(A_\theta),
\]
其中 $A_\theta$ 为命题 \ref{prop:fold-gauge-anomaly-pressure} 的 $4\times 4$ 加权邻接矩阵。并且端点的稀有事件指数具有完全闭式：
\[
\boxed{\ I(0)=\log(2/\varphi),\qquad I(1)=\log 2\ }.
\]
\end{proposition}

\begin{proof}[证明要点（可审计）]
由于 $g_t$ 是有限状态边移位上的可加观测量，G\"artner--Ellis 定理给出 $I$ 为 $P_G$ 的 Legendre--Fenchel 变换。端点闭式来自极限扭曲：当 $\theta\to-\infty$（即 $u=e^\theta\to 0$）时，失配边权趋于 $0$，谱半径退化为“全匹配子图”的 Perron 根，给出 $P_G(-\infty)=\log(\varphi/2)$，故 $I(0)=-P_G(-\infty)=\log(2/\varphi)$；当 $\theta\to+\infty$ 时，谱半径由全失配三环支配，满足 $P_G(\theta)=\theta-\log 2+o(1)$，故 $I(1)=\log 2$。
\end{proof}


\paragraph{Bernoulli\texorpdfstring{$(p)$}{(p)} 基线下的换能器核：失配加性泛函的显式极限律与代数压力}\label{par:fold-gauge-anomaly-bernoulli-p}
本节此前的“均匀基线”对应 IID Bernoulli\((1/2)\) 输入，其逐位失配指示可由 Berstel 的 \(4\) 状态不歧义 Zeckendorf 归一化换能器实现，并据此得到密度、CLT、压力函数与端点大偏差指数（定理 \ref{thm:fold-gauge-anomaly-density}--命题 \ref{prop:fold-gauge-anomaly-ldp}）。
下面给出一般 IID Bernoulli\((p)\)（\(p\in(0,1)\)）基线下的完全显式推广：失配密度、方差率、压力函数与端点指数均可写为有理/代数闭式；并且转移谱在 \(p=\tfrac12\) 处出现唯一 Jordan 共振，解释了均匀基线相关衰减中的 \(k2^{-k}\) 修正项（定理 \ref{thm:fold-gauge-anomaly-jordan-fingerprint} 的一般化）。

\begin{theorem}[Berstel--Zeckendorf 失配加性泛函：\(p\)-变形压力与一、二阶极限矩闭式]\label{thm:fold-gauge-anomaly-bernoulli-p-laws}
\leanverified{paper\_fold\_gauge\_anomaly\_bernoulli\_p\_laws}
取 \(p\in(0,1)\) 并记 \(q:=1-p\)。
在状态空间 \(V=\{a,b,c,d\}\) 上考虑平稳马尔可夫链 \((S_n)_{n\ge 0}\)，其转移矩阵为
\begin{equation}\label{eq:fold-gauge-anomaly-bernoulli-p-transition}
P(p)=
\begin{pmatrix}
q & p^2 & 0 & pq\\
0 & 0 & 1 & 0\\
q & p & 0 & 0\\
1 & 0 & 0 & 0
\end{pmatrix}.
\end{equation}
定义边势 \(g:V\times V\to\{0,1\}\) 为
\begin{equation}\label{eq:fold-gauge-anomaly-bernoulli-p-edge-potential}
g(a,b)=g(b,c)=g(c,a)=1,
\qquad
g(i,j)=0\ \ \text{（其余转移）}.
\end{equation}
并令长度 \(m\) 的失配计数为
\begin{equation}\label{eq:fold-gauge-anomaly-bernoulli-p-mismatch-sum}
G_m:=\sum_{n=0}^{m-1} g(S_n,S_{n+1}).
\end{equation}
则成立如下完全显式极限定律。
\begin{enumerate}
  \item \textbf{强大数定律与显式失配密度：}存在函数 \(g_\ast(p)\in(0,1)\) 使得
  \[
  \frac{G_m}{m}\xrightarrow[m\to\infty]{a.s.,\,L^1}g_\ast(p),
  \qquad
  \boxed{\;
  g_\ast(p)=\frac{p^2(3-2p)}{1+p^3}. \;}
  \]

  \item \textbf{中心极限定理与显式方差率：}存在常数 \(\sigma_G^2(p)>0\) 使得
  \[
  \frac{G_m-mg_\ast(p)}{\sqrt{m\,\sigma_G^2(p)}}\ \Rightarrow\ \mathcal N(0,1),
  \]
  且 \(\mathrm{Var}(G_m)=\sigma_G^2(p)\,m+o(m)\)。并且
  \[
  \boxed{\;
  \sigma_G^2(p)=
  \frac{p^2(1-p)\bigl(21 p^5 - 6 p^4 + 14 p^3 - 36 p^2 + 7 p + 9\bigr)}
  {(p+1)^3(p^2-p+1)^3}. \;}
  \]

  \item \textbf{三阶密度累积量：}压力函数在 \(\theta=0\) 的三阶导数（第三阶累积量密度）具有闭式
  \[
  \boxed{\;
  P_{G,p}^{(3)}(0)=\frac{p^2(p-1)\,\Pi_{11}(p)}{(p+1)^5(p^2-p+1)^5}\;},
  \]
  其中
  \[
  \Pi_{11}(p):=
  \begin{aligned}[t]
    &117p^{11}-450p^{10}+242p^9-1360p^8+2537p^7\\
    &\quad-833p^6+301p^5-1076p^4+188p^3+344p^2-37p-27.
  \end{aligned}
  \]

  \item \textbf{代数压力函数与 Perron 矩阵接口：}对任意 \(\theta\in\RR\) 令 \(u=e^{\theta}>0\)，并定义加权邻接矩阵
  \[
  A_{\theta,p}:=
  \begin{pmatrix}
  q & qu & 0 & p\\
  0 & 0 & pu & 0\\
  pu & 1 & 0 & 0\\
  q & 0 & 0 & 0
  \end{pmatrix}.
  \]
  定义极限压力（对数矩母函数的单位长度极限）
  \[
  P_{G,p}(\theta):=\lim_{m\to\infty}\frac{1}{m}\log\EE\!\left[e^{\theta G_m}\right]
  =\lim_{m\to\infty}\frac{1}{m}\log\EE\!\left[u^{G_m}\right].
  \]
  则该极限存在，并满足
  \[
  \boxed{\;P_{G,p}(\theta)=\log\rho(A_{\theta,p})\;}
  \]
  其中 \(\rho(\cdot)\) 为谱半径。
  等价地，对任意 \(u>0\) 定义倾斜矩阵
  \[
  Q_p(u)_{ij}:=P(p)_{ij}\,u^{g(i,j)}.
  \]
  则极限
  \[
  \Lambda_p(u):=\lim_{m\to\infty}\frac{1}{m}\log\EE\bigl[u^{G_m}\bigr]
  \]
  存在，并满足
  \[
  \boxed{\;\Lambda_p(u)=\log\rho\!\bigl(Q_p(u)\bigr)\;}
  \]
  其中 \(\rho(\cdot)\) 为谱半径。
  并且在 \(u=e^\theta\) 下有
  \(
  P_{G,p}(\theta)=\Lambda_p(e^\theta)=\log\rho(A_{\theta,p})=\log\rho\!\bigl(Q_p(e^\theta)\bigr)
  \)。
  进一步，记 \(\lambda(u,p):=\rho(Q_p(u))\)。
  则 \(\lambda(u,p)\) 为下述四次方程在 \((0,\infty)\) 上的唯一最大实根：
  \begin{equation}\label{eq:fold-gauge-anomaly-bernoulli-p-pressure-quartic}
  \boxed{\;
  \lambda^4-(1-p)\lambda^3+p(p-u-1)\lambda^2
  +p(1-p)u(1-pu^2)\lambda+p^2(1-p)u=0.\;}
  \end{equation}
\end{enumerate}
\end{theorem}
\begin{proof}[证明要点]
由 \eqref{eq:fold-gauge-anomaly-bernoulli-p-transition} 可见：从任一状态均可达 \(a\)（例如 \(b\to c\to a\)、\(d\to a\)），且 \(a\to a\) 的自环概率为 \(1-p>0\)，故链不可约且非周期，从而存在唯一平稳分布 \(\pi(p)\) 并满足遍历定理。
\par
设 \(\pi(p)=(\pi_a,\pi_b,\pi_c,\pi_d)\)，解线性方程 \(\pi^\top P(p)=\pi^\top\)、\(\sum_i\pi_i=1\) 得
\[
\pi(p)=\frac{1}{1+p^3}\bigl(1-p,\ p^2,\ p^2,\ p(1-p)^2\bigr).
\]
由遍历定理，
\(
m^{-1}G_m\to \EE_\pi[g(S_0,S_1)]
\)
几乎处处并在 \(L^1\) 收敛。
而
\[
\EE_\pi[g(S_0,S_1)]
=\pi_a P_{ab}+\pi_b P_{bc}+\pi_c P_{ca}
=\frac{p^2(3-2p)}{1+p^3},
\]
得第 (1) 条。
\par
第 (2) 条为有限状态马尔可夫链上有界加性泛函的经典 CLT。
方差率的显式闭式可由 Poisson 方程给出：记
\(
f(i):=\EE[g(S_0,S_1)\mid S_0=i]=\sum_jP_{ij}g(i,j)
=(p^2,1,1-p,0)^\top
\)，
并解 \((I-P(p))h=f-g_\ast(p)\mathbf 1\)（配以规范 \(\pi^\top h=0\) 以唯一化）。
将 \(h\) 代入边势型方差公式即可化简得到所示 \(\sigma_G^2(p)\)。
\par
第 (3) 条由压力四次方程 \eqref{eq:fold-gauge-anomaly-bernoulli-p-pressure-quartic} 在 \((u,\lambda)=(1,1)\) 处对 Perron 分支作三阶隐式求导并结合 \(u=e^\theta\) 的链式法则得到；
化简可得所示闭式表达。
\par
第 (4) 条给出两种等价的矩阵接口。
其一为加权邻接矩阵 \(A_{\theta,p}\)：在 Berstel 的四状态不歧义换能器上，把边标记 \((X,Y)\in\{0,1\}^2\) 的 Bernoulli\((p)\) 输入权重与逐位失配势 \(e^{\theta\mathbf 1\{X\neq Y\}}\) 相乘，并对同起讫点的平行边求和，即得到上述 \(A_{\theta,p}\)；
由不歧义性，\(\EE[e^{\theta G_m}]\) 可写为某个初末向量夹持下的 \(A_{\theta,p}^{\,m}\) 的线性型，从而其指数增长率由 \(\rho(A_{\theta,p})\) 给出。
其二为 Parry 归一化后的倾斜核 \(Q_p(u)\)：在 \(\theta=0\) 时，矩阵
\(
A_{0,p}=
\begin{psmallmatrix}
q & q & 0 & p\\
0 & 0 & p & 0\\
p & 1 & 0 & 0\\
q & 0 & 0 & 0
\end{psmallmatrix}
\)
的 Perron 本征值为 \(1\)，右本征向量可取
\(
r=(q,\ p^2,\ p,\ q^2)^{\mathsf T}
\)，并满足 \(P(p)_{ij}=A_{0,p}(i,j)\,r_j/r_i\)。
因此对任意 \(\theta\in\RR\)（记 \(u=e^\theta\)），有对角相似关系
\(
Q_p(e^\theta)=D_r^{-1}A_{\theta,p}D_r
\)
（\(D_r:=\mathrm{diag}(r)\)），从而 \(\rho(Q_p(e^\theta))=\rho(A_{\theta,p})\) 并得到
\(
P_{G,p}(\theta)=\log\rho(A_{\theta,p})=\Lambda_p(e^\theta)
\)。
四次方程 \eqref{eq:fold-gauge-anomaly-bernoulli-p-pressure-quartic} 为特征方程 \(\det(\lambda I-Q_p(u))=0\) 的显式展开。证毕。
\end{proof}

\begin{proposition}[四阶密度累积量的完全显式式与三阶符号翻转结构]\label{prop:fold-gauge-anomaly-bernoulli-p-kappa4-and-kappa3-signflip}
\leanverified{paper\_fold\_gauge\_anomaly\_bernoulli\_p\_kappa4\_and\_kappa3\_signflip}
在定理 \ref{thm:fold-gauge-anomaly-bernoulli-p-laws} 的记号下，记 \(P_{G,p}(\theta)=\log\rho(Q_p(e^\theta))\)。
则：
\begin{enumerate}
  \item 四阶密度累积量具有完全显式闭式
  \[
  \boxed{\;
  P_{G,p}^{(4)}(0)=\frac{p^2(1-p)\,\Pi_{17}(p)}{(p+1)^7(p^2-p+1)^7}\;},
  \]
  其中
  \[
  \Pi_{17}(p):=
  \begin{aligned}[t]
  &609p^{17}-7572p^{16}+16148p^{15}-30024p^{14}+126169p^{13}-182509p^{12}+114342p^{11}\\
  &\quad-197084p^{10}+260364p^9-65334p^8-20148p^7-40612p^6+12405p^5+19432p^4\\
  &\quad-3652p^3-2466p^2+175p+81.
  \end{aligned}
  \]
  \item 多项式 \(\Pi_{11}(p)\) 在区间 \((0,1)\) 内恰有两个零点 \(p_-<p_+\)，因此 \(P_{G,p}^{(3)}(0)\) 在 \((0,1)\) 内发生两次符号翻转。
  数值上
  \[
  p_-\approx 0.3929567984354658,\qquad
  p_+\approx 0.7651253794543234.
  \]
\end{enumerate}
\end{proposition}
\begin{proof}[证明要点]
由定理 \ref{thm:fold-gauge-anomaly-bernoulli-p-laws}(4) 的 Perron 分支 \(\lambda(u,p)=\rho(Q_p(u))\) 与恒等式
\(
P_{G,p}(\theta)=\log\lambda(e^\theta,p)
\)，对四次特征方程 \(\det(\lambda I-Q_p(u))=0\) 在 \((u,\lambda)=(1,1)\) 处作四阶隐式求导并整理为 \(\theta\)-导数，即得 \(\kappa_4(p)\) 的有理闭式；
对分母作因式化可化为所示 \((p+1)^7(p^2-p+1)^7\) 的结构。
关于 \(\Pi_{11}\) 在 \((0,1)\) 内零点个数的断言，可由 Sturm 序列对根计数得到（并可据此确定 \(\kappa_3\) 的符号翻转次数）；所列数值由高精度求根给出。证毕。
\end{proof}

\endinput


\begin{proposition}[全阶密度累积量的有理性与分母指数律]\label{prop:fold-gauge-anomaly-bernoulli-p-cumulant-denominator-law}
\leanverified{paper\_fold\_gauge\_anomaly\_bernoulli\_p\_cumulant\_denominator\_law}
在定理 \ref{thm:fold-gauge-anomaly-bernoulli-p-laws} 的记号下，记 \(\lambda(u,p):=\rho(Q_p(u))\)（\(u=e^\theta\)），并定义密度累积量
\[
\kappa_r(p):=\left.\frac{d^r}{d\theta^r}P_{G,p}(\theta)\right|_{\theta=0}\qquad(r\ge 1).
\]
则对任意整数 \(r\ge 1\)，\(\kappa_r(p)\) 为 \(p\) 的有理函数，并存在多项式 \(N_r(p)\in\ZZ[p]\) 使得
\[
\boxed{\;\kappa_r(p)=\frac{N_r(p)}{(1+p^3)^{2r-1}}.\;}
\]
更进一步，\(\lambda(u,p)\) 在 \(u=1\) 处的 \(u\)-导数亦满足同分母指数律：存在 \(M_r(p)\in\ZZ[p]\) 使
\[
\boxed{\;\partial_u^{\,r}\lambda(1,p)=\frac{M_r(p)}{(1+p^3)^{2r-1}}.\;}
\]
\end{proposition}
\begin{proof}
令 \(F(\lambda,u,p)\) 为 \eqref{eq:fold-gauge-anomaly-bernoulli-p-pressure-quartic} 左端多项式。
由于 \(Q_p(1)=P(p)\)，并由命题 \ref{prop:fold-gauge-anomaly-bernoulli-p-spectrum} 的谱分解，
\[
F(\lambda,1,p)=\det(\lambda I-P(p))=(\lambda-1)(\lambda+p)\bigl(\lambda^2-p(1-p)\bigr),
\]
从而
\(
\partial_\lambda F(1,1,p)=(1+p)\bigl(1-p(1-p)\bigr)=(1+p)(p^2-p+1)=1+p^3\neq 0
\)。
故由隐函数定理，存在在 \(u=1\) 邻域解析的 Perron 分支 \(u\mapsto \lambda(u,p)\)，且 \(\lambda(1,p)=1\)。
对恒等式 \(F(\lambda(u,p),u,p)\equiv 0\) 逐阶对 \(u\) 求导，可得 \(\partial_u^{\,r}\lambda(1,p)\) 表示为以 \(\partial_\lambda F(1,1,p)=1+p^3\) 为分母、以 \(p\) 的整系数多项式与低阶导数 \(\partial_u^{\,k}\lambda(1,p)\)（\(k<r\)）的有限多项式组合为分子。
据此可归纳得到：存在 \(M_r(p)\in\ZZ[p]\) 使 \(\partial_u^{\,r}\lambda(1,p)=M_r(p)/(1+p^3)^{2r-1}\)。
\par
最后，由 \(P_{G,p}(\theta)=\log\lambda(e^\theta,p)\)，\(\kappa_r(p)\) 是 \(\partial_u^{\,k}\lambda(1,p)\)（\(1\le k\le r\)）的 Fa\`a di Bruno 型有限组合；并且 \(\log\) 与 \(u=e^\theta\) 的链式法则不会引入除 \(1+p^3\) 之外的新分母因子。
因此存在 \(N_r(p)\in\ZZ[p]\) 使 \(\kappa_r(p)=N_r(p)/(1+p^3)^{2r-1}\)。证毕。
\end{proof}
\begin{remark}[换能器加权邻接矩阵口径与 Parry 归一化]\label{rem:fold-gauge-anomaly-bernoulli-p-adjacency-parry}
令 \(q:=1-p\) 且 \(u=e^\theta>0\)。
在 Berstel 的四状态不歧义 Zeckendorf 归一化换能器上，对每条边 \(e\) 赋权
\(
w_\theta(e):=u^{\mathbf 1_{\{X(e)\neq Y(e)\}}}\,\PP(X(e))
\)
（\(\PP(X=1)=p\)、\(\PP(X=0)=q\)），得到加权邻接矩阵（状态序 \(a,b,c,d\)）
\[
A_\theta(p)=
\begin{pmatrix}
q & q u & 0 & p\\
0 & 0 & p u & 0\\
p u & 1 & 0 & 0\\
q & 0 & 0 & 0
\end{pmatrix}.
\]
由有限状态边移位的 Perron--Frobenius 理论可得
\(
\lim_{m\to\infty}m^{-1}\log \EE[\exp(\theta G_m)]=\log\rho(A_\theta(p))
\)。
并且取 \(r=(q,p^2,p,q^2)^{\top}\) 有对角相似关系
\[
\mathrm{diag}(r)^{-1}A_\theta(p)\,\mathrm{diag}(r)=Q_p(u),
\]
从而 \(\rho(A_\theta(p))=\rho(Q_p(u))\)，并与 \eqref{eq:fold-gauge-anomaly-bernoulli-p-pressure-quartic} 的四次代数封闭一致。
\end{remark}
\begin{corollary}[失配密度的单峰性与唯一极大点]\label{cor:fold-gauge-anomaly-bernoulli-p-density-unimodal}
令 \(g_\ast(p)=\frac{p^2(3-2p)}{1+p^3}\)。
则
\begin{equation}\label{eq:fold-gauge-anomaly-bernoulli-p-density-derivative}
\boxed{\;
g_\ast'(p)=
-\frac{3p\,(p^3+2p-2)}{(p+1)^2(p^2-p+1)^2}\; }.
\end{equation}
因此在区间 \(p\in(0,1)\) 上存在唯一 \(p_\star\) 满足
\begin{equation}\label{eq:fold-gauge-anomaly-bernoulli-p-density-maximizer}
p_\star^3+2p_\star-2=0,
\end{equation}
并且 \(g_\ast\) 在 \((0,p_\star)\) 上严格递增、在 \((p_\star,1)\) 上严格递减，从而 \(p_\star\) 为 \(g_\ast\) 的全局唯一极大点。
并且在该极大点处满足精确恒等式
\[
\boxed{\;g_\ast(p_\star)=p_\star^2.\;}
\]
数值上
\[
p_\star=\sqrt[3]{1+\sqrt{35/27}}+\sqrt[3]{1-\sqrt{35/27}},
\qquad
p_\star\approx 0.770917,
\qquad
g_\ast(p_\star)=p_\star^2\approx 0.594313.
\]
\leanverified{paper\_fold\_gauge\_anomaly\_bernoulli\_p\_density\_unimodal}
\end{corollary}
\begin{proof}
由 \eqref{eq:fold-gauge-anomaly-bernoulli-p-density-derivative} 得
\(\mathrm{sgn}\,g_\ast'(p)=-\mathrm{sgn}(p^3+2p-2)\)。
多项式 \(p^3+2p-2\) 在 \((0,1)\) 上严格递增且在端点处取值异号，故其在 \((0,1)\) 内存在且仅存在一根 \(p_\star\)，并据此得到单峰性与唯一极大性。
由驻点方程 \(p_\star^3=2-2p_\star\) 得 \(1+p_\star^3=3-2p_\star\)，代回 \(g_\ast(p)=\frac{p^2(3-2p)}{1+p^3}\) 得 \(g_\ast(p_\star)=p_\star^2\)。
数值近似由求根计算给出。证毕。
\end{proof}

\begin{proposition}[全参数大偏差原理：压力的 Legendre--Fenchel 对偶与端点常数]\label{prop:fold-gauge-anomaly-bernoulli-p-ldp}
在定理 \ref{thm:fold-gauge-anomaly-bernoulli-p-laws} 的平稳设定下，函数 \(\theta\mapsto P_{G,p}(\theta)=\log\rho(A_{\theta,p})\) 在 \(\RR\) 上实解析且严格凸。
并且比例序列 \((G_m/m)_{m\ge 1}\) 在区间 \([0,1]\) 上满足大偏差原理：存在良率函数 \(I_p:[0,1]\to[0,\infty]\) 使得对任意 Borel 集 \(B\subset\RR\)，
\[
-\inf_{x\in B^\circ}I_p(x)\ \le\
\liminf_{m\to\infty}\frac1m\log\PP\!\left(\frac{G_m}{m}\in B\right)
\le
\limsup_{m\to\infty}\frac1m\log\PP\!\left(\frac{G_m}{m}\in B\right)
\le\
-\inf_{x\in \overline B}I_p(x),
\]
且 \(I_p\) 由 Legendre--Fenchel 变换给出：
\[
\boxed{\;
I_p(x)=\sup_{\theta\in\RR}\bigl(\theta x-P_{G,p}(\theta)\bigr). \;}
\]
端点零失配指数满足
\[
\boxed{\;
I_p(0)=-\lim_{\theta\to-\infty}P_{G,p}(\theta)
=-\log\!\Bigl(\frac{q+\sqrt{q^{2}+4pq}}{2}\Bigr). \;}
\]
\leanverified{paper\_fold\_gauge\_anomaly\_bernoulli\_p\_ldp}
\end{proposition}
\begin{proof}[证明要点]
由于 \(A_{\theta,p}\) 在 \(\theta\in\RR\) 上解析且对每个 \(\theta\) 为原始非负矩阵，Perron 根 \(\rho(A_{\theta,p})\) 随 \(\theta\) 实解析；
非平凡边势保证 \(P_{G,p}\) 严格凸。
G\"artner--Ellis 定理适用，从而得到大偏差原理与 Legendre--Fenchel 表达。
端点 \(I_p(0)\) 由 \(u=e^\theta\to 0^+\) 的退化给出：此时仅保留匹配子图，Perron 根满足二次方程 \(\lambda^2-q\lambda-pq=0\)，取正根即得到所示闭式。证毕。
\end{proof}

\begin{proposition}[仅基于折叠读出位的 Bayes 推断相变与代数阈值]\label{prop:fold-gauge-anomaly-bernoulli-p-bayes-phase-transition}
\leanverified{paper\_fold\_gauge\_anomaly\_bernoulli\_p\_bayes\_phase\_transition}
在命题 \ref{prop:fold-gauge-anomaly-bernoulli-p-bitpair-law} 的单点联合律下，仅观测 \(Y_0\) 来估计 \(X_0\)。
设 \(p_{\mathrm c}\in(0,1)\) 为四次方程
\[
p^4-3p^3+3p^2+p-1=0
\]
在 \((0,1)\) 内的唯一实根。则 Bayes 最优判决满足
\[
\widehat X^\star(y)
=\arg\max_{x\in\{0,1\}}\PP(X_0=x\mid Y_0=y)
=
\begin{cases}
y,& 0<p\le p_{\mathrm c},\\
1,& p_{\mathrm c}\le p<1.
\end{cases}
\]
其最小错误率 \(R(p):=\PP(\widehat X^\star(Y_0)\neq X_0)\) 具有分段闭式
\[
R(p)=
\begin{cases}
g_\ast(p)=\dfrac{p^2(3-2p)}{1+p^3},& 0<p\le p_{\mathrm c},\\[4pt]
1-p,& p_{\mathrm c}\le p<1,
\end{cases}
\]
并在临界点满足精确匹配 \(g_\ast(p_{\mathrm c})=1-p_{\mathrm c}\)。
\end{proposition}
\begin{proof}
由命题 \ref{prop:fold-gauge-anomaly-bernoulli-p-bitpair-law}，
\[
\PP(X_0=1,Y_0=1)-\PP(X_0=0,Y_0=1)
=\frac{p\,(p^3+2p^2-3p+1)}{1+p^3}>0,
\]
故在分支 \(Y_0=1\) 上最优判决恒为 \(\widehat X^\star=1\)。
在分支 \(Y_0=0\) 上，比较
\[
\PP(X_0=1,Y_0=0)-\PP(X_0=0,Y_0=0)
=\frac{p^4-3p^3+3p^2+p-1}{1+p^3},
\]
从而阈值由上述四次式决定。该四次式在 \((0,1)\) 仅有一根（例如由 Sturm 序列可直接计数），故 \(p_{\mathrm c}\) 唯一。低偏置相对应判决 \(\widehat X^\star(Y_0)=Y_0\)，高偏置相对应常值判决 \(\widehat X^\star\equiv1\)。分别代入即可得到
\(R(p)=\PP(X_0\neq Y_0)=g_\ast(p)\) 与 \(R(p)=\PP(X_0=0)=1-p\)；二式相等即给出匹配关系。证毕。
\end{proof}

\begin{proposition}[转移谱的完全分解与 \texorpdfstring{$p=\tfrac12$}{p=1/2} 的唯一 Jordan 共振]\label{prop:fold-gauge-anomaly-bernoulli-p-spectrum}
\leanverified{paper\_fold\_gauge\_anomaly\_bernoulli\_p\_spectrum}
对任意 \(p\in(0,1)\)，矩阵 \(P(p)\) 的特征多项式完全分解为
\[
\chi_{P(p)}(\lambda)=(\lambda-1)(\lambda+p)\bigl(\lambda^2-p(1-p)\bigr).
\]
因此
\[
\mathrm{Spec}(P(p))=\Bigl\{1,\ -p,\ \sqrt{p(1-p)},\ -\sqrt{p(1-p)}\Bigr\}.
\]
令 \(\rho_{\mathrm{mix}}(p):=\max\{p,\sqrt{p(1-p)}\}\in(0,1)\)，则 \(P(p)\) 的谱隙为 \(1-\rho_{\mathrm{mix}}(p)\)。
并且：
\begin{enumerate}
  \item 若 \(p\neq\tfrac12\)，则四个特征值两两不同，\(P(p)\) 在 \(\RR\) 上可对角化；
  因而对任意平方可积零均值可观测量，其相关与混合的指数尺度由 \(\rho_{\mathrm{mix}}(p)^n\) 控制。
  \item 若 \(p=\tfrac12\)，则 \(-p=-\sqrt{p(1-p)}=-\tfrac12\) 的代数重数为 \(2\)，而其特征子空间维数为 \(1\)，从而 \(P(\tfrac12)\) 存在唯一非平凡 Jordan 块（大小 \(2\)）。
  相应地，除指数因子 \(2^{-n}\) 外出现线性 \(n\) 修正，与均匀基线下的 Jordan 指纹一致。
\end{enumerate}
\end{proposition}
\begin{proof}
直接计算 \(\det(\lambda I-P(p))\) 得到上述因式分解。
当 \(p\neq\tfrac12\) 时谱单，故可对角化。
当 \(p=\tfrac12\) 时，\(\chi\) 在 \(-\tfrac12\) 处有二重根；并且可直接检验
\(\dim\ker(P(\tfrac12)+\tfrac12 I)=1\)，故几何重数为 \(1\)，从而存在大小 \(2\) 的 Jordan 块。证毕。
\end{proof}

\begin{proposition}[失配协方差序列的三阶递推、Hankel 秩与显式生成函数]\label{prop:fold-gauge-anomaly-bernoulli-p-covariance-recurrence}
\leanverified{paper\_fold\_gauge\_anomaly\_bernoulli\_p\_covariance\_recurrence}
在定理 \ref{thm:fold-gauge-anomaly-bernoulli-p-laws} 的平稳设定下，记
\[
\Delta_n:=g(S_n,S_{n+1})-g_\ast(p),
\qquad
c_k(p):=\mathrm{Cov}(\Delta_0,\Delta_k)\quad(k\ge 1).
\]
则协方差序列满足精确三阶线性递推
\begin{equation}\label{eq:fold-gauge-anomaly-bernoulli-p-covariance-recurrence}
c_{k+3}(p)
=-p\,c_{k+2}(p)+p(1-p)\,c_{k+1}(p)+p^2(1-p)\,c_k(p),
\qquad k\ge 1.
\end{equation}
从而 Hankel 矩阵 \(\bigl(c_{i+j}(p)\bigr)_{i,j\ge 1}\) 的秩至多为 \(3\)。
\par
并且生成函数 \(C_p(z):=\sum_{k\ge 1}c_k(p)z^k\) 为有理函数，具体为
\begin{equation}\label{eq:fold-gauge-anomaly-bernoulli-p-covariance-gf}
\boxed{\;
C_p(z)=
\frac{p^2 z\Bigl(
\bigl(3p^7-14p^6+19p^5-9p^4+2p^3-2p^2+p\bigr)z^2
\;+\;\bigl(-3p^6-2p^5+9p^4+p^3-7p^2+p+1\bigr)z
\;+\;\bigl(-p^6+p^5-6p^4+13p^3-8p^2-2p+2\bigr)
\Bigr)}{(1+p^3)^2 (1+p z)\bigl(1-p(1-p) z^2\bigr)}.\;}
\end{equation}
\end{proposition}
\begin{proof}[证明要点]
令 \(H:=P(p)\circ g\) 为逐项乘积矩阵（\((P\circ g)_{ij}=P_{ij}g_{ij}\)）。
对 \(k\ge 1\) 有标准恒等式
\[
\EE[g(S_0,S_1)g(S_k,S_{k+1})]=\pi^\top H\,P(p)^{k-1}H\,\mathbf 1,
\]
从而
\(
c_k(p)=\pi^\top H\,P(p)^{k-1}H\,\mathbf 1-g_\ast(p)^2
\)。
由命题 \ref{prop:fold-gauge-anomaly-bernoulli-p-spectrum}，
矩阵 \(P(p)\) 在去掉平稳特征方向后之最小多项式整除 \((\lambda+p)(\lambda^2-p(1-p))\)，因此任意序列 \(k\mapsto \alpha^\top P(p)^{k-1}\beta\) 必满足对应三阶递推；对 \(c_k(p)\) 得 \eqref{eq:fold-gauge-anomaly-bernoulli-p-covariance-recurrence}，Hankel 秩结论随即推出。
\par
生成函数由
\[
\sum_{k\ge 1}\pi^\top H\,P(p)^{k-1}H\,\mathbf 1\,z^k
=z\,\pi^\top H(I-zP(p))^{-1}H\mathbf 1
\]
与 \(\sum_{k\ge 1}g_\ast(p)^2z^k=g_\ast(p)^2\frac{z}{1-z}\) 合并，并对分母作代数化简得到 \eqref{eq:fold-gauge-anomaly-bernoulli-p-covariance-gf}。证毕。
\end{proof}
\begin{corollary}[协方差生成函数的部分分式分解与指数衰减率]\label{cor:fold-gauge-anomaly-bernoulli-p-covariance-partial-fractions}
\leanverified{paper\_fold\_gauge\_anomaly\_bernoulli\_p\_covariance\_partial\_fractions}
在命题 \ref{prop:fold-gauge-anomaly-bernoulli-p-covariance-recurrence} 的记号下，取 \(p\neq\tfrac12\)，并定义移位生成函数
\[
\widetilde C_p(z):=\sum_{n\ge 0}c_{n+1}(p)\,z^n=\frac{C_p(z)}{z}.
\]
则存在严格的部分分式分解
\begin{equation}\label{eq:fold-gauge-anomaly-bernoulli-p-covariance-gf-partial-fractions}
\boxed{\;
\widetilde C_p(z)
=\frac{B(p)}{1+p z}
\;+\;
\frac{E(p)+F(p)z}{1-p(1-p)z^2}\;},
\end{equation}
其中
\begin{align}
B(p)&=
-\frac{p^3(p^2-3p+1)^2}{(2p-1)(p+1)(1+p^3)},\label{eq:fold-gauge-anomaly-bernoulli-p-covariance-gf-B}\\
E(p)&=
\frac{p^2(1-p)\bigl(p^5+4p^4-4p^3-5p^2+7p-2\bigr)}{(2p-1)(1+p^3)(p^2-p+1)},\label{eq:fold-gauge-anomaly-bernoulli-p-covariance-gf-E}\\
F(p)&=
\frac{p^2(1-p)^2\bigl(p^5-p^4-6p^3+3p^2+2p-1\bigr)}{(2p-1)(1+p^3)(p^2-p+1)}.\label{eq:fold-gauge-anomaly-bernoulli-p-covariance-gf-F}
\end{align}
因此对每个 \(k\ge 1\)，\(c_k(p)\) 为 \(({-p})^{k-1}\) 与 \((\pm\sqrt{p(1-p)})^{k-1}\) 的显式线性组合。
并且当 \(p\neq\tfrac12\) 时协方差呈纯指数衰减，其衰减率由 \(\max\{p,\sqrt{p(1-p)}\}\) 控制。
点 \(p=\tfrac12\) 为唯一的极点合并位置，对应到命题 \ref{prop:fold-gauge-anomaly-bernoulli-p-spectrum} 的非半单退化；在该点自协方差出现线性多项式修正因子。
\end{corollary}
\begin{proof}
由命题 \ref{prop:fold-gauge-anomaly-bernoulli-p-covariance-recurrence} 的闭式 \eqref{eq:fold-gauge-anomaly-bernoulli-p-covariance-gf}，对 \(p\neq\tfrac12\) 作部分分式分解即可得到 \eqref{eq:fold-gauge-anomaly-bernoulli-p-covariance-gf-partial-fractions} 与 \eqref{eq:fold-gauge-anomaly-bernoulli-p-covariance-gf-B}--\eqref{eq:fold-gauge-anomaly-bernoulli-p-covariance-gf-F}。
指数衰减率与 \(p=\tfrac12\) 的极点合并由命题 \ref{prop:fold-gauge-anomaly-bernoulli-p-spectrum} 的谱分解及标准有理型生成函数系数抽取给出。证毕。
\end{proof}

\begin{corollary}[失配自协方差的显式奇偶闭式与 Jordan 退化阈值]\label{cor:fold-gauge-anomaly-bernoulli-p-covariance-explicit-even-odd}
\leanverified{paper\_fold\_gauge\_anomaly\_bernoulli\_p\_covariance\_explicit\_even\_odd}
保持命题 \ref{prop:fold-gauge-anomaly-bernoulli-p-covariance-recurrence} 的记号，并假设 \(p\neq\tfrac12\)。
定义
\[
A(p):=-\frac{p^3(p^2-3p+1)^2}{(p+1)^2(2p-1)(p^2-p+1)},
\]
\[
C_1(p):=-\frac{p(p-1)\bigl(p^5-p^4-6p^3+3p^2+2p-1\bigr)}{(p+1)(2p-1)(p^2-p+1)^2},
\]
\[
C_0(p):=-\frac{p^2(p-1)\bigl(p^5+4p^4-4p^3-5p^2+7p-2\bigr)}{(p+1)(2p-1)(p^2-p+1)^2}.
\]
则对任意整数 \(m\ge 0\) 有
\[
\boxed{\;
c_{2m+1}(p)\;:=\;A(p)(-p)^{2m}+\bigl(p(1-p)\bigr)^m C_0(p),
\;},
\]
并且对任意整数 \(m\ge 1\) 有
\[
\boxed{\;
c_{2m}(p)\;:=\;A(p)(-p)^{2m-1}+\bigl(p(1-p)\bigr)^m C_1(p)
\;}.
\]
因此当 \(p\neq\tfrac12\) 时，失配自协方差序列为两种指数模态之线性组合：一支由 \((-p)^k\) 给出，另一支由 \(\bigl(\pm\sqrt{p(1-p)}\bigr)^k\) 的奇偶分裂给出。
\end{corollary}
\begin{proof}
由推论 \ref{cor:fold-gauge-anomaly-bernoulli-p-covariance-partial-fractions} 的部分分式分解 \eqref{eq:fold-gauge-anomaly-bernoulli-p-covariance-gf-partial-fractions}，
对 \(\widetilde C_p(z)=\sum_{n\ge 0}c_{n+1}(p)\,z^n\) 作幂级数展开并按奇偶系数抽取即可得到所示形式。
系数的显式闭式由 \eqref{eq:fold-gauge-anomaly-bernoulli-p-covariance-gf-B}--\eqref{eq:fold-gauge-anomaly-bernoulli-p-covariance-gf-F} 直接化简得到，其中
\(
A(p)=-p\,B(p)
\)、\(
C_0(p)=E(p)
\)、\(
C_1(p)=F(p)/(p(1-p))
\)。
证毕。
\end{proof}

\begin{proposition}[端点大偏差指数：零失配与满失配的精确指数]\label{prop:fold-gauge-anomaly-bernoulli-p-endpoint-ldp}
\leanverified{paper\_fold\_gauge\_anomaly\_bernoulli\_p\_endpoint\_ldp}
在定理 \ref{thm:fold-gauge-anomaly-bernoulli-p-laws} 的平稳设定下，极限
\[
I_p(0):=\lim_{m\to\infty}-\frac{1}{m}\log \PP(G_m=0),
\qquad
I_p(1):=\lim_{m\to\infty}-\frac{1}{m}\log \PP(G_m=m)
\]
存在，并且具有显式表达式
\begin{equation}\label{eq:fold-gauge-anomaly-bernoulli-p-endpoint-rates}
\boxed{\;
I_p(0)=-\log\!\Bigl(\frac{(1-p)+\sqrt{(1-p)(1+3p)}}{2}\Bigr),
\qquad
I_p(1)=-\frac13\log\bigl(p^2(1-p)\bigr).\;}
\end{equation}
\end{proposition}
\begin{proof}[证明要点]
注意到 \(\PP(G_m=0)=\EE[0^{G_m}]\)，由定理 \ref{thm:fold-gauge-anomaly-bernoulli-p-laws}(4) 得
\(
\lim_{m\to\infty}\frac{1}{m}\log \PP(G_m=0)=\log\rho(Q_p(0))
\)。
将 \(u=0\) 代入 \(Q_p(u)\) 得到“仅保留匹配边”的非负矩阵，其 Perron 根为
\(\rho(Q_p(0))=\frac{(1-p)+\sqrt{(1-p)(1+3p)}}{2}\)，从而得到 \(I_p(0)\)。
\par
对 \(G_m=m\)，每一步都必须沿失配边前进。
失配边构成唯一 \(3\)-回路 \(a\to b\to c\to a\)，其三步转移概率乘积为 \(p^2(1-p)\)。
因此任意长段“满失配”路径的每步对数概率平均收敛到 \(\frac13\log(p^2(1-p))\)，从而得到 \(I_p(1)\)。证毕。
\end{proof}

\begin{proposition}[Gallavotti--Cohen 缺陷远场常数与黄金分割临界偏置]\label{prop:fold-gauge-anomaly-bernoulli-p-gc-defect-farfield-critical}
\leanverified{paper\_fold\_gauge\_anomaly\_bernoulli\_p\_gc\_defect\_farfield\_critical}
定义 Gallavotti--Cohen 缺陷函数
\[
\Delta_p(\theta):=P_{G,p}(\theta)-\theta-P_{G,p}(-\theta).
\]
则 \(\Delta_p\) 为奇函数，且存在远场极限
\[
\boxed{\;
\lim_{\theta\to+\infty}\Delta_p(\theta)=
\log\frac{(p^2(1-p))^{1/3}}{\alpha(p)},
\qquad
\alpha(p)=\frac{1-p+\sqrt{(1-p)(1+3p)}}{2}.
\;}
\]
并且该远场极限在 \(p\in(0,1)\) 上恰在唯一点
\[
\boxed{\;p=\frac{1+\sqrt5}{4}=\frac{\varphi}{2}\;}
\]
处为零；更精确地，
\[
\operatorname{sign}\Bigl(\lim_{\theta\to+\infty}\Delta_p(\theta)\Bigr)=\operatorname{sign}(4p^2-2p-1).
\]
\end{proposition}
\begin{proof}
奇性由定义直接给出：\(\Delta_p(-\theta)=-\Delta_p(\theta)\)。
\par
令 \(u=e^\theta\)。
当 \(\theta\to-\infty\)（即 \(u\to 0^+\)）时，矩阵 \(A_{\theta,p}\) 中所有失配边的权重趋于 \(0\)，从而 Perron 根退化为“仅保留匹配边”的可达强连通分量的谱半径：
\[
\lim_{\theta\to+\infty}P_{G,p}(-\theta)=\lim_{u\to 0^+}\log\rho(A_{\theta,p})=\log\alpha(p),
\]
其中 \(\alpha(p)\) 为命题 \ref{prop:fold-gauge-anomaly-bernoulli-p-endpoint-ldp} 所示闭式。
另一方面，当 \(\theta\to+\infty\) 时，主导贡献来自每步均取失配边的唯一三环 \(a\to b\to c\to a\)，其三步权重乘积为 \(p^2(1-p)u^3\)，故由非负矩阵谱半径的最大回路均值支配原理得到
\[
\lim_{\theta\to+\infty}\bigl(P_{G,p}(\theta)-\theta\bigr)=\frac13\log\bigl(p^2(1-p)\bigr).
\]
两式相减即得远场极限公式。
\par
远场极限为零当且仅当 \((p^2(1-p))^{1/3}=\alpha(p)\)，等价于 \(p^2(1-p)=\alpha(p)^3\)。
又 \(\alpha(p)\) 为二次方程 \(\alpha^2-(1-p)\alpha-p(1-p)=0\) 的正根，消元可化为 \(p^2(4p^2-2p-1)=0\)，从而在 \((0,1)\) 上唯一解为 \(p=\frac{1+\sqrt5}{4}=\frac{\varphi}{2}\)。
由于 \(4p^2-2p-1\) 在 \((0,1)\) 上仅有此一根且另一根为负，故其符号与 \(p-\varphi/2\) 在 \((0,1)\) 上一致；
并且远场极限在 \(p\downarrow 0\) 时趋于 \(-\infty\)、在 \(p\uparrow 1\) 时趋于 \(+\infty\)，于是得到所示符号律。证毕。
\end{proof}

\begin{proposition}[Lee--Yang 分岔点 \(u=1\) 的偏置唯一性]\label{prop:fold-gauge-anomaly-bernoulli-p-leyang-branch-u1}
令 \(u=e^\theta\)，并记
\[
\chi(\lambda;u,p):=\det\!\bigl(\lambda I-A_{\theta,p}\bigr).
\]
则 \(\chi\) 是关于 \(\lambda\) 的四次多项式，且其 \(\lambda\)-判别式满足
\[
\Disc_{\lambda}\!\bigl(\chi(\lambda;u,p)\bigr)=p^3(1-p)\,u\,\mathcal D(u,p),
\]
其中 \(\mathcal D(u,p)\in\ZZ[p,u]\) 关于 \(u\) 的次数为 \(11\)。
在 \(u=1\) 处有完全闭式
\[
\mathcal D(1,p)=4(p+1)^2(2p-1)^2(p^2-p+1)^2.
\]
因此在 \(p\in(0,1)\) 内，
\[
u=1\ \text{为谱分岔点候选}
\quad\Longleftrightarrow\quad
p=\frac12.
\]
\end{proposition}
\begin{proof}
对四次特征方程 \(\chi(\lambda;u,p)=0\) 直接计算 \(\Disc_\lambda(\chi)\) 并作整系数因式分解，即得
\(p^3(1-p)u\mathcal D(u,p)\) 的结构与 \(\deg_u\mathcal D=11\)。
再代入 \(u=1\) 并因式分解可得
\(\mathcal D(1,p)=4(p+1)^2(2p-1)^2(p^2-p+1)^2\)。
由于 \(p+1>0\) 且 \(p^2-p+1>0\) 在 \((0,1)\) 恒正，故仅当 \(2p-1=0\) 即 \(p=\tfrac12\) 时出现分岔。证毕。
\end{proof}
\leanverified{paper\_fold\_gauge\_anomaly\_bernoulli\_p\_leyang\_branch\_u1}


\begin{proposition}[输入--输出位对的平稳单点联合律（闭式表达）]\label{prop:fold-gauge-anomaly-bernoulli-p-bitpair-law}
\leanverified{paper\_fold\_gauge\_anomaly\_bernoulli\_p\_bitpair\_law}
在 Berstel--Zeckendorf 四状态图的边标记 \((X,Y)\in\{0,1\}^2\) 解释下，
令 \((X_0,Y_0)\) 表示平稳情形下某一典型位置的输入--输出比特对（输入为 Bernoulli\((p)\) 基线）。
则其联合分布为
\begin{align*}
\PP(X_0=0,Y_0=0)&=\frac{1 - p - p^2 + 2p^3 - p^4}{1+p^3},&
\PP(X_0=0,Y_0=1)&=\frac{p^2(1-p)}{1+p^3},\\
\PP(X_0=1,Y_0=0)&=\frac{p^2(2-p)}{1+p^3},&
\PP(X_0=1,Y_0=1)&=\frac{p\bigl(p^3+(1-p)^2\bigr)}{1+p^3}.
\end{align*}
特别地，输出比特密度为
\[
\PP(Y_0=1)=\frac{p(p^3-p+1)}{1+p^3},
\qquad
p-\PP(Y_0=1)=\frac{p^2}{1+p^3}>0.
\]
并且协方差闭式为
\[
\Cov(X_0,Y_0)
=-\frac{p(1-p)^2(p^2+p-1)}{1+p^3}.
\]
记 \(p_\phi:=\frac{\sqrt5-1}{2}\)，则
\[
\Cov(X_0,Y_0)
\begin{cases}
>0,& 0<p<p_\phi,\\
=0,& p=p_\phi,\\
<0,& p_\phi<p<1.
\end{cases}
\]
此外，失配密度满足 \(\PP(X_0\neq Y_0)=g_\ast(p)\)。
\end{proposition}
\begin{proof}[证明要点]
平稳一阶边分布由 \(\pi(p)\) 与转移概率给出；除状态 \(c\to b\) 上存在两条不同标记的平行边（分别对应 \((0,0)\) 与 \((1,1)\)）外，其余允许边的标记唯一。
对平行边按 Bernoulli\((p)\) 权重分解并逐边求和，即得所示四项闭式；对 \(Y_0=1\) 汇总并化简得到严格收缩恒等式
\(p-\PP(Y_0=1)=\frac{p^2}{1+p^3}\)；
再由 \(\Cov(X_0,Y_0)=\EE[X_0Y_0]-\EE[X_0]\EE[Y_0]\) 直接化简得到协方差闭式与符号阈值；最后将 \((0,1)\) 与 \((1,0)\) 两项相加即得 \(\PP(X_0\neq Y_0)=g_\ast(p)\)。证毕。
\end{proof}

\begin{corollary}[均匀基线作为特例：\texorpdfstring{$p=\tfrac12$}{p=1/2}]\label{cor:fold-gauge-anomaly-bernoulli-p-specialize-half}
\leanverified{paper\_fold\_gauge\_anomaly\_bernoulli\_p\_specialize\_half}
在 \(p=\tfrac12\) 时，
\[
g_\ast(p)=\frac49,\qquad \sigma_G^2(p)=\frac{118}{243},
\qquad
I_p(0)=\log\!\Bigl(\frac{2}{\varphi}\Bigr),\qquad I_p(1)=\log 2,
\]
并且命题 \ref{prop:fold-gauge-anomaly-bernoulli-p-bitpair-law} 的联合律退化为定理 \ref{thm:fold-gauge-anomaly-density} 所列的有理表。
\end{corollary}
\subparagraph{primorial 归一化下的 Gödel 外置素数寄存器：自平均、偏差与复杂度胀量}
在定理 \ref{thm:fold-gauge-anomaly-bernoulli-p-laws} 的平稳 Bernoulli\((p)\) 设定下，定义逐位失配指示
\[
M_k:=g(S_{k-1},S_k)\in\{0,1\},\qquad
G_m=\sum_{k=1}^{m}M_k.
\]
再令 \(\mathfrak p_k\) 为第 \(k\) 个素数，并记
\[
\mathfrak P_m:=\prod_{k=1}^{m}\mathfrak p_k,\qquad
A_m:=\log \mathfrak P_m=\sum_{k=1}^{m}\log \mathfrak p_k,
\]
\[
R_m:=\prod_{k=1}^{m}\mathfrak p_k^{M_k},\qquad
H_m:=\log R_m=\sum_{k=1}^{m}M_k\log \mathfrak p_k.
\]
下述结果把“投影失配外置到素数寄存器”的统计极限与复杂度代价转写为可检验闭式。

\begin{theorem}[primorial 归一化下的外置寄存器极限定律]\label{thm:fold-gauge-anomaly-bernoulli-p-primorial-godel-laws}
在上述设定下，令
\[
g_\ast(p)=\frac{p^2(3-2p)}{1+p^3},
\]
并令 \(\sigma_G^2(p)\) 为定理 \ref{thm:fold-gauge-anomaly-bernoulli-p-laws} 中的方差率闭式。
则有：
\begin{enumerate}
  \item \textbf{强自平均：}
  \[
  \frac{H_m}{A_m}\xrightarrow[m\to\infty]{a.s.}g_\ast(p).
  \]
  \item \textbf{中心极限定理（同方差率）：}
  \[
  \frac{H_m-g_\ast(p)A_m}{\sqrt m\,\log m}
  \ \Rightarrow\
  \mathcal N\!\bigl(0,\sigma_G^2(p)\bigr).
  \]
  \item \textbf{大偏差同率函数：}
  序列 \(\bigl(H_m/A_m\bigr)_{m\ge1}\) 在速度 \(m\) 下满足大偏差原理，其良率函数与
  \(\bigl(G_m/m\bigr)_{m\ge1}\) 相同，即命题
  \ref{prop:fold-gauge-anomaly-bernoulli-p-ldp} 的 \(I_p\)。
\end{enumerate}
\end{theorem}
\begin{proof}
\textbf{(1) 强自平均.}
记
\(
S_m:=\sum_{k=1}^{m}(M_k-g_\ast(p))
\)。
由定理 \ref{thm:fold-gauge-anomaly-bernoulli-p-laws}(1)，有 \(S_m=o(m)\) 几乎处处。
由 Abel 变换，
\[
H_m-g_\ast(p)A_m
=
S_m\log\mathfrak p_m
-\sum_{k=1}^{m-1}S_k\bigl(\log\mathfrak p_{k+1}-\log\mathfrak p_k\bigr).
\]
利用经典素数渐近 \(\log\mathfrak p_k=\log k+O(\log\log k)\)，得
\(
\log\mathfrak p_{k+1}-\log\mathfrak p_k=O(k^{-1})
\)、
\(
A_m\sim m\log m
\)。
于是第一项满足
\[
\frac{|S_m|\log\mathfrak p_m}{A_m}\to 0.
\]
第二项中，对任意 \(\varepsilon>0\)，取 \(K\) 使 \(|S_k|\le \varepsilon k\)（\(k\ge K\)）：
\[
\sum_{k=K}^{m-1}|S_k|\bigl|\log\mathfrak p_{k+1}-\log\mathfrak p_k\bigr|
\le
C\varepsilon\sum_{k=K}^{m-1}1
=
O(\varepsilon m),
\]
再除以 \(A_m\sim m\log m\) 得 \(O(\varepsilon/\log m)\to0\)；
有限前缀 \(k<K\) 贡献为 \(O(1/A_m)\to0\)。
故
\[
\frac{H_m-g_\ast(p)A_m}{A_m}\to0,
\]
即 \(H_m/A_m\to g_\ast(p)\) 几乎处处。
\par
\textbf{(2) 中心极限定理.}
分解
\[
H_m-g_\ast(p)A_m
=
(\log m)\bigl(G_m-mg_\ast(p)\bigr)
+\sum_{k=1}^{m}\bigl(M_k-g_\ast(p)\bigr)b_{m,k},
\]
其中
\(
b_{m,k}:=\log\mathfrak p_k-\log m
\)。
第一项由定理 \ref{thm:fold-gauge-anomaly-bernoulli-p-laws}(2) 直接给出
\[
\frac{(\log m)\bigl(G_m-mg_\ast(p)\bigr)}{\sqrt m\,\log m}
\Rightarrow
\mathcal N\!\bigl(0,\sigma_G^2(p)\bigr).
\]
对余项，令 \(\xi_k:=M_k-g_\ast(p)\)。
有限状态不可约非周期马尔可夫链上的有界可加观测给出协方差绝对可和，
故存在常数 \(C_0\) 使
\[
\operatorname{Var}\!\left(\sum_{k=1}^{m}\xi_k b_{m,k}\right)
\le
C_0\sum_{k=1}^{m}b_{m,k}^2.
\]
又由 \(\log\mathfrak p_k=\log k+O(\log\log k)\) 得
\[
b_{m,k}=\log(k/m)+O(\log\log m),
\]
从而
\[
\sum_{k=1}^{m}b_{m,k}^2=O\!\bigl(m(\log\log m)^2\bigr).
\]
于是
\[
\operatorname{Var}\!\left(
\frac{1}{\sqrt m\,\log m}\sum_{k=1}^{m}\xi_k b_{m,k}
\right)
=
O\!\left(\frac{(\log\log m)^2}{\log^2 m}\right)\to0,
\]
即该余项在 \(L^2\) 中收敛到 \(0\)。
由 Slutsky 定理得到所述极限正态律。
\par
\textbf{(3) 大偏差同率函数.}
有确定性估计
\[
\left|\frac{H_m}{A_m}-\frac{G_m}{m}\right|
\le
\frac{1}{A_m}\sum_{k=1}^{m}\left|\log\mathfrak p_k-\frac{A_m}{m}\right|.
\]
同样用 \(\log\mathfrak p_k=\log k+O(\log\log k)\) 与 \(A_m\sim m\log m\)，可得右端为
\(
O(1/\log m)\to0
\)。
因而对任意 \(\delta>0\)，充分大 \(m\) 时
\[
\PP\!\left(\left|\frac{H_m}{A_m}-\frac{G_m}{m}\right|>\delta\right)=0,
\]
即两序列在速度 \(m\) 下指数等价。
命题 \ref{prop:fold-gauge-anomaly-bernoulli-p-ldp} 已给出 \(G_m/m\) 的 LDP 与率函数 \(I_p\)；
指数等价保持率函数，故 \(H_m/A_m\) 共享同一 \(I_p\)。
证毕。
\end{proof}
\leanverified{paper\_fold\_gauge\_anomaly\_bernoulli\_p\_primorial\_godel\_laws}

\begin{proposition}[逐层赋值可审计强制互异素因子与 primorial 下界]\label{prop:fold-layerwise-auditable-primorial-necessity}
令 \(\Sigma\) 为有限字母表，\(T\ge 1\)。
设存在乘法外置编码
\[
\mathcal N(a_1\cdots a_T)=\prod_{t=1}^{T} q_t^{\,c_t(a_t)},
\qquad q_t\in\NN_{\ge 2},
\qquad c_t:\Sigma\hookrightarrow\{1,\dots,M\}\ \text{单射}.
\]
称其为\emph{逐层赋值可审计}，若对每个位置 \(t\) 存在素数 \(p_t\) 使对所有串 \(w=a_1\cdots a_T\) 都有
\[
v_{p_t}(\mathcal N(w))=c_t(a_t).
\]
则存在两两互异素数 \(p_1,\dots,p_T\) 满足：
\begin{enumerate}
  \item \(p_t\mid q_t\) 且 \(v_{p_t}(q_t)=1\)；
  \item 对 \(s\neq t\)，有 \(p_t\nmid q_s\)；
  \item 因而对所有 \(w\)，有 \(\mathcal N(w)\ge \prod_{t=1}^{T}p_t\)。
\end{enumerate}
特别地，在最小化最坏情形比特长度的意义下，最优选择为 \(q_t=p_t\) 且取前 \(T\) 个素数，从而强制出现 primorial 型增长。
\leanverified{paper\_fold\_layerwise\_auditable\_primorial\_necessity}
\end{proposition}

\begin{proof}
固定 \(t\)。
由赋值加法性与编码的乘法形式，对任意 \(w=a_1\cdots a_T\) 有
\[
v_{p_t}(\mathcal N(w))=\sum_{s=1}^{T} c_s(a_s)\,v_{p_t}(q_s).
\]
该值对所有 \(w\) 仅依赖 \(a_t\)，故若存在 \(s\neq t\) 使 \(v_{p_t}(q_s)>0\)，则改变 \(a_s\) 会改变右端，矛盾。
因此对 \(s\neq t\)，必有 \(v_{p_t}(q_s)=0\)，即 \(p_t\nmid q_s\)。
另一方面，若 \(v_{p_t}(q_t)=0\)，则上式右端与 \(a_t\) 无关，与可审计性矛盾；
故 \(v_{p_t}(q_t)\ge 1\)，即 \(p_t\mid q_t\)。
将 \(s\neq t\) 项已知为 \(0\) 代回上式可得
\(
v_{p_t}(\mathcal N(w))=c_t(a_t)\,v_{p_t}(q_t)
\)；
而该值对所有 \(a_t\) 恒等于 \(c_t(a_t)\)，故必有 \(v_{p_t}(q_t)=1\)。
由 (2) 立即推出 \(p_t\) 两两互异。
\par
最后，由 \(c_t(a_t)\ge 1\) 与 \(v_{p_t}(q_t)=1\) 得对所有 \(w\)，
\(
v_{p_t}(\mathcal N(w))\ge 1
\)，从而 \(p_t\mid \mathcal N(w)\)。
因此 \(\mathcal N(w)\ge \prod_{t=1}^{T}p_t\)。
在互异素数集合中，\(\prod_{t=1}^{T}p_t\) 的最小值由取前 \(T\) 个素数实现；同时 \(q_t\) 含额外素因子或更高幂只会增大最坏情形规模，故最优为 \(q_t=p_t\)。
\end{proof}

\begin{corollary}[有限字典 Gödel 码的比特长度为 \texorpdfstring{$\Theta(T\log T)$}{Theta(T log T)}]\label{cor:fold-godel-finite-dictionary-bitlength-theta-TlogT}
\leanverified{paper\_fold\_godel\_finite\_dictionary\_bitlength\_theta\_tlogt}
\leanverified{paper\_fold\_godel\_finite\_dictionary\_bitlength\_theta\_TlogT}
固定 \(B\in\NN_{\ge 1}\)，令字母表 \(A_B:=\{1,\dots,B\}\)。
对 \(T\ge 1\) 定义位址型 Gödel 编码
\[
G_{B,T}:A_B^T\longrightarrow \NN_{\ge 1},
\qquad
G_{B,T}(a_1,\dots,a_T):=\prod_{t=1}^{T}\mathfrak p_t^{a_t},
\]
其中 \(\mathfrak p_t\) 为第 \(t\) 个素数，并记 primorial \(\mathfrak P_T:=\prod_{t=1}^{T}\mathfrak p_t\)。
则对任意 \(a\in A_B^T\) 有
\[
\mathfrak P_T\le G_{B,T}(a)\le (\mathfrak P_T)^B.
\]
从而其二进制长度 \(L(a):=\lfloor \log_2 G_{B,T}(a)\rfloor+1\) 满足
\[
\log_2\mathfrak P_T\ \le\ L(a)\ \le\ B\log_2\mathfrak P_T+1.
\]
特别地，当 \(T\to\infty\) 时有 \(L(a)=\Theta(T\log T)\)，且常数因子落在区间 \([1,B]\)。
\end{corollary}
\begin{proof}
由 \(a_t\in\{1,\dots,B\}\) 直接得到
\(
G_{B,T}(a)\ge \prod_{t=1}^{T}\mathfrak p_t=\mathfrak P_T
\)
与
\(
G_{B,T}(a)\le \prod_{t=1}^{T}\mathfrak p_t^{B}=(\mathfrak P_T)^B
\)。
取 \(\log_2\) 即得长度夹逼。
最后由 primorial 渐近 \(\log\mathfrak P_T\sim T\log T\) 得 \(L(a)=\Theta(T\log T)\)。证毕。
\end{proof}

\begin{theorem}[有限字典 Gödel 像的 Dirichlet 分配函数与全纯收敛域]\label{thm:fold-godel-finite-dictionary-dirichlet-partition-holomorphic}
\leanverified{paper\_fold\_godel\_finite\_dictionary\_dirichlet\_partition\_holomorphic}
固定 \(B\in\NN_{\ge 1}\)，定义 Gödel 影像集合
\[
\mathcal G_B:=\bigcup_{T\ge 0}\left\{\ \prod_{t=1}^{T}\mathfrak p_t^{a_t}:\ a_t\in\{1,\dots,B\}\ \right\}\subset \NN_{\ge 1},
\]
其中 \(T=0\) 时乘积约定为 \(1\)。
定义其 Dirichlet 分配函数
\[
Z_B(s):=\sum_{n\in\mathcal G_B} n^{-s}\qquad (s\in\CC).
\]
则对所有满足 \(\Re(s)>0\) 的 \(s\)，级数 \(Z_B(s)\) 绝对收敛，并具有分层分解
\[
Z_B(s)=\sum_{T=0}^{\infty}\ \prod_{t=1}^{T}\left(\sum_{a=1}^{B}\mathfrak p_t^{-as}\right).
\]
因此 \(Z_B\) 在半平面 \(\Re(s)>0\) 上全纯。
\end{theorem}
\begin{proof}
设 \(\sigma=\Re(s)>0\)。
对每个 \(t\) 有
\[
\sum_{a=1}^{B}\left|\mathfrak p_t^{-as}\right|
=\sum_{a=1}^{B}\mathfrak p_t^{-a\sigma}
\le B\,\mathfrak p_t^{-\sigma}.
\]
从而对每个 \(T\ge 0\) 有
\[
\prod_{t=1}^{T}\left(\sum_{a=1}^{B}\mathfrak p_t^{-a\sigma}\right)
\le
B^T\prod_{t=1}^{T}\mathfrak p_t^{-\sigma}
=
B^T\,\mathfrak P_T^{-\sigma}.
\]
又由 \(\mathfrak p_t\ge t+1\) 得 \(\mathfrak P_T\ge (T+1)!\)，故
\[
B^T\,\mathfrak P_T^{-\sigma}\le B^T\,((T+1)!)^{-\sigma}.
\]
而 \(\sum_{T\ge 0} B^T/((T+1)!)^{\sigma}\) 由阶乘支配收敛，故 \(Z_B(s)\) 在 \(\Re(s)>0\) 绝对收敛。
对任意紧子集上一致收敛由同一估计给出，从而由 Weierstrass 判别法得 \(Z_B\) 在该半平面全纯。
分层分解由按长度 \(T\) 分组并使用有限乘积展开的交换性得到。证毕。
\end{proof}

\begin{theorem}[Gödel--\texorpdfstring{$\zeta$}{zeta} Gibbs 因子化：长度与位条件独立]\label{thm:fold-godel-finite-dictionary-gibbs-factorization}
取 \(\sigma>0\) 并令 \(\mathbf A_B^{\ast}:=\bigcup_{T\ge 0}A_B^{T}\) 为有限字典的自由幺半群。
按定理 \ref{thm:fold-godel-finite-dictionary-dirichlet-partition-holomorphic} 的 \(Z_B(\sigma)\) 定义概率测度
\[
\PP_{\sigma,B}(a_1,\dots,a_T)
:=
\frac{\left(\prod_{t=1}^{T}\mathfrak p_t^{a_t}\right)^{-\sigma}}{Z_B(\sigma)}
\qquad\bigl((a_1,\dots,a_T)\in A_B^T,\ T\ge 0\bigr),
\]
其中 \(T=0\) 对应空字。
记随机变量 \(T\) 为字长，并在条件 \(T=T_0\) 下记 \((A_t)_{t=1}^{T_0}\) 为位数字。
则：
\begin{enumerate}
  \item \(T\) 的分布为
  \[
  \PP_{\sigma,B}(T=T_0)=\frac{W_{T_0}(\sigma)}{Z_B(\sigma)},
  \qquad
  W_T(\sigma):=\prod_{t=1}^{T} w_t(\sigma),
  \qquad
  w_t(\sigma):=\sum_{a=1}^{B}\mathfrak p_t^{-a\sigma}.
  \]
  \item 在条件 \(T=T_0\) 下，\((A_t)_{t=1}^{T_0}\) 相互独立，且边际为
  \[
  \PP_{\sigma,B}(A_t=a\mid T=T_0)
  =
  \frac{\mathfrak p_t^{-a\sigma}}{\sum_{u=1}^{B}\mathfrak p_t^{-u\sigma}}
  \qquad (a\in\{1,\dots,B\}).
  \]
\end{enumerate}
\end{theorem}
\begin{proof}
对固定 \(T_0\) 有
\[
\sum_{(a_1,\dots,a_{T_0})\in A_B^{T_0}}\prod_{t=1}^{T_0}\mathfrak p_t^{-a_t\sigma}
\;=\;
\prod_{t=1}^{T_0}\left(\sum_{a=1}^{B}\mathfrak p_t^{-a\sigma}\right)
\;=\;
W_{T_0}(\sigma),
\]
故长度分布为 \(W_{T_0}(\sigma)/Z_B(\sigma)\)。
在条件 \(T=T_0\) 下，权重 \(\prod_{t=1}^{T_0}\mathfrak p_t^{-a_t\sigma}\) 对坐标 \(t\) 完全分离，归一化后即得到乘积分布与所示边际。证毕。
\end{proof}
\leanverified{paper\_fold\_godel\_finite\_dictionary\_gibbs\_factorization}

\begin{conjecture}[有限字典 Dirichlet 分配函数的自然边界与有效温度匹配]\label{conj:fold-godel-finite-dictionary-natural-boundary-temperature-match}
对任意 \(B\ge 2\)，猜想 \(Z_B(s)\) 在直线 \(\Re(s)=0\) 上存在自然边界：不存在穿越该直线的解析延拓。
此外，若某一可审计 prime-register 外置方案在宏观上产生证书整数族 \(N_m\) 并满足步数/密度律（例如定理 \ref{thm:fold-local-rewrite-ldp-lower-tail-barrier} 的线性下界口径），
则预期存在唯一 \(\sigma^\ast>0\) 使得经适当归一化后的 \(\log N_m\) 的偏差原理与 \(Z_B(\sigma)\) 的 Legendre 对偶在 \(\sigma=\sigma^\ast\) 处匹配；
等价地，可将 \(\sigma^\ast\) 视为使得外置证书的长度统计在主阶上逼近定理 \ref{thm:fold-godel-finite-dictionary-gibbs-factorization} 所诱导的 Gibbs 因子化机制的有效温度参数。
\end{conjecture}

\begin{theorem}[坐标可分 Gödel 整数化的不可压缩胀量]\label{thm:fold-gauge-anomaly-godel-coprime-inflation}
\leanverified{paper\_fold\_gauge\_anomaly\_godel\_coprime\_inflation}
令 \((q_k)_{k\ge1}\subset\NN_{\ge2}\) 两两互素。
对任意二进模式 \(b=(b_k)_{k\ge1}\in\{0,1\}^{\NN}\)，定义前缀 Gödel 整数
\[
N_m(b):=\prod_{k=1}^{m}q_k^{b_k},
\qquad
s_m:=\sum_{k=1}^{m}b_k.
\]
若 \(b\) 的 \(1\)-密度存在且为 \(\alpha\in(0,1)\)，即 \(s_m/m\to\alpha\)，则有
\[
\liminf_{m\to\infty}\frac{\log N_m(b)}{m\log m}\ge \alpha.
\]
此外，当 \(q_k=\mathfrak p_k\)（前缀素数寄存器）时，极限存在且
\[
\lim_{m\to\infty}\frac{\log N_m(b)}{m\log m}=\alpha.
\]
\end{theorem}
\begin{proof}
设
\(
I_m:=\{1\le k\le m:\ b_k=1\}
\)，则 \(|I_m|=s_m\)。
集合 \(\{q_k:\ k\in I_m\}\) 由 \(s_m\) 个两两互素整数构成，故其乘积不小于前 \(s_m\) 个素数的乘积：
\[
N_m(b)=\prod_{k\in I_m}q_k\ge \prod_{j=1}^{s_m}\mathfrak p_j=: \mathfrak P_{s_m}.
\]
于是
\[
\frac{\log N_m(b)}{m\log m}
\ge
\frac{\log \mathfrak P_{s_m}}{m\log m}.
\]
由 primorial 渐近 \(\log\mathfrak P_n\sim n\log n\)，再结合 \(s_m/m\to\alpha\)，得
\[
\liminf_{m\to\infty}\frac{\log N_m(b)}{m\log m}\ge \alpha.
\]
\par
若 \(q_k=\mathfrak p_k\)，则
\[
\log N_m(b)=\sum_{k=1}^{m}b_k\log\mathfrak p_k.
\]
记 \(B(x):=\sum_{k\le x}b_k=\alpha x+o(x)\)。
由部分求和与 \(\log\mathfrak p_k=\log k+O(\log\log k)\)：
\[
\sum_{k=1}^{m}b_k\log\mathfrak p_k
=
\sum_{k=1}^{m}b_k\log k
+O\!\left(\sum_{k=1}^{m}b_k\log\log k\right)
=
\alpha m\log m+o(m\log m).
\]
两边同除 \(m\log m\) 得极限为 \(\alpha\)。
证毕。
\end{proof}

\begin{corollary}[失配模式的 primorial 外置胀量]\label{cor:fold-gauge-anomaly-mismatch-primorial-growth}
\leanverified{paper\_fold\_gauge\_anomaly\_mismatch\_primorial\_growth}
取 \(b_k=M_k\)。
则在定理 \ref{thm:fold-gauge-anomaly-bernoulli-p-primorial-godel-laws} 的平稳设定下，几乎处处有
\[
\lim_{m\to\infty}\frac{\log R_m}{m\log m}=g_\ast(p).
\]
更一般地，任意坐标可分、两两互素的 Gödel 外置机制对同一失配模式都满足
\[
\liminf_{m\to\infty}\frac{\log N_m}{m\log m}\ge g_\ast(p)\qquad a.s.
\]
因此把长度 \(m\) 的失配模式外置为单整数时，位长主阶不可避免地带有 \(\log m\) 膨胀因子。
\end{corollary}
\begin{proof}
第一式由定理 \ref{thm:fold-gauge-anomaly-bernoulli-p-primorial-godel-laws}(1) 与
\(
A_m\sim m\log m
\) 直接推出。
第二式由定理 \ref{thm:fold-gauge-anomaly-godel-coprime-inflation} 代入
\(
\alpha=g_\ast(p)
\) 得到。证毕。
\end{proof}

\begin{corollary}[局部重写步数下界强制 \texorpdfstring{$m\log m$}{m log m} 级 prime-certificate 比特胀量]\label{cor:fold-local-rewrite-prime-certificate-bit-inflation}
在定理 \ref{thm:fold-local-rewrite-ldp-lower-tail-barrier} 的设定下，令 \(T_m(\mathsf S)\) 为局部重写正规化算法 \(\mathsf S\) 的步数。
设一类可审计 prime-register 证书把每一步编码为一位素数轴上的正指数，从而存在 \(a_t\ge 1\) 使证书整数可写为
\[
N_m=\prod_{t=1}^{T_m(\mathsf S)}\mathfrak p_t^{a_t}.
\]
则几乎处处成立下界
\[
\liminf_{m\to\infty}\frac{\log_2 N_m}{m\log_2 m}\ \ge\ \frac{4}{9R}.
\]
因此该类逐步可审计的整数证书在典型尺度上必为 \(\Omega(m\log m)\) 比特。
\end{corollary}
\begin{proof}
由 \(a_t\ge 1\) 得
\(
N_m\ge \prod_{t=1}^{T_m(\mathsf S)}\mathfrak p_t=\mathfrak P_{T_m(\mathsf S)}
\)，故 \(\log_2 N_m\ge \log_2\mathfrak P_{T_m(\mathsf S)}\)。
由 primorial 渐近 \(\log \mathfrak P_n\sim n\log n\) 得
\(
\log_2\mathfrak P_{T_m}\ge (1+o(1))\,T_m\log_2 T_m
\)。
另一方面，定理 \ref{thm:fold-local-rewrite-ldp-lower-tail-barrier}(2) 给出
\(
T_m(\mathsf S)\ge (\frac{4}{9R}+o(1))m
\)
几乎处处成立，从而 \(\log T_m\sim \log m\)。
合并即得
\(
\log_2 N_m\ge (\frac{4}{9R}+o(1))m\log_2 m
\)。证毕。
\end{proof}
\leanverified{paper\_fold\_local\_rewrite\_prime\_certificate\_bit\_inflation}

\begin{proposition}[算术可审计性与信息论旁信息下界之间的发散缺口]\label{prop:fold-prime-certificate-vs-it-gap-diverges}
令 \(\omega\sim\mathrm{Unif}(\Omega_m)\) 且令 \(x:=\Fold_m(\omega)\in X_m\) 为确定性折叠读出（其中 \(|\Omega_m|=2^m\)、\(|X_m|=F_{m+2}\)）。
则条件熵满足
\[
H(\omega\mid x)=H(\omega)-H(x)\le m\log 2-\log|X_m|=O(m).
\]
另一方面，在推论 \ref{cor:fold-local-rewrite-prime-certificate-bit-inflation} 的设定下，任一逐步可审计 prime-register 证书的比特长度为 \(\Omega(m\log m)\)。
因此其与信息论旁信息下界之比满足
\[
\frac{\mathrm{Len}_2(\mathrm{certificate})}{H(\omega\mid x)}\xrightarrow[m\to\infty]{}\infty,
\]
并且至少以 \(\Omega(\log m)\) 速度发散。
\end{proposition}
\begin{proof}
由 \(x\) 为 \(\omega\) 的确定函数，\(H(\omega\mid x)=H(\omega)-H(x)\)。
又 \(H(\omega)=m\log 2\) 且 \(H(x)\ge 0\)，并且 \(\log|X_m|=\Theta(m)\)（由 \(|X_m|=F_{m+2}\asymp\varphi^{m}\)），故 \(H(\omega\mid x)=O(m)\)。
另一方面，\(\Omega(m\log m)\) 的证书比特下界由推论 \ref{cor:fold-local-rewrite-prime-certificate-bit-inflation} 给出。
相除得到发散，并由主阶比较得到至少 \(\Omega(\log m)\) 的速度。证毕。
\end{proof}

\leanverified{paper\_fold\_prime\_certificate\_vs\_it\_gap\_diverges}
\leanverified{paper\_fold\_gauge\_anomaly\_godel\_ellipsoid\_volume}
\begin{proposition}[素数寄存器诱导的随机椭球体积律]\label{prop:fold-gauge-anomaly-godel-ellipsoid-volume}
对每个 \(m\)，定义线性算子
\[
D_m:=\mathrm{diag}\!\left(\mathfrak p_1^{M_1/2},\dots,\mathfrak p_m^{M_m/2}\right),
\]
并令
\[
B_m:=\{x\in\RR^m:\ |x|_2\le1\},\qquad E_m:=D_m(B_m).
\]
则体积比恒等式成立：
\[
\frac{\operatorname{Vol}(E_m)}{\operatorname{Vol}(B_m)}
=\det(D_m)
=\prod_{k=1}^{m}\mathfrak p_k^{M_k/2}
=R_m^{1/2}.
\]
因此在平稳 Bernoulli\((p)\) 设定下：
\begin{enumerate}
  \item \[
  \frac{\log\operatorname{Vol}(E_m)-\log\operatorname{Vol}(B_m)}{A_m}
  \xrightarrow[m\to\infty]{a.s.}
  \frac12 g_\ast(p).
  \]
  \item \[
  \frac{\log\operatorname{Vol}(E_m)-\log\operatorname{Vol}(B_m)-\frac12 g_\ast(p)A_m}
  {\sqrt m\,\log m}
  \Rightarrow
  \mathcal N\!\left(0,\frac14\sigma_G^2(p)\right).
  \]
  \item 变量
  \(
  \bigl(\log\operatorname{Vol}(E_m)-\log\operatorname{Vol}(B_m)\bigr)/A_m
  \)
  在速度 \(m\) 下满足 LDP，率函数为
  \[
  J_p(x)=I_p(2x),\qquad x\in[0,1/2].
  \]
\end{enumerate}
\end{proposition}
\begin{proof}
体积恒等式由线性代数公式
\(
\operatorname{Vol}(D(B_m))=\det(D)\operatorname{Vol}(B_m)
\)
直接给出。
其余三条分别由定理
\ref{thm:fold-gauge-anomaly-bernoulli-p-primorial-godel-laws}
在变换 \(x\mapsto x/2\) 下立即推出。证毕。
\end{proof}

\begin{theorem}[Gödel 素数外置的 Benford 极限定律]\label{thm:fold-godel-prime-register-benford}
\leanverified{paper\_fold\_godel\_prime\_register\_benford}
设整数底 \(b\ge 2\)，并取参数 \(q\in(0,1)\)。
令 \((B_k)_{k\ge 1}\) 为独立同分布的 Bernoulli\((q)\) 随机变量，并以第 \(k\) 个素数 \(\mathfrak p_k\) 定义 Gödel 化素数外置寄存器
\[
R_n:=\prod_{k=1}^{n}\mathfrak p_k^{B_k}\qquad(n\ge 1).
\]
记 \(\{x\}:=x-\lfloor x\rfloor\) 为小数部分。
则随机变量 \(\{\log_b R_n\}\) 在分布意义下收敛到 \([0,1)\) 上的均匀分布。
从而 \(R_n\) 在底 \(b\) 表示下的首位数字分布收敛到 Benford 律：对 \(d\in\{1,\dots,b-1\}\)，
\[
\lim_{n\to\infty}\PP(\text{首位数字}=d)=\log_b\!\Bigl(1+\frac{1}{d}\Bigr).
\]
\end{theorem}
\begin{proof}
令 \(\alpha_k:=\log_b\mathfrak p_k\)，则
\(
\log_b R_n=\sum_{k=1}^{n}B_k\alpha_k
\)。
由分布型 Weyl 判别准则，只需证明对任意非零整数 \(\ell\)，其 Fourier 系数
\[
\phi_n(\ell):=\EE\bigl[\exp(2\pi i\ell \log_b R_n)\bigr]
\]
满足 \(\phi_n(\ell)\to 0\)。
由独立性，
\[
\phi_n(\ell)
=\prod_{k=1}^{n}\EE\bigl[\exp(2\pi i\ell B_k\alpha_k)\bigr]
=\prod_{k=1}^{n}\Bigl((1-q)+q\,\exp(2\pi i\ell\alpha_k)\Bigr).
\]
对任意实数 \(\theta\) 有恒等式
\[
\bigl|(1-q)+q\,\exp(2\pi i\theta)\bigr|^2
=1-4q(1-q)\sin^2(\pi\theta),
\]
再用不等式 \(\sqrt{1-x}\le e^{-x/2}\)（\(x\in[0,1]\)）得
\[
\bigl|(1-q)+q\,\exp(2\pi i\theta)\bigr|
\le
\exp\!\bigl(-2q(1-q)\sin^2(\pi\theta)\bigr).
\]
于是
\[
|\phi_n(\ell)|
\le
\exp\!\Bigl(-2q(1-q)\sum_{k=1}^{n}\sin^2(\pi\ell\alpha_k)\Bigr).
\]
因此只需证级数 \(\sum_{k\ge 1}\sin^2(\pi\ell\log_b \mathfrak p_k)\) 发散。
\par
取 \(\eta>0\) 足够小使 \(0<|\ell|\log_b(1+2\eta)<1\)，并记
\[
\delta_1:=|\ell|\log_b(1+\eta),\qquad
\delta_2:=|\ell|\log_b(1+2\eta),
\]
则 \(0<\delta_1<\delta_2<1\)。
连续函数 \(t\mapsto \sin^2(\pi t)\) 在紧区间 \([\delta_1,\delta_2]\) 上取到正最小值
\[
c:=\min_{t\in[\delta_1,\delta_2]}\sin^2(\pi t)>0.
\]
对任意充分大的整数 \(m\)，考虑区间
\[
I_m:=[\,b^{m/|\ell|}(1+\eta),\ b^{m/|\ell|}(1+2\eta)\,].
\]
其长度为 \(\eta\,b^{m/|\ell|}\)，由素数定理可知 \(I_m\) 中的素数个数随 \(m\to\infty\) 发散；故对无穷多个 \(m\)，可取 \(\mathfrak p\in I_m\) 为素数。
写 \(\mathfrak p=b^{m/|\ell|}u\) 其中 \(u\in(1+\eta,1+2\eta)\)，则
\[
|\ell|\log_b \mathfrak p=m+|\ell|\log_b u,
\]
其小数部分落在 \([\delta_1,\delta_2]\)，从而
\(
\sin^2(\pi\ell\log_b \mathfrak p)\ge c
\)。
因此该发散级数有无穷多个项 \(\ge c\)，故必发散。
于是 \(|\phi_n(\ell)|\to 0\)，从而 \(\{\log_b R_n\}\Rightarrow \mathrm{Unif}[0,1)\)。
Benford 首位数字分布是均匀小数部分的直接推论。证毕。
\end{proof}

\endinput

\subparagraph{独立 Bernoulli 输入下的 Gödel 密度估计量：确定性近似、Cramér 偏差与迭代对数律}
上一段已在 primorial 归一化口径下给出 Gödel 外置寄存器的强自平均、中心极限定理与大偏差主链。若把驱动序列进一步退化为彼此独立的 Bernoulli 位本身，则 prime-register 加权密度估计量还会呈现一条更强的保持性：它不仅在强大数律、中心极限与大偏差层面与普通经验均值同型，而且在确定性逼近、最优阈值检验与迭代对数律层面也保持同一主阶结构。

在以下各条结果中，令 \(\mathfrak p_j\) 为第 \(j\) 个素数，并对二元样本 \(x=(x_1,\dots,x_n)\in\{0,1\}^n\) 定义
\[
G_n(x):=\prod_{j=1}^{n}\mathfrak p_j^{x_j},
\qquad
\Theta_n:=\sum_{j=1}^{n}\log \mathfrak p_j,
\qquad
\widehat d_n(x):=\frac{\log G_n(x)}{\Theta_n}.
\]
若 \((X_j)_{j\ge 1}\) 为随机样本，则记
\[
S_n:=\log G_n=\sum_{j=1}^{n}X_j\log \mathfrak p_j,
\qquad
\widehat p_n:=\frac{S_n}{\Theta_n}.
\]

\begin{theorem}[Gödel 密度估计量与经验均值的确定性逼近]\label{thm:fold-godel-density-estimator-deterministic-mean-approx}
\leanverified{paper\_fold\_godel\_density\_estimator\_deterministic\_mean\_approx}
对任意 \(x=(x_1,\dots,x_n)\in\{0,1\}^n\)，定义经验均值
\[
\overline x_n:=\frac1n\sum_{j=1}^{n}x_j.
\]
则存在绝对常数 \(C>0\)，使对充分大的 \(n\) 成立一致估计
\[
\sup_{x\in\{0,1\}^n}\bigl|\widehat d_n(x)-\overline x_n\bigr|
\le
C\,\frac{\log\log n}{\log n}.
\]
特别地，Gödel 密度估计量与普通经验均值在确定性意义下只相差 \(o(1)\)。
\end{theorem}
\begin{proof}
记
\[
w_{n,j}:=\frac{\log \mathfrak p_j}{\Theta_n},
\qquad 1\le j\le n.
\]
则 \(\sum_{j=1}^{n}w_{n,j}=1\)，并且
\[
\widehat d_n(x)-\overline x_n
=
\sum_{j=1}^{n}\left(w_{n,j}-\frac1n\right)x_j.
\]
因此
\[
\sup_{x\in\{0,1\}^n}\bigl|\widehat d_n(x)-\overline x_n\bigr|
\le
\sum_{j=1}^{n}\left|w_{n,j}-\frac1n\right|
=
\frac{1}{\Theta_n}\sum_{j=1}^{n}\left|\log \mathfrak p_j-\frac{\Theta_n}{n}\right|.
\]

由素数渐近
\[
\mathfrak p_j=j\bigl(\log j+O(\log\log j)\bigr),
\]
可得对所有 \(1\le j\le n\) 一致成立
\[
\log \mathfrak p_j=\log j+O(\log\log n).
\]
另一方面，
\[
\Theta_n=\sum_{j=1}^{n}\log \mathfrak p_j
=
n\log n+O(n\log\log n),
\]
故
\[
\frac{\Theta_n}{n}=\log n+O(\log\log n).
\]
于是
\[
\sum_{j=1}^{n}\left|\log \mathfrak p_j-\frac{\Theta_n}{n}\right|
\le
\sum_{j=1}^{n}\bigl|\log j-\log n\bigr|+O(n\log\log n).
\]
由于 \(j\le n\) 时 \(|\log j-\log n|=\log(n/j)\)，故
\[
\sum_{j=1}^{n}\bigl|\log j-\log n\bigr|
=
\sum_{j=1}^{n}\log\frac{n}{j}
=
n\log n-\log(n!)
=
O(n)
\]
（由 Stirling 公式）。再结合 \(\Theta_n\asymp n\log n\)，得到
\[
\sup_{x\in\{0,1\}^n}\bigl|\widehat d_n(x)-\overline x_n\bigr|
=
O\!\left(\frac{\log\log n}{\log n}\right).
\]
证毕。
\end{proof}

\begin{theorem}[Gödel 密度估计量的渐近正态性与一阶效率]\label{thm:fold-godel-density-estimator-clt-efficient}
\leanverified{paper\_fold\_godel\_density\_estimator\_clt\_efficient}
设 \(X_1,X_2,\dots\) 独立同分布，且
\[
X_j\sim \mathrm{Bernoulli}(p),
\qquad 0<p<1.
\]
则
\[
\sqrt n\,(\widehat p_n-p)\ \Rightarrow\ \mathcal N\!\bigl(0,\ p(1-p)\bigr).
\]
因此，尽管 \(\widehat p_n\) 使用的是素数对数加权而非等权平均，它与普通样本均值具有完全相同的一阶渐近方差。
\end{theorem}
\begin{proof}
令
\[
Y_j:=(X_j-p)\log \mathfrak p_j,
\qquad
T_n:=\sum_{j=1}^{n}Y_j=S_n-p\Theta_n.
\]
则 \((Y_j)_{j\ge 1}\) 相互独立、均值为 \(0\)，并且
\[
\Var(Y_j)=p(1-p)(\log \mathfrak p_j)^2.
\]
记
\[
s_n^2:=\sum_{j=1}^{n}\Var(Y_j)=p(1-p)\sum_{j=1}^{n}(\log \mathfrak p_j)^2.
\]

由
\[
\log \mathfrak p_j=\log j+O(\log\log n)
\qquad (1\le j\le n)
\]
一致成立，可得
\[
\Theta_n=\sum_{j=1}^{n}\log \mathfrak p_j
=
n\log n+O(n\log\log n),
\]
以及
\[
\sum_{j=1}^{n}(\log \mathfrak p_j)^2
=
n(\log n)^2+O\bigl(n\log n\log\log n\bigr).
\]
故
\[
s_n^2\sim p(1-p)\,n(\log n)^2.
\]

又因为
\[
\max_{1\le j\le n}|Y_j|
\le
\log \mathfrak p_n
=
O(\log n)
=
o(s_n),
\]
Lindeberg 条件自动成立，从而由 Lindeberg--Feller 中心极限定理，
\[
\frac{T_n}{s_n}\Rightarrow \mathcal N(0,1).
\]

另一方面，
\[
\sqrt n\,(\widehat p_n-p)
=
\frac{\sqrt n\,T_n}{\Theta_n}
=
\frac{\sqrt n\,s_n}{\Theta_n}\cdot \frac{T_n}{s_n},
\]
而
\[
\frac{\sqrt n\,s_n}{\Theta_n}
\to
\sqrt{p(1-p)}.
\]
由 Slutsky 定理即得
\[
\sqrt n\,(\widehat p_n-p)\Rightarrow \mathcal N\!\bigl(0,p(1-p)\bigr).
\]
证毕。
\end{proof}

\begin{theorem}[Gödel 密度估计量的 Cramér 型大偏差原理]\label{thm:fold-godel-density-estimator-cramer-ldp}
在定理 \ref{thm:fold-godel-density-estimator-clt-efficient} 的设定下，\((\widehat p_n)_{n\ge 1}\) 以速度 \(n\) 满足大偏差原理，其好速率函数为
\[
I_p(x)=
\begin{cases}
x\log\frac{x}{p}+(1-x)\log\frac{1-x}{1-p},& x\in[0,1],\\[2mm]
+\infty,& x\notin[0,1].
\end{cases}
\]
\end{theorem}
\begin{proof}
定义对数矩母函数
\[
\Lambda_n(t):=\frac1n\log \EE\exp\bigl(nt\widehat p_n\bigr).
\]
由独立性，
\[
\Lambda_n(t)
=
\frac1n\sum_{j=1}^{n}
\log\Bigl(1-p+p\,e^{t a_{n,j}}\Bigr),
\qquad
a_{n,j}:=\frac{n\log \mathfrak p_j}{\Theta_n}.
\]

由定理 \ref{thm:fold-godel-density-estimator-deterministic-mean-approx} 的证明可知
\[
\frac1n\sum_{j=1}^{n}|a_{n,j}-1|
=
\sum_{j=1}^{n}\left|\frac{\log \mathfrak p_j}{\Theta_n}-\frac1n\right|
=
O\!\left(\frac{\log\log n}{\log n}\right)\to 0.
\]
另一方面，
\[
a_{n,j}\le \frac{n\log \mathfrak p_n}{\Theta_n}=O(1),
\]
故对每个固定 \(t\in\RR\)，函数
\[
g_t(a):=\log\bigl(1-p+p\,e^{ta}\bigr)
\]
在包含全部 \(a_{n,j}\) 的某个紧区间上 Lipschitz。于是
\[
\left|\Lambda_n(t)-g_t(1)\right|
\le
\frac{L_t}{n}\sum_{j=1}^{n}|a_{n,j}-1|
\to 0,
\]
从而
\[
\Lambda_n(t)\to \Lambda(t):=\log(1-p+pe^t).
\]

又因为 \(\widehat p_n\in[0,1]\) 几乎处处成立，故指数紧性自动满足。函数 \(\Lambda\) 在 \(\RR\) 上有限且 \(C^\infty\)，因而可由 Gärtner--Ellis 定理推出 \((\widehat p_n)\) 在速度 \(n\) 下满足大偏差原理，其速率函数为 Legendre--Fenchel 变换
\[
I_p(x)=\sup_{t\in\RR}\{tx-\Lambda(t)\}.
\]
对 Bernoulli 对数矩母函数 \(\Lambda(t)=\log(1-p+pe^t)\) 直接计算其 Legendre 变换，即得上述显式二元相对熵公式。证毕。
\end{proof}
\leanverified{paper\_fold\_godel\_density\_estimator\_cramer\_ldp}

\begin{corollary}[基于 Gödel 密度估计量的最优阈值检验与平衡 Chernoff 指数]\label{cor:fold-godel-density-estimator-optimal-threshold-test}
固定 \(0<p_0<p_1<1\)。对阈值 \(\eta\in(p_0,p_1)\) 定义检验
\[
\phi_{n,\eta}:=\mathbf 1_{\{\widehat p_n\ge \eta\}}.
\]
则类型 I、II 误差满足
\[
\lim_{n\to\infty}\frac1n\log \PP_{p_0}(\phi_{n,\eta}=1)
=
-I_{p_0}(\eta),
\]
\[
\lim_{n\to\infty}\frac1n\log \PP_{p_1}(\phi_{n,\eta}=0)
=
-I_{p_1}(\eta).
\]
并且存在唯一阈值
\[
\eta^\star=
\frac{\log\!\bigl(\frac{1-p_0}{1-p_1}\bigr)}
{\log\!\bigl(\frac{p_1(1-p_0)}{p_0(1-p_1)}\bigr)}
\in(p_0,p_1)
\]
使得
\[
I_{p_0}(\eta^\star)=I_{p_1}(\eta^\star).
\]
因此，\(\eta^\star\) 在阈值检验类中给出唯一的平衡点，而共同值
\[
I_{p_0}(\eta^\star)=I_{p_1}(\eta^\star)
\]
可视为该阈值类的平衡 Chernoff 指数。
\end{corollary}
\begin{proof}
由定理 \ref{thm:fold-godel-density-estimator-cramer-ldp}，并注意到 \(I_{p_0}\) 在 \([p_0,1]\) 上严格递增，而 \(I_{p_1}\) 在 \([0,p_1]\) 上严格递减，可得
\[
\lim_{n\to\infty}\frac1n\log \PP_{p_0}(\widehat p_n\ge \eta)
=
-I_{p_0}(\eta),
\]
\[
\lim_{n\to\infty}\frac1n\log \PP_{p_1}(\widehat p_n\le \eta)
=
-I_{p_1}(\eta),
\]
即所示两条误差指数。

平衡条件为
\[
I_{p_0}(\eta)=I_{p_1}(\eta),
\]
等价于
\[
\eta\log\frac{p_1}{p_0}
=(1-\eta)\log\frac{1-p_0}{1-p_1}.
\]
直接解得
\[
\eta^\star=
\frac{\log\!\bigl(\frac{1-p_0}{1-p_1}\bigr)}
{\log\!\bigl(\frac{1-p_0}{1-p_1}\bigr)+\log\!\bigl(\frac{p_1}{p_0}\bigr)}
=
\frac{\log\!\bigl(\frac{1-p_0}{1-p_1}\bigr)}
{\log\!\bigl(\frac{p_1(1-p_0)}{p_0(1-p_1)}\bigr)}.
\]
又
\[
\frac{\mathrm d}{\mathrm d\eta}\bigl(I_{p_0}(\eta)-I_{p_1}(\eta)\bigr)
=
\log\!\bigl(\tfrac{p_1(1-p_0)}{p_0(1-p_1)}\bigr)>0,
\]
故该平衡阈值唯一。证毕。
\end{proof}
\leanverified{paper\_fold\_godel\_density\_estimator\_optimal\_threshold\_test}

\begin{theorem}[Gödel 密度估计量的 Hartman--Wintner 型迭代对数律]\label{thm:fold-godel-density-estimator-lil}
在定理 \ref{thm:fold-godel-density-estimator-clt-efficient} 的设定下，几乎处处成立
\[
\limsup_{n\to\infty}
\sqrt{\frac{n}{2p(1-p)\log\log n}}\,
(\widehat p_n-p)=1,
\]
\[
\liminf_{n\to\infty}
\sqrt{\frac{n}{2p(1-p)\log\log n}}\,
(\widehat p_n-p)=-1.
\]
\leanverified{paper\_fold\_godel\_density\_estimator\_lil}
\end{theorem}
\begin{proof}
仍记
\[
Y_j:=(X_j-p)\log \mathfrak p_j,
\qquad
T_n:=\sum_{j=1}^{n}Y_j=S_n-p\Theta_n,
\]
以及
\[
s_n^2:=\sum_{j=1}^{n}\Var(Y_j)
=
p(1-p)\sum_{j=1}^{n}(\log \mathfrak p_j)^2.
\]
由定理 \ref{thm:fold-godel-density-estimator-clt-efficient} 的证明，
\[
s_n^2\sim p(1-p)\,n(\log n)^2.
\]
并且
\[
\max_{1\le j\le n}|Y_j|
\le
\log \mathfrak p_n
\sim
\log n.
\]
因此
\[
\frac{\max_{1\le j\le n}|Y_j|}
{s_n/\sqrt{\log\log s_n^2}}
\ll
\sqrt{\frac{\log\log n}{n}}
\to 0.
\]
由独立非同分布中心变量的 Hartman--Wintner--Kolmogorov 迭代对数律，可得
\[
\limsup_{n\to\infty}
\frac{T_n}{\sqrt{2s_n^2\log\log s_n^2}}=1,
\qquad
\liminf_{n\to\infty}
\frac{T_n}{\sqrt{2s_n^2\log\log s_n^2}}=-1
\quad a.s.
\]

另一方面，
\[
\widehat p_n-p=\frac{T_n}{\Theta_n},
\qquad
\Theta_n\sim n\log n,
\qquad
s_n^2\sim p(1-p)\,n(\log n)^2,
\]
并且 \(\log\log s_n^2\sim \log\log n\)。故
\[
\frac{\sqrt{2s_n^2\log\log s_n^2}}{\Theta_n}
\sim
\sqrt{\frac{2p(1-p)\log\log n}{n}}.
\]
将上述等价关系代回，即得所示两条迭代对数律。证毕。
\end{proof}

\endinput

% (User input window)

\paragraph{均匀基线下失配指示过程的同步收缩与无限记忆}\label{par:fold-gauge-anomaly-sofic-memory}
在定理 \ref{thm:fold-gauge-anomaly-bernoulli-p-laws} 的 Markov 表示下取 \(p=\tfrac12\)。
沿用边势 \(g\)（式 \eqref{eq:fold-gauge-anomaly-bernoulli-p-edge-potential}），并把它视作二元观测过程
\[
g_t:=g(S_{t-1},S_t)\in\{0,1\},\qquad t\in\ZZ,
\]
其中 \((S_t)_{t\in\ZZ}\) 为该平稳链。
记 \(0^n,1^n\) 为长度 \(n\) 的常字串，并沿用 \S\ref{subsec:folding-fibonacci-stable-syntax} 的 Fibonacci 数列 \((F_k)\) 与黄金比例 \(\varphi\)。

\par
为固定记号，回到命题 \ref{prop:fold-gauge-anomaly-pressure} 的加权邻接矩阵表示：在均匀输入下
\[
g_t=\mathbf 1_{\{X_t\neq Y_t\}}\in\{0,1\},
\qquad
A_\theta=\frac12\begin{pmatrix}
1 & e^\theta & 0 & 1\\
0 & 0 & e^\theta & 0\\
e^\theta & 2 & 0 & 0\\
1 & 0 & 0 & 0
\end{pmatrix}
\]
（按状态次序 \((a,b,c,d)\)），其中恰有三条带因子 \(e^\theta\) 的边对应失配边。
在 \(\theta=0\) 处取 Perron--Frobenius/Parry 归一化得到基准动力学核，该核在文中用于相关闭式与有限长度方差的完全闭式化。

\begin{proposition}[Parry 转移核与失配边标记]\label{prop:fold-gauge-anomaly-parry-kernel-mismatch-label}
\leanverified{paper\_fold\_gauge\_anomaly\_parry\_kernel\_mismatch\_label}
令 \(M:=2A_0\)。则 \(\rho(M)=2\)，并且 \(v=(2,1,2,1)^{\mathsf T}\) 为一组右 Perron 向量。
定义 Parry 转移核
\[
P_{ij}:=\frac{M_{ij}v_j}{2v_i}\quad (i,j\in\{a,b,c,d\}),
\]
则
\[
P=
\begin{pmatrix}
\frac12&\frac14&0&\frac14\\[2pt]
0&0&1&0\\[2pt]
\frac12&\frac12&0&0\\[2pt]
1&0&0&0
\end{pmatrix}.
\]
并在允许边上定义输出字母
\[
g(i\to j)=1\iff (i,j)\in\{(a,b),(b,c),(c,a)\},
\]
其余允许边取 \(g=0\)。
令 \(g_t=g(S_{t-1}\to S_t)\) 为边标记输出，则该输出过程与均匀基线下的逐位失配指示一致。
\end{proposition}
\begin{proof}
直接计算得 \(Mv=2v\)，从而 \(\rho(M)=2\) 且 \(v\) 为右 Perron 向量；代入 Parry 公式得到所示 \(P\)。
最后的边标记恰与 \(A_\theta\) 中出现 \(e^\theta\) 的三处位置一致，故与失配指示同一。
证毕。
\end{proof}

\begin{theorem}[长度 \(3\) 同步词与可数 unifilar \texorpdfstring{$\varepsilon$}{ε}-机表示]\label{thm:fold-gauge-anomaly-epsilon-machine-synchronizing-word}
\leanverified{paper\_fold\_gauge\_anomaly\_epsilon\_machine\_synchronizing\_word}
在命题 \ref{prop:fold-gauge-anomaly-parry-kernel-mismatch-label} 的核 \(P\) 与边标记 \(g\) 下，单词 \(001\) 为同步词：若 \(g_{t-2}g_{t-1}g_t=001\)，则隐藏状态必为 \(S_t=b\)，且与更早历史无关。
\par
因此 \((g_t)\) 具有一个可数且 unifilar 的预测态表示（\(\varepsilon\)-机覆盖）。
其状态集可取为
\[
\mathcal S=\{B,C,A\}\cup\{R_n:n\in\NN_0\},
\]
其中 \(B,C,A\) 分别对应后验确定的隐藏状态 \(b,c,a\)，而 \(R_n\) 对应后验支撑于 \(\{a,d\}\) 的不确定预测态。
其转移（箭头标记输出符号）为
\[
B \xrightarrow{1} C,\qquad
C \xrightarrow{0} B,\qquad C \xrightarrow{1} A,
\]
\[
A \xrightarrow{1} B,\qquad A \xrightarrow{0} R_0,
\qquad
R_n \xrightarrow{1} B,\qquad R_n \xrightarrow{0} R_{n+1}\qquad(n\ge 0).
\]
\leanverified{paper\_fold\_epsilon\_machine\_synchronizing\_word\_seeds}
\leanverified{paper\_fold\_epsilon\_machine\_synchronizing\_word}
\end{theorem}
\begin{proof}
设 \(g_{t-2}g_{t-1}g_t=001\)。
由 \(g_{t-2}=0\) 与 \(P\) 的结构知不存在任何标记为 \(0\) 且终点为 \(c\) 的边，故 \(S_{t-2}\neq c\)；同理 \(S_{t-1}\neq c\)。
而 \(g_t=1\) 强制 \((S_{t-1},S_t)\in\{(a,b),(b,c),(c,a)\}\)，结合 \(S_{t-1}\neq c\) 排除 \((c,a)\)，仅余 \((a,b)\) 与 \((b,c)\)。
若 \((S_{t-1},S_t)=(b,c)\)，则 \(S_{t-1}=b\)；
但 \(g_{t-1}=0\) 时到达 \(b\) 的唯一 \(0\)-边为 \(c\to b\)，从而 \(S_{t-2}=c\)，矛盾。
故必为 \((S_{t-1},S_t)=(a,b)\)，从而 \(S_t=b\)，同步性成立。
\par
其余转移由 \(P\) 与边标记直接读出：
\(b\to c\) 为确定一步且输出 \(1\)，得 \(B\xrightarrow{1}C\)；
在 \(c\) 处，\(c\to b\) 与 \(c\to a\) 分别输出 \(0\) 与 \(1\) 且各以概率 \(1/2\)，得 \(C\xrightarrow{0}B\)、\(C\xrightarrow{1}A\)；
在 \(a\) 处，\(a\to b\) 输出 \(1\)，而 \(a\to a\)、\(a\to d\) 输出 \(0\)，得 \(A\xrightarrow{1}B\)、\(A\xrightarrow{0}R_0\)。
一旦后验支撑落在 \(\{a,d\}\)，输出 \(1\) 仅可能来自 \(a\to b\)，故观测到 \(1\) 即同步回 \(b\)，给出 \(R_n\xrightarrow{1}B\)；
观测到 \(0\) 则保持在 \(\{a,d\}\) 并更新不确定度，给出 \(R_n\xrightarrow{0}R_{n+1}\)。
证毕。
\end{proof}

\begin{theorem}[不确定预测态的 Fibonacci--M\"obius 完全闭式]\label{thm:fold-gauge-anomaly-epsilon-machine-fibonacci-mobius}
记 Fibonacci 数列 \(F_0=0,F_1=1,F_{n+2}=F_{n+1}+F_n\)。
在定理 \ref{thm:fold-gauge-anomaly-epsilon-machine-synchronizing-word} 的 \(\varepsilon\)-机中，存在显式参数 \(r_n>0\) 使得
\[
R_n:\quad \PP(S_t=d\mid R_n)=\frac{r_n}{1+r_n},\qquad
\PP(S_t=a\mid R_n)=\frac{1}{1+r_n},
\]
且
\[
r_n=\frac{F_{n+1}}{2F_{n+2}}\qquad(n\ge 0).
\]
从而在 \(R_n\) 处下一步输出为 \(1\) 的条件概率为
\[
\alpha_n:=\PP(g_{t+1}=1\mid R_n)=\frac{1}{4(1+r_n)}=\frac{F_{n+2}}{2F_{n+4}}.
\]
此外，\((r_n)\) 满足 M\"obius 迭代
\[
r_{n+1}=\frac{1}{4r_n+2},
\]
并以速度 \(\varphi^{-2n}\) 交错收敛到 \(r_\infty=\frac1{2\varphi}\)。
\leanverified{paper\_fold\_epsilon\_machine\_fibonacci\_mobius\_seeds}
\leanverified{paper\_fold\_epsilon\_machine\_fibonacci\_mobius}
\end{theorem}
\begin{proof}
由 \(A\xrightarrow{0}R_0\) 的 Bayes 更新：在 \(a\) 处条件于输出 \(0\)，两条 \(0\)-边 \(a\to a\)、\(a\to d\) 的相对概率为 \((1/2):(1/4)=2:1\)，故 \(r_0=\PP(d)/\PP(a)=1/2\)。
\par
若在 \(R_n\) 处再次观测到 \(0\)，则下一状态只可能为 \(a\) 或 \(d\)。
未归一化权重为：到 \(d\) 仅由 \(a\to d\)，权重 \(\frac{1}{1+r_n}\cdot\frac14\)；到 \(a\) 由 \(a\to a\) 与 \(d\to a\)，权重 \(\frac{1}{1+r_n}\cdot\frac12+\frac{r_n}{1+r_n}\cdot 1\)。
因此
\[
r_{n+1}
=\frac{\frac{1}{1+r_n}\cdot \frac14}{\frac{1}{1+r_n}\cdot\frac12+\frac{r_n}{1+r_n}\cdot 1}
=\frac{1}{4r_n+2}.
\]
归纳验证 Fibonacci 闭式：若 \(r_n=\frac{F_{n+1}}{2F_{n+2}}\)，则
\[
r_{n+1}=\frac{1}{4r_n+2}
=\frac{1}{2\frac{F_{n+1}}{F_{n+2}}+2}
=\frac{F_{n+2}}{2(F_{n+1}+F_{n+2})}
=\frac{F_{n+2}}{2F_{n+3}}.
\]
从而对一切 \(n\ge 0\) 成立。
由“输出 \(1\) 仅来自 \(a\to b\)”与 \(P_{ab}=1/4\) 得 \(\alpha_n=\frac{1}{4(1+r_n)}\)；
代入 \(r_n\) 并用恒等式 \(F_{n+4}=2F_{n+2}+F_{n+1}\) 得 \(\alpha_n=\frac{F_{n+2}}{2F_{n+4}}\)。
收敛性由 \(F_{n+1}/F_{n+2}\to\varphi^{-1}\) 立得。
证毕。
\end{proof}

\begin{theorem}[\texorpdfstring{$\varepsilon$}{ε}-机平稳分布的 Fibonacci 尾与完全归一化]\label{thm:fold-gauge-anomaly-epsilon-machine-stationary-fibonacci-tail}
该 \(\varepsilon\)-机存在唯一平稳分布 \(\pi\)，并满足
\[
\pi(B)=\pi(C)=\frac{2}{9},\qquad \pi(A)=\frac{1}{9},
\]
以及对一切 \(n\ge 0\),
\[
\pi(R_n)=\frac{F_{n+4}}{36\cdot 2^n}.
\]
特别地，\(\sum_{n\ge 0}\pi(R_n)=\frac49\)。
\leanverified{paper\_fold\_epsilon\_machine\_stationary\_fibonacci\_tail}
\leanverified{paper\_fold\_epsilon\_machine\_stationary\_fibonacci\_tail\_seeds}
\end{theorem}
\begin{proof}
由 \(B\xrightarrow{1}C\) 的确定性得 \(\pi(C)=\pi(B)\)，而进入 \(A\) 的唯一入流边为 \(C\xrightarrow{1}A\) 且概率 \(1/2\)，故 \(\pi(A)=\frac12\pi(C)=\frac12\pi(B)\)。
对 \(R_0\)，唯一入流边为 \(A\xrightarrow{0}R_0\) 且概率 \(3/4\)，故 \(\pi(R_0)=\frac34\pi(A)=\frac{3}{8}\pi(B)\)。
又由 \(R_n\xrightarrow{0}R_{n+1}\) 得 \(\pi(R_{n+1})=\pi(R_n)(1-\alpha_n)\)。
由定理 \ref{thm:fold-gauge-anomaly-epsilon-machine-fibonacci-mobius}，
\(
1-\alpha_n=\frac{F_{n+5}}{2F_{n+4}}
\)，从而
\[
\pi(R_n)=\pi(R_0)\prod_{k=0}^{n-1}\frac{F_{k+5}}{2F_{k+4}}
=\pi(R_0)\cdot\frac{F_{n+4}}{3\cdot 2^n}.
\]
归一化由生成函数恒等式 \(\sum_{n\ge0}F_{n+4}2^{-n}=16\) 给出：
\[
\sum_{n\ge0}\pi(R_n)=\pi(R_0)\cdot\frac13\sum_{n\ge0}\frac{F_{n+4}}{2^n}
=\pi(R_0)\cdot\frac{16}{3}
=2\pi(B).
\]
因此 \(1=\pi(B)+\pi(C)+\pi(A)+\sum_{n\ge0}\pi(R_n)=\frac92\pi(B)\)，从而 \(\pi(B)=2/9\)，代回即得全部闭式。
证毕。
\end{proof}

\begin{corollary}[不确定相位零串长度的精确分布、母函数与矩]\label{cor:fold-gauge-anomaly-epsilon-machine-zero-run-length}
条件于进入 \(R_0\)，记从该时刻起直至下一次输出 \(1\) 为止所观测到的连续 \(0\) 的总长度为 \(L\in\{1,2,\dots\}\)。
则对一切 \(n\ge0\),
\[
\PP(L=n+1)=\frac{F_{n+2}}{6\cdot 2^n},\qquad
\PP(L\ge n+1)=\frac{F_{n+4}}{3\cdot 2^n}.
\]
其概率母函数为
\[
\EE[z^L]=\frac{z(2+z)}{12-6z-3z^2},
\]
并且
\[
\EE[L]=\frac{16}{3},\qquad \mathrm{Var}(L)=\frac{200}{9}.
\]
\leanverified{paper\_fold\_gauge\_anomaly\_epsilon\_machine\_zero\_run\_length}
\end{corollary}
\begin{proof}
令 \(T:=\min\{n\ge0:\text{在 }R_n\text{ 处输出 }1\}\)，则 \(L=T+1\)。
由乘法公式
\[
\PP(T=n)=\alpha_n\prod_{k=0}^{n-1}(1-\alpha_k),
\qquad
\PP(T\ge n)=\prod_{k=0}^{n-1}(1-\alpha_k),
\]
并用 \(\alpha_n=\frac{F_{n+2}}{2F_{n+4}}\) 及望远镜积
\(
\prod_{k=0}^{n-1}(1-\alpha_k)=\frac{F_{n+4}}{3\cdot 2^n}
\)
得到两条分布闭式。
母函数由恒等式 \(\sum_{n\ge0}F_{n+2}x^n=\frac{1+x}{1-x-x^2}\) 在 \(x=z/2\) 处代入得到。
矩可由母函数求导或由尾和公式 \(\EE[L]=\sum_{n\ge0}\PP(L\ge n+1)\) 得到；二阶矩同理由母函数二阶导数闭式化得到 \(\mathrm{Var}(L)=200/9\)。
证毕。
\end{proof}

\begin{theorem}[Markov 阶无限与 \texorpdfstring{$\varepsilon$}{ε}-机可数无限]\label{thm:fold-gauge-anomaly-epsilon-machine-infinite-markov-order}
输出过程 \((g_t)\) 不是任何有限阶 \(k\) 的 Markov 链；等价地，其最小 Markov 阶为 \(\infty\)。
并且其 \(\varepsilon\)-机具有可数无限个状态 \(\{R_n\}_{n\ge0}\)（因此严格超出任意有限状态 Markov 近似）。
\leanverified{paper\_fold\_epsilon\_machine\_infinite\_markov\_order\_seeds}
\leanverified{paper\_fold\_gauge\_anomaly\_epsilon\_machine\_infinite\_markov\_order}
\end{theorem}
\begin{proof}
固定任意 \(k\in\NN\)，考虑两段具正概率的历史前缀
\[
h_k:=00111\,0^{k+1},\qquad h_k':=00111\,0^{k+2}.
\]
由定理 \ref{thm:fold-gauge-anomaly-epsilon-machine-synchronizing-word}，\(001\) 同步到 \(B\)，随后必经 \(B\xrightarrow{1}C\)，并以正概率 \(C\xrightarrow{1}A\)；
故前缀 \(00111\) 保证处于 \(A\)，继而 \(0^{k+1}\) 与 \(0^{k+2}\) 分别把预测态推进到 \(R_k\) 与 \(R_{k+1}\)。
因此
\[
\PP(g_{t+1}=1\mid \mathrm{past}=h_k)=\alpha_k,\qquad
\PP(g_{t+1}=1\mid \mathrm{past}=h_k')=\alpha_{k+1},
\]
其中 \(\alpha_n=\frac{F_{n+2}}{2F_{n+4}}\)（定理 \ref{thm:fold-gauge-anomaly-epsilon-machine-fibonacci-mobius}），故两者不相等。
另一方面，\(h_k\) 与 \(h_k'\) 的最后 \(k\) 个输出符号均为 \(0^k\)。
若 \((g_t)\) 为 \(k\) 阶 Markov，则上述两条件概率应相等，矛盾。
故 Markov 阶为 \(\infty\)。
\par
此外，\(\{R_n\}\) 给出不同的一步预测分布 \(\PP(g_{t+1}=1\mid R_n)=\alpha_n\)，且 \(\alpha_n\) 随 \(n\) 变化严格互异，因此这些预测态不可合并，故 \(\varepsilon\)-机状态数可数无限。
证毕。
\end{proof}

\begin{proposition}[熵率与统计复杂度的显式级数]\label{prop:fold-gauge-anomaly-epsilon-machine-entropy-rate-complexity}
\leanverified{paper\_fold\_gauge\_anomaly\_epsilon\_machine\_entropy\_rate\_complexity}
令二元熵函数 \(H_2(p):=-p\log_2 p-(1-p)\log_2(1-p)\)。
则熵率存在并由 \(\varepsilon\)-机平稳分布给出绝对收敛级数：
\[
h_\mu
=\frac{2}{9}\,H_2\!\Bigl(\frac12\Bigr)
\;+\;\frac{1}{9}\,H_2\!\Bigl(\frac14\Bigr)
\;+\;\sum_{n\ge0}\frac{F_{n+4}}{36\cdot 2^n}\,
H_2\!\Bigl(\frac{F_{n+2}}{2F_{n+4}}\Bigr),
\]
其中第一项来自状态 \(C\)、第二项来自状态 \(A\)，无穷和来自 \(\{R_n\}_{n\ge0}\)，而状态 \(B\) 因确定输出 \(1\) 不贡献熵。
该级数以比率 \((\varphi/2)^n\) 指数收敛，数值上
\[
h_\mu\approx 0.6215495141\ \text{bits/step}.
\]
此外，统计复杂度
\(
C_\mu:=-\sum_{s\in\mathcal S}\pi(s)\log_2\pi(s)
\)
有限，并由 \(\pi(B)=\pi(C)=2/9\)、\(\pi(A)=1/9\)、\(\pi(R_n)=F_{n+4}/(36\cdot2^n)\) 给出显式收敛级数。
\end{proposition}
\begin{proof}
由于 \(\varepsilon\)-机为 unifilar，熵率满足
\(
h_\mu=\sum_{s\in\mathcal S}\pi(s)\,H_2(\PP(g_{t+1}=1\mid s))
\)。
由定理 \ref{thm:fold-gauge-anomaly-epsilon-machine-synchronizing-word} 的转移结构与定理 \ref{thm:fold-gauge-anomaly-epsilon-machine-fibonacci-mobius} 的 \(\alpha_n\) 得到所示表达式。
收敛性由 \(F_{n+4}\sim \varphi^n/\sqrt5\) 推出 \(\pi(R_n)=O((\varphi/2)^n)\) 并结合 \(H_2\) 有界得到；\(C_\mu\) 的有限性同理由 \(\pi(R_n)=O((\varphi/2)^n)\) 与 \(|\log\pi(R_n)|=O(n)\) 得出。
证毕。
\end{proof}

\begin{theorem}[零段条件律的 Fibonacci 闭式与同步收缩]\label{thm:fold-gauge-anomaly-zero-run-fibonacci}
在上述均匀基线下，对任意整数 \(n\ge 2\) 成立
\[
\PP\!\bigl(g_{t+1}=1\mid g_{t-n+1}\cdots g_t=0^n\bigr)
=\frac{F_{n+3}}{2F_{n+5}}.
\]
并且该条件概率收敛到极限
\[
\lim_{n\to\infty}\PP(g_{t+1}=1\mid 0^n)=\frac{1}{2\varphi^2}=\frac{3-\sqrt5}{4},
\]
且振荡误差满足上界
\[
\Bigl|\PP(g_{t+1}=1\mid 0^n)-\frac{3-\sqrt5}{4}\Bigr|
<\frac{1}{2F_{n+5}F_{n+6}}
\qquad(n\ge 2).
\]
\leanverified{paper\_fold\_gauge\_anomaly\_zero\_run\_fibonacci}
\leanverified{paper\_fold\_gauge\_anomaly\_zero\_run\_fibonacci\_seeds}
\end{theorem}
\begin{proof}
在 \(p=\tfrac12\) 时，由 \eqref{eq:fold-gauge-anomaly-bernoulli-p-transition} 得到四态链的转移矩阵
\[
P=
\begin{pmatrix}
\frac12&\frac14&0&\frac14\\
0&0&1&0\\
\frac12&\frac12&0&0\\
1&0&0&0
\end{pmatrix},
\qquad
\pi=\Bigl(\frac49,\frac29,\frac29,\frac19\Bigr),
\]
其中 \(\pi\) 为平稳分布（定理 \ref{thm:fold-gauge-anomaly-bernoulli-p-laws} 的闭式在 \(p=\tfrac12\) 的特化）。
\par
由 \eqref{eq:fold-gauge-anomaly-bernoulli-p-edge-potential}，事件 \(g_t=0\) 等价于一步转移不在三条失配边
\(
a\to b,\ b\to c,\ c\to a
\)
之中；因此在该链上 \(g=0\) 的允许边恰为
\[
a\to a,\ a\to d,\ c\to b,\ d\to a.
\]
由此可见：若 \(g_{t-1}g_t=00\)，则第二步必须从 \(S_{t-1}\) 走一条 \(0\)-边，而状态 \(b\) 无 \(0\)-出边，故 \(S_{t-1}\neq b\)；
又因 \(c\) 无 \(0\)-入边（即上式无任何 \(0\)-边以 \(c\) 为终点），结合 \(g_t=0\) 可得
\[
\PP(S_t\in\{a,d\}\mid g_{t-1}g_t=00)=1.
\]
从而对一切 \(n\ge 2\)，在条件 \(g_{t-n+1}\cdots g_t=0^n\) 下亦有 \(\PP(S_t\in\{a,d\}\mid 0^n)=1\)。
\par
记
\[
p_n:=\PP(S_t=a\mid g_{t-n+1}\cdots g_t=0^n)\qquad(n\ge 2).
\]
由于在集合 \(\{a,d\}\) 上仅 \(a\) 能以概率 \(P_{ab}=\tfrac14\) 产生下一位 \(g_{t+1}=1\)，从而
\[
\PP(g_{t+1}=1\mid 0^n)=\frac{p_n}{4}.
\]
另一方面，在后验分布 \((p,1-p)\)（对应 \(\{a,d\}\)）下观测到下一位为 \(0\) 的 Bayes 更新为
\[
p\longmapsto f(p)
:=\frac{p\cdot P(a\to a)+(1-p)\cdot P(d\to a)}{p\cdot \PP(g_{t+1}=0\mid S_t=a)+(1-p)\cdot \PP(g_{t+1}=0\mid S_t=d)}
=\frac{p\cdot\frac12+(1-p)\cdot 1}{p\cdot\frac34+(1-p)\cdot 1}
=\frac{4-2p}{4-p},
\]
从而 \(p_{n+1}=f(p_n)\)（\(n\ge 2\)）。
由平稳分布与两步滤波直接计算得到初值 \(p_2=\frac{10}{13}\)。
\par
断言
\[
p_n=\frac{2F_{n+3}}{F_{n+5}}\qquad(n\ge 2)
\]
可由归纳法验证：基例 \(n=2\) 成立；
若对某个 \(n\ge 2\) 成立，则
\[
p_{n+1}=f(p_n)=\frac{4-2p_n}{4-p_n}
=\frac{4-\frac{4F_{n+3}}{F_{n+5}}}{4-\frac{2F_{n+3}}{F_{n+5}}}
=\frac{2F_{n+4}}{F_{n+6}}
\]
（其中用到 \(F_{n+5}-F_{n+3}=F_{n+4}\)、\(2F_{n+4}+F_{n+3}=F_{n+6}\)），得到归纳步。
于是
\[
\PP(g_{t+1}=1\mid 0^n)=\frac{p_n}{4}=\frac{F_{n+3}}{2F_{n+5}}.
\]
极限由 \(F_{n+3}/F_{n+5}\to\varphi^{-2}\) 给出。
误差界来自 Fibonacci 连分数收敛子的标准估计：\(\frac{F_{n+3}}{F_{n+5}}\) 为 \(\varphi^{-2}\) 的交错有界逼近且满足
\(
\bigl|\frac{F_{n+3}}{F_{n+5}}-\varphi^{-2}\bigr|<\frac{1}{F_{n+5}F_{n+6}}
\)；
两边乘以 \(\tfrac12\) 即得所示上界。证毕。
\end{proof}

\begin{theorem}[有限型支撑下的无限 Markov 阶：可计算的不可截断记忆]\label{thm:fold-gauge-anomaly-strictly-sofic-infinite-memory}
在定理 \ref{thm:fold-gauge-anomaly-zero-run-fibonacci} 的设定下，对任意整数 \(n\ge 2\) 成立
\[
\PP\!\bigl(g_{t+1}=1\mid g_{t-n}\cdots g_t=10^n\bigr)=\frac{F_{n+1}}{2F_{n+3}}.
\]
并且对任意 \(k\ge 1\)，过程 \((g_t)_{t\in\ZZ}\) 不可能是 \(k\) 阶 Markov 链。
更强地，对任意 \(n\ge \max\{2,k\}\)，两段共享同一长度 \(k\) 后缀的历史 \(0^n\) 与 \(10^n\) 给出严格不同的下一位条件律，且差值具有闭式：
\[
\PP(g_{t+1}=1\mid 10^n)-\PP(g_{t+1}=1\mid 0^n)
=\frac{(-1)^n}{2F_{n+3}F_{n+5}}\neq 0.
\]
\end{theorem}
\begin{proof}
对 \(n\ge 2\)，历史 \(10^n\) 仍以 \(00\) 结尾，因而同样将 \(S_t\) 同步到 \(\{a,d\}\)，并在零段内部满足同一 Möbius 更新 \(p\mapsto f(p)=\frac{4-2p}{4-p}\)。
令
\[
p^{(1)}_n:=\PP(S_t=a\mid g_{t-n}\cdots g_t=10^n)\qquad(n\ge 2),
\]
则对 \(n\ge 2\) 有 \(\PP(g_{t+1}=1\mid 10^n)=p^{(1)}_n/4\) 且 \(p^{(1)}_{n+1}=f(p^{(1)}_n)\)。
只需确定初值：当 \(n=2\) 时 \(10^n=100\)，直接由平稳链计算得 \(\PP(g_{t+1}=1\mid 100)=\frac15\)，从而 \(p^{(1)}_2=\frac45\)。
\par
由归纳法可得
\[
p^{(1)}_n=\frac{2F_{n+1}}{F_{n+3}}\qquad(n\ge 2),
\]
因而
\(
\PP(g_{t+1}=1\mid 10^n)=\frac{F_{n+1}}{2F_{n+3}}
\)。
\par
其次，Fibonacci 的广义 Cassini 恒等式给出
\[
F_{n+1}F_{n+5}-F_{n+3}^2 = (-1)^n,
\]
从而差值闭式成立：
\[
\frac{F_{n+1}}{2F_{n+3}}-\frac{F_{n+3}}{2F_{n+5}}
=\frac{F_{n+1}F_{n+5}-F_{n+3}^2}{2F_{n+3}F_{n+5}}
=\frac{(-1)^n}{2F_{n+3}F_{n+5}}.
\]
最后，若 \((g_t)\) 为某个有限阶 \(k\) 的 Markov 链，则任意两段历史只要共享长度 \(k\) 的末尾块，其下一位条件分布必相同。
取任意 \(n\ge k\)，则 \(0^n\) 与 \(10^n\) 共享后缀 \(0^k\)，但其下一位取值 \(1\) 的条件概率差值严格非零，矛盾。
因此不存在有限 \(k\) 使 \((g_t)\) 成为 \(k\) 阶 Markov。
由于 \((g_t)\) 为有限状态链在边标号下的输出过程，其支撑由定理 \ref{thm:fold-gauge-anomaly-g-support-sofic} 给出为有限型子移位，上述结论可视为“有限型支撑 + 无限记忆核”的可计算证书。证毕。
\end{proof}

\leanverified{paper\_fold\_gauge\_anomaly\_strictly\_sofic\_infinite\_memory}

\begin{theorem}[规范异常场支撑的有限型刻画与拓扑不变量]\label{thm:fold-gauge-anomaly-g-support-sofic}
设 \((X_t,Y_t)_{t\in\ZZ}\) 为定理 \ref{thm:fold-gauge-anomaly-bernoulli-p-laws} 在 \(p=\tfrac12\) 下的四状态联合过程，并记
\[
g_t:=\mathbf 1_{\{X_t\neq Y_t\}}\in\{0,1\}.
\]
在边势定义 \eqref{eq:fold-gauge-anomaly-bernoulli-p-edge-potential} 的标记图口径下，\(g_t=1\) 恰对应三条边
\(
a\to b,\ b\to c,\ c\to a
\)，其余可达边均标记为 \(0\)。

定义
\[
\Sigma_g:=\Bigl\{\mathbf g\in\{0,1\}^{\ZZ}:\ 
\text{任意长度 \(4\) 子块均不等于 }0010\text{ 或 }0100\Bigr\}.
\]
则 \(g\) 的拓扑支撑满足
\[
\operatorname{supp}(g)=\Sigma_g.
\]
并且
\[
h_{\mathrm{top}}(\Sigma_g)=\log\lambda,
\quad
\lambda>1\ \text{为}\ x^5-x^4-x^3-x-1=0\ \text{的 Perron 实根},
\]
其 Artin--Mazur \(\zeta\) 函数为
\[
\zeta_{\Sigma_g}(z)=\frac{1}{\det(I-zA)}
=\frac{1}{(1-z)(1-z-z^2-z^4-z^5)},
\]
其中 \(A\) 为以 \(\{0,1\}^3\) 为顶点、由允许长度 \(4\) 块诱导的三步 de Bruijn 呈示邻接矩阵。
\leanverified{paper\_fold\_gauge\_anomaly\_g\_support\_sofic}
\end{theorem}
\begin{proof}
\textbf{必要性。}
由边标记结构，\(1\) 仅可出现在三条三环边上；
并且从 \(b\) 的唯一出边为 \(b\to c\)（标记 \(1\)），从 \(d\) 的唯一出边为 \(d\to a\)（标记 \(0\)）。
对长度 \(4\) 局部块做穷尽可达性分析可得：模式 \(0010\) 与 \(0100\) 均与上述局部转移约束矛盾，故不可实现。
于是 \(\operatorname{supp}(g)\subseteq\Sigma_g\)。

\textbf{充分性。}
对四状态标记图施行标准确定化（子集构造）并去除瞬态状态，得到一个 \(8\) 状态再返呈示；
该呈示与长度 \(3\) 二元块集合 \(\{0,1\}^3\) 一一对应，且允许长度 \(4\) 块恰为
\[
\{0,1\}^4\setminus\{0010,0100\}.
\]
因此任意不含禁块 \(0010,0100\) 的双边序列都可在该确定型呈示上实现。
由确定化前后语言等价性，原四状态标记图亦可实现该序列，故
\(\Sigma_g\subseteq\operatorname{supp}(g)\)。
综上 \(\operatorname{supp}(g)=\Sigma_g\)。

\textbf{熵与 \(\zeta\) 函数。}
由三步 de Bruijn 呈示邻接矩阵 \(A\) 的特征多项式分解
\[
\chi_A(x)=x^2(x-1)\bigl(x^5-x^4-x^3-x-1\bigr),
\]
其 Perron 根即上述 \(\lambda\)，故
\(
h_{\mathrm{top}}(\Sigma_g)=\log\lambda
\)。
对 SFT 有标准公式
\(
\zeta_{\Sigma_g}(z)=1/\det(I-zA)
\)，直接计算得
\[
\det(I-zA)=(1-z)(1-z-z^2-z^4-z^5).
\]
命题成立。
\end{proof}

\begin{theorem}[再生分解与周期内异常数的显式母函数]\label{thm:fold-gauge-anomaly-regeneration-pgf}
在同一四状态链中，状态 \(d\) 满足 \(d\to a\) 为确定转移，且进入 \(d\) 的唯一方式为 \(a\to d\)（该边标记 \(0\)）。
令 \((\tau_k)_{k\ge0}\) 为 \(d\) 的相继到达时刻，并定义块
\[
B_k:=(g_{\tau_k},g_{\tau_k+1},\dots,g_{\tau_{k+1}-1}).
\]
则 \((B_k)_{k\ge0}\) 构成独立同分布再生块。
记
\[
L:=\tau_{k+1}-\tau_k,\qquad
S:=\sum_{t=\tau_k}^{\tau_{k+1}-1}g_t.
\]
则二元概率母函数为
\[
\EE[u^S z^L]
=\frac{u z^{4}-2 z^{2}}{u^{3}z^{3}-2u z^{3}+4u z^{2}+4z-8},
\qquad (|z|\ \text{充分小}).
\]
特别地，
\[
\EE[u^S]=\frac{u-2}{u^3+2u-4},
\]
并且
\[
\PP(S=1)=\PP(S=2)=0,
\]
且对 \(k\ge3\) 存在整数序列 \((a_k)\) 使
\[
\PP(S=k)=\frac{a_k}{2^k},\qquad
a_3=a_4=a_5=1,\quad
a_k=a_{k-1}+2a_{k-3}\ (k\ge6).
\]
\end{theorem}
\leanverified{paper\_fold\_gauge\_anomaly\_regeneration\_pgf}
\begin{proof}
\textbf{再生性。}
由 \(d\to a\) 的确定性与 \(a\to d\) 的唯一入射结构，\(d\) 为再生状态；
Markov 性立即给出 \((B_k)\) 独立同分布。

\textbf{母函数。}
从 \(d\) 出发第一步必为 \(d\to a\)，贡献一因子 \(z\)（长度加 \(1\)，异常计数加 \(0\)）。
令 \(F_s(u,z)\) 表示从 \(s\in\{a,b,c\}\) 出发直到首次到达 \(d\)（含吸收步）的 \(\EE[u^S z^L]\)。
由边标记与转移概率，得到线性系统
\[
\begin{aligned}
F_a&=\tfrac14 z+\tfrac12 z F_a+\tfrac14 z\,u\,F_b,\\
F_b&=z\,u\,F_c,\\
F_c&=\tfrac12 z\,u\,F_a+\tfrac12 z\,F_b.
\end{aligned}
\]
解得
\[
F_a(u,z)=\frac{u z^3-2z}{u^3 z^3-2u z^3+4u z^2+4z-8}.
\]
故
\[
\EE[u^S z^L]=zF_a(u,z)
=\frac{u z^{4}-2 z^{2}}{u^{3}z^{3}-2u z^{3}+4u z^{2}+4z-8}.
\]
取 \(z=1\) 得 \(\EE[u^S]=(u-2)/(u^3+2u-4)\)。

\textbf{系数递推。}
写
\(
\EE[u^S]=\sum_{k\ge0}p_k u^k
\)。
由
\[
(u^3+2u-4)\sum_{k\ge0}p_k u^k=u-2
\]
比较系数可得
\[
p_0=\tfrac12,\qquad p_1=p_2=0,\qquad
p_k=\tfrac12 p_{k-1}+\tfrac14 p_{k-3}\ (k\ge3).
\]
令 \(p_k=a_k/2^k\)（\(k\ge3\)），即得
\[
a_3=a_4=a_5=1,\qquad a_k=a_{k-1}+2a_{k-3}\ (k\ge6).
\]
命题成立。
\end{proof}

\begin{proposition}[柱事件 \(1^n\) 的三周期精确律]\label{prop:fold-gauge-anomaly-one-block-three-period-law}
\leanverified{paper\_fold\_gauge\_anomaly\_one\_block\_three\_period\_law}
在同一设定下，记柱事件概率
\[
p_n:=\PP(g_{t+1}\cdots g_{t+n}=1^n)\qquad(n\ge 1).
\]
则对任意整数 \(m\ge 0\) 成立
\[
p_{3m+1}=\frac{4}{9\cdot 8^{m}},\qquad
p_{3m+2}=\frac{1}{4\cdot 8^{m}},\qquad
p_{3m}=\frac{1}{9\cdot 8^{m-1}}\ (m\ge 1).
\]
因此下一位条件概率呈严格三周期：
\[
\PP(g_{t+1}=1\mid 1^{3m})=\frac12,\quad
\PP(g_{t+1}=1\mid 1^{3m+1})=\frac{9}{16},\quad
\PP(g_{t+1}=1\mid 1^{3m+2})=\frac{4}{9}.
\]
\end{proposition}
\begin{proof}
在 \eqref{eq:fold-gauge-anomaly-bernoulli-p-edge-potential} 的标号下，标号为 \(1\) 的边强迫隐藏态沿三环
\(
a\to b\to c\to a
\)
运动；并且对应转移概率分别为 \(\frac14,1,\frac12\)，三步一圈的乘积为 \(\frac18\)。
因此对任意 \(m\ge 0\)，在事件 \(1^{3m}\) 下隐藏态回到同一顶点并额外带来权重因子 \(8^{-m}\)。
再由平稳分布 \(\pi_a+\pi_b+\pi_c=\frac89\) 与
\(
p_1=\PP(g_{t+1}=1)=g_\ast=\frac49
\)
及
\(
p_2=\PP(g_{t+1}g_{t+2}=11)=\frac14
\)
（直接由 \(\pi\) 与 \(P\) 计算），即可得到三条闭式；条件概率由比值 \(p_{n+1}/p_n\) 立得。证毕。
\end{proof}

\begin{conjecture}[金属平均 Ostrowski 家族中的 continuant 支配律]\label{conj:fold-gauge-anomaly-metallic-continuant-law}
设 \(\alpha_A=[0;A,A,A,\dots]\) 为常数型连分数（\(A\in\NN_{\ge 1}\)），并假设相应的“取规范形”过程可由某个同步、确定性的有限状态转导器实现（Fibonacci 情形可参见 \cite{ITA_2001__35_6_491_0,LabbeLepsova2023FibonacciTwosComplement} 的 Berstel 加法器与不歧义换能器接口）。
令 \((g^{(A)}_t)_{t\in\ZZ}\) 为相应的规范差异（mismatch）指示过程。
\par
记 continuant（广义 Fibonacci）序列 \((K^{(A)}_n)_{n\ge 0}\) 由
\[
K^{(A)}_0=0,\quad K^{(A)}_1=1,\quad K^{(A)}_{n+2}=A K^{(A)}_{n+1}+K^{(A)}_{n}
\]
定义。
则存在有限阈值 \(n_0(A)\) 使得对一切 \(n\ge n_0(A)\) 成立闭式条件律
\[
\PP\!\bigl(g^{(A)}_{t+1}=1\mid 0^n\bigr)=\frac{K^{(A)}_{n+3}}{2K^{(A)}_{n+5}}.
\]
并且对任意 \(k\ge 1\)，过程 \((g^{(A)}_t)\) 不可能是 \(k\) 阶 Markov；其不可截断记忆差值应由 \(K^{(A)}\) 的广义 Cassini 型行列式恒等式给出显式闭式（在 \(A=1\) 时退化为定理 \ref{thm:fold-gauge-anomaly-strictly-sofic-infinite-memory} 的 Fibonacci 公式）。
\end{conjecture}

\endinput
\paragraph{有限上下文对数损失的不可消去缺口：由 \texorpdfstring{$\varepsilon$}{ε}-机尾部给出的显式下界}\label{par:fold-gauge-anomaly-sofic-memory-logloss}
沿用 \S\ref{par:fold-gauge-anomaly-sofic-memory} 的记号：\((g_t)\) 为均匀基线下的失配指示过程，其 \(\varepsilon\)-机状态集为 \(\mathcal S=\{B,C,A\}\cup\{R_n:n\in\NN_0\}\)，并满足
\[
\alpha_n:=\PP(g_{t+1}=1\mid R_n)=\frac{F_{n+2}}{2F_{n+4}},
\qquad
\pi(R_n)=\frac{F_{n+4}}{36\cdot 2^n},
\]
以及对 \(k\ge 2\),
\[
q_k:=\PP(g_{t+1}=1\mid g_{t-k+1}\cdots g_t=0^k)=\frac{F_{k+3}}{2F_{k+5}}
\]
（分别见定理 \ref{thm:fold-gauge-anomaly-epsilon-machine-fibonacci-mobius}、定理 \ref{thm:fold-gauge-anomaly-epsilon-machine-stationary-fibonacci-tail} 与定理 \ref{thm:fold-gauge-anomaly-zero-run-fibonacci}）。

\begin{theorem}[\texorpdfstring{$k$}{k}-上下文预测的对数损失缺口下界（Fibonacci 显式）]\label{thm:fold-gauge-anomaly-logloss-context-gap}
\leanverified{paper\_fold\_gauge\_anomaly\_logloss\_context\_gap}
令 \(k\ge 2\)，并令 \(\mathcal Q_k\) 为所有 \(k\)-上下文概率预测器
\[
\mathcal Q_k:=\Bigl\{q:\{0,1\}^k\to(0,1)\Bigr\}.
\]
对任意 \(q\in\mathcal Q_k\)，定义其单步期望对数损失（自然对数）为
\[
\mathcal L(q):=
\EE\Bigl[-\log q(g_{t-k+1}\cdots g_t)\cdot \mathbf 1_{\{g_{t+1}=1\}}
-\log\bigl(1-q(g_{t-k+1}\cdots g_t)\bigr)\cdot \mathbf 1_{\{g_{t+1}=0\}}\Bigr].
\]
令
\(
\mathcal L_k:=\inf_{q\in\mathcal Q_k}\mathcal L(q)
\)
并令 \(\mathcal L_\infty\) 表示允许使用全历史的 Bayes 最优对数损失（等于熵率，见命题 \ref{prop:fold-gauge-anomaly-epsilon-machine-entropy-rate-complexity} 的 unifilar 表达）。
则成立显式下界
\[
\boxed{\;
\mathcal L_k-\mathcal L_\infty
\ \ge\
\pi(R_{k-1})\cdot
D_{\mathrm{KL}}\!\Bigl(\mathrm{Bern}(\alpha_{k-1})\ \big\|\ \mathrm{Bern}(q_k)\Bigr)\;},
\]
其中
\[
\pi(R_{k-1})=\frac{F_{k+3}}{36\cdot 2^{k-1}},
\qquad
\alpha_{k-1}=\frac{F_{k+1}}{2F_{k+3}},
\qquad
q_k=\frac{F_{k+3}}{2F_{k+5}}.
\]
并且存在常数 \(c>0\) 与 \(k_0\) 使得对所有 \(k\ge k_0\) 成立指数下界
\[
\mathcal L_k-\mathcal L_\infty\ \ge\ c\cdot (2\varphi^3)^{-k}.
\]
\end{theorem}
\begin{proof}
对数损失的 Bayes 分解给出
\[
\mathcal L(q)-\mathcal L_\infty
=
\EE\Bigl[
D_{\mathrm{KL}}\!\Bigl(\PP(g_{t+1}\in\cdot\mid g_{-\infty}^t)\ \big\|\ \mathrm{Bern}(q(g_{t-k+1}\cdots g_t))\Bigr)
\Bigr],
\]
其中右端的“真条件律”由全历史决定。
对任意 \(k\ge 2\)，在事件 \(\{S_t\in R_{k-1}\}\) 上，\(\varepsilon\)-机给出
\(
\PP(g_{t+1}=1\mid g_{-\infty}^t)=\PP(g_{t+1}=1\mid R_{k-1})=\alpha_{k-1}
\)；
并且由 \(R_n\xrightarrow{0}R_{n+1}\) 的结构可知在该事件上最近 \(k\) 个输出为 \(0^k\)，从而 \(g_{t-k+1}\cdots g_t=0^k\)。
因此
\[
\mathcal L(q)-\mathcal L_\infty
\ \ge\
\PP(R_{k-1})\cdot
D_{\mathrm{KL}}\!\Bigl(\mathrm{Bern}(\alpha_{k-1})\ \big\|\ \mathrm{Bern}(q(0^k))\Bigr).
\]
当取 \(q\) 为 \(\mathcal Q_k\) 的最优解时，由对数损失的点态 Bayes 最优性必有 \(q(0^k)=\PP(g_{t+1}=1\mid 0^k)=q_k\)，得到第一条不等式。
代入 \(\PP(R_{k-1})=\pi(R_{k-1})\) 的闭式即得所示显式下界。
\par
最后，对 \(k\to\infty\) 由 Binet 渐近 \(F_n\asymp \varphi^n\) 得 \(\pi(R_{k-1})\asymp (\varphi/2)^k\)，并且 \(|\alpha_{k-1}-q_k|\asymp \varphi^{-2k}\)。
由于两者均收敛到 \(\frac{1}{2\varphi^2}\in(0,1)\)，Bernoulli 的 \(D_{\mathrm{KL}}\) 在小偏差下具有二次型下界 \(D_{\mathrm{KL}}(\mathrm{Bern}(p)\|\mathrm{Bern}(p+\Delta))\gtrsim \Delta^2\)；
故 \(\mathcal L_k-\mathcal L_\infty\gtrsim (\varphi/2)^k\cdot\varphi^{-4k}=(2\varphi^3)^{-k}\)，从而存在 \(c,k_0\) 使对所有 \(k\ge k_0\) 成立。证毕。
\end{proof}

\begin{corollary}[达到指定误差阈值所需的上下文深度]\label{cor:fold-gauge-anomaly-logloss-context-depth}
\leanverified{paper\_fold\_gauge\_anomaly\_logloss\_context\_depth}
在定理 \ref{thm:fold-gauge-anomaly-logloss-context-gap} 的指数下界下，对任意 \(\varepsilon\in(0,1)\)，若 \(\mathcal L_k-\mathcal L_\infty\le \varepsilon\) 且 \(k\ge k_0\)，则必有
\[
k\ \ge\ \frac{\log(c/\varepsilon)}{\log(2\varphi^3)}.
\]
\end{corollary}

\begin{conjecture}[对数损失缺口的精确指数率与主导项]\label{conj:fold-gauge-anomaly-logloss-gap-asymptotic}
沿用定理 \ref{thm:fold-gauge-anomaly-logloss-context-gap} 的记号。
猜想存在常数 \(C>0\) 使当 \(k\to\infty\) 时
\[
\mathcal L_k-\mathcal L_\infty\ \sim\ C\cdot (2\varphi^3)^{-k},
\]
并且主导贡献由 \(\varepsilon\)-机尾部单态 \(R_{k-1}\) 的项
\(
\pi(R_{k-1})\cdot D_{\mathrm{KL}}\bigl(\mathrm{Bern}(\alpha_{k-1})\|\mathrm{Bern}(q_k)\bigr)
\)
所承载。
\end{conjecture}

\paragraph{零块语境下的有限记忆极小化误差}
对数损失刻画的是平均意义下的不可压缩缺口；在最坏情形的绝对误差口径下，\(\varepsilon\)-机尾部 \(\{R_n\}\) 诱导出一族可精确求解的极小化问题。

\begin{theorem}[零块语境下的最优有限记忆绝对误差（Fibonacci 闭式）]\label{thm:fold-gauge-anomaly-zero-block-minimax-abs-error}
取整数 \(n\ge 2\)。
令 \(\mathcal E_n^{(0)}\) 表示如下零块条件下的一步预测极小化误差：
允许任意预测器 \(\varphi:\{0,1\}^n\to[0,1]\)，并定义
\[
\mathcal E_n^{(0)}
:=\inf_{\varphi}\ \sup_{\substack{u\in\operatorname{supp}(g)\\ u_{t-n+1}\cdots u_t=0^n}}
\left|\ \PP(g_{t+1}=1\mid g_{-\infty}^t=u)\;-\;\varphi(0^n)\ \right|.
\]
则成立严格等式
\[
\boxed{\ \mathcal E_n^{(0)}=\frac{1}{4F_{n+3}F_{n+4}}\ }.
\]
\leanverified{paper\_fold\_zero\_block\_minimax\_abs\_error}
\leanverified{paper\_fold\_gauge\_anomaly\_zero\_block\_minimax\_abs\_error}
\end{theorem}
\begin{proof}
由定理 \ref{thm:fold-gauge-anomaly-epsilon-machine-synchronizing-word}，在观测到至少两个连续 \(0\) 后，预测态必落入尾部族 \(\{R_m:m\ge 1\}\)，并且每额外观测一个 \(0\) 就把预测态推进一级 \(R_m\mapsto R_{m+1}\)。
因此对任意满足 \(u_{t-n+1}\cdots u_t=0^n\) 的可实现全历史 \(u\)，存在 \(m\in\{n-1,n,\dots\}\cup\{\infty\}\) 使得
\[
\PP(g_{t+1}=1\mid g_{-\infty}^t=u)=\PP(g_{t+1}=1\mid R_m)=\alpha_m,
\]
其中 \(\alpha_m=\frac{F_{m+2}}{2F_{m+4}}\)（定理 \ref{thm:fold-gauge-anomaly-epsilon-machine-fibonacci-mobius}），并约定 \(\alpha_\infty=\lim_{m\to\infty}\alpha_m=\frac{1}{2\varphi^2}\)。
\par
由同一定理中的 Möbius 迭代 \(r_{m+1}=\frac{1}{4r_m+2}\) 与单调性可知：区间
\(
[\,\min\{\alpha_{m},\alpha_{m+1}\},\max\{\alpha_{m},\alpha_{m+1}\}\,]
\)
随 \(m\) 增大严格嵌套收缩；因此集合 \(\{\alpha_j:\ j\ge n-1\}\cup\{\alpha_\infty\}\) 的直径由前两项实现，
\[
\sup_{j\ge n-1}\alpha_j-\inf_{j\ge n-1}\alpha_j=|\alpha_{n-1}-\alpha_n|.
\]
对任意预测器 \(\varphi\)，在零块 \(0^n\) 上其输出只能取某个常数 \(q:=\varphi(0^n)\)；
于是最坏误差至少为区间直径的一半，
\[
\sup_{\substack{u\in\operatorname{supp}(g)\\ u_{t-n+1}\cdots u_t=0^n}}
\bigl|\PP(g_{t+1}=1\mid u)-q\bigr|
\ge \frac12|\alpha_{n-1}-\alpha_n|.
\]
取 \(q=\tfrac12(\alpha_{n-1}+\alpha_n)\) 可使该下界取等，故
\(
\mathcal E_n^{(0)}=\frac12|\alpha_{n-1}-\alpha_n|
\)。
由 \(\alpha_{n-1}=\frac{F_{n+1}}{2F_{n+3}}\)、\(\alpha_{n}=\frac{F_{n+2}}{2F_{n+4}}\)（定理 \ref{thm:fold-gauge-anomaly-epsilon-machine-fibonacci-mobius}），并用 Cassini 型恒等式
\(
F_{n+2}F_{n+3}-F_{n+1}F_{n+4}=(-1)^{n+1}
\)，得到
\[
|\alpha_{n-1}-\alpha_n|
=\frac{1}{2F_{n+3}F_{n+4}},
\]
从而 \(\mathcal E_n^{(0)}=\frac{1}{4F_{n+3}F_{n+4}}\)。证毕。
\end{proof}

\begin{corollary}[达到给定误差阈值所需的零块记忆深度]\label{cor:fold-gauge-anomaly-zero-block-memory-depth}
令 \(\varepsilon\in(0,1/8)\)，并定义
\[
n_{\varepsilon}:=\min\Bigl\{n\ge 2:\ F_{n+3}F_{n+4}\ge \frac{1}{4\varepsilon}\Bigr\}.
\]
则存在且仅存在长度 \(n\ge n_{\varepsilon}\) 的有限记忆预测器 \(\varphi\) 使得 \(\mathcal E_n^{(0)}\le \varepsilon\)。
并且当 \(\varepsilon\downarrow 0\) 时有渐近
\[
n_{\varepsilon}=\frac{1}{2\log\varphi}\log\frac{1}{\varepsilon}+O(1).
\]
\leanverified{paper\_fold\_gauge\_anomaly\_zero\_block\_memory\_depth}
\end{corollary}
\begin{proof}
第一断言由定理 \ref{thm:fold-gauge-anomaly-zero-block-minimax-abs-error} 的闭式立即得到。
渐近式由 Binet 公式 \(F_n=\varphi^n/\sqrt5+O(\varphi^{-n})\) 推出
\(
F_{n+3}F_{n+4}=\varphi^{2n+7}/5\cdot(1+o(1))
\)，代回定义并取对数即可。证毕。
\end{proof}

\leanverified{paper\_fold\_gauge\_anomaly\_finite\_state\_minimax\_abs\_error}
\begin{theorem}[有限状态近似的最坏情形误差尺度律]\label{thm:fold-gauge-anomaly-finite-state-minimax-abs-error}
令 \(\mathscr P_N\) 表示如下有限状态概率预测器类：其内态空间大小为 \(N\)，读入输出符号后按确定性规则更新，并在每一步输出下一符号为 \(1\) 的预测概率。
对 \(\Pi\in\mathscr P_N\) 记其单步最坏情形绝对误差为
\[
\mathcal R(\Pi):=\sup_{u\in\operatorname{supp}(g)}\left|\PP(g_{t+1}=1\mid g_{-\infty}^t=u)-\Pi(u)\right|,
\]
并令
\[
\mathcal R_N:=\inf_{\Pi\in\mathscr P_N}\mathcal R(\Pi).
\]
则存在常数 \(c_1,c_2>0\) 与 \(N_0\) 使得对所有 \(N\ge N_0\) 成立
\[
\boxed{\ c_1\,\varphi^{-2N}\ \le\ \mathcal R_N\ \le\ c_2\,\varphi^{-2N}\ }.
\]
从而达到误差阈值 \(\varepsilon\) 的最小态数满足 \(N(\varepsilon)=\Theta\!\bigl(\log(1/\varepsilon)\bigr)\)。
另一方面，若用长度 \(n\) 的朴素 \(n\)-阶 Markov 预测器（状态数 \(2^n\)）实现同阶精度，则状态数满足
\[
2^{n(\varepsilon)}=\exp\!\bigl(\Theta(\log(1/\varepsilon))\bigr),
\]
与 \(N(\varepsilon)=\Theta(\log(1/\varepsilon))\) 形成指数分离。
\end{theorem}
\begin{proof}
\textbf{上界。}
由定理 \ref{thm:fold-gauge-anomaly-epsilon-machine-synchronizing-word} 的 \(\varepsilon\)-机结构，取整数 \(M\ge 0\)，保留有限状态集 \(\{B,C,A,R_0,\dots,R_M\}\) 并将尾部 \(\{R_{m}:m>M\}\) 合并为单态 \(\widehat R\)。
在 \(\widehat R\) 上输出取区间 \(\{\alpha_m:m\ge M+1\}\cup\{\alpha_\infty\}\) 的中点。
由于 \(R_m\xrightarrow{0}R_{m+1}\) 与 \(R_m\xrightarrow{1}B\) 的确定性，该合并保持确定更新。
由定理 \ref{thm:fold-gauge-anomaly-epsilon-machine-fibonacci-mobius} 的收缩性，尾部区间直径为 \(|\alpha_{M+1}-\alpha_{M+2}|\)，从而该合并引入的最坏误差为其一半：
\[
\mathcal R(\Pi)\le \frac12|\alpha_{M+1}-\alpha_{M+2}|=\frac{1}{4F_{M+5}F_{M+6}}=\Theta(\varphi^{-2M}).
\]
取 \(N=M+4\) 并吸收常数即得 \(\mathcal R_N\le c_2\varphi^{-2N}\)。
\par
\textbf{下界。}
对每个 \(m\ge 0\)，定理 \ref{thm:fold-gauge-anomaly-epsilon-machine-infinite-markov-order} 的构造给出可实现历史，使得预测态为 \(R_m\)，并且真实下一位条件概率为 \(\alpha_m\)。
对任意 \(N\)-态预测器 \(\Pi\)，考虑由上述构造得到的 \(N+1\) 个历史（分别把预测态固定为 \(R_0,\dots,R_N\)）。
由于内态只有 \(N\) 个，必存在 \(0\le i<j\le N\) 使两条历史被映射到同一内态，因此 \(\Pi\) 在这两条历史上的输出概率相同。
于是
\[
\mathcal R(\Pi)\ \ge\ \frac12\,|\alpha_i-\alpha_j|\ \ge\ \frac12\min_{0\le i<j\le N}|\alpha_i-\alpha_j|.
\]
由差值闭式
\(
\alpha_{m+1}-\alpha_{m}=\frac{(-1)^{m}}{2F_{m+4}F_{m+5}}
\)
可知 \((\alpha_{2m})_{m\ge 0}\) 严格递增、\((\alpha_{2m+1})_{m\ge 0}\) 严格递减，并且 \(\alpha_{2m}<\alpha_{2m+1}\)。
因此在集合 \(\{\alpha_0,\dots,\alpha_N\}\) 中，最小间距由最靠近极限 \(\alpha_\infty\) 的一对异奇偶项实现，即 \(|\alpha_{N-1}-\alpha_N|\)。
从而
\[
\min_{0\le i<j\le N}|\alpha_i-\alpha_j|=|\alpha_{N-1}-\alpha_N|
=\frac{1}{2F_{N+3}F_{N+4}}=\Theta(\varphi^{-2N}),
\]
故 \(\mathcal R(\Pi)\ge \Theta(\varphi^{-2N})\)，进而 \(\mathcal R_N\ge c_1\varphi^{-2N}\)。
两侧合并即得结论。证毕。
\end{proof}

\endinput

\begin{remark}[审计脚本]\label{rem:fold-gauge-anomaly-bernoulli-p-audit-script}
定理 \ref{thm:fold-gauge-anomaly-bernoulli-p-laws}--推论 \ref{cor:fold-gauge-anomaly-bernoulli-p-specialize-half} 中出现的闭式表达式（包括三阶密度累积量、协方差生成函数、压力四次方程以及 Lee--Yang 判别式阈值核验）可由脚本 \texttt{scripts/exp\_fold\_gauge\_anomaly\_bernoulli\_p\_closed\_form.py} 进行一致性核验。
\end{remark}
\endinput

\subsubsection{阈值折点的几何刚性}\label{subsubsec:fold-gauge-anomaly-bernoulli-p-kink-geometry}

\begin{proposition}[Bayes 阈值函数的不可微性与唯一极大点]\label{prop:fold-gauge-anomaly-bernoulli-p-kink-unique-max}
沿用命题 \ref{prop:fold-gauge-anomaly-bernoulli-p-bayes-phase-transition} 的记号，定义
\[
g^\star(p):=\min\!\left\{\,g_\ast(p),\ 1-p\,\right\},
\qquad
g_\ast(p)=\frac{p^2(3-2p)}{1+p^3},
\qquad p\in(0,1),
\]
并令 \(p_{\mathrm c}\in(0,1)\) 为四次方程
\(
p^4-3p^3+3p^2+p-1=0
\)
在 \((0,1)\) 内的唯一根。
则：
\begin{enumerate}
  \item \(g^\star\) 在 \((0,1)\) 上连续，并在 \(p=p_{\mathrm c}\) 处不可微；
  \item \(p_{\mathrm c}\) 是 \(g^\star\) 在 \((0,1)\) 上的唯一全局极大点；
  \item \(p=\tfrac12\) 满足 \(p< p_{\mathrm c}\)，因此 \(g^\star(\tfrac12)=g_\ast(\tfrac12)=\tfrac49\)，并且该常数落在解析分支 \(g_\ast\) 的内部而非线性分支 \(1-p\)。
\end{enumerate}
\end{proposition}
\begin{proof}
由命题 \ref{prop:fold-gauge-anomaly-bernoulli-p-bayes-phase-transition}，\(g^\star\) 在 \((0,1)\) 上满足分段闭式
\[
g^\star(p)=
\begin{cases}
g_\ast(p),& 0<p\le p_{\mathrm c},\\
1-p,& p_{\mathrm c}\le p<1,
\end{cases}
\qquad
g_\ast(p_{\mathrm c})=1-p_{\mathrm c}.
\]
因此 \(g^\star\) 作为两条连续函数的拼接在 \(p_{\mathrm c}\) 处连续。
\par
不可微性：由 \eqref{eq:fold-gauge-anomaly-bernoulli-p-density-derivative}，
\[
g_\ast'(p)= -\frac{3p\,(p^3+2p-2)}{(p+1)^2(p^2-p+1)^2}.
\]
设 \(h(p):=p^3+2p-2\)。则 \(h\) 在 \((0,1)\) 上严格递增且 \(h(\tfrac34)=-\tfrac{5}{64}<0\)、\(h(\tfrac45)=\tfrac{14}{125}>0\)，故其唯一根 \(p_\star\) 落在 \((\tfrac34,\tfrac45)\)。
另一方面，四次式 \(q(p):=p^4-3p^3+3p^2+p-1\) 满足 \(q(\tfrac12)=-\tfrac{1}{16}<0\)、\(q(\tfrac34)=\tfrac{125}{256}>0\)，并且在 \((0,1)\) 内仅有一根 \(p_{\mathrm c}\)，故 \(p_{\mathrm c}\in(\tfrac12,\tfrac34)\)。
于是 \(p_{\mathrm c}<p_\star\)，从而 \(h(p_{\mathrm c})<0\) 并给出 \(g_\ast'(p_{\mathrm c})>0\)。
而右侧分支 \(1-p\) 的导数恒为 \(-1\)，因此左右导数不相等，\(g^\star\) 在 \(p_{\mathrm c}\) 处不可微。
\par
唯一极大点：由上式 \(p_{\mathrm c}<p_\star\) 且 \(h<0\) 在 \((0,p_\star)\) 上成立，可知 \(g_\ast'(p)>0\) 对所有 \(p\in(0,p_{\mathrm c}]\) 成立，从而 \(g^\star=g_\ast\) 在 \((0,p_{\mathrm c}]\) 上严格递增。
而在 \([p_{\mathrm c},1)\) 上 \(g^\star(p)=1-p\) 严格递减。
两分支在 \(p_{\mathrm c}\) 处匹配，故 \(p_{\mathrm c}\) 为唯一全局极大点。
\par
最后，由 \(p_{\mathrm c}>\tfrac12\) 知 \(\tfrac12<p_{\mathrm c}\)，从而 \(g^\star(\tfrac12)=g_\ast(\tfrac12)\)。
代入 \(g_\ast(p)=\frac{p^2(3-2p)}{1+p^3}\) 得
\[
g_\ast\!\left(\tfrac12\right)
=\frac{\tfrac14\cdot 2}{1+\tfrac18}
=\frac{\tfrac12}{\tfrac98}
=\frac49.
\]
证毕。
\end{proof}

\endinput


% Mismatch language: word complexity + Parry MME + cohomology obstruction.

\paragraph{失配语言的词复杂度递推、极大熵态与上同调障碍}\label{par:fold-gauge-anomaly-mismatch-language-word-complexity}
沿用定理 \ref{thm:fold-gauge-anomaly-g-support-sofic} 的有限型支撑 \(\Sigma_g\subset\{0,1\}^{\ZZ}\)。
记其长度 \(n\) 的允许词集合为 \(\mathcal L_n(\Sigma_g)\)，并令
\[
N(n):=\#\mathcal L_n(\Sigma_g)\qquad(n\ge 1).
\]

\begin{proposition}[允许词计数的精确递推与指数谱分解]\label{prop:fold-gauge-anomaly-mismatch-language-word-count-recurrence}
对一切 \(n\ge 6\)，允许词计数满足线性递推
\[
\boxed{\ N(n)=N(n-1)+N(n-2)+N(n-4)+N(n-5)\ }.
\]
其初值可取为
\[
N(1)=2,\quad N(2)=4,\quad N(3)=8,\quad N(4)=14,\quad N(5)=25.
\]
令 \(\lambda>1\) 为定理 \ref{thm:fold-gauge-anomaly-g-support-sofic} 中多项式
\(
x^5-x^4-x^3-x-1
\)
的 Perron 实根。
则存在常数 \(C>0\) 与 \(0<\beta<\lambda\) 使
\[
N(n)=C\,\lambda^{n}+O(\beta^{n}).
\]
\leanverified{mismatchWordCount\_eleven}
\leanverified{mismatchWordCount\_twelve}
\leanverified{mismatchWordCount\_thirteen}
\leanverified{mismatchWordCount\_lt\_pow\_two}
\leanverified{paper\_mismatchWordCount\_extended}
\leanverified{mismatchWordCount\_fourteen}
\leanverified{mismatchWordCount\_fifteen}
\leanverified{mismatchWordCount\_sixteen}
\leanverified{paper\_mismatchWordCount\_14\_to\_16\_package}
\leanverified{paper\_fold\_gauge\_anomaly\_mismatch\_language\_word\_count\_recurrence}
\end{proposition}
\begin{proof}
由定理 \ref{thm:fold-gauge-anomaly-g-support-sofic}，\(\Sigma_g\) 为以禁块 \(0010,0100\) 定义的 \(3\)-步 SFT；
因此其语言为一个正则对象，可由“记忆 \(3\)”的确定型有限自动机计数（transfer-matrix 方法；参见 \cite{LindMarcus1995,Kitchens1998SymbolicDynamics}）。
对该自动机的可达态空间作最小化后得到一个有限维线性递推系统，其特征多项式的主因子即为
\(
x^5-x^4-x^3-x-1
\)；
从而 \(N(n)\) 满足所示五阶递推，并以 \(N(1),\dots,N(5)\) 唯一确定。
\par
谱分解部分由递推的特征根展开给出：主项由 Perron 根 \(\lambda\) 贡献，其余根的模严格小于 \(\lambda\)，故余项为 \(O(\beta^n)\)（可取 \(\beta\) 为其余特征根模的最大值）。
证毕。
\leanverified{paper\_mismatch\_word\_count\_initial\_values}
\leanverified{paper\_mismatch\_word\_count\_strict\_mono}
\leanverified{paper\_mismatch\_perron\_root\_bound}
\end{proof}

\begin{corollary}[有限型支撑上的内禀唯一极大熵态]\label{cor:fold-gauge-anomaly-mismatch-language-intrinsic-ergodic}
\(\Sigma_g\) 为不可约有限型子移位，因而存在唯一极大熵测度（Parry 测度）\(\mu_{\max}\)，并且 \((\Sigma_g,\sigma)\) 内禀可测。
取一个右解析且 follower-separated 的呈示（例如定理 \ref{thm:fold-gauge-anomaly-g-support-sofic} 的三步 de Bruijn 呈示），即可将其视为右 Fischer cover 的可计算接口。
该测度可由定理 \ref{thm:fold-gauge-anomaly-g-support-sofic} 的三步 de Bruijn 呈示邻接矩阵 \(A\) 的 Perron--Frobenius 本征结构显式给出：若 \(\lambda=\rho(A)\)、\(r,l\) 分别为右/左 Perron 向量并按 \(\langle l,r\rangle=1\) 归一化，则对任意允许边 \(i\to j\) 有 Parry 核
\[
P(i\to j)=\frac{A_{ij}r_j}{\lambda r_i},
\]
并由逐边取符号的 lift 求和得到柱集概率。
\leanverified{paper\_fold\_gauge\_anomaly\_mismatch\_language\_intrinsic\_ergodic}
\end{corollary}

\begin{proposition}[失配标记的非 coboundary 性]\label{prop:fold-gauge-anomaly-mismatch-language-non-coboundary}
在定理 \ref{thm:fold-gauge-anomaly-g-support-sofic} 的三步 de Bruijn 呈示上，把边的输出标记视为整数值函数 \(\gamma(e)\in\{0,1\}\)。
则不存在顶点函数 \(h\) 使得对任意边 \(e:v\to w\) 成立
\[
\gamma(e)=h(w)-h(v).
\]
等价地，\(\gamma\) 在图上同调中定义了非零的 \(1\)-上同调类。
\leanverified{paper\_mismatch\_label\_non\_coboundary}
\end{proposition}
\begin{proof}
若 \(\gamma\) 为 coboundary，则任意有向回路 \(C\) 上必有 \(\sum_{e\in C}\gamma(e)=0\)。
但 \(\Sigma_g\) 含周期点 \(1^{\ZZ}\)，其在呈示图上对应于标记恒为 \(1\) 的回路；
取其任意周期长度 \(n\) 的闭路段即得 \(\sum_{e\in C}\gamma(e)=n\neq 0\)，矛盾。
证毕。
\end{proof}

\endinput

% Threshold corollaries for Bernoulli(p) baseline (kept separate to avoid inflating the main baseline file).

\begin{corollary}[单点输入--输出位的唯一独立阈值]\label{cor:fold-gauge-anomaly-bernoulli-p-bitpair-independence-threshold}
在命题 \ref{prop:fold-gauge-anomaly-bernoulli-p-bitpair-law} 的平稳 Bernoulli\((p)\) 基线设定下，
令 \(p_\phi:=\frac{\sqrt5-1}{2}\in(0,1)\)。
则随机变量 \(X_0,Y_0\in\{0,1\}\) 严格独立当且仅当 \(p=p_\phi\)。
\leanverified{paper\_fold\_gauge\_anomaly\_bitpair\_independence\_threshold}
\end{corollary}
\begin{proof}
由命题 \ref{prop:fold-gauge-anomaly-bernoulli-p-bitpair-law}，
\[
\Cov(X_0,Y_0)
=P(X_0=1,Y_0=1)-P(X_0=1)P(Y_0=1)
=-\frac{p(1-p)^2(p^2+p-1)}{1+p^3}.
\]
对 \(p\in(0,1)\) 有 \(p(1-p)^2/(1+p^3)>0\)，故 \(\Cov(X_0,Y_0)=0\) 当且仅当 \(p^2+p-1=0\)；
其在 \((0,1)\) 的唯一根为 \(p_\phi\)。
由于 \(\{0,1\}\)-值随机变量的联合分布由边缘分布与 \(P(X_0=1,Y_0=1)\) 唯一确定，
\(\Cov(X_0,Y_0)=0\) 等价于 \(P(X_0=1,Y_0=1)=P(X_0=1)P(Y_0=1)\)，从而等价于独立性。
证毕。
\end{proof}

\begin{corollary}[失配密度极大点的 Cardano 闭式]\label{cor:fold-gauge-anomaly-bernoulli-p-density-maximizer-cardano}
令 \(g_\ast(p)=\frac{p^2(3-2p)}{1+p^3}\)。
记 \(p_\star\in(0,1)\) 为三次方程
\[
p^3+2p-2=0
\]
在 \((0,1)\) 中的唯一实根（推论 \ref{cor:fold-gauge-anomaly-bernoulli-p-density-unimodal}）。
则该根可由 Cardano 公式写为
\[
p_\star=\sqrt[3]{\,1+\sqrt{\tfrac{35}{27}}\,}+\sqrt[3]{\,1-\sqrt{\tfrac{35}{27}}\,},
\]
并且极大值满足恒等式 \(g_\ast(p_\star)=p_\star^2\)。
\leanverified{paper\_fold\_gauge\_anomaly\_density\_maximizer\_cardano}
\end{corollary}
\begin{proof}
由推论 \ref{cor:fold-gauge-anomaly-bernoulli-p-density-unimodal}，\(p_\star\) 为所示三次方程在 \((0,1)\) 的唯一根，
且 \(g_\ast(p_\star)=p_\star^2\) 已在该推论中给出。
将方程写为降次形式 \(p^3+2p=2\)，应用 Cardano 公式得到所示闭式表示。证毕。
\end{proof}

\begin{theorem}[Gallavotti--Cohen 缺陷的线性响应阈值]\label{thm:fold-gauge-anomaly-gc-defect-linear-response-threshold}
在定理 \ref{thm:fold-gauge-anomaly-bernoulli-p-laws} 的 Bernoulli\((p)\) 基线设定下，定义压力
\[
P_{G,p}(\theta)=\log\lambda(e^\theta,p),
\]
以及 Gallavotti--Cohen 缺陷函数
\[
\Delta_p(\theta):=P_{G,p}(\theta)-\theta-P_{G,p}(-\theta).
\]
则 \(\Delta_p\) 在 \(\theta=0\) 的泰勒展开满足
\[
\Delta_p(\theta)=\bigl(2g_\ast(p)-1\bigr)\theta+O(\theta^3),
\qquad \theta\to 0,
\]
其中失配密度 \(g_\ast(p)=\frac{p^2(3-2p)}{1+p^3}\)。
因此存在唯一阈值
\[
p_0=\frac{1+\sqrt{21}}{10}\in(0,1)
\]
使得当 \(p<p_0\) 时，\(\Delta_p(\theta)<0\) 对一切充分小的 \(\theta>0\) 成立；当 \(p>p_0\) 时，\(\Delta_p(\theta)>0\) 对一切充分小的 \(\theta>0\) 成立。
\leanverified{paper\_fold\_gauge\_anomaly\_gc\_defect\_linear\_response\_threshold}
\end{theorem}

\begin{proof}
由定理 \ref{thm:fold-gauge-anomaly-bernoulli-p-laws}(4)，\(P_{G,p}\) 在 \(\theta=0\) 的邻域实解析，并且
\(
P_{G,p}'(0)=g_\ast(p)
\)。
由定义，
\[
\Delta_p'(0)=P_{G,p}'(0)-1-(-P_{G,p}'(0))=2P_{G,p}'(0)-1=2g_\ast(p)-1.
\]
由于 \(\Delta_p\) 为奇函数（见命题 \ref{prop:fold-gauge-anomaly-bernoulli-p-gc-defect-farfield-critical} 的证明开头），其泰勒展开只含奇次项，从而
\(
\Delta_p(\theta)=\Delta_p'(0)\theta+O(\theta^3)
\)。
代入 \(g_\ast(p)=\frac{p^2(3-2p)}{1+p^3}\) 得
\[
\Delta_p'(0)=\frac{6p^2-5p^3-1}{1+p^3}
=\frac{(1-p)(5p^2-p-1)}{1+p^3}.
\]
在区间 \(p\in(0,1)\) 上，分母与因子 \(1-p\) 严格为正，故 \(\Delta_p'(0)\) 的符号由二次多项式 \(5p^2-p-1\) 的符号决定。
该二次方程在 \((0,1)\) 内存在且仅存在一根
\(
p_0=\frac{1+\sqrt{21}}{10}
\)。
结合 \(\Delta_p(\theta)=\Delta_p'(0)\theta+O(\theta^3)\) 即得小正 \(\theta\) 下的符号结论。证毕。
\end{proof}

\begin{proposition}[高场一阶指数渐近与黄金点闭式系数]\label{prop:fold-gauge-anomaly-gc-defect-farfield-expansion}
\leanverified{paper\_fold\_gauge\_anomaly\_gc\_defect\_farfield\_expansion}
在命题 \ref{prop:fold-gauge-anomaly-bernoulli-p-gc-defect-farfield-critical} 的记号下，令
\[
\alpha(p):=\frac{1-p+\sqrt{(1-p)(1+3p)}}{2},
\qquad
c(p):=\bigl(p^2(1-p)\bigr)^{1/3}.
\]
则存在显式函数 \(A(p)\) 使当 \(\theta\to+\infty\) 时
\[
\Delta_p(\theta)=\log\frac{c(p)}{\alpha(p)}+A(p)e^{-\theta}+O(e^{-2\theta}),
\]
并且可以取
\[
A(p)=\frac{d(p)}{c(p)},
\qquad
d(p)=\frac13\left(\left(\frac{p}{1-p}\right)^{1/3}+(1-p)\right).
\]
特别地，在黄金偏置点 \(p=\frac{1+\sqrt5}{4}\) 处有 \(\log\frac{c(p)}{\alpha(p)}=0\) 且
\[
\Delta_{p}(\theta)=\frac{5+\sqrt5}{6}e^{-\theta}+O(e^{-2\theta}),
\qquad \theta\to+\infty.
\]
\end{proposition}

\begin{proof}
令 \(u=e^\theta\) 并记 \(\lambda(u,p)\) 为定理 \ref{thm:fold-gauge-anomaly-bernoulli-p-laws} 中的 Perron 根（满足四次方程 \eqref{eq:fold-gauge-anomaly-bernoulli-p-pressure-quartic}）。
由定义
\[
\Delta_p(\theta)=\log\frac{\lambda(u,p)}{u\,\lambda(u^{-1},p)}.
\]
\par
\textbf{(i) 大场端 \(u\to\infty\)。}
设 \(\lambda(u,p)=u\,s(u,p)\)。
将其代入 \eqref{eq:fold-gauge-anomaly-bernoulli-p-pressure-quartic} 并以 \(u^4\) 归一化，可得
\[
s(u,p)^4-p^2(1-p)\,s(u,p)+O(u^{-1})=0,
\qquad u\to\infty.
\]
由 Perron 分支的正性与唯一性，\(s(u,p)\to c(p)\) 且满足 \(c(p)^3=p^2(1-p)\)。
进一步令
\(
\lambda(u,p)=c(p)u+d+O(u^{-1})
\)。
将该展开代回 \eqref{eq:fold-gauge-anomaly-bernoulli-p-pressure-quartic}，最高阶 \(u^4\) 项由 \(c(p)^3=p^2(1-p)\) 消去；比较 \(u^3\) 系数得到一条线性方程，从而唯一确定
\[
d(p)=\frac13\left(\left(\frac{p}{1-p}\right)^{1/3}+(1-p)\right).
\]
因此
\[
\log\frac{\lambda(u,p)}{u}
=\log\left(c(p)+\frac{d(p)}{u}+O(u^{-2})\right)
=\log c(p)+\frac{d(p)}{c(p)}u^{-1}+O(u^{-2}).
\]
\par
\textbf{(ii) 小场端 \(u\to 0^+\)。}
由命题 \ref{prop:fold-gauge-anomaly-bernoulli-p-endpoint-ldp}，有 \(\lambda(0,p)=\alpha(p)\)。
对四次方程 \eqref{eq:fold-gauge-anomaly-bernoulli-p-pressure-quartic} 记其左端为 \(F(\lambda,u,p)\)。
在 \((\lambda,u)=(\alpha(p),0)\) 处有 \(F(\alpha,0,p)=0\)。
利用 \(\alpha(p)\) 满足二次关系 \(\alpha^2-(1-p)\alpha-p(1-p)=0\)，可直接化简得到
\[
\partial_u F(\alpha(p),0,p)=0,
\]
从而 Perron 分支在 \(u=0\) 处的一阶项消失，并存在常数 \(C(p)<\infty\) 使
\[
\lambda(u,p)=\alpha(p)+O(u^2),
\qquad u\to 0^+.
\]
于是
\[
\log\lambda(u^{-1},p)=\log\alpha(p)+O(u^{-2}),
\qquad u\to\infty.
\]
\par
将 (i)(ii) 合并并取 \(u=e^\theta\) 即得
\[
\Delta_p(\theta)=\log\frac{c(p)}{\alpha(p)}+\frac{d(p)}{c(p)}e^{-\theta}+O(e^{-2\theta}).
\]
黄金偏置点由命题 \ref{prop:fold-gauge-anomaly-bernoulli-p-gc-defect-farfield-critical} 给出为 \(p=\frac{1+\sqrt5}{4}\)，此时 \(\log\frac{c(p)}{\alpha(p)}=0\) 且 \(p/(1-p)=\varphi^3\)，从而
\(
d(p)=\frac13(\varphi+(1-p))
\)、
\(
c(p)=\tfrac12
\)，化简得到 \(A(p)=d(p)/c(p)=\frac{5+\sqrt5}{6}\)。
证毕。
\end{proof}

\begin{corollary}[必然零点区间与区间夹逼]\label{cor:fold-gauge-anomaly-gc-defect-must-cross}
\leanverified{paper\_fold\_gauge\_anomaly\_gc\_defect\_must\_cross}
令
\(
p_0=\frac{1+\sqrt{21}}{10}
\)（定理 \ref{thm:fold-gauge-anomaly-gc-defect-linear-response-threshold}）
与
\(
p_\infty=\frac{1+\sqrt5}{4}
\)（命题 \ref{prop:fold-gauge-anomaly-bernoulli-p-gc-defect-farfield-critical}）。
则对任意 \(p\in(p_0,p_\infty)\)，存在 \(\theta_p>0\) 使得 \(\Delta_p(\theta_p)=0\)。
\end{corollary}

\begin{proof}
由定理 \ref{thm:fold-gauge-anomaly-gc-defect-linear-response-threshold}，当 \(p>p_0\) 时，存在 \(\varepsilon>0\) 使 \(\Delta_p(\theta)>0\) 对所有 \(\theta\in(0,\varepsilon)\) 成立。
另一方面，由命题 \ref{prop:fold-gauge-anomaly-bernoulli-p-gc-defect-farfield-critical}，当 \(p<p_\infty\) 时有
\(
\lim_{\theta\to+\infty}\Delta_p(\theta)<0
\)，故存在 \(T>0\) 使得对所有 \(\theta>T\) 有 \(\Delta_p(\theta)<0\)。
由于 \(\Delta_p\) 在 \((0,\infty)\) 上连续，介值定理给出某个 \(\theta_p\in(\varepsilon,T)\) 使 \(\Delta_p(\theta_p)=0\)。证毕。
\end{proof}

\begin{corollary}[黄金偏置处的零温中性与一阶指数衰减]\label{cor:fold-gauge-anomaly-gc-defect-golden-neutral}
\leanverified{paper\_fold\_gauge\_anomaly\_gc\_defect\_golden\_neutral}
取 \(p=\frac{1+\sqrt5}{4}\)。
则 \(\Delta_p(\theta)\to 0\) 且满足一阶指数衰减
\[
\Delta_p(\theta)=\frac{5+\sqrt5}{6}e^{-\theta}+O(e^{-2\theta}),
\qquad \theta\to+\infty.
\]
\end{corollary}

\begin{proof}
由命题 \ref{prop:fold-gauge-anomaly-gc-defect-farfield-expansion} 的特例直接得到。证毕。
\end{proof}

\endinput

\paragraph{Bernoulli\texorpdfstring{$(p)$}{(p)} 基线的归并闭式链：Parry 归一化、代数压力与端点指数}
以下结果在段落 \ref{par:fold-gauge-anomaly-bernoulli-p} 的同一换能器口径下整理，符号与
定理 \ref{thm:fold-gauge-anomaly-bernoulli-p-laws}、
命题 \ref{prop:fold-gauge-anomaly-bernoulli-p-bitpair-law}、
命题 \ref{prop:fold-gauge-anomaly-bernoulli-p-spectrum}、
命题 \ref{prop:fold-gauge-anomaly-bernoulli-p-endpoint-ldp}
保持一致。

\begin{theorem}[规范差压力与 Parry 归一化链的显式闭式]\label{thm:fold-bernoulli-p-parry-pressure-chain}
取 \(p\in(0,1)\)，记 \(q:=1-p\)。令失配计数
\[
G_m=\sum_{t=0}^{m-1}\mathbf 1_{\{X_t\neq Y_t\}},
\]
并定义扭曲矩阵
\begin{equation}\label{eq:fold-bernoulli-p-parry-Atheta}
A_{\theta}(p)=
\begin{pmatrix}
q&q e^\theta&0&p\\
0&0&p e^\theta&0\\
p e^\theta&1&0&0\\
q&0&0&0
\end{pmatrix}.
\end{equation}
记 \(\rho(\theta;p):=\rho(A_\theta(p))\)。则极限压力存在且
\begin{equation}\label{eq:fold-bernoulli-p-parry-pressure}
P_G(\theta;p):=\lim_{m\to\infty}\frac1m\log\EE_p[e^{\theta G_m}]
=\log \rho(\theta;p).
\end{equation}
在 \(\theta=0\) 处有 \(\rho(0;p)=1\)，并可取 Perron 右、左特征向量为
\begin{equation}\label{eq:fold-bernoulli-p-parry-eigenvectors}
r(p)=\bigl(q,\ p^2,\ p,\ q^2\bigr)^{\mathsf T},
\qquad
\ell(p)=\bigl(1,\ 1,\ p,\ p\bigr).
\end{equation}
由 Parry 归一化得到
\begin{equation}\label{eq:fold-bernoulli-p-parry-kernel}
P(p)=\mathrm{diag}(r(p))^{-1}A_0(p)\,\mathrm{diag}(r(p))
=
\begin{pmatrix}
q&p^2&0&pq\\
0&0&1&0\\
q&p&0&0\\
1&0&0&0
\end{pmatrix},
\end{equation}
其平稳分布
\begin{equation}\label{eq:fold-bernoulli-p-parry-stationary}
\pi(p)=\frac{\ell\odot r}{\langle \ell,r\rangle}
=\frac1{1+p^3}\bigl(q,\ p^2,\ p^2,\ p q^2\bigr).
\end{equation}
\end{theorem}
\begin{proof}
\eqref{eq:fold-bernoulli-p-parry-Atheta} 与
\eqref{eq:fold-bernoulli-p-parry-pressure} 即
定理 \ref{thm:fold-gauge-anomaly-bernoulli-p-laws} 的矩阵接口重写。
直接代入可验
\[
A_0(p)\,r(p)=r(p),\qquad
\ell(p)A_0(p)=\ell(p),
\]
得 \eqref{eq:fold-bernoulli-p-parry-eigenvectors}。
由公式
\(
P_{ij}=A_{0,ij}r_j/r_i
\)
得 \eqref{eq:fold-bernoulli-p-parry-kernel}。
再由
\(
\pi_i\propto \ell_i r_i
\)
并归一化，得到 \eqref{eq:fold-bernoulli-p-parry-stationary}。证毕。
\end{proof}
\leanverified{paper\_fold\_bernoulli\_p\_parry\_pressure\_chain}

\begin{proposition}[状态 \(a\) 的再生结构与回返时间的精确律]\label{prop:fold-bernoulli-p-return-time-a}
沿用定理 \ref{thm:fold-bernoulli-p-parry-pressure-chain} 的链 \((S_t)_{t\ge 0}\) 与边势 \(g\)。
从 \(a\) 出发令首次回返时刻
\[
\tau_a:=\inf\{t\ge 1:S_t=a\mid S_0=a\}.
\]
则 \(\tau_a\) 的支撑为 \(\{1,2\}\cup\{2n+3:n\ge 0\}\)，并且对 \(n\ge 0\) 有
\begin{equation}\label{eq:fold-bernoulli-p-return-time-a-law}
\PP_p(\tau_a=1)=q,\qquad
\PP_p(\tau_a=2)=pq,\qquad
\PP_p(\tau_a=2n+3)=p^{n+2}q.
\end{equation}
其概率母函数为
\begin{equation}\label{eq:fold-bernoulli-p-return-time-a-pgf}
\EE_p\!\left[z^{\tau_a}\right]
=qz+pq\,z^2+\frac{p^2q\,z^3}{1-pz^2}.
\end{equation}

进一步令一次回返段长度 \(L_a:=\tau_a\)，并记该段内失配边数
\[
M_a:=\sum_{t=0}^{L_a-1} g(S_t,S_{t+1}).
\]
则联合分布完全显式：对 \(n\ge 0\),
\begin{equation}\label{eq:fold-bernoulli-p-return-segment-a-joint}
\PP_p(L_a=1,M_a=0)=q,\quad
\PP_p(L_a=2,M_a=0)=pq,\quad
\PP_p(L_a=2n+3,M_a=n+3)=p^{n+2}q,
\end{equation}
并且在 \(\{L_a\ge 3\}\) 上存在确定的仿射约束
\begin{equation}\label{eq:fold-bernoulli-p-return-segment-a-affine}
L_a=2M_a-3.
\end{equation}
\leanverified{paper\_fold\_bernoulli\_p\_return\_time\_a}
\end{proposition}
\begin{proof}
由 \eqref{eq:fold-bernoulli-p-parry-kernel} 可见：从 \(a\) 出发第一步以概率 \(q\) 走 \(a\to a\)，从而 \(\tau_a=1\)；
以概率 \(pq\) 走 \(a\to d\) 且必有 \(d\to a\)，从而 \(\tau_a=2\)；
以概率 \(p^2\) 走 \(a\to b\) 且必有 \(b\to c\)。
在第三种分支上，处于 \(c\) 时以概率 \(q\) 走 \(c\to a\) 回返，以概率 \(p\) 走 \(c\to b\) 并必经 \(b\to c\) 回到 \(c\)。
因此在 \(c\) 处选择 \(c\to b\) 的次数 \(N\) 服从参数为 \(q\) 的几何分布 \(\PP(N=n)=p^n q\)，并且 \(\tau_a=3+2N\)。
与进入该分支的概率 \(p^2\) 相乘即得 \eqref{eq:fold-bernoulli-p-return-time-a-law}；母函数 \eqref{eq:fold-bernoulli-p-return-time-a-pgf} 由几何级数求和得到。

对联合律，注意 \(\tau_a\in\{1,2\}\) 时路径仅经匹配边，故 \(M_a=0\)。
当 \(N=n\) 时回返段轨道强制为
\(
a\to b\to c(\to b\to c)^n\to a
\)，其长度 \(L_a=2n+3\)。
由边势 \eqref{eq:fold-gauge-anomaly-bernoulli-p-edge-potential}，该段中失配边恰为 \(a\to b\)、全部 \(b\to c\) 以及末端 \(c\to a\)，故 \(M_a=1+(n+1)+1=n+3\)，从而得到 \eqref{eq:fold-bernoulli-p-return-segment-a-joint}。
消去 \(n\) 立即得到 \eqref{eq:fold-bernoulli-p-return-segment-a-affine}。证毕。
\end{proof}

\begin{theorem}[平稳联合律、失配密度与输出压缩闭式]\label{thm:fold-bernoulli-p-bitpair-gamma-q-closed}
在定理 \ref{thm:fold-bernoulli-p-parry-pressure-chain} 的平稳链上，\((X,Y)\) 的单点联合分布为
\begin{align}
\PP_p(X=0,Y=0)&=\frac{1-p-p^2+2p^3-p^4}{1+p^3},\nonumber\\
\PP_p(X=0,Y=1)&=\frac{p^2(1-p)}{1+p^3},\nonumber\\
\PP_p(X=1,Y=0)&=\frac{p^2(2-p)}{1+p^3},\label{eq:fold-bernoulli-p-bitpair-joint-law}\\
\PP_p(X=1,Y=1)&=\frac{p\bigl(p^3+(1-p)^2\bigr)}{1+p^3}. \nonumber
\end{align}
从而失配密度
\begin{equation}\label{eq:fold-bernoulli-p-gamma}
\gamma(p):=\lim_{m\to\infty}\frac1m\EE_p[G_m]
=\PP_p(X\neq Y)
=\frac{p^2(3-2p)}{1+p^3}.
\end{equation}
输出端“1”密度满足
\begin{equation}\label{eq:fold-bernoulli-p-output-density}
q(p):=\PP_p(Y=1)
=\frac{p(p^3-p+1)}{1+p^3},
\qquad
0<q(p)<p\quad(p\in(0,1)).
\end{equation}
\leanverified{paper\_fold\_bernoulli\_p\_bitpair\_gamma\_q\_closed}
\leanverified{paper\_fold\_bernoulli\_p\_bitpair\_gamma\_closed\_seeds}
\end{theorem}
\begin{proof}
\eqref{eq:fold-bernoulli-p-bitpair-joint-law} 与
\eqref{eq:fold-bernoulli-p-gamma}
由命题 \ref{prop:fold-gauge-anomaly-bernoulli-p-bitpair-law} 直接得到。
\eqref{eq:fold-bernoulli-p-output-density} 的第一式来自
\(
\PP(Y=1)=\PP(0,1)+\PP(1,1)
\)；
第二式由恒等式
\[
p-q(p)=\frac{p^2}{1+p^3}>0
\]
给出。证毕。
\end{proof}

\begin{theorem}[代数压力四次方程与 CLT 方差率闭式]\label{thm:fold-bernoulli-p-pressure-quartic-clt-variance}
令 \(u=e^\theta\)，记 \(\lambda(u,p):=\rho(A_{\log u}(p))\)。
则 \(\lambda(u,p)\) 是下列四次方程在 \(\lambda>0\) 上的唯一最大实根：
\begin{equation}\label{eq:fold-bernoulli-p-pressure-quartic-restated}
\lambda^4+(p-1)\lambda^3+\bigl(p^2-p(u+1)\bigr)\lambda^2
+\bigl(p^3u^3-p^2u^3-p^2u+pu\bigr)\lambda
+\bigl(-p^3u+p^2u\bigr)=0.
\end{equation}
并且 \(P_G(\theta;p)=\log\lambda(e^\theta,p)\) 在 \(\theta\in\RR\) 上解析。

进一步有中心极限定理
\begin{equation}\label{eq:fold-bernoulli-p-clt-restated}
\frac{G_m-m\gamma(p)}{\sqrt m}\ \Rightarrow\ \mathcal N\!\bigl(0,\sigma^2(p)\bigr),
\end{equation}
其中
\begin{equation}\label{eq:fold-bernoulli-p-variance-closed}
\sigma^2(p)=
\frac{p^2(1-p)\bigl(21 p^5-6 p^4+14 p^3-36 p^2+7 p+9\bigr)}
{(p+1)^3(p^2-p+1)^3}.
\end{equation}
\leanverified{paper\_fold\_bernoulli\_p\_pressure\_quartic\_clt\_variance}
\leanverified{paper\_fold\_bernoulli\_p\_pressure\_quartic\_seeds}
\end{theorem}
\begin{proof}
\eqref{eq:fold-bernoulli-p-pressure-quartic-restated} 与
\eqref{eq:fold-bernoulli-p-variance-closed}
与定理 \ref{thm:fold-gauge-anomaly-bernoulli-p-laws} 的四次特征方程与方差率闭式完全等价；
\eqref{eq:fold-bernoulli-p-clt-restated} 则由同一定理的 CLT 结论给出。证毕。
\end{proof}

\begin{proposition}[偏置拆解 Jordan 共振与协方差纯指数混合]\label{prop:fold-bernoulli-p-jordan-resonance-only-half}
矩阵 \(P(p)\) 的谱为
\begin{equation}\label{eq:fold-bernoulli-p-spectrum}
\mathrm{Spec}(P(p))
=\left\{1,\ -p,\ \pm\sqrt{p(1-p)}\right\}.
\end{equation}
仅当 \(p=\tfrac12\) 时发生非主特征值碰撞
\(
-p=-\sqrt{p(1-p)}
\)，并可能出现非半单 Jordan 块；
该点确有 Jordan 退化（见命题 \ref{prop:fold-gauge-anomaly-bernoulli-p-spectrum}）。

对任意 \(p\in(0,1)\setminus\{\tfrac12\}\)，特征值互异，\(P(p)\) 在 \(\RR\) 上可对角化。
因此对任意平方可积边观测
\(
H_t=H(S_t,S_{t+1})
\)
（平稳初值）存在常数 \(C_H(p)<\infty\) 使
\begin{equation}\label{eq:fold-bernoulli-p-covariance-exp-bound}
\bigl|\mathrm{Cov}(H_0,H_k)\bigr|
\le C_H(p)\,\Lambda(p)^k,\qquad
\Lambda(p):=\max\!\left\{p,\sqrt{p(1-p)}\right\}<1.
\end{equation}
特别地，对 \(g_t=\mathbf 1_{\{X_t\neq Y_t\}}\)，当 \(p\neq\tfrac12\) 时协方差无任何 \(k\) 多项式前因子。
\end{proposition}
\leanverified{paper\_fold\_bernoulli\_p\_jordan\_resonance\_only\_half}
\begin{proof}
\eqref{eq:fold-bernoulli-p-spectrum}
由命题 \ref{prop:fold-gauge-anomaly-bernoulli-p-spectrum} 直接成立。
当 \(p\neq\tfrac12\) 时可对角化，谱分解给出
\eqref{eq:fold-bernoulli-p-covariance-exp-bound}。
当 \(p=\tfrac12\) 时，\(-\tfrac12\) 二重根的几何重数为 \(1\)，故出现 Jordan 修正。证毕。
\end{proof}

\begin{theorem}[规范差密度的大偏差端点指数]\label{thm:fold-bernoulli-p-endpoint-ldp-restated}
定义
\[
I_p(s):=\sup_{\theta\in\RR}\bigl(s\theta-P_G(\theta;p)\bigr).
\]
则端点指数具有闭式
\begin{equation}\label{eq:fold-bernoulli-p-I1}
I_p(1)=\frac13\log\frac{1}{p^2(1-p)},
\end{equation}
\begin{equation}\label{eq:fold-bernoulli-p-I0}
I_p(0)=
-\log\!\left(\frac{(1-p)+\sqrt{(1-p)(1+3p)}}{2}\right).
\end{equation}
等价地，
\begin{equation}\label{eq:fold-bernoulli-p-endpoint-prob}
\PP_p(G_m=m)=\exp\bigl(-mI_p(1)+o(m)\bigr),\qquad
\PP_p(G_m=0)=\exp\bigl(-mI_p(0)+o(m)\bigr).
\end{equation}
\leanverified{paper\_fold\_bernoulli\_p\_endpoint\_ldp\_restated}
\end{theorem}
\begin{proof}
\eqref{eq:fold-bernoulli-p-I1}--\eqref{eq:fold-bernoulli-p-I0}
由命题 \ref{prop:fold-gauge-anomaly-bernoulli-p-endpoint-ldp} 直接给出；
\eqref{eq:fold-bernoulli-p-endpoint-prob} 是端点点质量的指数等价表述。证毕。
\end{proof}

\begin{theorem}[平稳态端点事件的有限长闭式]\label{thm:fold-bernoulli-p-endpoint-exact-finite}
在定理 \ref{thm:fold-bernoulli-p-parry-pressure-chain} 的平稳分布 \(\pi(p)\) 下，记
\[
\lambda_\pm(p):=\frac{q\pm \sqrt{q(1+3p)}}{2},
\qquad
t(p):=p^2q.
\]
则成立以下有限长精确闭式。

\medskip
\noindent\textbf{零失配事件。}
对任意 \(m\ge 2\) 有
\begin{equation}\label{eq:fold-bernoulli-p-stationary-zero-mismatch-exact}
\PP_{\pi(p)}(G_m=0)=C_+(p)\,\lambda_+(p)^m+C_-(p)\,\lambda_-(p)^m,
\end{equation}
其中
\[
S_0(p):=\pi_a+\pi_d,\qquad
S_1(p):=\pi_a\,q(1+p)+\pi_d,
\]
并且
\begin{equation}\label{eq:fold-bernoulli-p-stationary-zero-mismatch-Cpm}
C_+(p):=\frac{S_1(p)-\lambda_-(p)S_0(p)}{\lambda_+(p)-\lambda_-(p)},\qquad
C_-(p):=\frac{\lambda_+(p)S_0(p)-S_1(p)}{\lambda_+(p)-\lambda_-(p)}.
\end{equation}
此外，
\begin{equation}\label{eq:fold-bernoulli-p-stationary-zero-mismatch-m01}
\PP_{\pi(p)}(G_0=0)=1,\qquad
\PP_{\pi(p)}(G_1=0)=\frac{1-3p^2+3p^3}{1+p^3}.
\end{equation}

\medskip
\noindent\textbf{满失配事件。}
对任意 \(m\ge 1\)，写 \(m=3k+r\)（\(k\ge 0\)，\(r\in\{0,1,2\}\)），则
\begin{equation}\label{eq:fold-bernoulli-p-stationary-full-mismatch-exact}
\PP_{\pi(p)}(G_m=m)=t(p)^k\times
\begin{cases}
\dfrac{1-p+2p^2}{1+p^3}, & r=0,\\[1.0ex]
\dfrac{p^2(3-2p)}{1+p^3}, & r=1,\\[1.0ex]
\dfrac{p^2(1-p)(2+p^2)}{1+p^3}, & r=2.
\end{cases}
\end{equation}
\leanverified{paper\_fold\_bernoulli\_p\_endpoint\_exact\_finite}
\end{theorem}
\begin{proof}
由边势 \eqref{eq:fold-gauge-anomaly-bernoulli-p-edge-potential} 可知零失配等价于在前 \(m\) 步内从未经过失配边 \((a,b)\)、\((b,c)\)、\((c,a)\)。
当 \(m\ge 2\) 时，只要在时刻 \(<m\) 进入 \(b\) 则下一步必经 \(b\to c\) 产生失配，故零失配路径不可能在 \(\{0,\dots,m-1\}\) 进入 \(b\)，从而也不可能进入 \(c\)（唯一入边为 \(b\to c\)）。
因此在 \(m\ge 2\) 的零失配约束下，链在区间 \(\{0,\dots,m\}\) 上被迫限制在 \(\{a,d\}\)，并仅使用匹配边 \(a\to a\)、\(a\to d\)、\(d\to a\)。
令 \(\pi_{ad}:=(\pi_a,\pi_d)^{\mathsf T}\)，并记 \(\mathbf 1:=(1,1)^{\mathsf T}\)。
对应的 \(2\times 2\) 转移矩阵（状态次序 \((a,d)\)）为
\[
P_{ad}:=
\begin{pmatrix}
q & pq\\
1 & 0
\end{pmatrix},
\]
从而
\[
\PP_{\pi(p)}(G_m=0)=\pi_{ad}^{\mathsf T}P_{ad}^{\,m}\mathbf 1\qquad(m\ge 2).
\]
矩阵 \(P_{ad}\) 的特征根正是 \(\lambda_\pm(p)\)，故存在常数 \(C_\pm(p)\) 使得
\(\PP_{\pi(p)}(G_m=0)=C_+\lambda_+^m+C_-\lambda_-^m\)。
由 \(m=0,1\) 的两点插值（即 \(C_++C_-=S_0\) 与 \(C_+\lambda_++C_-\lambda_-=S_1\)）得到 \eqref{eq:fold-bernoulli-p-stationary-zero-mismatch-Cpm}，并推出 \eqref{eq:fold-bernoulli-p-stationary-zero-mismatch-exact}。
对 \(m=1\)，除 \(\{a,d\}\) 的贡献 \(S_1\) 外，还存在从 \(c\) 走匹配边 \(c\to b\) 的路径，其概率为 \(\pi_c\cdot p=p^3/(1+p^3)\)，从而得到 \eqref{eq:fold-bernoulli-p-stationary-zero-mismatch-m01}。

对满失配事件，若每一步均为失配边，则轨道必须沿唯一三循环 \(a\to b\to c\to a\) 运行。
三条失配边的概率依次为 \(a\to b:p^2\)、\(b\to c:1\)、\(c\to a:q\)，每三步乘积为 \(t(p)=p^2q\)。
当 \(m=3k+r\) 时，从起点处于 \(\{a,b,c\}\) 的三种情形逐一写出首 \(r\) 步因子并乘以 \(t(p)^k\)，再以平稳权 \(\pi_a,\pi_b,\pi_c\) 加权求和即可得到 \eqref{eq:fold-bernoulli-p-stationary-full-mismatch-exact}。证毕。
\end{proof}

\begin{corollary}[对称基线的根式闭式与相位振荡]\label{cor:fold-bernoulli-half-endpoints-oscillation}
取 \(p=\tfrac12\)（故 \(q=\tfrac12\)，\(\varphi=\frac{1+\sqrt5}{2}\)）。
则对任意 \(m\ge 2\)，零失配概率具有根式闭式
\begin{equation}\label{eq:fold-bernoulli-half-zero-mismatch}
\PP_{\pi(1/2)}(G_m=0)=\frac{25+11\sqrt5}{90}\left(\frac{\varphi}{2}\right)^m
+\frac{25-11\sqrt5}{90}\left(\frac{1-\sqrt5}{4}\right)^m,
\end{equation}
其中第二项为严格的奇偶振荡修正。
并且对任意 \(m\ge 1\)，写 \(m=3k+r\) 则满失配概率满足
\begin{equation}\label{eq:fold-bernoulli-half-full-mismatch}
\PP_{\pi(1/2)}(G_m=m)=\left(\frac18\right)^k\times
\begin{cases}
\dfrac{8}{9}, & r=0,\\[0.8ex]
\dfrac{4}{9}, & r=1,\\[0.8ex]
\dfrac{1}{4}, & r=2.
\end{cases}
\end{equation}
\leanverified{paper\_fold\_bernoulli\_half\_full\_mismatch\_seeds}
\end{corollary}
\begin{proof}
将 \(p=q=\tfrac12\) 代入定理 \ref{thm:fold-bernoulli-p-endpoint-exact-finite}。
此时 \(\lambda_+(1/2)=\varphi/2\)、\(\lambda_-(1/2)=(1-\sqrt5)/4<0\)，并且由
\eqref{eq:fold-bernoulli-p-stationary-zero-mismatch-Cpm} 直接化简得到 \eqref{eq:fold-bernoulli-half-zero-mismatch}；
\eqref{eq:fold-bernoulli-half-full-mismatch} 则由 \eqref{eq:fold-bernoulli-p-stationary-full-mismatch-exact} 代入并化简得到。证毕。
\end{proof}

\begin{corollary}[黄金根偏置处的等指数衰减与模 \(3\) 比值极限]\label{cor:fold-bernoulli-phi-half-ratio-limits}
取
\[
p=\frac{\varphi}{2}=\frac{1+\sqrt5}{4},\qquad q=1-p=\frac{3-\sqrt5}{4}.
\]
则 \(\lambda_+(p)=\tfrac12\) 且 \(t(p)=\tfrac18\)，从而
\[
\PP_{\pi(p)}(G_m=0)=C_+(p)\,2^{-m}+O\!\left(|\lambda_-(p)|^m\right),\qquad
\PP_{\pi(p)}(G_{3k+r}=3k+r)=d_r(p)\,2^{-(3k+r)}
\]
（其中 \(C_+(p)\) 与 \(d_r(p)\) 由定理 \ref{thm:fold-bernoulli-p-endpoint-exact-finite} 的闭式给出）。
并且存在严格的模 \(3\) 比值极限：
\begin{equation}\label{eq:fold-bernoulli-phi-half-ratio-limits}
\lim_{k\to\infty}\frac{\PP_{\pi(p)}(G_{3k}=3k)}{\PP_{\pi(p)}(G_{3k}=0)}=6(\sqrt5-1),\qquad
\lim_{k\to\infty}\frac{\PP_{\pi(p)}(G_{3k+1}=3k+1)}{\PP_{\pi(p)}(G_{3k+1}=0)}=4\sqrt5,\qquad
\lim_{k\to\infty}\frac{\PP_{\pi(p)}(G_{3k+2}=3k+2)}{\PP_{\pi(p)}(G_{3k+2}=0)}=\frac{9\sqrt5-7}{2}.
\end{equation}
\end{corollary}
\leanverified{paper\_fold\_bernoulli\_phi\_half\_ratio\_limits}
\begin{proof}
由定理 \ref{thm:fold-bernoulli-p-endpoint-exact-finite} 的零失配闭式，
\(\PP_{\pi(p)}(G_m=0)=C_+\lambda_+^m+C_-\lambda_-^m\)。
在所取参数下 \(\lambda_+(p)=1/2\) 且 \(|\lambda_-(p)|<1/2\)，从而
\(\PP_{\pi(p)}(G_m=0)=C_+(p)\,2^{-m}+O(|\lambda_-(p)|^m)\)。
另一方面，由 \eqref{eq:fold-bernoulli-p-stationary-full-mismatch-exact} 得
\(\PP_{\pi(p)}(G_{3k+r}=3k+r)=d_r(p)\,2^{-(3k+r)}\)（\(d_r\) 为显式代数常数）。
将两式相除并令 \(k\to\infty\) 消去误差项即得比值极限；在该 \(p\) 下代入闭式并化简得到 \eqref{eq:fold-bernoulli-phi-half-ratio-limits}。证毕。
\end{proof}

\begin{proposition}[失配密度的全局最大化与平方最大值恒等式]\label{prop:fold-bernoulli-p-gamma-global-max}
函数
\[
\gamma(p)=\frac{p^2(3-2p)}{1+p^3}\qquad(p\in(0,1))
\]
具有唯一全局极大点 \(p_\star\in(0,1)\)，其由
\begin{equation}\label{eq:fold-bernoulli-p-gamma-root}
p_\star^3+2p_\star-2=0
\end{equation}
唯一刻画。
并且极大值满足精确恒等式
\begin{equation}\label{eq:fold-bernoulli-p-gamma-max-square}
\max_{p\in(0,1)}\gamma(p)=\gamma(p_\star)=p_\star^2.
\end{equation}
\leanverified{paper\_fold\_bernoulli\_p\_gamma\_global\_max}
\end{proposition}
\begin{proof}
命题
\ref{prop:fold-gauge-anomaly-bernoulli-p-cumulant-denominator-law}
之后的推论 \ref{cor:fold-gauge-anomaly-bernoulli-p-density-unimodal}
已给出 \eqref{eq:fold-bernoulli-p-gamma-root} 与
\eqref{eq:fold-bernoulli-p-gamma-max-square}。证毕。
\end{proof}

\begin{theorem}[Fold 输出密度的偏置衰减、饱和值与唯一拐点]\label{thm:fold-bernoulli-p-output-density-concavity-switch}
令
\[
q(p):=\PP_p(Y=1)=\frac{p(p^3-p+1)}{1+p^3},\qquad p\in[0,1].
\]
则成立：
\begin{enumerate}
  \item 严格单调：\(q\) 在 \([0,1]\) 上严格递增；
  \item 偏置衰减与饱和界：
  \[
  0<q(p)<p\quad(p\in(0,1)),\qquad
  q(p)\le \frac12\quad(p\in[0,1]),
  \]
  且 \(q(p)=\tfrac12\) 当且仅当 \(p=1\)；
  \item 唯一拐点：存在唯一 \(p_{\mathrm{inf}}\in(0,1)\) 使
  \begin{equation}\label{eq:fold-bernoulli-p-output-inflection}
  p_{\mathrm{inf}}^3=\frac{7-3\sqrt5}{2},
  \end{equation}
  并且 \(p<p_{\mathrm{inf}}\) 时 \(q\) 严格凹，\(p>p_{\mathrm{inf}}\) 时 \(q\) 严格凸。
\end{enumerate}
\leanverified{paper\_fold\_bernoulli\_p\_output\_density\_concavity\_seeds}
\leanverified{paper\_fold\_bernoulli\_p\_output\_density\_concavity\_switch}
\end{theorem}
\begin{proof}
由有理函数微分得
\[
q'(p)=\frac{p^6+p^4+2p^3-2p+1}{(p+1)^2(p^2-p+1)^2}>0,
\]
故严格递增。

再由
\[
p-q(p)=\frac{p^2}{1+p^3}>0\quad(p\in(0,1))
\]
得 \(0<q(p)<p\)。
并且
\[
q(p)\le \frac12
\iff
2p(p^3-p+1)\le 1+p^3
\iff
(p-1)(2p^3+p^2-p+1)\le 0.
\]
由于 \(2p^3+p^2-p+1>0\) 在 \([0,1]\) 上恒成立，得 \(q(p)\le \tfrac12\)，且仅 \(p=1\) 取等。

二阶导数
\[
q''(p)=
-\frac{2\,(p^6-7p^3+1)}{(p+1)^3(p^2-p+1)^3}.
\]
令 \(t=p^3\)，则符号由 \(t^2-7t+1\) 决定。
在 \((0,1)\) 内仅有小根
\(
t_\ast=\frac{7-3\sqrt5}{2}
\)，对应唯一拐点
\eqref{eq:fold-bernoulli-p-output-inflection}，并给出凹凸区间。证毕。
\end{proof}

\leanverified{paper\_fold\_bernoulli\_p\_pressure\_quartic\_package}

\endinput

\paragraph{偏置族中的递归相关证书、状态消元曲线与分歧判别式}
以下结论在 Bernoulli\((p)\) 基线（\(p\in(0,1)\)）下，沿用
定理 \ref{thm:fold-bernoulli-p-parry-pressure-chain} 与
定理 \ref{thm:fold-bernoulli-p-bitpair-gamma-q-closed}
的记号：平稳四状态链 \((S_t)\)、转移核 \(P(p)\)、平稳分布 \(\pi(p)\)，以及边可观测量
\[
g_t:=\mathbf 1_{\{(S_t,S_{t+1})\in\{(a,b),(b,c),(c,a)\}\}}.
\]

\begin{theorem}[失配指示过程的全自协方差闭式与临界 Jordan 修正]\label{thm:fold-bernoulli-p-full-autocovariance-jordan}
设
\[
r(p):=p(1-p),
\]
并定义
\[
\alpha(p):=
-\frac{p^{3}(p^{2}-3p+1)^{2}}{(p+1)^{2}(2p-1)(p^{2}-p+1)},
\]
\[
\beta(p):=
\frac{p^{2}(1-p)\bigl(p^{5}+4p^{4}-4p^{3}-5p^{2}+7p-2\bigr)}
{(p+1)(2p-1)(p^{2}-p+1)^{2}},
\]
\[
\delta(p):=
\frac{p^{2}(1-p)^{2}\bigl(p^{5}-p^{4}-6p^{3}+3p^{2}+2p-1\bigr)}
{(p+1)(2p-1)(p^{2}-p+1)^{2}}.
\]
当 \(p\neq \tfrac12\) 时，对任意整数 \(m\ge 0\) 与 \(m\ge 1\) 分别有
\[
\operatorname{Cov}(g_0,g_{2m+1})
=\alpha(p)\,p^{2m}+\beta(p)\,r(p)^m,
\]
\[
\operatorname{Cov}(g_0,g_{2m})
=-\alpha(p)\,p^{2m-1}+\delta(p)\,r(p)^{m-1}.
\]
在临界点 \(p=\tfrac12\) 处，对任意 \(k\ge 1\) 成立
\[
\operatorname{Cov}(g_0,g_k)
=2^{-(k-1)}
\left[
\frac1{16}
+(-1)^{k-1}\!\left(-\frac{13}{1296}-\frac{k-1}{216}\right)
\right],
\]
其中 \((k-1)2^{-k}\) 项对应 \(-\tfrac12\) 处的 Jordan 退化贡献。
\leanverified{paper\_fold\_bernoulli\_p\_jordan\_resonance\_seeds}
\leanverified{paper\_fold\_bernoulli\_p\_jordan\_resonance\_package}
\leanverified{paper\_fold\_bernoulli\_p\_full\_autocovariance\_jordan}
\end{theorem}
\begin{proof}
第一部分由推论
\ref{cor:fold-gauge-anomaly-bernoulli-p-covariance-explicit-even-odd}
重参数化得到：其偶、奇协方差闭式与上式等价。
第二部分即定理 \ref{thm:fold-gauge-anomaly-cov-closed} 的协方差闭式重写。
两式在 \(p\to\tfrac12\) 极限下由特征值并合
\(
-p\to-\sqrt{p(1-p)}
\)
产生线性多项式修正。证毕。
\end{proof}

\begin{theorem}[有限谱隙 Jordan 临界点的双奇偶缩放极限]\label{thm:fold-bernoulli-p-jordan-critical-double-parity-scaling}
沿用推论 \ref{cor:fold-gauge-anomaly-bernoulli-p-covariance-explicit-even-odd} 的记号。固定 \(\xi\in\RR\)，并定义
\[
p_m:=\frac12+\frac{\xi}{m}\qquad (m\ge 1).
\]
则当 \(m\to\infty\) 时，
\[
\lim_{m\to\infty}\frac{4^m}{m}\,c_{2m+1}(p_m)
=-\frac{e^{4\xi}-1}{432\,\xi},
\]
\[
\lim_{m\to\infty}\frac{4^m}{m}\,c_{2m}(p_m)
=\frac{e^{4\xi}-1}{216\,\xi},
\]
其中右端在 \(\xi=0\) 处按连续延拓分别取值 \(-\frac1{108}\) 与 \(\frac1{54}\)。
\end{theorem}
\begin{proof}
先设 \(\xi\neq 0\)，并记
\[
\varepsilon:=\frac{\xi}{m},
\qquad
p_m=\frac12+\varepsilon.
\]
由推论 \ref{cor:fold-gauge-anomaly-bernoulli-p-covariance-explicit-even-odd}，
\[
c_{2m+1}(p_m)=A(p_m)(-p_m)^{2m}+\bigl(p_m(1-p_m)\bigr)^m C_0(p_m),
\]
\[
c_{2m}(p_m)=A(p_m)(-p_m)^{2m-1}+\bigl(p_m(1-p_m)\bigr)^m C_1(p_m).
\]
把 \(A,C_0,C_1\) 在 \(p=\frac12+\varepsilon\) 处作 Laurent 展开，可得
\[
A\!\left(\frac12+\varepsilon\right)
=-\frac{1}{432\,\varepsilon}+O(1),
\]
\[
C_0\!\left(\frac12+\varepsilon\right)
=\frac{1}{432\,\varepsilon}+O(1),
\qquad
C_1\!\left(\frac12+\varepsilon\right)
=-\frac{1}{216\,\varepsilon}+O(1).
\]
另一方面，
\[
p_m^{2m}
=\left(\frac12+\frac{\xi}{m}\right)^{2m}
=4^{-m}e^{4\xi}\bigl(1+O(m^{-1})\bigr),
\]
\[
p_m^{2m-1}
=2\cdot 4^{-m}e^{4\xi}\bigl(1+O(m^{-1})\bigr),
\]
并且
\[
\bigl(p_m(1-p_m)\bigr)^m
=\left(\frac14-\frac{\xi^2}{m^2}\right)^m
=4^{-m}\bigl(1+O(m^{-1})\bigr).
\]
代回上式并用 \(\varepsilon=\xi/m\)，得到
\[
c_{2m+1}(p_m)
=4^{-m}\left[
-\frac{m}{432\,\xi}e^{4\xi}
+\frac{m}{432\,\xi}
+O(1)
\right],
\]
\[
c_{2m}(p_m)
=4^{-m}\left[
\frac{2m}{432\,\xi}e^{4\xi}
-\frac{m}{216\,\xi}
+O(1)
\right].
\]
于是
\[
\frac{4^m}{m}\,c_{2m+1}(p_m)
=-\frac{e^{4\xi}-1}{432\,\xi}+o(1),
\qquad
\frac{4^m}{m}\,c_{2m}(p_m)
=\frac{e^{4\xi}-1}{216\,\xi}+o(1).
\]
这就证明了 \(\xi\neq 0\) 的情形。

当 \(\xi=0\) 时，\(p_m=\tfrac12\) 对所有 \(m\) 恒成立。由定理 \ref{thm:fold-bernoulli-p-full-autocovariance-jordan} 的临界闭式直接得到
\[
c_{2m+1}\!\left(\frac12\right)
=4^{-m}\left(-\frac{m}{108}+O(1)\right),
\qquad
c_{2m}\!\left(\frac12\right)
=4^{-m}\left(\frac{m}{54}+O(1)\right),
\]
故极限分别为 \(-\frac1{108}\) 与 \(\frac1{54}\)。另一方面
\[
\lim_{\xi\to 0}\frac{e^{4\xi}-1}{\xi}=4,
\]
故这些值恰与连续延拓一致。证毕。
\end{proof}

\begin{remark}[非半单而不闭隙的临界窗口]\label{rem:fold-bernoulli-p-jordan-critical-finite-gap}
由命题 \ref{prop:fold-gauge-anomaly-bernoulli-p-spectrum}，在 \(p=\tfrac12\) 处次主谱模始终等于 \(\frac12\)，故谱隙保持为 \(\frac12\) 而并未闭合。上一定理表明真正的临界性来自 \(-p\) 与 \(-\sqrt{p(1-p)}\) 的非半单并合：在 \(p=\tfrac12+\xi/m\) 的窗口中，两支贡献的简单极点与指数因子发生精确抵消，最后留下单一缩放轮廓
\[
\Xi(\xi):=\frac{e^{4\xi}-1}{\xi}.
\]
因此该窗口反映的是固定谱隙下的 Jordan 临界，而非由谱隙塌缩驱动的软边缘临界机制。
\end{remark}


\begin{theorem}[自协方差生成函数的完全有理性与临界二重极点]\label{thm:fold-bernoulli-p-autocovariance-generating-rational}
\leanverified{paper\_fold\_bernoulli\_p\_autocovariance\_generating\_rational}
令
\[
s_n:=\operatorname{Cov}(g_0,g_{n+1}),\qquad n\ge 0,
\]
并定义生成函数
\[
\mathcal S_p(z):=\sum_{n\ge 0}s_n z^n.
\]
则
\[
\mathcal S_p(z)=
\frac{s_0+(s_1+p s_0)z+\bigl(s_2+p s_1-p(1-p)s_0\bigr)z^2}
{(1+p z)\bigl(1-p(1-p)z^2\bigr)}.
\]
其中分母在 \(p\neq\tfrac12\) 时分解为一线性因子与一二次因子；在 \(p=\tfrac12\) 时发生并合
\[
(1+p z)\bigl(1-p(1-p)z^2\bigr)\Big|_{p=1/2}
=(1+z/2)^2(1-z/2),
\]
故 \(\mathcal S_{1/2}(z)\) 在 \(z=-2\) 处出现二重极点。
\end{theorem}
\begin{proof}
由命题 \ref{prop:fold-gauge-anomaly-bernoulli-p-covariance-recurrence}，
\[
s_{n+3}
=-p\,s_{n+2}+p(1-p)\,s_{n+1}+p^2(1-p)\,s_n\qquad(n\ge 0).
\]
对上式乘以 \(z^n\) 并求和，整理初值项即得所示有理表达。
当 \(p=\tfrac12\) 时分母因子在 \(-2\) 合并为二重因子，因此出现二重极点。证毕。
\end{proof}

\begin{theorem}[输出密度--失配密度状态方程与消元曲线]\label{thm:fold-bernoulli-p-output-mismatch-elimination-curve}
记
\[
\mathsf q(p):=\PP(Y_t=1),\qquad
\gamma(p):=\PP(X_t\neq Y_t).
\]
则有状态方程
\[
\gamma(p)=(3-2p)\,\bigl(p-\mathsf q(p)\bigr).
\]
此外，点 \((\mathsf q(p),\gamma(p))\) 落在平面四次代数曲线
\[
\begin{aligned}
0={}&\gamma^{4}+9\gamma^{3}\mathsf q+4\gamma^{3}
+27\gamma^{2}\mathsf q^{2}+18\gamma^{2}\mathsf q-23\gamma^{2}\\
&\quad+35\gamma\,\mathsf q^{3}+72\gamma\,\mathsf q^{2}-54\gamma\,\mathsf q+23\gamma
+70\mathsf q^{3}-69\mathsf q^{2}.
\end{aligned}
\]
\end{theorem}
\begin{proof}
由定理 \ref{thm:fold-bernoulli-p-bitpair-gamma-q-closed}，
\[
\mathsf q(p)=\frac{p(p^3-p+1)}{1+p^3},\qquad
\gamma(p)=\frac{p^2(3-2p)}{1+p^3},
\]
并且
\[
p-\mathsf q(p)=\frac{p^2}{1+p^3},
\]
故状态方程立即成立。
对 \((\mathsf q,\gamma)\) 的多项式关系，取
\[
\mathsf q(1+p^3)-p(p^3-p+1)=0,\qquad
\gamma(1+p^3)-p^2(3-2p)=0
\]
对参数 \(p\) 进行 Sylvester 结式消元，得到所示四次曲线。
把 \((\mathsf q,\gamma)=(\mathsf q(p),\gamma(p))\) 代回可得恒等为零。证毕。
\end{proof}
\leanverified{paper\_fold\_bernoulli\_p\_output\_mismatch\_elimination\_curve}

\begin{theorem}[失配指示过程的严格三阶 Markov 性与完整转移核]\label{thm:fold-bernoulli-p-mismatch-strict-third-order-markov}
记 \(G_t:=g_t\in\{0,1\}\)。
则在平稳分布下，\((G_t)\) 为三阶 Markov 过程：转移核
\[
\PP(G_{t+3}=1\mid G_t,G_{t+1},G_{t+2})
\]
仅依赖三位历史，且对允许的三比特上下文有闭式
\[
\PP(1\mid 001)=1,\qquad \PP(1\mid 010)=1,
\]
\[
\PP(1\mid 000)=\frac{p^{2}(2p+1)}{(1+p)(1+2p-p^{2})},
\]
\[
\PP(1\mid 011)=\frac{2p^{3}-p^{2}-p+1}{1+p-p^{2}},
\qquad
\PP(1\mid 100)=\frac{p^{2}}{1+p-p^{2}},
\]
\[
\PP(1\mid 101)=\frac{(1-p)(1+p-p^{2})}{1+p-2p^{2}+p^{3}},
\qquad
\PP(1\mid 110)=p,
\]
\[
\PP(1\mid 111)=\frac{p^{2}(3-2p)}{1-p+2p^{2}}.
\]
并且该过程不是二阶 Markov：对 \(p\in(0,1)\) 有
\[
\PP(1\mid 000)-\PP(1\mid 100)
=-\frac{p^5}{(1+p)(1+2p-p^2)(1+p-p^2)}\neq 0.
\]
\leanverified{paper\_fold\_bernoulli\_p\_mismatch\_strict\_third\_order\_markov}
\end{theorem}
\begin{proof}
令
\[
(B_1)_{ij}:=P_{ij}\,\mathbf 1_{\{g(i,j)=1\}},\qquad B_0:=P-B_1.
\]
对任意字 \(w=w_1\cdots w_n\in\{0,1\}^n\)，有标准矩阵表示
\[
\PP(G_0\cdots G_{n-1}=w)
=\pi^\top B_{w_1}\cdots B_{w_n}\mathbf 1.
\]
故
\[
\PP(G_3=d\mid G_0G_1G_2=abc)
=\frac{\PP(abcd)}{\PP(abc)}
\]
可由有限次有理运算显式计算，得到上述八个上下文公式。
其中 \((001)\) 与 \((010)\) 的后继为确定性，是由不可实现四字模式的分子为零导致。
最后由
\(
\PP(1\mid 000)\neq \PP(1\mid 100)
\)
可知同后缀 \((00)\) 不能决定下一位，故二阶 Markov 性不成立。证毕。
\end{proof}

\begin{theorem}[再生分解与周期二元概率母函数的严格闭式]\label{thm:fold-bernoulli-p-regeneration-bivariate-pgf}
沿用四态链 \((S_t)_{t\in\ZZ}\) 与边势 \(g_t\) 的记号，并令
\[
\tau_0:=0,\qquad \tau_{n+1}:=\inf\{t>\tau_n:\ S_t=d\}\qquad(n\ge 0)
\]
为状态 \(d\) 的相继返回时刻。定义第 \(n\) 个再生周期的长度与异常计数
\[
L_n:=\tau_{n+1}-\tau_n,\qquad
X_n:=\sum_{t=\tau_n}^{\tau_{n+1}-1} g_t .
\]
则在条件于 \(S_{\tau_n}=d\) 的再生版本下（尤其在初始 \(S_0=d\) 下），序列 \(\bigl(L_n,X_n\bigr)_{n\ge 0}\) 为独立同分布。
其二元概率母函数在 \(|z|\) 充分小时满足严格闭式
\begin{equation}\label{eq:fold-bernoulli-p-regeneration-bivariate-pgf}
\Phi_p(u,z):=\EE\!\left[u^{X_0}z^{L_0}\mid S_0=d\right]
=\frac{p(1-p)\,z^2\bigl(1-pu z^2\bigr)}
{\bigl(1-(1-p)z\bigr)\bigl(1-pu z^2\bigr)-p^2(1-p)\,u^3 z^3}.
\end{equation}
此外，存在相互独立的随机变量
\[
M\sim \mathrm{Geom}(1-p),\qquad
K_0,K_1,\dots\ \text{i.i.d.}\ \mathrm{Geom}(p),\qquad
N_1,N_2,\dots\ \text{i.i.d.}\ \mathrm{Geom}(1-p),
\]
其中 \(\mathrm{Geom}(q)\) 取值于 \(\NN_0\) 且 \(\PP(G=k)=(1-q)^k q\)，使得同分布意义下
\begin{equation}\label{eq:fold-bernoulli-p-regeneration-decomposition}
X_0=\sum_{i=1}^{M}\bigl(3+N_i\bigr),\qquad
L_0=2+\sum_{j=0}^{M}K_j+3M+2\sum_{i=1}^{M}N_i .
\end{equation}
\end{theorem}
\begin{proof}
由核 \(P(p)\) 的结构，\(d\to a\) 为确定转移且进入 \(d\) 的唯一方式为 \(a\to d\)；
故 \(d\) 为再生状态。
由强马尔可夫性质，\(\{(\tau_{n+1}-\tau_n,(g_t)_{\tau_n\le t<\tau_{n+1}})\}_{n\ge 0}\) 在条件于 \(S_{\tau_n}=d\) 下为 i.i.d.；
从而 \((L_n,X_n)\) 亦为 i.i.d.。

从 \(d\) 出发第一步必为 \(d\to a\) 且不贡献异常。
在 \(a\) 处，每一步以概率 \(1-p\) 自环、以概率 \(p\) 离开 \(a\)，故“离开 \(a\) 前自环次数”服从 \(\mathrm{Geom}(p)\)，记为 \(K\)。
条件于“离开 \(a\)”，离开方向满足
\[
\PP(a\to d\mid a\ \text{离开})=1-p,\qquad
\PP(a\to b\mid a\ \text{离开})=p,
\]
因此“在最终 \(a\to d\) 前发生的 \(a\to b\) 次数”服从 \(\mathrm{Geom}(1-p)\)，记为 \(M\)。
一旦发生 \(a\to b\)，随后 \(b\to c\) 为确定转移并贡献一次异常；
到达 \(c\) 后，以概率 \(p\) 走 \(c\to b\)（不异常），以概率 \(1-p\) 走 \(c\to a\)（异常并回到 \(a\)）。
于是“在首次 \(c\to a\) 前 \(c\to b\) 的次数”服从 \(\mathrm{Geom}(1-p)\)，记为 \(N\)。
该段从 \(a\) 离开并返回 \(a\) 的异常增量为 \(3+N\)，长度增量为 \(3+2N\)。
对每次返回 \(a\) 之后再次离开之前的自环次数引入独立拷贝 \(K_0,K_1,\dots\)，并对每次 \(a\to b\) 回路引入独立拷贝 \(N_1,N_2,\dots\)，即可得到分解式 \eqref{eq:fold-bernoulli-p-regeneration-decomposition}。
独立性来自强马尔可夫性质与几何分布的无记忆性。

将分解写为生成函数。
记
\[
\EE[z^K]=\frac{p}{1-(1-p)z},
\qquad
\EE\!\left[(u z^2)^N\right]=\frac{1-p}{1-pu z^2}.
\]
并对 \(M\) 的几何和求和，即得到 \eqref{eq:fold-bernoulli-p-regeneration-bivariate-pgf}。证毕。
\end{proof}
\leanverified{paper\_fold\_bernoulli\_p\_regeneration\_bivariate\_pgf}

\begin{theorem}[再生周期三元组 \texorpdfstring{$(M,X_0,L_0)$}{(M,X0,L0)} 的显式联合分布]\label{thm:fold-bernoulli-p-regeneration-triple-joint-law}
沿用定理 \ref{thm:fold-bernoulli-p-regeneration-bivariate-pgf} 的设定。
记 \(p\in(0,1)\)、\(q:=1-p\)，并令几何分解变量
\[
M\sim \mathrm{Geom}(1-p),\qquad
K_0,K_1,\dots\ \text{i.i.d.}\ \mathrm{Geom}(p),\qquad
N_1,N_2,\dots\ \text{i.i.d.}\ \mathrm{Geom}(1-p)
\]
相互独立，满足 \eqref{eq:fold-bernoulli-p-regeneration-decomposition}。
则对任意整数 \(\ell\ge 2\) 有
\[
\PP(M=0,\ X_0=0,\ L_0=\ell\mid S_0=d)
=p\,q^{\ell-1}.
\]
并且对任意整数 \(m\ge 1\)、\(k\ge 3m\)、\(\ell\ge 2+2k-3m\)，令
\[
t:=\ell-2-2k+3m\ (\ge 0),
\]
则
\[
\PP(M=m,\ X_0=k,\ L_0=\ell\mid S_0=d)
=\binom{k-2m-1}{m-1}\binom{t+m}{m}\,
p^{k-m+1}q^{t+m+1},
\]
并约定二项式系数在参数不合法时取值为 \(0\)。

从而对任意整数 \(k\ge 0\)、\(\ell\ge 2\)，有有限求和闭式
\[
\PP(X_0=k,\ L_0=\ell\mid S_0=d)
=\mathbf 1_{\{k=0\}}\,p\,q^{\ell-1}
+\sum_{m\ge 1}
\binom{k-2m-1}{m-1}\binom{\ell-2-2k+4m}{m}\,
p^{k-m+1}q^{\ell-1-2k+4m},
\]
其中求和仅在同时满足 \(k\ge 3m\) 与 \(\ell\ge 2+2k-3m\) 时贡献非零项，因而该和为有限和。
\end{theorem}
\leanverified{paper\_fold\_bernoulli\_p\_regeneration\_triple\_joint\_law}
\begin{proof}
对 \(m\ge 1\) 固定 \(M=m\)，记
\[
S_N:=\sum_{i=1}^{m}N_i,\qquad S_K:=\sum_{j=0}^{m}K_j.
\]
由 \eqref{eq:fold-bernoulli-p-regeneration-decomposition}，
\[
X_0=3m+S_N,\qquad
L_0=2+S_K+3m+2S_N.
\]
因此事件 \(\{M=m,X_0=k,L_0=\ell\}\) 等价于
\(
S_N=k-3m
\) 与
\(
S_K=\ell-2-2k+3m=t
\) 同时成立。
由于 \(N_i\sim\mathrm{Geom}(1-p)\)、\(K_j\sim\mathrm{Geom}(p)\) 取值于 \(\NN_0\) 且独立，
几何和服从负二项分布：
\[
\PP(S_N=s)=\binom{s+m-1}{m-1}\,q^m p^{s}\qquad(s\ge 0),
\]
\[
\PP(S_K=t)=\binom{t+m}{m}\,p^{m+1} q^{t}\qquad(t\ge 0).
\]
再乘上
\(
\PP(M=m)=p^{m}q
\)
并用独立性，即得 \(m\ge 1\) 的联合闭式。

当 \(m=0\) 时，\(X_0=0\) 且 \(L_0=2+K_0\)。于是
\[
\PP(M=0,X_0=0,L_0=\ell)
=\PP(M=0)\,\PP(K_0=\ell-2)
=q\cdot p q^{\ell-2}
=p q^{\ell-1}.
\]
最后，对 \(m\) 求和得到 \((X_0,L_0)\) 的有限和闭式。证毕。
\end{proof}

\begin{theorem}[给定 \texorpdfstring{$(M,X_0)$}{(M,X0)} 后 \texorpdfstring{$L_0$}{L0} 的条件分布]\label{thm:fold-bernoulli-p-regeneration-conditional-length-shifted-nb}
在定理 \ref{thm:fold-bernoulli-p-regeneration-triple-joint-law} 的设定下，
给定 \(M=m\) 与 \(X_0=k\)（其中 \(m=0\Rightarrow k=0\)，且 \(m\ge 1\Rightarrow k\ge 3m\)），有分布恒等式
\[
L_0\ \stackrel{d}{=}\ 2+2k-3m+S_K,
\]
其中 \(S_K\) 为 \(m+1\) 个独立 \(\mathrm{Geom}(p)\) 变量之和。
等价地，对任意整数 \(\ell\ge 2+2k-3m\) 有
\[
\PP(L_0=\ell\mid M=m,X_0=k)
=\binom{\ell-2-2k+4m}{m}\,q^{\ell-2-2k+3m}\,p^{m+1}.
\]
并且条件矩满足
\[
\EE[L_0\mid M=m,X_0=k]=2+2k-3m+\frac{(m+1)q}{p},
\qquad
\Var(L_0\mid M=m,X_0=k)=\frac{(m+1)q}{p^2}.
\]
\end{theorem}
\begin{proof}
由 \eqref{eq:fold-bernoulli-p-regeneration-decomposition}，
\[
L_0=2+\sum_{j=0}^{m}K_j+3m+2\sum_{i=1}^{m}N_i.
\]
在给定 \(M=m\) 与 \(X_0=k\) 时，\(\sum_{i=1}^{m}N_i=k-3m\) 被确定，从而
\[
L_0=2+\sum_{j=0}^{m}K_j+3m+2(k-3m)=2+2k-3m+S_K.
\]
由于 \((K_j)_{0\le j\le m}\) 与 \((N_i)_{1\le i\le m}\) 相互独立，条件下的 \(S_K\) 仍为 \(m+1\) 个独立 \(\mathrm{Geom}(p)\) 变量之和，故给出负二项质量函数与条件矩。证毕。
\end{proof}

\begin{theorem}[周期长度 \texorpdfstring{$L_0$}{L0} 的概率母函数、三阶递推与指数尾常数]\label{thm:fold-bernoulli-p-regeneration-length-pgf-recurrence-tail}
\leanverified{paper\_fold\_bernoulli\_p\_regeneration\_length\_pgf\_recurrence\_tail}
在定理 \ref{thm:fold-bernoulli-p-regeneration-bivariate-pgf} 的设定下，令
\[
f_\ell:=\PP(L_0=\ell\mid S_0=d)\qquad(\ell\in\NN).
\]
则其概率母函数 \(F_L(z):=\sum_{\ell\ge 0} f_\ell z^\ell\) 为有理函数，且满足
\[
F_L(z)=\EE[z^{L_0}\mid S_0=d]
=\Phi_p(1,z)
=\frac{p q\,z^2(1-p z^2)}{1-qz-pz^2+p q^2 z^3}.
\]
因此 \((f_\ell)\) 在 \(\ell\ge 5\) 上满足齐次三阶线性递推
\[
f_\ell=q f_{\ell-1}+p f_{\ell-2}-p q^2 f_{\ell-3}.
\]
并且分母三次多项式
\[
D_L(z):=1-qz-pz^2+p q^2 z^3
\]
之判别式为
\[
\Disc_z\!\bigl(D_L(z)\bigr)
=-p\bigl(31p^5-110p^4+147p^3-96p^2+28p-4\bigr),
\]
其在 \(p\in(0,1)\) 上恒不为零，因而 \(D_L\) 无重根退化并且三个根两两不同。

令 \(R_L(p)>1\) 为方程 \(D_L(z)=0\) 的最小正实根，则存在常数 \(C_L(p)>0\) 使得
\[
f_\ell\sim C_L(p)\,R_L(p)^{-\ell}\qquad(\ell\to\infty),
\]
其中常数可由留数给出
\[
C_L(p)=\frac{p q\,R_L(p)\bigl(1-p R_L(p)^2\bigr)}{q+2pR_L(p)-3p q^2 R_L(p)^2}.
\]
\end{theorem}
\begin{proof}
由 \eqref{eq:fold-bernoulli-p-regeneration-bivariate-pgf} 令 \(u=1\) 即得
\(
F_L(z)=\Phi_p(1,z)
\)，其分母恰为 \(D_L(z)\)。
将恒等式 \(D_L(z)F_L(z)=p q\,z^2-p^2 q\,z^4\) 比较系数可得从 \(\ell\ge 5\) 起右端系数为零，从而递推齐次化并得到所示三阶递推。

对根的单性，直接计算判别式得到
\(
\Disc_z(D_L(z))=-p(31p^5-110p^4+147p^3-96p^2+28p-4)
\)；
该五次多项式在 \((0,1)\) 内无实根（其唯一实根落在 \((1,\infty)\)），故在 \(p\in(0,1)\) 上判别式恒不为零，从而三根互异。
由于 \(f_\ell\ge 0\)，由 Pringsheim 定理可知收敛半径由正实轴上的最小奇点给出；
对有理型 \(F_L\)，该奇点为 \(D_L\) 的最小正实根 \(R_L(p)>1\)，且由于根为单根，对应于简单极点。
由简单极点的系数抽取公式得到
\[
f_\ell\sim -\frac{N(R_L)}{R_L D_L'(R_L)}\,R_L^{-\ell},
\qquad
N(z):=p q\,z^2(1-pz^2),
\]
并注意
\(
-D_L'(R_L)=q+2pR_L-3p q^2 R_L^2
\)，即可化简得到所示常数 \(C_L(p)\)。证毕。
\end{proof}

\endinput


\begin{corollary}[周期异常计数的显式分布、递推与精确指数尾]\label{cor:fold-bernoulli-p-cycle-mismatch-law-tail}
在定理 \ref{thm:fold-bernoulli-p-regeneration-bivariate-pgf} 的设定下，周期异常计数 \(X_0\) 的一元概率母函数为
\begin{equation}\label{eq:fold-bernoulli-p-cycle-mismatch-pgf}
\EE\!\left[u^{X_0}\mid S_0=d\right]
=\Phi_p(u,1)
=\frac{(1-p)(1-pu)}{1-pu-p(1-p)u^3}.
\end{equation}
令 \(w_k:=\PP(X_0=k\mid S_0=d)\)。
则 \(w_0=1-p\)、\(w_1=w_2=0\)，并且对所有 \(k\ge 3\) 有严格三阶递推
\begin{equation}\label{eq:fold-bernoulli-p-cycle-mismatch-recurrence}
w_k=p\,w_{k-1}+p(1-p)\,w_{k-3}.
\end{equation}
进一步，对任意整数 \(k\ge 3\) 有有限求和闭式
\begin{equation}\label{eq:fold-bernoulli-p-cycle-mismatch-pmf-sum}
w_k
=(1-p)^2 \sum_{r=0}^{\lfloor (k-3)/3\rfloor}
\binom{k-2r-3}{r}\,p^{k-2r-2}\,(1-p)^{r}.
\end{equation}
最后，令 \(r(p)>1\) 为三次方程
\begin{equation}\label{eq:fold-bernoulli-p-cycle-mismatch-tail-root}
p(1-p)r^3+pr-1=0
\end{equation}
在 \((1,\infty)\) 上的唯一实根，则存在常数 \(C(p)>0\) 使
\begin{equation}\label{eq:fold-bernoulli-p-cycle-mismatch-tail-asymp}
w_k \sim C(p)\,r(p)^{-k}\qquad(k\to\infty),
\qquad
C(p)=\frac{(1-p)\bigl(1-pr(p)\bigr)}{p\,r(p)\bigl(1+3(1-p)r(p)^2\bigr)}.
\end{equation}
\leanverified{paper\_fold\_bernoulli\_p\_cycle\_mismatch\_law\_tail}
\end{corollary}
\begin{proof}
式 \eqref{eq:fold-bernoulli-p-cycle-mismatch-pgf} 由 \eqref{eq:fold-bernoulli-p-regeneration-bivariate-pgf} 取 \(z=1\) 立得。
将
\(
\sum_{k\ge 0} w_k u^k
=\frac{(1-p)(1-pu)}{1-pu-p(1-p)u^3}
\)
两边清分母并比较系数，得到 \(w_0=1-p\)、\(w_1=w_2=0\) 与递推 \eqref{eq:fold-bernoulli-p-cycle-mismatch-recurrence}。

为得到 \eqref{eq:fold-bernoulli-p-cycle-mismatch-pmf-sum}，使用恒等式（\(|u|\) 充分小时）
\[
\frac{1}{1-au-bu^3}=\sum_{k\ge0}\left(\sum_{r=0}^{\lfloor k/3\rfloor}\binom{k-2r}{r}a^{k-3r}b^r\right)u^k,
\]
取 \(a=p\)、\(b=p(1-p)\)，并乘以 \((1-p)(1-pu)\) 后整理组合系数差分
\(
\binom{n}{r}-\binom{n-1}{r}=\binom{n-1}{r-1}
\)，即可得到所示有限求和闭式。

对尾部渐近，记分母 \(D(u):=1-pu-p(1-p)u^3\)。
由 \eqref{eq:fold-bernoulli-p-cycle-mismatch-tail-root}，点 \(u=r(p)\) 为 \(D\) 的简单零点，且由于系数非负，\(r(p)\) 为最小模的正奇点。
因此 \(u=r(p)\) 为主极点；由留数型系数渐近可得
\(
w_k\sim -\frac{N(r(p))}{r(p)D'(r(p))}\,r(p)^{-k}
\)
（其中 \(N(u)=(1-p)(1-pu)\)），化简即得常数 \(C(p)\) 的显式形式 \eqref{eq:fold-bernoulli-p-cycle-mismatch-tail-asymp}。证毕。
\end{proof}

\begin{theorem}[失配指示过程的 Shannon 熵率闭式]\label{thm:fold-bernoulli-p-mismatch-entropy-rate}
令 \(\nu_p\) 为在平稳链上由观测 \(g_t\) 诱导的平稳二值过程 \((g_t)_{t\in\ZZ}\) 的分布。
以自然对数计（单位：nats），其 Shannon 熵率存在，并满足显式闭式
\begin{equation}\label{eq:fold-bernoulli-p-mismatch-entropy-rate}
h(\nu_p)
=\lim_{n\to\infty}\frac1n H_{\nu_p}(g_0,g_1,\dots,g_{n-1})
=\frac{p(1-p)}{1+p^3}\Bigl(H_{\mathrm{geom}}(p)+H_{\mathrm{geom}}(1-p)\Bigr),
\end{equation}
其中
\begin{equation}\label{eq:fold-bernoulli-p-geom-entropy}
H_{\mathrm{geom}}(q):=-\log q-\frac{1-q}{q}\log(1-q)\qquad(q\in(0,1))
\end{equation}
为取值 \(\NN_0\) 的 \(\mathrm{Geom}(q)\) 分布的熵。
等价地，
\begin{equation}\label{eq:fold-bernoulli-p-mismatch-entropy-rate-expanded}
h(\nu_p)=\frac{p(1-p)}{1+p^3}\left[
-\log\bigl(p(1-p)\bigr)
-\frac{1-p}{p}\log(1-p)
-\frac{p}{1-p}\log p
\right].
\end{equation}
\leanverified{paper\_fold\_bernoulli\_p\_mismatch\_entropy\_rate}
\end{theorem}
\begin{proof}
由定理 \ref{thm:fold-bernoulli-p-regeneration-bivariate-pgf} 的再生分解，从 \(d\) 出发时，输出序列可分解为 i.i.d.\ 再生周期的拼接。
并且由转移图的确定性（\(d\to a\)、\(b\to c\) 为确定；每次出现 \(11\) 之后在 \(c\) 处由下一位 \(g\in\{0,1\}\) 唯一决定 \(c\to b\) 或 \(c\to a\)），每个再生周期的二值字串与随机变量族
\(
(M,(K_j)_{j=0}^{M},(N_i)_{i=1}^{M})
\)
之间存在一一对应的可测双射；因此周期字串的熵等于这些独立几何变量的总熵。
条件于 \(M\)，有
\[
H(\text{周期字串}\mid M)= (M+1)H_{\mathrm{geom}}(p)+M H_{\mathrm{geom}}(1-p),
\]
从而
\[
H(\text{周期字串})=H(M)+\EE[M+1]H_{\mathrm{geom}}(p)+\EE[M]H_{\mathrm{geom}}(1-p)
=\frac{H_{\mathrm{geom}}(p)+H_{\mathrm{geom}}(1-p)}{1-p},
\]
其中用到 \(M\sim\mathrm{Geom}(1-p)\) 给出 \(\EE[M]=p/(1-p)\) 以及 \(H(M)=H_{\mathrm{geom}}(1-p)\)。
另一方面，由 \eqref{eq:fold-bernoulli-p-regeneration-decomposition} 直接求期望得到
\(
\EE[L_0]=(1+p^3)/(p(1-p)^2)
\)。
对再生拼接体系，熵率等于“单位时间平均信息量”，即
\(
h(\nu_p)=H(\text{周期字串})/\EE[L_0]
\)，代入得到 \eqref{eq:fold-bernoulli-p-mismatch-entropy-rate}，而 \eqref{eq:fold-bernoulli-p-mismatch-entropy-rate-expanded} 由 \eqref{eq:fold-bernoulli-p-geom-entropy} 直接化简得到。证毕。
\end{proof}

\leanverified{paper\_fold\_bernoulli\_p\_finite\_time\_mismatch\_rational\_recurrence}
\begin{theorem}[有限时刻异常计数的有理型时间生成函数与四阶递推]\label{thm:fold-bernoulli-p-finite-time-mismatch-rational-recurrence}
从再生态 \(d\) 出发，令
\[
G_m:=\sum_{t=0}^{m-1} g_t \qquad (m\ge 0),
\qquad
M_m(u):=\EE\!\left[u^{G_m}\mid S_0=d\right].
\]
则其时间变量的常规生成函数为严格有理函数
\begin{equation}\label{eq:fold-bernoulli-p-finite-time-mgf-gf}
\sum_{m\ge0} M_m(u)\,z^m=\frac{N_p(u,z)}{D_p(u,z)},
\end{equation}
其中
\begin{align}
N_p(u,z)&=
1+pz+\bigl(p^2u-pu\bigr)z^2+\bigl(p^3u^3-p^2u^3+p^2u^2-p^2u\bigr)z^3,\label{eq:fold-bernoulli-p-finite-time-mgf-num}\\[2mm]
D_p(u,z)&=
1+(p-1)z+\bigl(p^2-p-pu\bigr)z^2+\bigl(pu-p^2u-p^2u^3+p^3u^3\bigr)z^3+\bigl(p^2u-p^3u\bigr)z^4.\label{eq:fold-bernoulli-p-finite-time-mgf-den}
\end{align}
因此对所有 \(m\ge4\)，序列 \(\bigl(M_m(u)\bigr)_{m\ge0}\) 满足齐次四阶线性递推
\begin{equation}\label{eq:fold-bernoulli-p-finite-time-mgf-recurrence}
M_m(u)=(1-p)M_{m-1}(u)+p(1+u-p)M_{m-2}(u)-p(1-p)u\bigl(1-pu^2\bigr)M_{m-3}(u)-p^2(1-p)u\,M_{m-4}(u),
\end{equation}
并具有显式初值
\begin{equation}\label{eq:fold-bernoulli-p-finite-time-mgf-initial}
M_0(u)=1,\quad M_1(u)=1,\quad M_2(u)=1+p^2(u-1),\quad M_3(u)=1+p^2(u-1)(u+2-p).
\end{equation}
\end{theorem}
\begin{proof}
令倾斜核 \(Q_p(u)\) 为
\[
Q_p(u)_{ij}:=P(p)_{ij}\,u^{g(i,j)} .
\]
则按定义
\[
M_m(u)=e_d^{\mathsf T}Q_p(u)^{m}\mathbf 1,
\qquad
\sum_{m\ge0}M_m(u)z^m=e_d^{\mathsf T}(I-zQ_p(u))^{-1}\mathbf 1,
\]
从而是矩阵有理函数。
对 \(4\times4\) 矩阵 \(I-zQ_p(u)\) 直接计算行列式与伴随矩阵，得到
\eqref{eq:fold-bernoulli-p-finite-time-mgf-gf}--\eqref{eq:fold-bernoulli-p-finite-time-mgf-den}。
有理函数系数序列满足由分母多项式 \eqref{eq:fold-bernoulli-p-finite-time-mgf-den} 给出的线性递推；
且因 \(\deg_z N_p\le3\)，从 \(m\ge4\) 起递推齐次化，即得 \eqref{eq:fold-bernoulli-p-finite-time-mgf-recurrence}。
将 \eqref{eq:fold-bernoulli-p-finite-time-mgf-gf} 在 \(z=0\) 处展开至三阶即可得到初值 \eqref{eq:fold-bernoulli-p-finite-time-mgf-initial}。证毕。
\end{proof}

\begin{theorem}[二维常系数递推闭合：\texorpdfstring{$G_m$}{Gm} 的全概率质量函数可由有限步精确生成]\label{thm:fold-bernoulli-p-finite-time-pmf-2d-recurrence}
\leanverified{paper\_fold\_bernoulli\_p\_finite\_time\_pmf\_2d\_recurrence}
从再生态 \(d\) 出发，令
\[
a_{m,k}:=\PP(G_m=k\mid S_0=d)\qquad(m\ge 0,\ k\ge 0),
\]
并约定当 \(m<0\) 或 \(k<0\) 时取 \(a_{m,k}=0\)。
则对任意 \(m\ge 4\)、\(k\ge 0\) 成立二维齐次递推
\begin{equation}\label{eq:fold-bernoulli-p-finite-time-pmf-2d-recurrence}
\begin{aligned}
a_{m,k}={}&(1-p)\,a_{m-1,k}
+p(1-p)\,a_{m-2,k}
+p\,a_{m-2,k-1}
-p(1-p)\,a_{m-3,k-1}\\
&\quad
+p^{2}(1-p)\,a_{m-3,k-3}
-p^{2}(1-p)\,a_{m-4,k-1}.
\end{aligned}
\end{equation}
其初值由 \eqref{eq:fold-bernoulli-p-finite-time-mgf-initial} 唯一确定，具体为
\[
a_{0,0}=1,\quad a_{1,0}=1,
\]
\[
a_{2,0}=1-p^{2},\quad a_{2,1}=p^{2},
\]
\[
a_{3,0}=1-2p^{2}+p^{3},\quad a_{3,1}=p^{2}-p^{3},\quad a_{3,2}=p^{2}.
\]
因此格点族 \(\{a_{m,k}\}_{m,k}\) 在上述递推与边界条件下被完全刻画，从而 \(G_m\) 的全概率质量函数可由有限步计算精确生成，而无需矩阵幂或谱分解。
\end{theorem}
\begin{proof}
令
\(
M_m(u)=\sum_{k\ge 0}a_{m,k}u^{k}
\)，并定义二元生成函数
\[
F(u,z):=\sum_{m\ge 0}M_m(u)\,z^{m}
=\sum_{m,k\ge 0}a_{m,k}u^{k}z^{m}.
\]
由定理 \ref{thm:fold-bernoulli-p-finite-time-mismatch-rational-recurrence}，
\(
F(u,z)=N_p(u,z)/D_p(u,z)
\)，其中 \(D_p(u,z)\) 由 \eqref{eq:fold-bernoulli-p-finite-time-mgf-den} 给出。
将恒等式 \(D_p(u,z)F(u,z)=N_p(u,z)\) 的两端按 \(z^{m}u^{k}\) 比较系数：
由于 \(\deg_z N_p\le 3\)，当 \(m\ge 4\) 时右端系数恒为零，从而得到二维齐次递推 \eqref{eq:fold-bernoulli-p-finite-time-pmf-2d-recurrence}。
初值由 \(M_0,M_1,M_2,M_3\) 的显式表达展开为 \(u\)-系数即可读出。证毕。
\end{proof}

\paragraph{Fibonacci 低缺陷阴影与极端态三周期}
沿用定理 \ref{thm:fold-bernoulli-p-finite-time-mismatch-rational-recurrence} 的 \(G_m\) 与 \(M_m(u)\)。
在均匀基线 \(p=\tfrac12\) 下，定义清分母计数型生成函数
\[
M_m^{\mathrm f}(u):=2^m M_m(u)\in\ZZ_{\ge 0}[u],
\qquad
a_{m,k}:=[u^k]\,M_m^{\mathrm f}(u)\qquad(m,k\ge 0).
\]
并记 Fibonacci 数列 \((F_n)_{n\ge 0}\) 由 \(F_0=0\)、\(F_1=1\)、\(F_{n+2}=F_{n+1}+F_n\) 定义。

\begin{theorem}[总规范差计数多项式的 Fibonacci 低缺陷阴影]\label{thm:fold-bernoulli-p-fibonacci-low-defect-shadow}
\leanverified{paper\_fold\_bernoulli\_p\_fibonacci\_low\_defect\_shadow}
对一切 \(m\ge 0\)，系数 \(a_{m,k}=[u^k]M_m^{\mathrm f}(u)\) 在固定小缺陷阶上具有如下封闭式：
\[
a_{m,0}=F_{m+2},
\]
\[
a_{m,1}=
\begin{cases}
0,& m=0,1,\\
F_{m-1},& m\ge 2,
\end{cases}
\qquad
a_{m,2}=
\begin{cases}
0,& m=0,1,2,\\
2F_{m-1},& m\ge 3.
\end{cases}
\]
\end{theorem}

\begin{proof}
把定理 \ref{thm:fold-bernoulli-p-finite-time-mismatch-rational-recurrence} 的四阶递推 \eqref{eq:fold-bernoulli-p-finite-time-mgf-recurrence} 取 \(p=\tfrac12\) 并乘以 \(2^m\)，得到对 \(m\ge 4\) 的计数型递推
\begin{equation}\label{eq:fold-bernoulli-p-half-Mf-order4}
M_m^{\mathrm f}(u)=M_{m-1}^{\mathrm f}(u)+(2u+1)M_{m-2}^{\mathrm f}(u)+(u^3-2u)M_{m-3}^{\mathrm f}(u)-2u\,M_{m-4}^{\mathrm f}(u).
\end{equation}
由初值 \eqref{eq:fold-bernoulli-p-finite-time-mgf-initial} 在 \(p=\tfrac12\) 下的展开可得
\[
M_0^{\mathrm f}(u)=1,\qquad M_1^{\mathrm f}(u)=2,\qquad M_2^{\mathrm f}(u)=3+u,\qquad M_3^{\mathrm f}(u)=5+u+2u^2.
\]
对 \eqref{eq:fold-bernoulli-p-half-Mf-order4} 作系数抽取得到对任意 \(k\ge 0\) 的二维递推
\[
\begin{aligned}
a_{m,k}={}&a_{m-1,k}+a_{m-2,k}+2a_{m-2,k-1}+a_{m-3,k-3}-2a_{m-3,k-1}-2a_{m-4,k-1},
\end{aligned}
\]
其中约定 \(a_{m,\ell}=0\) 当 \(\ell<0\)。
取 \(k=0\) 得 \(a_{m,0}=a_{m-1,0}+a_{m-2,0}\)（\(m\ge 4\)），并由上初值读出 \(a_{0,0}=1,a_{1,0}=2,a_{2,0}=3,a_{3,0}=5\)，从而 \(a_{m,0}=F_{m+2}\)。
\par
取 \(k=1\) 得
\[
a_{m,1}=a_{m-1,1}+a_{m-2,1}+2a_{m-2,0}-2a_{m-3,0}-2a_{m-4,0}.
\]
由 \(a_{m,0}\) 的 Fibonacci 递推可知后三项严格抵消，于是 \(a_{m,1}=a_{m-1,1}+a_{m-2,1}\)（\(m\ge 4\)）。
再由初值 \(a_{0,1}=a_{1,1}=0\)、\(a_{2,1}=a_{3,1}=1\)，得到所示封闭式。
\par
同理取 \(k=2\) 得
\[
a_{m,2}=a_{m-1,2}+a_{m-2,2}+2a_{m-2,1}-2a_{m-3,1}-2a_{m-4,1},
\]
而 \(a_{m,1}\) 亦满足 Fibonacci 递推，使后三项抵消。
由初值 \(a_{0,2}=a_{1,2}=a_{2,2}=0\)、\(a_{3,2}=2\) 得 \(a_{m,2}=2F_{m-1}\)（\(m\ge 3\)）。
证毕。
\end{proof}

\begin{proposition}[最大规范差与极端态数的三周期]\label{prop:fold-bernoulli-p-half-max-defect-degree-and-leading-period3}
在 \(p=\tfrac12\) 的均匀基线下，对所有 \(m\ge 2\) 有
\[
\deg M_m^{\mathrm f}(u)=m-1,
\]
等价地 \(G_m\le m-1\) 几乎处处成立。
此外，记顶次系数
\(
\ell_m:=[u^{m-1}]\,M_m^{\mathrm f}(u)
\)，则
\[
\ell_m=
\begin{cases}
1,& m\equiv 2\pmod 3,\\
2,& m\equiv 0,1\pmod 3.
\end{cases}
\]
\end{proposition}
\leanverified{paper\_fold\_bernoulli\_p\_half\_max\_defect\_degree\_and\_leading\_period3}

\begin{proof}
由 \eqref{eq:fold-bernoulli-p-half-Mf-order4} 与初值可得 \(\deg M_2^{\mathrm f}=1\)、\(\deg M_3^{\mathrm f}=2\)。
设 \(m\ge 2\) 且已知 \(\deg M_{m-1}^{\mathrm f}=m-2\)。
则 \eqref{eq:fold-bernoulli-p-half-Mf-order4} 右端仅 \((u^3-2u)M_{m-3}^{\mathrm f}(u)\) 的 \(u^3 M_{m-3}^{\mathrm f}\) 项可贡献次数 \((m-3-1)+3=m-1\)，其余各项次数至多为 \(m-2\)；
从而 \(\deg M_m^{\mathrm f}=m-1\)，归纳完成第一式。
\par
对顶次系数，注意到 \(\deg M_m^{\mathrm f}=m-1\) 后，\eqref{eq:fold-bernoulli-p-half-Mf-order4} 中 \(M_{m-1}^{\mathrm f}\)、\((2u+1)M_{m-2}^{\mathrm f}\) 与 \(-2u\,M_{m-4}^{\mathrm f}\) 的次数均严格小于 \(m-1\)，故顶次项仅来自 \(u^3M_{m-3}^{\mathrm f}(u)\)，从而
\[
\ell_m=\ell_{m-3}\qquad(m\ge 5).
\]
由初值 \(M_2^{\mathrm f}(u)=3+u\)、\(M_3^{\mathrm f}(u)=5+u+2u^2\) 以及由 \eqref{eq:fold-bernoulli-p-half-Mf-order4} 计算得到的
\(
\ell_4=[u^3]M_4^{\mathrm f}(u)=2
\)，得到 \(\ell_2=1,\ell_3=2,\ell_4=2\)，从而 \((\ell_m)\) 以 \(3\) 为周期并取值如上。证毕。
\end{proof}

\endinput


\paragraph{Perron 分支的端点展开与系数增长律}
在均匀基线 \(p=\tfrac12\) 下，四次递推 \eqref{eq:fold-bernoulli-p-half-Mf-order4} 的特征多项式可写为
\begin{equation}\label{eq:fold-bernoulli-half-perron-quartic}
F(\mu,u):=\mu^{4}-\mu^{3}-(2u+1)\mu^{2}+(2u-u^{3})\mu+2u.
\end{equation}
对 \(u>0\)，记 \(\mu(u)\) 为 \(F(\mu,u)=0\) 的 Perron--Frobenius 分支（在实轴上为正实且模最大根）。
并定义压力函数
\[
P_G(\theta):=\log\!\left(\frac{\mu(e^\theta)}{2}\right).
\]

\begin{theorem}[低温三次激活律：主导谱支在 \texorpdfstring{$u=0$}{u=0} 的前三阶消失]\label{thm:fold-bernoulli-half-low-temp-cubic-activation}
记黄金分割数 \(\varphi=(1+\sqrt5)/2\)。
则 \(\mu(u)\) 在 \(u=0\) 的邻域解析，并具有展开
\[
\mu(u)=\varphi+\Bigl(\tfrac12-\tfrac{\sqrt5}{10}\Bigr)u^3+\Bigl(2-\tfrac{4\sqrt5}{5}\Bigr)u^4+\Bigl(10-\tfrac{22\sqrt5}{5}\Bigr)u^5+\Bigl(\tfrac{105}{2}-\tfrac{1173\sqrt5}{50}\Bigr)u^6+O(u^7).
\]
特别地，
\[
\mu'(0)=\mu''(0)=0,\qquad \mu^{(3)}(0)=3-\tfrac{3\sqrt5}{5}>0,
\]
从而当 \(\theta\to-\infty\)（即 \(u=e^\theta\to0\)）时有严格三次激活主项
\[
P_G(\theta)=\log(\varphi/2)+\Bigl(-\tfrac12+\tfrac{3\sqrt5}{10}\Bigr)e^{3\theta}+O(e^{4\theta}).
\]
\end{theorem}
\leanverified{paper\_fold\_bernoulli\_half\_low\_temp\_cubic\_activation}

\begin{proof}
在 \(u=0\) 时，
\(
F(\mu,0)=\mu^2(\mu^2-\mu-1)
\)，故存在简单根 \(\mu(0)=\varphi\)。
并且
\(
\partial_\mu F(\varphi,0)\neq 0
\)，故由隐函数定理，存在唯一解析分支 \(\mu(u)\) 满足 \(F(\mu(u),u)\equiv 0\) 且 \(\mu(0)=\varphi\)。
将形式幂级数 \(\mu(u)=\varphi+\sum_{n\ge 1}c_nu^n\) 代入并比较系数即可解得 \(c_1=c_2=0\) 以及所列 \(c_3,\dots,c_6\)。
\par
最后由 \(P_G(\theta)=\log(\mu(e^\theta)/2)\) 并对 \(\log\) 作 Taylor 展开，
主导项系数为 \(c_3/\varphi=-\tfrac12+\tfrac{3\sqrt5}{10}\)，得到所示三次激活律。证毕。
\end{proof}

\begin{corollary}[固定缺陷阶的 \texorpdfstring{$\varphi^m m^{\lfloor k/3\rfloor}$}{phi^m m^{floor(k/3)}} 增长律与 \texorpdfstring{$3j$}{3j} 阶首项常数]\label{cor:fold-bernoulli-half-fixed-defect-growth}
保持 \(M_m^{\mathrm f}(u)\) 与 \(a_{m,k}=[u^k]M_m^{\mathrm f}(u)\) 的记号。
则对任意固定 \(k\ge 0\)，存在常数 \(C_k>0\) 使
\[
a_{m,k}=O\!\bigl(\varphi^m m^{\lfloor k/3\rfloor}\bigr)\qquad(m\to\infty).
\]
更强地，对任意固定 \(j\ge 0\)，成立显式首项常数的渐近式
\[
a_{m,3j}
=\frac{5+3\sqrt5}{10}\,\varphi^m\,
\Bigl(\frac{3\sqrt5-5}{10}\Bigr)^{j}\frac{m^{j}}{j!}
+O\!\bigl(\varphi^m m^{j-1}\bigr)\qquad(m\to\infty).
\]
\end{corollary}
\leanverified{paper\_fold\_bernoulli\_half\_fixed\_defect\_growth}

\begin{proof}
对 \(|u|\) 足够小，四阶递推 \eqref{eq:fold-bernoulli-p-half-Mf-order4} 的四个特征根在 \(u=0\) 处分裂为 \(\{\varphi,-\varphi^{-1},0,0\}\)，且 \(\varphi\) 为简单根。
因此存在 \(\varepsilon>0\)、\(\kappa\in(0,\varphi)\) 以及解析振幅 \(C(u)\)（在 \(|u|<\varepsilon\) 上解析）使
\begin{equation}\label{eq:fold-bernoulli-half-Mf-dominant-root-decomposition}
M_m^{\mathrm f}(u)=C(u)\mu(u)^m+O(\kappa^m)
\end{equation}
在 \(|u|<\varepsilon\) 内一致成立。
\par
由定理 \ref{thm:fold-bernoulli-p-fibonacci-low-defect-shadow}，\(M_m^{\mathrm f}(0)=a_{m,0}=F_{m+2}\)，故
\[
C(0)=\lim_{m\to\infty}\frac{M_m^{\mathrm f}(0)}{\mu(0)^m}
=\lim_{m\to\infty}\frac{F_{m+2}}{\varphi^m}
=\frac{\varphi^2}{\sqrt5}
=\frac{5+3\sqrt5}{10}.
\]
并且由定理 \ref{thm:fold-bernoulli-half-low-temp-cubic-activation}，
\[
\mu(u)=\varphi\bigl(1+\delta(u)\bigr),
\qquad
\delta(u)=\delta_3u^3+O(u^4),
\qquad
\delta_3=\frac{3\sqrt5-5}{10}\neq 0.
\]
将 \eqref{eq:fold-bernoulli-half-Mf-dominant-root-decomposition} 中主导项展开为
\[
\mu(u)^m=\varphi^m(1+\delta(u))^m=\varphi^m\sum_{r\ge 0}\binom{m}{r}\delta(u)^r,
\]
注意到 \(\delta(u)=O(u^3)\) 导致 \([u^k]\delta(u)^r=0\) 当 \(3r>k\)，从而对固定 \(k\) 仅 \(r\le \lfloor k/3\rfloor\) 的项可能贡献，
并结合 \(\binom{m}{r}=O(m^r)\) 得到上界 \(a_{m,k}=O(\varphi^m m^{\lfloor k/3\rfloor})\)。
\par
当 \(k=3j\) 时，\(\delta(u)^j\) 的首项为 \(\delta_3^ju^{3j}\)，因而
\[
[u^{3j}]\,\mu(u)^m=\varphi^m\binom{m}{j}\delta_3^j+O(\varphi^m m^{j-1}).
\]
再乘以 \(C(0)\) 并用 \(\binom{m}{j}=m^j/j!+O(m^{j-1})\) 即得所示首项常数渐近式。证毕。
\end{proof}

\begin{theorem}[高温反幂展开：\texorpdfstring{$u\to\infty$}{u->infty} 时 Perron 根的完整主导级数]\label{thm:fold-bernoulli-half-high-temp-inverse-power}
当 \(u\to+\infty\) 时，Perron 分支满足渐近展开
\[
\mu(u)=u+1-\frac{1}{3u^{2}}+\frac{8}{9u^{3}}+O(u^{-4}).
\]
因此当 \(\theta\to+\infty\)（即 \(u=e^\theta\to\infty\)）时有
\[
P_G(\theta)=\log(\mu(e^\theta)/2)
=\theta-\log2+e^{-\theta}-\frac12e^{-2\theta}+\frac{35}{36}e^{-4\theta}+O(e^{-5\theta}),
\]
并且 \(e^{-3\theta}\) 项严格消失。
\end{theorem}

\begin{proof}
设
\(
\mu(u)=u+c_0+c_1u^{-1}+c_2u^{-2}+c_3u^{-3}+O(u^{-4})
\)，代入 \(F(\mu,u)=0\) 并按 \(u\) 的降幂比较系数，可唯一解得
\(
c_0=1,\ c_1=0,\ c_2=-\tfrac13,\ c_3=\tfrac89
\)，从而得到 \(\mu(u)\) 的反幂展开。
\par
再令 \(u=e^\theta\)，写
\[
\mu(e^\theta)=e^\theta\Bigl(1+e^{-\theta}-\tfrac13e^{-3\theta}+\tfrac89e^{-4\theta}+O(e^{-5\theta})\Bigr),
\]
并对 \(\log(1+x)\) 作标准级数展开：
由于 \(x=e^{-\theta}-\tfrac13e^{-3\theta}+\tfrac89e^{-4\theta}+O(e^{-5\theta})\)，可见三次项在
\(
x-\tfrac12x^2+\tfrac13x^3
\)
中严格抵消，得到所示 \(P_G(\theta)\) 的展开。证毕。
\end{proof}

\endinput


\paragraph{固定缺陷阶的 Fibonacci--Jordan 过滤与极点阶跃迁}
均匀基线 \(p=\tfrac12\) 下，计数型母函数 \(M_m^{\mathrm f}(u)\) 满足四阶递推 \eqref{eq:fold-bernoulli-p-half-Mf-order4}，其系数 \(a_{m,k}=[u^k]M_m^{\mathrm f}(u)\) 因此在 \(m\)-方向携带一条由 \((E^2-E-1)\) 生成的 Jordan 型过滤层级；该过滤的阶数每三个缺陷阶跃迁一次，并在生成函数上表现为 \((1-z-z^2)^{-s}\) 型极点阶跃迁。

\begin{theorem}[固定缺陷阶系数的 Fibonacci--Jordan 过滤与有理型生成函数]\label{thm:fold-bernoulli-half-fixed-defect-fibonacci-jordan-filter}
沿用 \(p=\tfrac12\) 的记号：\(M_m^{\mathrm f}(u)=2^m\EE[u^{G_m}\mid S_0=d]\in\ZZ_{\ge 0}[u]\)，
\(
a_{m,k}=[u^k]M_m^{\mathrm f}(u)
\)
（当 \(m<0\) 或 \(k<0\) 时约定 \(a_{m,k}=0\)）。
对每个 \(k\ge 0\) 定义关于 \(m\) 的常规生成函数
\[
A_k(z):=\sum_{m\ge 0}a_{m,k}z^m,
\]
并在序列空间上定义 Fibonacci 差分算子（约定 \(x_m=0\) 当 \(m<0\)）
\[
(\Delta x)_m:=x_m-x_{m-1}-x_{m-2},\qquad \Delta^{(r)}:=\underbrace{\Delta\circ\cdots\circ\Delta}_{r\ \text{次}}.
\]
记
\[
s(k):=\left\lfloor \frac{k}{3}\right\rfloor+1,
\qquad
\rho:=\frac{\sqrt5-1}{2}=\varphi^{-1}.
\]
则成立：
\begin{enumerate}
  \item \textbf{（有理性）}存在多项式 \(P_k(z)\in\ZZ[z]\) 使
  \[
  A_k(z)=\frac{P_k(z)}{(1-z-z^2)^{s(k)}}.
  \]
  \item \textbf{（极点阶的精确性）}\(A_k\) 在 \(z=\rho\) 处具有恰为 \(s(k)\) 的极点；等价地，\(P_k(\rho)\neq 0\)，并且 \(\rho\) 为 \(A_k\) 的收敛半径。
  \item \textbf{（Jordan 递推刻画）}令移位算子 \((Ex)_m:=x_{m+1}\)。
  则存在 \(M(k)\) 使对一切 \(m\ge M(k)\) 有
  \[
  (E^2-E-1)^{s(k)}a_{\bullet,k}(m)=0,
  \]
  并且幂指数 \(s(k)\) 为最小可能者：\((E^2-E-1)^{s(k)-1}a_{\bullet,k}\) 不可能最终恒为零。
\end{enumerate}
\end{theorem}
\begin{proof}
由 \eqref{eq:fold-bernoulli-p-half-Mf-order4} 取系数（见定理 \ref{thm:fold-bernoulli-p-fibonacci-low-defect-shadow} 的证明），对 \(m\ge 4\) 有二维递推
\[
a_{m,k}=a_{m-1,k}+a_{m-2,k}+2a_{m-2,k-1}+a_{m-3,k-3}-2a_{m-3,k-1}-2a_{m-4,k-1}.
\]
移项得
\begin{equation}\label{eq:fold-bernoulli-half-delta-recursion}
(\Delta a_{\bullet,k})_m
=2(\Delta a_{\bullet,k-1})_{m-2}+a_{m-3,k-3}\qquad(m\ge 4).
\end{equation}

\textbf{（有理性）}
对 \(k=0,1,2\)，定理 \ref{thm:fold-bernoulli-p-fibonacci-low-defect-shadow} 给出 \(a_{m,k}\) 的 Fibonacci 闭式，从而 \((\Delta a_{\bullet,k})_m\) 在 \(m\) 上具有有限支撑。
现对 \(k\ge 3\) 作归纳。
对 \eqref{eq:fold-bernoulli-half-delta-recursion} 两端作用 \(\Delta^{(s(k)-1)}\) 并注意 \(s(k)-1=s(k-3)\)，得
\begin{equation}\label{eq:fold-bernoulli-half-delta-s-recursion}
(\Delta^{(s(k))}a_{\bullet,k})_m
=2(\Delta^{(s(k))}a_{\bullet,k-1})_{m-2}+(\Delta^{(s(k)-1)}a_{\bullet,k-3})_{m-3}.
\end{equation}
由归纳假设，\(\Delta^{(s(k))}a_{\bullet,k-1}\) 与 \(\Delta^{(s(k-3))}a_{\bullet,k-3}=\Delta^{(s(k)-1)}a_{\bullet,k-3}\) 在 \(m\) 上均为有限支撑。
故 \eqref{eq:fold-bernoulli-half-delta-s-recursion} 右端为有限支撑序列，从而 \(\Delta^{(s(k))}a_{\bullet,k}\) 亦为有限支撑。

由差分与生成函数的标准对应（在 \(x_m=0\)（\(m<0\)）的约定下）
\[
\sum_{m\ge 0}(\Delta x)_m z^m=(1-z-z^2)\sum_{m\ge 0}x_m z^m,
\]
迭代得到 \((1-z-z^2)^{s(k)}A_k(z)\) 为多项式，因而 \(A_k\) 具有所述有理表达式。

\textbf{（极点阶与最小性）}
由分解式 \eqref{eq:fold-bernoulli-half-Mf-dominant-root-decomposition}，存在 \(\varepsilon>0\) 使当 \(|u|<\varepsilon\) 时
\[
M_m^{\mathrm f}(u)=C(u)\mu(u)^m+O(\kappa^m),
\]
其中 \(\mu(0)=\varphi\) 且 \(C\) 解析。
由定理 \ref{thm:fold-bernoulli-half-low-temp-cubic-activation}，
\[
\mu(u)=\varphi\bigl(1+\delta_3u^3+O(u^4)\bigr),
\qquad
\delta_3=\frac{3\sqrt5-5}{10}>0.
\]
将
\(
C(u)=c_0+c_1u+c_2u^2+O(u^3)
\)
代入，并注意 \(\mu'(0)=\mu''(0)=0\)，则抽取 \(u^0,u^1,u^2\) 系数得
\[
a_{m,0}=c_0\varphi^m+O(\kappa^m),\qquad
a_{m,1}=c_1\varphi^m+O(\kappa^m),\qquad
a_{m,2}=c_2\varphi^m+O(\kappa^m).
\]
另一方面，定理 \ref{thm:fold-bernoulli-p-fibonacci-low-defect-shadow} 给出 \(a_{m,0}=F_{m+2}\) 以及 \(a_{m,1}=F_{m-1}\)（\(m\ge 2\)）、\(a_{m,2}=2F_{m-1}\)（\(m\ge 3\)），故由 \(F_m\sim \varphi^m/\sqrt5\) 得
\[
c_0=\frac{\varphi^2}{\sqrt5},\qquad c_1=\frac{1}{\varphi\sqrt5},\qquad c_2=\frac{2}{\varphi\sqrt5},
\]
从而 \(c_0,c_1,c_2>0\)。
于是对任意分解 \(k=3j+r\)（\(r\in\{0,1,2\}\)），由
\[
\mu(u)^m=\varphi^m(1+\delta_3u^3+O(u^4))^m
=\varphi^m\sum_{t\ge 0}\binom{m}{t}\delta_3^t u^{3t}+O(\varphi^m m^j u^{3j+1})
\]
抽取 \(u^k\) 系数得到
\[
a_{m,k}
=\widetilde C_k\,m^{j}\varphi^m+O(\varphi^m m^{j-1}),
\qquad \widetilde C_k>0,
\]
从而 \(a_{m,k}\) 的 \(m\)-增长阶恰为 \(m^{\lfloor k/3\rfloor}\varphi^m\)。
因此 \(A_k\) 的收敛半径为 \(\rho=\varphi^{-1}\)，并且在 \(z=\rho\) 的主奇点处必出现阶 \(s(k)=j+1\) 的极点；否则将导致 \(a_{m,k}\) 的多项式阶严格低于 \(j\)，矛盾。

\textbf{（Jordan 递推刻画）}
由 \((1-z-z^2)^{s(k)}A_k(z)\in\ZZ[z]\) 等价于 \(\Delta^{(s(k))}a_{\bullet,k}\) 有限支撑，故 \(a_{\bullet,k}\) 从某个起点起满足常系数线性递推，其特征多项式为 \((\lambda^2-\lambda-1)^{s(k)}\)，即 \((E^2-E-1)^{s(k)}\) 最终消去该序列。
最小性由极点阶精确性等价推出：若幂指数更小，则对应极点阶严格更小，与上一步矛盾。证毕。
\end{proof}

\leanverified{paper\_fold\_bernoulli\_half\_fixed\_defect\_fibonacci\_jordan\_filter}

\leanverified{paper\_fold\_bernoulli\_half\_defect\_3\_6\_explicit}
\begin{proposition}[缺陷阶 \(k=3,4,5,6\) 的显式闭式]\label{prop:fold-bernoulli-half-defect-3-6-explicit}
在定理 \ref{thm:fold-bernoulli-half-fixed-defect-fibonacci-jordan-filter} 的记号下，令 \((F_m)_{m\ge 0}\) 为 Fibonacci 数列（\(F_0=0,F_1=1\)）。
则成立（并与计数意义一致）：
\begin{enumerate}
  \item 对一切 \(m\ge 5\)，
  \[
  a_{m,3}
  =\frac{1}{5}(m-26)F_m+\frac{2}{5}(m+20)F_{m-1}.
  \]
  \item 对一切 \(m\ge 7\)，
  \[
  a_{m,4}
  =\frac{1}{5}(m-136)F_m+\frac{2}{5}(m+105)F_{m-1}.
  \]
  \item 对一切 \(m\ge 9\)，
  \[
  a_{m,5}
  =\frac{1}{5}\Bigl(4(m-181)F_m+(1150-2m)F_{m-1}\Bigr).
  \]
  \item 对一切 \(m\ge 11\)，
  \[
  a_{m,6}
  =\frac{1}{50}\Bigl((5m^2+393m-40048)F_m+(-5m^2-619m+64550)F_{m-1}\Bigr).
  \]
\end{enumerate}
\end{proposition}
\begin{proof}
由定理 \ref{thm:fold-bernoulli-half-fixed-defect-fibonacci-jordan-filter}，
当 \(k=3,4,5\) 时 \(s(k)=2\)，故 \((a_{m,k})_{m\ge 0}\) 从某个起点起满足特征多项式 \((\lambda^2-\lambda-1)^2\) 的四阶常系数递推；
于是对充分大 \(m\) 存在 \(\alpha_k,\beta_k,\gamma_k,\delta_k\in\QQ\) 使
\[
a_{m,k}=(\alpha_k+\beta_k m)F_m+(\gamma_k+\delta_k m)F_{m-1}.
\]
同理，\(k=6\) 时 \(s(6)=3\)，故存在次数至多 \(2\) 的 \(p_6(m),q_6(m)\in\QQ[m]\) 使
\(
a_{m,6}=p_6(m)F_m+q_6(m)F_{m-1}
\)
对充分大 \(m\) 成立。

另一方面，\((M_m^{\mathrm f})\) 由 \eqref{eq:fold-bernoulli-p-half-Mf-order4} 与初值唯一确定，
从而每个固定 \(k\) 的有限多个初值 \(a_{m,k}\) 可由有限步递推计算得到。
把这些初值代入上述有限维参数化并解线性方程组，得到所示有理系数。
再由同一常系数递推的唯一性，可知相应闭式在给定阈值之后恒成立。证毕。
\end{proof}

\begin{corollary}[固定缺陷阶的边界大偏差精确型：指数率与多项式阶]\label{cor:fold-bernoulli-half-fixed-defect-boundary-ldp}
\leanverified{paper\_fold\_bernoulli\_half\_fixed\_defect\_boundary\_ldp}
对任意固定 \(k\ge 0\)，记 \(s=s(k)=\lfloor k/3\rfloor+1\)。
则存在常数 \(C_k>0\) 使
\[
\PP(G_m=k\mid S_0=d)=\frac{a_{m,k}}{2^m}
\sim
C_k\,m^{s-1}\Bigl(\frac{\varphi}{2}\Bigr)^{m}
\qquad(m\to\infty),
\]
从而
\[
\log \PP(G_m=k\mid S_0=d)= -m\log\!\Bigl(\frac{2}{\varphi}\Bigr)+(s-1)\log m+O(1).
\]
\end{corollary}
\begin{proof}
由定理 \ref{thm:fold-bernoulli-half-fixed-defect-fibonacci-jordan-filter} 的极点阶精确性与非负系数排除主部抵消，标准奇点分析给出
\(
a_{m,k}\sim \widetilde C_k\,m^{s-1}\varphi^m
\)
（\(\widetilde C_k>0\)）。
再除以 \(2^m\) 即得所示渐近与对数展开。证毕。
\end{proof}

\endinput


\paragraph{再生母函数的行列式接口与周期倾斜：更新方程、周期方程与 Doob 归一化}
沿用定理 \ref{thm:fold-bernoulli-p-regeneration-bivariate-pgf} 与定理 \ref{thm:fold-bernoulli-p-finite-time-mismatch-rational-recurrence} 的记号。
\providecommand{\Den}{\operatorname{Den}}
对 \(u>0\) 与 \(m\ge 0\) 记
\[
M_m(u):=\EE\!\left[u^{G_m}\mid S_0=d\right],
\qquad
F_p(u,z):=\sum_{m\ge 0}M_m(u)\,z^m.
\]

\begin{theorem}[行列式—再生母函数的严格分解与分母几何化]\label{thm:fold-bernoulli-p-renewal-determinant-factorization}
\leanverified{paper\_fold\_bernoulli\_p\_renewal\_determinant\_factorization}
取 \(p\in(0,1)\)，令 \(q:=1-p\)。
对 \(u>0\) 定义倾斜核
\(
Q_p(u)_{ij}:=P(p)_{ij}\,u^{g(i,j)}
\)。
并令 \(\Phi_p(u,z)\) 为定理 \ref{thm:fold-bernoulli-p-regeneration-bivariate-pgf} 给出的再生周期二元母函数：
\[
\Phi_p(u,z)=\EE\!\left[u^{X_0}z^{L_0}\mid S_0=d\right]
=\frac{p q\,z^2\bigl(1-pu z^2\bigr)}{\Den_\Phi(u,z)},
\qquad
\Den_\Phi(u,z):=\bigl(1-qz\bigr)\bigl(1-pu z^2\bigr)-p^2 q\,u^3 z^3.
\]
则成立以下两条严格恒等式。
\begin{enumerate}
  \item \textbf{严格再生分解（无卷积残差项）}：
  存在解析函数 \(\Psi_p(u,z)\)（\(|z|\) 充分小时）使得
  \begin{equation}\label{eq:fold-bernoulli-p-renewal-F-Psi-Phi}
  \boxed{\;
  F_p(u,z)=\frac{\Psi_p(u,z)}{1-\Phi_p(u,z)}\; }.
  \end{equation}
  并且 \(\Psi_p\) 可取为有理函数
  \begin{equation}\label{eq:fold-bernoulli-p-renewal-Psi-explicit}
  \boxed{\;
  \Psi_p(u,z)=1+\frac{z\bigl(1+p^2uz-puz^2+p^2u^2z^2\bigr)}{\Den_\Phi(u,z)}\;}
  \end{equation}
  （与定理 \ref{thm:fold-bernoulli-p-finite-time-mismatch-rational-recurrence} 中的显式分子 \(N_p(u,z)\) 等价）。

  \item \textbf{分母几何化恒等式：}对所有 \((u,z)\) 在收敛域内有
  \begin{equation}\label{eq:fold-bernoulli-p-renewal-den-geometrization}
  \boxed{\;
  1-\Phi_p(u,z)=\frac{\det(I-zQ_p(u))}{\Den_\Phi(u,z)}\;}
  \end{equation}
  从而 \(\det(I-zQ_p(u))\) 恰为 \(1-\Phi_p(u,z)\) 的分子。
  特别地，
  \begin{equation}\label{eq:fold-bernoulli-p-renewal-det-equals-Dp}
  \det(I-zQ_p(u))=D_p(u,z),
  \end{equation}
  其中 \(D_p(u,z)\) 为定理 \ref{thm:fold-bernoulli-p-finite-time-mismatch-rational-recurrence} 的分母多项式 \eqref{eq:fold-bernoulli-p-finite-time-mgf-den}。
\end{enumerate}
\end{theorem}
\begin{proof}
令 \(\tau_1:=\inf\{t>0:S_t=d\}\)。
由强马尔可夫性质，对每个 \(m\ge 0\) 成立严格更新方程
\[
M_m(u)=\EE\!\left[u^{G_m};\,m<\tau_1\mid S_0=d\right]
\;+\;\sum_{\ell=1}^{m}\EE\!\left[u^{X_0};\,L_0=\ell\mid S_0=d\right]\,M_{m-\ell}(u).
\]
对 \(m\) 作生成函数变换得到
\(
F_p(u,z)=\Psi_p(u,z)+\Phi_p(u,z)F_p(u,z)
\)，其中
\(
\Psi_p(u,z):=\sum_{m\ge 0}\EE[u^{G_m};m<\tau_1\mid S_0=d]z^m
\)。
整理即得 \eqref{eq:fold-bernoulli-p-renewal-F-Psi-Phi}。
对 \(\Psi_p\) 的显式表达，可将
\(
F_p(u,z)=N_p(u,z)/D_p(u,z)
\)（定理 \ref{thm:fold-bernoulli-p-finite-time-mismatch-rational-recurrence}）
与
\(
1-\Phi_p(u,z)=\det(I-zQ_p(u))/\Den_\Phi(u,z)
\)
联立，从而取 \(\Psi_p=N_p/\Den_\Phi\) 并化简得到 \eqref{eq:fold-bernoulli-p-renewal-Psi-explicit}。

对 \eqref{eq:fold-bernoulli-p-renewal-den-geometrization}，由 \(\Phi_p(u,z)=\frac{p q\,z^2(1-pu z^2)}{\Den_\Phi(u,z)}\) 直接得
\[
1-\Phi_p(u,z)=\frac{\Den_\Phi(u,z)-p q\,z^2(1-pu z^2)}{\Den_\Phi(u,z)}.
\]
而分子多项式逐项展开恰与 \(\det(I-zQ_p(u))\) 一致；另一方面，由 \(\sum_{m\ge0}M_m(u)z^m=e_d^{\mathsf T}(I-zQ_p(u))^{-1}\mathbf 1\) 可知
\(\det(I-zQ_p(u))\) 与 \(F_p\) 的分母一致，从而得到 \eqref{eq:fold-bernoulli-p-renewal-det-equals-Dp}。证毕。
\end{proof}

\begin{theorem}[压力的周期方程刻画与导数的周期倾斜公式]\label{thm:fold-bernoulli-p-pressure-cycle-equation}
对任意 \(\theta\in\RR\) 令 \(u=e^\theta\)，并记压力函数
\(
P_{G,p}(\theta)=\log\rho(Q_p(u))
\)
（定理 \ref{thm:fold-gauge-anomaly-bernoulli-p-laws}）。
令 \(z(\theta):=e^{-P_{G,p}(\theta)}=\rho(Q_p(u))^{-1}\)。
则：
\begin{enumerate}
  \item \(z(\theta)\) 是方程
  \begin{equation}\label{eq:fold-bernoulli-p-cycle-equation}
  \Phi_p\!\left(e^\theta,z\right)=1
  \end{equation}
  的最小正解；等价地，
  \[
  z(\theta)=\min\{z>0:1-\Phi_p(e^\theta,z)=0\}
  =\min\{z>0:\det(I-zQ_p(e^\theta))=0\}.
  \]
  \item 在 \(\Phi_p(u,z(\theta))\equiv 1\) 的定义域内，隐函数微分给出导数公式
  \begin{equation}\label{eq:fold-bernoulli-p-pressure-derivative-Phi}
  \boxed{\;
  P_{G,p}'(\theta)=
  \frac{u\,\partial_u\Phi_p(u,z)}{z\,\partial_z\Phi_p(u,z)}
  \Bigg|_{(u,z)=(e^\theta,z(\theta))}\; }.
  \end{equation}
  \item 定义单周期倾斜测度 \(\PP_\theta\)（作用于 \((L_0,X_0)\)）为
  \[
  \frac{\dd \PP_\theta}{\dd \PP}\ \propto\ \exp\!\bigl(\theta X_0-P_{G,p}(\theta)\,L_0\bigr)
  \quad\text{并在 }\Phi_p(e^\theta,e^{-P_{G,p}(\theta)})=1\text{ 下归一化}.
  \]
  则有严格识别式
  \begin{equation}\label{eq:fold-bernoulli-p-pressure-derivative-cycle-tilt}
  \boxed{\;
  P_{G,p}'(\theta)=\frac{\EE_\theta[X_0]}{\EE_\theta[L_0]}\; }.
  \end{equation}
  并且二阶导数满足
  \begin{equation}\label{eq:fold-bernoulli-p-pressure-second-derivative-cycle-tilt}
  \boxed{\;
  P_{G,p}''(\theta)=
  \frac{\Var_\theta\!\left(X_0-P_{G,p}'(\theta)L_0\right)}{\EE_\theta[L_0]}>0\; }.
  \end{equation}
\end{enumerate}
\leanverified{paper\_fold\_bernoulli\_p\_pressure\_cycle\_equation}
\end{theorem}
\begin{proof}
由定理 \ref{thm:fold-bernoulli-p-renewal-determinant-factorization} 的 \eqref{eq:fold-bernoulli-p-renewal-den-geometrization}，
\(
1-\Phi_p(u,z)=0
\)
与
\(
\det(I-zQ_p(u))=0
\)
在正实根上等价。
并且由 Pringsheim 定理，\(F_p(u,\cdot)\) 的主奇点为最小正根，因而 \(z(\theta)\) 由 \eqref{eq:fold-bernoulli-p-cycle-equation} 唯一刻画。

对恒等式 \(\Phi_p(u,z(\theta))\equiv 1\) 作微分，注意 \(u'=u\)、\(z'=-zP_{G,p}'(\theta)\)，得到
\(
0=\partial_u\Phi_p\cdot u+\partial_z\Phi_p\cdot z'
\)，整理即得 \eqref{eq:fold-bernoulli-p-pressure-derivative-Phi}。

最后，由 \(\Phi_p(u,z)=\EE[u^{X_0}z^{L_0}\mid S_0=d]\)，可识别
\(
u\partial_u\Phi_p=\EE[X_0u^{X_0}z^{L_0}]
\) 与
\(
z\partial_z\Phi_p=\EE[L_0u^{X_0}z^{L_0}]
\)；
在 \(\Phi_p(u,z)=1\) 的归一化下它们分别是 \(\PP_\theta\) 下的 \(\EE_\theta[X_0]\)、\(\EE_\theta[L_0]\)，从而得到 \eqref{eq:fold-bernoulli-p-pressure-derivative-cycle-tilt}。
再微分并整理为中心化形式即可得 \eqref{eq:fold-bernoulli-p-pressure-second-derivative-cycle-tilt}；严格正性由周期统计的非退化性给出。证毕。
\end{proof}

\input{subsubsec__fold-gauge-anomaly-bernoulli-p-pressure-ldp-compressed-cubic}
\input{subsubsec__fold-gauge-anomaly-bernoulli-p-pressure-curve-moments-local-ldp}

\leanverified{paper\_fold\_bernoulli\_p\_doob\_transform\_closed}
\begin{theorem}[指数倾斜的 Doob 变换在本模型中的闭式形式]\label{thm:fold-bernoulli-p-doob-transform-closed}
取 \(u>0\)，记 \(\rho(u;p):=\rho(Q_p(u))\)。
令 \(r(u)\)、\(\ell(u)\) 分别为 \(Q_p(u)\) 的 Perron 右、左特征向量：
\(
Q_p(u)r(u)=\rho(u;p)r(u)
\)、
\(
\ell(u)^{\mathsf T}Q_p(u)=\rho(u;p)\ell(u)^{\mathsf T}
\)。
在规范化 \(r_c(u)=1\)、\(\ell_a(u)=1\) 下有完全闭式。
\begin{enumerate}
  \item \textbf{右特征向量：}
  \[
  r_b(u)=\frac{u}{\rho},\qquad
  r_d(u)=\frac{\rho^2-pu}{(1-p)u\,\rho^2},\qquad
  r_a(u)=\rho\,r_d(u)=\frac{\rho^2-pu}{(1-p)u\,\rho}.
  \]
  \item \textbf{左特征向量：}
  \[
  \ell_d(u)=\frac{p(1-p)}{\rho},\qquad
  \ell_c(u)=\frac{u^2p^2}{\rho^2-pu},\qquad
  \ell_b(u)=\frac{\rho\,u p^2}{\rho^2-pu}.
  \]
  \item \textbf{Doob \texorpdfstring{$(h)$}{(h)}-变换核：}定义
  \[
  \widetilde P^{(u)}_{ij}:=\frac{1}{\rho(u;p)}\frac{r_j(u)}{r_i(u)}\,Q_p(u)_{ij}.
  \]
  则 \(\widetilde P^{(u)}\) 为同一支持图上的不可约马尔可夫核，其全部非零转移概率为
  \[
  \widetilde P^{(u)}_{d,a}=1,\quad \widetilde P^{(u)}_{b,c}=1,
  \]
  \[
  \widetilde P^{(u)}_{c,b}=\frac{pu}{\rho^2},\quad
  \widetilde P^{(u)}_{c,a}=1-\frac{pu}{\rho^2},
  \]
  \[
  \widetilde P^{(u)}_{a,a}=\frac{1-p}{\rho},\quad
  \widetilde P^{(u)}_{a,d}=\frac{p(1-p)}{\rho^2},\quad
  \widetilde P^{(u)}_{a,b}=\frac{p^2(1-p)\,u^3}{\rho(\rho^2-pu)}.
  \]
  \item \textbf{平稳分布：}平稳分布满足 \(\widetilde\pi^{(u)}_b=\widetilde\pi^{(u)}_c\)，且
  \[
  \widetilde\pi^{(u)}_i\ \propto\ \ell_i(u)\,r_i(u),
  \]
  即可由未归一化权重给出
  \[
  \widetilde\pi^{(u)}_a\ \propto\ \frac{\rho^2-pu}{(1-p)u\,\rho},\qquad
  \widetilde\pi^{(u)}_b=\widetilde\pi^{(u)}_c\ \propto\ \frac{u^2p^2}{\rho^2-pu},\qquad
  \widetilde\pi^{(u)}_d\ \propto\ \frac{p(\rho^2-pu)}{u\,\rho^3}.
  \]
\end{enumerate}
\end{theorem}
\begin{proof}
将 \(Q_p(u)r=\rho r\) 与 \(Q_p(u)^{\mathsf T}\ell=\rho\ell\) 逐行写出。
由确定性边 \(d\to a\)、\(b\to c\) 得
\(
r_a=\rho r_d
\) 与
\(
r_b=(u/\rho)r_c
\)，并由 \(c\)-行解出 \(r_d\)，再由 \(r_c=1\) 得 (1)。
左向量同理由稀疏结构解出并取 \(\ell_a=1\) 得 (2)。
将 (1) 代入 Doob 变换定义并逐项化简得到 (3)。
最后用一般事实 \(\widetilde\pi_i\propto \ell_i r_i\) 得 (4)。证毕。
\end{proof}

\begin{proposition}[周期倾斜下几何分解的闭包性与 Doob 变换一致性]\label{prop:fold-bernoulli-p-cycle-tilt-geometric-closure}
\leanverified{paper\_fold\_bernoulli\_p\_cycle\_tilt\_geometric\_closure}
沿用定理 \ref{thm:fold-bernoulli-p-regeneration-bivariate-pgf} 的几何分解变量
\(
(M,(K_j)_{j\ge 0},(N_i)_{i\ge 1})
\)。
对任意 \((u,z)\) 处于 \(\Phi_p(u,z)\) 的收敛域内，定义周期倾斜测度
\[
\frac{\dd \PP^{(u,z)}}{\dd \PP}\ \propto\ u^{X_0}z^{L_0}.
\]
则该倾斜族在几何分解类中闭包，具体为：在 \(\PP^{(u,z)}\) 下，
\begin{enumerate}
  \item 条件于 \(M=m\)，族 \((K_0,\dots,K_m)\) 仍为独立同分布几何变量，且
  \[
  K_j\sim \mathrm{Geom}\!\bigl(1-(1-p)z\bigr)\qquad(0\le j\le m).
  \]
  \item 条件于 \(M=m\)，族 \((N_1,\dots,N_m)\) 仍为独立同分布几何变量，且
  \[
  N_i\sim \mathrm{Geom}\!\bigl(1-pu z^2\bigr)\qquad(1\le i\le m).
  \]
  \item 随机变量 \(M\) 的边缘分布在倾斜下仍为几何分布：
  \[
  M\sim \mathrm{Geom}\!\bigl(1-q_M(u,z)\bigr),
  \qquad
  q_M(u,z)=\frac{p^2(1-p)\,u^3z^3}{(1-(1-p)z)(1-pu z^2)}.
  \]
  \item 若取 \(z=\rho(u;p)^{-1}\)，则上述倾斜更新与定理 \ref{thm:fold-bernoulli-p-doob-transform-closed} 的 Doob 变换核完全一致：在 \(a\)、\(c\) 两处的更新概率分别等于
  \[
  1-(1-p)z=1-\frac{1-p}{\rho},
  \qquad
  1-pu z^2=1-\frac{pu}{\rho^2},
  \]
  并且 \(q_M(u,z)\) 对应于倾斜链中“失配模块个数”的几何公比。
\end{enumerate}
\end{proposition}
\begin{proof}
由定理 \ref{thm:fold-bernoulli-p-regeneration-bivariate-pgf} 的条件独立结构与分解式
\eqref{eq:fold-bernoulli-p-regeneration-decomposition}，
倾斜权重 \(u^{X_0}z^{L_0}\) 在 \((M,K_j,N_i)\) 上按
\(
z^{\sum K_j}\cdot (u^3z^3)^M\cdot (u z^2)^{\sum N_i}
\)
逐因子分离。
对几何分布施加幂权重仍为几何分布，从而得到 (1)(2)。
对 \(M\) 的边缘分布，把条件矩母函数代入并对 \(m\) 求和可见其概率质量为幂几何型，公比即为 \(q_M(u,z)\)，得到 (3)。
最后令 \(z=\rho^{-1}\) 并对照定理 \ref{thm:fold-bernoulli-p-doob-transform-closed}(3) 的闭式转移概率，即得 (4)。证毕。
\end{proof}

\endinput


\leanverified{paper\_fold\_bernoulli\_p\_pressure\_discriminant\_degree11}
\begin{theorem}[压力分歧判别式的全分解与 \(u=1\) 临界并合条件]\label{thm:fold-bernoulli-p-pressure-discriminant-degree11}
令 \(u=e^\theta\)，并记 \(\lambda\) 为特征值变量。
把特征方程写成四次多项式
\[
\begin{aligned}
F(\lambda;u,p):={}&\lambda^{4}+(p-1)\lambda^{3}
+\bigl(p^{2}-p(u+1)\bigr)\lambda^{2}\\
&\quad+\bigl(p^{3}u^{3}-p^{2}u^{3}-p^{2}u+pu\bigr)\lambda
+p^{2}u(1-p).
\end{aligned}
\]
则其 \(\lambda\)-判别式满足
\[
\Disc_{\lambda}F(\lambda;u,p)=p^{3}u(1-p)\,D(u,p),
\]
其中 \(D(u,p)\in\ZZ[p,u]\) 为关于 \(u\) 的 \(11\) 次多项式：
\[
\begin{aligned}
D(u,p)={}&27p^{8}u^{11}-14p^{8}u^{8}+3p^{8}u^{5}-81p^{7}u^{11}-90p^{7}u^{9}+52p^{7}u^{8}
+170p^{7}u^{6}-11p^{7}u^{5}-68p^{7}u^{3}+12p^{7}\\
&+81p^{6}u^{11}+270p^{6}u^{9}-68p^{6}u^{8}-25p^{6}u^{7}-536p^{6}u^{6}
+14p^{6}u^{5}+36p^{6}u^{4}+252p^{6}u^{3}+24p^{6}u-44p^{6}\\
&-27p^{5}u^{11}-270p^{5}u^{9}+28p^{5}u^{8}-57p^{5}u^{7}+576p^{5}u^{6}
+126p^{5}u^{5}-68p^{5}u^{4}-328p^{5}u^{3}+36p^{5}u^{2}-40p^{5}u+56p^{5}\\
&+90p^{4}u^{9}+6p^{4}u^{8}+189p^{4}u^{7}-204p^{4}u^{6}-213p^{4}u^{5}
+176p^{4}u^{3}-36p^{4}u^{2}-16p^{4}u-24p^{4}\\
&-4p^{3}u^{8}-107p^{3}u^{7}-18p^{3}u^{6}+29p^{3}u^{5}+60p^{3}u^{4}
+20p^{3}u^{3}-16p^{3}u^{2}+48p^{3}u-4p^{3}\\
&+12p^{2}u^{6}+52p^{2}u^{5}-44p^{2}u^{3}-8p^{2}u+4p^{2}
-12pu^{4}-8pu^{3}+12pu^{2}-8pu+4u^{2}.
\end{aligned}
\]
并且
\[
D(1,p)=4(1+p)^2(2p-1)^2(p^2-p+1)^2.
\]
因此在 \(p\in(0,1)\) 内，
\[
u=1\ \text{为分歧点}
\quad\Longleftrightarrow\quad
p=\frac12.
\]
\end{theorem}
\begin{proof}
命题 \ref{prop:fold-gauge-anomaly-bernoulli-p-leyang-branch-u1}
已给出判别式分解结构与 \(u=1\) 判据。
将四次多项式 \(F\) 代入四次判别式闭式并在 \(\ZZ[p,u]\) 中化简，即得上式中 \(D(u,p)\) 的显式系数。
代入 \(u=1\) 并因式分解，得到
\(
D(1,p)=4(1+p)^2(2p-1)^2(p^2-p+1)^2
\)。
由于 \(1+p>0\) 且 \(p^2-p+1>0\) 在 \((0,1)\) 恒正，结论等价于 \(2p-1=0\)。证毕。
\end{proof}

\begin{conjecture}[主导分歧根唯一性与高阶累积量振荡律的参数族延拓]\label{conj:fold-bernoulli-p-dominant-branch-cumulant-oscillation}
对每个 \(p\in(0,1)\setminus\{\tfrac12\}\)，记
\[
\mathcal B(p):=\{u\in\mathbb C: D(u,p)=0\},
\]
其中 \(D(u,p)\) 由定理 \ref{thm:fold-bernoulli-p-pressure-discriminant-degree11} 给出。
猜想：
\begin{enumerate}
  \item \(\mathcal B(p)\) 存在唯一一对共轭根 \(u_\ast(p),\overline{u_\ast(p)}\)，满足
  \[
  |u_\ast(p)|=\min\{|u|:u\in\mathcal B(p)\},\qquad \Im u_\ast(p)\neq 0,
  \]
  且该根为单根；
  \item 取分支 \(\theta_\ast(p)=\log u_\ast(p)\)（\(\Re\theta_\ast(p)<0\)），则失配计数每步归一化累积量 \(\kappa_n(p)\) 在 \(n\to\infty\) 时满足
  \[
  \kappa_n(p)\sim
  C(p)\,n!\,|\theta_\ast(p)|^{-n}
  \cos\!\bigl(n\arg\theta_\ast(p)+\phi(p)\bigr),
  \]
  其中 \(C(p)\neq 0\)、\(\phi(p)\in\mathbb R\) 随 \(p\) 实解析变化。
\end{enumerate}
\end{conjecture}

\endinput


% Arithmetic-geometry layer for the fold gauge-anomaly spectral curve.

\paragraph{谱曲线的算术几何与分歧多项式的同域结构}\label{par:fold-gauge-anomaly-spectral-curve-arithmetic}
命题 \ref{prop:fold-gauge-anomaly-discriminant-factorization} 给出谱四次 \(F(\mu,u)\) 的判别式分解与分支集合，
命题 \ref{prop:fold-gauge-anomaly-rate-curve-param} 给出速率曲线的消元模型 \(R(\alpha,u)=0\) 及其可行分支选择。
基于同一组代数对象，可以进一步抽取出谱曲线、分歧多项式与速率曲线之间的几何与算术对齐关系。
以下条目主要围绕特征多项式 \(F(\mu,u)\) 与分歧十次 \(P_{10}(u)\) 的算术几何展开；
其中涉及的若干群论与密度型推论可视为标准框架下的直接推演，而外部文献的先后性仍需作系统核验。

\begin{proposition}[谱四次曲线的平面模型、亏格与简单分歧]\label{prop:fold-gauge-anomaly-spectral-quartic-geometry}
沿用命题 \ref{prop:fold-gauge-anomaly-discriminant-factorization} 的记号。
令 \(F^{\mathrm{hom}}(\mu,u,w)\) 为 \(F(\mu,u)\) 的总次数齐次化，并记平面曲线
\[
C:=\left\{[\mu:u:w]\in \PP^{2}:\ F^{\mathrm{hom}}(\mu,u,w)=0\right\},
\]
其中
\[
F^{\mathrm{hom}}(\mu,u,w)
=\mu^{4}-\mu^{3}w-2\mu^{2}uw-\mu^{2}w^{2}-\mu u^{3}+2\mu u w^{2}+2u w^{3}.
\]
则 \(C\) 为光滑平面四次曲线，因而其亏格为 \(3\)，并且为非超椭圆曲线。
投影
\[
\pi:C\to\PP^{1},\qquad [\mu:u:w]\mapsto [u:w],
\]
为次数 \(4\) 的覆盖；其分支集合由 \(\mathrm{Disc}_{\mu}(F(\mu,u))=0\) 给出，等价地由
\(
u(u-1)P_{10}(u)=0
\)
支撑。
并且在全部有限分支点处均为单分歧（分歧指数为 \(2\)），而在 \(u=\infty\) 处不分歧。
\end{proposition}
\begin{proof}[证明要点]
判别式分解 \(\mathrm{Disc}_{\mu}(F)=-u(u-1)P_{10}(u)\)（命题 \ref{prop:fold-gauge-anomaly-discriminant-factorization}）表明：有限分歧像恰为
\(\{u=0\}\cup\{u=1\}\cup\{u:\,P_{10}(u)=0\}\)。
由结论 \ref{con:fold-gauge-anomaly-disc-factorization-simple-branch} 的平方因子自由性可知，\(\mathrm{Disc}_{\mu}(F)\) 在每个有限分歧像处均为一次零点；
因此对每个有限分歧值，\(\mu\)-纤维中恰有一对根作二重并合且无更高阶并合，从而分歧指数为 \(2\)。
\par
对 \(u=\infty\) 的分歧性，令 \(t=1/u\) 并作变量替换 \(\nu=t\mu\)。
将方程写成
\[
\nu^{4}-\nu+t(-\nu^{3}-2\nu^{2})+t^{2}(2\nu-\nu^{2})+2t^{3}=0.
\]
在 \(t=0\) 处化为 \(\nu(\nu^{3}-1)=0\)，给出四个互异解 \(\nu=0,1,\omega,\omega^{2}\)；
并且 \(\partial_{\nu}(\nu^{4}-\nu)=4\nu^{3}-1\) 在上述四点均非零，故由隐函数定理得到四条解析分支。
因此 \(\pi^{-1}(\infty)\) 含四个互异点，且以 \(t\) 为局部参数时 \(u=1/t\) 仅呈一阶极点，故 \(u=\infty\) 上方无分歧。
\par
由此 \(\pi\) 的分支值总数为 \(12\)，且全部为简单分歧，故总分歧量为 \(R=12\)。
对正规化 \(X\) 应用 Riemann--Hurwitz 公式得
\(
2g(X)-2=4\cdot(-2)+R=-8+12=4
\)，从而 \(g(X)=3\)。
另一方面，平面四次曲线的算术亏格为 \((4-1)(4-2)/2=3\)，故 \(C\) 的 \(\delta\)-不变量为零，曲线 \(C\) 光滑且与其正规化一致；
平面光滑四次曲线亏格为 \(3\) 且必非超椭圆。证毕。
\end{proof}

% Spectral quartic automorphisms.

\begin{theorem}[\texorpdfstring{$\QQ$}{Q}-自同构群的平凡性]\label{thm:fold-gauge-anomaly-spectral-quartic-autQ-trivial}
令 \(\mathcal C_F\subset\PP^2_{\QQ}\) 为命题 \ref{prop:fold-gauge-anomaly-spectral-quartic-geometry} 的光滑平面四次曲线。
则
\[
\operatorname{Aut}_{\QQ}(\mathcal C_F)=\{1\}.
\]
\leanverified{rat\_pow\_seven\_eq\_one\_iff}
\leanverified{rat\_fixed\_case\_forces\_identity}
\leanverified{paper\_fold\_gauge\_anomaly\_spectral\_quartic\_autQ\_trivial}
\end{theorem}
\begin{proof}
由于 \(\mathcal C_F\) 为非超椭圆的属 \(3\) 曲线，其典范嵌入即为该平面四次模型。
因此任一 \(\QQ\)-自同构由某个 \(A\in\mathrm{PGL}_3(\QQ)\) 在 \(\PP^2\) 上诱导，且存在 \(\kappa\in\QQ^\times\) 使
\[
F^{\mathrm{hom}}\!\bigl(A\cdot(\mu,u,w)\bigr)=\kappa\,F^{\mathrm{hom}}(\mu,u,w).
\]
\par
在直线 \(\{w=0\}\) 上有 \(F^{\mathrm{hom}}(\mu,u,0)=\mu(\mu^3-u^3)\)，故 \(\mathcal C_F\cap\{w=0\}\) 由四点组成。
其中恰有两点为 \(\QQ\)-有理点：
\[
P_0=[0:1:0],\qquad P_1=[1:1:0],
\]
其余两点的 \(\mu/u\) 比值为非平凡三次单位根。
因而 \(A\) 必保持集合 \(\{P_0,P_1\}\)，并由 \(P_0,P_1\) 所张成的唯一直线推出 \(A\) 保持 \(\{w=0\}\)。
\par
\emph{交换情形的排除。}
若 \(A\) 交换 \(P_0\leftrightarrow P_1\)，则可在 \(\mathrm{PGL}_3(\QQ)\) 中取代表使
\[
A=
\begin{pmatrix}
-1&1&c\\
d&1&f\\
0&0&1
\end{pmatrix},
\qquad c,d,f\in\QQ,
\]
并且存在 \(\kappa\in\QQ^\times\) 满足上述齐次恒等式。
将该矩阵代入并比较单项式系数可得：
由 \(u^{3}w\) 项系数得到 \(c=1\)；
由 \(\mu u^{3}\) 项系数得到 \(\kappa=3\)；
由 \(\mu^{3}u\) 项系数得到 \(3d-2=0\)，即 \(d=2/3\)。
另一方面，由 \(\mu^4\) 项系数得到
\(
-3d^{3}+6d^{2}-4d-\kappa+1=0
\)，
代入 \(d=2/3\) 与 \(\kappa=3\) 得 \(-8/9-2\neq 0\)，矛盾。
故交换情形不可能发生。
\par
\emph{固定情形。}
于是 \(A\) 必分别固定 \(P_0,P_1\)，并保持直线 \(\{w=0\}\)。
可取代表矩阵
\[
A=
\begin{pmatrix}
a& b& c\\
d& e& f\\
0& 0& 1
\end{pmatrix}\in\mathrm{GL}_3(\QQ).
\]
由 \(A(P_0)=P_0\) 得 \(b=0\)；
又 \(P_0\) 处切线为 \(\{\mu=0\}\)，故切线在 \(A\) 下保持，从而 \(c=0\)。
于是
\[
A=
\begin{pmatrix}
a& 0& 0\\
d& e& f\\
0& 0& 1
\end{pmatrix}.
\]
将其代入齐次恒等式并比较 \(uw^{3}\)、\(u^{4}\)、\(\mu^{3}u\)、\(\mu^{3}w\) 的系数，得到
\[
f=0,\qquad \kappa=e,\qquad d=0,\qquad e=a^{3},\qquad a e^{2}=1.
\]
从而 \(a^{7}=1\)，在 \(\QQ\) 中只能 \(a=1\)，继而 \(e=1\) 且 \(\kappa=1\)。
故 \(A\) 在 \(\mathrm{PGL}_3(\QQ)\) 中为恒等元，结论成立。证毕。
\end{proof}

\endinput

\begin{proposition}[谱四次曲线 Jacobian 的几何单纯性：不可约 Frobenius 多项式证书]\label{prop:fold-gauge-anomaly-spectral-quartic-jacobian-simple}
\leanverified{paper\_fold\_gauge\_anomaly\_spectral\_quartic\_jacobian\_simple}
沿用命题 \ref{prop:fold-gauge-anomaly-spectral-quartic-geometry} 的光滑平面四次曲线 \(C/\QQ\)，并令 \(J:=\operatorname{Jac}(C)\)。
则 \(J\) 在 \(\overline{\QQ}\) 上单纯。
更具体地，在素数 \(p=7\) 处 \(C\) 具有良好约化，且其约化曲线 \(C/\FF_{7}\) 的局部 \(L\)-多项式满足
\[
\boxed{\ L_{7}(T)=1+4T+9T^{2}+9T^{3}+63T^{4}+196T^{5}+343T^{6}\ }
\]
并在 \(\QQ[T]\) 中不可约。
\end{proposition}
\begin{proof}[证明要点]
脚本 \texttt{scripts/exp\_fold\_gauge\_anomaly\_spectral\_quartic\_jacobian\_L7\_audit.py} 以有限域穷举实现如下核验：
\begin{enumerate}
  \item 在 \(p=7\) 上对齐次模型 \(F^{\mathrm{hom}}(\mu,u,w)=0\) 检验无奇点，从而 \(C\) 在 \(p=7\) 处良好约化。
  \item 记 \(N_n:=\#C(\FF_{7^{n}})\ (n=1,2,3)\)，并由点数恒等式
  \(
  N_n=7^n+1-\sum_{i=1}^{6}\alpha_i^{n}
  \)
  回收幂和 \(S_n:=\sum_{i=1}^{6}\alpha_i^{n}\)；
  再由 Newton 恒等式确定 \(L_7(T)=\prod_{i=1}^{6}(1-\alpha_iT)\) 的系数，并得到上述闭式。
  \item 将 \(L_7(T)\) 模 \(13\) 化后在 \(\FF_{13}[T]\) 中保持不可约，从而 \(L_7(T)\) 在 \(\QQ[T]\) 中不可约。
\end{enumerate}
相应的可引用 TeX 证书片段为：
\input{sections/generated/eq_fold_gauge_anomaly_spectral_quartic_jacobian_L7_audit}
\par
由 Honda--Tate 理论，\(L_7(T)\) 的不可约性蕴含约化雅可比簇 \(J_{\FF_7}\) 在 \(\overline{\FF}_7\) 上为简单阿贝尔三维簇，故其端同态代数无非平凡幂等元。
另一方面，若 \(J_{\overline{\QQ}}\) 可分，则 \(\mathrm{End}^0_{\overline{\QQ}}(J)\) 含非平凡幂等元；由良好约化处端同态专门化映射的单射性（Tate--Faltings 型结果），该幂等元将注入到 \(\mathrm{End}^0_{\overline{\FF}_{7}}(J_{\FF_7})\)，矛盾。
故 \(J\) 在 \(\overline{\QQ}\) 上单纯。证毕。
\end{proof}

\begin{proposition}[第二良素处的局部 \texorpdfstring{$L$}{L}-多项式]\label{prop:fold-gauge-anomaly-spectral-quartic-jacobian-L13}
\leanverified{paper\_fold\_gauge\_anomaly\_spectral\_quartic\_jacobian\_l13}
在素数 \(p=13\) 处，曲线 \(C\) 亦具有良好约化，且其约化曲线 \(C/\FF_{13}\) 的局部 \(L\)-多项式满足
\[
\boxed{\ L_{13}(T)=1+3T+7T^{2}+4T^{3}+91T^{4}+507T^{5}+2197T^{6}\ }
\]
并在 \(\QQ[T]\) 中不可约；其三次系数 \(4\) 不被 \(13\) 整除，从而约化雅可比簇 \(J_{\FF_{13}}\) 为普通阿贝尔三维簇。
\end{proposition}
\begin{proof}[证明要点]
脚本 \texttt{scripts/exp\_fold\_gauge\_anomaly\_spectral\_quartic\_jacobian\_L13\_audit.py} 以有限域穷举核验良好约化与点计数，并由 Newton 恒等式回收 \(L_{13}(T)\) 的闭式；
其在 \(\QQ[T]\) 中的不可约性由模 \(43\) 的不可约性证书给出。
相应可引用的 TeX 片段为：
\input{sections/generated/eq_fold_gauge_anomaly_spectral_quartic_jacobian_L13_audit}
证毕。
\end{proof}

\begin{theorem}[几何端同态环的极小性与绝对单纯性]\label{thm:fold-gauge-anomaly-spectral-quartic-jacobian-endomorphism-Z}
\leanverified{paper\_fold\_gauge\_anomaly\_spectral\_quartic\_jacobian\_endomorphism\_z}
在上述记号下，
\[
\operatorname{End}_{\overline{\QQ}}(J)\cong \ZZ,
\qquad
\text{因而 }J\text{ 在 }\overline{\QQ}\text{ 上绝对单纯。}
\]
\end{theorem}
\begin{proof}
取 \(p\in\{7,13\}\)。
由命题 \ref{prop:fold-gauge-anomaly-spectral-quartic-jacobian-simple} 与命题 \ref{prop:fold-gauge-anomaly-spectral-quartic-jacobian-L13}，\(C\) 在 \(p\) 处良好约化且 \(L_p(T)\) 在 \(\QQ[T]\) 中不可约，并且 \(T^{3}\) 系数不被 \(p\) 整除。
由 Honda--Tate 理论，约化雅可比簇 \(J_{\FF_p}\) 在 \(\overline{\FF}_p\) 上为普通且简单的阿贝尔三维簇，并满足
\[
\operatorname{End}^0_{\overline{\FF}_p}(J_{\FF_p})\cong K_p,
\qquad
K_p:=\QQ(\pi_p),\qquad [K_p:\QQ]=6,
\]
其中 \(\pi_p\) 为 Frobenius 元。
\par
另一方面，良好约化处端同态的专门化给出单射
\[
\operatorname{End}^0_{\overline{\QQ}}(J)\hookrightarrow \operatorname{End}^0_{\overline{\FF}_p}(J_{\FF_p})\cong K_p,
\]
从而
\[
\operatorname{End}^0_{\overline{\QQ}}(J)\subseteq K_7\cap K_{13}.
\]
\par
记
\[
f_7(x):=x^6+4x^5+9x^4+9x^3+63x^2+196x+343,
\qquad
f_{13}(x):=x^6+3x^5+7x^4+4x^3+91x^2+507x+2197,
\]
它们分别为 \(L_7,L_{13}\) 的互反多项式，生成相应的 CM 域 \(K_7,K_{13}\)。
直接计算得
\[
\gcd\!\bigl(\mathrm{Disc}(f_7),\mathrm{Disc}(f_{13})\bigr)=9.
\]
其中 \(\mathrm{Disc}(f_p)\) 为整系数多项式判别式，等于整环阶 \(\ZZ[\pi_p]\) 的判别式，因而被 \(\mathrm{disc}(K_p)\) 所整除。
若 \(K_7\cap K_{13}\neq\QQ\)，则其非平凡子域的域判别式同时整除 \(\mathrm{disc}(K_7)\) 与 \(\mathrm{disc}(K_{13})\)，从而亦整除 \(9\)，故只能为 \(\QQ(\sqrt{-3})\)。
\par
但对素数 \(\ell=41\equiv 2\pmod 3\)，多项式 \(f_7\bmod 41\) 在 \(\FF_{41}[x]\) 中分解型为 \((3,3)\)，从而 \(41\) 在 \(K_7\) 中存在残基次数为 \(3\) 的素因子；
若 \(\QQ(\sqrt{-3})\subseteq K_7\)，则 \(\ell\equiv 2\pmod 3\) 在 \(\QQ(\sqrt{-3})\) 中惰性，故 \(K_7\) 上一切残基次数应为偶数，矛盾。
同理，对 \(\ell=11\equiv 2\pmod 3\)，有 \(f_{13}\bmod 11\) 的分解型为 \((1,1,2,2)\)，从而 \(11\) 在 \(K_{13}\) 中存在残基次数为 \(1\) 的素因子，亦排除 \(\QQ(\sqrt{-3})\subseteq K_{13}\)。
故必有 \(K_7\cap K_{13}=\QQ\)。
\par
于是 \(\operatorname{End}^0_{\overline{\QQ}}(J)\subseteq\QQ\)，从而 \(\operatorname{End}^0_{\overline{\QQ}}(J)=\QQ\) 且 \(\operatorname{End}_{\overline{\QQ}}(J)=\ZZ\)。
最后，若 \(J\) 在 \(\overline{\QQ}\) 上可分，则 \(\operatorname{End}^0_{\overline{\QQ}}(J)\) 含非平凡幂等元，与其等于 \(\QQ\) 矛盾。
故 \(J\) 绝对单纯。证毕。
\end{proof}

\begin{corollary}[低属靶曲线的态射排除与 N\'eron--Severi 秩]\label{cor:fold-gauge-anomaly-spectral-quartic-no-low-genus-maps-ns-rank1}
对任意 \(\overline{\QQ}\)-定义的属 \(1\) 或属 \(2\) 曲线 \(D\)，不存在非恒定态射 \(C\to D\)。
并且 \(\mathrm{NS}(J)\) 的秩为 \(1\)。
\end{corollary}
\leanverified{paper\_fold\_gauge\_anomaly\_spectral\_quartic\_no\_low\_genus\_maps\_ns\_rank1}
\begin{proof}
若存在非恒定态射 \(C\to D\) 且 \(g(D)\in\{1,2\}\)，则诱导非零同态 \(J\to \operatorname{Jac}(D)\)，从而 \(J\) 存在维数 \(1\) 或 \(2\) 的非平凡阿贝尔商，进而在 \(\overline{\QQ}\) 上非单纯，矛盾于定理 \ref{thm:fold-gauge-anomaly-spectral-quartic-jacobian-endomorphism-Z}。
故第一断言成立。
\par
第二断言由标准同一 \(\mathrm{NS}(J)\otimes\QQ\cong \operatorname{End}^0_{\overline{\QQ}}(J)^{\dagger}\)（Rosati 对称端同态）与
\(\operatorname{End}^0_{\overline{\QQ}}(J)=\QQ\) 直接推出。证毕。
\end{proof}

\begin{conjecture}[最大 Mumford--Tate 与 \texorpdfstring{$\ell$}{l}-进满像]\label{conj:fold-gauge-anomaly-spectral-quartic-mumford-tate-gsp6}
设 \(\rho_{\ell}:\mathrm{Gal}(\overline{\QQ}/\QQ)\to \mathrm{GSp}_{6}(\ZZ_\ell)\) 为 \(J\) 的 \(\ell\)-进 Tate 模表示。
预期存在有限集合 \(S\) 使对一切 \(\ell\notin S\)，\(\rho_\ell\) 的像为有限指数子群，且其 Zariski 闭包为 \(\mathrm{GSp}_{6}\)。
等价地，\(J\) 的 Mumford--Tate 群为 \(\mathrm{GSp}_6\)，并呈 \(\mathrm{USp}(6)\) 型的 Sato--Tate 行为。
\end{conjecture}

\endinput


\begin{proposition}[判别式十次因子的满对称伽罗瓦性与判别式分解]\label{prop:fold-gauge-anomaly-P10-galois-discriminant}
令 \(u=e^\theta\)，并记 \(\mu(u):=2\rho(A_\theta)\)。
由命题 \ref{prop:fold-gauge-anomaly-pressure}，\(\mu(u)\) 满足四次方程
\[
P(\mu,u):=\mu^4-\mu^3-2u\mu^2-\mu^2-u^3\mu+2u\mu+2u=0.
\]
其关于 \(\mu\) 的判别式分解为
\[
\mathrm{Disc}_{\mu}P(\mu,u)=-u(u-1)D(u),
\]
其中
\[
\begin{aligned}
D(u)=\;&27u^{10}+27u^9-153u^8-163u^7+793u^6+1061u^5\\
&-832u^4-816u^3+1320u^2-440u+40.
\end{aligned}
\]
并记 \(P_{10}(u):=D(u)\)。
并且 \(D(u)\) 在 \(\QQ\) 上不可约，且
\[
\mathrm{Gal}\bigl(D/\QQ\bigr)\cong S_{10}.
\]
此外，
\[
\boxed{\ \mathrm{Disc}\bigl(D\bigr)=2^{38}3^{24}5^{2}\cdot 1931\cdot 9013^{2}\cdot 34319^{3}\ }.
\]
其平方自由部分为 \(1931\cdot 34319=66269989\)，从而 \(D\) 的分裂域 \(L/\QQ\) 具有唯一的非平凡正规子域
\[
L^{A_{10}}=\QQ\!\left(\sqrt{66269989}\right)
=\QQ\!\left(\sqrt{\mathrm{Disc}(D)}\right).
\]
并且 \(L\) 的最大阿贝尔子扩张恰为上述二次子域（亦即任意阿贝尔正规子扩张均包含于其中）。
\leanverified{paper\_fold\_gauge\_anomaly\_p10\_galois\_discriminant}
\end{proposition}
\begin{proof}[证明要点]
取素数 \(p\nmid 27\cdot\mathrm{Disc}(D)\)。
由 Dedekind 分解型与 Frobenius 循环型的对应，\(D\bmod p\) 的不可约因子次数分解等于某个 Frobenius 元在 \(S_{10}\) 中的循环分解型。
直接计算给出如下三种分解型：
\[
D\bmod 13\ \text{不可约}\ (10),\qquad
D\bmod 7\ \text{分解型为}\ (9,1),\qquad
D\bmod 11\ \text{分解型为}\ (7,3).
\]
其中模 \(13\) 不可约性首先推出 \(D\) 在 \(\QQ\) 上不可约。
因此 \(\mathrm{Gal}(D/\QQ)\) 含一 \(10\)-循环，从而传递且含奇置换；并含一 \((9)(1)\)-型元素，从而该传递子群为本原群。
又由 \((7)(3)\)-型元素可取其三次幂得到一条 \(7\)-循环；Jordan 定理断言：含素数长度 \(p\le 10-3\) 的 \(p\)-循环之本原子群必包含 \(A_{10}\)。
故 \(A_{10}\subseteq \mathrm{Gal}(D/\QQ)\)。
再由群中含奇置换，推出 \(\mathrm{Gal}(D/\QQ)=S_{10}\)。
\par
判别式分解可由 Sylvester--resultant 定义直接计算并在 \(\ZZ\) 中分解得到。
其平方自由部分为 \(1931\cdot 34319=66269989\equiv 1\pmod 4\)，故二次域 \(\QQ(\sqrt{66269989})\) 的域判别式即为 \(66269989\)，并恰在素数 \(1931,34319\) 处分歧。
由于 \(S_{10}\) 的唯一非平凡真正规范子群为 \(A_{10}\)，分裂域 \(L/\QQ\) 的唯一指数 \(2\) 子扩张为 \(L^{A_{10}}\)；
而判别式的平方根生成符号表示的二次子域，故 \(L^{A_{10}}=\QQ(\sqrt{\mathrm{Disc}(D)})=\QQ(\sqrt{66269989})\)。证毕。
\end{proof}

\begin{conclusion}[分歧十次多项式的满对称伽罗瓦群判定]\label{con:fold-gauge-anomaly-p10-galois-s10}
设
\[
P_{10}(u)=27u^{10}+27u^9-153u^8-163u^7+793u^6+1061u^5-832u^4-816u^3+1320u^2-440u+40.
\]
则其分裂域在 \(\QQ\) 上的伽罗瓦群满足
\[
\mathrm{Gal}(P_{10}/\QQ)\cong S_{10}.
\]
并且
\[
\mathrm{Disc}(P_{10})=2^{38}\cdot 3^{24}\cdot 5^{2}\cdot 1931\cdot 9013^{2}\cdot 34319^{3}
\]
非平方，故该群不包含于 \(A_{10}\)。
\end{conclusion}
\begin{proof}[证明要点]
结论 \ref{con:fold-gauge-anomaly-disc-factorization-simple-branch} 已给出 \(P_{10}\) 的不可约性，故伽罗瓦群传递。
命题 \ref{prop:fold-gauge-anomaly-P10-galois-discriminant} 给出模 \(7\) 分解型 \((9,1)\)、模 \(11\) 分解型 \((7,3)\) 与模 \(13\) 不可约型 \((10)\)。
据 Dedekind--Frobenius 对应，群中含 \((9)(1)\) 与 \((7)(3)\) 型元素；由 \((7)(3)\) 的三次幂得到 \(7\)-循环。
传递子群含 \((9)(1)\) 可推出本原性；由 Jordan 定理（\(n=10,p=7\le n-3\)）得 \(A_{10}\subseteq G\)。
再由判别式非平方（命题 \ref{prop:fold-gauge-anomaly-P10-galois-discriminant}）排除 \(G\subseteq A_{10}\)，故 \(G=S_{10}\)。证毕。
\end{proof}

\begin{conclusion}[压力谱四次方程在 \texorpdfstring{$\QQ(u)$}{Q(u)} 上的泛型伽罗瓦群]\label{con:fold-gauge-anomaly-quartic-galois-s4}
令
\[
f(\mu,u):=\mu^4-\mu^3-(2u+1)\mu^2+(2u-u^3)\mu+2u.
\]
则在函数域 \(\QQ(u)\) 上有
\[
\mathrm{Gal}\bigl(f(\mu,u)/\QQ(u)\bigr)\cong S_4.
\]
\end{conclusion}
\begin{proof}[证明要点]
首先，\(u=2\) 特化给出四次多项式
\(
\mu^4-\mu^3-5\mu^2-4\mu+4
\)
在 \(\QQ\) 上不可约，从而泛型不可约。
其次，命题 \ref{prop:fold-gauge-anomaly-cubic-resolvent-discriminant} 给出
\(
\mathrm{Disc}_\mu(f)=-u(u-1)P_{10}(u)
\)
非平方，故伽罗瓦群不包含于 \(A_4\)。
再者，三次预解式
\[
\mathcal R(y,u)=y^3+(1+2u)y^2+(u^3-10u)y-(u^6-4u^4+20u^2+10u)
\]
在 \(u=2\) 处化为
\(
y^3+5y^2-12y-100
\)
且在 \(\QQ\) 上无有理根，故预解式泛型不可约。
四次群分类遂给出 \(S_4\)。
该结论亦可直接由结论 \ref{con:fold-gauge-anomaly-generic-galois-s4} 读取。证毕。
\end{proof}

\begin{conclusion}[谱四次覆盖与 \texorpdfstring{$S_4$}{S4} 伽罗瓦闭包的亏格公式]\label{con:fold-gauge-anomaly-genus-g3-g49}
设 \(\mathcal C_F\) 为谱四次 \(f(\mu,u)=0\) 的光滑模型，\(\pi:\mathcal C_F\to\PP^1_u\) 为投影。
则：
\begin{enumerate}
  \item \(\pi\) 为次数 \(4\) 覆盖，分支值恰为 \(u=0,1\) 与 \(P_{10}(u)=0\) 的十个根，共 \(12\) 个有限分支点，且全部简单分支；
  \item \(g(\mathcal C_F)=3\)；
  \item 其 \(S_4\)-伽罗瓦闭包 \(\mathcal X\to\PP^1_u\)（次数 \(24\)）满足 \(g(\mathcal X)=49\)。
\end{enumerate}
\end{conclusion}
\begin{proof}[证明要点]
第 (1) 条与第 (2) 条见命题 \ref{prop:fold-gauge-anomaly-spectral-quartic-geometry} 及结论 \ref{con:fold-gauge-anomaly-disc-factorization-simple-branch}。
对第 (3) 条，惯性均为换位（阶 \(2\)），每个分支点对 Riemann--Hurwitz 的贡献为 \(24(1-\tfrac12)=12\)。
故
\[
2g(\mathcal X)-2=24(-2)+12\cdot 12=96,
\]
从而 \(g(\mathcal X)=49\)。
亦与结论 \ref{con:fold-gauge-anomaly-s4-hurwitz-tower-genus} 一致。证毕。
\end{proof}

\begin{conclusion}[全纯微分的 Chevalley--Weil 精确分解]\label{con:fold-gauge-anomaly-chevalley-weil-s4}
在结论 \ref{con:fold-gauge-anomaly-genus-g3-g49} 的 \(\mathcal X\) 上，令 \(G=S_4\)。
则
\[
H^0(\mathcal X,\Omega^1)\cong
5\cdot \mathbf{sgn}\ \oplus\ 4\cdot \rho_{2}\ \oplus\ 3\cdot \rho_{3}\ \oplus\ 9\cdot \rho_{3'},
\]
其中 \(\mathbf{sgn}\) 为符号表示，\(\rho_2\) 为 \(2\) 维不可约表示，\(\rho_3\) 为标准 \(3\) 维表示，\(\rho_{3'}=\rho_3\otimes\mathbf{sgn}\)。
\end{conclusion}
\begin{proof}[证明要点]
分支点数 \(r=12\)，且全部惯性由换位生成。
Chevalley--Weil 公式给出
\[
m_\rho=-\dim\rho+\frac r2\bigl(\dim\rho-\dim\rho^\tau\bigr),
\]
其中 \(\tau\) 为换位。
由 \(S_4\) 角色表：
\(
\chi_{\mathbf{sgn}}(\tau)=-1,\chi_{\rho_2}(\tau)=0,\chi_{\rho_3}(\tau)=1,\chi_{\rho_{3'}}(\tau)=-1
\)，
故
\(
\dim(\mathbf{sgn})^\tau=0,\dim(\rho_2)^\tau=1,\dim(\rho_3)^\tau=2,\dim(\rho_{3'})^\tau=1
\)。
代入 \(r=12\) 得重数
\[
m_{\mathbf{sgn}}=5,\quad m_{\rho_2}=4,\quad m_{\rho_3}=3,\quad m_{\rho_{3'}}=9.
\]
维数核对：
\(
5\cdot1+4\cdot2+3\cdot3+9\cdot3=49=g(\mathcal X)
\)。
证毕。
\end{proof}

\begin{conclusion}[Jacobian 的 \texorpdfstring{$S_4$}{S4}-等变同源分解与 \texorpdfstring{$L$}{L}-函数因子化]\label{con:fold-gauge-anomaly-jacobian-equivariant-factorization}
设 \(J:=\mathrm{Jac}(\mathcal X)\)。
则存在定义于 \(\QQ\) 上的阿贝尔簇 \(A_{\mathbf{sgn}},A_2,A_3,A_{3'}\)，使
\[
J\sim_{\QQ}A_{\mathbf{sgn}}\times A_2\times A_3\times A_{3'}.
\]
其维数分别为
\[
\dim A_{\mathbf{sgn}}=5,\qquad \dim A_2=8,\qquad \dim A_3=9,\qquad \dim A_{3'}=27.
\]
并且在良约化素数处有 Euler 因子乘法分解，从而
\[
L(J,s)=L(A_{\mathbf{sgn}},s)\,L(A_2,s)\,L(A_3,s)\,L(A_{3'},s).
\]
\end{conclusion}
\begin{proof}[证明要点]
由结论 \ref{con:fold-gauge-anomaly-chevalley-weil-s4} 的表示分解，群代数中心幂等元在 \(J\) 上给出等变同源分解（Kani--Rosen 型机制）。
各分量维数由“表示维数 \(\times\) 重数”直接给出，即 \((5,8,9,27)\)。
\(\ell\)-进同调的半单分解与同源不变性给出 \(L\)-函数因子化。证毕。
\end{proof}

\begin{theorem}[\texorpdfstring{$S_4$}{S4} 正则闭包 Jacobian 的 N\'eron--Severi 秩下界]\label{thm:fold-gauge-anomaly-s4-jacobian-picard-rank-lower-bound-76}
令 \(J=\mathrm{Jac}(\mathcal X)\)，并在 \(\overline{\QQ}\) 上取其 N\'eron--Severi 群 \(\mathrm{NS}(J_{\overline{\QQ}})\)。
则
\[
\rho(J_{\overline{\QQ}}):=\mathrm{rank}\,\mathrm{NS}(J_{\overline{\QQ}})\ \ge\ 76.
\]
\leanverified{paper\_fold\_gauge\_anomaly\_s4\_jacobian\_picard\_rank\_lower\_bound\_76}
\end{theorem}
\begin{proof}
令
\(
V:=H^{0}(\mathcal X,\Omega^{1})
\)。
由结论 \ref{con:fold-gauge-anomaly-chevalley-weil-s4}，
\[
V\cong
5\cdot \mathbf{sgn}\ \oplus\ 4\cdot \rho_{2}\ \oplus\ 3\cdot \rho_{3}\ \oplus\ 9\cdot \rho_{3'}.
\]
由于 \(S_4\) 的不可约表示均可在 \(\QQ\) 上实现，且对每个不可约 \(\rho\) 有 \(\operatorname{End}_{S_4}(\rho)\cong\QQ\)，
故其交换子代数满足
\[
\operatorname{End}_{S_4}(V)\ \cong\ M_{5}(\QQ)\times M_{4}(\QQ)\times M_{3}(\QQ)\times M_{9}(\QQ).
\]
\par
另一方面，\(S_4\) 作用于 \(\mathcal X\) 诱导 \(\QQ[S_4]\) 在 \(J\) 上的对应作用，并在
\(
H^{0}(\mathcal X,\Omega^{1})
\)
上给出与上式兼容的表示；因此 \(\operatorname{End}_{S_4}(V)\) 可视为 \(\operatorname{End}^{0}(J)\) 中的一族 \(\QQ\)-子代数（在同源范畴中由等变分解的重数空间给出矩阵单元）。
\par
取 \(J\) 的主极化所诱导的 Rosati 反自同态 \(\dagger\)。
在上述分解下，\(\dagger\) 在每个矩阵因子 \(M_m(\QQ)\) 上诱导转置反演；因此 Rosati 不变子空间维数至少为
\[
\sum_{m\in\{5,4,3,9\}}\dim_{\QQ}\mathrm{Sym}_{m}(\QQ)
\;=\;
\sum_{m\in\{5,4,3,9\}}\frac{m(m+1)}{2}
\;=\;76.
\]
最后，对任意主极化阿贝尔簇 \(A\) 有标准同构
\[
\mathrm{NS}(A)_{\QQ}\ \cong\ \operatorname{End}^{0}(A)^{\dagger=1},
\]
从而 \(\rho(J_{\overline{\QQ}})=\dim_{\QQ}\mathrm{NS}(J_{\overline{\QQ}})_{\QQ}\ge 76\)。证毕。
\end{proof}

\begin{theorem}[换位分歧下的驯 Artin 导数与函数域导子度]\label{thm:fold-gauge-anomaly-s4-regular-closure-tame-artin-conductor-degrees}
沿用结论 \ref{con:fold-gauge-anomaly-genus-g3-g49} 的正则闭包覆盖
\(\mathcal X\to\PP^1_u\)。
其有限分歧点数为 \(12\)，且每个惰性群由一换位 \(\tau\in S_4\) 生成。
对 \(S_4\) 的四个非平凡不可约表示 \(\rho\in\{\mathbf{sgn},\rho_2,\rho_3,\rho_{3'}\}\)，定义驯 Artin 导子指数
\[
a(\rho;\tau):=\dim(\rho)-\dim\bigl(\rho^{\langle\tau\rangle}\bigr).
\]
则
\[
a(\mathbf{sgn};\tau)=1,\qquad
a(\rho_{2};\tau)=1,\qquad
a(\rho_{3};\tau)=1,\qquad
a(\rho_{3'};\tau)=2.
\]
从而相应的全局导子度（导子除子之次数）
\(
A(\rho):=\sum_{b}a(\rho;I_b)
\)
满足
\[
A(\mathbf{sgn})=12,\qquad
A(\rho_{2})=12,\qquad
A(\rho_{3})=12,\qquad
A(\rho_{3'})=24,
\]
其中求和遍及全部 \(12\) 个分歧点。
\leanverified{paper\_fold\_gauge\_anomaly\_s4\_artin\_conductor\_degrees}
\end{theorem}
\begin{proof}
由于覆盖在每个分歧点处为驯分歧且惰性群为阶 \(2\) 的循环群 \(\langle\tau\rangle\)，有
\(
a(\rho;\tau)=\dim\rho-\dim\rho^{\langle\tau\rangle}
\)。
而对阶 \(2\) 作用，\(\rho^{\langle\tau\rangle}\) 的维数由角色平均给出：
\[
\dim\rho^{\langle\tau\rangle}=\frac{\dim\rho+\operatorname{tr}(\rho(\tau))}{2}.
\]
对 \(S_4\) 的换位 \(\tau\)，角色值为
\[
\operatorname{tr}(\mathbf{sgn}(\tau))=-1,\qquad
\operatorname{tr}(\rho_2(\tau))=0,\qquad
\operatorname{tr}(\rho_3(\tau))=1,\qquad
\operatorname{tr}(\rho_{3'}(\tau))=-1.
\]
代入得到
\[
\dim(\mathbf{sgn})^{\langle\tau\rangle}=0,\quad
\dim(\rho_2)^{\langle\tau\rangle}=1,\quad
\dim(\rho_3)^{\langle\tau\rangle}=2,\quad
\dim(\rho_{3'})^{\langle\tau\rangle}=1,
\]
从而 \(a(\rho;\tau)\) 如陈述所示。
由于共有 \(12\) 个分歧点且局部数据一致，故 \(A(\rho)=12\cdot a(\rho;\tau)\)，得到全局导子度。证毕。
\end{proof}

\leanverified{paper\_fold\_gauge\_anomaly\_p10\_unique\_quadratic\_subfield}
\begin{corollary}[分裂域的唯一二次子域]\label{cor:fold-gauge-anomaly-p10-unique-quadratic-subfield}
设 \(L\) 为 \(P_{10}\) 的分裂域。
则 \(L\) 存在且仅存在一个二次子域，且
\[
L^{A_{10}}=\QQ\!\bigl(\sqrt{\mathrm{Disc}(P_{10})}\bigr)=\QQ\!\bigl(\sqrt{66269989}\bigr).
\]
\end{corollary}
\begin{proof}
由结论 \ref{con:fold-gauge-anomaly-p10-galois-s10}，\(\mathrm{Gal}(L/\QQ)\cong S_{10}\)。
\(S_{10}\) 的唯一指数 \(2\) 正规子群为 \(A_{10}\)，故二次子域唯一。
判别式平方自由部分为 \(1931\cdot 34319=66269989\)，即得所述表达。证毕。
\end{proof}

\begin{corollary}[分裂域次数与分支根不可根式性]\label{cor:fold-gauge-anomaly-p10-degree-and-unsolvability}
对 \(P_{10}\) 的分裂域 \(L\) 有
\[
[L:\QQ]=|S_{10}|=10!=3{,}628{,}800.
\]
从而 \(P_{10}(u)\) 的任一根均不可由根式表达。
\end{corollary}
\begin{proof}
由结论 \ref{con:fold-gauge-anomaly-p10-galois-s10}，分裂域次数等于伽罗瓦群阶，故为 \(10!\)。
又 \(S_{10}\) 非可解，依 Galois--Abel 判别，根式解不存在。证毕。
\leanverified{paper\_fold\_gauge\_anomaly\_p10\_degree\_and\_unsolvability}
\end{proof}

\begin{corollary}[两个正规商曲线的 Jacobian 等型实现]\label{cor:fold-gauge-anomaly-quotient-jacobians-isotypic}
设
\[
\mathcal C_{\mathrm{disc}}:=\mathcal X/A_4=\mathcal H,
\qquad
\mathcal Y_{V_4}:=\mathcal X/V_4.
\]
则有同源分解
\[
\mathrm{Jac}(\mathcal C_{\mathrm{disc}})\sim_{\QQ}A_{\mathbf{sgn}},
\qquad
\mathrm{Jac}(\mathcal Y_{V_4})\sim_{\QQ}A_{\mathbf{sgn}}\times A_2.
\]
\end{corollary}
\begin{proof}
对任意正规子群 \(N\triangleleft S_4\)，有
\[
H^0(\mathcal X/N,\Omega^1)\cong H^0(\mathcal X,\Omega^1)^N.
\]
由结论 \ref{con:fold-gauge-anomaly-chevalley-weil-s4}，
\[
H^0(\mathcal X,\Omega^1)\cong
5\cdot \mathbf{sgn}\oplus 4\cdot \rho_2\oplus 3\cdot \rho_3\oplus 9\cdot \rho_{3'}.
\]
对 \(A_4\) 取不变向量时，仅 \(\mathbf{sgn}\) 在 \(A_4\) 上平凡，故
\(
H^0(\mathcal X,\Omega^1)^{A_4}\cong 5\cdot\mathbf{sgn}
\)，
从而 \(\mathrm{Jac}(\mathcal X/A_4)\sim_{\QQ}A_{\mathbf{sgn}}\)。
对 \(V_4\) 取不变向量时，\(\mathbf{sgn}\) 与 \(\rho_2\) 平凡而 \(\rho_3,\rho_{3'}\) 非平凡，故
\(
H^0(\mathcal X,\Omega^1)^{V_4}\cong 5\cdot\mathbf{sgn}\oplus 4\cdot\rho_2
\)，
进而
\(
\mathrm{Jac}(\mathcal X/V_4)\sim_{\QQ}A_{\mathbf{sgn}}\times A_2
\)。
证毕。
\leanverified{paper\_fold\_gauge\_anomaly\_quotient\_jacobians\_isotypic}
\end{proof}

\begin{theorem}[\texorpdfstring{$A_2$}{A2} 因子的 Prym 实现]\label{thm:fold-gauge-anomaly-a2-prym-realization}
由 \(V_4\triangleleft A_4\triangleleft S_4\) 得到自然覆盖
\[
\pi:\mathcal Y_{V_4}=\mathcal X/V_4\longrightarrow \mathcal C_{\mathrm{disc}}=\mathcal X/A_4,
\qquad
\mathrm{Gal}(\pi)\cong A_4/V_4\cong C_3.
\]
其 Prym 簇
\[
\operatorname{Prym}(\pi):=\ker\!\bigl(\operatorname{Nm}_{\pi}:\mathrm{Jac}(\mathcal Y_{V_4})\to \mathrm{Jac}(\mathcal C_{\mathrm{disc}})\bigr)^0
\]
满足
\[
\operatorname{Prym}(\pi)\sim_{\QQ}A_2,
\qquad
\dim \operatorname{Prym}(\pi)=8.
\]
\end{theorem}
\begin{proof}
在结论 \ref{con:fold-gauge-anomaly-disc-factorization-simple-branch} 的 \(S_4\)-闭包中，有限分支点的惯性均由换位生成；
换位属于 \(S_4\setminus A_4\)，故其与 \(A_4\)（进而与 \(V_4\)）的交均为平凡。
因此 \(\pi\) 为循环三次非分歧覆盖。
对非分歧循环三次覆盖，Jacobian 有标准分解
\[
\mathrm{Jac}(\mathcal Y_{V_4})\sim_{\QQ}\mathrm{Jac}(\mathcal C_{\mathrm{disc}})\times \operatorname{Prym}(\pi).
\]
再由推论 \ref{cor:fold-gauge-anomaly-quotient-jacobians-isotypic}，
\(
\mathrm{Jac}(\mathcal Y_{V_4})\sim_{\QQ}A_{\mathbf{sgn}}\times A_2
\)
且
\(
\mathrm{Jac}(\mathcal C_{\mathrm{disc}})\sim_{\QQ}A_{\mathbf{sgn}}
\)，
比较即得 \(\operatorname{Prym}(\pi)\sim_{\QQ}A_2\)。
维数由
\(
g(\mathcal Y_{V_4})-g(\mathcal C_{\mathrm{disc}})=13-5=8
\)
给出。证毕。
\end{proof}

\leanverified{paper\_fold\_gauge\_anomaly\_a2\_prym\_realization}

\leanverified{paper\_fold\_gauge\_anomaly\_a2\_eisenstein\_multiplication}
\begin{theorem}[\texorpdfstring{$A_2$}{A2} 的 Eisenstein 整数乘法]\label{thm:fold-gauge-anomaly-a2-eisenstein-multiplication}
在定理 \ref{thm:fold-gauge-anomaly-a2-prym-realization} 的记号下，存在环嵌入
\[
\ZZ[\zeta_3]\hookrightarrow \operatorname{End}\!\bigl(\operatorname{Prym}(\pi)\bigr),
\qquad
\QQ(\zeta_3)\hookrightarrow \operatorname{End}^0\!\bigl(\operatorname{Prym}(\pi)\bigr).
\]
等价地，
\[
\QQ(\sqrt{-3})\subseteq \operatorname{End}^0(A_2).
\]
\end{theorem}
\begin{proof}
令 \(\sigma\) 为 \(\pi\) 的阶 \(3\) 甲板变换生成元。
\(\sigma\) 作用于 \(\mathrm{Jac}(\mathcal Y_{V_4})\)，并与范数映射 \(\operatorname{Nm}_{\pi}=1+\sigma+\sigma^2\) 交换，故保持
\(\operatorname{Prym}(\pi)=\ker(\operatorname{Nm}_{\pi})^0\)。
在 Prym 上，关系 \(\sigma^3=1\) 与 \(1+\sigma+\sigma^2=0\) 同时成立，故
\(\sigma\) 满足最小多项式 \(x^2+x+1\)，从而诱导
\(
\ZZ[\zeta_3]\cong \ZZ[\sigma]\subseteq \operatorname{End}(\operatorname{Prym}(\pi))
\)。
张量 \(\QQ\) 即得 \(\QQ(\zeta_3)\subseteq \operatorname{End}^0(\operatorname{Prym}(\pi))\)。
再由定理 \ref{thm:fold-gauge-anomaly-a2-prym-realization} 的同源识别
\(\operatorname{Prym}(\pi)\sim_{\QQ}A_2\) 得到结论。证毕。
\end{proof}

\begin{conjecture}[等型因子的绝对单纯与内自同态环极小化]\label{conj:fold-gauge-anomaly-isotypic-absolute-simplicity}
在结论 \ref{con:fold-gauge-anomaly-jacobian-equivariant-factorization} 的分解
\[
\mathrm{Jac}(\mathcal X)\sim_{\QQ}A_{\mathbf{sgn}}\times A_2\times A_3\times A_{3'}
\]
中，预期成立：
\begin{enumerate}
  \item \(A_2\) 在 \(\overline{\QQ}\) 上绝对单纯，且
  \[
  \operatorname{End}^0_{\overline{\QQ}}(A_2)=\QQ(\zeta_3).
  \]
  \item \(A_3,A_{3'}\) 在 \(\overline{\QQ}\) 上绝对单纯，且
  \[
  \operatorname{End}^0_{\overline{\QQ}}(A_3)=\operatorname{End}^0_{\overline{\QQ}}(A_{3'})=\QQ.
  \]
\end{enumerate}
该预期与通用简单分歧 \(S_4\)-覆盖的 Mumford--Tate 极大性一致；其可通过多组良素数 Frobenius 特征多项式与 Honda--Tate 判别进行有效核验。
\end{conjecture}

\endinput


\begin{corollary}[唯一二次子域与黄金二次域的线性互素性]\label{cor:fold-gauge-anomaly-P10-disjoint-sqrt5}
\leanverified{paper\_fold\_gauge\_anomaly\_p10\_disjoint\_sqrt5}
令 \(L/\QQ\) 为命题 \ref{prop:fold-gauge-anomaly-P10-galois-discriminant} 中 \(P_{10}\) 的分裂域。
则
\[
L\cap \QQ(\sqrt5)=\QQ,
\qquad
\mathrm{Gal}\bigl(L\cdot \QQ(\sqrt5)/\QQ\bigr)\cong S_{10}\times C_2.
\]
\end{corollary}
\begin{proof}[证明要点]
若 \(\QQ(\sqrt5)\subseteq L\)，则 \(\QQ(\sqrt5)\) 必为 \(L/\QQ\) 的指数 \(2\) 子域；
但命题 \ref{prop:fold-gauge-anomaly-P10-galois-discriminant} 已给出该二次子域唯一且等于 \(\QQ(\sqrt{66269989})\)，矛盾。
故 \(L\cap \QQ(\sqrt5)=\QQ\)。
由于 \(L/\QQ\) 与 \(\QQ(\sqrt5)/\QQ\) 均为伽罗瓦扩张，且交域平凡，故线性互素并给出直积型伽罗瓦群同构。证毕。
\end{proof}

\begin{corollary}[Chebotarev 等分布：模素分解型的显式密度]\label{cor:fold-gauge-anomaly-P10-modp-splitting-densities}
设 \(p\nmid \mathrm{Disc}(P_{10})\) 为素数。
则 \(P_{10}\) 在 \(\mathbb F_p[u]\) 中的不可约分解型按 \(S_{10}\) 的循环型共轭类等分布：
对任意 \(10\) 的分拆 \(\lambda\)，满足“模 \(p\) 分解型为 \(\lambda\)”的素数集合存在自然密度，且该密度等于
\[
\frac{|C_{\lambda}|}{10!},
\]
其中 \(C_{\lambda}\subset S_{10}\) 为循环型为 \(\lambda\) 的共轭类。
特别地，
\[
(10):\ \frac1{10},\qquad (9,1):\ \frac1{9},\qquad (7,3):\ \frac1{21},\qquad
\text{（含一次因子）}:\ 1-\frac{!10}{10!}=\frac{2293839}{3628800}=\frac{28319}{44800}.
\]
\end{corollary}
\begin{proof}[证明要点]
由命题 \ref{prop:fold-gauge-anomaly-P10-galois-discriminant}，\(P_{10}\) 的分裂域在 \(\QQ\) 上的伽罗瓦群为 \(S_{10}\)。
对 \(p\nmid \mathrm{Disc}(P_{10})\) 的非分歧素数，Frobenius 元在 \(S_{10}\) 的共轭类按 Chebotarev 密度定理等分布；
而其循环型与 \(P_{10}\) 在 \(\mathbb F_p[u]\) 的分解型一致，故一般陈述成立。
\par
对所列密度，分别计数相应循环型共轭类占比即可：\((10)\) 型元素占比为 \(1/10\)，\((9)(1)\) 型元素占比为 \(1/9\)，\((7)(3)\) 型元素个数为 \(10!/(7\cdot 3)\)，故占比为 \(1/21\)。
“含一次因子”等价于 Frobenius 在 \(S_{10}\) 中存在不动点；
其补集为错排，占比由容斥原理给出 \(\sum_{k=0}^{10}(-1)^k/k!\)，故所求密度为
\(
1-\sum_{k=0}^{10}(-1)^k/k!=1-\frac{!10}{10!}=\frac{2293839}{3628800}
\)。证毕。
\end{proof}

\begin{proposition}[分裂域与导子 \texorpdfstring{$37$}{37} 椭圆曲线挠域的线性互素性]\label{prop:fold-gauge-anomaly-P10-elliptic-torsion-disjointness}
设 \(L/\QQ\) 为命题 \ref{prop:fold-gauge-anomaly-P10-galois-discriminant} 中 \(P_{10}\) 的分裂域（\(\mathrm{Gal}(L/\QQ)\cong S_{10}\)）。
令
\[
E/\QQ:\ y^{2}+y=x^{3}-x
\]
并记其 \(\ell\)-挠域 \(M_\ell:=\QQ(E[\ell])\)。
则对任意素数 \(\ell\notin\{1931,34319\}\)，有
\[
L\cap M_\ell=\QQ,
\qquad\text{从而}\qquad
\mathrm{Gal}(LM_\ell/\QQ)\cong S_{10}\times \mathrm{Gal}(M_\ell/\QQ).
\]
\end{proposition}
\begin{proof}
由命题 \ref{prop:fold-gauge-anomaly-P10-galois-discriminant}，\(L\) 的唯一非平凡指数 \(2\) 子域为
\(
K=L^{A_{10}}=\QQ(\sqrt{66269989})
\)，且 \(K/\QQ\) 恰在素数 \(1931,34319\) 处分歧。
由于 \(L/\QQ\) 为伽罗瓦扩张，交域 \(L\cap M_\ell\) 必为 \(\QQ\) 的伽罗瓦扩张并包含于 \(L\)；
而在 \(S_{10}\) 下，非平凡真正规范子扩张仅能是 \(K/\QQ\)，故 \(L\cap M_\ell\in\{\QQ,K\}\)。
\par
另一方面，由 N\'eron--Ogg--Shafarevich 判据，\(M_\ell/\QQ\) 仅可能在 \(E\) 的坏约化素数与 \(\ell\) 处分歧；
对该曲线坏素数仅为 \(37\)，故当 \(\ell\notin\{1931,34319\}\) 时，\(M_\ell\) 在 \(1931,34319\) 处不分歧，从而不可能包含在这些素数处分歧的二次域 \(K\)。
于是 \(L\cap M_\ell=\QQ\)。
两者同为伽罗瓦扩张且交域平凡，推出线性互素性与直积型伽罗瓦群同构。证毕。
\end{proof}

\begin{proposition}[\texorpdfstring{$P_{10}$}{P10} 十次分裂域与 Lee--Yang 三次分裂域的线性互素性]\label{prop:fold-gauge-anomaly-P10-leyang-linear-disjointness}
\leanverified{paper\_fold\_gauge\_anomaly\_P10\_leyang\_linear\_disjointness}
令 \(L/\QQ\) 为命题 \ref{prop:fold-gauge-anomaly-P10-galois-discriminant} 中分歧多项式 \(P_{10}(u)\) 的分裂域，则 \(\mathrm{Gal}(L/\QQ)\cong S_{10}\) 且其唯一指数 \(2\) 子域为
\(
L^{A_{10}}=\QQ(\sqrt{66269989})
\)。
令
\[
P_{\mathrm{LY}}(y):=256y^{3}+411y^{2}+165y+32\in\ZZ[y],
\qquad
L_{\mathrm{LY}}:=\mathrm{Spl}\bigl(P_{\mathrm{LY}}\bigr),
\]
则 \(\mathrm{Gal}(L_{\mathrm{LY}}/\QQ)\cong S_3\)，且其唯一指数 \(2\) 子域为
\(
L_{\mathrm{LY}}^{A_3}=\QQ(\sqrt{-111})
\)
（结论 \ref{con:xi-terminal-zm-leyang-cubic-arithmetic}）。
因此
\[
L\cap L_{\mathrm{LY}}=\QQ,
\qquad\text{从而}\qquad
\mathrm{Gal}(LL_{\mathrm{LY}}/\QQ)\cong S_{10}\times S_{3}.
\]
\end{proposition}
\begin{proof}
由于 \(L/\QQ\) 与 \(L_{\mathrm{LY}}/\QQ\) 均为伽罗瓦扩张，交域 \(L\cap L_{\mathrm{LY}}\) 亦为 \(\QQ\) 上伽罗瓦扩张。
作为 \(L\) 的伽罗瓦子扩张，\(L\cap L_{\mathrm{LY}}\in\{\QQ,\QQ(\sqrt{66269989})\}\)；
作为 \(L_{\mathrm{LY}}\) 的伽罗瓦子扩张，\(L\cap L_{\mathrm{LY}}\in\{\QQ,\QQ(\sqrt{-111})\}\)。
但 \(\QQ(\sqrt{66269989})\) 为实二次域而 \(\QQ(\sqrt{-111})\) 为虚二次域，二者不相等，故只能 \(L\cap L_{\mathrm{LY}}=\QQ\)。
交域平凡且两边皆为伽罗瓦扩张，推出线性互素并得到直积型伽罗瓦群同构。证毕。
\end{proof}

\begin{corollary}[十次分歧域与 Lee--Yang 三次域的联合 Chebotarev 乘法律]\label{cor:fold-gauge-anomaly-P10-leyang-chebotarev-product-law}
\leanverified{paper\_fold\_gauge\_anomaly\_P10\_leyang\_chebotarev\_product\_law}
沿用命题 \ref{prop:fold-gauge-anomaly-P10-leyang-linear-disjointness} 的记号，并写
\[
L_{10}:=L.
\]
对任意满足
\[
p\nmid \mathrm{Disc}(P_{10})\mathrm{Disc}(P_{\mathrm{LY}})
\]
的素数 \(p\)，记
\[
\sigma_{10}(p):=\mathrm{Frob}_p(L_{10}/\QQ)\in S_{10},
\qquad
\sigma_{3}(p):=\mathrm{Frob}_p(L_{\mathrm{LY}}/\QQ)\in S_{3}.
\]
则对任意共轭类
\[
C_{10}\subset S_{10},
\qquad
C_{3}\subset S_{3},
\]
满足
\[
\sigma_{10}(p)\in C_{10},
\qquad
\sigma_{3}(p)\in C_{3}
\]
的素数集合具有自然密度
\[
\boxed{\;
\delta\Bigl(
\sigma_{10}(p)\in C_{10},\ \sigma_{3}(p)\in C_{3}
\Bigr)
=
\frac{|C_{10}|}{|S_{10}|}\cdot\frac{|C_{3}|}{|S_{3}|}.\;}
\]
其中由
\[
L_{10}^{A_{10}}=\QQ(\sqrt{66269989}),
\qquad
L_{\mathrm{LY}}^{A_3}=\QQ(\sqrt{-111})
\]
可见，显式暴露出的二次分歧支撑至少包含四个基本坏素
\[
3,\qquad 37,\qquad 1931,\qquad 34319.
\]
等价地，\(P_{10}\bmod p\) 在 \(\FF_p[u]\) 上的分解型与 \(P_{\mathrm{LY}}\bmod p\) 在 \(\FF_p[y]\) 上的分解型在自然密度意义下严格独立。
\end{corollary}
\begin{proof}[证明要点]
由命题 \ref{prop:fold-gauge-anomaly-P10-leyang-linear-disjointness}，
\[
\mathrm{Gal}(L_{10}L_{\mathrm{LY}}/\QQ)\cong S_{10}\times S_3.
\]
Chebotarev 密度定理说明：对共同不分歧素数，Frobenius 共轭类在该直积群中按共轭类大小等分布。而直积群中的共轭类正是笛卡尔积 \(C_{10}\times C_3\)，其比例为
\[
\frac{|C_{10}|\,|C_3|}{|S_{10}|\,|S_3|}
=
\frac{|C_{10}|}{|S_{10}|}\cdot\frac{|C_3|}{|S_3|}.
\]
对非分歧素数，Frobenius 循环型与模素分解型一致，故得到所示乘法律。证毕。
\end{proof}

\leanverified{paper\_fold\_gauge\_anomaly\_p10\_leyang\_irreducible\_root\_split\_densities}
\begin{corollary}[十次不可约与 Lee--Yang 单根/全分裂的显式正密度]\label{cor:fold-gauge-anomaly-P10-leyang-irreducible-root-split-densities}
\leanverified{paper\_fold\_gauge\_anomaly\_p10\_leyang\_irreducible\_root\_split\_densities}
在推论 \ref{cor:fold-gauge-anomaly-P10-leyang-chebotarev-product-law} 的设定下，记
\[
g(y):=P_{\mathrm{LY}}(y)=256y^{3}+411y^{2}+165y+32.
\]
对共同不分歧素数 \(p\)，定义事件
\[
\mathcal I:=\{\text{\(P_{10}\bmod p\) 在 \(\FF_p[u]\) 上不可约}\},
\]
\[
\mathcal R:=\{\text{\(g\bmod p\) 在 \(\FF_p[y]\) 上恰有一个 \(\FF_p\)-根}\},
\qquad
\mathcal S:=\{\text{\(g\bmod p\) 在 \(\FF_p[y]\) 上完全分裂}\}.
\]
则有显式正密度
\[
\boxed{\;
\delta(\mathcal I\cap \mathcal R)
=
\frac{1}{10}\cdot\frac12
=
\frac1{20}.\;}
\]
\[
\boxed{\;
\delta(\mathcal I\cap \mathcal S)
=
\frac{1}{10}\cdot\frac16
=
\frac1{60}.\;}
\]
\end{corollary}
\begin{proof}
由推论 \ref{cor:fold-gauge-anomaly-P10-modp-splitting-densities}，事件 \(\mathcal I\) 对应 \(S_{10}\) 中的 \(10\)-循环共轭类，故
\[
\delta(\mathcal I)=\frac1{10}.
\]
另一方面，由结论 \ref{con:xi-terminal-zm-leyang-cubic-modp-splitting-density}，\(g\bmod p\) 的分解型与 \(S_3\) 的循环型一一对应：换位类对应 \((1)(2)\)，故
\[
\delta(\mathcal R)=\frac{3}{6}=\frac12;
\]
恒等类对应 \((1)(1)(1)\)，故
\[
\delta(\mathcal S)=\frac{1}{6}.
\]
将这些边际密度代入推论 \ref{cor:fold-gauge-anomaly-P10-leyang-chebotarev-product-law} 的乘法律，即得两式。证毕。
\end{proof}

\begin{corollary}[最大阿贝尔子扩张：双二次判别角色]\label{cor:fold-gauge-anomaly-P10-leyang-max-abelian-subextension}
\leanverified{paper\_fold\_gauge\_anomaly\_p10\_leyang\_max\_abelian\_subextension}
令 \(F:=LL_{\mathrm{LY}}\)。
则 \(F/\QQ\) 的最大阿贝尔子扩张为
\[
F^{\mathrm{ab}}
=\QQ\!\left(\sqrt{66269989},\sqrt{-111}\right),
\qquad
\mathrm{Gal}(F^{\mathrm{ab}}/\QQ)\cong (\ZZ/2\ZZ)^2.
\]
并且对 \(p\nmid 3\cdot 37\cdot 1931\cdot 34319\) 的素数，\(p\) 在 \(F^{\mathrm{ab}}\) 中的分解行为由两条二次特征
\[
\left(\frac{66269989}{p}\right),\qquad \left(\frac{-111}{p}\right)
\]
完全刻画；四种符号组合在自然密度意义下各占 \(1/4\)。
\end{corollary}
\begin{proof}
由命题 \ref{prop:fold-gauge-anomaly-P10-leyang-linear-disjointness} 得 \(\mathrm{Gal}(F/\QQ)\cong S_{10}\times S_3\)。
而 \(S_{10}^{\mathrm{ab}}\cong C_2\)、\(S_3^{\mathrm{ab}}\cong C_2\)，故
\(
(S_{10}\times S_3)^{\mathrm{ab}}\cong C_2\times C_2
\)；
对应的最大阿贝尔子扩张由两因子的唯一二次子域生成，即为所示双二次域。
最后的 \(1/4\) 等分布由双二次扩张的 Chebotarev 定理直接推出。证毕。
\end{proof}

\begin{theorem}[Frobenius 类函数统计的严格张量独立]\label{thm:fold-gauge-anomaly-P10-leyang-class-function-tensor-independence}
沿用命题 \ref{prop:fold-gauge-anomaly-P10-leyang-linear-disjointness} 的记号。对任意类函数
\[
f:S_{10}\to\CC,
\qquad
g:S_3\to\CC,
\]
记群平均
\[
\langle f\rangle_{S_{10}}:=\frac{1}{10!}\sum_{\tau\in S_{10}}f(\tau),
\qquad
\langle g\rangle_{S_3}:=\frac{1}{6}\sum_{\upsilon\in S_3}g(\upsilon).
\]
再记
\[
\pi_{10,\mathrm{LY}}^\ast(x):=
\#\left\{p\le x:\ p\nmid \Disc(P_{10})\Disc(P_{\mathrm{LY}})\right\}.
\]
则有极限恒等式
\[
\boxed{\;
\lim_{x\to\infty}
\frac{1}{\pi_{10,\mathrm{LY}}^\ast(x)}
\sum_{\substack{p\le x\\ p\nmid \Disc(P_{10})\Disc(P_{\mathrm{LY}})}}
f\!\bigl(\sigma_{10}(p)\bigr)\,
g\!\bigl(\sigma_{3}(p)\bigr)
=
\langle f\rangle_{S_{10}}\,
\langle g\rangle_{S_3}.\;}
\]
特别地，若 \(\chi^\lambda\) 与 \(\psi^\mu\) 分别为 \(S_{10}\) 与 \(S_3\) 的不可约特征标，则
\[
\boxed{\;
\lim_{x\to\infty}
\frac{1}{\pi_{10,\mathrm{LY}}^\ast(x)}
\sum_{\substack{p\le x\\ p\nmid \Disc(P_{10})\Disc(P_{\mathrm{LY}})}}
\chi^\lambda\!\bigl(\sigma_{10}(p)\bigr)\,
\psi^\mu\!\bigl(\sigma_3(p)\bigr)
=
\begin{cases}
1,& \chi^\lambda=\mathbf 1,\ \psi^\mu=\mathbf 1,\\
0,& \text{否则}.
\end{cases}\;}
\]
\end{theorem}
\leanverified{paper\_fold\_gauge\_anomaly\_p10\_leyang\_class\_function\_tensor\_independence}
\begin{proof}
把 \(f\) 与 \(g\) 分别写成共轭类指示函数的线性组合：
\[
f=\sum_C a_C\,\mathbf 1_C,
\qquad
g=\sum_D b_D\,\mathbf 1_D,
\]
其中 \(C\subset S_{10}\)、\(D\subset S_3\) 遍历共轭类。于是
\[
f\!\bigl(\sigma_{10}(p)\bigr)\,g\!\bigl(\sigma_3(p)\bigr)
=
\sum_{C,D}a_Cb_D\,
\mathbf 1_{\{\sigma_{10}(p)\in C,\ \sigma_3(p)\in D\}}.
\]
对非分歧素数取平均，并应用推论 \ref{cor:fold-gauge-anomaly-P10-leyang-chebotarev-product-law}，得到
\[
\lim_{x\to\infty}
\frac{1}{\pi_{10,\mathrm{LY}}^\ast(x)}
\sum_{\substack{p\le x\\ p\nmid \Disc(P_{10})\Disc(P_{\mathrm{LY}})}}
f\!\bigl(\sigma_{10}(p)\bigr)\,g\!\bigl(\sigma_3(p)\bigr)
=
\sum_{C,D}a_Cb_D\frac{|C|}{10!}\frac{|D|}{6}.
\]
右端按 \(C\) 与 \(D\) 分组即化为
\[
\left(\sum_C a_C\frac{|C|}{10!}\right)
\left(\sum_D b_D\frac{|D|}{6}\right)
=
\langle f\rangle_{S_{10}}\langle g\rangle_{S_3}.
\]
这就证明了第一式。

对不可约特征标，使用有限群的标准平均恒等式
\[
\frac{1}{|G|}\sum_{g\in G}\chi(g)=
\begin{cases}
1,& \chi=\mathbf 1,\\
0,& \chi\neq \mathbf 1,
\end{cases}
\]
分别作用于 \(G=S_{10}\) 与 \(G=S_3\)，即可得到第二式。证毕。
\end{proof}

\begin{corollary}[十次判别式符号与 Lee--Yang 三次分解型的零互信息]\label{cor:fold-gauge-anomaly-P10-leyang-zero-mutual-information}
\leanverified{paper\_fold\_gauge\_anomaly\_p10\_leyang\_zero\_mutual\_information}
对不分歧素数 \(p\)，定义
\[
\chi_{10}(p):=\mathrm{sgn}\!\bigl(\sigma_{10}(p)\bigr)
=
\left(\frac{66269989}{p}\right)\in\{\pm 1\},
\]
并记 \(\mathcal T_{\mathrm{LY}}(p)\) 为 \(P_{\mathrm{LY}}\bmod p\) 的分解型随机变量，其取值于
\[
\{(1)(1)(1),\ (1)(2),\ (3)\}.
\]
将不分歧素数集合赋予自然密度概率测度，则 \(\chi_{10}\) 与 \(\mathcal T_{\mathrm{LY}}\) 严格独立。等价地，对任意
\[
\tau\in\{(1)(1)(1),(1)(2),(3)\},
\qquad
\varepsilon\in\{\pm 1\},
\]
都有
\[
\boxed{\;
\mathbb{P}\!\bigl(\mathcal T_{\mathrm{LY}}=\tau,\ \chi_{10}=\varepsilon\bigr)
=
\mathbb{P}\!\bigl(\mathcal T_{\mathrm{LY}}=\tau\bigr)\,
\mathbb{P}\!\bigl(\chi_{10}=\varepsilon\bigr).\;}
\]
从而
\[
\boxed{\;
I\!\bigl(\chi_{10};\mathcal T_{\mathrm{LY}}\bigr)=0.\;}
\]
更一般地，对任意函数
\[
h:\{(1)(1)(1),(1)(2),(3)\}\to\CC
\]
都有
\[
\boxed{\;
\mathbb{E}\!\left[h(\mathcal T_{\mathrm{LY}})\mid \chi_{10}=1\right]
=
\mathbb{E}\!\left[h(\mathcal T_{\mathrm{LY}})\mid \chi_{10}=-1\right]
=
\frac16 h((1)(1)(1))+\frac12 h((1)(2))+\frac13 h((3)).\;}
\]
\end{corollary}
\begin{proof}
取 \(f_\varepsilon:S_{10}\to\{0,1\}\) 为符号类函数
\[
f_\varepsilon(\tau):=\mathbf 1_{\{\mathrm{sgn}(\tau)=\varepsilon\}},
\]
并取 \(g_\tau:S_3\to\{0,1\}\) 为分解型 \(\tau\) 所对应的共轭类指示函数。由定理
\ref{thm:fold-gauge-anomaly-P10-leyang-class-function-tensor-independence}，
\[
\mathbb{P}\!\bigl(\mathcal T_{\mathrm{LY}}=\tau,\ \chi_{10}=\varepsilon\bigr)
=
\langle f_\varepsilon\rangle_{S_{10}}\,
\langle g_\tau\rangle_{S_3}
=
\mathbb{P}\!\bigl(\chi_{10}=\varepsilon\bigr)\,
\mathbb{P}\!\bigl(\mathcal T_{\mathrm{LY}}=\tau\bigr).
\]
故两随机变量独立。

独立的离散随机变量互信息为零，得第二式。最后，条件期望公式只是把独立性代入
\[
\mathbb{E}\!\left[h(\mathcal T_{\mathrm{LY}})\mid \chi_{10}=\varepsilon\right]
=
\sum_{\tau}h(\tau)\,
\mathbb{P}\!\bigl(\mathcal T_{\mathrm{LY}}=\tau\mid \chi_{10}=\varepsilon\bigr)
\]
并使用 \(S_3\) 三个共轭类大小为 \(1,3,2\)，故其边际分布分别为 \(1/6,1/2,1/3\)。证毕。
\end{proof}

\begin{proposition}[\texorpdfstring{$P_{10}$}{P10} 十次分裂域与指数平方因子 \texorpdfstring{$H$}{H} 的线性互素性]\label{prop:fold-gauge-anomaly-P10-H-linear-disjointness}
令 \(L/\QQ\) 为命题 \ref{prop:fold-gauge-anomaly-P10-galois-discriminant} 中 \(P_{10}(u)\) 的分裂域。
令 \(H(u)\) 为命题 \ref{prop:fold-gauge-anomaly-H-galois-discriminant} 中次数 \(19\) 的不可约多项式，并记其分裂域为 \(L_H:=\mathrm{Spl}(H)\)。
则
\[
L\cap L_H=\QQ,
\qquad\text{从而}\qquad
\mathrm{Gal}(LL_H/\QQ)\cong S_{10}\times S_{19}.
\]
\end{proposition}
\begin{proof}
记 \(E:=L\cap L_H\)。
由于 \(L/\QQ\) 与 \(L_H/\QQ\) 均为伽罗瓦扩张，\(E/\QQ\) 亦为伽罗瓦扩张。
\par
由命题 \ref{prop:fold-gauge-anomaly-P10-galois-discriminant}，\(L\) 的唯一指数 \(2\) 子域为
\(
K_{10}:=L^{A_{10}}=\QQ(\sqrt{66269989})
\)，故 \(E\in\{\QQ,K_{10}\}\)。
由命题 \ref{prop:fold-gauge-anomaly-H-galois-discriminant}，\(\mathrm{Gal}(L_H/\QQ)\cong S_{19}\) 且其唯一指数 \(2\) 子域为
\[
K_H:=L_H^{A_{19}}=\QQ\!\left(\sqrt{\mathrm{Disc}(H)}\right),
\]
故 \(E\in\{\QQ,K_H\}\)。
\par
由命题 \ref{prop:fold-gauge-anomaly-H-galois-discriminant} 的判别式分解可见素数 \(101\) 在 \(\mathrm{Disc}(H)\) 中以奇指数出现，因而 \(101\) 在二次域 \(K_H\) 中分歧。
而 \(101\nmid 66269989\)，故 \(101\) 在 \(K_{10}\) 中不分歧，从而 \(K_{10}\neq K_H\)。
于是 \(E\) 不可能为任何非平凡二次子域，必有 \(E=\QQ\)。
交域平凡且两边皆为伽罗瓦扩张，推出线性互素并得到直积型伽罗瓦群同构。证毕。
\leanverified{paper\_fold\_gauge\_anomaly\_P10\_H\_linear\_disjointness}
\leanverified{paper\_fold\_gauge\_anomaly\_p10\_h\_linear\_disjointness}
\end{proof}

\begin{corollary}[合成分裂域的直积分解]\label{cor:fold-gauge-anomaly-P10H-splitting-field-product}
记 \(D(u):=P_{10}(u)\,H(u)\in\ZZ[u]\)。
则 \(D\) 的分裂域满足
\[
\mathrm{Spl}(D)=LL_H,
\qquad
\mathrm{Gal}\bigl(\mathrm{Spl}(D)/\QQ\bigr)\cong S_{10}\times S_{19}.
\]
\end{corollary}
\begin{proof}
\(\mathrm{Spl}(D)\) 由 \(P_{10}\) 与 \(H\) 的根集合共同生成，故 \(\mathrm{Spl}(D)=LL_H\)。
伽罗瓦群同构由命题 \ref{prop:fold-gauge-anomaly-P10-H-linear-disjointness} 直接推出。证毕。
\end{proof}
\leanverified{paper\_fold\_gauge\_anomaly\_P10H\_splitting\_field\_product}
\leanverified{paper\_fold\_gauge\_anomaly\_p10h\_splitting\_field\_product}

\begin{corollary}[Chebotarev 乘积律：\texorpdfstring{$P_{10}$}{P10} 与 \texorpdfstring{$H$}{H} 的分解型完全解耦]\label{cor:fold-gauge-anomaly-P10-H-chebotarev-product-law}
\leanverified{paper\_fold\_gauge\_anomaly\_P10\_H\_chebotarev\_product\_law}
在命题 \ref{prop:fold-gauge-anomaly-P10-H-linear-disjointness} 的记号下，取素数 \(p\nmid \mathrm{Disc}(P_{10})\mathrm{Disc}(H)\)。
令 \(\lambda\vdash 10\)、\(\mu\vdash 19\) 为分拆，并记 \(C_\lambda\subset S_{10}\)、\(C_\mu\subset S_{19}\) 为相应循环型的共轭类。
则满足 \(P_{10}\bmod p\) 的分解型为 \(\lambda\) 且 \(H\bmod p\) 的分解型为 \(\mu\) 的素数集合具有自然密度
\[
\delta(\lambda,\mu)=\frac{|C_\lambda|}{10!}\cdot \frac{|C_\mu|}{19!}.
\]
\end{corollary}
\begin{proof}
由命题 \ref{prop:fold-gauge-anomaly-P10-H-linear-disjointness}，\(\mathrm{Gal}(LL_H/\QQ)\cong S_{10}\times S_{19}\)。
对非分歧素数 \(p\)，Frobenius 共轭类落在 \(C_\lambda\times C_\mu\) 中当且仅当两边模素分解型分别为 \(\lambda\) 与 \(\mu\)。
Chebotarev 密度定理给出密度为
\(
|C_\lambda\times C_\mu|/|S_{10}\times S_{19}|
=\frac{|C_\lambda|}{10!}\cdot\frac{|C_\mu|}{19!}
\)。
证毕。
\end{proof}

\begin{corollary}[若干可核验的显式密度常数]\label{cor:fold-gauge-anomaly-P10-H-explicit-densities}
\leanverified{paper\_fold\_gauge\_anomaly\_p10\_h\_explicit\_densities}
在推论 \ref{cor:fold-gauge-anomaly-P10-H-chebotarev-product-law} 的条件下，以下密度均指除去有限个分歧素数后的自然密度：
\begin{enumerate}
  \item 同时不可约：
  \[
  \delta\bigl((10),(19)\bigr)=\frac1{10}\cdot\frac1{19}=\frac1{190}.
  \]
  \item \(P_{10}\) 含一次因子且 \(H\) 不可约：
  \[
  \delta\bigl(\text{$P_{10}$ 含一次因子},(19)\bigr)
  =
  \Bigl(1-\frac{!10}{10!}\Bigr)\cdot\frac1{19}
  =
  \frac{28319}{44800}\cdot\frac1{19}.
  \]
  \item \(P_{10}\) 与 \(H\) 同时含一次因子：
  \[
  \delta\bigl(\text{$P_{10}$ 含一次因子},\text{$H$ 含一次因子}\bigr)
  =
  \Bigl(1-\frac{!10}{10!}\Bigr)\cdot \Bigl(1-\frac{!19}{19!}\Bigr).
  \]
\end{enumerate}
\end{corollary}
\begin{proof}
由推论 \ref{cor:fold-gauge-anomaly-P10-H-chebotarev-product-law}，密度分解为边际密度之乘积。
\(\delta((10))=1/10\)、\(\delta((19))=1/19\)。
并且 \(\delta(\text{$P_{10}$ 含一次因子})=1-!10/10!=28319/44800\) 由推论 \ref{cor:fold-gauge-anomaly-P10-modp-splitting-densities}。
而 \(\delta(\text{$H$ 含一次因子})=1-!19/19!\) 为 \(S_{19}\) 中存在不动点的比例。
证毕。
\end{proof}

\begin{corollary}[最大阿贝尔子扩张与双二次特征的等分布]\label{cor:fold-gauge-anomaly-P10-H-max-abelian-and-quadratic-equidistribution}
\leanverified{paper\_fold\_gauge\_anomaly\_p10\_h\_max\_abelian\_and\_quadratic\_equidistribution}
设 \(F:=LL_H\)。
则 \(F/\QQ\) 的最大阿贝尔子扩张为
\[
F^{\mathrm{ab}}=K_{10}K_H,
\qquad
\mathrm{Gal}(F^{\mathrm{ab}}/\QQ)\cong (\ZZ/2\ZZ)^2,
\]
其中 \(K_{10}=\QQ(\sqrt{66269989})\)、\(K_H=\QQ(\sqrt{\mathrm{Disc}(H)})\)。
并且对 \(p\) 不在 \(K_{10}K_H\) 中分歧的素数，二次特征对
\[
\chi_{10}(p):=\left(\frac{66269989}{p}\right),
\qquad
\chi_H(p):=\left(\frac{\mathrm{Disc}(H)}{p}\right)
\]
的四种符号组合在自然密度意义下各占 \(1/4\)。
\end{corollary}
\begin{proof}
由命题 \ref{prop:fold-gauge-anomaly-P10-H-linear-disjointness}，
\(\mathrm{Gal}(F/\QQ)\cong S_{10}\times S_{19}\)。
而 \(S_{10}^{\mathrm{ab}}\cong C_2\)、\(S_{19}^{\mathrm{ab}}\cong C_2\)，故
\(
(S_{10}\times S_{19})^{\mathrm{ab}}\cong C_2\times C_2
\)；
对应的最大阿贝尔子扩张由两因子的唯一二次子域生成，即 \(K_{10}K_H\)。
对双二次扩张应用 Chebotarev 定理得到 Frobenius 在 \((\ZZ/2\ZZ)^2\) 上等分布，从而四种符号组合密度均为 \(1/4\)。证毕。
\end{proof}

\begin{proposition}[固定一次因子数的 Rencontres 密度律]\label{prop:fold-gauge-anomaly-H-rencontres-fixed-point-density}
\leanverified{paper\_fold\_gauge\_anomaly\_H\_rencontres\_fixed\_point\_density}
\leanverified{paper\_fold\_gauge\_anomaly\_h\_rencontres\_fixed\_point\_density}
对 \(r\in\{0,1,\dots,19\}\)，记事件
\[
\mathcal E_r(p):=\text{$H\bmod p$ 恰有 \(r\) 个一次因子}.
\]
则除去有限个分歧素数后，\(\mathcal E_r\) 的自然密度为
\[
\delta(\mathcal E_r)
=\frac{\binom{19}{r}\,!(19-r)}{19!}
=\frac{1}{r!}\sum_{k=0}^{19-r}\frac{(-1)^k}{k!},
\]
其中 \(!(m)\) 为 \(m\) 阶错排数。
进一步，若记事件 \(\mathcal I(p):=\text{$P_{10}\bmod p$ 不可约}\)，则有联合密度
\[
\delta(\mathcal I\cap \mathcal E_r)=\frac1{10}\cdot \delta(\mathcal E_r).
\]
\end{proposition}
\begin{proof}
由命题 \ref{prop:fold-gauge-anomaly-H-galois-discriminant}，\(\mathrm{Gal}(H/\QQ)\cong S_{19}\)。
对非分歧素数 \(p\)，\(H\bmod p\) 的一次因子个数等于 Frobenius 置换在 \(19\) 个根上不动点数。
恰有 \(r\) 个不动点的置换个数为 \(\binom{19}{r}!(19-r)\)，故密度为 \(\binom{19}{r}!(19-r)/19!\)。
等价表达式由错排数的容斥公式化简得到。
\par
联合密度由推论 \ref{cor:fold-gauge-anomaly-P10-H-chebotarev-product-law} 的直积独立性给出：\(\delta(\mathcal I)=1/10\) 且与 \(\mathcal E_r\) 独立，故
\(
\delta(\mathcal I\cap\mathcal E_r)=\delta(\mathcal I)\delta(\mathcal E_r)=\frac1{10}\delta(\mathcal E_r)
\)。
证毕。
\end{proof}

\begin{theorem}[Lee--Yang 三次域、\texorpdfstring{$P_{10}$}{P10} 十次域与指数域的三重直积解耦]\label{thm:fold-gauge-anomaly-P10-leyang-H-triple-product}
\leanverified{paper\_fold\_gauge\_anomaly\_p10\_leyang\_h\_triple\_product}
在命题 \ref{prop:fold-gauge-anomaly-P10-leyang-linear-disjointness} 与命题 \ref{prop:fold-gauge-anomaly-P10-H-linear-disjointness} 的记号下，有
\[
\mathrm{Gal}(LL_{\mathrm{LY}}L_H/\QQ)\cong S_{10}\times S_{3}\times S_{19}.
\]
\end{theorem}
\begin{proof}
记 \(F:=LL_{\mathrm{LY}}\) 并由命题 \ref{prop:fold-gauge-anomaly-P10-leyang-linear-disjointness} 得 \(\mathrm{Gal}(F/\QQ)\cong S_{10}\times S_3\)。
由推论 \ref{cor:fold-gauge-anomaly-P10-leyang-max-abelian-subextension}，\(F\) 的最大阿贝尔子扩张为
\(
F^{\mathrm{ab}}=\QQ(\sqrt{66269989},\sqrt{-111})
\)。
\par
设 \(E:=F\cap L_H\)。
由于 \(\mathrm{Gal}(L_H/\QQ)\cong S_{19}\)，其伽罗瓦子扩张只有 \(\QQ\) 与唯一二次子域 \(K_H=L_H^{A_{19}}\)，故 \(E\in\{\QQ,K_H\}\)。
由命题 \ref{prop:fold-gauge-anomaly-H-galois-discriminant}，\(101\) 在 \(K_H\) 中分歧。
但 \(101\) 在 \(F^{\mathrm{ab}}\) 中不分歧（其分歧素数只可能来自 \(66269989\) 与 \(-111\) 的素因子），故 \(K_H\not\subseteq F\)，从而 \(E=\QQ\)。
于是 \(F\) 与 \(L_H\) 线性互素并得到
\(
\mathrm{Gal}(FL_H/\QQ)\cong \mathrm{Gal}(F/\QQ)\times \mathrm{Gal}(L_H/\QQ)
\cong S_{10}\times S_3\times S_{19}
\)。
证毕。
\end{proof}
\leanverified{paper\_fold\_gauge\_anomaly\_P10\_leyang\_H\_triple\_product}

\begin{corollary}[三重分解型的完全乘法密度]\label{cor:fold-gauge-anomaly-P10-leyang-H-chebotarev-triple-product}
在定理 \ref{thm:fold-gauge-anomaly-P10-leyang-H-triple-product} 的设定下，
对任意共轭类 \(C_\lambda\subset S_{10}\)、\(C_\nu\subset S_3\)、\(C_\mu\subset S_{19}\)，满足 \(\mathrm{Frob}_p\in C_\lambda\times C_\nu\times C_\mu\) 的素数集合（除去有限坏素集合后）具有自然密度
\[
\delta(\lambda,\nu,\mu)=\frac{|C_\lambda|}{10!}\cdot\frac{|C_\nu|}{3!}\cdot\frac{|C_\mu|}{19!}.
\]
\end{corollary}
\begin{proof}
由定理 \ref{thm:fold-gauge-anomaly-P10-leyang-H-triple-product}，伽罗瓦群为直积。
对直积群应用 Chebotarev 密度定理即可得到共轭类分布的乘积分解。证毕。
\end{proof}
\leanverified{paper\_fold\_gauge\_anomaly\_P10\_leyang\_H\_chebotarev\_triple\_product}
\leanverified{paper\_fold\_gauge\_anomaly\_p10\_leyang\_h\_chebotarev\_triple\_product}

\begin{proposition}[Lee--Yang 三次分裂域与 \texorpdfstring{$\ell$}{l}-挠域的线性互素性]\label{prop:fold-gauge-anomaly-leyang-elliptic-torsion-disjointness}
\leanverified{paper\_fold\_gauge\_anomaly\_leyang\_elliptic\_torsion\_disjointness}
沿用命题 \ref{prop:fold-gauge-anomaly-P10-elliptic-torsion-disjointness} 的椭圆曲线 \(E/\QQ\) 与挠域 \(M_\ell=\QQ(E[\ell])\) 记号，并记 \(L_{\mathrm{LY}}=\mathrm{Spl}(P_{\mathrm{LY}})\)。
则对任意素数 \(\ell\neq 3\)，有
\[
L_{\mathrm{LY}}\cap M_\ell=\QQ,
\qquad\text{从而}\qquad
\mathrm{Gal}(L_{\mathrm{LY}}M_\ell/\QQ)\cong S_{3}\times \mathrm{Gal}(M_\ell/\QQ).
\]
\end{proposition}
\begin{proof}
交域 \(L_{\mathrm{LY}}\cap M_\ell\) 为 \(\QQ\) 上伽罗瓦扩张，且作为 \(L_{\mathrm{LY}}\) 的伽罗瓦子扩张只能为 \(\QQ\) 或其唯一二次子域 \(\QQ(\sqrt{-111})\)（结论 \ref{con:xi-terminal-zm-leyang-cubic-arithmetic}）。
但 \(\QQ(\sqrt{-111})\) 在素数 \(3\) 处分歧，而当 \(\ell\neq 3\) 时，N\'eron--Ogg--Shafarevich 判据给出 \(M_\ell/\QQ\) 仅可能在 \(37\) 与 \(\ell\) 处分歧（命题 \ref{prop:fold-gauge-anomaly-P10-elliptic-torsion-disjointness}），故其任意子扩张在 \(3\) 处亦不分歧。
矛盾，遂交域只能为 \(\QQ\)；两边皆为伽罗瓦扩张且交域平凡，得到线性互素与直积型伽罗瓦群同构。证毕。
\end{proof}

\begin{corollary}[三通道直积独立性：\texorpdfstring{$P_{10}$}{P10}、Lee--Yang 与 \texorpdfstring{$\ell$}{l}-挠域]\label{cor:fold-gauge-anomaly-P10-leyang-elliptic-triple-independence}
\leanverified{paper\_fold\_gauge\_anomaly\_p10\_leyang\_elliptic\_triple\_independence}
在命题 \ref{prop:fold-gauge-anomaly-P10-galois-discriminant}、命题 \ref{prop:fold-gauge-anomaly-P10-leyang-linear-disjointness}、命题 \ref{prop:fold-gauge-anomaly-P10-elliptic-torsion-disjointness} 与命题 \ref{prop:fold-gauge-anomaly-leyang-elliptic-torsion-disjointness} 的记号下，
若 \(\ell\notin\{3,1931,34319\}\)，则
\[
\mathrm{Gal}(LL_{\mathrm{LY}}M_\ell/\QQ)
\ \cong\
S_{10}\times S_{3}\times \mathrm{Gal}(M_\ell/\QQ).
\]
因此对除去有限坏素集合后的素数 \(p\)，下列三类数据在自然密度意义下相互独立：
\begin{enumerate}
  \item \(P_{10}\bmod p\) 的分解型（对应 \(S_{10}\) 的 Frobenius 循环型）；
  \item \(P_{\mathrm{LY}}\bmod p\) 的分解型（对应 \(S_{3}\) 的 Frobenius 循环型）；
  \item \(\rho_{E,\ell}(\mathrm{Frob}_p)\) 在 \(\mathrm{Gal}(M_\ell/\QQ)\) 中的共轭类。
\end{enumerate}
\end{corollary}
\begin{proof}[证明要点]
令 \(F:=LL_{\mathrm{LY}}\) 并由命题 \ref{prop:fold-gauge-anomaly-P10-leyang-linear-disjointness} 得 \(\mathrm{Gal}(F/\QQ)\cong S_{10}\times S_3\)。
设 \(K_1=\QQ(\sqrt{66269989})\)、\(K_2=\QQ(\sqrt{-111})\)。
由推论 \ref{cor:fold-gauge-anomaly-P10-leyang-max-abelian-subextension}，\(F\) 的全部二次子域恰为 \(K_1\)、\(K_2\) 与 \(\QQ(\sqrt{66269989\cdot(-111)})\) 三者。
\par
当 \(\ell\notin\{3,1931,34319\}\) 时，N\'eron--Ogg--Shafarevich 判据表明 \(M_\ell/\QQ\) 在 \(3,1931,34319\) 处不分歧（命题 \ref{prop:fold-gauge-anomaly-P10-elliptic-torsion-disjointness}），从而 \(M_\ell\) 不可能包含上述任一二次子域。
另一方面，由于 \(S_{10}\) 与 \(S_3\) 的全部非平凡正规商都经由符号特征产生的 \(C_2\)（命题 \ref{prop:fold-gauge-anomaly-P10-galois-discriminant} 与结论 \ref{con:xi-terminal-zm-leyang-cubic-arithmetic}），任意非平凡伽罗瓦子扩张 \(E/\QQ\) 满足 \(E\subseteq F\) 时必含某一二次子域。
据此 \(F\cap M_\ell\) 只能为 \(\QQ\)，从而 \(F\) 与 \(M_\ell\) 线性互素，给出所示直积同构。
独立性为 Chebotarev 密度定理在直积群上的乘积分布之直接推论。证毕。
\end{proof}

\paragraph{有限域审计接口}
脚本 \texttt{scripts/exp\_fold\_gauge\_anomaly\_s4\_hurwitz\_tower\_audit.py} 在同一 JSON 证书中并列给出 \(P_{10}\) 与 \(P_{\mathrm{LY}}\) 的判别式平方自由核、二次子域分离，以及对一段有限素数范围的联合分解型计数（输出：\texttt{artifacts/export/fold\_gauge\_anomaly\_s4\_hurwitz\_tower\_audit.json}）。

\endinput


\begin{proposition}[规范差谱四次的分歧—本征值对偶消元与同域证书]\label{prop:fold-gauge-anomaly-Q10-tschirnhaus}
考虑系统
\[
F(\mu,u)=0,\qquad \partial_{\mu}F(\mu,u)=0,
\]
其中 \(F(\mu,u)=\det(\mu I-2A_{\theta})\in\ZZ[\mu,u]\) 且 \(u=e^{\theta}\)（命题 \ref{prop:fold-gauge-anomaly-discriminant-factorization}）。
令消元理想
\[
I:=(F,\partial_{\mu}F)\subset \QQ[\mu,u].
\]
则对字典序 \(u\succ \mu\) 的消元，有结果式分解
\[
\Res_{u}\bigl(F(\mu,u),\partial_{\mu}F(\mu,u)\bigr)=-\mu(\mu+1)\,Q_{10}(\mu),
\]
其中
\[
\begin{aligned}
Q_{10}(\mu)=\;&27\mu^{10}-81\mu^{9}+90\mu^{8}-46\mu^{7}+11\mu^{6}-\mu^{5}\\
&-32\mu^{4}+32\mu^{3}+16\mu+16.
\end{aligned}
\]
因此，除去 \(\mu\in\{0,-1\}\) 的平凡因子外，\(Q_{10}\) 的根精确刻画所有非平凡分支点处的重根谱值（即同时满足 \(F=\partial_{\mu}F=0\) 的 \(\mu\)-坐标）。
并且在同一消元计算中得到 \(u\) 关于 \(\mu\) 的线性回代约束
\[
\begin{aligned}
u=\frac{1}{200}\bigl(&-189\mu^{9}+540\mu^{8}-360\mu^{7}-308\mu^{6}+329\mu^{5}\\
&+304\mu^{4}-104\mu^{3}-196\mu^{2}-128\mu-16\bigr),
\end{aligned}
\]
从而对任意满足 \(Q_{10}(\mu)=0\) 的根，上式给出的 \(u\) 自动满足 \(P_{10}(u)=0\)（命题 \ref{prop:fold-gauge-anomaly-discriminant-factorization}），并在代数闭包中给出非平凡分歧集合的双向参数化。
此外，
\[
\mathrm{Gal}\bigl(Q_{10}/\QQ\bigr)\cong S_{10},\qquad
\mathrm{Disc}\bigl(Q_{10}\bigr)=2^{42}3^{16}5^{8}\cdot 1931\cdot 34319.
\]
\leanverified{paper\_fold\_gauge\_anomaly\_q10\_tschirnhaus}
\end{proposition}
\begin{proof}[证明要点（消元）]
消元定理给出 \(I\cap\QQ[\mu]\) 由 \(\Res_{u}(F,\partial_{\mu}F)\) 生成。
对 \((F,\partial_{\mu}F)\) 作字典序 \(u\succ\mu\) 的 Gr\"obner 消元（或直接计算结果式）得到所示分解
\(
\Res_{u}(F,\partial_{\mu}F)=-\mu(\mu+1)Q_{10}(\mu)
\)。
同一组消元基在 \(I\) 中产生一个对 \(u\) 的一次多项式，从而给出所示回代公式。
将该回代式代入 \(P_{10}(u)\) 并在商环 \(\QQ[\mu]/(Q_{10}(\mu))\) 中化简为零，即得 \(Q_{10}\)-根到 \(P_{10}\)-根的代数对应。
\par
关于 \(\mathrm{Gal}(Q_{10}/\QQ)\cong S_{10}\) 与判别式分解，取不整除 \(\mathrm{Disc}(Q_{10})\) 的素数并比较模素分解型得到一组 Frobenius 循环型证书，从而推出 \(A_{10}\subseteq \mathrm{Gal}(Q_{10}/\QQ)\)；
再由判别式非平方排除 \(\mathrm{Gal}(Q_{10}/\QQ)\subseteq A_{10}\)，得到满对称群。判别式因式分解由符号计算得到。
\end{proof}

\begin{proposition}[\texorpdfstring{$P_{10}$}{P10} 与 \texorpdfstring{$Q_{10}$}{Q10} 的显式 Tschirnhaus 互逆与分裂域一致性]\label{prop:fold-gauge-anomaly-P10-Q10-tschirnhaus-inverse}
设
\[
K_u:=\QQ[u]/(P_{10}(u)),\qquad K_{\mu}:=\QQ[\mu]/(Q_{10}(\mu)).
\]
定义 \(\QQ\)-代数同态
\[
\Psi:K_u\to K_{\mu},\qquad
\Psi(u)=\frac{1}{200}\bigl(-189\mu^{9}+540\mu^{8}-360\mu^{7}-308\mu^{6}+329\mu^{5}+304\mu^{4}-104\mu^{3}-196\mu^{2}-128\mu-16\bigr),
\]
以及
\[
\Phi:K_{\mu}\to K_u,\qquad
\Phi(\mu)=\frac{1}{14847223056}\sum_{j=0}^{9}c_j u^{j},
\]
其中系数为
\[
\begin{aligned}
(c_0,c_1,c_2,c_3,c_4)=\;&(-4448138872,\,-39661981520,\,83566151680,\,-13806152072,\,-124713321760),\\
(c_5,c_6,c_7,c_8,c_9)=\;&(-47968965903,\,21106705790,\,9579962376,\,-3730846095,\,-2044918440).
\end{aligned}
\]
则 \(\Phi,\Psi\) 互为逆同构。
特别地，对任意非平凡分歧对 \((u_0,\mu_0)\)（即 \(P_{10}(u_0)=0\) 且 \(Q_{10}(\mu_0)=0\) 且满足命题 \ref{prop:fold-gauge-anomaly-Q10-tschirnhaus} 的回代关系），有
\[
\QQ(u_0)=\QQ(\mu_0),
\qquad
\mathrm{Spl}(P_{10})=\mathrm{Spl}(Q_{10}).
\]
\leanverified{paper\_fold\_gauge\_anomaly\_P10\_Q10\_tschirnhaus\_inverse}
\end{proposition}
\begin{proof}
由命题 \ref{prop:fold-gauge-anomaly-Q10-tschirnhaus} 的回代关系知，\(\Psi(u)\) 在 \(K_{\mu}\) 中满足 \(P_{10}(\Psi(u))=0\)，故 \(\Psi\) 良定义。
另一方面，\(\Phi(\mu)\) 的选择满足 \(Q_{10}(\Phi(\mu))=0\) 在 \(K_u\) 中成立，从而 \(\Phi\) 良定义。
对复合 \(\Psi\circ\Phi\) 与 \(\Phi\circ\Psi\) 作多项式代入并分别在商环 \(K_{\mu}\)、\(K_u\) 中约化，可得二者皆为恒等映射，故互为逆同构。
由根域同构立得 \(\QQ(u_0)=\QQ(\mu_0)\)。
并且由于两侧在代数闭包中可由显式多项式互相回代得到全部共轭根，故生成的分裂域相同。
\end{proof}

\begin{corollary}[判别式核不变量与唯一二次子域]\label{cor:fold-gauge-anomaly-Q10-unique-quadratic-subfield}
\leanverified{paper\_fold\_gauge\_anomaly\_q10\_unique\_quadratic\_subfield}
设 \(L\) 为 \(Q_{10}\) 的分裂域。
则 \(L/\QQ\) 的唯一指数 \(2\) 子域为
\[
L^{A_{10}}=\QQ\!\left(\sqrt{66269989}\right),
\qquad
66269989=1931\cdot 34319=\mathrm{sf}\!\bigl(\mathrm{Disc}(Q_{10})\bigr).
\]
因此 \(Q_{10}\)（亦即 \(P_{10}\)）在根式意义下不可解。
\end{corollary}
\begin{proof}
由命题 \ref{prop:fold-gauge-anomaly-Q10-tschirnhaus}，\(\mathrm{Gal}(L/\QQ)\cong S_{10}\) 且 \(\mathrm{Disc}(Q_{10})\) 的平方自由部分为 \(66269989\)。
符号表示对应的二次子域由 \(\sqrt{\mathrm{Disc}(Q_{10})}\) 生成，且 \(S_{10}\) 的唯一指数 \(2\) 子群为 \(A_{10}\)，故所示二次子域唯一。
由于 \(S_{10}\) 不可解，分裂域不可由根式塔得到，从而多项式不可由根式求解。
\end{proof}


\paragraph{分歧点的有理降阶、斜率平方不变量与 Lee--Yang 边界的普适定标核}
本段以一参数变形谱四次
\[
F_t(\mu,u):=\mu^4-\mu^3-(tu+1)\mu^2+\mu(tu-u^3)+tu
\]
及其分歧系统
\(
F_t(\mu,u)=\partial_\mu F_t(\mu,u)=0
\)
为背景记号，抽取两类内禀对象：由双根约束诱导的有理消元律与平面代数曲线化表述，以及由 Puiseux 分歧诱导的 Lee--Yang 斜率平方不变量 \(b_\ast\)。
在主导并合分支点处，它们进一步给出 \(\widetilde M_m(u)\) 的两类普适定标核与二次晶格型零点律。
\par
在折叠语义的 Bernoulli 半 Perron 谱模型中，沿用 \eqref{eq:fold-bernoulli-half-perron-quartic} 的记号并取特化 \(t=2\)，记
\(
F(\mu,u):=F_2(\mu,u)
\)；
本段后续关于 \(Q_{10}\)、\(P_{10}\) 与 \(b_\ast\) 的对象均在该特化下讨论。

\begin{proposition}[双根约束的线性消元闭式与降阶 Tschirnhaus 同余]\label{prop:fold-gauge-anomaly-branchpoint-rational-reduction}
设 \(u\neq 0\)。若存在 \(\mu\in\CC\) 使得 \(\mu\) 为 \(F_t(\cdot,u)\) 的重根，即
\[
F_t(\mu,u)=0,\qquad \partial_\mu F_t(\mu,u)=0,
\]
则必有 \(\mu\neq \pm i\)。令
\[
X:=tu,\qquad Y:=u^3,
\]
则 \((X,Y)\) 由 \(\mu\) 唯一确定，并满足显式有理公式
\[
\boxed{\;
X=\frac{\mu^2(3\mu^2-2\mu-1)}{\mu^2+1},\qquad
Y=-\frac{2\mu(\mu^2-\mu-1)^2}{\mu^2+1}\; }.
\]
反之，给定任意 \(\mu\neq \pm i\)，若存在 \((t,u)\) 使得 \(tu=X\) 且 \(u^3=Y\)，则 \(\mu\) 必为 \(F_t(\cdot,u)\) 的重根。
\par
在特化 \(t=2\) 且 \(\mu\neq \pm i\) 时，由 \(u=X/t\) 得到降阶有理消元律
\begin{equation}\label{eq:fold-gauge-anomaly-u-rational-reduction}
\boxed{\;
u=\mathcal R(\mu):=\frac{\mu^2(3\mu^2-2\mu-1)}{2(\mu^2+1)}\; }.
\end{equation}
进一步，令 \(Q_{10}(\mu)\) 与 \(P_{10}(u)\) 如命题 \ref{prop:fold-gauge-anomaly-Q10-tschirnhaus}。
则对任意根 \(\mu_\ast\) 满足 \(Q_{10}(\mu_\ast)=0\)，都有 \(u_\ast=\mathcal R(\mu_\ast)\) 满足 \(P_{10}(u_\ast)=0\)；
并且 \(\mathcal R\) 与命题 \ref{prop:fold-gauge-anomaly-Q10-tschirnhaus} 中的 Tschirnhaus 变换 \(T(\mu)\) 在 \(Q_{10}=0\) 的根集上严格一致：存在整系数恒等式
\begin{equation}\label{eq:fold-gauge-anomaly-R-minus-T-congruence}
\boxed{\;
\mathcal R(\mu)-T(\mu)=\frac{(7\mu+1)\,Q_{10}(\mu)}{200(\mu^2+1)}\; }.
\end{equation}
\end{proposition}
\leanverified{paper\_fold\_gauge\_anomaly\_branchpoint\_rational\_reduction}
\begin{proof}
将 \(F_t(\mu,u)=0\) 与 \(\partial_\mu F_t(\mu,u)=0\) 中出现的 \(tu\) 与 \(u^3\) 记为未知量 \(X,Y\)。
两式等价改写为关于 \((X,Y)\) 的二元一次方程组
\[
\mu^4-\mu^3-\mu^2-\mu Y+(-\mu^2+\mu+1)X=0,
\qquad
4\mu^3-3\mu^2-2\mu-Y+(1-2\mu)X=0.
\]
其系数行列式为 \(\mu^2+1\)，故当且仅当 \(\mu\neq \pm i\) 时该方程组有唯一解。
直接消元得到所示 \(X(\mu),Y(\mu)\) 的闭式表达；再将解代回方程组即得“反向”结论。
\par
对 \eqref{eq:fold-gauge-anomaly-u-rational-reduction}，取特化 \(t=2\) 并由 \(u=X/t\) 立即得到。
\par
对 \eqref{eq:fold-gauge-anomaly-R-minus-T-congruence}，将 \(T(\mu)\) 取为命题 \ref{prop:fold-gauge-anomaly-Q10-tschirnhaus} 的显式九次多项式除以 \(200\)，通分并化简即可得到所示整系数恒等式；
右端显式含因子 \(Q_{10}(\mu)\)，故在 \(Q_{10}(\mu)=0\) 的根集上两者一致。
从而对任意 \(Q_{10}\)-根 \(\mu_\ast\)，有 \(\mathcal R(\mu_\ast)=T(\mu_\ast)\)，并由命题 \ref{prop:fold-gauge-anomaly-Q10-tschirnhaus} 得 \(P_{10}(\mathcal R(\mu_\ast))=0\)。证毕。
\end{proof}

\begin{corollary}[\texorpdfstring{$\mu=-1$}{mu=-1} 分歧共振的全复域分类与实刚性]\label{cor:fold-gauge-anomaly-mu-minus1-branch-classification}
在命题 \ref{prop:fold-gauge-anomaly-branchpoint-rational-reduction} 的设定下，若重根取 \(\mu=-1\)，则双根条件等价于
\[
tu=2,\qquad u^3=1.
\]
因此 \(\mu=-1\) 的重根事件在复参数域中恰好发生于三点族
\[
\boxed{\ (t,u)=\bigl(2/\zeta,\ \zeta\bigr),\qquad \zeta^3=1\ },
\]
并且在实参数截面上仅有 \((t,u)=(2,1)\)。
\end{corollary}
\begin{proof}
由命题 \ref{prop:fold-gauge-anomaly-branchpoint-rational-reduction} 的显式公式直接代入 \(\mu=-1\) 得 \(X=2\)、\(Y=1\)，等价于 \(tu=2\)、\(u^3=1\)。
反向由定义代回可知此时 \(\mu=-1\) 同时使 \(F_t\) 与 \(\partial_\mu F_t\) 为零。
实参数情形 \(u^3=1\Rightarrow u=1\) 从而 \(t=2\)。证毕。
\end{proof}
\leanverified{paper\_fold\_gauge\_anomaly\_mu\_minus1\_branch\_classification}

\begin{proposition}[\texorpdfstring{$u$}{u}-判别式的 \texorpdfstring{$\mu$}{mu}-降维与有理回收]\label{prop:fold-gauge-anomaly-branchpoint-mu-reduction-Et}
\leanverified{paper\_fold\_gauge\_anomaly\_branchpoint\_mu\_reduction\_et}
设 \(t\neq 0\)。则 \(u\) 为多项式 \(F_t(\cdot,u)\) 在 \(\mu\)-方向出现重根的分歧点当且仅当存在 \(\mu\neq \pm i\) 满足单一方程
\[
\boxed{\;
E_t(\mu):=
2t^3(\mu^2+1)^2(\mu^2-\mu-1)^2+\mu^5(3\mu^2-2\mu-1)^3=0\; }.
\]
在该情形下分歧点可由 \(\mu\) 有理回收为
\[
\boxed{\;
u=\frac{\mu^2(3\mu^2-2\mu-1)}{t(\mu^2+1)}\; }.
\]
特别地，当 \(t=2\) 时，
\[
E_2(\mu)=(\mu+1)\,Q_{10}(\mu),
\]
其中 \(Q_{10}\) 与命题 \ref{prop:fold-gauge-anomaly-Q10-tschirnhaus} 的定义一致，并且 \(\mu=-1\) 对应分歧点 \(u=1\)。
\end{proposition}
\begin{proof}
由命题 \ref{prop:fold-gauge-anomaly-branchpoint-rational-reduction}，在 \(\mu\neq \pm i\) 下双根条件等价于 \(X=tu\) 与 \(Y=u^3\) 取为所给有理表达。
由 \(t\neq 0\) 得
\(
u=X/t=\mu^2(3\mu^2-2\mu-1)/(t(\mu^2+1))
\)。
将其立方并与 \(Y\) 相等，清分母即得 \(E_t(\mu)=0\)。
反向：若 \(E_t(\mu)=0\)，以上定义的 \(u\) 满足 \(u^3=Y\)，从而 \((X,Y)\) 同时满足二元一次方程组，等价于 \(F_t(\mu,u)=\partial_\mu F_t(\mu,u)=0\)。
\par
当 \(t=2\) 时，对 \(E_2(\mu)\) 作整式因式分解得到 \((\mu+1)Q_{10}(\mu)\)；并由 \(\mu=-1\) 时的有理回收立得 \(u=1\)。证毕。
\end{proof}

\begin{proposition}[分歧曲线在单项式坐标 \texorpdfstring{$(X,Y)=(tu,u^3)$}{(X,Y)=(tu,u^3)} 下的显式代数方程与有理性]\label{prop:fold-gauge-anomaly-branch-curve-XY-explicit}
\leanverified{paper\_fold\_gauge\_anomaly\_branch\_curve\_xy\_explicit}
令 \(X=tu\)、\(Y=u^3\)。则 \(F_t(\cdot,u)\) 在 \(\mu\)-方向出现重根的充要条件等价于平面代数曲线
\[
R(X,Y)=0,
\]
其中 \(R\in\ZZ[X,Y]\) 可取为如下显式不可约多项式：
\[
\begin{aligned}
R(X,Y)=\;&20X^5-8X^4Y-120X^4+4X^3Y^2+236X^3Y+220X^3\\
&-239X^2Y^2-164X^2Y-120X^2+90XY^3-70XY^2-108XY+20X\\
&-27Y^4-22Y^3+5Y^2.
\end{aligned}
\]
并且该曲线为有理曲线：由 \(\mu\neq \pm i\) 给出有理参数化
\[
X(\mu)=\frac{\mu^2(3\mu^2-2\mu-1)}{\mu^2+1},\qquad
Y(\mu)=-\frac{2\mu(\mu^2-\mu-1)^2}{\mu^2+1}.
\]
\end{proposition}
\begin{proof}
命题 \ref{prop:fold-gauge-anomaly-branchpoint-rational-reduction} 给出 \((X,Y)\) 关于 \(\mu\) 的有理表达；消去参数 \(\mu\)（例如取两式的结果式）得到某个非零整系数消元多项式 \(R\) 使 \(R(X(\mu),Y(\mu))\equiv 0\)。
对所给显式 \(R\) 直接检验该恒等式成立。
由于给出了有理参数化，曲线函数域同构于 \(\CC(\mu)\)，从而曲线为有理曲线；不可约性可由参数化的函数域同构立即推出。证毕。
\end{proof}

\leanverified{paper\_fold\_gauge\_anomaly\_resonant\_branchpoint\_series}
\begin{proposition}[\texorpdfstring{$(t,u)=(2,1)$}{(t,u)=(2,1)} 共振点附近的解析分歧支与三阶展开]\label{prop:fold-gauge-anomaly-resonant-branchpoint-series}
考虑推论 \ref{cor:fold-gauge-anomaly-mu-minus1-branch-classification} 的实共振点 \((t,u)=(2,1)\)。
则存在唯一解析函数 \(u_b(t)\)（定义在 \(t=2\) 的充分小邻域内）满足 \(u_b(2)=1\) 且 \(F_t(\cdot,u_b(t))\) 在 \(\mu\)-方向出现重根；
相应的重根位置 \(\mu_b(t)\) 亦解析，且局部展开满足
\[
\boxed{\;
u_b(t)=1+(t-2)-\frac12(t-2)^2+\frac78(t-2)^3+O\!\bigl((t-2)^4\bigr)\;},
\]
\[
\boxed{\;
\mu_b(t)=-1-\frac12(t-2)+\frac{3}{16}(t-2)^2-\frac{23}{64}(t-2)^3+O\!\bigl((t-2)^4\bigr)\; }.
\]
从而 \(u_b'(2)=1\)。
\end{proposition}
\begin{proof}[证明要点]
由命题 \ref{prop:fold-gauge-anomaly-branchpoint-rational-reduction}，分歧曲线可在 \(\mu=-1\) 的邻域用参数 \(\mu\) 表示为 \(X(\mu)=tu\)、\(Y(\mu)=u^3\)。
在 \(Y(-1)=1\) 的邻域取满足 \(u(-1)=1\) 的解析分支 \(u(\mu)=Y(\mu)^{1/3}\)，并令 \(t(\mu)=X(\mu)/u(\mu)\)；
则 \((t(\mu),u(\mu))\) 参数化了穿过 \((2,1)\) 的分歧曲线局部支。
对 \(t(\mu)\)、\(u(\mu)\) 在 \(\mu=-1\) 处作泰勒展开并反演级数得到所示三阶展开；\(\mu_b(t)\) 由同一反演过程给出。证毕。
\end{proof}

\leanverified{paper\_fold\_gauge\_anomaly\_leyang\_slope\_square\_minpoly}
\begin{proposition}[Lee--Yang 斜率平方不变量与十次整系数极小多项式]\label{prop:fold-gauge-anomaly-leyang-slope-square-minpoly}
令 \((\mu_\ast,u_\ast)\) 为任一非平凡分歧点，满足
\[
u_\ast\notin\{0,1\},\qquad
F(\mu_\ast,u_\ast)=\partial_\mu F(\mu_\ast,u_\ast)=0,
\qquad
\partial_uF(\mu_\ast,u_\ast)\neq 0,\ \partial_{\mu\mu}^2F(\mu_\ast,u_\ast)\neq 0.
\]
定义 Puiseux 分歧诱导的斜率平方不变量
\begin{equation}\label{eq:fold-gauge-anomaly-bstar-definition}
\boxed{\;
b_\ast:=
\frac{-2\,\partial_u F(\mu_\ast,u_\ast)}{\partial_{\mu\mu}^2F(\mu_\ast,u_\ast)\,\mu_\ast^{2}}\; }.
\end{equation}
则 \(b_\ast\) 为下列十次整系数多项式的根：
\begin{equation}\label{eq:fold-gauge-anomaly-bstar-minpoly}
\boxed{\;
\begin{aligned}
0=\;&316283904\, b^{10}-52713984\, b^{9}+379196672\, b^{8}+406374016\, b^{7}\\
&+670701296\, b^{6}+1305085464\, b^{5}+761320448\, b^{4}+12917152\, b^{3}\\
&+53372880\, b^{2}+7820550\, b-1303425.
\end{aligned}}
\end{equation}
\end{proposition}
\begin{proof}
由命题 \ref{prop:fold-gauge-anomaly-branchpoint-rational-reduction}，任一非平凡分歧点满足 \(u_\ast=\mathcal R(\mu_\ast)\)，故 \eqref{eq:fold-gauge-anomaly-bstar-definition} 的 \(b_\ast\) 可化为 \(\mu_\ast\) 的有理函数值。
另一方面，\(\mu_\ast\) 必满足命题 \ref{prop:fold-gauge-anomaly-Q10-tschirnhaus} 的约束 \(Q_{10}(\mu_\ast)=0\)。
将方程 \(b-\mathcal B(\mu)=0\)（其中 \(\mathcal B\) 为上述有理表达）与 \(Q_{10}(\mu)=0\) 对 \(\mu\) 作结果式消元，得到仅含 \(b\) 的十次整系数多项式；化简后即为 \eqref{eq:fold-gauge-anomaly-bstar-minpoly}。证毕。
\end{proof}

\leanverified{paper\_fold\_gauge\_anomaly\_leyang\_universal\_kernels}
\begin{theorem}[Lee--Yang 边界的两类普适定标核：\texorpdfstring{$\sinh(w)/w$}{sinh(w)/w} 与 \texorpdfstring{$\cosh(w)$}{cosh(w)}]\label{thm:fold-gauge-anomaly-leyang-universal-kernels}
沿用命题 \ref{prop:fold-gauge-anomaly-mgf-order4-recurrence} 的记号，令 \(\widetilde M_m(u):=2^mM_m(u)\)。
取一非平凡分歧点 \((\mu_\ast,u_\ast)\) 位于主导谱支并合处，即在 \(u=u_\ast\) 时四根中存在两根作并合且其模长达到最大值 \(|\mu_\ast|\)（例如命题 \ref{prop:fold-gauge-anomaly-branch-point-classification} 的 Lee--Yang edge 共轭对）。
令 \(b_\ast\) 为 \eqref{eq:fold-gauge-anomaly-bstar-definition}。
则存在常数 \(C_1,C_0\in\CC\)（仅依赖于 \(\widetilde M_m\) 的初值与端点泛函，而与 \(m\) 无关），使得对任意紧集 \(K\subset\CC\)，当 \(m\to\infty\) 时在 \(w\in K\) 上一致成立以下二择一结论：
\begin{enumerate}
  \item 若 \(\displaystyle \lim_{m\to\infty}\frac{\widetilde M_m(u_\ast)}{m\,\mu_\ast^{m}}=C_1\neq 0\)，则
  \begin{equation}\label{eq:fold-gauge-anomaly-leyang-sinh-kernel}
  \boxed{\;
  \lim_{m\to\infty}\frac{\widetilde M_m\!\left(u_\ast+\frac{w^2}{b_\ast m^2}\right)}{m\,\mu_\ast^{m}}
  =2C_1\,\frac{\sinh(w)}{w}\; }.
  \end{equation}
  \item 若 \(\displaystyle \lim_{m\to\infty}\frac{\widetilde M_m(u_\ast)}{\mu_\ast^{m}}=C_0\neq 0\)，则
  \begin{equation}\label{eq:fold-gauge-anomaly-leyang-cosh-kernel}
  \boxed{\;
  \lim_{m\to\infty}\frac{\widetilde M_m\!\left(u_\ast+\frac{w^2}{b_\ast m^2}\right)}{\mu_\ast^{m}}
  =2C_0\,\cosh(w)\; }.
  \end{equation}
\end{enumerate}
\end{theorem}
\begin{proof}[证明要点]
在 \(u_\ast\) 的去心邻域内取局部参数 \(s=\sqrt{u-u_\ast}\)，由分歧非退化性得到两条主导谱支的 Puiseux 展开
\(
\mu_\pm(s)=\mu_\ast\exp(\pm \sqrt{b_\ast}\,s+O(s^2))
\)。
由四阶递推的谱分解，\(\widetilde M_m(u)\) 在该邻域可写为四个谱支贡献之和；其中两项为 \(\mu_\pm(s)^m\)，其余两项在 \(u_\ast\) 处模长严格小于 \(|\mu_\ast|\)，从而在 \(m\to\infty\) 的主导定标下指数可忽略。
由于 \(\widetilde M_m(u(s))\) 为 \(s\) 的偶函数，其在 \(\mu_\pm(s)^m\) 的系数允许出现至多一阶极点，故可整理为
\[
\widetilde M_m(u(s))
=a_{-1}(s^2)\frac{\mu_+(s)^m-\mu_-(s)^m}{s}
+a_0(s^2)\bigl(\mu_+(s)^m+\mu_-(s)^m\bigr)
+O(\rho^m),
\]
其中 \(|\rho|<|\mu_\ast|\) 且 \(a_{-1},a_0\) 在 \(s^2\) 上解析。
令 \(w=m\sqrt{b_\ast}\,s\)（等价于 \(u=u_\ast+w^2/(b_\ast m^2)\)），并用指数近似
\(
\mu_\pm(s)^m=\mu_\ast^m e^{\pm w}(1+o(1))
\)
即可得到 \eqref{eq:fold-gauge-anomaly-leyang-sinh-kernel} 与 \eqref{eq:fold-gauge-anomaly-leyang-cosh-kernel} 的两种主导情形。证毕。
\end{proof}

\begin{corollary}[Lee--Yang 零点的二次晶格律与方向不变量]\label{cor:fold-gauge-anomaly-leyang-quadratic-lattice-zeros}
在定理 \ref{thm:fold-gauge-anomaly-leyang-universal-kernels} 的设定下，存在 \(\delta_0>0\) 使得当 \(m\) 足够大时，\(\widetilde M_m(u)\) 在圆盘 \(D(u_\ast,\delta_0)\) 内的零点满足二次晶格定律：
\begin{enumerate}
  \item 在 \eqref{eq:fold-gauge-anomaly-leyang-sinh-kernel} 的情形下，对每个固定 \(k\in\NN\) 存在零点 \(u_{m,k}\in D(u_\ast,\delta_0)\) 使
  \[
  u_{m,k}=u_\ast-\frac{\pi^2 k^2}{b_\ast m^2}+O\!\left(\frac{k^3}{m^3}\right)\qquad(m\to\infty).
  \]
  \item 在 \eqref{eq:fold-gauge-anomaly-leyang-cosh-kernel} 的情形下，对每个固定 \(k\in\NN\) 存在零点 \(u_{m,k}\in D(u_\ast,\delta_0)\) 使
  \[
  u_{m,k}=u_\ast-\frac{\pi^2 (k+\tfrac12)^2}{b_\ast m^2}+O\!\left(\frac{k^3}{m^3}\right)\qquad(m\to\infty).
  \]
\end{enumerate}
因此零点以 \(m^{-2}\) 速度逼近 \(u_\ast\)，其逼近方向由复向量
\[
\boxed{\ \mathrm{dir}(u_\ast)=-\frac{1}{b_\ast}\ }
\]
完全决定。
\leanverified{paper\_fold\_gauge\_anomaly\_leyang\_quadratic\_lattice\_zeros}
\end{corollary}
\begin{proof}
定理 \ref{thm:fold-gauge-anomaly-leyang-universal-kernels} 给出经归一化后的解析函数在紧集上一致收敛到
\(
2C_1\sinh(w)/w
\)
或
\(
2C_0\cosh(w)
\)。
两者零点分别为 \(w=i\pi k\ (k\in\ZZ\setminus\{0\})\) 与 \(w=i\pi(k+\tfrac12)\ (k\in\ZZ)\)。
由一致收敛与 Rouch\'e 定理，极限零点在 \(m\) 足够大时稳定地产生对应的 \(w_{m,k}\)，再由解析反变换 \(u=u_\ast+w^2/(b_\ast m^2)\) 得到所示渐近式。证毕。
\end{proof}

\begin{corollary}[数域同构：\texorpdfstring{$u_\ast,\mu_\ast,b_\ast$}{u*,mu*,b*} 生成同一十次数域]\label{cor:fold-gauge-anomaly-field-equality-u-mu-b}
\leanverified{paper\_fold\_gauge\_anomaly\_field\_equality\_u\_mu\_b}
取任一满足 \(P_{10}(u_\ast)=0\) 的非平凡分歧点，并令 \(\mu_\ast\) 为相应重根，\(b_\ast\) 如 \eqref{eq:fold-gauge-anomaly-bstar-definition}。
则
\[
\boxed{\ \QQ(u_\ast)=\QQ(\mu_\ast)=\QQ(b_\ast)\ } ,
\]
且该公共数域为 \(\QQ\) 的次数 \(10\) 扩张。
\end{corollary}
\begin{proof}
由命题 \ref{prop:fold-gauge-anomaly-P10-galois-discriminant} 与命题 \ref{prop:fold-gauge-anomaly-Q10-tschirnhaus}，多项式 \(P_{10}\) 与 \(Q_{10}\) 均为次数 \(10\) 的不可约整系数多项式，故 \([\QQ(u_\ast):\QQ]=[\QQ(\mu_\ast):\QQ]=10\)。
由命题 \ref{prop:fold-gauge-anomaly-branchpoint-rational-reduction} 的有理表达 \(u_\ast=\mathcal R(\mu_\ast)\) 得 \(\QQ(u_\ast)\subseteq \QQ(\mu_\ast)\)，从而两者必相等。
另一方面由 \eqref{eq:fold-gauge-anomaly-bstar-definition} 得 \(b_\ast\in\QQ(\mu_\ast)\)；
而命题 \ref{prop:fold-gauge-anomaly-leyang-slope-square-minpoly} 给出 \(b_\ast\) 满足次数 \(10\) 的不可约整系数多项式，故 \([\QQ(b_\ast):\QQ]=10\)。
于是 \(\QQ(b_\ast)\subseteq\QQ(\mu_\ast)\) 强制 \(\QQ(b_\ast)=\QQ(\mu_\ast)\)。证毕。
\end{proof}

% --- Distillation writeback ---
\begin{proposition}[Lee--Yang 分歧定标的 Puiseux 缺口滤过与终端残数判据]\label{distill:grothendieck-cohomological_gap_filtration-folding-leyang-puiseux-terminal-residue}
沿用定理 \ref{thm:fold-gauge-anomaly-leyang-universal-kernels} 的非退化主导分歧点 \((\mu_\ast,u_\ast)\)，并取一支平方根 \(\sqrt{b_\ast}\)。令
\[
u(s):=u_\ast+s^2,
\qquad
\mu_\pm(s)=\mu_\ast\exp\bigl(\pm\sqrt{b_\ast}\,s+O(s^2)\bigr)
\]
为并合的两条 Puiseux 谱支。设一族局部审计量 \(G_m\) 在 \(s=0\) 的去心邻域内具有有限 Puiseux 缺口展开
\[
\begin{aligned}
G_m(\nu(s))=\;&a_{-1}(s^2)\frac{\mu_+(s)^m-\mu_-(s)^m}{s}
+a_0(s^2)\bigl(\mu_+(s)^m+\mu_-(s)^m\bigr)\\
&+\mu_\ast^m\sum_{\ell=1}^{L}s^\ell a_\ell(s^2)+R_m(s),
\end{aligned}
\]
其中 \(a_{-1},a_0,a_1,\ldots,a_L\) 在 \(0\) 邻域解析，且对任意紧集 \(K\subset\CC\) 有
\[
\sup_{w\in K}\left|\frac{R_m\bigl(w/(m\sqrt{b_\ast})\bigr)}{\mu_\ast^m}\right|\longrightarrow0.
\]
定义两条终端传播算子
\[
\partial_{-1}(c)(w):=2\sqrt{b_\ast}\,c\,\frac{\sinh w}{w},
\qquad
\partial_0(c)(w):=2c\cosh w .
\]
则在任意紧集 \(K\subset\CC\) 上一致成立
\[
\lim_{m\to\infty}
\frac{G_m\!\left(u_\ast+\frac{w^2}{b_\ast m^2}\right)}{m\mu_\ast^m}
=\partial_{-1}\bigl(a_{-1}(0)\bigr)(w).
\]
若进一步 \(a_{-1}(0)=0\)，则
\[
\lim_{m\to\infty}
\frac{G_m\!\left(u_\ast+\frac{w^2}{b_\ast m^2}\right)}{\mu_\ast^m}
=\partial_0\bigl(a_0(0)\bigr)(w).
\]
因此所有正 Puiseux 阶 \(s^\ell a_\ell(s^2)\) 都是 Lee--Yang 定标下的有限部伪层；它们在终端核中塌缩为零。相反，极点残数层 \(a_{-1}(0)\) 不被任何有限正阶 jet 的消失所控制，并且在 \(\mu_\ast^m\)-归一化下必须先被杀死，方可进入 \(\cosh\)-型终端层。
\end{proposition}
\begin{proof}
取 \(s=w/(m\sqrt{b_\ast})\)。由 Puiseux 展开式，紧集上一致有
\[
\frac{\mu_\pm(s)^m}{\mu_\ast^m}
=\exp\bigl(\pm w+mO(s^2)\bigr)=e^{\pm w}(1+o(1)),
\]
因为 \(m s^2=O(m^{-1})\)。于是
\[
\frac{1}{m\mu_\ast^m}\,
 a_{-1}(s^2)\frac{\mu_+(s)^m-\mu_-(s)^m}{s}
=a_{-1}(s^2)\frac{e^w-e^{-w}+o(1)}{ms}
\longrightarrow 2\sqrt{b_\ast}\,a_{-1}(0)\frac{\sinh w}{w}.
\]
同一归一化下，\(a_0\)-项为 \(O(m^{-1})\)，正阶 Puiseux 项为 \(O(m^{-1})\)，余项由假设为 \(o(m^{-1})\)，故得到第一条极限。

若 \(a_{-1}(0)=0\)，则 \(a_{-1}(s^2)=O(s^2)\)。在 \(\mu_\ast^m\)-归一化下，极点项为
\[
O(s^2)\frac{e^w-e^{-w}+o(1)}{s}=O(s)\longrightarrow0.
\]
而
\[
\frac{a_0(s^2)(\mu_+(s)^m+\mu_-(s)^m)}{\mu_\ast^m}
\longrightarrow a_0(0)(e^w+e^{-w})=2a_0(0)\cosh w.
\]
每个正阶项含因子 \(s^\ell=O(m^{-\ell})\)，故消失；余项仍由假设消失。若 \(a_{-1}(0)\ne0\)，则同一计算显示 \(\mu_\ast^m\)-归一化下极点项在一般 \(w\) 上具有 \(m\)-级增长，故不能被正阶有限 jet 抵消。命题得证。
\end{proof}

% --- End distillation writeback ---

\endinput


\begin{corollary}[判别式平方比值的同域证书]\label{cor:fold-gauge-anomaly-disc-square-ratio}
\leanverified{paper\_fold\_gauge\_anomaly\_disc\_square\_ratio}
在命题 \ref{prop:fold-gauge-anomaly-P10-galois-discriminant} 与命题 \ref{prop:fold-gauge-anomaly-Q10-tschirnhaus} 的记号下，
\[
\boxed{\ \frac{\mathrm{Disc}(P_{10})}{\mathrm{Disc}(Q_{10})}=\Bigl(\frac{25054688907}{500}\Bigr)^{2}\ }.
\]
因此，即使不显式写出极大整环的整基，上式仍在纯判别式层面给出同域证书：两者所诱导的整环阶之差异仅体现在指数的平方因子上。
\end{corollary}

\begin{proposition}[速率曲线消元模型的判别式平方因子分解]\label{prop:fold-gauge-anomaly-rate-curve-discriminant-squarefactor}
\leanverified{paper\_fold\_gauge\_anomaly\_rate\_curve\_discriminant\_squarefactor}
令 \(R(\alpha,u)\in\ZZ[\alpha,u]\) 为命题 \ref{prop:fold-gauge-anomaly-rate-curve-param} 的消元证书（表 \ref{tab:fold_gauge_anomaly_rate_curve_resultant_degree}），其定义使得速率曲线由 \(R(\alpha,u)=0\) 刻画且 \(\deg_{\alpha}R=4\)。
则其关于 \(\alpha\) 的判别式满足精确分解
\[
\boxed{\ \mathrm{Disc}_{\alpha}\bigl(R(\alpha,u)\bigr)=-(u^{5}H(u))^{2}\cdot u(u-1)P_{10}(u)\ },
\]
其中 \(H(u)\in\ZZ[u]\) 为次数 \(19\) 的不可约多项式
\[
\begin{aligned}
H(u)=\;&137781u^{19}+629856u^{18}+934578u^{17}-1546209u^{16}-6187752u^{15}+1402704u^{14}\\
&+15349014u^{13}-3368736u^{12}-17688806u^{11}+2216732u^{10}+17425320u^{9}-4765670u^{8}\\
&-11620472u^{7}+7448180u^{6}+2529720u^{5}-4736240u^{4}+2386300u^{3}-612800u^{2}+81800u-4500.
\end{aligned}
\]
因此，速率曲线作为 \(u\)-直线上的次数 \(4\) 覆盖，其平方自由判别式与谱四次曲线的平方自由判别式一致，分歧支撑同为 \(u(u-1)P_{10}(u)=0\)；
而 \(H(u)^{2}\) 以纯平方因子出现，反映了消元坐标环与其整闭包之间的非平凡指数现象，而非分歧的消失。
\end{proposition}

\begin{proposition}[指数平方因子的算术一般性：\texorpdfstring{$H$}{H} 的满对称伽罗瓦群与判别式]\label{prop:fold-gauge-anomaly-H-galois-discriminant}
令 \(H(u)\) 如命题 \ref{prop:fold-gauge-anomaly-rate-curve-discriminant-squarefactor}。
则
\[
\mathrm{Gal}\bigl(H/\QQ\bigr)\cong S_{19},
\]
并且其判别式分解为
\[
\begin{aligned}
\mathrm{Disc}(H)=\;&2^{40}3^{91}5^{32}\cdot 101\cdot 173^{2}\cdot 1931^{10}\cdot 9013^{3}\\
&\cdot 795449^{6}\cdot 15208978734786979108724207.
\end{aligned}
\]
该性质表明：\(\mathrm{Disc}_{\alpha}(R)\) 中的平方因子并非来自低阶可解对称性，而在伽罗瓦意义上同样处于满对称的一般情形。
\end{proposition}

\begin{conjecture}[平方因子的满对称泛型性与 Chebotarev 独立性]\label{conj:fold-gauge-anomaly-squarefactor-genericity-independence}
设 \(R(\alpha,u)\in\ZZ[\alpha,u]\) 为一类消元证书，并假设其关于 \(\alpha\) 的判别式在 \(\ZZ[u]\) 中分解为
\[
\mathrm{Disc}_{\alpha}(R)=\textup{(显式幂因子)}\cdot B(u)^{2e}\cdot C(u)^{2f},
\]
其中 \(B(u)\) 为平方自由因子（承载实际分歧支撑），而 \(C(u)\) 为以纯平方指数出现的因子（反映整闭性或指数障碍）。
在不可约且非退化的泛型条件下，预期成立：
\begin{enumerate}
  \item \(\mathrm{Gal}(B/\QQ)\cong S_{\deg B}\) 且 \(\mathrm{Gal}(C/\QQ)\cong S_{\deg C}\)；
  \item \(\mathrm{Spl}(B)\cap \mathrm{Spl}(C)=\QQ\)（等价于线性互素性）；
  \item 因而 \(B\bmod p\) 与 \(C\bmod p\) 的分解型在 Chebotarev 意义下渐近独立。
\end{enumerate}
该猜想在本节的显式样本 \(B=P_{10}\)、\(C=H\) 上由命题 \ref{prop:fold-gauge-anomaly-P10-galois-discriminant}、命题 \ref{prop:fold-gauge-anomaly-H-galois-discriminant} 与命题 \ref{prop:fold-gauge-anomaly-P10-H-linear-disjointness} 实现为可核验实例。
\end{conjecture}

\begin{proposition}[谱四次的三次预解式与判别式恒等]\label{prop:fold-gauge-anomaly-cubic-resolvent-discriminant}
\leanverified{paper\_fold\_gauge\_anomaly\_cubic\_resolvent\_discriminant}
对谱四次 \(F(\mu,u)\) 取标准三次预解式，得到
\[
\mathcal{R}(y,u)=y^{3}+(1+2u)y^{2}+(u^{3}-10u)y-\bigl(u^{6}-4u^{4}+20u^{2}+10u\bigr).
\]
则存在判别式恒等式
\[
\boxed{\ \mathrm{Disc}_{y}\bigl(\mathcal{R}(y,u)\bigr)=\mathrm{Disc}_{\mu}\bigl(F(\mu,u)\bigr)=-u(u-1)P_{10}(u)\ }.
\]
因此谱四次的 \(S_{4}\) 伽罗瓦闭包在 \(\QQ(u)\) 上携带由 \(\mathcal{R}\) 显式给出的规范 \(S_{3}\) 中间扩张，其分歧位置与谱四次曲线完全一致。
\end{proposition}

\begin{proposition}[一参数变形的分歧共振唯一性与泛性伽罗瓦群]\label{prop:fold-gauge-anomaly-branch-resonance-one-parameter}
将命题 \ref{prop:fold-gauge-anomaly-pressure} 中矩阵 \(2A_{\theta}\) 的 \((3,2)\) 元从 \(2\) 变形为参数 \(t\)，得到谱多项式
\[
F_{t}(\mu,u)=\mu^{4}-\mu^{3}-(tu+1)\mu^{2}+\mu(tu-u^{3})+tu.
\]
其关于 \(\mu\) 的判别式可写为
\[
\mathrm{Disc}_{\mu}\bigl(F_{t}(\mu,u)\bigr)=u\cdot D_{t}(u),
\]
其中 \(D_{t}(u)\in\ZZ[t][u]\) 为关于 \(u\) 的次数 \(11\) 多项式。
并且满足以下刚性算术性质：
\begin{enumerate}
\item 在实参数 \(t\in\RR\) 中，\(u=1\) 为分支点（等价于 \(D_{t}(1)=0\)）的唯一取值为 \(t=2\)；且在该点出现分解
\[
\boxed{\ D_{2}(u)=-(u-1)P_{10}(u)\ }.
\]
\item 在 \(t=1\) 时，\(D_{1}(u)\) 在 \(\QQ\) 上不可约且
\(
\mathrm{Gal}\bigl(D_{1}/\QQ\bigr)\cong S_{11}
\)；
由 Hilbert 不可约性可推出该一参数族在 \(\QQ(t)\) 上的通用分歧多项式具有泛性伽罗瓦群 \(S_{11}\)。
\end{enumerate}
因此 \(t=2\) 情形出现的异常可约分歧具有严格的唯一性框架：它是该一参数族中的唯一实分歧共振点。
\end{proposition}

\paragraph{单值群、符号二次商与 \texorpdfstring{$S_4$}{S4}-Hurwitz 塔}\label{par:fold-gauge-anomaly-s4-hurwitz-tower}
以下把命题 \ref{prop:fold-gauge-anomaly-spectral-quartic-geometry}--\ref{prop:fold-gauge-anomaly-cubic-resolvent-discriminant} 中给出的“谱四次—判别式—三次预解式”数据，转写为覆叠的 Hurwitz 型、泛型伽罗瓦群以及关键中间商曲线的亏格。
为与外部记号对齐，记
\[
D_{10}(u):=P_{10}(u),\qquad D_{19}(u):=H(u).
\]

\begin{conclusion}[规范异常四次谱曲线的几何型]\label{con:fold-gauge-anomaly-spectral-quartic-geomtype}
令 \(\mathcal C_F\subset\PP^2\) 为谱四次 \(F(\mu,u)=0\) 的总次数齐次化闭包（命题 \ref{prop:fold-gauge-anomaly-spectral-quartic-geometry} 中的曲线 \(C\)）。
则 \(\mathcal C_F\) 在特征 \(0\) 上无奇点，因而为非超椭圆的平面四次曲线，并满足
\[
g(\mathcal C_F)=3.
\]
\end{conclusion}

\begin{conclusion}[判别式的完全分解与简单分支]\label{con:fold-gauge-anomaly-disc-factorization-simple-branch}
将 \(F(\mu,u)\) 视为关于 \(\mu\) 的四次多项式，则其判别式满足
\[
\mathrm{Disc}_{\mu}\bigl(F(\mu,u)\bigr)=-u(u-1)D_{10}(u),
\qquad D_{10}(u)=27u^{10}+27u^{9}-153u^{8}-163u^{7}+793u^{6}+1061u^{5}-832u^{4}-816u^{3}+1320u^{2}-440u+40,
\]
并且 \(u(u-1)D_{10}(u)\) 在 \(\ZZ[u]\) 中平方因子自由。
因此投影 \(\pi:\mathcal C_F\to\PP^1_u\) 为次数 \(4\) 覆叠，其分支集合恰为
\[
\{u=0\}\cup\{u=1\}\cup\{u:\ D_{10}(u)=0\},
\]
且在每个分支值之上均为简单分支（惯性型为单一换位）。
\end{conclusion}
\begin{proof}[证明要点]
判别式因式分解见命题 \ref{prop:fold-gauge-anomaly-discriminant-factorization}。
平方因子自由性等价于
\(
\gcd\!\bigl(D_{10},D_{10}'\bigr)=\gcd(D_{10},u)=\gcd(D_{10},u-1)=1
\)
在 \(\ZZ[u]\) 中成立；从而各分支值均对应 \(\mu\)-纤维中的二重并合且无更高阶并合，故为简单分支。
\end{proof}

\begin{conclusion}[特征多项式的泛型伽罗瓦群为 \texorpdfstring{$S_4$}{S4}]\label{con:fold-gauge-anomaly-generic-galois-s4}
函数域扩张 \(\QQ(u)\subset \QQ(u,\mu)\)（由 \(F(\mu,u)=0\) 给出）之泛型伽罗瓦群满足
\[
\mathrm{Gal}\bigl(F(\mu,u)/\QQ(u)\bigr)\cong S_4.
\]
等价地，覆叠 \(\pi:\mathcal C_F\to\PP^1_u\) 的几何单值群为满对称群 \(S_4\)。
更进一步，设 \(L\) 为 \(F(\mu,u)\) 在 \(\QQ(u)\) 上的分裂域，则该扩张为正则：
\[
L\cap \overline{\QQ}=\QQ,
\]
亦即其常数域等于 \(\QQ\)。
\end{conclusion}
\begin{proof}[证明要点（可核验特化证书）]
由命题 \ref{prop:fold-gauge-anomaly-cubic-resolvent-discriminant}，谱四次的三次预解式 \(\mathcal R(y,u)\) 满足
\(
\mathrm{Disc}_y(\mathcal R)=\mathrm{Disc}_{\mu}(F)=-u(u-1)D_{10}(u)
\)，其在 \(\QQ(u)\) 中非平方，故伽罗瓦群不包含于 \(A_4\)。
取 \(u=2\) 的特化，四次多项式 \(F(\mu,2)\in\ZZ[\mu]\) 在模 \(3\) 下不可约，而在模 \(11\) 下分解型为 \((3,1)\)，从而特化伽罗瓦群含有 \(3\)-循环并因此包含 \(A_4\)。
合并得特化伽罗瓦群为 \(S_4\)，从而泛型伽罗瓦群亦为 \(S_4\)。
此外，由 \(\mathrm{Disc}_{\mu}(F)=-u(u-1)D_{10}(u)\) 在分支集合 \(B=\{0,1\}\cup\{D_{10}=0\}\) 的每一点均具有奇阶零点可知，
\(\mathrm{Disc}_{\mu}(F)\) 在 \(\overline{\QQ}(u)\) 中亦非平方；
因此几何伽罗瓦群不包含于 \(A_4\)，从而与算术伽罗瓦群同为 \(S_4\)，并推出常数域为 \(\QQ\)。
\end{proof}

\begin{corollary}[甲板自同构的平凡性与分歧集的 \texorpdfstring{M\"obius}{Mobius} 刚性]\label{cor:fold-gauge-anomaly-deck-trivial-mobius-rigidity}
设 \(\pi:\mathcal C_F\to\PP^1_u\) 为结论 \ref{con:fold-gauge-anomaly-disc-factorization-simple-branch} 的次数 \(4\) 覆叠，并记其有限分支集合
\[
B:=\{0,1\}\cup\{u\in\PP^1(\overline{\QQ}):\ D_{10}(u)=0\}.
\]
则成立：
\begin{enumerate}
  \item 覆叠 \(\pi\) 在 \(\overline{\QQ}\) 上的甲板自同构群为平凡群。
  \item 若 \(\varphi\in\mathrm{PGL}_2(\QQ)\) 满足 \(\varphi(B)=B\)，则 \(\varphi=\mathrm{id}\)。
\end{enumerate}
\end{corollary}
\begin{proof}
(1) 由结论 \ref{con:fold-gauge-anomaly-generic-galois-s4}，\(\pi\) 的几何单值群为 \(S_4\)，并通过一般纤维的四点集给出自然置换表示 \(S_4=\mathrm{Sym}(4)\)。
甲板自同构在一般纤维上诱导置换，且与单值群作用可交换，因而属于中心化子 \(C_{\mathrm{Sym}(4)}(S_4)\)。
由于 \(S_4=\mathrm{Sym}(4)\) 且 \(\mathrm{Sym}(4)\) 的中心平凡，故 \(C_{\mathrm{Sym}(4)}(S_4)=\{1\}\)，从而甲板自同构群平凡。
\par
(2) 由命题 \ref{prop:fold-gauge-anomaly-P10-galois-discriminant}，\(D_{10}\) 在 \(\QQ\) 上不可约，因而其零点均非 \(\QQ\)-有理点。
故任意 \(\varphi\in\mathrm{PGL}_2(\QQ)\) 若满足 \(\varphi(B)=B\)，必将 \(\{0,1\}\) 保持为集合，且将
\(
\Sigma:=\{u\in\PP^1(\overline{\QQ}):\ D_{10}(u)=0\}
\)
保持为集合。
于是 \(\varphi\) 在 \(\Sigma\) 上诱导一置换 \(\sigma\in\mathrm{Sym}(\Sigma)\)。
由于 \(\varphi\) 定义于 \(\QQ\)，对任意 \(\tau\in\mathrm{Gal}(\overline{\QQ}/\QQ)\) 与任意 \(\alpha\in\Sigma\) 有
\(
\tau(\varphi(\alpha))=\varphi(\tau(\alpha))
\)，
从而 \(\sigma\) 与 \(\mathrm{Gal}(\overline{\QQ}/\QQ)\) 在 \(\Sigma\) 上的置换作用可交换。
而该置换作用之像为 \(\mathrm{Gal}(D_{10}/\QQ)\cong S_{10}\)；由于满对称群在其自然作用下的中心化子为平凡群，故 \(\sigma=\mathrm{id}\)。
于是 \(\varphi\) 逐点固定 \(\Sigma\)。
由于 \(\Sigma\) 至少包含三个互异点，M\"obius 变换在 \(\PP^1\) 上由三点像唯一确定，故 \(\varphi=\mathrm{id}\)。\qedhere
\end{proof}

\begin{conclusion}[十次分歧根域的签名与不可根式性]\label{con:fold-gauge-anomaly-P10-rootfield-signature-nonradical}
令 \(u_0\) 为 \(D_{10}(u)=P_{10}(u)\) 的任一根，并记 \(K:=\QQ(u_0)\)。
则 \([K:\QQ]=10\)，且其正规闭包的伽罗瓦群同构于 \(S_{10}\)。
因此 \(u_0\) 不可由根式表示。
此外，由 \(D_{10}\) 的实根计数（命题 \ref{prop:fold-gauge-anomaly-branch-point-classification} 或结论 \ref{con:fold-gauge-anomaly-real-complex-phase}），\(K\) 的签名为
\[
(r_1,r_2)=(2,4).
\]
特别地，对命题 \ref{prop:fold-gauge-anomaly-branch-point-classification} 所选主导共轭分支点 \(u_\ast\) 的根域 \(\QQ(u_\ast)\)，上述结论成立。
\end{conclusion}
\begin{proof}[证明要点]
\([K:\QQ]=10\) 由命题 \ref{prop:fold-gauge-anomaly-discriminant-factorization} 中 \(P_{10}\) 的不可约性给出；
正规闭包的伽罗瓦群为 \(S_{10}\) 由命题 \ref{prop:fold-gauge-anomaly-P10-galois-discriminant}。
由于 \(S_{10}\) 不可解，故根不可由根式表示。
签名由 \(D_{10}\) 的两条实根与四对共轭复根直接推出。证毕。
\end{proof}

\begin{conclusion}[分支点对数参数的超越性]\label{con:fold-gauge-anomaly-log-transcendence}
设 \(u_0\) 为 \(D_{10}(u)=0\) 的任一根，并取任意 \(\theta_0\in\CC\) 满足
\[
e^{\theta_0}=u_0.
\]
则 \(\theta_0\) 为超越数。
因而，对应于 \(\mathrm{Disc}_{\mu}P(\mu,u)=0\) 的任意有限 \(\theta\)-平面分支点（即 \(u_0\neq0\)）均为超越点。
\end{conclusion}
\begin{proof}
由结论 \ref{con:fold-gauge-anomaly-disc-factorization-simple-branch} 与命题 \ref{prop:fold-gauge-anomaly-P10-galois-discriminant}，\(u_0\) 为非零代数数。
若 \(\theta_0\) 为代数数且 \(\theta_0\neq0\)，则 Lindemann--Weierstrass 定理给出 \(e^{\theta_0}\) 为超越数，与 \(e^{\theta_0}=u_0\) 矛盾。
故 \(\theta_0\) 必为超越数。证毕。
\end{proof}

\begin{conclusion}[由谱四次专化生成无穷多 \texorpdfstring{$S_4$}{S4} 数域]\label{con:fold-gauge-anomaly-s4-number-fields-specialization}
令 \(B:=\{0,1\}\cup\{u:\ D_{10}(u)=0\}\) 为结论 \ref{con:fold-gauge-anomaly-disc-factorization-simple-branch} 的分支集合。
则存在无穷多 \(u_0\in\QQ\setminus B\)，使得专化四次多项式
\[
F(\mu,u_0)\in\QQ[\mu]
\]
在 \(\QQ\) 上的伽罗瓦群同构于 \(S_4\)。
对任意此类 \(u_0\)，其判别式满足
\[
\mathrm{Disc}_{\mu}\bigl(F(\mu,u_0)\bigr)=-u_0(u_0-1)D_{10}(u_0),
\]
并且其 \(V_4\)-不变子域可由三次预解式 \(\mathcal R(y,u)\)（命题 \ref{prop:fold-gauge-anomaly-cubic-resolvent-discriminant}）在 \(u=u_0\) 处的根域显式给出。
\end{conclusion}
\begin{proof}[证明要点]
由结论 \ref{con:fold-gauge-anomaly-generic-galois-s4} 与 Hilbert 不可约性定理，存在无穷多 \(u_0\in\QQ\) 使得特化伽罗瓦群保持为 \(S_4\)；
再要求 \(u_0\notin B\) 以避免判别式为零即可。
判别式恒等式由结论 \ref{con:fold-gauge-anomaly-disc-factorization-simple-branch} 逐点专化得到；
\(V_4\)-不变子域由命题 \ref{prop:fold-gauge-anomaly-cubic-resolvent-discriminant} 的三次预解式给出。证毕。
\end{proof}

\begin{conclusion}[符号二次子覆盖：判别式双覆盖的亏格]\label{con:fold-gauge-anomaly-sign-quadratic-cover}
对应于 \(\mathrm{Gal}(F/\QQ(u))\cong S_4\) 的符号表示之二次子域由 \(\sqrt{\mathrm{Disc}_{\mu}(F)}\) 生成，其几何模型可取为双覆盖
\[
\mathcal H:\quad w^2=-u(u-1)D_{10}(u).
\]
由于右端次数为 \(12\)，故 \(g(\mathcal H)=5\)，并且 \(\mathcal H\to\PP^1_u\) 的分支集合与结论 \ref{con:fold-gauge-anomaly-disc-factorization-simple-branch} 完全一致。
此外，由右端首项系数为 \(-27\) 可知：\(\mathcal H\) 在无穷远处的两点之最小分裂域为 \(\QQ(\sqrt{-3})\)。
\end{conclusion}

\begin{conclusion}[显式三次分辨：同一判别式核]\label{con:fold-gauge-anomaly-cubic-resolvent}
与 \(F(\mu,u)\) 相关联的三次预解式（命题 \ref{prop:fold-gauge-anomaly-cubic-resolvent-discriminant} 中的 \(\mathcal R(y,u)\)）可记为
\[
\Phi(y,u):=\mathcal R(y,u)=y^3+(1+2u)y^2+(u^3-10u)y-\bigl(u^6-4u^4+20u^2+10u\bigr),
\]
并满足判别式恒等式
\[
\mathrm{Disc}_{y}\bigl(\Phi(y,u)\bigr)=-u(u-1)D_{10}(u).
\]
因此 \(\Phi(y,u)=0\) 在 \(\PP^1_u\) 上给出一个次数 \(3\) 的不可约覆叠，其分支参数与 \(\pi:\mathcal C_F\to\PP^1_u\) 完全同位，且在每一分支值处呈现简单分支。
\end{conclusion}

\begin{conclusion}[本征谱的实–复相图由 \texorpdfstring{$D_{10}$}{D10} 的符号刻画]\label{con:fold-gauge-anomaly-real-complex-phase}
设 \(u>0\) 为实参数，并将 \(F(\mu,u)\) 在 \(\RR[\mu]\) 中的根型按实根/复根计数分类。
则 \(D_{10}(u)=0\) 在区间 \((0,1)\) 内恰有两条实根
\(
0<u_-<u_+<1
\)，并且：
\begin{enumerate}
  \item 当 \(0<u<u_-\) 或 \(u_+<u<1\) 时，\(F(\mu,u)\) 具有四条互异实根；
  \item 当 \(u_-<u<u_+\) 时，\(F(\mu,u)\) 具有两条互异实根与一对共轭非实复根；
  \item 当 \(u>1\) 时，\(F(\mu,u)\) 恒具有两条互异实根与一对共轭非实复根。
\end{enumerate}
此外，\(D_{10}(u)\) 在 \((-\infty,0)\cup(1,\infty)\) 无实根，故上述相图不存在其他实临界点。
\end{conclusion}
\begin{proof}[证明要点]
由结论 \ref{con:fold-gauge-anomaly-disc-factorization-simple-branch} 与平方因子自由性可知：除去 \(\{0,1\}\cup\{D_{10}=0\}\) 外，根型在各连通区间内常值。
又由于在 \(u\in(0,1)\) 上有 \(\mathrm{Disc}_{\mu}(F)=u(1-u)D_{10}(u)\)，故判别式符号等价于 \(D_{10}\) 的符号。
选取代表点（例如 \(u=\frac1{20},\frac15,\frac12,\frac65\)）作一次实根计数核验即可确定三类区间的根型；而 \(D_{10}\) 在 \((-\infty,0)\cup(1,\infty)\) 无实根见命题 \ref{prop:fold-gauge-anomaly-branch-point-classification} 的根枚举。
\end{proof}

\begin{conclusion}[速率曲线判别式分解与新的 \texorpdfstring{$19$}{19} 次平方因子]\label{con:fold-gauge-anomaly-rate-curve-discriminant-D19}
将消元证书 \(R(\alpha,u)\) 视为关于 \(\alpha\) 的四次多项式，则其判别式满足
\[
\mathrm{Disc}_{\alpha}\bigl(R(\alpha,u)\bigr)= -u^{11}(u-1)D_{10}(u)\,D_{19}(u)^2,
\]
其中 \(D_{19}(u)\in\ZZ[u]\) 可取为命题 \ref{prop:fold-gauge-anomaly-rate-curve-discriminant-squarefactor} 中的 \(H(u)\)。
并且 \(D_{19}(u)\) 在实轴上恰有 \(7\) 条实根，且全部落于区间 \((0,1)\)；在 \((-\infty,0)\cup(1,\infty)\) 无实根。
因此 \(D_{19}(u)^2\) 仅以纯平方因子形式进入判别式分解：它刻画了消元坐标环的指数型退化位形，而不改变速率曲线覆叠在平方自由意义下的分支支撑 \(\{u=0,1\}\cup\{D_{10}(u)=0\}\)。
\end{conclusion}
\begin{proof}[证明要点]
判别式分解见命题 \ref{prop:fold-gauge-anomaly-rate-curve-discriminant-squarefactor} 并注意 \((u^5D_{19}(u))^2=u^{10}D_{19}(u)^2\)。
实根计数与区间定位可由 Sturm 计数或高精度根枚举复核。
\end{proof}

\begin{conclusion}[速率曲线的泛型单值群同为 \texorpdfstring{$S_4$}{S4}]\label{con:fold-gauge-anomaly-rate-curve-generic-galois-s4}
域扩张 \(\QQ(u)\subset\QQ(u,\alpha)\)（由 \(R(\alpha,u)=0\) 定义）之泛型伽罗瓦群满足
\[
\mathrm{Gal}\bigl(R(\alpha,u)/\QQ(u)\bigr)\cong S_4.
\]
\end{conclusion}
\begin{proof}[证明要点（可核验特化证书）]
取 \(u=2\) 的特化并由结果式消元得到一元四次多项式 \(R(\alpha,2)\in\ZZ[\alpha]\)。
该多项式在模 \(3\) 下不可约且在模 \(11\) 下分解型为 \((3,1)\)，并且判别式非平方；
故其特化伽罗瓦群为 \(S_4\)，从而泛型亦为 \(S_4\)。
\end{proof}

\paragraph{速率消元曲线的正规化、奇点分解与指数理想闭合}\label{par:fold-gauge-anomaly-rate-curve-normalization-singularity-index}
沿用命题 \ref{prop:fold-gauge-anomaly-rate-curve-param}、命题 \ref{prop:fold-gauge-anomaly-rate-curve-discriminant-squarefactor} 与
结论 \ref{con:fold-gauge-anomaly-rate-curve-discriminant-D19} 的记号，记
\[
F(\mu,u)=\mu^{4}-\mu^{3}-\mu^{2}(2u+1)+\mu(-u^{3}+2u)+2u,
\]
\[
\alpha=\frac{u\,\mu_u}{\mu}
\quad\Longleftrightarrow\quad
\alpha\mu\,\partial_{\mu}F(\mu,u)+u\,\partial_uF(\mu,u)=0.
\]
并设
\[
R(\alpha,u):=\frac{1}{u^{2}}\operatorname{Res}_{\mu}\!\Bigl(F(\mu,u),\ \alpha\mu\,\partial_{\mu}F(\mu,u)+u\,\partial_uF(\mu,u)\Bigr)\in\QQ[\alpha,u].
\]

\begin{theorem}[速率消元曲线与谱四次的同函数域性]\label{thm:fold-gauge-anomaly-rate-curve-function-field-identity}
\leanverified{paper\_fold\_gauge\_anomaly\_rate\_curve\_function\_field\_identity}
令 \(\overline C\subset\PP^2\) 为 \(F(\mu,u)\) 齐次化零轨迹的射影闭包，令 \(\Gamma\subset\mathbb A^2_{\alpha,u}\) 为 \(R(\alpha,u)=0\)。
则
\[
\QQ(u,\alpha)=\QQ(u,\mu)=\QQ(\overline C),
\]
从而 \(\overline C\) 与 \(\overline\Gamma\)（\(\Gamma\) 在 \(\PP^1_{\alpha}\times\PP^1_u\) 中的闭包）双有理等价，且 \(\overline C\) 是 \(\overline\Gamma\) 的正规化。
\end{theorem}

\begin{proof}
由定义 \(\alpha\in\QQ(u,\mu)\)，故 \(\QQ(u,\alpha)\subseteq \QQ(u,\mu)\)。
由命题 \ref{prop:fold-gauge-anomaly-spectral-quartic-geometry}，\(\overline C\to\PP^1_u\) 为次数 \(4\) 覆盖，故
\(
[\QQ(u,\mu):\QQ(u)]=4
\)。
另一方面，表 \ref{tab:fold_gauge_anomaly_rate_curve_resultant_degree} 给出 \(\deg_{\alpha}R=4\)；
并且 \(u=2\) 特化时
\[
R(\alpha,2)=-2\bigl(25372\alpha^4-25372\alpha^3-31541\alpha^2+79560\alpha-36800\bigr),
\]
其模 \(3\) 约化为
\(
\alpha^4-\alpha^3+\alpha^2+1
\)
且在 \(\mathbb F_3[\alpha]\) 中不可约。
故 \(R\) 在 \(\QQ(u)[\alpha]\) 中不可约，进而
\(
[\QQ(u,\alpha):\QQ(u)]=4
\)。
两侧扩张次数相同且有包含关系，遂得 \(\QQ(u,\alpha)=\QQ(u,\mu)\)。
后半结论由同函数域曲线的一般事实与 \(\overline C\) 的光滑性（命题 \ref{prop:fold-gauge-anomaly-spectral-quartic-geometry}）给出。
\end{proof}

\begin{corollary}[双次数 \texorpdfstring{$(4,11)$}{(4,11)} 模型的亏格塌缩与 \texorpdfstring{$\delta$}{delta} 缺陷守恒]\label{cor:fold-gauge-anomaly-rate-curve-delta-defect-27}
\leanverified{paper\_fold\_gauge\_anomaly\_rate\_curve\_delta\_defect\_27}
\(\overline\Gamma\subset\PP^1_{\alpha}\times\PP^1_u\) 的双次数为 \((4,11)\)，故其算术亏格
\[
p_a(\overline\Gamma)=(4-1)(11-1)=30.
\]
由定理 \ref{thm:fold-gauge-anomaly-rate-curve-function-field-identity}，\(\overline\Gamma\) 的正规化亏格为 \(g(\overline C)=3\)。
因此
\[
\sum_{P\in\mathrm{Sing}(\overline\Gamma)}\delta_P
=p_a(\overline\Gamma)-g(\overline C)=30-3=27.
\]
\end{corollary}

\begin{proof}
双次数由 \(\deg_{\alpha}R=4,\deg_uR=11\)（表 \ref{tab:fold_gauge_anomaly_rate_curve_resultant_degree}）直接给出。
\(\PP^1\times\PP^1\) 中双次数 \((a,b)\) 曲线满足 \(p_a=(a-1)(b-1)\)。
再用 \(p_a-g=\sum\delta_P\) 即得。
\end{proof}

\begin{theorem}[奇点消元律与 \texorpdfstring{$D_{19}$}{D19} 纤维定位]\label{thm:fold-gauge-anomaly-rate-curve-singular-elimination-u5d19}
在仿射图 \(\Gamma\subset\mathbb A^2_{\alpha,u}\) 上，奇点集由
\[
R(\alpha,u)=\partial_{\alpha}R(\alpha,u)=\partial_uR(\alpha,u)=0
\]
刻画。消去 \(\alpha\) 后得到
\[
u^5D_{19}(u)=0,
\]
其中 \(D_{19}\) 为结论 \ref{con:fold-gauge-anomaly-rate-curve-discriminant-D19} 的十九次多项式。
并且：
\begin{enumerate}
  \item 在 \(u=0\) 纤维上恰有两点
  \[
  (\alpha,u)=(0,0),\qquad (\alpha,u)=\Bigl(\frac12,0\Bigr),
  \]
  因为
  \[
  R(\alpha,0)=10\alpha^2(2\alpha-1)^2.
  \]
  \item 对每个 \(D_{19}(u_0)=0\)（\(u_0\neq0\)），恰有唯一奇点 \((\alpha_0,u_0)\in\Gamma\) 与之对应。
\end{enumerate}
\end{theorem}

\begin{proof}
命题 \ref{prop:fold-gauge-anomaly-rate-curve-discriminant-squarefactor} 给出
\(
\Disc_{\alpha}(R)=-u^{11}(u-1)D_{10}(u)D_{19}(u)^2
\)；
但奇点投影由理想 \(\langle R,R_{\alpha},R_u\rangle\) 的消元控制。
对该理想作 lex 消元可得
\[
\langle R,R_{\alpha},R_u\rangle\cap\QQ[u]=\langle u^5D_{19}(u)\rangle,
\]
其计算证书见
\texttt{scripts/exp\_fold\_gauge\_anomaly\_rate\_curve\_singularity\_audit.py}
与注 \ref{rem:fold-gauge-anomaly-rate-curve-singularity-audit-block}。
第一条由 \(R(\alpha,0)\) 分解直接得到。

对第二条，同一 Gröbner 基同时给出一条线性约束
\[
\alpha u^4-\Phi(u)=0,
\]
故 \(u_0\neq0\) 时
\(
\alpha_0=\Phi(u_0)/u_0^4
\)
唯一确定。再由 \(D_{19}\) 在特征 \(0\) 上不可约（命题 \ref{prop:fold-gauge-anomaly-H-galois-discriminant}）可知其平方自由，故 \(19\) 个根互异并给出 \(19\) 个互异奇点。
\end{proof}
\leanverified{paper\_fold\_gauge\_anomaly\_rate\_curve\_singular\_elimination\_u5d19}

\begin{proposition}[三处高阶奇点的局部型与 \texorpdfstring{$\delta$}{delta} 值]\label{prop:fold-gauge-anomaly-rate-curve-three-special-singularities}
除定理 \ref{thm:fold-gauge-anomaly-rate-curve-singular-elimination-u5d19} 的 \(D_{19}\)-纤维点外，\(\overline\Gamma\) 还存在三处高阶奇点：
\begin{enumerate}
  \item \(P_0=(0,0)\)：存在两条光滑分支
  \[
  \alpha=c_{\pm}u^3+O(u^4),\qquad
  c_{\pm}=\frac{-15\pm 9\sqrt5}{10},
  \]
  且二分支交数为 \(3\)，故 \(\delta_{P_0}=3\)。
  \item \(P_{1/2}=(\frac12,0)\)：令 \(b=\alpha-\frac12\)，则存在 Puiseux 参数
  \[
  u=s^2,\qquad b=\kappa s^5+O(s^6),\qquad \kappa^2=\frac{25}{32},
  \]
  为 \((2,5)\)-尖点，故 \(\delta_{P_{1/2}}=2\)。
  \item \(P_{\infty}=(\alpha,u)=(1,\infty)\)：取 \(v=1/u,\ \beta=\alpha-1\)，则
  \[
  v^{11}R(1+\beta,1/v)
  =-27(\beta+v)\bigl(3\beta^2-3\beta v+v^2\bigr)+\text{(高阶项)},
  \]
  在代数闭包上给出三条互异切线，故为普通三重点，\(\delta_{P_{\infty}}=3\)。
\end{enumerate}
此外，\(D_{19}\)-纤维上的 \(19\) 个奇点均为普通结点。
\end{proposition}

\begin{proof}
上述三处局部初形由
\texttt{scripts/exp\_fold\_gauge\_anomaly\_rate\_curve\_singularity\_audit.py}
给出的可复算展开式直接读出（见注 \ref{rem:fold-gauge-anomaly-rate-curve-singularity-audit-block}）。
于是：
\begin{itemize}
  \item 在 \(P_0\) 处，主导方程 \(5c^2+15c-9=0\) 给出两分支首项系数 \(c_{\pm}\)；
  两分支首差为 \((c_+-c_-)u^3\)，故交数为 \(3\)。
  \item 在 \(P_{1/2}\) 处，主导平衡 \(32k^2-25=0\) 对应 Puiseux 指数 \(5/2\)，即 \((2,5)\)-尖点。
  \item 在 \(P_{\infty}\) 处，最低齐次项分解为三条互异直线，故为普通三重点。
\end{itemize}
对 \(D_{19}\)-纤维点，审计脚本在全部 \(19\) 点上验证了 Hessian 行列式
\(
R_{\alpha\alpha}R_{uu}-R_{\alpha u}^2
\)
非零（最小绝对值 \(>0.3\)），故均为普通二重点，即普通结点。
最后，由推论 \ref{cor:fold-gauge-anomaly-rate-curve-delta-defect-27} 的总缺陷 \(27\) 与
\(
3+2+3=8
\)
可知结点总贡献为 \(19\)，与点数一致。
\end{proof}
\leanverified{paper\_fold\_gauge\_anomaly\_rate\_curve\_three\_special\_singularities}

\begin{theorem}[判别式平方因子的指数理想公式]\label{thm:fold-gauge-anomaly-rate-curve-index-ideal-formula}
\leanverified{paper\_fold\_gauge\_anomaly\_rate\_curve\_index\_ideal\_formula}
令
\[
\mathcal O_{\mu}:=\ZZ[u][\mu]/(F(\mu,u)),\qquad
\mathcal O_{\alpha}:=\ZZ[u][\alpha]/(R(\alpha,u)).
\]
二者均嵌入定理 \ref{thm:fold-gauge-anomaly-rate-curve-function-field-identity} 的同一函数域。
记 \(D_{19}(u)=H(u)\)。则有
\[
\Disc(\mathcal O_{\alpha}/\ZZ[u])
=\Disc(\mathcal O_{\mu}/\ZZ[u])\cdot\bigl(u^5D_{19}(u)\bigr)^2,
\]
其中
\[
\Disc(\mathcal O_{\mu}/\ZZ[u])=-u(u-1)D_{10}(u),
\]
\[
\Disc(\mathcal O_{\alpha}/\ZZ[u])=-u^{11}(u-1)D_{10}(u)D_{19}(u)^2.
\]
因此 \(u^5D_{19}(u)\) 精确刻画了 \(\mathcal O_{\alpha}\) 到其整闭包（同函数域中的正规阶）之间的指数理想。
\end{theorem}

\begin{proof}
两条判别式分解分别来自
命题 \ref{prop:fold-gauge-anomaly-discriminant-factorization}
与命题 \ref{prop:fold-gauge-anomaly-rate-curve-discriminant-squarefactor}。
相除得
\[
\frac{\Disc(\mathcal O_{\alpha}/\ZZ[u])}{\Disc(\mathcal O_{\mu}/\ZZ[u])}
=u^{10}D_{19}(u)^2
=\bigl(u^5D_{19}(u)\bigr)^2.
\]
对同一可分函数域中的两阶，判别式之比等于指数理想的平方，故结论成立。
\end{proof}

\begin{remark}[局部审计公式块]\label{rem:fold-gauge-anomaly-rate-curve-singularity-audit-block}
\input{sections/generated/eq_fold_gauge_anomaly_rate_curve_singularity_audit}
\end{remark}


\begin{conclusion}[\texorpdfstring{$\mathcal C_F$}{CF} 的无穷远有理点与三重覆盖结构]\label{con:fold-gauge-anomaly-trigonal-structure}
曲线 \(\mathcal C_F\subset\PP^2_{[\mu:u:w]}\) 在无穷远直线 \(w=0\) 上恰含两点 \(\QQ\)-有理点
\[
P_{\infty}^{(0)}=[0:1:0],\qquad P_{\infty}^{(1)}=[1:1:0],
\]
且二者皆为光滑点。
因此，从任一 \(P_{\infty}^{(i)}\) 作线性投影均可在 \(\QQ\) 上诱导一个次数 \(3\) 的态射
\(
\varphi_i:\mathcal C_F\to\PP^1
\)，从而 \(\mathcal C_F\) 具有 \(\QQ\)-三重覆盖结构（trigonal structure）。
\end{conclusion}
\begin{proof}[证明要点]
由齐次方程 \(F^{\mathrm{hom}}(\mu,u,w)=0\) 在 \(w=0\) 上化为 \(\mu(\mu^3-u^3)=0\) 可得无穷远点集合；
其中 \(\mu=0\) 与 \(\mu=u\) 给出上述两点的有理性。
又由 \(\partial_{\mu}F^{\mathrm{hom}}|_{w=0}=4\mu^3-u^3\) 在两点处非零知其为光滑点。
线性投影给出的线系在光滑点处可延拓为正则态射，且对平面四次曲线其次数为 \(3\)。
\end{proof}

\paragraph{三重覆盖的对偶判别式与分歧几何}\label{par:fold-gauge-anomaly-trigonal-mu-ramification}
沿用结论 \ref{con:fold-gauge-anomaly-trigonal-structure} 的平面四次曲线 \(\mathcal C_F\subset\PP^2_{[\mu:u:w]}\)，并在仿射片 \(w=1\) 上以
\[
F(\mu,u)=\mu^4-\mu^3-(2u+1)\mu^2+(2u-u^3)\mu+2u=0
\]
记其仿射模型。
固定投影中心 \(P_\infty^{(0)}=[0:1:0]\)。
由穿过 \(P_\infty^{(0)}\) 的直线丛
\(
\{\mu=t w\}_{t\in\PP^1}
\)
在 \(\mathcal C_F\) 上截得线性系 \(|H-P_\infty^{(0)}|\)，并诱导次数 \(3\) 的态射
\[
\varphi_0:\mathcal C_F\to\PP^1_{\mu},
\qquad
(\mu,u)\longmapsto \mu
\quad(w=1\ \text{仿射片上}),
\]
其纤维由把 \(F(\mu,u)\) 视为关于 \(u\) 的三次方程所给出。

\begin{proposition}[\texorpdfstring{$u$}{u}-判别式的完全分解：\texorpdfstring{$\mu$}{mu}-直线上的对偶分歧因子]\label{prop:fold-gauge-anomaly-trigonal-disc-u-factorization}
\leanverified{paper\_fold\_gauge\_anomaly\_trigonal\_disc\_u\_factorization}
将 \(F(\mu,u)\) 视为关于 \(u\) 的三次多项式，则其判别式在 \(\ZZ[\mu]\) 中分解为
\[
\boxed{\ \mathrm{Disc}_{u}\bigl(F(\mu,u)\bigr)
=-\mu(3\mu+2)(\mu^{2}-\mu-1)^{2}\,q_{4}(\mu)\ },
\]
其中
\[
q_{4}(\mu):=9\mu^{4}-6\mu^{3}+4\mu^{2}+8\mu-16\in\ZZ[\mu]
\]
为四次因子。
因此 \((\mu^{2}-\mu-1)\) 仅以纯平方因子出现，而其余不可约因子均为平方自由。
\end{proposition}
\begin{proof}
把 \(F(\mu,u)\) 写为关于 \(u\) 的三次式：
\[
F(\mu,u)=-\mu u^{3}+\bigl(-2\mu^{2}+2\mu+2\bigr)u+\bigl(\mu^{4}-\mu^{3}-\mu^{2}\bigr).
\]
其二次项系数为零；令 \(A(\mu):=\mu^{2}-\mu-1\)，则
\(
-2\mu^{2}+2\mu+2=-2A(\mu)
\)
且
\(
\mu^{4}-\mu^{3}-\mu^{2}=\mu^{2}A(\mu)
\)。
对一般三次 \(au^{3}+bu^{2}+cu+d\) 的判别式公式在 \(b=0\) 情形化简为
\(
\mathrm{Disc}=-4ac^{3}-27a^{2}d^{2}
\)。
代入 \(a=-\mu\)、\(c=-2A(\mu)\)、\(d=\mu^{2}A(\mu)\) 得
\[
\mathrm{Disc}_{u}(F)=-\mu A(\mu)^{2}\bigl(32A(\mu)+27\mu^{5}\bigr).
\]
直接恒等分解
\(
32A(\mu)+27\mu^{5}=(3\mu+2)q_{4}(\mu)
\)
即得所示完全分解。证毕。
\end{proof}

\begin{proposition}[两条有理直线的切触构型与显式有理点]\label{prop:fold-gauge-anomaly-trigonal-rational-tangency}
\leanverified{paper\_fold\_gauge\_anomaly\_trigonal\_rational\_tangency}
\(\mathcal C_F\) 的仿射片 \(w=1\) 上包含如下 \(\QQ\)-有理点：
\[
(0,0),\quad \Bigl(-\frac23,\frac13\Bigr),\quad \Bigl(-\frac23,-\frac23\Bigr),\quad (-1,1),\quad (1,1),\quad (2,1).
\]
并且：
\begin{enumerate}
  \item 直线 \(u=1\) 与 \(\mathcal C_F\) 的交点分解满足
  \[
  F(\mu,1)=(\mu-2)(\mu-1)(\mu+1)^{2},
  \]
  因而 \(u=1\) 在点 \((\mu,u)=(-1,1)\) 处为切线（交重数为 \(2\)）。
  \item 直线 \(3\mu+2=0\)（即 \(\mu=-\tfrac23\)）与 \(\mathcal C_F\) 的交点分解满足
  \[
  F\Bigl(-\frac23,u\Bigr)=\frac{2}{81}\,(3u-1)^{2}(3u+2),
  \]
  因而 \(\mu=-\tfrac23\) 在点 \((-\tfrac23,\tfrac13)\) 处为切线，并给出额外有理交点 \((-\tfrac23,-\tfrac23)\)。
\end{enumerate}
\end{proposition}
\begin{proof}
逐一代入并在 \(\QQ[\mu]\) 或 \(\QQ[u]\) 中因式分解即可得到所示分解式，从而得到相应有理点。
切触陈述由：平面曲线与直线在某点交重数大于 \(1\) 当且仅当该直线为该点切线。证毕。
\end{proof}

\begin{theorem}[投影 \texorpdfstring{$\varphi_0$}{phi0} 的分歧值全集与分歧指数]\label{thm:fold-gauge-anomaly-trigonal-ramification-mu}
\leanverified{paper\_fold\_gauge\_anomaly\_trigonal\_ramification\_mu}
令 \(\varphi_0:\mathcal C_F\to\PP^1_{\mu}\) 为上文由 \(P_\infty^{(0)}\) 投影诱导的次数 \(3\) 态射。
则其有限分歧值集合恰为
\[
\mu=0,\qquad \mu=-\frac23,\qquad \mu^{2}-\mu-1=0,\qquad q_{4}(\mu)=0,
\]
并且 \(\mu=\infty\) 不分歧。
其分歧指数分布如下：
\begin{itemize}
  \item 在 \(\mu^{2}-\mu-1=0\) 的两个分歧值上，\(\varphi_0\) 发生全分歧：纤维中唯一点为 \((\mu,u)=(\mu,0)\)，且分歧指数 \(e=3\)。
  \item 在 \(\mu=-\tfrac23\) 上存在唯一简单分歧点 \((-\tfrac23,\tfrac13)\)，其分歧指数 \(e=2\)。
  \item 在 \(q_4(\mu)=0\) 的四个根上，各存在且仅存在一个简单分歧点，分歧指数均为 \(e=2\)。
  \item 在 \(\mu=0\) 上，分歧点为无穷远点 \(P_\infty^{(0)}\)，且 \(e(P_\infty^{(0)})=2\)；同一纤维中另含仿射点 \((0,0)\) 为非分歧点。
\end{itemize}
\end{theorem}
\begin{proof}
对有限分歧值（除 \(\mu=0\) 外），在仿射片 \(w=1\) 上 \(\varphi_0\) 的纤维由三次方程 \(F(\mu,u)=0\) 给出；
因此 \(\varphi_0\) 在有限参数 \(\mu=\mu_0\neq 0\) 处分歧当且仅当 \(F(\mu_0,u)\) 在 \(u\) 上有重根，即当且仅当 \(\mathrm{Disc}_{u}(F)(\mu_0)=0\)。
由命题 \ref{prop:fold-gauge-anomaly-trigonal-disc-u-factorization}，这等价于
\(
\mu_0\in\{-\tfrac23\}\cup\{\mu^{2}-\mu-1=0\}\cup\{q_4(\mu)=0\}
\)。
\par
当 \(\mu^{2}-\mu-1=0\) 时，\(F(\mu,u)=-\mu u^{3}\)，故纤维唯一且为三重根，分歧指数为 \(3\)。
当 \(\mu=-\tfrac23\) 时由命题 \ref{prop:fold-gauge-anomaly-trigonal-rational-tangency} 得到 \(u=\tfrac13\) 为二重根，故出现唯一简单分歧点且 \(e=2\)。
当 \(q_4(\mu_0)=0\) 且 \(\mu_0\notin\{0,-\tfrac23\}\) 时，\(\mathrm{Disc}_{u}(F)\) 在 \(\mu=\mu_0\) 为一次零点，故 \(F(\mu_0,u)\) 恰含一个二重根与一个单根，从而纤维上恰有一处简单分歧点且 \(e=2\)。
\par
对 \(\mu=0\) 的分歧性需回到射影模型。
在仿射图 \(u=1\) 中取局部坐标 \(m=\mu/u\)、\(t=w/u\)；
将 \(F^{\mathrm{hom}}(\mu,u,w)=0\) 写成
\[
m^{4}-m^{3}t-2m^{2}t-m^{2}t^{2}-m+2mt^{2}+2t^{3}=0.
\]
在点 \(P_\infty^{(0)}\) 处 \((m,t)=(0,0)\)。
该局部方程的线性主项为 \(-m\)，故切线为 \(m=0\)，即 \(\mu=0\)。
并且沿切线 \(m=0\) 有方程化为 \(2t^{3}=0\)，从而切线与曲线在 \(P_\infty^{(0)}\) 处交重数为 \(3\)。
因此在由 \(|H-P_\infty^{(0)}|\) 给出的三重覆盖中，参数 \(\mu=0\) 的纤维包含点 \(P_\infty^{(0)}\) 以重数 \(2\) 出现，即 \(e(P_\infty^{(0)})=2\)。
另一方面，代入 \(\mu=0\) 于仿射方程得 \(2u=0\)，故另有仿射点 \((0,0)\) 且为单点。
\par
对 \(\mu=\infty\)，在图 \(u=1\) 上令基底局部坐标为 \(s=w/\mu=t/m\)（\(\mu\neq 0\) 的无穷远点皆满足此条件）。
在 \(t=0\)（即 \(w=0\)）时局部方程化为 \(m(m^{3}-1)=0\)，给出四个互异无穷远点，且对每一根 \(\partial_{m}(m^{4}-m)=4m^{3}-1\neq 0\)；
因此 \(w=0\) 与 \(\mathcal C_F\) 在上述点处均横截，故 \(s\) 在这些点处零阶为 \(1\)，\(\mu=\infty\) 上方不分歧。
\par
最后，由 \(g(\mathcal C_F)=3\) 与 \(\deg(\varphi_0)=3\) 的 Riemann--Hurwitz 公式知总分歧量为 \(10\)；
上面的分歧指数贡献为 \(2+2+6=10\)，故无其他分歧值。证毕。
\end{proof}

\begin{proposition}[四次分歧因子 \texorpdfstring{$q_4$}{q4} 的伽罗瓦群与判别式]\label{prop:fold-gauge-anomaly-trigonal-q4-galois}
\leanverified{paper\_fold\_gauge\_anomaly\_trigonal\_q4\_galois}
令
\[
q_4(x)=9x^4-6x^3+4x^2+8x-16\in\ZZ[x].
\]
则 \(q_4\) 在 \(\QQ\) 上不可约，且其分裂域的伽罗瓦群同构于 \(S_4\)。
并且其判别式为
\[
\boxed{\ \mathrm{Disc}(q_4)=-2^{12}\cdot 3^{4}\cdot 1931\ },
\]
从而 \(\mathrm{Disc}(q_4)\) 非平方，且分裂域的唯一二次子域为 \(\QQ(\sqrt{-1931})\)。
\end{proposition}
\begin{proof}
由于 \(\mathrm{Disc}(q_4)=-2^{12}\cdot 3^{4}\cdot 1931\)，故 \(13,17\nmid \mathrm{Disc}(q_4)\)，从而可用模素分解型读出 Frobenius 循环型。
在模 \(17\) 下有
\(
q_4(x)\equiv x^4+5x^3+8x^2-x+2
\pmod{17}
\)
且其在 \(\FF_{17}[x]\) 中不可约，从而 \(q_4\) 在 \(\QQ\) 上不可约。
又在模 \(13\) 下分解为
\[
q_4(x)\equiv (x+4)(x-3)(x^2-6x+4)\pmod{13},
\]
其分解型为 \((1,1,2)\)，对应于伽罗瓦群中存在换位。
不可约性给出传递性，而传递子群若同时含 \(4\)-循环与换位则必为 \(S_4\)，故 \(\mathrm{Gal}(q_4/\QQ)\cong S_4\)。
判别式由一元判别式公式直接计算并在 \(\ZZ\) 中分解得到，且其平方自由部分为 \(-1931\)，从而唯一二次子域为 \(\QQ(\sqrt{-1931})\)。证毕。
\end{proof}

\begin{theorem}[三重覆盖的 \texorpdfstring{$S_3$}{S3} 伽罗瓦闭包：二次--三次可解塔与亏格]\label{thm:fold-gauge-anomaly-trigonal-s3-galois-closure}
\leanverified{paper\_fold\_gauge\_anomaly\_trigonal\_s3\_galois\_closure}
定义属 \(2\) 的超椭圆曲线
\[
H_\mu:\quad v^2=-\mu(3\mu+2)q_4(\mu),
\]
并令 \(\widetilde{\mathcal C}\) 为仿射曲线
\[
v^2=-\mu(3\mu+2)q_4(\mu),\qquad F(\mu,u)=0
\]
之射影闭包的正规化。
则：
\begin{enumerate}
  \item 投影 \(\pi:\widetilde{\mathcal C}\to H_\mu\) 为循环三重覆盖（伽罗瓦群同构于 \(C_3\)）；
  \item 复合映射 \(\widetilde{\mathcal C}\xrightarrow{\pi}H_\mu\to\PP^1_{\mu}\) 给出 \(\varphi_0\) 的 \(S_3\)-伽罗瓦闭包；
  \item \(\widetilde{\mathcal C}\) 的亏格为 \(g(\widetilde{\mathcal C})=8\)。
\end{enumerate}
\end{theorem}
\begin{proof}
由命题 \ref{prop:fold-gauge-anomaly-trigonal-disc-u-factorization}，
\[
\mathrm{Disc}_{u}(F)= (\mu^2-\mu-1)^2\cdot v^2
\]
在函数域 \(\QQ(H_\mu)=\QQ(\mu,v)\) 中为平方。
另一方面，特化 \(\mu=3\) 时三次多项式 \(F(3,u)\in\ZZ[u]\) 无有理根，故 \(F(\mu,u)\) 在 \(\QQ(\mu)[u]\) 中不可约，因而 \(\QQ(\mu)\subset\QQ(\mu,u)\) 为次数 \(3\) 扩张。
对不可约三次而言，判别式在基域中为平方当且仅当其伽罗瓦群包含于 \(A_3\simeq C_3\)，从而在 \(\QQ(H_\mu)\) 上该扩张为循环伽罗瓦扩张，得到第 (1) 条。
\par
由 \(H_\mu\to\PP^1_\mu\) 的超椭圆对合给出一个阶 \(2\) 自同构，而上一步的三次循环扩张给出阶 \(3\) 自同构；
两者合成产生半直积 \(C_3\rtimes C_2\simeq S_3\)，并与三次覆盖 \(\varphi_0\) 的分裂域相一致，从而得到第 (2) 条。
\par
第 (3) 条由 Riemann--Hurwitz 公式计算：由定理 \ref{thm:fold-gauge-anomaly-trigonal-ramification-mu}，\(\varphi_0\) 的分歧值中有 \(6\) 个为简单分歧（惯性阶 \(2\)），并有 \(2\) 个为全分歧（惯性阶 \(3\)）。
对次数 \(6\) 的 \(S_3\)-伽罗瓦覆盖，得到
\[
2g(\widetilde{\mathcal C})-2
=6(-2)+6\cdot 6\Bigl(1-\frac12\Bigr)+2\cdot 6\Bigl(1-\frac13\Bigr)
=14,
\]
故 \(g(\widetilde{\mathcal C})=8\)。证毕。
\end{proof}

\begin{proposition}[\texorpdfstring{$\mu$}{mu}-分歧集合的分裂域：\texorpdfstring{$48$}{48} 次直积型伽罗瓦群]\label{prop:fold-gauge-anomaly-trigonal-branch-splitting-field}
\leanverified{paper\_fold\_gauge\_anomaly\_trigonal\_branch\_splitting\_field}
令
\[
B(\mu):=\mu(3\mu+2)(\mu^{2}-\mu-1)\,q_4(\mu)\in\ZZ[\mu]
\]
为命题 \ref{prop:fold-gauge-anomaly-trigonal-disc-u-factorization} 中判别式的平方自由支撑。
记 \(L\) 为 \(B(\mu)\) 在 \(\QQ\) 上的分裂域，则
\[
[L:\QQ]=48,\qquad \mathrm{Gal}(L/\QQ)\cong S_4\times C_2.
\]
\end{proposition}
\begin{proof}
由命题 \ref{prop:fold-gauge-anomaly-trigonal-q4-galois}，\(q_4\) 的分裂域 \(K/\QQ\) 满足 \(\mathrm{Gal}(K/\QQ)\cong S_4\)，并且其唯一二次子域为 \(\QQ(\sqrt{-1931})\)。
另一方面，\(\mu^{2}-\mu-1\) 的根域为 \(\QQ(\sqrt5)\)，为实二次域，故 \(\QQ(\sqrt5)\not\subset K\)。
由于 \(K/\QQ\) 与 \(\QQ(\sqrt5)/\QQ\) 皆为伽罗瓦扩张且交域平凡，二者线性互素，从而复合域 \(K\cdot\QQ(\sqrt5)\) 的次数为 \(24\cdot 2=48\)，且伽罗瓦群为直积 \(S_4\times C_2\)。
线性因子 \(\mu\)、\(3\mu+2\) 不引入新的扩张，故该复合域即为 \(B(\mu)\) 的分裂域。证毕。
\end{proof}

\begin{conjecture}[有理点的完备刻画]\label{conj:fold-gauge-anomaly-trigonal-rational-points}
设 \(\mathcal C_F(\QQ)\) 表示平面四次曲线 \(\mathcal C_F\) 的 \(\QQ\)-有理点集合。
则猜想
\[
\mathcal C_F(\QQ)=
\Bigl\{
P_\infty^{(0)},\ P_\infty^{(1)},\ (0,0),\ \bigl(-\tfrac23,\tfrac13\bigr),\ \bigl(-\tfrac23,-\tfrac23\bigr),\ (-1,1),\ (1,1),\ (2,1)
\Bigr\}.
\]
\end{conjecture}
\begin{remark}[计算证据接口]\label{rem:fold-gauge-anomaly-trigonal-mu-audit}
命题 \ref{prop:fold-gauge-anomaly-trigonal-disc-u-factorization}--\ref{prop:fold-gauge-anomaly-trigonal-branch-splitting-field} 的因式分解、判别式与模素分解型核验，以及对既约分式 \(\mu=p/q\) 满足 \(\max(|p|,q)\le 50\) 的有理特化枚举所见的全部仿射有理点列表，可由脚本 \texttt{scripts/exp\_fold\_gauge\_anomaly\_trigonal\_mu\_audit.py} 复现。
\end{remark}

\paragraph{双无穷远三重结构：参数判别式、闭包亏格与刚性指纹}

\begin{theorem}[第一三重结构的参数判别式与黄金比三重分歧]\label{thm:fold-gauge-anomaly-first-trigonal-structure-golden-ratio}
\leanverified{paper\_fold\_gauge\_anomaly\_first\_trigonal\_structure\_golden\_ratio}
设
\[
\mathcal C_F=\{F^{\mathrm{hom}}(\mu,u,w)=0\}\subset\PP^2_{[\mu:u:w]},
\qquad
P_\infty^{(0)}=[0:1:0].
\]
由过 \(P_\infty^{(0)}\) 的直线族 \(\mu=t\,w\) 诱导三重映射
\[
\phi_0:\mathcal C_F\to\PP^1_t,\qquad t=\mu/w.
\]
在仿射片 \(w=1\) 上，除去中心点交叉后其余三交点由三次方程
\[
g_t(u)=t u^3+(2t^2-2t-2)u+(-t^4+t^3+t^2)=0
\]
刻画，且
\[
\mathrm{Disc}_u(g_t)= -t(3t+2)(t^2-t-1)^2\bigl(9t^4-6t^3+4t^2+8t-16\bigr).
\]
因此分歧参数由
\[
t=0,\quad t=-\frac23,\quad 9t^4-6t^3+4t^2+8t-16=0,\quad t^2-t-1=0
\]
给出；其中 \(t^2-t-1=0\) 的两根为三重分歧（局部分歧指数 \(e=3\)）。
\end{theorem}
\begin{proof}
在仿射片 \(w=1\) 上代入 \(\mu=t\) 得
\[
F(t,u)=-t u^3+(-2t^2+2t+2)u+(t^4-t^3-t^2).
\]
取 \(g_t(u):=-F(t,u)\) 即得到所示三次式。
三次判别式在常数倍下不变（仅乘四次幂因子，符号不变），故
\[
\mathrm{Disc}_u(g_t)=\mathrm{Disc}_u(F(t,u)),
\]
由命题 \ref{prop:fold-gauge-anomaly-trigonal-disc-u-factorization} 把 \(\mu\) 替换为 \(t\) 即得分解式。

当 \(t^2-t-1=0\) 时，
\[
g_t(u)=t u^3
\]
为三重根并合，故对应分歧指数 \(e=3\)。
其余一次因子参数对应一处二重根并合，故为简单分歧。证毕。
\end{proof}

\begin{proposition}[第一三重结构四次分歧因子的伽罗瓦群]\label{prop:fold-gauge-anomaly-first-trigonal-q4-s4}
\leanverified{paper\_fold\_gauge\_anomaly\_first\_trigonal\_q4\_s4}
设
\[
Q(t):=9t^4-6t^3+4t^2+8t-16.
\]
则 \(Q(t)\) 在 \(\QQ\) 上不可约，且
\[
\mathrm{Gal}(Q/\QQ)\cong S_4.
\]
其判别式满足
\[
\mathrm{Disc}(Q)=-2^{12}\cdot 3^4\cdot 1931.
\]
\end{proposition}
\begin{proof}
这是命题 \ref{prop:fold-gauge-anomaly-trigonal-q4-galois} 的同一四次因子重述（变量由 \(\mu\) 改写为 \(t\)），结论逐字同构。证毕。
\end{proof}

\begin{corollary}[第一三重结构的几何单值群与伽罗瓦闭包亏格]\label{cor:fold-gauge-anomaly-first-trigonal-monodromy-genus}
\leanverified{paper\_fold\_gauge\_anomaly\_first\_trigonal\_monodromy\_genus}
映射 \(\phi_0\) 的几何单值群同构于 \(S_3\)；其伽罗瓦闭包为次数 \(6\) 的 \(S_3\)-伽罗瓦覆盖，且闭包曲线亏格为
\[
g=8.
\]
\end{corollary}
\begin{proof}
由定理 \ref{thm:fold-gauge-anomaly-first-trigonal-structure-golden-ratio}，\(\phi_0\) 同时存在三重分歧（给出 \(3\)-循环）与简单分歧（给出换位）；在次数 \(3\) 的连通覆盖中，含换位的传递子群只能是 \(S_3\)。
闭包亏格结论与定理 \ref{thm:fold-gauge-anomaly-trigonal-s3-galois-closure} 一致。
\end{proof}

\begin{theorem}[第二三重结构的参数判别式]\label{thm:fold-gauge-anomaly-second-trigonal-structure-discriminant}
设 \(P_\infty^{(1)}=[1:1:0]\in\mathcal C_F\)。
由过 \(P_\infty^{(1)}\) 的直线族 \(\mu-u=t\,w\) 诱导三重映射
\[
\phi_1:\mathcal C_F\to\PP^1_t,\qquad t=\frac{\mu-u}{w}.
\]
代入 \(\mu=u+t w\) 后，\(F^{\mathrm{hom}}\) 自动含因子 \(w\)；除去该中心因子并在仿射片 \(w=1\) 上，剩余三交点满足
\[
h_t(u)= (t^4-t^3-t^2)+ (4t^3-5t^2+2)u+(6t^2-7t+1)u^2+(3t-3)u^3=0.
\]
其判别式为
\[
\mathrm{Disc}_u(h_t)=-(t-1)\,P_9(t),
\]
其中
\[
P_9(t)=3t^9+3t^8-31t^7-83t^6+100t^5+424t^4-16t^3-436t^2-52t+100.
\]
并且全部分歧参数均为简单分歧参数。
\end{theorem}
\begin{proof}
把 \(\mu=u+t w\) 代入 \(F^{\mathrm{hom}}\) 并展开，得到
\[
F^{\mathrm{hom}}(u+t w,u,w)=w\cdot \widetilde h_t(u,w),
\]
其中 \(\widetilde h_t\) 关于 \(u\) 为三次齐次式。
在 \(w=1\) 上得到所示 \(h_t(u)\)。

直接计算 \(\mathrm{Disc}_u(h_t)\) 得分解
\[
\mathrm{Disc}_u(h_t)=-(t-1)P_9(t).
\]
再计算
\[
\gcd(P_9,P_9')=1,\qquad P_9(1)=12\neq 0,
\]
故 \((t-1)P_9(t)\) 平方自由，从而所有分歧参数均为简单分歧参数。证毕。
\end{proof}
\leanverified{paper\_fold\_gauge\_anomaly\_second\_trigonal\_structure\_discriminant}

\input{subsubsec__fold-gauge-anomaly-second-trigonal-p9-branchfield-arithmetic}

\begin{corollary}[第二三重结构的几何单值群与伽罗瓦闭包亏格]\label{cor:fold-gauge-anomaly-second-trigonal-monodromy-genus}
映射 \(\phi_1\) 的几何单值群同构于 \(S_3\)。
其伽罗瓦闭包是次数 \(6\) 的 \(S_3\)-伽罗瓦覆盖，并满足
\[
g=10.
\]
\end{corollary}
\begin{proof}
由定理 \ref{thm:fold-gauge-anomaly-second-trigonal-structure-discriminant}，\(\phi_1\) 存在简单分歧，故单值群含换位。
覆盖次数为 \(3\) 且连通，故几何单值群为 \(S_3\)。

伽罗瓦闭包次数为 \(6\)。共有 \(10\) 个简单分歧参数（\(t=1\) 与 \(P_9(t)=0\) 的九根），每个在 \(S_3\)-伽罗瓦覆盖中的贡献为
\[
6\left(1-\frac12\right)=3.
\]
由 Riemann--Hurwitz 公式
\[
2g-2=6(-2)+10\cdot 3=18,
\]
从而 \(g=10\)。证毕。
\leanverified{paper\_fold\_gauge\_anomaly\_second\_trigonal\_monodromy\_genus}
\end{proof}

\begin{conjecture}[黄金比三重分歧的结构刚性]\label{conj:fold-gauge-anomaly-golden-ratio-triple-branch-rigidity}
考虑保持以下数据不变的代数形变族：谱四次覆盖的 \(S_4\) 一般性、以及两无穷远点
\[
P_\infty^{(0)}=[0:1:0],\qquad P_\infty^{(1)}=[1:1:0].
\]
则由 \(P_\infty^{(0)}\) 诱导的三重映射在形变后仍保持恰好两个三重分歧参数，并可在适当参数归一化下写成
\[
t^2-t-1=0
\]
的两根；亦即黄金比三重分歧是该谱几何机制的刚性指纹。
\end{conjecture}


\begin{conclusion}[三次分辨曲线的平面六次模型与正规化亏格]\label{con:fold-gauge-anomaly-resolvent-sextic-normalization-genus}
令 \(\mathcal C_3\subset\PP^2_{[y:u:w]}\) 为 \(\Phi(y,u)=0\) 的射影闭包。
则 \(\mathcal C_3\) 为一条平面六次曲线，且无穷远处仅含一点 \([1:0:0]\)；其仿射部分为光滑曲线，而 \([1:0:0]\) 为唯一奇点。
记 \(\widetilde{\mathcal C}_3\) 为其正规化，则
\[
g(\widetilde{\mathcal C}_3)=4,
\]
并且覆叠 \(\widetilde{\mathcal C}_3\to\PP^1_u\) 为次数 \(3\) 的映射，在结论 \ref{con:fold-gauge-anomaly-disc-factorization-simple-branch} 的 \(12\) 个分支参数处均为简单分支，而在 \(u=\infty\) 处不分支。
等价地，\([1:0:0]\) 的 \(\delta\)-不变量为 \(6\)。
\end{conclusion}
\begin{proof}[证明要点]
齐次化可见 \(\mathcal C_3\) 的无穷远点满足 \(u=0=w\)，故唯一为 \([1:0:0]\)。
由于 \(\mathrm{Disc}_y(\Phi)=-u(u-1)D_{10}(u)\) 且平方因子自由，有限分支值均为简单分支，故除无穷远处外曲线无奇点。
进一步，由 \(u\to\infty\) 时的主导平衡 \(y^3-u^6=0\) 得到三条互异分支 \(y\sim \zeta u^2\)（\(\zeta^3=1\)），从而 \(u=\infty\) 上方无分歧。
因此 \(\widetilde{\mathcal C}_3\to\PP^1_u\) 为次数 \(3\) 的覆叠且仅在 \(12\) 个有限点简单分支，
由 Riemann--Hurwitz 公式得 \(g(\widetilde{\mathcal C}_3)=4\)。
平面六次的算术亏格为 \(10\)，故 \(\delta=10-4=6\)。
\end{proof}

\begin{conclusion}[\texorpdfstring{$S_4$}{S4}-Hurwitz 塔：伽罗瓦闭包与中间商曲线的亏格]\label{con:fold-gauge-anomaly-s4-hurwitz-tower-genus}
令 \(\mathcal X\to\PP^1_u\) 表示由谱四次 \(F(\mu,u)=0\) 的 \(S_4\) 伽罗瓦闭包所对应的正规伽罗瓦覆叠（次数 \(24\)，群 \(S_4\)），并以结论 \ref{con:fold-gauge-anomaly-disc-factorization-simple-branch} 的 \(12\) 个简单分支点作为全部分支数据。
则 \(\mathcal X\) 的亏格为
\[
g(\mathcal X)=49.
\]
并且对以下中间子群的商曲线具有完全显式的亏格分层：
\begin{enumerate}
  \item 对 \(A_4\) 的商曲线同构于结论 \ref{con:fold-gauge-anomaly-sign-quadratic-cover} 的判别式双覆盖 \(\mathcal H\)，故亏格为 \(5\)；
  \item 对任一根稳定子 \(S_3\) 的商曲线同构于 \(\mathcal C_F\)，故亏格为 \(3\)；
  \item 对任一稳定配对结构的二面体子群 \(D_8\)（阶 \(8\)）的商曲线同构于结论 \ref{con:fold-gauge-anomaly-resolvent-sextic-normalization-genus} 的正规化 \(\widetilde{\mathcal C}_3\)，故亏格为 \(4\)；
  \item 对正规 Klein 四子群 \(V_4\) 的商覆叠次数为 \(6\)，且每一分支值的惯性在该 \(6\) 叶作用下为 \((2,2,2)\)，从而该商曲线的亏格为 \(13\)。
\end{enumerate}
\end{conclusion}
\begin{proof}[证明要点]
在 \(S_4\) 伽罗瓦覆叠中，每个分支点的惯性群由换位生成，故惯性阶为 \(2\)。
由 Riemann--Hurwitz 公式，
\(
2g(\mathcal X)-2=|S_4|\cdot(-2)+12\cdot|S_4|\,(1-\tfrac12)
\)
立得 \(g(\mathcal X)=49\)。
各中间商的亏格由同一公式对相应的置换表示（等价于对商覆叠的次数与惯性分解型）逐一计算得到；
其中 \(S_4/D_8\) 给出三次预解式的根域覆盖，而 \(S_4/V_4\cong S_3\) 给出其 \(S_3\) 伽罗瓦闭包。
\end{proof}

\begin{remark}[计算审计来源]\label{rem:fold-gauge-anomaly-spectral-curve-arithmetic-audit}
上述伽罗瓦群判定、判别式因子分解与 Tschirnhaus 变换系数可由计算代数数论的标准算法复现；其审计记录见 \cite{Internal2026FoldGaugeAnomalySpectralCurveArithmeticAudit20260218}。
\end{remark}

% Spectral / recurrence / index certificates for the fold gauge anomaly.
% This file is intentionally hand-written (not auto-generated) to keep the
% main text coherent while reusing generated algebraic fragments.

\paragraph{谱—递推—指数闭合：规范差机制的可核验三角}\label{par:fold-gauge-anomaly-spectral-recurrence-triangle}
上一段输入的生成片段已给出规范差逐位过程的压力函数（命题 \ref{prop:fold-gauge-anomaly-pressure}）、Jordan 相关闭式（定理 \ref{thm:fold-gauge-anomaly-autocov-closed}）、
谱指纹与 Hankel 证书（定理 \ref{thm:fold-gauge-anomaly-hankel-rank3}），以及大偏差端点与代数速率曲线（命题 \ref{prop:fold-gauge-anomaly-ldp} 与命题 \ref{prop:fold-gauge-anomaly-rate-curve-param}）。
本段把这些结果进一步压缩为三类\emph{有限数据即可核验}的接口：频域功率谱的有理指纹、有限长度矩生成函数的四阶封闭递推，以及折叠条件期望的有限指数（最大纤维多重度）。

\begin{theorem}[局部重写步数的指数下尾与几乎处处线性下界]\label{thm:fold-local-rewrite-ldp-lower-tail-barrier}
沿用均匀基线 \(\omega\sim\mathrm{Unif}(\Omega_{m+1})\)，令 \(G_m(\omega)\) 为规范差总和（定义 \ref{def:fold-gauge-anomaly}）。
设 \(G_m/m\) 满足大偏差原理，其速率函数记为 \(I\)（命题 \ref{prop:fold-gauge-anomaly-ldp}），并记均值密度极限
\[
g_\ast:=\lim_{m\to\infty}\frac{\EE G_m}{m}=\frac49
\qquad(\text{定理 \ref{thm:fold-gauge-anomaly-density}}).
\]
考虑任意一类局部重写正规化算法 \(\mathsf S\)，并设存在常数 \(R\in\NN\) 使得对任意输入、任意一步，算法最多消除 \(R\) 个单位规范差；记长度 \(m\) 的步数为 \(T_m(\mathsf S)\)。
若 \(G_m\le R\,T_m(\mathsf S)\) 恒成立，则：
\begin{enumerate}
  \item 对任意 \(a<g_\ast/R\)，有指数下尾
  \[
  \limsup_{m\to\infty}\frac{1}{m}\log \PP\!\left(\frac{T_m(\mathsf S)}{m}\le a\right)
  \le -I(Ra)\ <0.
  \]
  \item 线性下界几乎处处成立：
  \[
  \liminf_{m\to\infty}\frac{T_m(\mathsf S)}{m}\ \ge\ \frac{g_\ast}{R}=\frac{4}{9R}
  \qquad\text{几乎必然}.
  \]
\end{enumerate}
\end{theorem}
\begin{proof}
由 \(G_m\le R\,T_m(\mathsf S)\) 得事件包含
\[
\left\{\frac{T_m(\mathsf S)}{m}\le a\right\}
\subseteq
\left\{\frac{G_m}{m}\le Ra\right\}.
\]
对右侧事件应用命题 \ref{prop:fold-gauge-anomaly-ldp} 的大偏差上界即得 (1)。
由 (1) 与 Borel--Cantelli 引理可知：对任意 \(\varepsilon>0\)，事件
\(
\{T_m/m\le g_\ast/R-\varepsilon\}
\)
仅有限次发生，从而得到 (2)。证毕。
\end{proof}
\leanverified{paper\_fold\_local\_rewrite\_ldp\_lower\_tail\_barrier}

\begin{theorem}[Bernoulli\texorpdfstring{$(p)$}{(p)} 基线下的局部重写步数下界与指数下尾]\label{thm:fold-local-rewrite-ldp-lower-tail-barrier-bernoulli-p}
取 \(p\in(0,1)\)，并令输入 \(\omega\in\Omega_{m+1}=\{0,1\}^{m+1}\) 按 IID Bernoulli\((p)\) 采样。
令 \(G_m(\omega)\) 为 Berstel--Zeckendorf 失配加性泛函（定理 \ref{thm:fold-gauge-anomaly-bernoulli-p-laws}，式 \eqref{eq:fold-gauge-anomaly-bernoulli-p-mismatch-sum}），并记其均值密度
\[
g_\ast(p):=\lim_{m\to\infty}\frac{\EE G_m}{m}=\frac{p^2(3-2p)}{1+p^3}
\qquad(\text{定理 \ref{thm:fold-gauge-anomaly-bernoulli-p-laws}}),
\]
以及其大偏差率函数 \(I_p\)（命题 \ref{prop:fold-gauge-anomaly-bernoulli-p-ldp}）。
考虑任意局部重写正规化算法 \(\mathsf S\)，并设存在常数 \(R\in\NN\) 使得对任意输入、任意一步，算法最多消除 \(R\) 个单位规范差；记步数为 \(T_m(\mathsf S)\)。
若 \(G_m\le R\,T_m(\mathsf S)\) 恒成立，则：
\begin{enumerate}
  \item 对任意 \(a<g_\ast(p)/R\)，有指数下尾
  \[
  \limsup_{m\to\infty}\frac{1}{m}\log \PP\!\left(\frac{T_m(\mathsf S)}{m}\le a\right)
  \le -I_p(Ra)\ <0.
  \]
  \item 线性下界在 Bernoulli\((p)\) 基线下几乎处处成立：
  \[
  \liminf_{m\to\infty}\frac{T_m(\mathsf S)}{m}\ \ge\ \frac{g_\ast(p)}{R}
  \qquad\text{几乎必然}.
  \]
\end{enumerate}
\end{theorem}
\begin{proof}
与定理 \ref{thm:fold-local-rewrite-ldp-lower-tail-barrier} 同理，由 \(G_m\le R\,T_m(\mathsf S)\) 得
\[
\left\{\frac{T_m(\mathsf S)}{m}\le a\right\}
\subseteq
\left\{\frac{G_m}{m}\le Ra\right\}.
\]
对右侧事件应用命题 \ref{prop:fold-gauge-anomaly-bernoulli-p-ldp} 的大偏差上界即得 (1)。
再由 Borel--Cantelli 引理推出：对任意 \(\varepsilon>0\)，事件
\(
\{T_m/m\le g_\ast(p)/R-\varepsilon\}
\)
仅有限次发生，从而得到 (2)。证毕。
\end{proof}
\leanverified{paper\_fold\_local\_rewrite\_ldp\_lower\_tail\_barrier\_bernoulli\_p}

\leanverified{paper\_fold\_local\_rewrite\_bernoulli\_p\_optimally\_biased\_lower\_bound}
\begin{corollary}[可调偏置下的最强步数密度下界]\label{cor:fold-local-rewrite-bernoulli-p-optimally-biased-lower-bound}
在定理 \ref{thm:fold-local-rewrite-ldp-lower-tail-barrier-bernoulli-p} 的设定下，若允许在 IID Bernoulli\((p)\) 输入族上调节 \(p\in(0,1)\)，则由推论 \ref{cor:fold-gauge-anomaly-bernoulli-p-density-unimodal} 的唯一极大点 \(p_\star\) 得到最强线性下界常数：
\[
\liminf_{m\to\infty}\frac{T_m(\mathsf S)}{m}\ \ge\ \frac{\max_{p\in(0,1)}g_\ast(p)}{R}
\ =\ \frac{p_\star^2}{R}
\qquad\text{（在选取 }p=p_\star\text{ 的基线下几乎处处）}.
\]
\end{corollary}

\begin{theorem}[\texorpdfstring{$4/9$}{4/9}--\texorpdfstring{$R$}{R} 饱和下界的步效率刚性]\label{thm:fold-local-rewrite-saturation-step-efficiency-rigidity}
设 \(\mathsf S\) 为局部重写正规化算法，并存在 \(R\in\NN\) 使得在任意输入上每一步至多消去 \(R\) 个单位规范差。
对每个 \(m\) 与输入 \(\omega\in\Omega_{m+1}\)，记算法步数为 \(T_m^{\mathsf S}(\omega)\)。
令 \(G_{m}^{(0)}(\omega):=G_m(\omega)\)，并令 \(G_m^{(i)}(\omega)\) 为执行第 \(i\) 步后（\(0\le i\le T_m^{\mathsf S}(\omega)\)）的剩余规范差总量；定义单步净消去量
\[
\Delta_{m,i}(\omega):=G_m^{(i-1)}(\omega)-G_m^{(i)}(\omega)\in\{0,1,\dots,R\}.
\]
则 \(\sum_{i=1}^{T_m^{\mathsf S}(\omega)}\Delta_{m,i}(\omega)=G_m(\omega)\)。
若存在一列输入 \(\omega_m\in\Omega_{m+1}\) 使得
\[
\lim_{m\to\infty}\frac{G_m(\omega_m)}{m}=g_\ast=\frac49,
\qquad
\lim_{m\to\infty}\frac{T_m^{\mathsf S}(\omega_m)}{m}=\frac{g_\ast}{R},
\]
则对任意 \(\varepsilon>0\) 都有
\[
\lim_{m\to\infty}\frac{1}{T_m^{\mathsf S}(\omega_m)}\,
\#\Bigl\{1\le i\le T_m^{\mathsf S}(\omega_m):\ \Delta_{m,i}(\omega_m)\le R-\varepsilon\Bigr\}=0.
\]
\leanverified{paper\_fold\_local\_rewrite\_saturation\_step\_efficiency\_rigidity}
\end{theorem}
\begin{proof}
由 \(\sum_{i=1}^{T_m}\Delta_{m,i}=G_m\) 得步均净消去率
\[
\frac{1}{T_m}\sum_{i=1}^{T_m}\Delta_{m,i}=\frac{G_m}{T_m}\xrightarrow[m\to\infty]{}\frac{g_\ast}{g_\ast/R}=R,
\]
其中 \(T_m:=T_m^{\mathsf S}(\omega_m)\)、\(G_m:=G_m(\omega_m)\)。
又恒有 \(\Delta_{m,i}\le R\)。
若存在 \(\varepsilon>0\) 与 \(\eta>0\) 及子列 \(m_k\) 使得
\[
\frac{1}{T_{m_k}}\#\{i:\ \Delta_{m_k,i}\le R-\varepsilon\}\ge \eta,
\]
则
\[
\frac{1}{T_{m_k}}\sum_{i=1}^{T_{m_k}}\Delta_{m_k,i}
\le (R-\varepsilon)\eta+R(1-\eta)=R-\varepsilon\eta<R,
\]
与步均净消去率极限为 \(R\) 矛盾。
故所示占比极限为 \(0\)。证毕。
\end{proof}

\leanverified{paper\_fold\_local\_rewrite\_saturation\_failure\_quantitative\_cost}
\begin{corollary}[饱和失败的定量代价]\label{cor:fold-local-rewrite-saturation-failure-quantitative-cost}
在定理 \ref{thm:fold-local-rewrite-saturation-step-efficiency-rigidity} 的设定下，
若存在 \(\varepsilon>0\)、\(\eta>0\) 使得对无穷多 \(m\) 有
\[
\frac{1}{T_m^{\mathsf S}(\omega_m)}\#\Bigl\{1\le i\le T_m^{\mathsf S}(\omega_m):\ \Delta_{m,i}(\omega_m)\le R-\varepsilon\Bigr\}\ge \eta,
\]
则必有严格次优的线性下界
\[
\liminf_{m\to\infty}\frac{T_m^{\mathsf S}(\omega_m)}{m}\ \ge\ \frac{g_\ast}{R-\varepsilon\eta}\ >\ \frac{g_\ast}{R}=\frac{4}{9R}.
\]
\end{corollary}
\begin{proof}
由占比条件与 \(\Delta_{m,i}\le R\) 得
\[
G_m(\omega_m)=\sum_{i=1}^{T_m}\Delta_{m,i}
\le (R-\varepsilon)\eta\,T_m+R(1-\eta)\,T_m=(R-\varepsilon\eta)\,T_m,
\]
其中 \(T_m:=T_m^{\mathsf S}(\omega_m)\)。
故 \(T_m\ge G_m(\omega_m)/(R-\varepsilon\eta)\)。
两边除以 \(m\) 并取 \(\liminf\)，再用 \(\lim_{m\to\infty}G_m(\omega_m)/m=g_\ast\) 即得结论。证毕。
\end{proof}

\leanverified{paper\_fold\_local\_rewrite\_reversible\_event\_alphabet\_lower\_bound}
\begin{corollary}[速度饱和与可逆事件字母表的下界]\label{cor:fold-local-rewrite-reversible-event-alphabet-lower-bound}
令 \(\Sigma\) 为有限字母表，\(|\Sigma|=M\ge 2\)。
设存在一族单射提升
\[
\Psi_m:\Omega_m\hookrightarrow X_m\times \Sigma^{L_m},
\qquad
\operatorname{pr}_1\circ\Psi_m=\Fold_m,
\]
其中 \(\Sigma^{L_m}\) 表示长度恰为 \(L_m\) 的事件词集合。
记最大纤维
\[
D_m:=\max_{x\in X_m}\card{\Fold_m^{-1}(x)}.
\]
则对所有 \(m\) 有
\[
L_m\ \ge\ \left\lceil \log_M D_m\right\rceil.
\]
并且由定理 \ref{thm:pom-max-fiber} 的渐近 \(D_m=\Theta(\varphi^{m/2})\) 得
\[
\liminf_{m\to\infty}\frac{L_m}{m}\ \ge\ \frac{\tfrac12\log\varphi}{\log M}.
\]
若进一步满足
\[
\frac{L_m}{m}\ \le\ \frac{4}{9R}+o(1),
\]
则必有
\[
M\ \ge\ \varphi^{\frac{9R}{8}}.
\]
在 \(R=3\) 时，\(\varphi^{27/8}\approx 5.07\)，从而 \(M\ge 6\)。
\end{corollary}
\begin{proof}
由推论 \ref{cor:pom-injectivization-mary-auxlength-exact}（取 \(f=\Fold_m\)）得
\(
L_m\ge \lceil \log_M D_m\rceil
\)。
由定理 \ref{thm:pom-max-fiber} 得 \(\log D_m=\frac{m}{2}\log\varphi+O(1)\)，两边除以 \(m\) 并取 \(\liminf\) 即得第二式。
\par
在附加上界下，取 \(\limsup\) 得 \(\limsup_m L_m/m\le 4/(9R)\)，从而
\[
\frac{\tfrac12\log\varphi}{\log M}\ \le\ \liminf_{m\to\infty}\frac{L_m}{m}
\ \le\ \limsup_{m\to\infty}\frac{L_m}{m}\ \le\ \frac{4}{9R}.
\]
整理得到 \(\log M\ge \frac{9R}{8}\log\varphi\)，即 \(M\ge \varphi^{9R/8}\)。证毕。
\end{proof}

\begin{remark}[变长事件词的计数版本]\label{rem:fold-local-rewrite-variable-length-event-words}
在推论 \ref{cor:fold-local-rewrite-reversible-event-alphabet-lower-bound} 的同一设定下，若将事件词允许为变长集合 \(\Sigma^{\le L_m}\)，即
\[
\Psi_m:\Omega_m\hookrightarrow X_m\times \Sigma^{\le L_m},
\qquad
\operatorname{pr}_1\circ\Psi_m=\Fold_m,
\]
则同一计数法给出
\[
L_m\ \ge\ \left\lceil \log_M D_m\right\rceil-1.
\]
因此
\(
\liminf_{m\to\infty}L_m/m\ge \frac{\tfrac12\log\varphi}{\log M}
\)
仍成立，并且由 \(L_m/m\le 4/(9R)+o(1)\) 所诱导的字母表门槛保持不变。
\end{remark}
\begin{proof}
固定 \(x\in X_m\)。
由于 \(\Psi_m\) 为单射且第一坐标固定为 \(\Fold_m\)，在同一纤维 \(\Fold_m^{-1}(x)\) 上第二坐标必须两两不同，故
\[
d_m(x)\le |\Sigma^{\le L_m}|=\sum_{\ell=0}^{L_m}M^\ell=\frac{M^{L_m+1}-1}{M-1}<M^{L_m+1}.
\]
取 \(x\) 为最大纤维达到者即得 \(D_m\le M^{L_m+1}\)，从而 \(L_m\ge \lceil\log_M D_m\rceil-1\)。
渐近式与门槛结论直接由推论 \ref{cor:fold-local-rewrite-reversible-event-alphabet-lower-bound} 的证明中同一极限比较得到。证毕。
\end{proof}

\begin{theorem}[有限字母表下最大纤维的强逆界]\label{thm:fold-local-rewrite-maxfiber-strong-converse-finite-alphabet}
令 \(\Sigma\) 为有限字母表，\(|\Sigma|=M\ge 2\)。
对每个 \(m\) 给定任意编码 \(\mathrm{enc}_m:\Omega_m\to \Sigma^{\le L_m}\) 与任意译码器
\[
\mathrm{dec}_m:\ X_m\times \Sigma^{\le L_m}\longrightarrow \Omega_m,
\qquad
\Fold_m(\mathrm{dec}_m(x,s))=x
\]
（对所有 \((x,s)\) 满足 \(\mathrm{dec}_m(x,s)\) 定义）。
取任意最大纤维达到者 \(x_m^\star\in X_m\) 满足 \(d_m(x_m^\star)=D_m\)，并在条件分布
\(
\omega\sim \mathrm{Unif}(\Fold_m^{-1}(x_m^\star))
\)
下考虑成功概率。
则对所有 \(m\) 有上界
\[
\PP\!\left(\mathrm{dec}_m(x_m^\star,\mathrm{enc}_m(\omega))=\omega\ \middle|\ X=x_m^\star\right)
\le
\frac{|\Sigma^{\le L_m}|}{D_m}
<
\frac{M^{L_m+1}}{D_m}.
\]
若进一步满足 \(L_m\le \frac{4}{27}m+O(1)\)，则由定理 \ref{thm:pom-max-fiber} 得存在常数 \(C>0\) 使
\[
\PP\!\left(\mathrm{dec}_m(x_m^\star,\mathrm{enc}_m(\omega))=\omega\ \middle|\ X=x_m^\star\right)
\le
\exp\!\left(-\Bigl(\frac12\log\varphi-\frac{4}{27}\log M\Bigr)m+C\right).
\]
特别地，当 \(M=5\) 时指数率常数
\[
c_5:=\frac12\log\varphi-\frac{4}{27}\log 5
\]
严格为正，从而最大纤维上的条件成功率以 \(e^{-c_5m+O(1)}\) 衰减。
\end{theorem}
\leanverified{paper\_fold\_local\_rewrite\_maxfiber\_strong\_converse\_finite\_alphabet}
\begin{proof}
对固定 \(x=x_m^\star\)，把纤维 \(\Fold_m^{-1}(x)\) 按码字 \(s=\mathrm{enc}_m(\omega)\) 分组。
对每个给定的 \((x,s)\)，译码器 \(\mathrm{dec}_m(x,s)\) 至多输出一个候选微态；因此在同一组内至多有一个元素能被正确恢复。
故在该纤维上可被正确恢复的微态总数不超过该纤维上出现的不同码字数，进而不超过 \(|\Sigma^{\le L_m}|\)。
在 \(\mathrm{Unif}(\Fold_m^{-1}(x))\) 下得到
\[
\PP\!\left(\mathrm{dec}_m(x,\mathrm{enc}_m(\omega))=\omega\ \middle|\ X=x\right)
\le \frac{|\Sigma^{\le L_m}|}{d_m(x)}
=\frac{|\Sigma^{\le L_m}|}{D_m}.
\]
又有 \(|\Sigma^{\le L_m}|<M^{L_m+1}\)，得到第一条上界。
\par
若 \(L_m\le \frac{4}{27}m+O(1)\)，则 \(M^{L_m+1}\le \exp((\frac{4}{27}\log M)m+O(1))\)。
另一方面，定理 \ref{thm:pom-max-fiber} 给出 \(\log D_m=\frac{m}{2}\log\varphi+O(1)\)。
代入并整理即得指数上界与常数 \(c_5>0\) 的结论。证毕。
\end{proof}

\begin{proposition}[秩编码给出的可逆事件词提升]\label{prop:fold-local-rewrite-rank-encoding-reversible-lift}
\leanverified{paper\_fold\_local\_rewrite\_rank\_encoding\_reversible\_lift}
令 \(\Sigma\) 为有限字母表，\(|\Sigma|=M\ge 2\)，并令
\(
D_m=\max_{x\in X_m}|\Fold_m^{-1}(x)|
\)。
取
\[
L_m:=\left\lceil \log_M D_m\right\rceil.
\]
则存在一族单射提升
\[
\Psi_m:\Omega_m\hookrightarrow X_m\times \Sigma^{L_m},
\qquad
\operatorname{pr}_1\circ\Psi_m=\Fold_m.
\]
特别地，当 \(M=6\) 时有
\[
\limsup_{m\to\infty}\frac{L_m}{m}\ \le\ \frac{\tfrac12\log\varphi}{\log 6},
\qquad
\frac{\tfrac12\log\varphi}{\log 6}\ <\ \frac{4}{27}.
\]
\end{proposition}
\begin{proof}
对每个 \(m\) 与每个 \(x\in X_m\)，固定纤维 \(\Fold_m^{-1}(x)\) 的一个全序，并令 \(\mathrm{rank}_x:\Fold_m^{-1}(x)\xrightarrow{\sim}\{0,1,\dots,d_m(x)-1\}\) 为对应秩标号。
令 \(\mathrm{Enc}_M:[0,M^{L_m}-1]\to \Sigma^{L_m}\) 为长度 \(L_m\) 的 \(M\)-进制确定性编码（对不足位数作左端补零），它是单射。
定义
\[
\Psi_m(\omega):=\bigl(\Fold_m(\omega),\ \mathrm{Enc}_M(\mathrm{rank}_{\Fold_m(\omega)}(\omega))\bigr).
\]
对固定的 \(x\)，\(\mathrm{rank}_x\) 为双射且 \(d_m(x)\le D_m\le M^{L_m}\)，故纤维内单射成立；从而整体 \(\Psi_m\) 为单射，并满足 \(\operatorname{pr}_1\circ\Psi_m=\Fold_m\)。
\par
当 \(M=6\) 时，定理 \ref{thm:pom-max-fiber} 给出 \(\log D_m=\frac{m}{2}\log\varphi+O(1)\)，代入 \(L_m=\lceil \log_6 D_m\rceil\) 得 \(\limsup L_m/m\le \frac{\tfrac12\log\varphi}{\log 6}\)。
最后由直接常数比较 \(\frac{\tfrac12\log\varphi}{\log 6}<\frac{4}{27}\) 得结论。证毕。
\end{proof}

\begin{corollary}[重写密度强制 prime-register 位长的 \texorpdfstring{$\Omega(m\log m)$}{Omega(m log m)} 下界]\label{cor:fold-prime-register-bitlength-omega-mlogm}
沿用定理 \ref{thm:fold-local-rewrite-ldp-lower-tail-barrier} 的设定，并进一步假设存在一类 prime-register 外置制度满足：
对每个分辨率 \(m\) 与输入 \(\omega\in\Omega_{m+1}\)，其日志支持满足
\[
\#\mathrm{supp}\bigl(r_{\mathsf S}(\omega)\bigr)=T_{m+1}^{\mathsf S}(\omega),
\]
并且支持中的每个坐标对应一个模数 \(q_{m,t}\ge 2\)，且 \(\{q_{m,t}\}_{t\in \mathrm{supp}(r_{\mathsf S}(\omega))}\) 两两互素。
令总模数乘积
\[
P_m(\omega):=\prod_{t\in \mathrm{supp}(r_{\mathsf S}(\omega))} q_{m,t}.
\]
则几乎处处成立
\[
\boxed{\;
\liminf_{m\to\infty}\frac{\log_2 P_m(\omega)}{m\log_2 m}\ \ge\ \frac{4}{9R}\; }.
\]
特别地在 \(R=3\) 时，
\(
\liminf_{m\to\infty}\frac{\log_2 P_m(\omega)}{m\log_2 m}\ge \frac{4}{27}
\)。
\end{corollary}
\begin{proof}
记 \(t_m(\omega):=\#\mathrm{supp}(r_{\mathsf S}(\omega))=T_{m+1}^{\mathsf S}(\omega)\)。
由两两互素性，每个 \(q_{m,t}\) 至少贡献一个互异素因子，故
\[
P_m(\omega)\ \ge\ \prod_{j=1}^{t_m(\omega)} p_j,
\]
其中 \(p_j\) 为第 \(j\) 个素数（右端为 primorial）。
由素数定理的等价形式，
\(
\log\prod_{j=1}^{t}p_j=(1+o(1))\,t\log t
\)，
从而
\[
\log_2 P_m(\omega)
\ge (1+o(1))\,t_m(\omega)\,\log_2 t_m(\omega).
\]
另一方面，由定理 \ref{thm:fold-local-rewrite-ldp-lower-tail-barrier}(2) 得几乎处处
\(
\liminf_{m\to\infty}\frac{t_m(\omega)}{m}\ge \frac{4}{9R}
\)。
并且当 \(t_m(\omega)\asymp m\) 时有 \(\log_2 t_m(\omega)/\log_2 m\to 1\)。
将上述不等式除以 \(m\log_2 m\) 并取 \(\liminf\)，得到所示结论。证毕。
\end{proof}
\leanverified{paper\_fold\_prime\_register\_bitlength\_omega\_mlogm}

\begin{remark}[平均缺口驱动的字母表下界]\label{rem:fold-local-rewrite-average-gap-alphabet-bound}
由恒等式 \(\sum_{x\in X_m}d_m(x)=2^m\) 与 \(d_m(x)\le D_m\) 得
\[
D_m\ \ge\ \frac{2^m}{|X_m|}=\frac{2^m}{F_{m+2}}.
\]
因此若 \(\Psi_m\) 如推论 \ref{cor:fold-local-rewrite-reversible-event-alphabet-lower-bound}，则亦有
\[
\liminf_{m\to\infty}\frac{L_m}{m}\ \ge\ \frac{\log(2/\varphi)}{\log M}.
\]
另一方面，若在某一基线分布下满足步数密度上界
\[
\frac{L_m}{m}\ \le\ \frac{g}{R}+o(1)
\]
（例如：每步输出一符号且总输出长度不超过渐近最优步数密度，其中 \(g\) 为对应基线下的 \(\lim_{m\to\infty}\EE[G_m]/m\)），
则联立上述平均缺口下界给出字母表门槛
\[
\boxed{\;
\frac{\log(2/\varphi)}{\log M}\ \le\ \frac{g}{R}
\qquad\Longleftrightarrow\qquad
M\ \ge\ \Bigl(\frac{2}{\varphi}\Bigr)^{R/g}\; }.
\]
在 IID Bernoulli\((p)\) 基线下，\(g=g_\ast(p)\)（定理 \ref{thm:fold-gauge-anomaly-bernoulli-p-laws}），从而
\[
M\ \ge\ \Bigl(\frac{2}{\varphi}\Bigr)^{R/g_\ast(p)}.
\]
并且由推论 \ref{cor:fold-gauge-anomaly-bernoulli-p-density-unimodal}，当 \(p=p_\star\)（即 \(g_\ast(p_\star)=\max_{p\in(0,1)}g_\ast(p)=p_\star^2\)）时门槛最小：
\[
M\ \ge\ \Bigl(\frac{2}{\varphi}\Bigr)^{R/p_\star^2}.
\]
例如在 \(R=3\) 时有 \(\bigl(\frac{2}{\varphi}\bigr)^{3/p_\star^2}\approx 2.915\)，故任何满足上述“渐近步数最优 + 一步一符号 + 单射提升”的方案都强制 \(M\ge 3\)；
而在 \(p=\tfrac12\)（\(g_\ast(1/2)=4/9\)）时得到 \(M\ge 5\)。
\end{remark}

\begin{theorem}[均匀基线下的 \texorpdfstring{$B$}{B}-比特反演上界与阈值分离]\label{thm:fold-bbit-inversion-avg-worst-gap}
固定分辨率 \(m\ge 1\)，令 \(\omega\sim \mathrm{Unif}(\Omega_m)\)，并记 \(X=\Fold_m(\omega)\)。
给定任意 \(B\ge 0\)，令编码 \(b:\Omega_m\to\{0,1\}^B\) 与译码器
\[
\widehat\omega:\ X_m\times\{0,1\}^B\longrightarrow \Omega_m
\]
满足 \(\Fold_m(\widehat\omega(x,u))=x\)（对所有 \((x,u)\) 使 \(\widehat\omega(x,u)\) 定义）。
则成功概率满足粗界
\[
\PP(\widehat\omega(X,b(\omega))=\omega)\ \le\ \frac{|X_m|\cdot 2^B}{|\Omega_m|}
=\frac{F_{m+2}\cdot 2^B}{2^m}.
\]
因此若 \(B\le (m-\log_2|X_m|)-t\)，则 \(\PP(\widehat\omega=\omega)\le 2^{-t+O(1)}\)，从而要使成功率不趋于 \(0\) 必须满足
\[
\frac{B}{m}\ \ge\ \log_2\!\left(\frac{2}{\varphi}\right)+o(1).
\]
另一方面，若要求存在零误差可逆提升 \(\Omega_m\hookrightarrow X_m\times\{0,1\}^B\) 且第一坐标等于 \(\Fold_m\)，则必有
\[
B\ \ge\ \left\lceil\log_2 D_m\right\rceil
=\Bigl(\tfrac12\log_2\varphi+o(1)\Bigr)m.
\]
两条阈值的差为严格正的常数：
\[
\frac12\log_2\varphi-\log_2\!\left(\frac{2}{\varphi}\right)
=\log_2\!\left(\frac{\varphi^{3/2}}{2}\right)
\approx 0.04136,
\]
与注 \ref{rem:pom-defect-gap} 的“最坏/平均偏离基数”一致。
\end{theorem}
\leanverified{paper\_fold\_local\_rewrite\_bbit\_inversion\_avg\_worst\_gap}
\leanverified{paper\_fold\_bbit\_inversion\_avg\_worst\_gap}
\begin{proof}
对每个输入对 \((x,u)\in X_m\times\{0,1\}^B\)，译码器至多输出一个微态。
因此可被正确恢复的微态总数不超过可见输入对的总数 \(|X_m|\cdot 2^B\)，从而在 \(\Omega_m\) 上均匀抽样得到
\[
\PP(\widehat\omega=\omega)\le \frac{|X_m|\cdot 2^B}{|\Omega_m|}.
\]
代入 \(|\Omega_m|=2^m\)、\(|X_m|=F_{m+2}\) 得第一式；其余不等式由直接移项得到。
再由 \(F_{m+2}\sim \varphi^m/\sqrt5\) 得 \(\log_2|X_m|=m\log_2\varphi+O(1)\)，从而 \(m-\log_2|X_m|=m\log_2(2/\varphi)+O(1)\)，给出所示必要下界。
\par
对零误差要求，取最大纤维达到者 \(x^\star\) 满足 \(d_m(x^\star)=D_m\)。
在纤维 \(\Fold_m^{-1}(x^\star)\) 上实现单射化需为每个微态分配互异的 \(B\)-比特标签，故 \(2^B\ge D_m\)，即 \(B\ge \lceil\log_2 D_m\rceil\)。
由定理 \ref{thm:pom-max-fiber} 得 \(\log_2 D_m=\frac{m}{2}\log_2\varphi+O(1)\)，从而得到最坏情形阈值。
最后常数差值恒等式为直接代数化简。证毕。
\end{proof}

\begin{conjecture}[三局部约束下的最小可逆事件字母表]\label{conj:fold-local-rewrite-minimal-reversible-event-alphabet}
在 \(R=3\) 的局部约束下，存在一类可逆局部重写策略，使得其事件日志以每步一符号的方式外置，并且其步数密度达到 \(\frac{4}{27}\)。
并且在同一类协议约束下，\(\frac{4}{27}\) 为可达到的最优线性密度常数；同时事件字母表的最小基数为 \(6\)。
\end{conjecture}

\endinput


\leanverified{paper\_fold\_gauge\_anomaly\_psd\_rational}
\begin{theorem}[Jordan 频谱指纹：逐位规范差指示过程的功率谱为有理函数]\label{thm:fold-gauge-anomaly-psd-rational}
在命题 \ref{prop:fold-gauge-anomaly-pressure} 的均匀基线与平稳口径下，
令 \(g_t=\mathbf 1\{X_t\neq Y_t\}\in\{0,1\}\)，并记
\(
C_k:=\mathrm{Cov}(g_t,g_{t+k})
\)（\(k\in\ZZ\)）。
定义其离散功率谱（Fourier--Stieltjes 变换）
\[
\mathcal S_g(\omega):=\sum_{k\in\ZZ} C_k e^{-ik\omega}\qquad(\omega\in\RR).
\]
则 \(\mathcal S_g\) 为严格有理函数；写 \(s:=e^{i\omega}\) 时有完全闭式
\begin{equation}\label{eq:fold-gauge-anomaly-psd-rational}
\boxed{\;
\mathcal S_g(\omega)=\frac{2}{9}\cdot
\frac{4s^{6}+3s^{5}-46s^{4}-99s^{3}-46s^{2}+3s+4}{8s^{6}+20s^{5}-26s^{4}-85s^{3}-26s^{2}+20s+8}\; }.
\end{equation}
\end{theorem}

\begin{proof}[证明要点（由相关生成函数直接求和）]
由推论 \ref{cor:fold-gauge-anomaly-cov-genfun}，滞后协方差在 \(k\ge 1\) 的生成函数
\(
C(z)=\sum_{k\ge 1}C_k z^k
\)
为有理函数。
又因 \(C_{-k}=C_k\) 且 \(C_0=\mathrm{Var}(g_t)=20/81\)（定理 \ref{thm:fold-gauge-anomaly-autocov-closed}），
故
\(
\mathcal S_g(\omega)=C_0+C(e^{i\omega})+C(e^{-i\omega})
\)。
将 \(z=s\) 与 \(z=s^{-1}\) 代入并作代数化简，即得 \eqref{eq:fold-gauge-anomaly-psd-rational}。证毕。
\end{proof}

\begin{corollary}[分母的完全因子分解与协方差闭式]\label{cor:fold-gauge-anomaly-psd-denominator-factorization-cov}
\leanverified{paper\_fold\_gauge\_anomaly\_psd\_denominator\_factorization\_cov}
在定理 \ref{thm:fold-gauge-anomaly-psd-rational} 的记号下，功率谱 \eqref{eq:fold-gauge-anomaly-psd-rational} 的分母多项式满足完全分解
\[
8s^6+20s^5-26s^4-85s^3-26s^2+20s+8
=(s-2)(s+2)^2(2s-1)(2s+1)^2.
\]
因此 \(\mathcal S_g\) 在单位圆内的极点仅为 \(s=\tfrac12\) 的一重极点与 \(s=-\tfrac12\) 的二重极点。
由留数抽取可得协方差闭式：对一切 \(k\ge 1\)，
\[
C_k=\mathrm{Cov}(g_t,g_{t+k})
=\frac18\Big(\frac12\Big)^k+\Big(\frac{7}{648}+\frac{k}{108}\Big)\Big(-\frac12\Big)^k,
\qquad
C_0=\frac{20}{81}.
\]
\end{corollary}
\begin{proof}
对分母多项式直接代数展开验证分解式。
由 \(|s|=1\) 的 Fourier 系数抽取公式，
\(
C_k=\frac{1}{2\pi i}\oint_{|s|=1}\mathcal S_g(\omega)\,s^{k-1}\,\dd s
\)，
而分解式表明单位圆内仅有 \(s=\tfrac12\) 与 \(s=-\tfrac12\) 两类极点，其中后者为二重极点。
应用留数定理并对二重极点使用导数留数公式即可得到所示闭式；
其结果与定理 \ref{thm:fold-gauge-anomaly-autocov-closed} 的协方差闭式一致。证毕。
\end{proof}

\begin{proposition}[协方差序列的三阶整系数递推与 Hankel 秩证书]\label{prop:fold-gauge-anomaly-covariance-recurrence-integers-hankel-rank3}
\leanverified{paper\_fold\_gauge\_anomaly\_covariance\_recurrence\_integers\_hankel\_rank3}
令 \(C_k=\mathrm{Cov}(g_t,g_{t+k})\)（\(k\ge 0\)）。
则对一切 \(k\ge 1\) 成立三阶整系数递推
\[
8C_{k+3}+4C_{k+2}-2C_{k+1}-C_k=0.
\]
从而无限 Hankel 矩阵 \(H=(C_{i+j+1})_{i,j\ge 0}\) 的秩恰为 \(3\)；
等价地，任意 \(4\times 4\) Hankel 子式行列式恒为零，而存在某个 \(3\times 3\) Hankel 子式非零。
\end{proposition}
\begin{proof}
由定理 \ref{thm:fold-gauge-anomaly-hankel-rank3}，令 \(d_k:=\mathrm{Cov}(g_t,g_{t+k})\)（\(k\ge 1\)），则
\(
d_{k+3}=-\frac12d_{k+2}+\frac14d_{k+1}+\frac18d_k
\)。
对该递推两边同乘 \(8\) 即得所述整系数形式（并将 \(d_k\) 记回 \(C_k\)）。
Hankel 秩结论与定理 \ref{thm:fold-gauge-anomaly-hankel-rank3} 等价。证毕。
\end{proof}

\begin{proposition}[总规范差的四阶递推：概率生成函数由固定特征多项式封闭]\label{prop:fold-gauge-anomaly-mgf-order4-recurrence}
令 \(G_m:=\sum_{t=1}^{m} g_t\)，并对 \(u>0\) 定义概率生成函数
\[
M_m(u):=\mathbb E\!\left[u^{G_m}\right]
=\mathbb E\exp\!\bigl((\log u)\,G_m\bigr).
\]
引入清分母的计数型生成函数 \(\widetilde M_m(u):=2^{m}M_m(u)\)。
则对所有 \(m\ge 0\)，序列 \(\bigl(\widetilde M_m(u)\bigr)_{m\ge 0}\) 满足四阶线性递推
\begin{equation}\label{eq:fold-gauge-anomaly-Mm-order4}
\boxed{\;
\widetilde M_{m+4}(u)=\widetilde M_{m+3}(u)+(2u+1)\widetilde M_{m+2}(u)+(u^{3}-2u)\widetilde M_{m+1}(u)-2u\,\widetilde M_m(u)\; }.
\end{equation}
等价地，\(\bigl(M_m(u)\bigr)\) 同样满足四阶递推，但系数整体按 \(2^{-j}\) 比例缩放（对应于 \(A_\theta=\tfrac12(2A_\theta)\) 的尺度差）。
在 \(u=1\) 处，该递推的特征因子为
\(
(\mu-2)(\mu-1)(\mu+1)^2
\)，
其中 \((\mu+1)^2\) 的 Jordan 二重根对应于定理 \ref{thm:fold-gauge-anomaly-A0-jordan} 的交错相位残差，
从而解释有限长统计量中出现的 \(m2^{-m}\) 级振荡修正项（定理 \ref{thm:fold-gauge-anomaly-var-finite-closed}）。
\end{proposition}
\leanverified{paper\_fold\_gauge\_anomaly\_mgf\_order4\_recurrence}

\begin{proof}[证明要点（Cayley--Hamilton）]
令 \(\theta=\log u\)，并取命题 \ref{prop:fold-gauge-anomaly-pressure} 的矩阵 \(A_\theta\)。
记 \(B_u:=2A_\theta\)，则其特征多项式正是
\(
F(\mu,u)=\det(\mu I-B_u)
\)，
其显式四次多项式写法已在命题 \ref{prop:fold-gauge-anomaly-pressure} 中给出。
由 Cayley--Hamilton 定理得 \(F(B_u,u)=0\)。
而对固定 \(u\)，计数型生成函数 \(\widetilde M_m(u)\) 可写为有限状态换能器的路径权重和，等价地为某个线性泛函作用在 \(B_u^m\) 上；
将 \(F(B_u,u)=0\) 左右乘以相应向量并取该线性泛函，得到 \eqref{eq:fold-gauge-anomaly-Mm-order4}。证毕。
\end{proof}

\begin{theorem}[生成函数的显式有理性与整系数递推]\label{thm:fold-gauge-anomaly-gf-rational-integral-recurrence}
在命题 \ref{prop:fold-gauge-anomaly-parry-kernel-mismatch-label} 的 Parry 转移核 \(P\) 与边标记 \(g\) 下，
令
\[
D_m(u):=\EE\!\left[u^{G_m}\right]=\sum_{k=0}^{m}p_{m,k}u^k,\qquad
p_{m,k}:=\PP(G_m=k),
\]
并定义步长母函数
\[
\mathcal F(z,u):=\sum_{m\ge 0} D_m(u)\,z^m.
\]
则 \(\mathcal F\) 为显式有理函数，且满足
\[
\mathcal F(z,u)=\frac{N(z,u)}{9\,\Delta(z,u)},
\]
其中
\[
\Delta(z,u)=8-4z-(4u+2)z^2+(2u-u^3)z^3+u z^4,
\]
\[
\begin{aligned}
N(z,u)=\;&72+4z-12z^2-2z^3+32u z-24u z^2-10u z^3\\
&\quad+18u^2 z^2+4u^2 z^3-u^3 z^3.
\end{aligned}
\]
等价地，对一切 \(m\ge 4\) 有四阶递推
\[
8D_m(u)=4D_{m-1}(u)+(4u+2)D_{m-2}(u)+(u^3-2u)D_{m-3}(u)-uD_{m-4}(u).
\]
更强地，令
\[
H_m(u):=9\cdot 2^m D_m(u)=9\,\widetilde M_m(u),
\]
则对所有 \(m\ge 0\) 有 \(H_m(u)\in\ZZ[u]\)，并且对一切 \(m\ge 4\) 有整系数递推
\[
H_m(u)=H_{m-1}(u)+(2u+1)H_{m-2}(u)+(u^3-2u)H_{m-3}(u)-2u\,H_{m-4}(u).
\]
因此对每个 \(m,k\) 皆有 \(9\cdot 2^m\,p_{m,k}\in\ZZ\)。
\end{theorem}

\begin{proof}
由 Markov 结构与倾斜核 \(Q(u)_{ij}:=P_{ij}u^{g(i,j)}\)，有
\(
D_m(u)=\pi^\top Q(u)^m\mathbf 1
\)，其中 \(\pi\) 为 \(P\) 的平稳分布。
因此
\(
\mathcal F(z,u)=\pi^\top (I-zQ(u))^{-1}\mathbf 1
\)
为有理函数。
对本 \(4\) 态模型直接展开 \(\det(I-zQ(u))\) 得到 \(\Delta(z,u)\)，并由伴随矩阵条目计算得到 \(N(z,u)\)。
将恒等式 \(9\Delta(z,u)\mathcal F(z,u)=N(z,u)\) 逐项比较 \(z^m\) 系数得到 \(D_m\) 的四阶递推；该递推与命题 \ref{prop:fold-gauge-anomaly-mgf-order4-recurrence} 对 \(\widetilde M_m\) 的递推在缩放下等价。
将递推乘以 \(9\cdot 2^m\) 并整理得到 \(H_m\) 的整系数递推。
由 \(H_0,\dots,H_3\) 的有限枚举可知其属于 \(\ZZ[u]\)，递推闭合推出 \(H_m\in\ZZ[u]\) 对所有 \(m\) 成立，从而 \(9\cdot 2^m p_{m,k}\in\ZZ\)。证毕。
\end{proof}
\leanverified{paper\_fold\_gauge\_anomaly\_gf\_rational\_integral\_recurrence}

\begin{theorem}[端点事件的精确闭式：Fibonacci 端点与 \texorpdfstring{$\bmod 3$}{mod 3} 律]\label{thm:fold-gauge-anomaly-endpoint-fibonacci-mod3}
沿用定理 \ref{thm:fold-gauge-anomaly-gf-rational-integral-recurrence} 的记号。
令 Fibonacci 数列满足 \(F_0=0,F_1=1,F_{n+2}=F_{n+1}+F_n\)。
则端点事件概率满足
\[
\PP(G_m=0)=
\begin{cases}
1,& m=0,\\[2pt]
\frac59,& m=1,\\[6pt]
\displaystyle \frac{F_{m+5}}{9\cdot 2^m},& m\ge 2,
\end{cases}
\]
并且对任意 \(m\ge 1\)，满长端点满足
\[
\PP(G_m=m)=
\begin{cases}
2^{-m},& m\equiv 2\pmod 3,\\[4pt]
\frac89\,2^{-m},& m\equiv 0,1\pmod 3.
\end{cases}
\]
\leanverified{fib\_seven}
\leanverified{fib\_eight\_fold}
\leanverified{fib\_nine\_eq\_fib\_eight\_add\_fib\_seven}
\leanverified{abstract\_endpoint\_fibonacci}
\leanverified{paper\_fold\_gauge\_anomaly\_endpoint\_fibonacci\_mod3}
\end{theorem}

\begin{proof}
由定义 \(\PP(G_m=0)=D_m(0)\)。
将定理 \ref{thm:fold-gauge-anomaly-gf-rational-integral-recurrence} 的四阶递推代入 \(u=0\) 并化简可得：对 \(m\ge 4\)，序列
\(
U_m:=9\cdot 2^m D_m(0)
\)
满足二阶递推 \(U_m=U_{m-1}+U_{m-2}\)。
直接计算 \(U_2=13\)、\(U_3=21\) 得 \(U_m=F_{m+5}\)（\(m\ge 2\)），从而得到 \(\PP(G_m=0)=F_{m+5}/(9\cdot 2^m)\)。
\(m=0\) 时 \(G_0=0\) 给出 \(\PP(G_0=0)=1\)；\(m=1\) 时由 \(\EE[g_t]=4/9\) 得 \(\PP(G_1=0)=1-\EE[g_1]=5/9\)。
\par
对 \(\PP(G_m=m)\)，由命题 \ref{prop:fold-gauge-anomaly-parry-kernel-mismatch-label} 知 \(g=1\) 的允许边仅为
\((a,b)\)、\((b,c)\)、\((c,a)\)，故事件 \(\{G_m=m\}\) 强制隐藏链沿 \(a\to b\to c\to a\) 的 \(3\)-循环确定性行走。
该循环上一步概率依次为 \(P_{ab}=\frac14\)、\(P_{bc}=1\)、\(P_{ca}=\frac12\)，每三步乘积为 \(2^{-3}\)。
按起点 \(S_0\in\{a,b,c\}\) 与余数 \(m\bmod 3\) 分情形计算满循环乘积，并以平稳分布
\(\pi=(\pi_a,\pi_b,\pi_c,\pi_d)=(\frac49,\frac29,\frac29,\frac19)\) 加权求和，即得所示 \(\bmod 3\) 分段公式。证毕。
\end{proof}

\begin{theorem}[单位圆扭结的 Perron--Frobenius 刚性与严格谱隙]\label{thm:fold-gauge-anomaly-unit-circle-rigidity}
令
\[
B(u):=2A_{\log u}=
\begin{pmatrix}
1&u&0&1\\
0&0&u&0\\
u&2&0&0\\
1&0&0&0
\end{pmatrix},
\]
其中 \(A_\theta\) 由命题 \ref{prop:fold-gauge-anomaly-pressure} 给出，\(\log u\) 取任意分支使 \(e^{\log u}=u\)。
则对任意 \(|u|=1\) 有
\[
\rho(B(u))\le \rho(B(1))=2,
\]
且当 \(u\neq 1\) 时严格成立 \(\rho(B(u))<2\)。
等价地，对任意 \(t\in\RR\setminus 2\pi\ZZ\)，有
\[
\rho(A_{it})<1.
\]
\leanverified{paper\_fold\_gauge\_anomaly\_unit\_circle\_rigidity}
\end{theorem}

\begin{proof}
当 \(|u|=1\) 时逐项取绝对值得 \(|B(u)|=B(1)\)（按元素绝对值）。
矩阵 \(B(1)\) 非负且原始：其存在自环 \(1\to 1\)，并且从任一状态都可达状态 \(1\) 且可由状态 \(1\) 到达，故强连通且非周期。
由 Frobenius--Wielandt 比较定理，
\[
\rho(B(u))\le \rho(B(1))=2.
\]
若等号成立，则存在 \(\varphi\in\RR\) 与对角酉矩阵 \(D=\mathrm{diag}(d_1,d_2,d_3,d_4)\) 使
\[
B(u)=e^{i\varphi}D\,B(1)\,D^{-1}.
\]
比较 \((1,1)\) 元得 \(e^{i\varphi}=1\)，从而 \(B(u)=D B(1)D^{-1}\)。
对每个 \(B(1)_{ij}\neq 0\) 有
\[
\frac{d_j}{d_i}=\frac{B(u)_{ij}}{B(1)_{ij}}.
\]
由 \((1,2)\)、\((2,3)\)、\((3,1)\) 元分别得
\[
\frac{d_2}{d_1}=u,\qquad \frac{d_3}{d_2}=u,\qquad \frac{d_1}{d_3}=u,
\]
相乘得到 \(u^3=1\)。
另一方面，\((3,2)\) 元给出 \(d_2/d_3=1\)，与 \(d_3/d_2=u\) 合并即 \(u=1\)。
故 \(u\neq 1\) 时不可能取等，必有 \(\rho(B(u))<2\)。
最后令 \(u=e^{it}\)，则 \(B(u)=2A_{it}\)，于是
\[
\rho(A_{it})=\frac{\rho(B(e^{it}))}{2}<1\qquad (t\notin 2\pi\ZZ).
\]
证毕。
\end{proof}

\begin{corollary}[规范差同余分布的指数均匀化]\label{cor:fold-gauge-anomaly-congruence-mixing}
记
\[
M_m(u):=\mathbb E\!\left[u^{G_m}\right],\qquad \widetilde M_m(u):=2^m M_m(u).
\]
则 \(\widetilde M_m\) 满足四阶递推
\[
\widetilde M_{m+4}
=\widetilde M_{m+3}
+(2u+1)\widetilde M_{m+2}
+(u^3-2u)\widetilde M_{m+1}
-2u\,\widetilde M_m,
\]
其特征多项式即
\[
F(\mu,u)=\mu^4-\mu^3-\mu^2(2u+1)+\mu(-u^3+2u)+2u.
\]
固定整数 \(q\ge 2\)，令 \(\omega_q=e^{2\pi i/q}\)，并记
\[
R_{m,q}:=G_m\bmod q\in \ZZ/q\ZZ.
\]
定义
\[
\delta_q:=\max_{1\le j\le q-1}\frac{\rho(B(\omega_q^{\,j}))}{2}\in(0,1).
\]
则存在常数 \(C_q>0\)，使对任意 \(m\ge 1\) 与 \(r\in\ZZ/q\ZZ\) 都有
\[
\left|\mathbb P(R_{m,q}=r)-\frac1q\right|
\le C_q\,m^3\,\delta_q^{\,m}.
\]
特别地，\(R_{m,q}\) 以指数速度趋于 \(\ZZ/q\ZZ\) 上的均匀分布。
当 \(q=2\) 时，
\[
\delta_2=\frac{\rho(B(-1))}{2},
\]
其中 \(\rho(B(-1))\) 为多项式
\[
\mu^4-\mu^3+\mu^2-\mu-2
\]
的最大实根。
\end{corollary}
\leanverified{paper\_fold\_gauge\_anomaly\_congruence\_mixing}

\begin{proof}
离散 Fourier 反演给出
\[
\mathbb P(R_{m,q}=r)
=\frac1q\sum_{j=0}^{q-1}\omega_q^{-jr}\,\mathbb E\!\left[(\omega_q^{\,j})^{G_m}\right]
=\frac1q+\frac1q\sum_{j=1}^{q-1}\omega_q^{-jr}M_m(\omega_q^{\,j}),
\]
从而
\[
\left|\mathbb P(R_{m,q}=r)-\frac1q\right|
\le \frac1q\sum_{j=1}^{q-1}\left|M_m(\omega_q^{\,j})\right|.
\]
对固定 \(u\in\CC\)，由命题 \ref{prop:fold-gauge-anomaly-mgf-order4-recurrence} 的四阶递推与 Jordan 分解，
\[
\widetilde M_m(u)=\sum_{\ell} p_\ell(m)\lambda_\ell(u)^m,
\]
其中 \(\lambda_\ell(u)\) 是 \(B(u)\) 的特征值，且 \(\deg p_\ell\le 3\)。
故存在 \(C(u)>0\) 使
\[
|\widetilde M_m(u)|\le C(u)\,m^3\,\rho(B(u))^m,
\qquad
|M_m(u)|\le C(u)\,m^3\!\left(\frac{\rho(B(u))}{2}\right)^m.
\]
当 \(u=\omega_q^{\,j}\) 且 \(1\le j\le q-1\) 时，\(u\neq 1\) 且 \(|u|=1\)，由定理 \ref{thm:fold-gauge-anomaly-unit-circle-rigidity} 得 \(\rho(B(u))<2\)。
取
\[
\delta_q=\max_{1\le j\le q-1}\frac{\rho(B(\omega_q^{\,j}))}{2}<1
\]
并吸收有限个 \(j\) 的常数，即得所述不等式。
\(q=2\) 情形由 \(u=-1\) 代入 \(F(\mu,u)\) 直接得到。
证毕。
\end{proof}

\begin{theorem}[大模数同余谱隙的扩散标度]\label{thm:fold-gauge-anomaly-large-q-delta}
\leanverified{paper\_fold\_gauge\_anomaly\_large\_q\_delta}
设 \(g_\ast=\frac49\)、\(\sigma_G^2=\frac{118}{243}\)（定理 \ref{thm:fold-gauge-anomaly-clt}）。
则存在 \(q_0\ge 2\)，使对所有 \(q\ge q_0\) 有
\[
\delta_q
=\exp\!\left(-\frac{2\pi^2\sigma_G^2}{q^2}+O(q^{-4})\right)
=\exp\!\left(-\frac{236\pi^2}{243\,q^2}+O(q^{-4})\right),
\]
等价地，
\[
\delta_q
=1-\frac{2\pi^2\sigma_G^2}{q^2}+O(q^{-4}).
\]
\end{theorem}

\begin{proof}
写
\[
F(\mu,u)=\mu^4-\mu^3-\mu^2(2u+1)+\mu(-u^3+2u)+2u.
\]
由命题 \ref{prop:fold-gauge-anomaly-pressure}，\(F(2,1)=0\) 且
\[
\partial_\mu F(2,1)=9\neq 0.
\]
由隐函数定理，存在 \(u=1\) 邻域上的解析主导根分支 \(\mu(u)\)（\(\mu(1)=2\)）。
令
\[
\Lambda(\theta):=\frac{\mu(e^\theta)}{2}.
\]
在 \(\theta=0\) 邻域内，\(\log \Lambda(\theta)=P_G(\theta)\)（实轴上一致），故
\[
P_G(it)=ig_\ast t-\frac{\sigma_G^2}{2}t^2+O(t^3),
\qquad
\Re P_G(it)=-\frac{\sigma_G^2}{2}t^2+O(t^4),
\]
其中第二式用到实系数解析函数在虚轴上的偶性。
因此当 \(|t|\) 足够小时，
\[
\rho(A_{it})
=\exp\!\bigl(\Re P_G(it)\bigr)
=\exp\!\left(-\frac{\sigma_G^2}{2}t^2+O(t^4)\right).
\]
取 \(t=2\pi/q\) 得
\[
\rho\!\left(A_{2\pi i/q}\right)
=\exp\!\left(-\frac{2\pi^2\sigma_G^2}{q^2}+O(q^{-4})\right).
\]
又由定理 \ref{thm:fold-gauge-anomaly-unit-circle-rigidity}，所有非零频率均严格小于 \(1\)；
当 \(q\) 足够大时，\(\delta_q=\max_{1\le j\le q-1}\rho(A_{2\pi i j/q})\) 的最大项由最小非零频率 \(j=1\)（等价地 \(j=q-1\)）取得，从而得到所述展开。
证毕。
\end{proof}

\begin{theorem}[随窗口增长的同余阈值相变]\label{thm:fold-gauge-anomaly-congruence-phase-transition}
令 \(q=q(m)\ge 2\) 随 \(m\to\infty\) 变化，并记
\[
R_{m,q(m)}:=G_m\bmod q(m).
\]
则有：
\begin{enumerate}
  \item 若 \(\frac{m}{q(m)^2}\to\infty\)，则
  \[
  \sup_{r\in\ZZ/q(m)\ZZ}
  \left|\mathbb P(R_{m,q(m)}=r)-\frac{1}{q(m)}\right|
  \longrightarrow 0,
  \]
  且该差异可压到 \(\exp(-c\,m/q(m)^2)\) 量级（某个 \(c>0\)）。
  \item 若 \(\frac{q(m)}{\sqrt m}\to\infty\)，则
  \[
  \left\|\mathcal L(R_{m,q(m)})-\mathrm{Unif}(\ZZ/q(m)\ZZ)\right\|_{\mathrm{TV}}
  \longrightarrow 1.
  \]
\end{enumerate}
\end{theorem}

\begin{proof}
对任意 \(q\ge 2\)，Fourier 反演给出
\[
\sup_r\left|\mathbb P(R_{m,q}=r)-\frac1q\right|
\le \frac1q\sum_{j=1}^{q-1}\left|\mathbb E\!\left[e^{2\pi i j G_m/q}\right]\right|.
\]
由推论 \ref{cor:fold-gauge-anomaly-congruence-mixing} 的 Jordan 上界与定理 \ref{thm:fold-gauge-anomaly-large-q-delta} 的小频率展开，可得
\[
\left|\mathbb E\!\left[e^{2\pi i j G_m/q}\right]\right|
\le C\,m^3\exp\!\left(-c_0\,\frac{j^2m}{q^2}\right)
\]
（常数与 \(q\) 足够大时的局部频率区间统一可取；高频段由定理 \ref{thm:fold-gauge-anomaly-unit-circle-rigidity} 的严格谱隙给出更强指数衰减）。
故
\[
\sup_r\left|\mathbb P(R_{m,q}=r)-\frac1q\right|
\le C' q\,m^3\,\exp\!\left(-c_0\frac{m}{q^2}\right).
\]
当 \(m/q(m)^2\to\infty\) 时，前因子 \(q\,m^3\) 可被指数项吸收，得到 (1)。

对 (2)，由定理 \ref{thm:fold-gauge-anomaly-clt} 的方差线性增长，
对任意 \(C>0\) 存在 \(m_0(C)\) 使 \(m\ge m_0(C)\) 时
\[
\mathbb P\!\left(|G_m-g_\ast m|\le C\sqrt m\right)\ge 1-\varepsilon_C,
\qquad
\varepsilon_C:=\frac{2\sigma_G^2}{C^2}.
\]
记
\[
S_m:=\left\{r\in\ZZ/q\ZZ:\exists n\in\ZZ,\ |n-g_\ast m|\le C\sqrt m,\ n\equiv r\!\!\!\pmod q\right\},
\]
则 \(|S_m|\le 2C\sqrt m+1\)，且 \(\mathbb P(R_{m,q}\in S_m)\ge 1-\varepsilon_C\)。
于是
\[
\left\|\mathcal L(R_{m,q})-\mathrm{Unif}(\ZZ/q\ZZ)\right\|_{\mathrm{TV}}
\ge \mathbb P(R_{m,q}\in S_m)-\frac{|S_m|}{q}
\ge (1-\varepsilon_C)-\frac{2C\sqrt m+1}{q}.
\]
若 \(q/\sqrt m\to\infty\)，则对任意固定 \(C\)，上式下极限不小于 \(1-\varepsilon_C\)；
再令 \(C\to\infty\)（故 \(\varepsilon_C\to 0\)）得到总变差极限为 \(1\)。
证毕。
\leanverified{paper\_fold\_gauge\_anomaly\_congruence\_phase\_transition}
\end{proof}

\begin{proposition}[非渐近 Chernoff 轮廓：尾概率指数由代数压力函数唯一决定]\label{prop:fold-gauge-anomaly-chernoff-nonasymptotic}
\leanverified{paper\_fold\_gauge\_anomaly\_chernoff\_nonasymptotic}
令压力函数 \(P_G(\theta)=\log\rho(A_\theta)\) 如命题 \ref{prop:fold-gauge-anomaly-pressure}，并令
\[
I(a)=\sup_{\theta\in\RR}\bigl(a\theta-P_G(\theta)\bigr)\qquad(a\in[0,1])
\]
为其 Legendre--Fenchel 变换（命题 \ref{prop:fold-gauge-anomaly-ldp}）。
则对任意 \(a\in(0,1)\) 存在常数 \(C(a)>0\) 与整数 \(r(a)\ge 0\)，使得对所有 \(m\ge 1\) 有统一上界
\begin{equation}\label{eq:fold-gauge-anomaly-chernoff-two-sided}
\boxed{\;
\mathbb P\!\left(\frac{G_m}{m}\ge a\right)\le C(a)\,m^{r(a)}e^{-m I(a)},\qquad
\mathbb P\!\left(\frac{G_m}{m}\le a\right)\le C(a)\,m^{r(a)}e^{-m I(a)}.\;}
\end{equation}
其中鞍点参数 \(u(a)=e^{\theta(a)}\) 可由消元证书 \(R(a,u)=0\) 的可行分支唯一确定（命题 \ref{prop:fold-gauge-anomaly-rate-curve-param}），故指数轮廓 \(I(a)\) 为可审计的“代数—对数型”对象；
而多项式前因子 \(m^{r(a)}\) 的出现仅由有限维转移算子的 Jordan 缺陷决定（在本 \(4\) 态接口下可取 \(r(a)\le 3\)）。
\end{proposition}

\begin{proof}[证明要点（矩阵 Chernoff + 有限维谱分解）]
对任意 \(\theta\in\RR\)，Chernoff 不等式给出
\(
\mathbb P(G_m\ge am)\le e^{-am\theta}\,\mathbb E e^{\theta G_m}
\)；
而 \(\mathbb E e^{\theta G_m}\) 是有限维加权邻接矩阵 \(A_\theta\) 的某个矩阵元（或其线性组合），故可由有限维谱/Jordan 分解界为
\(
\mathbb E e^{\theta G_m}\le C(\theta)\,m^{r(\theta)}\rho(A_\theta)^m
\)。
优化 \(\theta\) 得指数项 \(I(a)\)，并将剩余常数与多项式因子吸收为 \(C(a),r(a)\)。证毕。
\end{proof}

\begin{theorem}[Gallavotti--Cohen 缺陷的严格符号律与远场饱和]\label{thm:fold-gauge-anomaly-gc-defect-sign}
定义 GC 缺陷函数
\[
\Delta(\theta):=P_G(\theta)-\theta-P_G(-\theta).
\]
则：
\begin{enumerate}
  \item \(\Delta\) 为奇函数，\(\Delta(0)=0\)；
  \item 其线性化偏置满足
  \begin{equation}\label{eq:fold-gauge-anomaly-gc-defect-slope}
  \boxed{\;
  \Delta'(0)=2P_G'(0)-1=2\cdot\frac49-1=-\frac19\;};
  \end{equation}
  \item 其远场极限满足
  \begin{equation}\label{eq:fold-gauge-anomaly-gc-defect-saturation}
  \boxed{\;
  \lim_{\theta\to+\infty}\Delta(\theta)=-\log\varphi\;};
  \end{equation}
  \item 严格符号律成立：
  \begin{equation}\label{eq:fold-gauge-anomaly-gc-defect-sign}
  \boxed{\;
  \Delta(\theta)<0\ (\theta>0),\qquad \Delta(\theta)>0\ (\theta<0)\;}.
  \end{equation}
\end{enumerate}
\leanverified{paper\_fold\_gauge\_anomaly\_gc\_defect\_sign}
\end{theorem}

\begin{proof}[证明要点（可审计）]
（1）由定义直接计算得 \(\Delta(-\theta)=-\Delta(\theta)\)。
（2）由 \(\Delta'(\theta)=P_G'(\theta)-1+P_G'(-\theta)\) 取 \(\theta=0\)，并用命题 \ref{prop:fold-gauge-anomaly-pressure} 的闭式 \(P_G'(0)=4/9\)。
（3）由命题 \ref{prop:fold-gauge-anomaly-ldp} 的端点闭式：当 \(\theta\to+\infty\) 时 \(P_G(\theta)=\theta-\log 2+o(1)\)，而 \(P_G(-\theta)\to \log(\varphi/2)\)，故 \(\Delta(\theta)\to -\log\varphi\)。
（4）对该最小 \(4\) 态接口，\(\Delta\) 在 \(\theta>0\) 上严格递减（等价于 \(\Delta'(\theta)<0\)），可由命题 \ref{prop:fold-gauge-anomaly-rate-curve-param} 的代数参数化把 \(\Delta'(\theta)\) 化为代数分支上的显式表达并作符号审计；于是由 \(\Delta(0)=0\) 得 \(\theta>0\) 时 \(\Delta(\theta)<0\)。\(\theta<0\) 的符号由奇性得到。证毕。
\end{proof}

\begin{corollary}[两常数闭合：\texorpdfstring{$\frac49$}{4/9} 与 \texorpdfstring{$\log\varphi$}{log phi} 锁定同一缺陷曲线的局部与远场]\label{cor:fold-gauge-anomaly-two-constants-closure}
\leanverified{paper\_fold\_gauge\_anomaly\_two\_constants\_closure}
在定理 \ref{thm:fold-gauge-anomaly-gc-defect-sign} 的记号下，\(\Delta'(0)=-\frac19\) 与
\(\lim_{\theta\to+\infty}\Delta(\theta)=-\log\varphi\) 构成一对局部—远场规范条件：
它们分别由 \(P_G'(0)=4/9\) 的均值常数与端点指数 \(I(0)=\log(2/\varphi)\)（命题 \ref{prop:fold-gauge-anomaly-ldp}）诱导。
并且在原点处存在严格的三阶校正（偶次项完全相消）：
\begin{equation}\label{eq:fold-gauge-anomaly-delta-taylor3}
\boxed{\;
\Delta(\theta)= -\frac{\theta}{9}\;-\;\frac{1174}{6561}\,\theta^{3}\;+\;O(\theta^{5})\qquad(\theta\to 0).}
\end{equation}
因此 \(\Delta\) 的“非自对偶”指纹不止于三阶累积量非零，而是以单参数缺陷曲线的局部斜率与远场饱和值同时封口。
\end{corollary}

\begin{proof}
由命题 \ref{prop:fold-gauge-anomaly-pressure}，\(P_G\) 在 \(\theta=0\) 处的前三阶导数均有闭式，
并且 \(P_G\) 为光滑函数。
将 \(P_G(\theta)\) 与 \(P_G(-\theta)\) 作 Taylor 展开，代入 \(\Delta(\theta)=P_G(\theta)-\theta-P_G(-\theta)\)；
注意偶次项相消、奇次项加倍，即得 \eqref{eq:fold-gauge-anomaly-delta-taylor3}。证毕。
\end{proof}

\begin{proposition}[折叠条件期望的有限指数：最大纤维多重度给出最坏信息损失的对数尺度]\label{prop:fold-condexp-index-maxfiber}
固定分辨率 \(m\)。
沿用小节 \ref{subsec:pom-max-fiber} 的记号，记纤维多重度
\(
d_m(x)=|\Fold_m^{-1}(x)|
\)
与最大多重度
\(
D_m=\max_{x\in X_m} d_m(x)
\)。
取有限维交换代数 \(A_m=\CC^{\Omega_m}\) 及其子代数 \(B_m\subset A_m\)（由 \(X_m\) 上函数拉回得到的纤维常值函数），
并令 \(E_m:A_m\to B_m\) 为按纤维均匀平均的条件期望。
则其 Pimsner--Popa 指数满足严格等式
\begin{equation}\label{eq:fold-condexp-index-maxfiber}
\boxed{\;
\mathrm{Ind}(E_m)=D_m\; }.
\end{equation}
并且由定理 \ref{thm:dephys-index-max-relative-entropy-loss}（有限指数的相对熵损失极值公式）得到有限维离散原型：
\begin{equation}\label{eq:fold-condexp-index-entropy-loss}
\boxed{\;
\sup_{\rho,\sigma}\Bigl[D(\rho\|\sigma)-D(\rho\circ E_m\ \|\ \sigma\circ E_m)\Bigr]
=\log \mathrm{Ind}(E_m)=\log D_m.\;}
\end{equation}
\leanverified{paper\_fold\_condexp\_index\_maxfiber\_support}
\leanverified{paper\_fold\_condexp\_index\_maxfiber}
\end{proposition}

\begin{proof}[证明要点]
对任意非负函数 \(f\in A_m\) 与任意 \(a\in\Omega_m\)，若 \(x=\Fold_m(a)\) 且 \(d=d_m(x)\)，则
\(
(E_m f)(a)=\frac1d\sum_{b\in\Fold_m^{-1}(x)} f(b)\ge \frac1d f(a)
\)。
因此 \(E_m(f)\ge D_m^{-1}f\)，即 \(\mathrm{Ind}(E_m)\le D_m\)。
反过来，取某个达到最大多重度的 \(x_\star\) 及其纤维中的一点 \(a_\star\)，令 \(f=\mathbf 1_{\{a_\star\}}\)。
则在该纤维上 \(E_m f \equiv 1/D_m\)，从而若 \(E_m(f)\ge \Lambda^{-1}f\) 成立必有 \(\Lambda\ge D_m\)，故 \(\mathrm{Ind}(E_m)\ge D_m\)。
合并得 \(\mathrm{Ind}(E_m)=D_m\)。
式 \eqref{eq:fold-condexp-index-entropy-loss} 由定理 \ref{thm:dephys-index-max-relative-entropy-loss} 直接推出。证毕。
\end{proof}

\begin{remark}[视界极值态的图论普适性：最大信息损失归约为独立集计数优化]\label{rem:fold-horizon-extremal-independent-set}
命题 \ref{prop:fold-condexp-index-maxfiber} 表明：在该离散视界原型中，“最坏外部化损失”的对数尺度由最大纤维多重度 \(D_m\) 唯一锁定，并且其证书见证可取自某条达到 \(D_m\) 的极大纤维（证明中取 \(f=\mathbf 1_{\{a_\star\}}\) 即给出指数的饱和下界见证）。
\par
另一方面，由纤维的独立集坐标化定理 \ref{thm:pom-fiber-indset-factorization}，对任意 \(x\in X_m\) 存在候选位点诱导图 \(\Gamma_x\) 使得
\[
\Fold_m^{-1}(x)\cong \mathsf{Ind}(\Gamma_x),
\qquad
d_m(x)=|\mathsf{Ind}(\Gamma_x)|=Z_x(1),
\]
并且当 \(\Gamma_x\cong\bigsqcup_i P_{\ell_i}\) 时有严格乘法分解
\(
d_m(x)=\prod_i F_{\ell_i+2}
\)。
因此
\[
D_m=\max_{x\in X_m} d_m(x)=\max_{x\in X_m}|\mathsf{Ind}(\Gamma_x)|
\]
把“最大信息损失（最大纤维）”严格等价为一个纯组合优化问题：在族 \(\{\Gamma_x\}_{x\in X_m}\) 中最大化独立集总数（等价地最大化路径分量长度谱所诱导的 Fibonacci 乘积因子）。
该归约的输入仅为 \(\Gamma_x\) 的分量结构，因而与其余微观实现细节解耦，形成一条可对外引用的图论普适性结论。
\end{remark}

\begin{corollary}[纤维熵势的饱和判据：平均对数多重度与平均多重度之间的严格缺口]\label{cor:fold-fiber-entropy-saturation}
\leanverified{paper\_fold\_fiber\_entropy\_saturation}
对任意 \(\pi\in\Delta(X_m)\)，定义纤维熵势
\[
S_{\mathrm{fib}}(\pi):=\mathbb E_{\pi}[\log d_m(X)].
\]
则恒有 Jensen 上界
\begin{equation}\label{eq:fold-fiber-entropy-jensen-gap}
\boxed{\;
S_{\mathrm{fib}}(\pi)\le \log\mathbb E_{\pi}[d_m(X)]\;}
\end{equation}
且等号成立当且仅当 \(d_m(X)\) 在 \(\pi\)-几乎处处为常数（即所有可见类型承载相同的纤维多重度）。
特别地，对均匀可见基线 \(\pi=\mathrm{Unif}(X_m)\) 有 \(\mathbb E_{\pi}[d_m(X)]=|\Omega_m|/|X_m|=I_m\)（定义 \ref{def:pom-fold-index}），从而
\[
\boxed{\;
\mathbb E_{\pi}[\log d_m(X)]\le \log\frac{|\Omega_m|}{|X_m|}=S_Z(m)\; }.
\]
\end{corollary}

\begin{proof}
由对数的凹性（Jensen 不等式）得 \eqref{eq:fold-fiber-entropy-jensen-gap}。
等号条件为 \(d_m(X)\) 在 \(\pi\)-支撑上为常数。
最后一段由
\(
\mathbb E_{\mathrm{Unif}(X_m)}[d_m(X)]=\frac{1}{|X_m|}\sum_{x\in X_m} d_m(x)=\frac{|\Omega_m|}{|X_m|}
\)
得出。证毕。
\end{proof}
\leanverified{paper\_fold\_fiber\_entropy\_saturation}

\begin{proposition}[可审计的“谱—层析—证书”三角：功率谱与四阶递推共同锁定最小识别复杂度]\label{prop:fold-gauge-anomaly-spectrum-tomography-certificate-triangle}
设观测者仅能访问两类有限数据：
\begin{enumerate}
  \item[(i)] 功率谱在有限频点网格上的采样 \(\{\mathcal S_g(\omega_\ell)\}_{\ell=1}^{L}\)；
  \item[(ii)] 对某个固定 \(u\neq 1\)，有限长度生成函数值 \(\{M_m(u)\}_{m=0}^{N}\)（其中 \(N\ge 3\)）。
\end{enumerate}
则：
\begin{enumerate}
  \item 由命题 \ref{prop:fold-gauge-anomaly-mgf-order4-recurrence}，数据 \(\{M_m(u)\}_{m=0}^{N}\) 在代数上必须落在四阶递推生成的四维线性空间中；
  任意声称“同机制”的实现都必须满足该递推的有限数据一致性约束。
  \item 由定理 \ref{thm:fold-gauge-anomaly-psd-rational}，\(\mathcal S_g\) 具有固定的有理极点结构；
  任意机制替换若改变 \( (-1/2)\) 交错模态或其 Jordan 残差（定理 \ref{thm:fold-gauge-anomaly-A0-jordan}），则必在有限个频点上产生可检测偏差。
  \item 两类数据联合后，最小识别复杂度被“阶数 \(4\)”强制：
  任意试图用阶数 \(<4\) 的线性状态模型同时拟合 (i)--(ii) 两类数据，将在严格代数意义下失败（无法同时满足四阶递推与谱有理式的极点/Jordan 约束）。
\end{enumerate}
因此，规范差机制的最低阶可核验承载维数被严格锁定为 \(4\)，并可由“递推一致性 + 频域极点结构”双通道证书同时审计。
\end{proposition}
\leanverified{paper\_fold\_gauge\_anomaly\_spectrum\_tomography\_certificate\_triangle}

\subsubsection{均匀基线下规范差的有限长度矩闭式}\label{subsubsec:fold-gauge-anomaly-uniform-finite-moment-closures}
在定义 \ref{def:fold-gauge-anomaly} 的符号下，取
\[
G_m(\omega)=\bigl|\,r_{m+1\to m}(\omega)\oplus \widetilde r_{m+1\to m}(\omega)\,\bigr|_0,\qquad
\omega\sim\mathrm{Unif}(\Omega_{m+1}).
\]
并记
\[
\bar G_m:=\mathbb E[G_m],\qquad
M_m(u):=\mathbb E[u^{G_m}],\qquad
\widetilde M_m(u):=2^m M_m(u).
\]
采用空窗约定 \(m=0\) 时 \(G_0\equiv 0\)，以下闭式对一切 \(m\ge 0\) 成立。

\begin{theorem}[均值的精确有限长度闭式]\label{thm:fold-gauge-anomaly-mean-finite-closed}
规范差均值满足
\[
\boxed{\;
\bar G_m
=\frac49\,m-\frac{29}{54}
+\frac58\,2^{-m}
+\Bigl(\frac{m}{36}-\frac{19}{216}\Bigr)\Bigl(-\frac12\Bigr)^m
\; }.
\]
因此
\[
\bar G_m=\frac49\,m-\frac{29}{54}+O(m2^{-m}),
\qquad
\frac{\bar G_m}{m}=\frac49-\frac{29}{54m}+O(2^{-m}).
\]
\leanverified{paper\_fold\_gauge\_anomaly\_mean\_finite\_closed}
\end{theorem}

\begin{proof}
由有限状态换能器路径权重求和，\(\widetilde M_m(u)\) 满足四阶递推
\[
\widetilde M_{m+4}
=\widetilde M_{m+3}+(2u+1)\widetilde M_{m+2}+(u^3-2u)\widetilde M_{m+1}-2u\,\widetilde M_m.
\]
对 \(u\) 求导并代入 \(u=1\)。由 \(M_m(1)=1\) 得 \(\widetilde M_m(1)=2^m\)。设
\[
B_m:=\widetilde M_m'(1)=2^m\bar G_m,
\]
则
\[
B_{m+4}=B_{m+3}+3B_{m+2}-B_{m+1}-2B_m+2^{m+3}.
\]
其齐次特征多项式为
\[
\mu^4-\mu^3-3\mu^2+\mu+2=(\mu-2)(\mu-1)(\mu+1)^2,
\]
故通解型为
\[
B_m=(\alpha m+\beta)2^m+\gamma+(\delta+\varepsilon m)(-1)^m.
\]
由定义直接枚举 \(m=0,1,2,3,4\) 可得
\[
\bar G_0=0,\quad
\bar G_1=\frac14,\quad
\bar G_2=\frac12,\quad
\bar G_3=\frac78,\quad
\bar G_4=\frac{41}{32},
\]
解得常数并化简为
\[
B_m=\Bigl(\frac49m-\frac{29}{54}\Bigr)2^m+\frac58+\Bigl(\frac{m}{36}-\frac{19}{216}\Bigr)(-1)^m.
\]
两边除以 \(2^m\) 即得结论。
\end{proof}

\begin{proposition}[均值密度的全局单调性]\label{prop:fold-gauge-anomaly-mean-density-monotone}
序列 \(\bigl(\bar G_m/m\bigr)_{m\ge 1}\) 单调不减，且
\[
\frac{\bar G_{m+1}}{m+1}>\frac{\bar G_m}{m}\quad(m\ge 2),
\]
并有
\[
\lim_{m\to\infty}\frac{\bar G_m}{m}=\frac49.
\]
\end{proposition}

\begin{proof}
由定理 \ref{thm:fold-gauge-anomaly-mean-finite-closed} 直接计算
\[
\frac{\bar G_{m+1}}{m+1}-\frac{\bar G_m}{m}
=
\frac{
232\cdot 2^m-135m-270+(-1)^m(-18m^2+39m+38)
}{
432\cdot 2^m\,m(m+1)
}.
\]
分母恒正，只需考察分子 \(N_m\)。

当 \(m\) 为偶数时，
\[
N_m=232\cdot 2^m-18m^2-96m-232.
\]
对 \(m\ge 2\) 有 \(2^m\ge 4\)，故
\[
N_m\ge 232\cdot 4-18\cdot 2^2-96\cdot 2-232>0.
\]

当 \(m\) 为奇数时，
\[
N_m=232\cdot 2^m+18m^2-174m-308.
\]
代入 \(m=1\) 得 \(N_1=0\)，即 \(\bar G_2/2=\bar G_1\)；对 \(m\ge 3\) 有 \(N_m>0\)。
于是得到“单调不减且 \(m\ge 2\) 后严格递增”。
极限由定理 \ref{thm:fold-gauge-anomaly-mean-finite-closed} 的渐近式立即给出。
\leanverified{paper\_fold\_gauge\_anomaly\_mean\_density\_monotone}
\end{proof}

\begin{theorem}[二阶阶乘矩的精确闭式]\label{thm:fold-gauge-anomaly-second-factorial-finite-closed}
对一切 \(m\ge 0\)，有
\[
\boxed{\;
\mathbb E\!\bigl[G_m(G_m-1)\bigr]
=\frac{16}{81}m^2-\frac{106}{243}m+\frac{443}{729}
-\Bigl(\frac{5}{16}m+\frac12\Bigr)2^{-m}
+\Bigl(-\frac{m^3}{648}+\frac{m^2}{27}-\frac{m}{432}-\frac{157}{1458}\Bigr)\Bigl(-\frac12\Bigr)^m
\; }.
\]
\end{theorem}

\begin{proof}
令
\[
C_m:=\widetilde M_m''(1)=2^m\,\mathbb E\!\bigl[G_m(G_m-1)\bigr].
\]
对 \(\widetilde M_m\) 的四阶递推二次求导并代入 \(u=1\)，得到
\[
C_{m+4}=C_{m+3}+3C_{m+2}-C_{m+1}-2C_m
+4B_{m+2}+2B_{m+1}-4B_m+6\cdot 2^{m+1},
\]
其中 \(B_m=\widetilde M_m'(1)\)。
由定理 \ref{thm:fold-gauge-anomaly-mean-finite-closed} 可将右端完全改写为 \(m\) 的显式线性组合。
该非齐次递推的齐次特征因子仍为 \((\mu-2)(\mu-1)(\mu+1)^2\)，故 \(C_m\) 必可写为
\[
C_m
=2^m(\alpha m^2+\beta m+\gamma)
+\delta m+\kappa
+(-1)^m(\varepsilon m^3+\zeta m^2+\eta m+\theta).
\]
代回递推确定 \(\alpha,\beta,\delta,\varepsilon,\zeta\)，再由 \(m=0,1,2,3\) 的直接枚举确定其余常数，整理并除以 \(2^m\) 即得结论。
\end{proof}

\leanverified{paper\_fold\_gauge\_anomaly\_second\_factorial\_finite\_closed}

\begin{theorem}[方差的精确有限长度闭式]\label{thm:fold-gauge-anomaly-variance-finite-window-closed}
对一切 \(m\ge 0\)，规范差方差满足
\[
\boxed{\;
\begin{aligned}
\mathrm{Var}(G_m)
={}&\frac{118}{243}\,m-\frac{635}{2916}
+\Bigl(\frac{43}{54}-\frac{125}{144}m\Bigr)2^{-m}\\
&+\Bigl(-\frac{m^3}{648}+\frac{m^2}{81}+\frac{173}{1296}m-\frac{47}{162}\Bigr)\Bigl(-\frac12\Bigr)^m\\
&+\Bigl(-\frac{m^2}{1296}+\frac{19}{3888}m-\frac{9293}{23328}\Bigr)4^{-m}
+\Bigl(-\frac{5}{144}m+\frac{95}{864}\Bigr)\Bigl(-\frac14\Bigr)^m.
\end{aligned}
\; }.
\]
因此
\[
\mathrm{Var}(G_m)=\frac{118}{243}\,m-\frac{635}{2916}+O(m^3 2^{-m}).
\]
\end{theorem}

\begin{proof}
由定理 \ref{thm:fold-gauge-anomaly-second-factorial-finite-closed}，
\[
\mathbb E[G_m^2]=\mathbb E[G_m(G_m-1)]+\mathbb E[G_m].
\]
再由定理 \ref{thm:fold-gauge-anomaly-mean-finite-closed}，
\[
\mathrm{Var}(G_m)=\mathbb E[G_m^2]-\bigl(\mathbb E[G_m]\bigr)^2.
\]
将两式代入并作代数化简即可得到所述闭式与渐近式。
\end{proof}

\leanverified{paper\_fold\_gauge\_anomaly\_variance\_finite\_window\_closed}

\begin{remark}[与平稳和口径的关系]\label{rem:fold-gauge-anomaly-finite-vs-stationary-moments}
本小节闭式针对有限窗口均匀基线 \(\omega\sim\mathrm{Unif}(\Omega_{m+1})\) 下由定义 \ref{def:fold-gauge-anomaly} 直接给出的 \(G_m(\omega)\)。
定理 \ref{thm:fold-gauge-anomaly-var-finite-closed} 采用的是平稳逐位过程口径 \(G_m=\sum_{t=1}^m g_t\)。
二者共享相同的一阶与二阶线性系数 \(\frac49\) 与 \(\frac{118}{243}\)，而有限长度修正项反映了不同边界条件下的可观测差异。
\end{remark}

\paragraph{均值与方差序列的 Hankel 刚性与坏素数}
本小节的均值/方差闭式不仅给出渐近常数，还强制这些序列在代数与模素数口径下呈现严格的 Hankel 刚性。

\begin{theorem}[均值序列的五阶 Hankel 刚性与坏素数 \texorpdfstring{$\{2,5\}$}{\{2,5\}}]\label{thm:fold-gauge-anomaly-mean-hankel-rigidity-rank5-badprimes}
\leanverified{paper\_fold\_gauge\_anomaly\_mean\_hankel\_rigidity\_rank5\_badprimes}
在定理 \ref{thm:fold-gauge-anomaly-mean-finite-closed} 的记号下，定义整数序列
\[
B_m^{\mathrm{int}}:=2^{m+1}\bar G_m\in\ZZ_{\ge 0}\qquad(m\ge 0).
\]
则成立：
\begin{enumerate}
  \item \(B^{\mathrm{int}}\) 满足极小阶为 \(5\) 的常系数线性递推；其特征多项式为
  \[
  \chi_{B}(\lambda)=(\lambda-2)^2(\lambda-1)(\lambda+1)^2.
  \]
  \item 因而无穷 Hankel 矩阵 \(\bigl(B^{\mathrm{int}}_{i+j}\bigr)_{i,j\ge 0}\) 的 Hankel 秩恰为 \(5\)；等价地，任意 \(6\times 6\) Hankel 子式行列式恒为零。
  \item 令
  \[
  H^{(5)}_{\ell}(B^{\mathrm{int}}):=\bigl(B^{\mathrm{int}}_{\ell+i+j}\bigr)_{0\le i,j\le 4}\in M_5(\ZZ)
  \qquad(\ell\ge 0)
  \]
  为平移 Hankel 块，则
  \[
  \det H^{(5)}_{\ell}(B^{\mathrm{int}})=1280\cdot 4^{\ell}\qquad(\ell\ge 0).
  \]
  特别地，\(\det H^{(5)}_{\ell}(B^{\mathrm{int}})\) 的全部素因子集合恒等于 \(\{2,5\}\)。
  因此对任意素数 \(p\notin\{2,5\}\)，在 \(\FF_p\) 上的 Hankel 秩稳定为 \(5\)；而在 \(p\in\{2,5\}\) 时该见证块对所有 \(\ell\) 均退化。
\end{enumerate}
\end{theorem}
\begin{proof}
由定理 \ref{thm:fold-gauge-anomaly-mean-finite-closed} 的闭式，直接得到
\[
B_m^{\mathrm{int}}
=\Bigl(\frac{8}{9}m-\frac{29}{27}\Bigr)2^{m}+\frac54+\Bigl(\frac{m}{18}-\frac{19}{108}\Bigr)(-1)^m.
\]
因此 \(B_m^{\mathrm{int}}\) 属于向量空间
\(
\mathrm{span}_{\QQ}\{2^m,\ m2^m,\ 1,\ (-1)^m,\ m(-1)^m\}
\)，从而被差分算子 \(E\)（\(Ex_m=x_{m+1}\)）的多项式
\(
\chi_B(E)=(E-2)^2(E-1)(E+1)^2
\)
所湮灭，给出阶 \(\le 5\) 的递推。
对不同特征根的指数族用算子 \((E-2)\) 与 \((E+1)\) 可逐一分离，遂上述五个基函数在 \(\QQ\) 上线性无关，极小阶即为 \(5\)。
\par
由 \(\chi_B\) 构造伴随矩阵 \(A_B\in M_5(\ZZ)\)，使状态向量
\(
s_\ell:=(B^{\mathrm{int}}_{\ell},\dots,B^{\mathrm{int}}_{\ell+4})^{\mathsf T}
\)
满足 \(s_{\ell+1}=A_B s_\ell\)。
注意到
\(
H^{(5)}_{\ell}(B^{\mathrm{int}})=[\,s_\ell\ s_{\ell+1}\ \cdots\ s_{\ell+4}\,]
\)，
从而
\(
H^{(5)}_{\ell+1}(B^{\mathrm{int}})=H^{(5)}_{\ell}(B^{\mathrm{int}})\,A_B
\)
并给出
\[
\det H^{(5)}_{\ell+1}(B^{\mathrm{int}})=\det H^{(5)}_{\ell}(B^{\mathrm{int}})\cdot \det(A_B).
\]
由于 \(\chi_B\) 的常数项为 \(-4\)，伴随矩阵满足 \(\det(A_B)=4\)，故
\(
\det H^{(5)}_{\ell}(B^{\mathrm{int}})=\det H^{(5)}_{0}(B^{\mathrm{int}})\,4^\ell
\)。
由闭式生成 \(B^{\mathrm{int}}_0,\dots,B^{\mathrm{int}}_8\) 并对 \(H^{(5)}_0\) 作精确行列式计算得 \(\det H^{(5)}_0(B^{\mathrm{int}})=1280\)，从而得到 (3)。
由递推的整系数性知对任意域上 Hankel 秩 \(\le 5\)；
而当 \(p\notin\{2,5\}\) 时，\(\det H^{(5)}_0(B^{\mathrm{int}})\not\equiv 0\ (\mathrm{mod}\ p)\) 强制秩 \(\ge 5\)，故模 \(p\) Hankel 秩恰为 \(5\)；当 \(p\in\{2,5\}\) 时所有平移主块行列式皆为零，故秩必降。证毕。\
可复现核验脚本：\texttt{scripts/exp\_fold\_gauge\_anomaly\_hankel\_mean\_variance\_certificates.py}。
\end{proof}

\begin{theorem}[方差序列的十三阶 Hankel 刚性与振幅素数 \texorpdfstring{$59$}{59}]\label{thm:fold-gauge-anomaly-variance-hankel-rigidity-rank13-prime59}
在定理 \ref{thm:fold-gauge-anomaly-variance-finite-window-closed} 的记号下，定义整数序列
\[
V_m:=2^{2m+2}\mathrm{Var}(G_m)\in\ZZ_{\ge 0}\qquad(m\ge 0).
\]
\leanverified{paper\_fold\_gauge\_anomaly\_variance\_hankel\_rigidity\_rank13\_prime59}
则成立：
\begin{enumerate}
  \item \(V\) 满足极小阶为 \(13\) 的常系数线性递推；其特征多项式完全分解为
  \[
  \chi_V(\lambda)=(\lambda-4)^2(\lambda-2)^2(\lambda-1)^3(\lambda+1)^2(\lambda+2)^4.
  \]
  \item 因而 \(V\) 的 Hankel 秩恰为 \(13\)。
  \item 令
  \[
  H^{(13)}_{\ell}(V):=\bigl(V_{\ell+i+j}\bigr)_{0\le i,j\le 12}\in M_{13}(\ZZ)
  \qquad(\ell\ge 0),
  \]
  则
  \[
  \det H^{(13)}_{0}(V)=-2^{81}\,3^{18}\,5^{16}\,59^{2},
  \qquad
  \det H^{(13)}_{\ell}(V)=\det H^{(13)}_{0}(V)\cdot 1024^{\ell}.
  \]
  \item 因此在且仅在 \(p\in\{2,3,5,59\}\) 时，\(\FF_p\) 上的 Hankel 秩发生退化；对所有其余素数 \(p\)，模 \(p\) Hankel 秩稳定为 \(13\)。
\end{enumerate}
\end{theorem}
\begin{proof}
将定理 \ref{thm:fold-gauge-anomaly-variance-finite-window-closed} 的闭式两边乘以 \(2^{2m+2}=4^{m+1}\)，得到 \(V_m\) 为指数--多项式有限和，其指数基底为
\(
4,\ 2,\ -2,\ 1,\ -1
\)
并分别伴随次数 \(1,1,3,2,1\) 的多项式因子；
这等价于特征根重数分别为 \(2,2,4,3,2\)，从而给出 \(\chi_V(E)V\equiv 0\) 与极小阶 \(13\)（不同 \(\alpha\) 的族 \(\{m^k\alpha^m\}\) 线性无关）。
因此在任意域上 Hankel 秩 \(\le 13\)。
\par
同理由 \(\chi_V\) 构造伴随矩阵 \(A_V\in M_{13}(\ZZ)\) 并令 \(s_\ell=(V_\ell,\dots,V_{\ell+12})^{\mathsf T}\)，则 \(s_{\ell+1}=A_V s_\ell\) 且
\(
H^{(13)}_{\ell+1}(V)=H^{(13)}_{\ell}(V)\,A_V
\)。
由于 \(\chi_V\) 的常数项为 \(-1024\)，伴随矩阵满足 \(\det(A_V)=1024\)，从而
\(
\det H^{(13)}_{\ell}(V)=\det H^{(13)}_0(V)\,1024^\ell
\)。
首块行列式 \(\det H^{(13)}_0(V)\) 由闭式生成 \(V_0,\dots,V_{24}\) 后对 \(13\times 13\) Hankel 主子式作精确代数计算得到，并完全分解如上式所示。
于是对 \(p\notin\{2,3,5,59\}\) 有 \(\det H^{(13)}_0(V)\not\equiv 0\ (\mathrm{mod}\ p)\)，强制 \(\FF_p\) 上 Hankel 秩 \(\ge 13\)，与 \(\le 13\) 合并即得稳定为 \(13\)；
而对 \(p\in\{2,3,5,59\}\)，由 (3) 知所有平移主块行列式皆被 \(p\) 整除，从而秩必降。证毕。\
可复现核验脚本：\texttt{scripts/exp\_fold\_gauge\_anomaly\_hankel\_mean\_variance\_certificates.py}。
\end{proof}

\begin{corollary}[谱素数与振幅素数的可分离性：\texorpdfstring{$59$}{59} 的纯振幅性]\label{cor:fold-gauge-anomaly-spectral-vs-amplitude-primes}
设
\[
s_n=\sum_{i=1}^{r}\sum_{k=0}^{m_i-1}c_{i,k}\,n^k\,\alpha_i^n
\qquad(\alpha_i\ \text{两两不同})
\]
为指数--多项式序列，并令 \(d=\sum_i m_i\)。
则存在仅依赖于 \(\{\alpha_i\}\) 与 \(\{m_i\}\) 的汇合 Vandermonde 型矩阵 \(V(\alpha)\) 与仅依赖于系数 \(\{c_{i,k}\}\) 的块对角矩阵 \(W(c)\)，使得见证 Hankel 块满足分解
\[
H^{(d)}_{0}(s)=V(\alpha)^{\mathsf T}\,W(c)\,V(\alpha),
\qquad
\det H^{(d)}_{0}(s)=\det(V(\alpha))^{2}\cdot \det(W(c)).
\]
因此 \(\det(V(\alpha))\) 的素因子（由 \(\alpha_i-\alpha_j\) 的差值控制）构成“谱素数”，而 \(\det(W(c))\) 的素因子构成“振幅素数”。
\par
将其应用于定理 \ref{thm:fold-gauge-anomaly-mean-hankel-rigidity-rank5-badprimes} 与定理 \ref{thm:fold-gauge-anomaly-variance-hankel-rigidity-rank13-prime59}：
\begin{enumerate}
  \item 对 \(B_m^{\mathrm{int}}\) 而言，特征根集合为 \(\{2,1,-1\}\)，差值仅产生谱素数 \(\{2,3\}\)，而 \(\det H^{(5)}_0(B^{\mathrm{int}})=2^{8}\cdot 5\) 额外含素数 \(5\)，故 \(5\) 为振幅素数。
  \item 对 \(V_m\) 而言，特征根集合为 \(\{4,2,1,-1,-2\}\)，其差值至多为 \(6\)，谱素数至多为 \(\{2,3,5\}\)，而 \(\det H^{(13)}_0(V)\) 含素数 \(59\)，故 \(59\) 为纯振幅素数（不可由根差解释）。
\end{enumerate}
\end{corollary}
\begin{proof}
对每个 \(\alpha_i\)，序列 \(\{n^k\alpha_i^n\}_{0\le k\le m_i-1}\) 可实现为在 \(\alpha_i\) 处的有限阶导数评估，从而把 Hankel 块写成“节点/重数”决定的汇合 Vandermonde 乘以“系数”块对角再乘以汇合 Vandermonde 的转置；行列式分解随之成立。
两条应用由定理 \ref{thm:fold-gauge-anomaly-mean-hankel-rigidity-rank5-badprimes} 与定理 \ref{thm:fold-gauge-anomaly-variance-hankel-rigidity-rank13-prime59} 的显式素因子分解直接读出。证毕。
\end{proof}
\leanverified{paper\_fold\_gauge\_anomaly\_spectral\_vs\_amplitude\_primes}

\begin{conjecture}[高阶阶乘矩的二进模态普适性]\label{conj:fold-gauge-anomaly-factorial-moment-binary-modal-universality}
令 \((G_m)_r:=G_m(G_m-1)\cdots(G_m-r+1)\) 为降阶阶乘，并定义整数序列
\[
F_r(m):=2^{rm+r}\,\mathbb E\!\bigl[(G_m)_r\bigr]\in\ZZ\qquad(r\ge 1,\ m\ge 0).
\]
猜想对每个 \(r\ge 1\) 成立：
\begin{enumerate}
  \item \(F_r\) 的极小特征多项式在 \(\ZZ[\lambda]\) 中完全分解，其全部根均属于集合 \(\{\pm 2^{k}:0\le k\le r\}\)，并带有由 \(r\) 控制的显式重数结构；
  \item \(F_r\) 的某个极小 Hankel 见证行列式之坏素数集合有限，且可分解为“谱素数”（来自根差）与“振幅素数”（来自系数结构）。
\end{enumerate}
其中 \(r=1,2\) 的坏素数集合分别由定理 \ref{thm:fold-gauge-anomaly-mean-hankel-rigidity-rank5-badprimes} 与定理 \ref{thm:fold-gauge-anomaly-variance-hankel-rigidity-rank13-prime59} 给出为 \(\{2,5\}\) 与 \(\{2,3,5,59\}\)，为该猜想提供了可审计的低阶证据链。
\end{conjecture}


\begin{theorem}[多步降阶的规范差密度不变性]\label{thm:fold-gauge-anomaly-kstep-invariance}
固定整数 \(k\ge 1\)。
对任意 \(\omega\in\Omega_{m+k}=\{0,1\}^{m+k}\)，定义多步直接截断与折叠感知 restriction
\[
r_{m+k\to m}(\omega):=\omega_{1}\cdots\omega_{m},
\qquad
\widetilde r_{m+k\to m}(\omega):=\pi_{m+k\to m}\bigl(\Fold_{m+k}(\omega)\bigr)\in X_m,
\]
并定义 \(k\)-步规范差
\[
G_m^{(k)}(\omega):=\bigl|\,r_{m+k\to m}(\omega)\ \oplus\ \widetilde r_{m+k\to m}(\omega)\,\bigr|_0,
\]
其中 \(\oplus\) 为逐位异或、\(|\cdot|_0\) 为汉明重量。
在均匀基线 \(\omega\sim\mathrm{Unif}(\Omega_{m+k})\) 下，有
\[
\boxed{\ \lim_{m\to\infty}\frac{1}{m}\,\mathbb E\!\left[G_m^{(k)}\right]=\frac49\ },
\qquad
\boxed{\ \frac{G_m^{(k)}-\frac49 m}{\sqrt m}\Rightarrow \mathcal N\!\left(0,\frac{118}{243}\right)\ }.
\]
并且对任意 \(\theta\in\RR\)，其每步对数矩母函数极限与 \(k=1\) 情形相同：存在与 \(k\) 无关的压力函数 \(P_G(\theta)\) 使得
\[
\lim_{m\to\infty}\frac{1}{m}\log\mathbb E\exp\!\bigl(\theta\,G_m^{(k)}\bigr)=P_G(\theta)=\log\rho(A_\theta)
\]
其中 \(A_\theta\) 为命题 \ref{prop:fold-gauge-anomaly-pressure} 的同一 \(4\times 4\) 加权邻接矩阵。
因此由命题 \ref{prop:fold-gauge-anomaly-ldp}，\(G_m^{(k)}/m\) 的大偏差率函数端点闭式 \(I(0)=\log(2/\varphi), I(1)=\log 2\) 亦不随 \(k\) 改变。
\leanverified{paper\_fold\_gauge\_anomaly\_kstep\_invariance}
\end{theorem}

\begin{proof}[证明要点（结构性）]
折叠 \(\Fold_\infty\) 具有有限状态的在线实现：在 IID Bernoulli\((1/2)\) 输入下，输出与逐位失配指示可被表示为一个原始有限状态 Markov 链上的可加观测量（定理 \ref{thm:fold-gauge-anomaly-clt} 与命题 \ref{prop:fold-gauge-anomaly-pressure} 的证明要点）。
对固定 \(k\)，有限词 \(\omega\in\Omega_{m+k}\) 仅改变该双向同步实现的右边界条件；由于链的原始性，右边界条件对内点统计的影响沿距离指数衰减，故其对总和 \(G_m^{(k)}\) 的期望、方差与指数矩仅造成 \(O(1)\) 的边界修正。
将该边界修正除以 \(m\)（或以 \(\sqrt m\) 归一化）即消失，从而密度极限、CLT 常数以及压力函数 \(P_G\)（因此 LDP 率函数）均与 \(k=1\) 情形一致。证毕。
\end{proof}

\begin{remark}[数值审计：固定超前位数下的常数密度与方差率]\label{rem:fold-gauge-anomaly-kstep-mc}
为补充上述“边界影响指数衰减”的结构论证，表 \ref{tab:fold_gauge_anomaly_kstep_mc} 对若干固定 \(k\) 与较大 \(m\) 给出 Monte Carlo 审计，显示 \(\mathbb E[G_m^{(k)}]/m\) 与 \(\mathrm{Var}(G_m^{(k)})/m\) 均稳定对齐于 \(\frac49\) 与 \(\frac{118}{243}\)。
\end{remark}

% Auto-generated by scripts/exp_fold_gauge_anomaly_kstep_mc.py
\begin{table}[H]
\centering
\small
\setlength{\tabcolsep}{6pt}
\renewcommand{\arraystretch}{1.12}
\begin{tabular}{rrrrrrrr}
\toprule
$k$ & $m$ & samples & $\widehat{\mathbb{E}}[G_m^{(k)}/m]$ & SE & $\,|\Delta|$ & $\widehat{\mathrm{Var}}(G_m^{(k)})/m$ & $\,|\Delta|$ \\
\midrule
1 & 200 & 6000 & 0.442476 & 6.30e-04 & 0.001969 & 0.475526 & 0.010071 \\
1 & 600 & 6000 & 0.444109 & 3.61e-04 & 0.000336 & 0.469476 & 0.016121 \\
1 & 1200 & 6000 & 0.443808 & 2.57e-04 & 0.000637 & 0.473880 & 0.011717 \\
2 & 200 & 6000 & 0.441257 & 6.33e-04 & 0.003187 & 0.481455 & 0.004142 \\
2 & 600 & 6000 & 0.443572 & 3.72e-04 & 0.000873 & 0.498006 & 0.012409 \\
2 & 1200 & 6000 & 0.444169 & 2.63e-04 & 0.000275 & 0.498704 & 0.013107 \\
3 & 200 & 6000 & 0.442386 & 6.32e-04 & 0.002059 & 0.478692 & 0.006905 \\
3 & 600 & 6000 & 0.443820 & 3.68e-04 & 0.000624 & 0.487064 & 0.001467 \\
3 & 1200 & 6000 & 0.443822 & 2.61e-04 & 0.000622 & 0.489897 & 0.004300 \\
5 & 200 & 6000 & 0.441727 & 6.43e-04 & 0.002718 & 0.496808 & 0.011211 \\
5 & 600 & 6000 & 0.443351 & 3.63e-04 & 0.001094 & 0.474074 & 0.011523 \\
5 & 1200 & 6000 & 0.443836 & 2.63e-04 & 0.000608 & 0.496864 & 0.011267 \\
\bottomrule
\end{tabular}
\caption{Monte Carlo audit for the $k$-step gauge anomaly $G_m^{(k)}$ under the uniform baseline $\omega\sim\mathrm{Unif}(\Omega_{m+k})$. We report the sample mean density $\widehat{\mathbb{E}}[G_m^{(k)}/m]$ with its standard error (SE), and the sample variance density $\widehat{\mathrm{Var}}(G_m^{(k)})/m$. Both are compared to the $k$-independent targets $4/9$ and $118/243$ from Theorem~\ref{thm:fold-gauge-anomaly-kstep-invariance}. (Targets: $4/9\approx 0.444444$, $118/243\approx 0.485597$.)}
\label{tab:fold_gauge_anomaly_kstep_mc}
\end{table}



\subsection{分辨率折叠映射（folding）}\label{subsec:folding-map}

折叠的目标在于将任意窗口微态 $\omega=\omega_1\cdots \omega_m\in\Omega_m$ 规范化为一个合法词 $x\in X_m$。

为避免“人为编号”，本文直接以 Fibonacci 权重赋予微态可审计整数值，再取 Zeckendorf 规范形：
\begin{definition}[Fibonacci 加权值与折叠算子]\label{def:fold-word}
对 $\omega=\omega_1\cdots\omega_m\in\Omega_m$ 定义其 Fibonacci 加权值
\[
N(\omega):=\sum_{k=1}^{m}\omega_k F_{k+1}\ \in\ \NN.
\]
令 $Z(N(\omega))=(c_1,c_2,\dots)$ 为 $N(\omega)$ 的 Zeckendorf 数字（\cite{Zeckendorf1972}），定义分辨率 $m$ 的折叠
\[
\Fold_m:\Omega_m\to X_m,\qquad
\Fold_m(\omega):=\pi_m\!\bigl(Z(N(\omega))\bigr).
\]
\leanverified{Fold}
\end{definition}

\begin{remark}[折叠映射应读作“合法构造”，不只是取商]\label{rem:fold-map-admissible-construction}
定义 \ref{def:fold-word} 的重点不只是把 \(\Omega_m\) 压到 \(X_m\)，而是给出一条可审计的对象生成链：先取 Fibonacci 加权值，再进入 Zeckendorf 规范形，再截断到分辨率 \(m\) 的可见层。按 Euclid 式口径，这一步属于构造层，而不是刚性层。后面的幂等、满射、纤维规模与残差定理，才是在“对象已被合法构造出来”之后，进一步描述它有多唯一、还剩多少不可见自由度。
\end{remark}

\subsubsection{截断 Zeckendorf 位的模分解与顶位统计}\label{subsubsec:fold-zeckendorf-mod-topbits}

在定义 \ref{def:fold-word} 的记号下，记
\[
v_m(x):=\sum_{k=1}^{m}x_kF_{k+1}\qquad(x\in X_m),
\qquad
M_m:=F_{m+2},
\qquad
T_m:=\sum_{k=1}^{m}F_{k+1}=F_{m+3}-2.
\]
并令 $Z(N(\omega))=(c_1,c_2,\dots)\in X_\infty$ 为 Zeckendorf 数字序列；为强调“顶位”含义，下文记
\[
(\zeta_m,\zeta_{m+1}):=(c_m,c_{m+1})\in\{(0,0),(1,0),(0,1)\}.
\]

\begin{theorem}[截断位的模分解与唯一性]\label{thm:fold-zeckendorf-mod-decomposition-unique}
对任意 $m\ge 1$ 与任意 $\omega\in\Omega_m$，存在且仅存在一对 $(x,c)\in X_m\times\{0,1\}$ 使
\[
N(\omega)=v_m(x)+c\,M_m,
\qquad 0\le v_m(x)<M_m,
\]
并且该 $x$ 必为 $\Fold_m(\omega)$。等价地，
\[
v_m(\Fold_m(\omega))\equiv N(\omega)\pmod{M_m},
\qquad
0\le v_m(\Fold_m(\omega))<M_m,
\qquad
c=\mathbf 1_{\{N(\omega)\ge M_m\}}.
\]
\leanverified{stableValue\_injective}
\leanverified{stableAdd\_value\_mod}
\leanverified{paper\_fold\_zeckendorf\_mod\_decomposition\_unique}
\end{theorem}

\begin{proof}
由 $0\le N(\omega)\le T_m=F_{m+3}-2<M_m+F_{m+1}\le 2M_m$，可知
$c:=\mathbf 1_{\{N(\omega)\ge M_m\}}\in\{0,1\}$ 唯一确定。令 $r:=N(\omega)-cM_m$，则 $0\le r<M_m$。
在区间 $[0,M_m)$ 上，Zeckendorf 表示不会使用权重 $F_{m+2}$，因此存在唯一 $x\in X_m$ 使 $v_m(x)=r$。
另一方面，$Z(N(\omega))$ 的前 $m$ 位正是 $\Fold_m(\omega)$；当且仅当 $c=c_{m+1}=1$ 时出现顶位 $F_{m+2}$，截断后余数即为 $r$，从而 $v_m(\Fold_m(\omega))=r$。唯一性由 $c$ 与 Zeckendorf 表示唯一性推出。
\end{proof}

\begin{proposition}[补位对称诱导的 Fourier 相位刚性]\label{prop:fold-fiber-fourier-phase-rigidity}
沿用命题 \ref{prop:fold-fiber-count-reciprocity} 的记号，并令
\[
\alpha_m:=F_{m+1}-2\in\ZZ/M_m\ZZ.
\]
将 \(f_m\) 视为加法群 \(\ZZ/M_m\ZZ\) 上的类函数，并定义其离散 Fourier 变换
\[
\widehat f_m(k):=\sum_{r\in\ZZ/M_m\ZZ} f_m(r)\exp\!\left(-\frac{2\pi i kr}{M_m}\right),
\qquad k\in\ZZ/M_m\ZZ.
\]
则对任意 \(k\in\ZZ/M_m\ZZ\) 有严格恒等式
\[
\boxed{\;
\widehat f_m(k)
=\exp\!\left(-\frac{2\pi i k\alpha_m}{M_m}\right)\overline{\widehat f_m(k)}\;}
\]
从而
\[
\boxed{\;
\exp\!\left(\frac{\pi i k\alpha_m}{M_m}\right)\widehat f_m(k)\in\RR\;}
\]
即每个 Fourier 系数的相位被锁定在一条固定的实直线子空间中。
\end{proposition}
\begin{proof}
由命题 \ref{prop:fold-fiber-count-reciprocity} 的反射对称性 \(f_m(r)=f_m(\alpha_m-r)\)，有
\[
\widehat f_m(k)
=\sum_{r} f_m(\alpha_m-r)\exp\!\left(-\frac{2\pi i kr}{M_m}\right)
=\exp\!\left(-\frac{2\pi i k\alpha_m}{M_m}\right)\sum_{r} f_m(r)\exp\!\left(\frac{2\pi i kr}{M_m}\right).
\]
最后一项为 \(\overline{\widehat f_m(k)}\)，故得第一式；第二式由移项并乘以半相位因子得到。证毕。
\end{proof}
\leanverified{paper\_fold\_fiber\_fourier\_phase\_rigidity}

\begin{corollary}[模余类刚性与有限分辨率谓词的周期性]\label{cor:fold-mod-rigidity-periodic-predicates}
固定 $m\ge 1$，定义剩余映射
\[
r_m:\Omega_m\to \ZZ/M_m\ZZ,\qquad r_m(\omega):=N(\omega)\bmod M_m.
\]
令 \(s_m:\ZZ/M_m\ZZ\to X_m\) 为如下规范截面：对任意剩余类 \(\bar r\)，取其唯一代表元 \(r\in\{0,1,\dots,M_m-1\}\)，并令 \(s_m(\bar r)\) 为满足 \(v_m(x)=r\) 的唯一 \(x\in X_m\)。
则
\[
\Fold_m\ =\ s_m\circ r_m.
\]
因此对任意集合 \(Y\) 与任意映射 \(g:X_m\to Y\)，复合 \(g\circ \Fold_m\) 仅依赖于 \(r_m(\omega)\)。
特别地，对任意谓词 \(P\subseteq X_m\)，存在余类子集 \(R_P\subseteq \ZZ/M_m\ZZ\) 使得
\[
\Fold_m^{-1}(P)=r_m^{-1}(R_P).
\]
从而把 \(\ind_{\{\Fold_m(\omega)\in P\}}\) 视为 \(N(\omega)\) 的函数时，其必为周期 \(M_m\) 的周期函数。
\end{corollary}

\begin{proof}
由定理 \ref{thm:fold-zeckendorf-mod-decomposition-unique}，对任意 \(\omega\in\Omega_m\) 有
\(
v_m(\Fold_m(\omega))\equiv N(\omega)\ (\mathrm{mod}\ M_m)
\)
且 \(v_m(\Fold_m(\omega))\in\{0,\dots,M_m-1\}\)。
而 \(s_m\) 恰按该余数选择唯一的 \(X_m\) 元素，因此 \(\Fold_m(\omega)=s_m(r_m(\omega))\)。
其余断言由因子化直接推出。证毕。
\leanverified{fold\_eq\_of\_weight\_mod\_eq}
\leanverified{paper\_fold\_mod\_rigidity\_periodic\_predicates}
\end{proof}

\begin{proposition}[端点 $M_m-1$ 的子集和表示唯一性]\label{prop:fold-endpoint-Mm-minus-one-unique}
对任意 $m\ge 2$，方程
\[
\sum_{k=1}^{m}\omega_kF_{k+1}=M_m-1=F_{m+2}-1,\qquad \omega\in\Omega_m
\]
恰有一个解，并由交错模式给出
\[
(\omega_m,\omega_{m-1},\dots,\omega_1)=
\begin{cases}
(1,0,1,0,\dots,1,0), & m\ \text{为偶数},\\
(1,0,1,0,\dots,0,1), & m\ \text{为奇数}.
\end{cases}
\]
等价地，
\[
F_{m+2}-1=F_{m+1}+F_{m-1}+F_{m-3}+\cdots.
\]
\end{proposition}

\begin{proof}
设 $S=\sum_{k=1}^{m}\omega_kF_{k+1}=F_{m+2}-1$。注意 $\sum_{k=1}^{m-1}F_{k+1}=F_{m+2}-2$，故若 $\omega_m=0$ 则 $S\le F_{m+2}-2$ 矛盾，因而 $\omega_m=1$。
于是
\[
\sum_{k=1}^{m-1}\omega_kF_{k+1}=S-F_{m+1}=F_m-1.
\]
同理 $\sum_{k=1}^{m-2}F_{k+1}=F_m-2$，故 $\omega_{m-1}=1$ 不可能，从而 $\omega_{m-1}=0$。继续对剩余等式递归同样论证，得到 $\omega_{m-2}=1,\omega_{m-3}=0,\dots$ 的强制交错结构，因而解唯一，并给出所述展开。
\leanverified{alternating\_fib\_sum\_even}
\leanverified{alternating\_fib\_sum\_odd}
\leanverified{paper\_fold\_endpoint\_seeds}
\leanverified{paper\_fold\_endpoint\_Mm\_minus\_one\_unique\_seeds}
\leanverified{paper\_fold\_endpoint\_Mm\_minus\_one\_unique\_package}
\leanverified{paper\_fold\_endpoint\_Mm\_minus\_one\_unique}
\leanverified{paper\_fold\_endpoint\_mm\_minus\_one\_unique}
\end{proof}

\begin{theorem}[顶两位 Zeckendorf 统计的三分律与有限修正]\label{thm:fold-top-two-zeckendorf-trisect}
在 $\Omega_m$ 上取均匀测度，并令
\[
A_m:=\#\{\omega\in\Omega_m:\ N(\omega)\ge M_m\}.
\]
则对一切 $m\ge 2$，
\[
\#\{(\zeta_m,\zeta_{m+1})=(0,1)\}=A_m,\qquad
\#\{(\zeta_m,\zeta_{m+1})=(1,0)\}=2^m-2A_m-1,\qquad
\#\{(\zeta_m,\zeta_{m+1})=(0,0)\}=A_m+1,
\]
且
\[
A_m=
\begin{cases}
\dfrac{4^{m/2}-1}{3}, & m\ \text{为偶数},\\[6pt]
\dfrac{2\bigl(4^{(m-1)/2}-1\bigr)}{3}, & m\ \text{为奇数}.
\end{cases}
\]
因此极限分布满足
\[
\lim_{m\to\infty}\PP((\zeta_m,\zeta_{m+1})=(0,0))
=\lim_{m\to\infty}\PP((\zeta_m,\zeta_{m+1})=(1,0))
=\lim_{m\to\infty}\PP((\zeta_m,\zeta_{m+1})=(0,1))=\frac13,
\]
并且有限尺寸修正为
\[
\PP((\zeta_m,\zeta_{m+1})=(0,1))=
\begin{cases}
\dfrac{1-2^{-m}}{3}, & m\ \text{为偶数},\\[6pt]
\dfrac{1-2^{-(m-1)}}{3}, & m\ \text{为奇数},
\end{cases}
\]
\[
\PP((\zeta_m,\zeta_{m+1})=(1,0))=
\begin{cases}
\dfrac{1-2^{-m}}{3}, & m\ \text{为偶数},\\[6pt]
\dfrac{1+2^{-m}}{3}, & m\ \text{为奇数},
\end{cases}
\]
\[
\PP((\zeta_m,\zeta_{m+1})=(0,0))=
\begin{cases}
\dfrac{1+2^{-(m-1)}}{3}, & m\ \text{为偶数},\\[6pt]
\dfrac{1+2^{-m}}{3}, & m\ \text{为奇数}.
\end{cases}
\]
\leanverified{hiddenBitCount\_even\_closed}
\leanverified{hiddenBitCount\_odd\_closed}
\leanverified{complement\_hiddenBitCount}
\leanverified{hiddenBitCount\_lt\_pow}
\leanverified{hiddenBitCount\_pos}
\leanverified{paper\_fold\_top\_two\_zeckendorf\_trisect}
\end{theorem}

\begin{proof}
先给出 $A_m$ 的递推。若 $\omega_m=0$，则 $N(\omega)\le \sum_{k=1}^{m-1}F_{k+1}=F_{m+2}-2<M_m$，故事件 $\{N(\omega)\ge M_m\}$ 蕴含 $\omega_m=1$。于是
\[
A_m=\#\Bigl\{\omega'\in\{0,1\}^{m-1}:\ \sum_{k=1}^{m-1}\omega'_kF_{k+1}\ge F_m\Bigr\}.
\]
取补集并对 $\{0,1\}^{m-1}$ 施行逐位取反对合 $\varepsilon\mapsto \bar\varepsilon$，由
$\sum_{k=1}^{m-1}F_{k+1}=F_{m+2}-2$ 得
\[
A_m=\#\Bigl\{\varepsilon\in\{0,1\}^{m-1}:\ \sum_{k=1}^{m-1}\varepsilon_kF_{k+1}\le F_{m+1}-2\Bigr\}.
\]
按 $\varepsilon_{m-1}$ 分解：若 $\varepsilon_{m-1}=0$，则前 $m-2$ 位任意皆满足不等式（因 $\sum_{k=1}^{m-2}F_{k+1}=F_{m+1}-2$），贡献 $2^{m-2}$；
若 $\varepsilon_{m-1}=1$，则剩余需满足
\[
\sum_{k=1}^{m-2}\varepsilon_kF_{k+1}\le (F_{m+1}-2)-F_m=F_{m-1}-2,
\]
其计数恰为 $A_{m-2}$。因此得到递推
\[
A_m=2^{m-2}+A_{m-2}\qquad(m\ge 4),
\]
并由直接核算 $A_2=1,A_3=2$ 解得所示闭式。

其次把 $A_m$ 与顶两位联系。由 Zeckendorf 约束 $c_kc_{k+1}=0$，$(\zeta_m,\zeta_{m+1})$ 仅有三态，并且按权重区间划分有
\[
(\zeta_m,\zeta_{m+1})=(0,1)\iff N(\omega)\ge F_{m+2}=M_m,
\]
\[
(\zeta_m,\zeta_{m+1})=(1,0)\iff F_{m+1}\le N(\omega)<F_{m+2}=M_m,
\qquad
(\zeta_m,\zeta_{m+1})=(0,0)\iff N(\omega)<F_{m+1}.
\]
第一式给出 $\#\{(\zeta_m,\zeta_{m+1})=(0,1)\}=A_m$。由逐位取反对合 $\omega\mapsto \bar\omega$ 满足
$N(\bar\omega)=T_m-N(\omega)$，可得
\[
\#\{N(\omega)<F_{m+1}\}=\#\{N(\omega)\ge T_m-(F_{m+1}-1)+1\}=\#\{N(\omega)\ge F_{m+2}-1\}.
\]
右端等于 $A_m+\#\{N(\omega)=F_{m+2}-1\}$，而 $\#\{N(\omega)=F_{m+2}-1\}=1$ 由命题 \ref{prop:fold-endpoint-Mm-minus-one-unique} 给出，从而
\(
\#\{(\zeta_m,\zeta_{m+1})=(0,0)\}=A_m+1
\)。
再由总数 $2^m$ 得到
\[
\#\{(\zeta_m,\zeta_{m+1})=(1,0)\}=2^m-\bigl(A_m+(A_m+1)\bigr)=2^m-2A_m-1.
\]
将 $A_m$ 的闭式代入即可得到两条概率修正公式。
\end{proof}

\begin{corollary}[折叠溢出率与缺口期望的闭式]\label{cor:fold-overflow-gap-rate-expectation}
在 $\Omega_m$ 的均匀测度下，令
\[
\Delta_m(\omega):=N(\omega)-v_m(\Fold_m(\omega))\in\{0,M_m\}
\]
为折叠缺口，则对一切 $m\ge 2$，
\[
\PP(\Delta_m=M_m)=\frac{A_m}{2^m}=
\begin{cases}
\dfrac{1-2^{-m}}{3}, & m\ \text{为偶数},\\[6pt]
\dfrac{1-2^{-(m-1)}}{3}, & m\ \text{为奇数},
\end{cases}
\qquad
\EE[\Delta_m]=M_m\,\PP(\Delta_m=M_m).
\]
并且
\[
\PP\bigl(v_m(\Fold_m(\omega))<F_{m+1}\bigr)=
\begin{cases}
\dfrac{1+2^{-(m-1)}}{3}, & m\ \text{为偶数},\\[6pt]
\dfrac{1+2^{-m}}{3}, & m\ \text{为奇数}.
\end{cases}
\]
\[
\PP\bigl(F_{m+1}\le N(\omega)<M_m\bigr)=
\begin{cases}
\dfrac{1-2^{-m}}{3}, & m\ \text{为偶数},\\[6pt]
\dfrac{1+2^{-m}}{3}, & m\ \text{为奇数}.
\end{cases}
\]
\end{corollary}

\begin{proof}
由定理 \ref{thm:fold-zeckendorf-mod-decomposition-unique}，$\Delta_m=cM_m$ 且 $\Delta_m=M_m\iff N(\omega)\ge M_m$，故 $\PP(\Delta_m=M_m)=A_m/2^m$。代入定理 \ref{thm:fold-top-two-zeckendorf-trisect} 的 $A_m$ 化简得到所示闭式，期望等式由 $\Delta_m\in\{0,M_m\}$ 立得。第二个概率等式等价于事件 $(\zeta_m,\zeta_{m+1})=(0,0)$，第三个概率等式等价于事件 $(\zeta_m,\zeta_{m+1})=(1,0)$，均由定理 \ref{thm:fold-top-two-zeckendorf-trisect} 给出。
\leanverified{three\_mul\_hiddenBitCount\_add}
\leanverified{paper\_fold\_overflow\_gap\_rate\_expectation}
\end{proof}

\begin{proposition}[零剩余纤维的线性退化]\label{prop:fold-zero-fiber-linear}
令
\[
f_m(0):=\#\{\omega\in\Omega_m:\ \Fold_m(\omega)=0^m\},
\]
则对任意 $m\ge 2$，
\[
f_m(0)=1+\Bigl\lfloor\frac{m}{2}\Bigr\rfloor.
\]
\end{proposition}

\begin{proof}
由定理 \ref{thm:fold-zeckendorf-mod-decomposition-unique}，$\Fold_m(\omega)=0^m$ 当且仅当 $v_m(\Fold_m(\omega))=0$，等价于 $N(\omega)\equiv 0\pmod{M_m}$。
又因 $0\le N(\omega)<2M_m$，上述同余等价于 $N(\omega)\in\{0,M_m\}$。其中 $N(\omega)=0$ 仅由全零 $\omega=0^m$ 给出，贡献 $1$。
剩余需计数
\[
B_m:=\#\{\omega\in\Omega_m:\ N(\omega)=M_m=F_{m+2}\}.
\]
若 $\omega_m=0$ 则 $N(\omega)\le F_{m+2}-2$ 矛盾，故必有 $\omega_m=1$，从而
\[
\sum_{k=1}^{m-1}\omega_kF_{k+1}=F_{m+2}-F_{m+1}=F_m.
\]
在右侧表示中，若 $\omega_{m-1}=1$，则左侧含项 $F_m$，其余位必须为 $0$，得到一种表示；若 $\omega_{m-1}=0$，则需用 $\{F_2,\dots,F_{m-1}\}$ 表示 $F_m$，这与参数 $m-2$ 的同类计数等价，故
\[
B_m=1+B_{m-2}\qquad(m\ge 4),
\]
并由 $B_2=1,B_3=1$ 解得 $B_m=\lfloor m/2\rfloor$。因此 $f_m(0)=1+B_m$。
\leanverified{paper\_fold\_zero\_fiber\_linear}
\leanverified{paper\_fold\_zero\_fiber\_linear\_extended}
\leanverified{weightSumAtMm}
\leanverified{weightSumAtMm\_two}
\leanverified{weightSumAtMm\_three}
\leanverified{weightSumAtMm\_four}
\leanverified{weightSumAtMm\_five}
\leanverified{weightSumAtMm\_eq\_div\_two\_small}
\end{proof}

\begin{proposition}[纤维计数的互易律与反射对称性]\label{prop:fold-fiber-count-reciprocity}
\leanverified{paper\_fold\_fiber\_count\_reciprocity}
对任意 $m\ge 2$ 与任意 $r\in\{0,1,\dots,M_m-1\}$，令
\[
f_m(r):=\#\{\omega\in\Omega_m:\ v_m(\Fold_m(\omega))=r\},
\qquad
a_m(s):=\#\{\omega\in\Omega_m:\ N(\omega)=s\}\quad(0\le s\le T_m),
\]
则成立互易恒等式
\[
f_m(r)=a_m(r)+a_m(M_m+r),
\]
并且
\[
f_m(r)=a_m(r)+a_m(F_{m+1}-2-r),
\qquad
f_m(r)=f_m(F_{m+1}-2-r).
\]
\leanpartial{weightCongruenceCount\_complement}{条件 m≥2，r≤F\_{m+1}-2；互易恒等式 f\_m(r)=a\_m(r)+a\_m(M\_m+r) 未形式化}
\leanverified{fiberMultiplicity\_complement}
\leanverified{stableValue\_Fold\_add\_complement}
\leanverified{complementAction}
\leanverified{complementAction\_involutive}
\leanverified{fiberMultiplicity\_complementAction}
\leanverified{fiberMultiplicity\_complement\_parity}
\leanverified{popcount\_complement}
\leanverified{paper\_pom\_fiber\_reciprocity\_package}
\leanverified{paper\_pom\_fiber\_histogram\_palindrome}
\leanverified{paper\_folding\_weight\_count\_palindrome}
\end{proposition}

% --- Distillation writeback ---
\begin{definition}[fold-aware 跨分辨率限制]\label{distill:wang_zahl-sticky_reduction_and_scale_saturated_closure-fold-aware-cross-resolution}
\textbf{结果 17；日期：2026-04-23T03:31:39+08:00。依赖状态：construction.}
设 \(m\ge h\ge 1\)，并令 \(Z(N(\omega))=(c_1,c_2,\ldots)\) 为 \(N(\omega)\) 的 Zeckendorf 数字序列。定义 fold-aware 跨分辨率限制
\[
\Fold^{\rm aw}_{m\to h}:\Omega_m\longrightarrow X_h,
\qquad
\Fold^{\rm aw}_{m\to h}(\omega):=(c_1,\ldots,c_h).
\]
同时记稳定前缀投影为
\[
\pr_{m,h}:X_m\longrightarrow X_h,
\qquad
\pr_{m,h}(x_1,\ldots,x_m):=(x_1,\ldots,x_h).
\]
该对象不是把原始输入字 \(\omega\) 截成前 \(h\) 位后再施行 \(\Fold_h\)；它先在同一个整数 \(N(\omega)\) 上取 Zeckendorf 稳定展开，再取稳定前缀。因此它避免了把朴素输入截断误认为跨分辨率不变量的幻觉。
\end{definition}

\begin{proposition}[fold-aware 限制与稳定前缀交换]\label{distill:wang_zahl-sticky_reduction_and_scale_saturated_closure-fold-aware-prefix-commutation}
\textbf{结果 18；日期：2026-04-23T03:31:39+08:00。依赖状态：rigidity.}
对任意 \(m\ge h\ge 1\)，有严格交换关系
\[
\Fold^{\rm aw}_{m\to h}=\pr_{m,h}\circ\Fold_m.
\]
因此，若 \(\gamma_h:X_h\to Y_h\) 是任意 coarse carrier 映射，则
\[
\gamma_h\circ\Fold^{\rm aw}_{m\to h}
=
(\gamma_h\circ\pr_{m,h})\circ\Fold_m.
\]
特别地，所有跨分辨率 sticky/non-sticky 测试若经由 \(\gamma_h\) 读出，都可完全下降到稳定代表 \(\Fold_m(\omega)\) 的前缀，而不依赖原始微态字的朴素前缀。
\end{proposition}

% --- Distillation writeback ---
\begin{proposition}[投影 fold congruence 的尾部 lift 审计]\label{prop:distill-euclid-elements-fold-prefix-tail-lift-audit}
\textbf{Dependency status: rigidity.}
设 \(m\ge h\ge1\)。沿用定义 \ref{distill:wang_zahl-sticky_reduction_and_scale_saturated_closure-fold-aware-cross-resolution} 的
\(\Fold^{\rm aw}_{m\to h}\) 与稳定前缀投影 \(\pr_{m,h}\)。再定义尾部读出
\[
\operatorname{Tail}_{m,h}:X_m\longrightarrow \{0,1\}^{m-h},
\qquad
\operatorname{Tail}_{m,h}(x_1,\ldots,x_m):=(x_{h+1},\ldots,x_m),
\]
其中 \(h=m\) 时右端理解为单点集合。则对任意 \(\omega,\omega'\in\Omega_m\)，有
\[
\Fold_m(\omega)=\Fold_m(\omega')
\]
当且仅当同时满足
\[
\Fold^{\rm aw}_{m\to h}(\omega)=\Fold^{\rm aw}_{m\to h}(\omega')
\]
与
\[
\operatorname{Tail}_{m,h}(\Fold_m(\omega))
=
\operatorname{Tail}_{m,h}(\Fold_m(\omega')).
\]
因此，深度 \(h\) 的 projected congruence 只证明稳定前缀相等；它成为深度 \(m\) 的 lifted equality，当且仅当尾部 residual 同时消失。
\end{proposition}

% --- Distillation writeback ---
% Distillation timestamp: 2026-04-23T12:04:19+08:00; pipeline index: 35.
\begin{theorem}[投影折叠同余的最小尾部严格化]\label{distill:euclid_elements-lifted_congruence_and_residual_rigidity_families-fold-projected-congruence-minimal-tail-strictification}
\textbf{Dependency status: rigidity.}
设 \(m\ge h\ge 1\)，并取非空有限族 \(E\subseteq\Omega_m\)。沿用
\(\Fold^{\rm aw}_{m\to h}\)、\(\pr_{m,h}\) 与
\(\operatorname{Tail}_{m,h}\) 的记号。对每个 \(p\in X_h\)，定义投影纤维中的尾部残余集合
\[
\mathcal T_E(p)
:=
\left\{
\operatorname{Tail}_{m,h}\bigl(\Fold_m(\omega)\bigr):
\omega\in E,\ \Fold^{\rm aw}_{m\to h}(\omega)=p
\right\},
\]
并令
\[
\mu_{m,h}(E):=\max_{p\in X_h}\bigl|\mathcal T_E(p)\bigr|.
\]
称一个有限寄存器 \(\rho:E\to R\) 为该投影同余的尾部 lift 审计，若对任意
\(\omega,\omega'\in E\) 都有
\[
\Fold^{\rm aw}_{m\to h}(\omega)=\Fold^{\rm aw}_{m\to h}(\omega'),\qquad
\rho(\omega)=\rho(\omega')
\quad\Longrightarrow\quad
\Fold_m(\omega)=\Fold_m(\omega').
\]
则：
\begin{enumerate}
  \item 任意尾部 lift 审计都满足
  \[
  |R|\ge \mu_{m,h}(E),
  \]
  且存在达到等号的尾部 lift 审计；
  \item 深度 \(h\) 的 projected congruence 在 \(E\) 上已经推出深度 \(m\) 的 lifted equality，当且仅当
  \[
  \mu_{m,h}(E)=1;
  \]
  \item 若 \(\mu_{m,h}(E)>1\)，则任意使该投影同余失败的坏对
  \[
  \Fold^{\rm aw}_{m\to h}(\omega)=\Fold^{\rm aw}_{m\to h}(\omega'),\qquad
  \Fold_m(\omega)\ne \Fold_m(\omega')
  \]
  必且只必携带非零尾部 residual：
  \[
  \operatorname{Tail}_{m,h}\bigl(\Fold_m(\omega)\bigr)
e
  \operatorname{Tail}_{m,h}\bigl(\Fold_m(\omega')\bigr).
  \]
  因而任何仅经由 \(\Fold^{\rm aw}_{m\to h}\) 因子化的 coarse readout 都不能区分这样的坏对。
\end{enumerate}
\end{theorem}
\begin{proof}
固定 \(p\in X_h\)。若 \(\tau,\tau'\in\mathcal T_E(p)\) 且 \(\tau\ne\tau'\)，则可取
\(\omega,\omega'\in E\) 使
\[
\Fold^{\rm aw}_{m\to h}(\omega)=
\Fold^{\rm aw}_{m\to h}(\omega')=p,
\]
且两者尾部读出分别为 \(\tau,\tau'\)。由命题
\ref{distill:wang_zahl-sticky_reduction_and_scale_saturated_closure-fold-aware-prefix-commutation}，投影同余正是稳定前缀相等；再由命题
\ref{prop:distill-euclid-elements-fold-prefix-tail-lift-audit}，在同一前缀下，深度 \(m\) 的折叠值相等当且仅当尾部也相等。因此上述两点的 lifted equality 失败。若 \(\rho\) 是 lift 审计，则它必须给这些不同尾部赋予不同寄存器值。于是
\(|R|\ge |\mathcal T_E(p)|\)。对所有 \(p\) 取最大值得到
\(|R|\ge\mu_{m,h}(E)\)。

反过来，取一个集合 \(R\) 满足 \(|R|=\mu_{m,h}(E)\)。对每个 \(p\) 任选单射
\[
\iota_p:\mathcal T_E(p)\hookrightarrow R,
\]
并定义
\[
\rho(\omega):=\iota_{\Fold^{\rm aw}_{m\to h}(\omega)}
\left(
\operatorname{Tail}_{m,h}(\Fold_m(\omega))
\right).
\]
若两点具有相同 projected congruence 且 \(\rho\) 值相同，则它们位于同一个 \(p\)-纤维内，并因 \(\iota_p\) 单射而具有相同尾部。命题
\ref{prop:distill-euclid-elements-fold-prefix-tail-lift-audit} 遂给出 \(\Fold_m(\omega)=\Fold_m(\omega')\)。故等号可达。

第二条是第一条的退化情形。若 \(\mu_{m,h}(E)=1\)，则每个投影前缀下只有一个尾部 residual，故 projected congruence 自动给出前缀与尾部同时相等，从而给出 lifted equality。若 \(\mu_{m,h}(E)>1\)，则存在同一 \(p\) 下两个不同尾部，取相应代表即得到 projected congruent 但不 lifted equal 的坏对。

第三条再次由 \(X_m\subseteq\{0,1\}^m\) 中一个词由其前缀与尾部唯一决定得到：在前缀已经相等的前提下，折叠值不等价于尾部不等。若 \(g:E\to Y\) 仅经由 \(\Fold^{\rm aw}_{m\to h}\) 因子化，则存在 \(\bar g:X_h\to Y\) 使
\(g=\bar g\circ\Fold^{\rm aw}_{m\to h}\)，故上述坏对必有相同 \(g\)-值。这正是投影同余不能无审计转移为规范相等的尾部 residual 机制。证毕。
\end{proof}

% --- End distillation writeback ---

\begin{proof}
由命题 \ref{distill:wang_zahl-sticky_reduction_and_scale_saturated_closure-fold-aware-prefix-commutation}，
\[
\Fold^{\rm aw}_{m\to h}=\pr_{m,h}\circ\Fold_m.
\]
于是 projected congruence 等价于
\[
\pr_{m,h}(\Fold_m(\omega))=\pr_{m,h}(\Fold_m(\omega')).
\]
另一方面，任意 \(x\in X_m\subseteq\{0,1\}^m\) 由其前缀 \(\pr_{m,h}(x)\) 与尾部 \(\operatorname{Tail}_{m,h}(x)\) 唯一决定。因此两个折叠值 \(\Fold_m(\omega)\) 与 \(\Fold_m(\omega')\) 相等，当且仅当前缀和尾部均相等。证毕。
\end{proof}

% --- End distillation writeback ---

\begin{proof}
由定义，\(\Fold_m(\omega)\) 正是 \(Z(N(\omega))\) 的前 \(m\) 个 Zeckendorf 稳定位，即
\[
\Fold_m(\omega)=(c_1,\ldots,c_m).
\]
对其施加稳定前缀投影 \(\pr_{m,h}\)，得到
\[
\pr_{m,h}(\Fold_m(\omega))=(c_1,\ldots,c_h).
\]
右端按定义正是 \(\Fold^{\rm aw}_{m\to h}(\omega)\)。故第一式成立。第二式仅是在第一式两端复合 \(\gamma_h\)。最后一句说明：若 carrier 读出为 \(\gamma_h\)，则其跨尺度值已经由稳定前缀决定；这里未使用、也不需要使用原始输入字的前 \(h\) 位，因此不会把 \(\Fold_h(\omega_1,\ldots,\omega_h)\) 误作跨分辨率不变量。证毕。
\end{proof}

% --- End distillation writeback ---

\begin{proof}
由定理 \ref{thm:fold-zeckendorf-mod-decomposition-unique}，$v_m(\Fold_m(\omega))$ 是 $N(\omega)$ 模 $M_m$ 的规范代表，且 $0\le N(\omega)<2M_m$，故事件 $\{v_m(\Fold_m(\omega))=r\}$ 等价于 $\{N(\omega)\in\{r,M_m+r\}\}$，从而得到第一式。
对合 $\omega\mapsto\bar\omega$ 满足 $N(\bar\omega)=T_m-N(\omega)$，故 $a_m(s)=a_m(T_m-s)$。注意
\[
T_m-(M_m+r)=(F_{m+3}-2)-(F_{m+2}+r)=F_{m+1}-2-r,
\]
代入 $a_m(M_m+r)=a_m(F_{m+1}-2-r)$ 即得第二式。将第二式在 $r$ 与 $F_{m+1}-2-r$ 间交换即可推出反射对称性。
\end{proof}


\subsubsection{折叠的局部重写实现：终止、合流与顺序无关}\label{subsec:fold-local-rewrite}
定义 \ref{def:fold-word} 以“先取值、再取 Zeckendorf 规范形”给出 $\Fold_m$。为将其落实为\textbf{局部、可执行}的归一化过程，并澄清“归约顺序不影响正规形”，我们把输入窗口嵌入到允许中间态出现冗余数字的配置空间上实施值保持重写。

这一实现顺序也很适合主论文的整体叙述：先说明对象可由哪些局部许可重写得到，再说明为何不同归约顺序会坍缩到同一正规形。前者是 constructibility，后者是 rigidity；两层若不拆开，折叠就会被误读成“定义一个商映射”，而不是“通过合法局部作图获得规范对象”。

\paragraph{配置空间与值函数}
令
\[
\mathcal{D}:=\Bigl\{a=(a_1,a_2,\dots): a_k\in\NN,\ \text{仅有限个 }a_k\neq 0\Bigr\}
\]
为有限支撑的非负整数数字配置。定义值同态
\[
\mathrm{Val}(a):=\sum_{k\ge 1} a_k F_{k+1}.
\]
对给定 $\omega\in\Omega_m$，记其嵌入为
\[
\iota_m(\omega):=(\omega_1,\dots,\omega_m,0,0,\dots)\in\mathcal{D},
\]
则 $\mathrm{Val}(\iota_m(\omega))=N(\omega)$。

\paragraph{局部重写规则（值保持）}
对 $a\in\mathcal{D}$，允许在满足条件时对局部模式做如下替换（其中 $e_k$ 为第 $k$ 位单位向量）：
\begin{enumerate}
  \renewcommand{\theenumi}{R\arabic{enumi}}
  \renewcommand{\labelenumi}{(\theenumi)}
  \item \label{rule:fold:adj}\textbf{相邻合并：}若 $a_k\ge 1$ 且 $a_{k+1}\ge 1$，则
  \[
  a\ \to\ a-e_k-e_{k+1}+e_{k+2}.
  \]
  \item \label{rule:fold:dedup1}\textbf{底位去重：}若 $a_1\ge 2$，则
  \[
  a\ \to\ a-2e_1+e_2.
  \]
  \item \label{rule:fold:dedup2}\textbf{第二位去重：}若 $a_2\ge 2$，则
  \[
  a\ \to\ a-2e_2+e_1+e_3.
  \]
  \item \label{rule:fold:dedupk}\textbf{一般位去重（$k\ge 3$）：}若 $a_k\ge 2$，则
  \[
  a\ \to\ a-2e_k+e_{k-2}+e_{k+1}.
  \]
\end{enumerate}
每条规则都只改动有限个坐标，因此是严格的\textbf{局部}重写；并且它们分别对应恒等式
\(
F_{k+1}+F_{k+2}=F_{k+3},
2F_2=F_3,
2F_3=F_2+F_4,
2F_{k+1}=F_{k-1}+F_{k+2}
\)，因此每一步都保持 $\mathrm{Val}$ 不变。

\paragraph{不可约项与 Zeckendorf 规范形}
记 $X_\infty=\{(c_1,c_2,\dots)\in\{0,1\}^{\mathbb{N}}:\ c_kc_{k+1}=0\}$。上述规则恰在出现相邻 $11$ 或某位数字 $\ge 2$ 时可用，因此不可约项当且仅当属于 $X_\infty$。

\begin{proposition}[分辨率稳定类型的逆极限本体]\label{prop:fold-inverse-limit-xm-xinfty}
对每个 \(m\ge 1\) 令 \(r_{m+1\to m}:X_{m+1}\to X_m\) 为前缀截断 \(r_{m+1\to m}(a_1\cdots a_{m+1})=a_1\cdots a_m\)。
则 \((X_m,r_{m+1\to m})\) 构成逆系统，并存在自然同构
\[
\varprojlim_m X_m\ \cong\ X_\infty,
\]
其中 \(X_\infty\to X_m\) 的结构映射为取前 \(m\) 位前缀。
此外，对任意 \(m\ge 1\)，Zeckendorf 截断映射
\(
N\mapsto \pi_m(Z(N))
\)
在区间 \(\{0,1,\dots,F_{m+2}-1\}\) 上与 \(X_m\) 构成双射，其逆映射为 \(w\mapsto V_m(w)\)（\S\ref{subsec:emergent-arithmetic-zeckendorf-value-address}；并见引理 \ref{lem:zeckendorf-range-bijection}）。
\leanverified{paper\_fold\_inverse\_limit\_xm\_xinfty}
\leanverified{prefixWord\_stableValue\_coherent}
\leanverified{prefixWord\_surjective}
\end{proposition}
\begin{proof}
由于“无相邻 \(11\)”对前缀封闭，\(r_{m+1\to m}\) 良定义，从而得到逆系统。
给定相容族 \((w_m)_{m\ge 1}\in\varprojlim_m X_m\)，对每个坐标 \(i\ge 1\) 令 \(a_i\) 为任意 \(m\ge i\) 时 \(w_m\) 的第 \(i\) 位。
相容性保证该定义与 \(m\) 的选择无关，因此得到唯一的无限词 \(a=(a_i)_{i\ge 1}\in\{0,1\}^{\NN}\)；
又因每个 \(w_m\) 无相邻 \(11\)，故 \(a\in X_\infty\)。
反之，对任意 \(a\in X_\infty\)，其前缀 \(\pi_m(a)\in X_m\) 形成相容族 \((\pi_m(a))_m\in\varprojlim X_m\)。
两构造互为逆，得到同构。
最后一条双射性由 \(V_m\) 的范围双射与 Zeckendorf 唯一性给出（引理 \ref{lem:zeckendorf-range-bijection}）。证毕。
\end{proof}

\begin{proposition}[一致纤维集的前缀覆盖数与外置比特率上界]\label{prop:fold-consistent-fiber-hausdorff-prefix-count}
令 \(\Omega_\infty:=\{0,1\}^{\NN}\)，并以 \(\tau_{\infty\to m}:\Omega_\infty\to\Omega_m\) 表示取前 \(m\) 位前缀。
给定 \(x_\infty\in X_\infty\cong\varprojlim X_m\)，记 \(x_m:=\pi_m(x_\infty)\in X_m\)。
定义与 \(x_\infty\) 相容的一致纤维集
\[
\mathcal F(x_\infty)
:=\Bigl\{\omega\in\Omega_\infty:\ \Fold_m\bigl(\tau_{\infty\to m}(\omega)\bigr)=x_m\ \text{对所有 }m\ge 1\Bigr\}.
\]
对每个 \(m\ge 1\) 令
\[
N_m(x_\infty)
:=\#\Bigl\{a\in\Omega_m:\ [a]\cap \mathcal F(x_\infty)\neq\varnothing\Bigr\}
=\#\Bigl\{\tau_{\infty\to m}(\omega):\ \omega\in\mathcal F(x_\infty)\Bigr\},
\]
其中 \([a]\subset\Omega_\infty\) 为柱集。
则有纯计数上界
\[
N_m(x_\infty)\ \le\ d_m(x_m):=\bigl|\Fold_m^{-1}(x_m)\bigr|.
\]
并且对任意 \(\omega\in\mathcal F(x_\infty)\) 都存在确定性解码使
\[
K\bigl(\tau_{\infty\to m}(\omega)\mid x_\infty,m\bigr)
\ \le\
\log_2 d_m(x_m)+O(\log m).
\]
因此
\[
\liminf_{m\to\infty}\frac{K(\tau_{\infty\to m}(\omega)\mid x_\infty,m)}{m}
\ \le\
\liminf_{m\to\infty}\frac{\log_2 d_m(x_m)}{m}.
\]
\leanverified{paper\_fold\_consistent\_fiber\_hausdorff\_prefix\_count}
\end{proposition}
\begin{proof}
由定义，若 \([a]\cap\mathcal F(x_\infty)\neq\varnothing\)，则存在 \(\omega\in\mathcal F(x_\infty)\) 使 \(\tau_{\infty\to m}(\omega)=a\)，并满足 \(\Fold_m(a)=x_m\)；
故所有可达前缀集合包含于纤维 \(\Fold_m^{-1}(x_m)\)，得到 \(N_m(x_\infty)\le d_m(x_m)\)。
\par
对复杂度上界：给定 \((x_\infty,m)\) 可恢复 \(x_m=\pi_m(x_\infty)\)。
由 \(\Fold_m\) 可计算且纤维有限可枚举，可对集合 \(\Fold_m^{-1}(x_m)\) 取字典序枚举并用秩 \(i\in\{0,\dots,d_m(x_m)-1\}\) 编码元素；
该秩可用 \(\lceil\log_2 d_m(x_m)\rceil\) 比特描述，而解释器与枚举程序只引入 \(O(\log m)\) 的开销。
故得所示上界；最后一式由两边除以 \(m\) 并取 \(\liminf\) 得到。
证毕。
\end{proof}

\begin{proposition}[路径分量纤维的 rank/unrank 线性时间实现]\label{prop:fold-fiber-rank-unrank-linear}
对整数 \(\ell\ge 0\)，令
\[
\mathcal{L}_\ell:=\{s\in\{0,1\}^{\ell}:\ s\ \text{不含相邻子串 }11\},
\qquad
\card{\mathcal{L}_\ell}=F_{\ell+2}.
\]
存在确定性算法 \(\mathrm{rank}_\ell:\mathcal{L}_\ell\to[0,F_{\ell+2}-1]\) 与 \(\mathrm{unrank}_\ell:[0,F_{\ell+2}-1]\to\mathcal{L}_\ell\)，互为逆，并且二者均在 \(O(\ell)\) 时间内可计算。
进一步，若某纤维在某一规范双射下分解为有限直积
\[
\Fold_m^{-1}(x)\ \cong\ \prod_{i=1}^{r}\mathcal{L}_{\ell_i},
\qquad
d_m(x)=\prod_{i=1}^{r}F_{\ell_i+2},
\]
则可用混合进制分解把纤维内的 rank/unrank 归约为分量 rank/unrank，从而在 \(O(\sum_i\ell_i)=O(m)\) 时间内实现纤维内的 rank/unrank。
\end{proposition}
\begin{proof}
对 \(\mathcal{L}_\ell\) 的 rank/unrank 用 Fibonacci 递推实现。
对任意 \(s=s_1\cdots s_\ell\in\mathcal{L}_\ell\)，若 \(s_1=0\)，则其在字典序下与 \(\mathcal{L}_{\ell-1}\) 的尾串秩一致；若 \(s_1=1\)，则必有 \(s_2=0\)，并且以 \(0\) 起始的合法串共有 \(\card{\mathcal{L}_{\ell-1}}=F_{\ell+1}\) 个，故
\[
\mathrm{rank}_\ell(s)=F_{\ell+1}+\mathrm{rank}_{\ell-2}(s_3\cdots s_\ell).
\]
反秩同理：给定 \(j\in[0,F_{\ell+2}-1]\)，若 \(j<F_{\ell+1}\) 则取首位为 \(0\) 并递归于 \(\ell-1\)；否则取前两位为 \(10\) 并递归于 \(\ell-2\) 的余数 \(j-F_{\ell+1}\)。
每步长度减少 \(1\) 或 \(2\)，故时间为 \(O(\ell)\)。
\par
对直积情形，令 \(N_i:=F_{\ell_i+2}\)。
任意 \(j\in[0,\prod_i N_i-1]\) 唯一分解为混合进制坐标
\(
j=\sum_i j_i\prod_{k<i}N_k
\)
（其中 \(0\le j_i<N_i\)）。
对每个分量执行 \(\mathrm{unrank}_{\ell_i}(j_i)\) 并拼接即得全局 unrank；rank 由逐分量 rank 得到 \((j_i)\) 后再作同一混合进制合成。
总时间为 \(O(\sum_i\ell_i)=O(m)\)。证毕。
\end{proof}

\begin{proposition}[终止性、局部合流与合流（纽曼引理）]\label{prop:fold-rewrite-newman}
上述重写系统在 $\mathcal{D}$ 上强终止，且局部合流；由纽曼引理可得其合流性。因此每个 $a\in\mathcal{D}$ 的正规形存在且唯一，记为 $\mathrm{NF}(a)\in X_\infty$，并满足
\[
\mathrm{Val}(\mathrm{NF}(a))=\mathrm{Val}(a).
\]
\leanverified{step\_stronglyTerminating}
\leanverified{step\_confluent}
\leanverified{paper\_fold\_rewrite\_newman}
\end{proposition}

\begin{proof}
\textbf{（终止性）}令 $A(a)=\sum_{k\ge 1} a_k$、$B(a)=\sum_{k\ge 1} k\,a_k$、$C(a)=a_2$，并设 $\mu(a)=(A(a),B(a),C(a))$ 按字典序比较。则规则 (\ref{rule:fold:adj}) 与 (\ref{rule:fold:dedup1}) 使 $A$ 减少；规则 (\ref{rule:fold:dedupk})（$k\ge 3$）使 $B$ 减少；规则 (\ref{rule:fold:dedup2}) 保持 $A,B$ 不变但使 $C$ 减少。故每步重写都严格降低 $\mu$，系统强终止。

\textbf{（局部合流）}若 $a$ 有两条一步归约 $a\to a_1,a_2$，则两侧继续归约到不可约形后都落在 $X_\infty$，且由于每步保持 $\mathrm{Val}$，两不可约形同值；由 Zeckendorf 唯一性它们相同，于是两支可汇合。

\textbf{（合流）}由纽曼引理（强终止 + 局部合流 $\Rightarrow$ 合流）即得结论。
\end{proof}

\begin{corollary}[\texorpdfstring{$\Fold_m$}{Fold\_m} 的归约顺序无关实现]\label{cor:foldm-order-indep}
对任意 $\omega\in\Omega_m$，
\[
\Fold_m(\omega)\ =\ \pi_m\bigl(\mathrm{NF}(\iota_m(\omega))\bigr),
\]
且右端与归约时选取规则/位置的顺序无关。
\leanverified{normalPrefix\_iota\_eq\_Fold}
\leanverified{paper\_foldm\_order\_indep}
\end{corollary}

\begin{remark}[规范截面与可计算性]\label{rem:fold-normal-form-computable-section}
命题 \ref{prop:fold-rewrite-newman} 给出一个可计算的规范截面：对任意 $a\in\mathcal{D}$，正规形 $\mathrm{NF}(a)\in X_\infty$ 唯一存在且与归约顺序无关。因此，当“第三参数”仅对应于等价类中代表的选择（规范冗余）时，无需显式记录该参数；可通过求正规形在多项式时间内选出唯一代表（分辨率口径为线性，见命题 \ref{prop:fold-complexity}）。需要强调的是，该消去仅恢复稳定代表，不恢复同一纤维内的微态残差。
\end{remark}

\begin{theorem}[折叠的构造--刚性--障碍三分]\label{thm:fold-constructibility-rigidity-obstruction}
\leanverified{paper\_fold\_constructibility\_rigidity\_obstruction}
固定分辨率 \(m\ge 1\) 与稳定词 \(x\in X_m\)，记
\[
d_m(x):=\bigl|\Fold_m^{-1}(x)\bigr|.
\]
则折叠 \(\Fold_m\) 的对象论地位可分为三层：
\begin{enumerate}
  \item \textbf{构造层：}每个稳定词 \(x\) 都有一个显式 admissible witness，事实上 \(x\in\Omega_m\) 且 \(\Fold_m(x)=x\)。
  \item \textbf{刚性层：}若 \(a,b\in\mathcal D\) 满足 \(\mathrm{Val}(a)=\mathrm{Val}(b)\)，则
  \[
  \pi_m(\mathrm{NF}(a))=\pi_m(\mathrm{NF}(b)).
  \]
  换言之，分辨率 \(m\) 的稳定代表只依赖于值类，而不依赖于局部重写路径或初始冗余表示。
  \item \textbf{障碍层：}从 \(x\) 恢复唯一微态当且仅当 \(d_m(x)=1\)。
  若 \(d_m(x)>1\)，则以 \(\pi=\delta_x\) 代入定理 \ref{thm:fold-microstate-residual-window-reachability}，可实现的纤维残差恰覆盖闭区间
  \[
  [0,\log d_m(x)],
  \]
  等价地，可实现的纤维熵恰覆盖同一区间。
\end{enumerate}
\end{theorem}
\begin{proof}
\textbf{(1)} 由命题 \ref{prop:fold-basic} 的幂等性，任意 \(x\in X_m\subset\Omega_m\) 都满足 \(\Fold_m(x)=x\)，故稳定对象先以显式输入的形式被合法构造出来。

\textbf{(2)} 由命题 \ref{prop:fold-rewrite-newman}，正规形 \(\mathrm{NF}(a)\in X_\infty\) 唯一存在并仅由 \(\mathrm{Val}(a)\) 决定；故若 \(\mathrm{Val}(a)=\mathrm{Val}(b)\)，则 \(\mathrm{NF}(a)=\mathrm{NF}(b)\)，再取前 \(m\) 位即得
\(
\pi_m(\mathrm{NF}(a))=\pi_m(\mathrm{NF}(b))
\)。
这与推论 \ref{cor:foldm-order-indep} 一致地说明：折叠读出在值类上是刚性的。

\textbf{(3)} 若 \(d_m(x)=1\)，则纤维 \(\Fold_m^{-1}(x)\) 仅含一个元素，故微态恢复唯一。
反之若 \(d_m(x)>1\)，则唯一恢复显然失败。
将定理 \ref{thm:fold-microstate-residual-window-reachability} 中的宏观分布取为 \(\pi=\delta_x\)，便得
\[
\{\mathsf R(\mu):\ (\Fold_m)_{\#}\mu=\delta_x\}=[0,\log d_m(x)],
\]
且对应的熵谱为
\[
\{H(\mu):\ (\Fold_m)_{\#}\mu=\delta_x\}=[0,\log d_m(x)].
\]
因此当纤维非单点时，障碍不是“没有稳定对象”，而是存在一个不可由 \(x\) 本身消去的纤维残差窗口。
证毕。
\end{proof}

\begin{theorem}[纤维残差窗口的全区间可达性与熵谱完备]\label{thm:fold-microstate-residual-window-reachability}
设 \(\Omega\) 为有限集合，\(\Fold:\Omega\twoheadrightarrow X\) 为满射。
对每个 \(x\in X\) 记纤维 \(F_x:=\Fold^{-1}(x)\) 与多重度 \(d(x):=|F_x|\)。
给定任意宏观分布 \(\pi\in\Delta(X)\)，定义其\textbf{纤维均匀提升} \(\mu^{\mathrm{e}}_\pi\in\Delta(\Omega)\) 为
\[
\mu^{\mathrm{e}}_\pi(\omega):=\frac{\pi(x)}{d(x)},\qquad \omega\in F_x.
\]
并令
\[
\mathcal{C}_\pi:=\{\mu\in\Delta(\Omega):\ \Fold_{\#}\mu=\pi\}.
\]
定义相对于纤维均匀提升的\textbf{微态残差}
\[
\mathsf{R}(\mu):=D_{\mathrm{KL}}\!\left(\mu\ \middle\|\ \mu^{\mathrm{e}}_\pi\right),
\]
以及\textbf{可逆窗口长度}
\[
W_{\mathrm{inv}}(\pi):=\sum_{x\in X}\pi(x)\log d(x).
\]
则成立全区间可达性：
\[
\boxed{\ \{\mathsf{R}(\mu):\ \mu\in\mathcal{C}_\pi\}=[0,W_{\mathrm{inv}}(\pi)]\ }.
\]
等价地，\(\mathcal{C}_\pi\) 中可实现的 Shannon 熵恰覆盖闭区间
\[
\boxed{\ \{H(\mu):\ \mu\in\mathcal{C}_\pi\}=\bigl[H(\pi),\ H(\pi)+W_{\mathrm{inv}}(\pi)\bigr]\ }.
\]
其中右端点由 \(\mu=\mu^{\mathrm{e}}_\pi\) 取到（纤维均匀提升为最大熵 lift），左端点由每条纤维取原子条件分布取到。
\end{theorem}
\begin{proof}
取任意 \(\mu\in\mathcal{C}_\pi\)，令 \(\mu_x\in\Delta(F_x)\) 表示 \(\mu\) 在纤维上的条件分布，则对 \(\omega\in F_x\) 有
\[
\mu(\omega)=\pi(x)\mu_x(\omega),
\qquad
\mu^{\mathrm{e}}_\pi(\omega)=\pi(x)u_x(\omega),
\]
其中 \(u_x\) 为 \(F_x\) 上的均匀分布。
由 KL 的链式分解（对有限分布为直接恒等式）得
\[
D_{\mathrm{KL}}\!\left(\mu\ \middle\|\ \mu^{\mathrm{e}}_\pi\right)
=\sum_{x\in X}\pi(x)\,D_{\mathrm{KL}}(\mu_x\|u_x).
\]
并且对每个 \(x\)，有 \(D_{\mathrm{KL}}(\mu_x\|u_x)=\log d(x)-H(\mu_x)\)，因而
\[
0\le D_{\mathrm{KL}}(\mu_x\|u_x)\le \log d(x),
\]
且当 \(\mu_x\) 在 \(\Delta(F_x)\) 上变化时，\(H(\mu_x)\) 连续覆盖区间 \([0,\log d(x)]\)，从而 \(D_{\mathrm{KL}}(\mu_x\|u_x)\) 连续覆盖 \([0,\log d(x)]\)。
因此
\[
\{\mathsf{R}(\mu):\mu\in\mathcal{C}_\pi\}
=\sum_{x\in X}\pi(x)\,[0,\log d(x)]
=[0,\sum_{x\in X}\pi(x)\log d(x)]
=[0,W_{\mathrm{inv}}(\pi)],
\]
其中“和”按 Minkowski 和理解；最后一步由区间加法闭包性直接给出。
\par
另一方面，直接计算得
\[
\mathsf{R}(\mu)=D_{\mathrm{KL}}\!\left(\mu\ \middle\|\ \mu^{\mathrm{e}}_\pi\right)
=H(\mu^{\mathrm{e}}_\pi)-H(\mu),
\qquad
H(\mu^{\mathrm{e}}_\pi)=H(\pi)+W_{\mathrm{inv}}(\pi),
\]
从而熵谱结论与残差谱结论互相等价。证毕。
\end{proof}

\endinput


\begin{proposition}[折叠的规范性、幂等性与满射性]\label{prop:fold-basic}
折叠算子 $\Fold_m$ 良定义且满足：
\begin{enumerate}
\normalfont
  \item \textbf{（规范性）} $\Fold_m(\omega)\in X_m$；
  \item \textbf{（幂等性）} 若 $\omega\in X_m$，则 $\Fold_m(\omega)=\omega$；
  \item \textbf{（满射性）} $\Fold_m:\Omega_m\twoheadrightarrow X_m$ 为满射。
\end{enumerate}
\leanverified{Fold\_stable}
\leanverified{Fold\_idempotent}
\leanverified{Fold\_surjective}
\leanverified{paper\_fold\_surjective}
\leanverified{stableValue\_Fold\_lt}
\leanverified{exactWeightCount\_one\_eq}
\leanverified{paper\_fold\_basic}
\leanverified{paper\_allTrue\_not\_no11}
\end{proposition}

\begin{proof}
\textbf{（规范性）}由定义 \ref{def:fold-word}，$Z(N(\omega))\in X_\infty$，故 $\pi_m(Z(N(\omega)))\in X_m$。

\textbf{（幂等性）}若 $\omega\in X_m$，令 $N=\sum_{k=1}^m \omega_k F_{k+1}$。由于 $\omega$ 已是合法 Zeckendorf 数字的长度-$m$ 前缀，且 Zeckendorf 表示唯一，故 $Z(N)$ 的前 $m$ 位必等于 $\omega$，从而 $\Fold_m(\omega)=\omega$。

\textbf{（满射性）}对任意 $x\in X_m$，取 $\omega=x\in\Omega_m$，由幂等性得 $\Fold_m(\omega)=x$，故 $\Fold_m$ 满射。
\end{proof}


\subsection{稳定压缩率与分辨率不足的形式化}\label{subsec:folding-stable-compression-resolution}

折叠压缩率由
\[
\frac{|\Omega_m|}{|X_m|}=\frac{2^m}{F_{m+2}}
\]
刻画。由于 $F_{m+2}\asymp \varphi^m$，平均折叠多重度具有指数增长率 $(2/\varphi)^m$。因此，尽管稳定类型数量随 $m$ 指数增长，微态数量增长更快，折叠必然为多对一：不同微态在有限分辨率下不可区分而被归入同一稳定类型。该结论给出“分辨率不足导致规律不可读出”的严格表述：不可读出并非缺乏规律，而是稳定投影将大量微差异压缩到同一类中，从而在微观层面丢失可区分信息。

\begin{remark}[构造层与刚性层的区分]\label{rem:fold-constructibility-rigidity-split}
在方法上，应把 \(X_m\) 视为 Euclid 式“可构造可见层”：经过允许的折叠、归一化与规范回落后，真正进入对象语言的稳定代表。相对地，关于纤维大小、最大纤维、全息深度与残差账本的结论，则属于刚性层或障碍层：它们说明对象一旦被合法构造出来后，还有多少微观自由度被强制坍缩、多少信息必须外置承载。这样，“对象存在”“对象唯一”“可逆恢复失败”三类命题便可从一开始清楚分离。
\end{remark}

\subsubsection{信息二分与纤维残差}
在 Zeckendorf 语法约束下，折叠 $\Fold_m:\Omega_m\to X_m$ 将可观测微态信息分解为两类：其一为稳定类型 $x:=\Fold_m(\omega)\in X_m$，作为等价类的规范代表；其二为纤维残差 $r\in \Fold_m^{-1}(x)$，刻画被压缩抹去但在可逆重建口径下必须外置承载的自由度。压缩率 $|\Omega_m|/|X_m|=2^m/F_{m+2}$ 表明第二类自由度在分辨率增长时不可忽略；相应的纤维规模以 $d_m(x):=|\Fold_m^{-1}(x)|$ 定量刻画，并在小节 \ref{subsec:pom-information-ledger} 的信息账本中以 $\mathbb{E}_{\pi}[\log d_m(X)]$ 的形式进入熵分解。更精细的纤维内部结构可在投影本体数学部分直接闭式读出：纤维的分量分解与“逆代换次数分布”见定理 \ref{thm:pom-fiber-indset-factorization}--\ref{thm:pom-path-indset-poly-closed}；带符号指数与 \(\bmod 3\) 阻碍见定理 \ref{thm:pom-fiber-signed-index-periodicity}；纤维指数的热力学完成（LDP/CLT）见定理 \ref{thm:pom-fiber-index-ldp-thermo}--\ref{thm:pom-fiber-index-clt-thermo}。

\begin{theorem}[最大纤维内的全息深度可任意放大]\label{thm:fold-maxfiber-holographic-depth-amplification}
沿用折叠 \(\Fold_m:\Omega_m\to X_m\) 的记号，并记纤维多重度 \(d_m(x):=\card{\Fold_m^{-1}(x)}\) 与最大纤维
\(
D_m:=\max_{x\in X_m}d_m(x)
\)。
令 \(t:\NN\to\NN\) 为任意可计算且无界的时间界函数。
则存在无穷多个尺度 \(m\) 与最大达到者 \(x_m\in X_m\)（满足 \(d_m(x_m)=D_m\)），使得在纤维 \(\Fold_m^{-1}(x_m)\) 内存在微态 \(\omega_m\) 满足：
\begin{enumerate}
  \item \(K(\omega_m\mid x_m,m)=O(\log m)\)；
  \item 存在常数 \(c>0\) 使
  \[
  K^{t(m)}(\omega_m\mid x_m,m)\ \ge\ K(\omega_m\mid x_m,m)+c,
  \]
  其中 \(K\) 为固定通用机下的前缀条件 Kolmogorov 复杂度，\(K^{t(m)}\) 为时间受限的条件复杂度（要求在 \(t(m)\) 步内输出）。
\end{enumerate}
\end{theorem}

\begin{proof}
由定理 \ref{thm:pom-max-fiber} 的最大纤维渐近 \(D_m=\Theta(\varphi^{m/2})\)，可取无穷多个 \(m\) 使得对任意预先给定常数 \(A>0\) 有 \(D_m>m^{A}\)。
对每个这样的 \(m\)，取 \(x_m\) 为最大达到者。
\par
令 \(L:=\lceil A\log_2 m\rceil\)。
枚举所有长度不超过 \(L\) 的前缀程序 \(p\)，在条件输入 \((x_m,m)\) 下并行模拟每个 \(p\) 至多 \(t(m)\) 步。
令 \(S\) 为在该时间内确实输出、且输出落在纤维 \(\Fold_m^{-1}(x_m)\) 内的结果集合，则 \(\card{S}\le 2^{L}\le m^{A}\)。
由于 \(d_m(x_m)=D_m>m^A\)，存在 \(\omega_m\in \Fold_m^{-1}(x_m)\setminus S\)。
取字典序最小的此类 \(\omega_m\)。
\par
一方面，给定 \((x_m,m)\)，由固定算法（依赖于 \(t\) 与 \(A\)）执行上述并行枚举并取首个未在 \(t(m)\) 步内被长度 \(\le L\) 程序输出的纤维元素，即可确定 \(\omega_m\)，故 \(K(\omega_m\mid x_m,m)=O(\log m)\)。
另一方面，若存在长度 \(\le L\) 的程序能在 \(t(m)\) 步内输出 \(\omega_m\)，则 \(\omega_m\in S\)，与构造矛盾；因此 \(K^{t(m)}(\omega_m\mid x_m,m)\ge L+1\)。
取 \(A\) 充分大即可使 \(L+1-K(\omega_m\mid x_m,m)\) 被某个与 \(m\) 无关的常数 \(c>0\) 下界，从而得到所示不等式。证毕。
\end{proof}

\begin{proposition}[次指数纤维多重度下的全息码长率守恒]\label{prop:fold-subexp-multiplicity-holographic-rate-conservation}
设 \(\Pi_n:\Omega_n\to\bar\Omega_n\) 为任意多对一投影，并记最大多重度
\[
M_n:=\max_{\bar\omega\in\bar\Omega_n}\card{\Pi_n^{-1}(\bar\omega)}.
\]
给定任意码长函数 \(m:\Omega_n\to\RR_{\ge 0}\)。
对每个 \(\bar\omega\in\bar\Omega_n\) 定义最短码长与聚合码长
\[
m_{\min}(\bar\omega):=\min_{\omega\in\Pi_n^{-1}(\bar\omega)}m(\omega),
\qquad
m_{\mathrm{agg}}(\bar\omega):=-\log_2\sum_{\omega\in\Pi_n^{-1}(\bar\omega)}2^{-m(\omega)}.
\]
则存在与 \(n\) 无关的常数 \(C\) 使对任意 \(\bar\omega\in\bar\Omega_n\) 有
\[
\Bigl|m_{\mathrm{agg}}(\bar\omega)-m_{\min}(\bar\omega)\Bigr|
\ \le\ \log_2 M_n + C.
\]
特别地，若 \(\log M_n=o(n)\)，则
\[
\frac{1}{n}\Bigl|m_{\mathrm{agg}}(\bar\omega)-m_{\min}(\bar\omega)\Bigr|
\ \xrightarrow[n\to\infty]{}\ 0
\]
一致成立。
\end{proposition}

\begin{proof}
对任意 \(\bar\omega\)，由定义知 \(\sum_{\omega\in\Pi_n^{-1}(\bar\omega)}2^{-m(\omega)}\ge 2^{-m_{\min}(\bar\omega)}\)，从而
\(
m_{\mathrm{agg}}(\bar\omega)\le m_{\min}(\bar\omega)
\)。
另一方面，由 \(\card{\Pi_n^{-1}(\bar\omega)}\le M_n\) 与 \(m(\omega)\ge m_{\min}(\bar\omega)\) 得
\[
\sum_{\omega\in\Pi_n^{-1}(\bar\omega)}2^{-m(\omega)}
\le
\card{\Pi_n^{-1}(\bar\omega)}\,2^{-m_{\min}(\bar\omega)}
\le
M_n\,2^{-m_{\min}(\bar\omega)},
\]
从而 \(m_{\mathrm{agg}}(\bar\omega)\ge m_{\min}(\bar\omega)-\log_2 M_n\)。
将常数规范（由对数底与码长约定诱导）吸收到 \(C\) 即得第一式。
两端同除以 \(n\) 并用 \(\log M_n=o(n)\) 得一致收敛。证毕。
\end{proof}


\input{subsec__Phi_m-sliding-block-factor}

\subsection{诱导过程与转移结构：显式有限标记图}\label{subsec:Phi_m-graph}

下面给出 $Y_m$ 的一个完全显式的 sofic 图表示，从而刻画“将 $\Fold_m$ 沿长轨迹连续重叠窗口施加后，输出过程的状态/转移结构”。

\begin{definition}[长度 \texorpdfstring{$m-1$}{m-1} 后缀状态图]\label{def:G_m}
令顶点集为
\[
V_m:=\{0,1\}^{m-1}.
\]
对任意顶点 $v=v_1\cdots v_{m-1}\in V_m$ 与任意 $b\in\{0,1\}$，定义一条有向边
\[
e=(v,b):\ v\to v':=v_2\cdots v_{m-1}b,
\]
并用 $\ell_m(e)\in X_m$ 标记该边：
\[
\ell_m(e):=\Fold_m(v_1\cdots v_{m-1}b)\in X_m.
\]
由此得到一个有限标记图 $G_m=(V_m,E_m,\ell_m)$，其中 $|V_m|=2^{m-1}$，$|E_m|=2^{m}$。
\leanverified{phiState\_card}
\end{definition}

\begin{theorem}[显式 sofic 表示与可计算诱导结构]\label{thm:Phi_m-sofic-graph}
令 $Z_m\subset E_m^{\ZZ}$ 为图 $G_m$ 的双边边移位（即双边路径空间），并令
\[
(\mathcal{L}_m(e))_t:=\ell_m(e_t)\in X_m
\]
为逐边取标记得到的序列。则
\[
Y_m=\mathcal{L}_m(Z_m),
\]
从而 $Y_m$ 是一个由 $G_m$ 给出的 sofic 子移位；并且 $G_m$ 给出一个可审计的有限状态诱导结构：状态为长度 $m-1$ 的输入后缀，转移为追加一比特，输出为对应窗口经 $\Fold_m$ 的稳定类型。
\leanverified{paper\_phi\_m\_sofic\_graph}
\leanverified{phiEdge\_card}
\leanverified{phiGraph\_outdegree\_two}
\end{theorem}

\begin{proof}
\textbf{（包含 $\subset$）} 给定任意输入序列 $s\in\{0,1\}^{\ZZ}$，定义
\[
v_t:=s_t s_{t+1}\cdots s_{t+m-2}\in V_m,\qquad
e_t:=(v_t,s_{t+m-1})\in E_m.
\]
则 $e=(e_t)_{t\in\ZZ}$ 是 $G_m$ 上的一条双边路径（因为 $v_{t+1}$ 恰为 $v_t$ 去首加尾的更新）。并且由标记定义，
\[
(\mathcal{L}_m(e))_t=\ell_m(e_t)=\Fold_m(s_t\cdots s_{t+m-1})=(\Phi_m(s))_t.
\]
因此 $\Phi_m(s)\in \mathcal{L}_m(Z_m)$，即 $Y_m\subset \mathcal{L}_m(Z_m)$。

\textbf{（包含 $\supset$）} 反之，取任意双边路径 $e=(e_t)\in Z_m$。每条边 $e_t$ 具有形式 $(v_t,b_t)$，且路径条件保证 $v_{t+1}$ 等于 $v_t$ 去首加尾的更新。由此可一致地构造一个二元序列 $s\in\{0,1\}^{\ZZ}$：令 $s_{t+m-1}:=b_t$，并用 $v_t$ 提供窗口的其余 $m-1$ 位（路径相容性保证该定义不矛盾）。于是对该 $s$ 有
\[
(\Phi_m(s))_t=\Fold_m(s_t\cdots s_{t+m-1})=\Fold_m(v_t b_t)=\ell_m(e_t)=(\mathcal{L}_m(e))_t.
\]
因此 $\mathcal{L}_m(e)\in Y_m$，从而 $\mathcal{L}_m(Z_m)\subset Y_m$。两边相等，证毕。
\end{proof}

\begin{corollary}[序列层折叠的熵不降：每步保留 \texorpdfstring{$1$}{1} 比特自由度]\label{cor:Phi_m-entropy-no-drop}
对任意 \(m\ge 2\)，令 \(Z_m\subset E_m^{\ZZ}\) 为图 \(G_m\) 的双边边移位、\(\mathcal{L}_m:Z_m\to X_m^{\ZZ}\) 为逐边取标记映射（定义 \ref{def:G_m}），并记 \(Y_m=\mathcal{L}_m(Z_m)\)（定理 \ref{thm:Phi_m-sofic-graph}）。
则：
\begin{enumerate}
  \item \(Z_m\) 与二元全移位 \(\{0,1\}^{\ZZ}\) 共轭，因而 \(h_{\mathrm{top}}(Z_m)=\log 2\)；
  \item 标记图 \(G_m\) 右解析（同一顶点出发的两条出边标记不同；定理 \ref{thm:pom-right-resolving}），从而 \(\mathcal{L}_m\) 是有限对一因子码（每个输出序列的 lift 数至多为 \(|V_m|=2^{m-1}\)）；
  \item 因而拓扑熵保持：
  \[
  \boxed{\ h_{\mathrm{top}}(Y_m)=h_{\mathrm{top}}(Z_m)=\log 2\ }.
  \]
\end{enumerate}
该结论给出“窗口层多对一压缩”与“序列层记忆补偿”的严格命题：\(\Phi_m\) 将窗口层丢失的纤维自由度迁移到时间方向的有限状态记忆中，而不在熵意义下压缩每步 \(1\) 比特的自由度。
\leanverified{LabeledPath}
\leanverified{labeledPath\_initial\_injective}
\leanverified{labeledPath\_card\_le\_state\_card}
\leanverified{labeledPath\_count\_le\_states}
\leanverified{paper\_labeled\_path\_lift\_bound}
\leanverified{paper\_Phi\_m\_entropy\_no\_drop}
\end{corollary}
\begin{proof}
\textbf{(1)}--\textbf{(2)} 见定理 \ref{thm:Phi_m-sofic-graph}（共轭）与右解析性；\textbf{(3)} 有限对一因子码熵保持可参见 \cite{LindMarcus1995,Kitchens1998SymbolicDynamics}。
\end{proof}

\begin{remark}[时间记忆作为隐寄存器]\label{rem:Phi-m-memory-implicit-register}
窗口层的多对一压缩并不等价于时间层的信息损失：当 $m\ge 3$ 时，定理 \ref{thm:Phi_m-conjugacy-threshold} 给出有限窗逆码 $\Psi_m$，从而被折叠抹去的自由度可迁移到时间方向的有限状态记忆并在有限延迟后恢复。这提供了“第三寄存器”的另一种实现方式：其不必以静态轴显式存储，而可作为同步滤波器的内部状态隐式承载。在规范群胚口径下，该隐寄存器对应局部截面之间的粘接数据（见小节 \ref{subsec:pom-gauge-groupoid}）。
\end{remark}


\subsubsection{可逆性阈值：\texorpdfstring{$\Phi_m$}{Phi\_m} 的共轭与纤维坍缩}\label{subsec:Phi_m-conjugacy}

上一节已将 $Y_m$ 组织为标记图 $G_m$ 的像（定理 \ref{thm:Phi_m-sofic-graph}），并证明熵率不降（推论 \ref{cor:Phi_m-entropy-no-drop}）。
本节给出可逆性与纤维结构的确定结论：在 Zeckendorf 折叠 $\Fold_m$ 的口径下，除去 $m=2$ 的两点退化外，$\Phi_m$ 在 $m\ge 3$ 时为共轭。

\paragraph{可终止的局部判定口径（子集滤波）}
沿定义 \ref{def:G_m}，对任意候选后缀集合 $S\subseteq V_m$ 与输出符号 $a\in X_m$，定义一步更新
$$
\delta_m(S,a):=\{v'\in V_m:\exists\,e:v\to v'\ \text{s.t. }v\in S,\ \ell_m(e)=a\}.
$$
从初始集合 $S_0:=V_m$ 出发，对任意输出块 $a_0a_1\cdots a_{t-1}$ 递推
$$
S_{j+1}:=\delta_m(S_j,a_j)\qquad(0\le j<t).
$$
该递推可视为“仅由输出块驱动”的有限状态滤波器：$S_j$ 记录读入前 $j$ 个输出符号后，内部后缀（顶点）仍可能取到的候选集。
当 $S_t$ 为单点集时，意味着该长度-$t$ 输出块已将内部后缀同步到唯一值。

\begin{theorem}[可逆性阈值与有限窗逆码]\label{thm:Phi_m-conjugacy-threshold}
令 $\Phi_m$ 为定义 \ref{def:Phi_m} 的滑动块码，$Y_m$ 为定义 \ref{def:Y_m} 的像子移位。
则：
\begin{enumerate}
  \item 当 $m\ge 3$ 时，$\Phi_m:\{0,1\}^{\ZZ}\to Y_m$ 为拓扑共轭；更具体地，存在一个滑动块码
  $$
  \Psi_m:Y_m\to\{0,1\}^{\ZZ},
  $$
  其\textbf{记忆}为 $m-1$、\textbf{预期性}为 $0$，并满足
  $$
  \Psi_m\circ\Phi_m=\mathrm{id}_{\{0,1\}^{\ZZ}},\qquad
  \Phi_m\circ\Psi_m=\mathrm{id}_{Y_m}.
  $$
  \item 当 $m=2$ 时，$\Phi_2$ 恰好在两条常值序列 $0^{\ZZ}$ 与 $1^{\ZZ}$ 上出现 $2$-对-$1$ 退化，其余点一一。
\end{enumerate}
\end{theorem}
\leanverified{paper\_Phi\_m\_conjugacy\_threshold}

\begin{corollary}[像子移位的共轭类塌缩（\(m\ge 3\)）]\label{cor:Ym-conjugacy-all-mge3}
对任意 $m,n\ge 3$，移位系统 $(Y_m,\sigma)$ 与 $(Y_n,\sigma)$ 拓扑共轭。
更具体地，若取定理 \ref{thm:Phi_m-conjugacy-threshold} 中的逆码 $\Psi_m:Y_m\to\{0,1\}^{\ZZ}$，则
\[
h_{m\to n}:=\Phi_n\circ \Psi_m:\ Y_m\longrightarrow Y_n
\]
为同胚且满足 \(h_{m\to n}\circ\sigma=\sigma\circ h_{m\to n}\)。
\end{corollary}
\begin{proof}
由定理 \ref{thm:Phi_m-conjugacy-threshold}，当 \(m\ge 3\) 时 \(\Phi_m\) 为 \(\{0,1\}^{\ZZ}\) 与 \(Y_m\) 的共轭，其逆映射为 \(\Psi_m\)。
因此 \(h_{m\to n}\) 为同胚；又由于 \(\Phi_n\) 与 \(\Psi_m\) 均与移位交换，故 \(h_{m\to n}\) 与移位交换。
证毕。
\end{proof}
\leanverified{paper\_Ym\_conjugacy\_all\_mge3}
\leanverified{paper\_ym\_conjugacy\_all\_mge3}

\begin{proof}
证明思路建立在“有限可判定”的滤波器上。

\paragraph{局部同步：长度-$m$ 块把候选集压到单点}
对固定 $m$，考虑上述滤波更新 $S_{j+1}=\delta_m(S_j,a_j)$，并记
$$
R_m:=\{S_m:\exists\,a_0\cdots a_{m-1}\in\mathcal{L}_m(Y_m)\ \text{s.t. 递推得到 }S_m\}.
$$
若 $R_m$ 中每个元素都是单点集，则任意长度-$m$ 的输出块都把内部后缀同步到唯一顶点。
该命题对每个固定 $m$ 是一个\textbf{完全有限}的判定问题：状态空间为 $V_m$ 的幂集可达部分，转移由 $G_m$ 给出。
脚本 \texttt{scripts/exp\_phi\_m\_conjugacy\_threshold.py} 对 $m\in[2,12]$ 的范围给出可复算证书，输出汇总见表 \ref{tab:phi_m_conjugacy_threshold}。
\begin{table}[H]
\centering
\scriptsize
\setlength{\tabcolsep}{6pt}
\caption{Conjugacy threshold certificate for the sliding-block factor $\Phi_m$. We build the subset-filter automaton from the labeled graph $G_m$ and compute $R_{m-1}$ and $R_m$, the sets of reachable filter states after exactly $m-1$ and $m$ output symbols. If all states in $R_m$ are singletons, then each length-$m$ output word has a unique length-$(2m-1)$ preimage block, yielding a finite-window inverse (memory $m-1$). The non-singleton count in $R_{m-1}$ provides a sharpness witness for the sync delay. We also cross-check periodic-point counts up to a small $n$.}
\label{tab:phi_m_conjugacy_threshold}
\begin{tabular}{r r r r r r r l r}
\toprule
$m$ & det states & det edges & $|R_{m-1}|$ & non-singleton ends $(m-1)$ & $|R_m|$ & non-singleton ends $(m)$ & periodic $\#\mathrm{Fix}_n$ & checked $n$\\
\midrule
2 & 3 & 7 & 3 & 1 & 3 & 1 & $2^n-1$ & 8\\
3 & 7 & 20 & 5 & 1 & 4 & 0 & $2^n$ & 8\\
4 & 16 & 50 & 10 & 2 & 8 & 0 & $2^n$ & 8\\
5 & 32 & 95 & 20 & 4 & 16 & 0 & $2^n$ & 8\\
6 & 69 & 204 & 40 & 8 & 32 & 0 & $2^n$ & 8\\
7 & 142 & 404 & 80 & 16 & 64 & 0 & $2^n$ & 8\\
8 & 296 & 827 & 160 & 32 & 128 & 0 & $2^n$ & 8\\
9 & 607 & 1656 & 320 & 64 & 256 & 0 & $2^n$ & 8\\
10 & 1244 & 3342 & 640 & 128 & 512 & 0 & $2^n$ & 8\\
11 & 2532 & 6703 & 1280 & 256 & 1024 & 0 & $2^n$ & 8\\
12 & 5141 & 13464 & 2560 & 512 & 2048 & 0 & $2^n$ & 8\\
\bottomrule
\end{tabular}
\end{table}

其中 $m\ge 3$ 时 $R_m$ 的“non-singleton ends $(m)$”为 $0$，而 $m=2$ 时在任何步数都出现唯一一个非单点端态，对应常值输出块所导致的两点退化。
同时表中 “non-singleton ends $(m-1)$” 在 $m\ge 3$ 时严格为正，这给出同步延迟的尖锐性见证：一般需要完整的长度-$m$ 输出块才能把候选集压到单点（对应逆码记忆为 $m-1$）。

\paragraph{从单点端态到唯一 lift：利用双向解析性}
由定理 \ref{thm:pom-right-resolving} 与定理 \ref{thm:pom-left-resolving}，标记图 $G_m$ 在 Zeckendorf 折叠下同时右解析与左解析：同一顶点出发（或到达）的一步扩展可由标记唯一判别。
因此一旦长度-$m$ 输出块把端态同步为单点 $\{v'\}$，则该块在 $G_m$ 上的长度-$m$ lift 路径不仅存在，而且由左解析性可从终点 $v'$ 逐步唯一回溯得到；从而该输出块对应唯一的长度-$(2m-1)$ 输入块。

\paragraph{构造逆码 \texorpdfstring{$\Psi_m$}{Psi\_m}}
当 $m\ge 3$ 时，由上述两步可知：对任意允许的长度-$m$ 输出块
$$
a_{t-m+1}\cdots a_t\in\mathcal{L}_m(Y_m),
$$
存在唯一的 $G_m$-路径
$$
e_{t-m+1},\dots,e_t\in E_m
$$
使得其逐边标记恰为该输出块。
令 $e_{t-m+1}=(v,b)$（其中 $b\in\{0,1\}$ 是该边对应的追加比特），定义
$$
(\Psi_m(y))_t:=b,
$$
其中 $y\in Y_m$ 且 $a_j=y_j$。
该定义仅依赖 $y_{t-m+1},\dots,y_t$，故 $\Psi_m$ 的记忆为 $m-1$、预期性为 $0$。
由构造立得 $\Psi_m\circ\Phi_m=\mathrm{id}$，并因 $\Phi_m$ 满射到 $Y_m$ 得 $\Phi_m\circ\Psi_m=\mathrm{id}_{Y_m}$。

\paragraph{$m=2$ 的两点退化}
脚本证书给出：唯一的非单点端态对应常值输出块，其恰来自两条输入常值序列 $0^{\ZZ},1^{\ZZ}$；
此外对任意其它输出点滤波在有限步内同步到单点，从而 lift 唯一。
综上得到 $m=2$ 的“仅两点 $2$-对-$1$”退化结论。
\end{proof}

\begin{remark}[可逆性阈值的同调解释：\(m=2\) 的唯一分支点对应不可消去的 \texorpdfstring{$\mathbb{Z}/2$}{Z/2} 隐变量]\label{rem:Phi_m-conjugacy-threshold-z2-homology}
定理 \ref{thm:Phi_m-conjugacy-threshold} 可被理解为：一个 \(\mathbb{Z}/2\)-值的纤维标签是否能被长度-\(m\) 的输出窗口制度性读出，从而被“规范化”为同调零（coboundary）的判定。
滤波器递推 \(S_{j+1}=\delta_m(S_j,a_j)\) 记录在给定输出块下内部后缀的候选集合；当 \(m\ge 3\) 时，任意长度-\(m\) 输出块将候选集合同步到单点（表 \ref{tab:phi_m_conjugacy_threshold}），因而该 \(\mathbb{Z}/2\) 标签在窗口内可被完全同步并由逆码 \(\Psi_m\) 消去，表现为共轭。
相反，当 \(m=2\) 时，唯一的非单点端态恰对应两条常值轨道 \(0^{\ZZ},1^{\ZZ}\)，等价于：在该分支轨道上存在一个不可由长度-\(2\) 输出决定的 \(\mathbb{Z}/2\) 隐变量；它给出唯一的拓扑分支点并构成共轭失败的唯一障碍。
\end{remark}

\begin{corollary}[\texorpdfstring{$\zeta_{Y_m}$}{zetaYm} 的封口与“折叠因子不产出非平凡 prime shadow”]\label{cor:Phi_m-zeta-closed}
对 $m\ge 3$，由于 $Y_m$ 与二元全移位共轭（定理 \ref{thm:Phi_m-conjugacy-threshold}），其周期点计数满足
$$
\#\mathrm{Fix}(\sigma^n|_{Y_m})=2^n\qquad(n\ge 1),
$$
从而内禀 Artin--Mazur \(\zeta\)-函数为
$$
\zeta_{Y_m}(z)=\frac{1}{1-2z}.
$$
而 $m=2$ 时有
$$
\#\mathrm{Fix}(\sigma^n|_{Y_2})=2^n-1,\qquad
\zeta_{Y_2}(z)=\frac{1-z}{1-2z}.
$$
因此当 $m\ge 3$ 时，$Y_m$ 自身的 prime-orbit/\(\zeta\) 指纹退化为全移位骨架；本文中出现的非平凡 prime shadow 必然来自”核 + 势/权重 + 扭曲通道”（而非单纯来自折叠因子 $Y_m$）。
\leanverified{paper\_fold\_image\_subshift\_zeta\_closed}
\end{corollary}


\begin{remark}[关于 Markov 性与可计算近似]\label{rem:Phi_m-not-markov}
定理 \ref{thm:Phi_m-sofic-graph} 给出的是\emph{拓扑/组合}层面的有限状态表示：$Y_m$ 的所有允许序列由有限图刻画。但对任意输入过程（例如第 \ref{sec:experiments} 节的旋转扫描产生的 Sturmian 序列），其推前测度 $(\Phi_m)_*\mu$ 在 $Y_m$ 上未必是一阶马尔可夫测度。若需要“诱导转移结构”的数值近似，可在图 $G_m$ 上用长样本估计边/状态频率，从而得到一个可审计的有限状态 Markov 近似；该近似与第 \ref{sec:quant-fit} 节的直方图/偏差指标对齐。
\end{remark}


\subsection{从 sofic 图像到可计算的 Markov 覆盖：邻接矩阵/熵/最大熵测度与 \texorpdfstring{$\zeta$}{zeta}}\label{subsec:Phi_m-fischer}
上一小节将 $Y_m$ 表示为显式标记图 $G_m$ 的像，从而给出“$\Phi_m$ 的像空间为 sofic”的可核验表示。为将该表示提升为 ETDS 语境下常用的可计算不变量（邻接矩阵、熵、最大熵测度以及 \(\zeta\) 的行列式接口），需要引入 sofic shift 的标准 Markov 覆盖（cover）理论。

\subsubsection{右解析（right-resolving）表示与 Fischer cover（只保留可计算接口）}\label{subsubsec:Phi_m-right-resolving-fischer-cover}
对 sofic shift，存在\textbf{右解析}标记图表示，并可进一步最小化为右 Fischer cover。本文仅需其可计算接口与两个结构性结论；相关构造背景见 \cite{LindMarcus1995,Kitchens1998SymbolicDynamics}。

\begin{definition}[右解析表示的确定化口径]\label{def:Phi_m-right-resolving}
令 $G_m=(V_m,E_m,\ell_m)$ 为定义 \ref{def:G_m} 的标记图（字母表 $X_m$）。对任意子集状态 $S\subset V_m$ 与输出符号 $a\in X_m$，定义转移
\[
\delta_m(S,a):=\{v'\in V_m:\exists\,e:v\to v'\ \text{s.t. } v\in S,\ \ell_m(e)=a\}.
\]
从初始状态 $V_m$ 出发，取所有可达的 $S$ 作为状态集并按 $S\xrightarrow{a}\delta_m(S,a)$ 加边，得到一个确定性右解析标记图表示，记为 $H_m^{\mathrm{det}}$。
\end{definition}

\begin{proposition}[右解析表示的存在与等价]\label{prop:Phi_m-right-resolving}
由定义 \ref{def:Phi_m-right-resolving} 的确定化构造得到的标记图 $H_m^{\mathrm{det}}$ 是右解析的，并呈现与 $G_m$ 相同的标记语言；因此其标记移位等于 $Y_m$。
\end{proposition}

\begin{proof}
确定化保证“给定状态 $S$ 与标记 $a$ 时后继唯一”，故右解析；而 $H_m^{\mathrm{det}}$ 与 $G_m$ 识别相同正则语言可由自动机确定化直接得到。参见 \cite{LindMarcus1995,Kitchens1998SymbolicDynamics}。
\end{proof}
\leanverified{paper\_Phi\_m\_right\_resolving}
\leanverified{paper\_phi\_m\_right\_resolving}


\subsubsection{歧义壳的瞬态结构：单调性、无环性与同步延迟}\label{subsec:Ym-amb-transient}
本节把推论 \ref{cor:Ym-ambiguity-shell-low-entropy} 的“歧义壳低熵”进一步细化为一个可审计的结构命题：在本文的 Zeckendorf 折叠下（$m\ge 3$），歧义并非可循环的“二义语法”，而是\textbf{有限步同步的纯瞬态层}；因此它不承载任何周期点，也不改变内禀 \(\zeta_{Y_m}\) 的周期点谱（见推论 \ref{cor:Phi_m-zeta-closed}）。

\begin{proposition}[候选集合单调不增（右解析的集合动力学）]\label{prop:det-candidate-set-monotone}
在定义 \ref{def:G_m} 的标记图 $G_m$ 上，若同一顶点的两条出边标记不同（即 $G_m$ 右解析；见定理 \ref{thm:pom-right-resolving}），则对任意候选集 $S\subset V_m$ 与任意输出符号 $a\in X_m$，若 $\delta_m(S,a)\neq\varnothing$ 则
\[
|\delta_m(S,a)|\le |S|.
\]
并且对任意单点集 $\{v\}$，其后继 $\delta_m(\{v\},a)$ 要么为空、要么仍为单点；换言之，$H_m^{\mathrm{det}}$ 中的单点态集合在前向读入下是封闭的：一旦进入单点态，就不可能回到非单点态。
\end{proposition}
\begin{proof}
右解析性意味着：对每个 $v\in V_m$ 与每个 $a\in X_m$，从 $v$ 出发、标记为 $a$ 的出边至多一条。于是可定义一个部分函数
\[
f_a:V_m\rightharpoonup V_m,\qquad
f_a(v):=
\begin{cases}
v' & \exists\,e:v\to v'\ \text{s.t. }\ell_m(e)=a,\\
\text{undefined} & \text{otherwise}.
\end{cases}
\]
按定义 \ref{def:Phi_m-right-resolving}，
\(
\delta_m(S,a)=\{f_a(v):v\in S,\ f_a(v)\text{ defined}\}
\)
是 $S$ 在 $f_a$ 下的像，因此其基数不超过 $|S|$。当 $S=\{v\}$ 时，$\delta_m(\{v\},a)$ 至多包含一个元素，故为单点或空集，从而单点态对前向读入封闭。证毕。
\leanverified{succSet}
\leanverified{option\_toFinset\_card\_le\_one}
\leanverified{succSet\_singleton}
\leanverified{succSet\_singleton\_of\_some}
\leanverified{succSet\_singleton\_of\_none}
\leanverified{paper\_det\_candidate\_set\_monotone}
\end{proof}

\begin{theorem}[歧义壳无环：纯瞬态且无周期点]\label{thm:Ym-ambiguity-shell-dag}
令 $H_m^{\mathrm{det}}$ 为定义 \ref{def:Phi_m-right-resolving} 的确定化右解析表示，其状态集记为 $\mathcal{S}_m^{\mathrm{det}}$。定义歧义态集合
\[
\mathcal{S}_m^{\mathrm{amb}}:=\{S\in\mathcal{S}_m^{\mathrm{det}}:\ |S|\ge 2\},
\]
并令 $H_m^{\mathrm{amb}}$ 为 $H_m^{\mathrm{det}}$ 在 $\mathcal{S}_m^{\mathrm{amb}}$ 上的诱导子图。
则当 $m\ge 3$ 时，$H_m^{\mathrm{amb}}$ 为 DAG（不含任何有向环）。因此：
\begin{enumerate}
  \item $H_m^{\mathrm{amb}}$ 不承载任何周期轨道；等价地，$Y_m$ 的任何周期点都不可能在确定化滤波中永久停留在非单点态。
  \item 对任意 $y\in Y_m$，从初始候选集 $S_0:=V_m$ 出发递推 $S_{t+1}:=\delta_m(S_t,y_t)$，则在有限步后必进入单点态并从此保持单点；更具体地，对所有 $t\ge m$ 有 $|S_t|=1$。
\end{enumerate}
\end{theorem}
\begin{proof}
由定理 \ref{thm:Phi_m-conjugacy-threshold} 的“局部同步”部分（并见表 \ref{tab:phi_m_conjugacy_threshold}），当 $m\ge 3$ 时：对任意允许的长度-$m$ 输出块 $a_0\cdots a_{m-1}$，从 $S_0=V_m$ 出发按 $S_{j+1}=\delta_m(S_j,a_j)$ 递推得到的 $S_m$ 必为单点集。等价地，沿 $H_m^{\mathrm{det}}$ 从初始状态 $V_m$ 出发的任意长度-$m$ 标记路径，其终点都落在单点态集合中。

结合命题 \ref{prop:det-candidate-set-monotone} 的“单点态前向封闭性”，立刻得到：对任意 $y\in Y_m$ 的滤波轨道 $(S_t)_{t\ge 0}$，一旦 $S_m$ 为单点，则对所有 $t\ge m$ 都保持单点，故结论 (2) 成立。

为证 $H_m^{\mathrm{amb}}$ 无环，反设存在一个有向环完全落在 $\mathcal{S}_m^{\mathrm{amb}}$ 中。由于 $\mathcal{S}_m^{\mathrm{det}}$ 按定义是从初始状态 $V_m$ 可达的状态集合，该环上任取一状态必可由某条从 $V_m$ 出发的有限标记路径到达。随后沿该环循环可生成任意长的、始终停留在非单点态中的标记路径；这与“任意长度-$m$ 路径终点必落入单点态”矛盾。故 $H_m^{\mathrm{amb}}$ 不含有向环，为 DAG，从而也不承载周期轨道，结论 (1) 随之成立。证毕。
\end{proof}
\leanverified{paper\_Ym\_ambiguity\_shell\_dag}
\leanverified{paper\_ym\_ambiguity\_shell\_dag}

\begin{theorem}[歧义壳的 \texorpdfstring{$\zeta$}{zeta} 因子平凡性与单点态骨架封口]\label{thm:Ym-ambiguity-shell-zeta-factor-trivial}
设 \(m\ge 3\)。
沿用定义 \ref{def:Phi_m-right-resolving} 与定理 \ref{thm:Ym-ambiguity-shell-dag} 的记号。
记
\[
\mathcal S_m^{\mathrm{ess}}
:=
\bigl\{\{v\}:v\in V_m\bigr\}\cap \mathcal S_m^{\mathrm{det}}
\]
为 \(H_m^{\mathrm{det}}\) 中全部单点态所成的状态集，并令 \(H_m^{\mathrm{ess}}\) 表示其诱导子图。
再分别记
\[
A_m^{\mathrm{det}},\qquad A_m^{\mathrm{ess}}
\]
为 \(H_m^{\mathrm{det}}\) 与 \(H_m^{\mathrm{ess}}\) 的按边条数计数的邻接矩阵。
若按“先歧义态、后单点态”的顺序排列 \(\mathcal S_m^{\mathrm{det}}\)，则存在某个矩阵 \(R_m\) 使
\[
\boxed{\;
A_m^{\mathrm{det}}
=
\begin{pmatrix}
N_m & R_m\\
0 & A_m^{\mathrm{ess}}
\end{pmatrix},
\qquad
N_m^{\,m}=0.\;}
\]
因此
\[
\boxed{\;
\det(I-zA_m^{\mathrm{det}})
=
\det(I-zA_m^{\mathrm{ess}}).\;}
\]
若进一步定义表示图层面的 edge-\(\zeta\) 函数
\[
\zeta_{H_m^{\mathrm{det}}}(z):=\frac{1}{\det(I-zA_m^{\mathrm{det}})},
\qquad
\zeta_{H_m^{\mathrm{ess}}}(z):=\frac{1}{\det(I-zA_m^{\mathrm{ess}})},
\]
则有
\[
\boxed{\;
\zeta_{H_m^{\mathrm{det}}}(z)=\zeta_{H_m^{\mathrm{ess}}}(z).\;}
\]
换言之，歧义壳对表示图层面的 \(\zeta\)-函数贡献恒等因子 \(1\)。
结合推论 \ref{cor:Phi_m-zeta-closed} 可知：\(Y_m\) 的全部周期谱由单点态骨架独自承载，并满足
\[
\zeta_{Y_m}(z)=\frac{1}{1-2z}.
\]
\end{theorem}
\begin{proof}
由命题 \ref{prop:det-candidate-set-monotone}，单点态集合在前向读入下封闭，因此从单点态不可能返回歧义态。
这正给出所示块上三角形状中的左下零块。

接着证明 \(N_m^{\,m}=0\)。
对任意允许的长度-\(m\) 输出块 \(a_0\cdots a_{m-1}\in\mathcal L_m(Y_m)\)，定理 \ref{thm:Phi_m-conjugacy-threshold} 的证明已经表明：该块对应唯一的长度-\((2m-1)\) lift 块，因而其终端顶点唯一确定。
等价地，对任意可达候选集 \(S\subseteq V_m\)，沿递推
\[
S_{j+1}:=\delta_m(S_j,a_j),
\qquad
S_0:=S,
\]
在读入这 \(m\) 个输出符号后，所得 \(S_m\) 至多为单点。
因此不可能存在完全停留在歧义态集合 \(\mathcal S_m^{\mathrm{amb}}\) 内的长度-\(m\) 路径。
这正等价于 \(N_m^{\,m}=0\)。

于是 \(N_m\) 为幂零矩阵，故
\[
\det(I-zN_m)=1.
\]
再由块上三角行列式分解，
\[
\det(I-zA_m^{\mathrm{det}})
=
\det(I-zN_m)\det(I-zA_m^{\mathrm{ess}})
=
\det(I-zA_m^{\mathrm{ess}}),
\]
从而得到所示 edge-\(\zeta\) 恒等式。

最后，推论 \ref{cor:Phi_m-zeta-closed} 已给出 \(m\ge 3\) 时
\(
\zeta_{Y_m}(z)=1/(1-2z)
\)；
而定理 \ref{thm:Ym-ambiguity-shell-dag} 表明全部周期点都不可能永久驻留于歧义壳，故内禀周期谱完全由单点态骨架承载。证毕。
\end{proof}
\leanverified{paper\_ym\_ambiguity\_shell\_zeta\_factor\_trivial}

\begin{corollary}[周期点刚性与内禀 \texorpdfstring{$\zeta_{Y_m}$}{zetaYm} 的退化]\label{cor:Ym-periodic-rigidity}
当 $m\ge 3$ 时，对所有 $n\ge 1$ 有
\[
\#\mathrm{Fix}(\sigma^n|_{Y_m})=2^n,
\qquad
\zeta_{Y_m}(z)=\frac{1}{1-2z}.
\]
\end{corollary}
\begin{proof}
这是推论 \ref{cor:Phi_m-zeta-closed} 的直接重述（其证明基于定理 \ref{thm:Phi_m-conjugacy-threshold} 的共轭）。定理 \ref{thm:Ym-ambiguity-shell-dag} 进一步给出一个结构性解释：歧义壳为纯瞬态层，不贡献任何周期点，因此周期点谱完全由单点态骨架决定。证毕。
\end{proof}
\leanverified{paper\_Ym\_periodic\_rigidity}
\leanverified{paper\_ym\_periodic\_rigidity}

\begin{corollary}[同步延迟的闭式与尖锐性]\label{cor:Ym-sync-delay}
对 $y\in Y_m$ 定义确定化滤波轨道 $S_0:=V_m$、$S_{t+1}:=\delta_m(S_t,y_t)$（$t\ge 0$），并定义“最坏同步延迟”
\[
D_{\mathrm{sync}}(m):=\max\{\,t\ge 0:\exists\,y\in Y_m\ \text{s.t. }|S_t|>1\,\}.
\]
则当 $m\ge 3$ 时有
\[
\boxed{\ D_{\mathrm{sync}}(m)=m-1\ }.
\]
\leanverified{paper\_fold\_sync\_delay}
\end{corollary}
\begin{proof}
由定理 \ref{thm:Phi_m-conjugacy-threshold}（并见表 \ref{tab:phi_m_conjugacy_threshold} 中 $R_m$ 的“non-singleton ends $(m)$”为 $0$），对任意 $y\in Y_m$ 都有 $|S_m|=1$，结合命题 \ref{prop:det-candidate-set-monotone} 的单调性知对所有 $t\ge m$ 亦有 $|S_t|=1$，从而 $D_{\mathrm{sync}}(m)\le m-1$。
另一方面，同一张表同时给出 $R_{m-1}$ 的“non-singleton ends $(m-1)$”严格为正（$m\ge 3$），这意味着存在某个允许输出块使得 $|S_{m-1}|>1$，故 $D_{\mathrm{sync}}(m)\ge m-1$。两边合并即得等式。证毕。
\end{proof}

\begin{remark}[歧义壳作为子系统在 \texorpdfstring{$m\ge 3$}{m>=3} 时为空]\label{rem:Ym-amb-empty}
在推论 \ref{cor:Ym-ambiguity-shell-low-entropy} 中，我们把“完全不进入单点本质分量”的点集记为歧义壳子系统 \(Y_m^{\mathrm{amb}}\)。
而定理 \ref{thm:Ym-ambiguity-shell-dag} 的结论 (2) 表明：当 \(m\ge 3\) 时，从任意输出 \(y\in Y_m\) 出发的确定化滤波在至多 \(m\) 步内必进入并永远停留在单点态。
因此在本文的 Zeckendorf 折叠模型里，\(Y_m^{\mathrm{amb}}=\varnothing\)（两边移位意义下不存在“永远保持歧义”的轨道）；所谓“歧义壳”真正表达的是一个\emph{有限步同步延迟层}，其最坏延迟由推论 \ref{cor:Ym-sync-delay} 精确给出。
\par
进一步，条件 \(Y_m^{\mathrm{amb}}=\varnothing\) 可在确定化状态图上提升为一个\textbf{可计算的同步延迟定理}：
若在 $H_m^{\mathrm{det}}$ 中把所有单点态及其出边删去、得到非单点态诱导子图 $H_m^{\mathrm{amb}}$，则
\[
Y_m^{\mathrm{amb}}=\varnothing
\quad\Longleftrightarrow\quad
H_m^{\mathrm{amb}}\ \text{不含任何有向环}.
\]
事实上，若 $H_m^{\mathrm{amb}}$ 含有向环，则从初始状态可达性出发，沿该环可生成一条永不进入单点态的双边标记序列，从而得到 \(Y_m^{\mathrm{amb}}\neq\varnothing\)，矛盾；反向蕴含由定义即得。
因此在 \(Y_m^{\mathrm{amb}}=\varnothing\) 时，$H_m^{\mathrm{amb}}$ 为 DAG，从而存在统一有限同步步数
\[
L_m:=D_{\mathrm{sync}}(m)+1<\infty,
\]
使得任意 \(y\in Y_m\) 的确定化滤波轨道在不超过 \(L_m\) 步内进入单点态并从此保持单点。并且 $L_m$ 可在有限图上直接审计：对 DAG 做最长路径 DP 即得最小的 \(D_{\mathrm{sync}}(m)\)（从而得到 \(L_m\)）。脚本 \texttt{scripts/exp\_phi\_m\_sync\_theory.py} 输出该计算的可复现证书（表 \ref{tab:phi_m_sync_theory}）。
\end{remark}

\begin{table}[H]
\centering
\scriptsize
\setlength{\tabcolsep}{6pt}
\caption{Synchronization constants from the subset-filter automaton for $\Phi_m$. $H_m^{\mathrm{amb}}$ is the induced subgraph on non-singleton filter states. If $H_m^{\mathrm{amb}}$ has a directed cycle, then $Y_m^{\mathrm{amb}}\neq\varnothing$ and the worst-case sync delay is infinite; otherwise $H_m^{\mathrm{amb}}$ is a DAG and $D_{\mathrm{sync}}$ is obtained by a longest-path DP, with $L_m=D_{\mathrm{sync}}+1$. When $Y_m^{\mathrm{amb}}\neq\varnothing$, we also report $h_{\mathrm{top}}(Y_m^{\mathrm{amb}})$ and the entropy gap $\varepsilon_m=\log 2-h_{\mathrm{top}}(Y_m^{\mathrm{amb}})$.}
\label{tab:phi_m_sync_theory}
\begin{tabular}{r r r r r r r r}
\toprule
$m$ & det states & amb states & amb cycle? & $D_{\mathrm{sync}}$ & $L_m$ & $h_{\mathrm{top}}(Y_m^{\mathrm{amb}})$ & $\varepsilon_m$\\
\midrule
2 & 3 & 1 & yes & $\infty$ & $\infty$ & 0 & 0.693147\\
3 & 7 & 3 & no & 2 & 3 & -- & --\\
4 & 16 & 8 & no & 3 & 4 & -- & --\\
5 & 32 & 16 & no & 4 & 5 & -- & --\\
6 & 69 & 37 & no & 5 & 6 & -- & --\\
7 & 142 & 78 & no & 6 & 7 & -- & --\\
8 & 296 & 168 & no & 7 & 8 & -- & --\\
9 & 607 & 351 & no & 8 & 9 & -- & --\\
10 & 1244 & 732 & no & 9 & 10 & -- & --\\
11 & 2532 & 1508 & no & 10 & 11 & -- & --\\
12 & 5141 & 3093 & no & 11 & 12 & -- & --\\
\bottomrule
\end{tabular}
\end{table}


\begin{proposition}[主极点留数比作为常数级核指纹]\label{prop:Ym-zeta-residue-ratio}
设 cover 的 \(\zeta\) 与内禀 \(\zeta\) 在共同的主极点模长 \(z_\star=\rho_m^{-1}\)（其中 \(\rho_m:=\rho(A_m)=\rho(B_m)\)）处均为一阶极点。将其 Laurent 主部写为
\[
\zeta_{X_{A_m}}(z)\sim \frac{R_m^{\mathrm{cov}}}{1-\rho_m z},
\qquad
\zeta_{Y_m}(z)\sim \frac{R_m^{\mathrm{int}}}{1-\rho_m z},
\qquad (z\to z_\star),
\]
并定义常数比值
\[
\boxed{\ \eta_m:=\frac{R_m^{\mathrm{int}}}{R_m^{\mathrm{cov}}}\ }.
\]
则 \(\eta_m\) 仅由 \((A_m,B_m)\) 的解析结构决定，因而是一个常数级的“折叠核指纹”。并且由命题 \ref{prop:periodic-count-sandwich} 的上界可推出
\[
0<\eta_m\le 1.
\]
\leanverified{paper\_Ym\_zeta\_residue\_ratio}
\end{proposition}
\begin{definition}[歧义壳的 valuation 同步骨架]\label{distill:khovanskii_newton_polyhedra_fewnomials_and_tropical_skeletons-support_to_skeleton_reduction-ambiguity-valuation-skeleton-def}
设 \(m\ge3\)，并沿用 \(H_m^{\mathrm{det}}\)、\(\mathcal S_m^{\mathrm{det}}\)、\(\mathcal S_m^{\mathrm{amb}}\) 与 \(H_m^{\mathrm{amb}}\) 的记号。对 \(S\in\mathcal S_m^{\mathrm{amb}}\)，定义其歧义寿命
\[
 h_m(S):=\max\{\ell\ge0:\ \exists\ S=S_0\to S_1\to\cdots\to S_\ell\ \text{为 }H_m^{\mathrm{amb}}\text{ 中的有向路径}\}.
\]
由定理 \ref{thm:Ym-ambiguity-shell-dag}，该最大值有限。若 \(h_m(S)>0\)，定义从 \(S\) 出发的 valuation-selected 面为
\[
 \operatorname{Face}_m(S):=\{\,e:S\to T\text{ 为 }H_m^{\mathrm{amb}}\text{ 的边}:\ h_m(T)=h_m(S)-1\,\};
\]
若 \(h_m(S)=0\)，令 \(\operatorname{Face}_m(S)=\varnothing\)。由所有顶点 \(\mathcal S_m^{\mathrm{amb}}\)、权重 \(h_m\)、以及边集 \(\bigcup_S\operatorname{Face}_m(S)\) 组成的有限加权有向复形记为 \(\mathfrak T_m^{\mathrm{amb}}\)。
\end{definition}

\begin{theorem}[最坏歧义块的 valuation-selected 骨架分类]\label{distill:khovanskii_newton_polyhedra_fewnomials_and_tropical_skeletons-support_to_skeleton_reduction-valuation-selected-ambiguity-synchronization-skeleton}
设 \(m\ge3\)。则 \(\mathfrak T_m^{\mathrm{amb}}\) 是有限 DAG，且沿每条边 \(S\to T\) 有
\[
 h_m(T)=h_m(S)-1.
\]
特别地，每条 \(\mathfrak T_m^{\mathrm{amb}}\)-路径长度至多 \(m-1\)。此外，对任意 \(S\in\mathcal S_m^{\mathrm{amb}}\) 与任意允许输出块 \(a_0\cdots a_{L-1}\)，令
\[
 S_{i+1}:=\delta_m(S_i,a_i),\qquad S_0:=S.
\]
若 \(L=h_m(S)\)，则以下两条件等价：
\begin{enumerate}
  \item \(S_0,S_1,\ldots,S_L\) 全部属于 \(\mathcal S_m^{\mathrm{amb}}\)；
  \item 对每个 \(0\le i<L\)，边 \(S_i\to S_{i+1}\) 属于 \(\operatorname{Face}_m(S_i)\)。
\end{enumerate}
因此，任何达到最大同步延迟的歧义坏块都被唯一地压缩到有限骨架 \(\mathfrak T_m^{\mathrm{amb}}\) 上的一条 valuation-selected 路径；非选中边必然严格缩短剩余歧义寿命，不能出现在最坏块中。
\end{theorem}
\begin{proof}
定理 \ref{thm:Ym-ambiguity-shell-dag} 已说明 \(H_m^{\mathrm{amb}}\) 为有限 DAG，故从任一 \(S\in\mathcal S_m^{\mathrm{amb}}\) 出发的有向路径长度有最大值，\(h_m(S)\) 良定义。若 \(S\to T\) 是 \(H_m^{\mathrm{amb}}\) 中任意边，则从 \(T\) 出发的任一路径可在前端接上该边而成为从 \(S\) 出发的路径，所以
\[
 h_m(S)\ge h_m(T)+1.
\]
由 \(\operatorname{Face}_m(S)\) 的定义，\(\mathfrak T_m^{\mathrm{amb}}\) 中的每条边均满足 \(h_m(T)=h_m(S)-1\)。于是 \(h_m\) 沿 \(\mathfrak T_m^{\mathrm{amb}}\) 严格下降，故 \(\mathfrak T_m^{\mathrm{amb}}\) 不含有向环。

再证长度上界。任取 \(S\in\mathcal S_m^{\mathrm{amb}}\)。由于 \(\mathcal S_m^{\mathrm{det}}\) 按定义由初始候选集 \(V_m\) 可达，存在从 \(V_m\) 到 \(S\) 的某条有限路径。若从 \(S\) 出发还存在长度 \(\ell\ge m\) 的路径完全停留在 \(\mathcal S_m^{\mathrm{amb}}\)，则把该路径接在一条到达 \(S\) 的路径之后，并从其中取连续长度为 \(m\) 的后缀，可得到一段从某个可达候选态继续读入 \(m\) 个符号后仍处于非单点态的确定化滤波。将此前缀并入从 \(V_m\) 出发的读入过程，便得到一条从初始状态出发、在某个时刻 \(t\ge m\) 仍停留于 \(\mathcal S_m^{\mathrm{amb}}\) 的滤波轨道。这与定理 \ref{thm:Ym-ambiguity-shell-dag} 中“对所有 \(t\ge m\) 有 \(|S_t|=1\)”矛盾。因此 \(h_m(S)\le m-1\)，从而每条 \(\mathfrak T_m^{\mathrm{amb}}\)-路径长度至多 \(m-1\)。

最后证明等价性。设 \(L=h_m(S)\)，并假定 \(S_0,S_1,\ldots,S_L\) 全部属于 \(\mathcal S_m^{\mathrm{amb}}\)。对任意 \(0\le i<L\)，尾段
\[
 S_{i+1}\to S_{i+2}\to\cdots\to S_L
\]
说明 \(h_m(S_{i+1})\ge L-i-1\)。另一方面，前段
\[
 S_i\to S_{i+1}\to\cdots\to S_L
\]
 已经是从 \(S_i\) 出发、长度 \(L-i\) 的歧义路径；若 \(h_m(S_{i+1})>L-i-1\)，则可把从 \(S_{i+1}\) 出发的更长歧义路径接在边 \(S_i\to S_{i+1}\) 后，得到从 \(S_i\) 出发长于 \(L-i\) 的歧义路径，进而接在 \(S_0\to\cdots\to S_i\) 后得到从 \(S\) 出发长于 \(L\) 的歧义路径，矛盾。因此 \(h_m(S_{i+1})=L-i-1=h_m(S_i)-1\)，故 \(S_i\to S_{i+1}\in\operatorname{Face}_m(S_i)\)。这证明了 (1) 推出 (2)。

反过来，若每条边 \(S_i\to S_{i+1}\) 都属于 \(\operatorname{Face}_m(S_i)\)，则按定义这些边均为 \(H_m^{\mathrm{amb}}\) 中的边，故所有 \(S_i\) 均属于 \(\mathcal S_m^{\mathrm{amb}}\)。于是 (2) 推出 (1)。

该证明同时给出坏例子的结构分类：最大延迟块不能任意游走于整个确定化图，而只能沿权重 \(h_m\) 每步下降一的 valuation-selected 面前进；任何偏离该面选择的边都会使剩余可保持歧义的长度小于最坏值。故歧义壳在同步问题中的有效剩余对象正是有限的带权入射骨架 \(\mathfrak T_m^{\mathrm{amb}}\)。
\end{proof}

\begin{proof}
\(\eta_m\) 的不变量性来自定义：\(R_m^{\mathrm{cov}},R_m^{\mathrm{int}}\) 是两类 \(\zeta\) 在主极点处的主部系数。

由 \(\zeta\) 的周期点展开式（定义 \(\zeta_{Y_m}(z)=\exp(\sum_{n\ge 1}\#\mathrm{Fix}(\sigma^n|_{Y_m})z^n/n)\)）与命题 \ref{prop:periodic-count-sandwich} 的逐 \(n\) 上界
\(\#\mathrm{Fix}(\sigma^n|_{Y_m})\le \#\mathrm{Fix}(\sigma^n|_{X_{A_m}})\)
得到对任意 \(0\le z<z_\star\) 的比较：
\[
\log\zeta_{Y_m}(z)
=\sum_{n\ge 1}\frac{\#\mathrm{Fix}(\sigma^n|_{Y_m})}{n}z^n
\le
\sum_{n\ge 1}\frac{\#\mathrm{Fix}(\sigma^n|_{X_{A_m}})}{n}z^n
=\log\zeta_{X_{A_m}}(z),
\]
从而 \(\zeta_{Y_m}(z)\le \zeta_{X_{A_m}}(z)\)。令 \(z\uparrow z_\star\) 并乘以 \(1-\rho_m z\)，结合两者主部展开即得 \(R_m^{\mathrm{int}}\le R_m^{\mathrm{cov}}\)，从而 \(0<\eta_m\le 1\)。证毕。
\end{proof}

右 Fischer cover 的存在性、唯一性与最小性可由“确定化 + follower 最小化”构造给出（参见 \cite{LindMarcus1995,Kitchens1998SymbolicDynamics}）。


\subsubsection{邻接关系、拓扑熵与 Parry（最大熵）测度}\label{subsubsec:Ym-adjacency-entropy-parry}
取一个右解析且 follower-separated 的表示（例如上面的 Fischer cover），记为
\[
H_m=(\mathcal{S}_m,\mathcal{E}_m,\lambda_m),
\]
其中每条边 $e\in\mathcal{E}_m$ 具有起点 $s(e)\in\mathcal{S}_m$、终点 $t(e)\in\mathcal{S}_m$ 与标记 $\lambda_m(e)\in X_m$。

\begin{definition}[覆盖 SFT 的邻接矩阵（边移位）]\label{def:Phi_m-adjacency}
令 $X_{A_m}$ 为图 $H_m$ 的\textbf{边移位}（把边当作字母的 SFT）。其邻接矩阵 $A_m$ 为 $0$--$1$ 矩阵，索引集为 $\mathcal{E}_m$，并由
\[
(A_m)_{e,e'}:=\mathbf{1}\{\,t(e)=s(e')\,\}
\]
给出。这一邻接关系是完全显式的：它只依赖于 $H_m$ 的端点相容性。
\end{definition}

\begin{proposition}[拓扑熵的可计算公式]\label{prop:Phi_m-entropy}
设 $H_m$ 不可约（或取其最大熵不可约分量），则
\[
h_{\mathrm{top}}(Y_m)=h_{\mathrm{top}}(X_{A_m})=\log\rho(A_m).
\]
\end{proposition}

\begin{proof}
$X_{A_m}$ 通过逐边取标记映射到 $Y_m$（滑动块因子）。在 Fischer cover 上该因子码有限对一，因此熵保持；而 SFT 的熵为 $\log\rho(A_m)$。参见 \cite{LindMarcus1995,Kitchens1998SymbolicDynamics}。
\leanverified{dfa\_density\_dichotomy\_golden\_mean}
\leanverified{paper\_Phi\_m\_entropy}
\leanverified{paper\_phi\_m\_entropy}
\end{proof}

\leanverified{paper\_phi\_m\_parry}
\begin{proposition}[Parry 最大熵测度（覆盖 \texorpdfstring{$X_{A_m}$}{XAm} 上）]\label{prop:Phi_m-parry}
令 $\lambda=\rho(A_m)$，取 Perron--Frobenius 左/右特征向量 $l,r>0$ 满足
\(
l^{\mathsf T}A_m=\lambda l^{\mathsf T},\ A_m r=\lambda r
\)
并规范化为 $l^{\mathsf T}r=1$。定义状态权重
\[
\pi(e):=l_e r_e\qquad (e\in\mathcal{E}_m),
\]
以及转移概率
\[
P_{e\to e'}:=\frac{(A_m)_{e,e'}\,r_{e'}}{\lambda\, r_e}.
\]
则由 $(\pi,P)$ 给出的平稳马尔可夫测度是 $X_{A_m}$ 上唯一的最大熵测度（Parry 测度），记为 $\mu_{A_m}^{\mathrm{Parry}}$。
其对边柱集 $[e_0e_1\cdots e_{n-1}]$ 的概率满足
\[
\mu_{A_m}^{\mathrm{Parry}}([e_0\cdots e_{n-1}])=\pi(e_0)\prod_{j=0}^{n-2} P_{e_j\to e_{j+1}}.
\]
\end{proposition}
\leanverified{paper\_phi\_m\_parry\_measure}

\begin{proposition}[把 Parry 测度推前到 \texorpdfstring{$Y_m$}{Ym}]\label{prop:Phi_m-parry-pushforward}
令 $\pi_m:X_{A_m}\to Y_m$ 为逐边取标记的因子映射
\[
(\pi_m((e_t)_{t\in\ZZ}))_t:=\lambda_m(e_t)\in X_m.
\]
则 $\nu_m:=({\pi_m})_*\mu_{A_m}^{\mathrm{Parry}}$ 是 $Y_m$ 的最大熵测度（当 $Y_m$ 不可约时唯一）。对任意输出柱集 $[a_0\cdots a_{n-1}]$（$a_j\in X_m$），有完全显式的求和公式
\[
\nu_m([a_0\cdots a_{n-1}])
=
\sum_{\substack{e_0,\dots,e_{n-1}\in\mathcal{E}_m\\ \lambda_m(e_j)=a_j}}
\mu_{A_m}^{\mathrm{Parry}}([e_0\cdots e_{n-1}]),
\]
其中 $\mu_{A_m}^{\mathrm{Parry}}([e_0\cdots e_{n-1}])$ 由命题 \ref{prop:Phi_m-parry} 的 PF 数据给出。
\end{proposition}
\leanverified{paper\_phi\_m\_parry\_pushforward}

\begin{proof}
Fischer cover 的 Parry 测度推前为 sofic shift 的最大熵测度；柱集公式只是“推前测度 = 对 lift 求和”的展开。参见 \cite{LindMarcus1995,Kitchens1998SymbolicDynamics}。
\end{proof}
\leanverified{paper\_Phi\_m\_parry\_pushforward}

\begin{remark}[可复现的邻接关系与熵计算（脚本）]\label{rem:phi_m_entropy_script}
上文给出的是“如何从右解析表示得到邻接关系/熵/Parry 测度”的抽象接口。为回应“需要可计算邻接关系与显式熵值”的审计需求，仓库提供脚本 \texttt{scripts/exp\_phi\_m\_sofic\_entropy.py}：它从定义 \ref{def:G_m} 的标记图 $G_m$ 出发，经子集确定化与右语言最小化构造一个右解析表示，并输出其（最大熵本质分量上的）邻接计数矩阵的谱半径，从而得到 $h_{\mathrm{top}}(Y_m)$ 的数值值。表 \ref{tab:phi_m_sofic_entropy} 为该脚本的默认输出摘要（见 \texttt{artifacts/export/phi\_m\_sofic\_entropy.csv} 与 \texttt{sections/generated/tab\_phi\_m\_sofic\_entropy.tex}）。
\end{remark}

\begin{table}[H]
\centering
\scriptsize
\setlength{\tabcolsep}{4pt}
\caption{Sofic 图像 $Y_m$ 的右解析最小化表示规模与拓扑熵（脚本生成）。}
\label{tab:phi_m_sofic_entropy}
\begin{tabular}{rrrrrr}
\toprule
$m$ & det states & min states & ess states & det edges & $h_{\mathrm{top}}$\\
\midrule
2 & 3 & 4 & 2 & 7 & 0.693147\\
3 & 7 & 8 & 4 & 20 & 0.693147\\
4 & 16 & 17 & 8 & 50 & 0.693147\\
5 & 32 & 33 & 16 & 95 & 0.693147\\
6 & 69 & 70 & 32 & 204 & 0.693147\\
7 & 142 & 143 & 64 & 404 & 0.693147\\
8 & 296 & 297 & 128 & 827 & 0.693147\\
\bottomrule
\end{tabular}
\end{table}


\begin{remark}[为什么在本文中用“覆盖的 \texorpdfstring{$\zeta$}{zeta}”作为接口]\label{rem:sofic-zeta-vs-cover}
需要强调：sofic shift 的\emph{内禀} Artin--Mazur \(\zeta_{Y_m}\) 虽然是有理函数（可由自动机/cover 理论计算），但它\textbf{并不}在有限对一因子下保持不变，因此不能简单地由 \(\det(I-zA_m)\) 直接给出。本文把 \(\det(I-zA)\) 与 Fredholm 行列式作为“语法层可计算接口”时，所对应的是覆盖 SFT（即 Fischer cover）的 \(\zeta\)，这也是 Cuntz--Krieger/群胚 C*-代数与 PF/Parry 理论的对接对象。
\end{remark}


\subsubsection{滑动块因子 \texorpdfstring{$Y_m$}{Ym} 的谱/\texorpdfstring{$\zeta$}{zeta}：cover 与内禀对象的精确分离}\label{subsec:Ym-zeta-spectrum}
上一段已解释“为何本文主要使用 cover 的 \(\zeta\)”。为满足“折叠诱导滑动块因子的谱/\(\zeta\) 计算应明确到可执行层”的可核验性要求，本小节把两类对象的计算链条明确列出，并给出 \(Y_m\) 的内禀 \(\zeta_{Y_m}\) 的\textbf{可计算}表述，同时强调其与覆盖 SFT 行列式 \(\det(I-zA_m)\) 的定义差异。

\paragraph{cover-SFT 的谱与 \texorpdfstring{$\zeta$}{zeta}}
取 \(Y_m\) 的右 Fischer cover（或任一 follower-separated 右解析 cover）\(H_m\)，按定义 \ref{def:Phi_m-adjacency} 得到边移位 \(X_{A_m}\) 的邻接矩阵 \(A_m\)。则：
\begin{itemize}
  \item \textbf{谱/熵：} \(h_{\mathrm{top}}(Y_m)=\log\rho(A_m)\)（命题 \ref{prop:Phi_m-entropy}）；
  \item \textbf{最大熵测度：} Parry 测度由 \(A_m\) 的 Perron--Frobenius 左/右本征向量显式给出（命题 \ref{prop:Phi_m-parry}），并可推前得到 \(Y_m\) 的最大熵测度（命题 \ref{prop:Phi_m-parry-pushforward}）；
  \item \textbf{\(\zeta\)：} 覆盖 SFT 的 Artin--Mazur \(\zeta\) 满足
  \[
  \zeta_{X_{A_m}}(z)=\frac{1}{\det(I-zA_m)},
  \qquad
  \#\mathrm{Fix}(\sigma^n|_{X_{A_m}})=\mathrm{Tr}(A_m^n),
  \]
  并因此
  \(
  \#\mathrm{Fix}(\sigma^n|_{X_{A_m}})=\rho(A_m)^n+O(\rho_2^n)
  \)
  （\(\rho_2<\rho(A_m)\) 为次大谱半径，取不可约非周期分量时成立）。
\end{itemize}
这套链条之所以“可审计”，在于一切输入仅是有限图/矩阵：所有数值不变量（熵、主特征值、留数常数、primitive 轨道谱等）都可由线性代数直接复现。

\paragraph{sofic 因子 \texorpdfstring{$Y_m$}{Ym} 的内禀 \texorpdfstring{$\zeta_{Y_m}$}{zetaYm}：可计算但不等于 cover 行列式}
内禀 \(\zeta_{Y_m}\) 由周期点计数定义：
\[
\zeta_{Y_m}(z):=\exp\left(\sum_{n\ge 1}\frac{\#\mathrm{Fix}(\sigma^n|_{Y_m})}{n}z^n\right).
\]
由于 \(Y_m\) 是 sofic shift（有限标记图像；见 \cite{Weiss1973SFTandSofic,CovenPaul1975SoficSystems,LindMarcus1995,Kitchens1998SymbolicDynamics}），其周期点集是正则对象，因此 \(\zeta_{Y_m}\) 为有理函数（参见 \cite{Manning1971RationalZeta}）。关键是：\(\#\mathrm{Fix}(\sigma^n|_{Y_m})\) 计数的是\textbf{标记序列}的周期点，而 \(\mathrm{Tr}(A_m^n)\) 计数的是 cover 上\textbf{边序列}的周期点；当标记因子不是一一对应时，这两者不同。

\paragraph{一个可执行的计数口径（转移半群自动机）}
给定任一右解析表示（例如 $H_m^{\mathrm{det}}$ 或 Fischer cover），可把每个字母 $a\in X_m$ 诱导为状态集上的部分转移 $\tau_a$，并由可达的部分映射半群构造一个有限自动机；其邻接矩阵记为 $B_m$。则
\[
\#\mathrm{Fix}(\sigma^n|_{Y_m})=\mathrm{Tr}(B_m^n),\qquad
\zeta_{Y_m}(z)=\frac{1}{\det(I-zB_m)}.
\]
该口径见 \cite{Nasu1995TextileSystems}，以及 \cite{LindMarcus1995,Kitchens1998SymbolicDynamics} 关于 sofic/cover 的章节。

\begin{proposition}[内禀 \texorpdfstring{$\zeta_{Y_m}$}{zetaYm} 的主极点与熵：与 cover 同一指数尺度]\label{prop:Ym-zeta-leading-pole}
\leanverified{paper\_Ym\_zeta\_leading\_pole}
设 \(Y_m\) 不可约（或取其最大熵本质分量）。则 \(Y_m\) 的周期点增长率满足
\[
\limsup_{n\to\infty}\frac{1}{n}\log\#\mathrm{Fix}(\sigma^n|_{Y_m})
\ =\ h_{\mathrm{top}}(Y_m),
\]
从而内禀 \(\zeta_{Y_m}(z)\) 的收敛半径为
\[
R_{Y_m}=e^{-h_{\mathrm{top}}(Y_m)}.
\]
另一方面，cover 接口给出 \(h_{\mathrm{top}}(Y_m)=\log\rho(A_m)\)（命题 \ref{prop:Phi_m-entropy}），因此
\[
R_{Y_m}=\rho(A_m)^{-1}.
\]
换言之：尽管 \(\zeta_{Y_m}\) 与 \(\zeta_{X_{A_m}}=1/\det(I-zA_m)\) 一般不相等，但二者在主极点模长（即指数尺度）上由同一熵决定。
\end{proposition}

\begin{proof}
sofic shift 的周期点计数指数增长率与拓扑熵一致，参见 \cite{Weiss1973SFTandSofic,CovenPaul1975SoficSystems,LindMarcus1995,Kitchens1998SymbolicDynamics}。
另一方面由上段口径 \(\#\mathrm{Fix}(\sigma^n|_{Y_m})=\mathrm{Tr}(B_m^n)\)，故收敛半径为 \(\rho(B_m)^{-1}=e^{-h_{\mathrm{top}}(Y_m)}\)。再由命题 \ref{prop:Phi_m-entropy} 的 \(h_{\mathrm{top}}(Y_m)=\log\rho(A_m)\) 得 \(R_{Y_m}=\rho(A_m)^{-1}\)。
\end{proof}

\begin{proposition}[周期点计数的显式夹逼：cover \texorpdfstring{$\Rightarrow$}{=>} sofic]\label{prop:periodic-count-sandwich}
取 \(Y_m\) 的一个右解析标记图表示 \(H_m=(\mathcal{S}_m,\mathcal{E}_m,\lambda_m)\)（例如定义 \ref{def:Phi_m-right-resolving} 的 \(H_m^{\mathrm{det}}\)，或进一步最小化得到的 Fischer cover）。令 \(X_{A_m}\) 为其边移位（定义 \ref{def:Phi_m-adjacency}），并令 \(\pi_m:X_{A_m}\to Y_m\) 为逐边取标记的因子映射（定义同命题 \ref{prop:Phi_m-parry-pushforward}）。则对任意 \(n\ge 1\) 有
\[
\#\mathrm{Fix}(\sigma^n|_{Y_m})
\ \le\
\#\mathrm{Fix}(\sigma^n|_{X_{A_m}})
\ =\
\mathrm{Tr}(A_m^n),
\]
并且还存在完全显式的下界
\[
\#\mathrm{Fix}(\sigma^n|_{Y_m})
\ \ge\
\frac{1}{|\mathcal{S}_m|}\,\#\mathrm{Fix}(\sigma^n|_{X_{A_m}})
\ =\
\frac{1}{|\mathcal{S}_m|}\,\mathrm{Tr}(A_m^n).
\]
因此 \(\#\mathrm{Fix}(\sigma^n|_{Y_m})\) 与 \(\mathrm{Tr}(A_m^n)\) 在指数尺度上一致（差异至多为常数因子），并与命题 \ref{prop:Ym-zeta-leading-pole} 的“同一主极点模长”结论一致。
\leanverified{paper\_Ym\_periodic\_count\_sandwich}
\end{proposition}

\begin{proof}
上界源于因子映射 \(\pi_m\) 将 \(X_{A_m}\) 的 \(n\)-周期点映为 \(Y_m\) 的 \(n\)-周期点，故 \(\#\mathrm{Fix}(\sigma^n|_{Y_m})\le \#\mathrm{Fix}(\sigma^n|_{X_{A_m}})\)。
下界由右解析性给出前像多重度的统一控制：对任意 $y\in Y_m$，其 lift 由“起始状态 $S\in\mathcal{S}_m$ + 标记序列”唯一确定，故 $|\pi_m^{-1}(y)|\le |\mathcal{S}_m|$。对 $n$-周期点求和即得
\(\#\mathrm{Fix}(X_{A_m})\le |\mathcal{S}_m|\#\mathrm{Fix}(Y_m)\)。
最后 \(\#\mathrm{Fix}(X_{A_m})=\mathrm{Tr}(A_m^n)\) 由 SFT 周期点计数恒等式给出（\cite{LindMarcus1995}）。
\end{proof}

\begin{corollary}[长度-\texorpdfstring{$n$}{n} 周期压缩因子与常数级缺陷]\label{cor:Ym-periodic-compression-ratio}
沿用命题 \ref{prop:periodic-count-sandwich} 的记号。对任意满足 \(\#\mathrm{Fix}(\sigma^n|_{Y_m})>0\) 的 \(n\ge 1\)，定义长度-\(n\) 周期压缩因子
\[
R_m(n):=\frac{\#\mathrm{Fix}(\sigma^n|_{X_{A_m}})}{\#\mathrm{Fix}(\sigma^n|_{Y_m})}
=\frac{\mathrm{Tr}(A_m^n)}{\#\mathrm{Fix}(\sigma^n|_{Y_m})}.
\]
则对所有这样的 \(n\) 都有统一界
\[
\boxed{\ 1\ \le\ R_m(n)\ \le\ |\mathcal{S}_m|\ }.
\]
特别地，定义常数级“sofic 缺陷常数”
\[
\delta_m:=\limsup_{n\to\infty}\log R_m(n),
\]
则
\[
0\ \le\ \delta_m\ \le\ \log|\mathcal{S}_m|.
\]
该量度在指数尺度上不改变熵/主极点模长，但以可计算的常数级方式量化“cover 周期点被因子码压缩”的强度。
\leanverified{compressionRatio}
\leanverified{one\_le\_div\_of\_le}
\leanverified{div\_le\_of\_le\_mul}
\leanverified{compressionRatio\_bound}
\leanverified{paper\_Ym\_periodic\_compression\_ratio}
\end{corollary}

\begin{proof}
由命题 \ref{prop:periodic-count-sandwich} 得
\(
\#\mathrm{Fix}(\sigma^n|_{Y_m})\le \mathrm{Tr}(A_m^n)
\)
与
\(
\#\mathrm{Fix}(\sigma^n|_{Y_m})\ge |\mathcal{S}_m|^{-1}\mathrm{Tr}(A_m^n)
\)。
对两式同除以 \(\#\mathrm{Fix}(\sigma^n|_{Y_m})\) 即得 \(1\le R_m(n)\le |\mathcal{S}_m|\)。
再取对数并取 \(\limsup_{n\to\infty}\) 得 \(0\le \delta_m\le \log|\mathcal{S}_m|\)。证毕。
\end{proof}

\begin{remark}[机制：窗口层多对一不等于时间层熵损失]\label{rem:Ym-no-entropy-loss-mechanism}
命题 \ref{prop:periodic-count-sandwich} 给出一个精确的“常数因子压缩”律：尽管窗口层折叠 \(\Fold_m:\Omega_m\to X_m\) 典型为多对一，从而序列级因子 \(\Phi_m\) 也可能把多个 cover 周期轨道折叠到同一标记周期点，但在时间尺度的周期点增长上，这种压缩最多只体现为 \(|\mathcal{S}_m|\) 的常数因子，而不会改变指数斜率（熵/主极点模长）。

更具体地，由
\[
\frac{1}{|\mathcal{S}_m|}\Tr(A_m^n)\le \#\mathrm{Fix}(\sigma^n|_{Y_m})\le \Tr(A_m^n)
\]
立刻推出
\[
\limsup_{n\to\infty}\frac{1}{n}\log\#\mathrm{Fix}(\sigma^n|_{Y_m})
=\limsup_{n\to\infty}\frac{1}{n}\log\Tr(A_m^n)
=\log\rho(A_m),
\]
即 \(h_{\mathrm{top}}(Y_m)=\log\rho(A_m)\)（命题 \ref{prop:Ym-zeta-leading-pole}）。
因此“丢失的信息”不是被指数压扁为熵损失，而是迁移进右解析状态集 \(\mathcal{S}_m\) 所携带的记忆结构：窗口层的多重度对应的是有限记忆的重编码，而非指数增长率的衰减。

在本文的具体模型（Zeckendorf 折叠）中，上述机制还能进一步“封口”：定理 \ref{thm:Ym-ambiguity-shell-dag} 表明 $m\ge 3$ 时非单点歧义层是纯瞬态 DAG，完全不承载周期点；其造成的时间代价不是熵损失，而是一个可计入的有限同步延迟 \(D_{\mathrm{sync}}(m)=m-1\)（推论 \ref{cor:Ym-sync-delay}）。
\end{remark}

\begin{corollary}[单点本质态 \texorpdfstring{$\Rightarrow$}{=>} 有限延迟右解码]\label{cor:Ym-singleton-essential-finite-delay}
沿用定义 \ref{def:G_m} 与定义 \ref{def:Phi_m-right-resolving} 的记号。设通过“子集确定化 + 右语言最小化”得到的一个右解析最小表示（例如右 Fischer cover 或其等价的最小确定性表示）满足如下\textbf{单点本质态}条件：
\begin{equation}\label{eq:Ym-singleton-essential-singletons}
\text{其最大熵本质强连通分量的状态可辨识为全部单点子集 }\ \{\{v\}:v\in V_m\}.
\end{equation}
则在该本质分量上成立：
\begin{enumerate}
  \item \textbf{（无标记碰撞）}对任意 \(v\in V_m\)，两条出边 \((v,0)\)、\((v,1)\) 的标记不同：
  \[
  \Fold_m(v0)\neq \Fold_m(v1).
  \]
  \item \textbf{（逐步解码）}存在一个有限状态右逆转导器（解码器），其状态可取为 \(V_m\)。当处于状态 \(v_t\in V_m\) 且读入输出符号 \(a_t\in X_m\) 时，唯一确定 \(b_t\in\{0,1\}\) 使
  \[
  \Fold_m(v_t b_t)=a_t,
  \]
  并更新 \(v_{t+1}=v_t\) 去首加尾 \((b_t)\)。因此一旦输出轨道进入上述单点本质分量，输入比特序列从该时刻起可\textbf{零额外预期性}在线恢复。
\end{enumerate}
\end{corollary}
\begin{proof}
\textbf{(1)} 在右解析确定化表示中，对任意单点状态 \(\{v\}\) 与输出符号 \(a\)，后继由
\(
\delta_m(\{v\},a)
\)
给出。若存在 \(\Fold_m(v0)=\Fold_m(v1)=a\)，则从 \(\{v\}\) 读入 \(a\) 时两条不同的 \(b\)-边都会被纳入，从而
\(
\delta_m(\{v\},a)
\)
将包含至少两个顶点，迫使确定化状态离开单点类；这与“单点态在本质强连通分量上封闭”（条件 \eqref{eq:Ym-singleton-essential-singletons}）矛盾。故 \(\Fold_m(v0)\neq \Fold_m(v1)\)。

\textbf{(2)} 由 (1) 可知对每个 \(v_t\) 与允许的输出 \(a_t\)，在 \(b_t\in\{0,1\}\) 中恰有一个满足 \(\Fold_m(v_t b_t)=a_t\)。更新规则 \(v_{t+1}=v_t\) 去首加尾由定义 \ref{def:G_m} 的边结构给出，因此该过程在单点本质分量内闭合并确定一个有限状态转导器；其逐步输出 \(b_t\) 给出一条右逆（在该分量上逐时刻反演标记）。证毕。
\end{proof}
\leanverified{paper\_Ym\_singleton\_essential\_finite\_delay}

\begin{corollary}[歧义壳的熵严格低于 \texorpdfstring{$\log 2$}{log 2}]\label{cor:Ym-ambiguity-shell-low-entropy}
在推论 \ref{cor:Ym-singleton-essential-finite-delay} 的条件下，设右解析最小表示的状态图按强连通分量分解，其中最大熵本质分量为单点态分量（其熵为 \(\log 2\)，参见表 \ref{tab:phi_m_sofic_entropy} 的数值证书）。
令 \(Y_m^{\mathrm{amb}}\subset Y_m\) 为那些在该右解析最小表示中\textbf{完全不进入}单点本质分量的标记序列所构成的 sofic 子系统（等价地：其 lift 的当前候选集合在一侧读入过程中始终为非单点）。
在本文的 Zeckendorf 折叠模型中，当 \(m\ge 3\) 时由定理 \ref{thm:Ym-ambiguity-shell-dag} 进一步得到 \(Y_m^{\mathrm{amb}}=\varnothing\)，并参见备注 \ref{rem:Ym-amb-empty}。下述推论刻画一般 sofic 场景下歧义壳的严格低熵结构，并在本模型中作为一致性检查保留。
则有严格不等式
\[
h_{\mathrm{top}}(Y_m^{\mathrm{amb}})<\log 2.
\]
因此“窗口层的多对一”在时间层不会产生指数级丢熵；它只对应一个\textbf{严格低熵}的歧义壳层，而最大熵部分在有限状态记忆中可右解码。
\leanverified{paper\_Ym\_ambiguity\_shell\_low\_entropy}
\end{corollary}
\begin{proof}
由假设，单点态本质分量承载 \(Y_m\) 的最大熵，其邻接计数矩阵的谱半径为 \(2\)，故 \(h_{\mathrm{top}}=\log 2\)。
而 \(Y_m^{\mathrm{amb}}\) 由\emph{剔除该本质分量后的其余强连通分量}所呈现：这些分量的谱半径都严格小于最大值 \(2\)（否则将出现另一个同熵本质分量，与“最大熵分量由单点态承载”的设定矛盾）。
因此 \(Y_m^{\mathrm{amb}}\) 的拓扑熵严格小于 \(\log 2\)。证毕。
\end{proof}

\begin{corollary}[歧义壳指数稀薄：块计数比例具有显式指数缺口]\label{cor:Ym-ambiguity-shell-exponentially-thin}
令 \(N_{\mathrm{amb}}(n)\) 表示歧义壳子系统 \(Y_m^{\mathrm{amb}}\) 允许的长度 \(n\) 块数（即其语言在长度 \(n\) 的不同子词数）。则
\[
N_{\mathrm{amb}}(n)=\exp\!\bigl(h_{\mathrm{top}}(Y_m^{\mathrm{amb}})\,n+o(n)\bigr),
\]
而由 \(h_{\mathrm{top}}(Y_m)=\log 2\) 可得 \(Y_m\) 的长度 \(n\) 块数至少为 \(2^n\)。因此存在 \(\varepsilon_m>0\) 使得
\[
\frac{N_{\mathrm{amb}}(n)}{2^n}
=
\exp\!\bigl(-\varepsilon_m\,n+o(n)\bigr),
\qquad
\varepsilon_m:=\log 2-h_{\mathrm{top}}(Y_m^{\mathrm{amb}})>0.
\]
换言之，在块计数意义下，右解码失败所对应的歧义壳并非”同阶复杂”的部分，而是带有显式指数缺口的稀薄层。
\leanverified{stable\_language\_exponentially\_sparse}
\leanverified{density\_ratio\_decreasing\_instances}
\leanverified{paper\_Ym\_ambiguity\_shell\_exponentially\_thin}
\end{corollary}

\begin{proof}
对任意子移位 \(Z\)，其拓扑熵可按块计数定义为
\(
h_{\mathrm{top}}(Z)=\lim_{n\to\infty}\frac{1}{n}\log N_Z(n)
\)
（其中 \(N_Z(n)\) 为长度 \(n\) 允许块数），从而 \(N_Z(n)=\exp(h_{\mathrm{top}}(Z)\,n+o(n))\)。对 \(Z=Y_m^{\mathrm{amb}}\) 得第一式。
\par
另一方面，\(h_{\mathrm{top}}(Y_m)=\log 2\) 等价于 \(Y_m\) 的块增长率至少为 \(2^n\)（例如取熵定义的下极限即可），故比值满足所示指数衰减，其中 \(\varepsilon_m>0\) 由推论 \ref{cor:Ym-ambiguity-shell-low-entropy} 给出。证毕。
\end{proof}

\begin{definition}[同步时间与同步尾]\label{def:Ym-sync-time}
对任意 \(y\in Y_m\)，从初始候选集 \(S_0:=V_m\) 出发递推
\[
S_{t+1}:=\delta_m(S_t,y_t)\qquad (t\ge 0).
\]
定义\textbf{同步时间}（第一次进入单点态的时刻）
\[
T_m(y):=\inf\{\,t\ge 0:\ |S_t|=1\,\}\in\NN\cup\{\infty\}.
\]
若 \(T_m(y)=\infty\)，则称 \(y\) 永不同步（对应歧义壳轨道）。
\end{definition}

\begin{remark}[统计型 \texorpdfstring{$\mathsf{NULL}$}{NULL} 的时间层刻画：未同步即无地址]\label{rem:Ym-stat-null-unsync}
在确定化表示 \(H_m^{\mathrm{det}}\) 中，状态 \(S_t\subseteq V_m\) 是“给定过去输出前缀时仍可能的内部后缀集合”。
因此：
\(|S_t|=1\) 表示此时刻的内部后缀（从而局部 lift）被输出前缀\emph{唯一}确定；
\(|S_t|>1\) 则表示存在不可消去的右解码歧义。
在“地址先于概念值”的语义约定（见第 \ref{sec:spg} 节）下，这种歧义可被等价理解为时间层的统计型 \(\mathsf{NULL}\)：输出前缀不足以生成唯一地址，因此相应读出在该时刻不定义。
同步时间 \(T_m\) 即“统计型 \(\mathsf{NULL}\) 首次消失”的时刻；而定义 \ref{def:eta_m_hat_singleton_rate} 的 \(\widehat\eta_m\) 与命题 \ref{prop:eta_m_hat_ergodic_limit} 的 \(\eta_m^{\mathrm{obs}}\) 则给出其在线可观测频率。
\end{remark}

\begin{theorem}[同步尾界：最大熵测度下的指数型同步时间预算]\label{thm:Ym-sync-tail}
\leanverified{paper\_ym\_sync\_tail}
设 \(Y_m\) 的最大熵测度为 \(\nu_m\)（命题 \ref{prop:Phi_m-parry-pushforward}）。若歧义壳子系统 \(Y_m^{\mathrm{amb}}\neq\varnothing\)，令
\[
\varepsilon_m:=\log 2-h_{\mathrm{top}}(Y_m^{\mathrm{amb}})>0
\]
为推论 \ref{cor:Ym-ambiguity-shell-exponentially-thin} 给出的熵缺口。则存在常数 \(C_m>0\) 与 \(n_0\in\NN\)，使得对所有 \(n\ge n_0\) 有
\[
\boxed{\ \nu_m(T_m>n)\ \le\ C_m\,e^{-\varepsilon_m n}\ }.
\]
从而：
\begin{enumerate}
  \item \textbf{期望同步时间有限：}\(\mathbb{E}_{\nu_m}[T_m]<\infty\)，并且规模为 \(O(1/\varepsilon_m)\)。
  \item \textbf{同步时间预算：}对任意 \(0<\delta<1\)，若
  \[
  n\ \ge\ \frac{1}{\varepsilon_m}\log\frac{C_m}{\delta},
  \]
  则 \(\nu_m(T_m>n)\le \delta\)。
\end{enumerate}
当 \(Y_m^{\mathrm{amb}}=\varnothing\) 时，\(T_m(y)\le L_m\) 对所有 \(y\in Y_m\) 成立，其中 \(L_m<\infty\) 由备注 \ref{rem:Ym-amb-empty} 给出，因此同步尾在有限处截断。
\end{theorem}

\begin{proof}
记 \(H_m^{\mathrm{det}}\) 为确定化右解析表示（定义 \ref{def:Phi_m-right-resolving}），并以 \(S_t\) 表示其状态轨道。
由命题 \ref{prop:det-candidate-set-monotone} 的单点态前向封闭性，一旦 \(|S_t|=1\) 则对所有 \(s\ge t\) 都有 \(|S_s|=1\)。因此事件 \(\{T_m>n\}\) 等价于“读入前 \(n\) 步后仍处于非单点态”，并可写成长度-\(n\) 柱集的并：
\[
\{T_m>n\}=\bigcup_{w\in\mathcal{W}_{m}^{\mathrm{unsync}}(n)}[w],
\]
其中 \(\mathcal{W}_{m}^{\mathrm{unsync}}(n)\subset X_m^n\) 为那些使得从初始状态 \(S_0=V_m\) 出发读入 \(w\) 后得到 \(|S_n|>1\) 的输出块集合。
\par
\textbf{(i) 柱集概率的统一上界。}
由命题 \ref{prop:Phi_m-parry-pushforward}，\(\nu_m\) 是覆盖 SFT 上 Parry 测度的推前。Parry 测度对任意长度-\(n\) lift 柱集具有 \(\asymp e^{-h_{\mathrm{top}}(Y_m)n}=2^{-n}\) 的 Gibbs 型界；推前至 \(Y_m\) 后仍得到统一上界：存在常数 \(C_m^{\mathrm{cyl}}>0\) 与 \(n_1\in\NN\)，使得对所有允许的长度-\(n\) 输出柱集 \([w]\)（\(n\ge n_1\)）有
\[
\nu_m([w])\le C_m^{\mathrm{cyl}}\,2^{-n}.
\]
\par
\textbf{(ii) 未同步前缀数目的指数界。}
集合 \(\mathcal{W}_{m}^{\mathrm{unsync}}(n)\) 正是 $H_m^{\mathrm{det}}$ 的非单点态子图 $H_m^{\mathrm{amb}}$ 中、从初始状态出发长度为 \(n\) 的标记路径的标记集合。
由于 $H_m^{\mathrm{amb}}$ 为有限图，其长度-\(n\) 标记数满足
\[
|\mathcal{W}_{m}^{\mathrm{unsync}}(n)|
\ \le\
C_m^{\mathrm{lang}}\exp\!\bigl(h_{\mathrm{top}}(Y_m^{\mathrm{amb}})\,n\bigr)
\]
对某个常数 \(C_m^{\mathrm{lang}}>0\) 与所有充分大 \(n\) 成立（指数率由歧义壳本质分量的谱半径给出，等于 \(h_{\mathrm{top}}(Y_m^{\mathrm{amb}})\)）。
\par
将 (i)(ii) 代入并用并集上界，得对充分大 \(n\)
\[
\begin{aligned}
\nu_m(T_m>n)
&\le \sum_{w\in\mathcal{W}_{m}^{\mathrm{unsync}}(n)}\nu_m([w])\\
&\le |\mathcal{W}_{m}^{\mathrm{unsync}}(n)|\cdot C_m^{\mathrm{cyl}}2^{-n}\\
&\le (C_m^{\mathrm{cyl}}C_m^{\mathrm{lang}})\exp\!\bigl(-(\log 2-h_{\mathrm{top}}(Y_m^{\mathrm{amb}}))n\bigr),
\end{aligned}
\]
即所示尾界，其中 \(\varepsilon_m=\log 2-h_{\mathrm{top}}(Y_m^{\mathrm{amb}})\)。
\par
期望界由
\(
\mathbb{E}[T_m]=\sum_{n\ge 0}\nu_m(T_m>n)
\)
与指数尾收敛得到；同步预算由解不等式 \(C_m e^{-\varepsilon_m n}\le\delta\) 得到。证毕。
\end{proof}

\begin{theorem}[同步尾的 Perron--Frobenius 尖锐渐近]\label{thm:Ym-sync-tail-pf-sharp}
沿用定义 \ref{def:Ym-sync-time}--定理 \ref{thm:Ym-sync-tail} 的记号，并假设歧义壳子系统 \(Y_m^{\mathrm{amb}}\neq\varnothing\)。
令 \(H_m^{\mathrm{amb}}\) 为定理 \ref{thm:Ym-ambiguity-shell-dag} 定义的非单点态诱导子图，记其（按边条数计数的）邻接矩阵为
\[
(A_m^{\mathrm{amb}})_{S,S'}:=\#\{\,a\in X_m:\ \delta_m(S,a)=S'\,\},\qquad
S,S'\in\mathcal{S}_m^{\mathrm{amb}}.
\]
设 \(H_m^{\mathrm{amb}}\) 的最大本质不可约分量的 Perron--Frobenius 主特征值为 \(\rho_{\mathrm{amb}}(m)>1\)，从而
\[
h_{\mathrm{top}}(Y_m^{\mathrm{amb}})=\log\rho_{\mathrm{amb}}(m),
\qquad
\varepsilon_m=\log 2-\log\rho_{\mathrm{amb}}(m).
\]
若该本质分量的周期为 \(p\ge 1\)（原始情形即 \(p=1\)），则存在常数族
\(\{c_r(m)\}_{r=0}^{p-1}\subset(0,\infty)\)
使得在最大熵测度 \(\nu_m\) 下
\[
\nu_m(T_m>n)=c_{n\bmod p}(m)\,e^{-\varepsilon_m n}\,(1+o(1)),
\qquad n\to\infty.
\]
从而极限指数存在且为
\[
\lim_{n\to\infty}\frac{1}{n}\log \nu_m(T_m>n)=-\varepsilon_m.
\]
\end{theorem}
\leanverified{paper\_ym\_sync\_tail\_pf\_sharp}

\begin{proof}
仍记 \(\mathcal{W}_{m}^{\mathrm{unsync}}(n)\) 如定理 \ref{thm:Ym-sync-tail} 的证明中所定义。
由确定性右解析性（定义 \ref{def:Phi_m-right-resolving}），当起点固定为 \(S_0=V_m\) 时，每个标记词决定唯一确定化滤波轨道，故 \(\mathcal{W}_{m}^{\mathrm{unsync}}(n)\) 与 \(H_m^{\mathrm{amb}}\) 中从 \(V_m\) 出发的长度 \(n\) 标记路径集合一一对应。
因此 \(h_{\mathrm{top}}(Y_m^{\mathrm{amb}})=\log\rho_{\mathrm{amb}}(m)\) 且 \(\varepsilon_m=\log 2-\log\rho_{\mathrm{amb}}(m)\)。
\par
下面把同步尾写成一个有限矩阵幂的加权读出。
取一个右解析且 follower-separated 的表示
\(
H_m=(\mathcal{S}_m,\mathcal{E}_m,\lambda_m)
\)
（例如右 Fischer cover），令 \(X_{A_m}\) 为其边移位、\(\mu_{A_m}^{\mathrm{Parry}}\) 为命题 \ref{prop:Phi_m-parry} 的 Parry 测度，并记 \(\nu_m=(\pi_m)_*\mu_{A_m}^{\mathrm{Parry}}\)（命题 \ref{prop:Phi_m-parry-pushforward}）。
由命题 \ref{prop:Phi_m-entropy} 及 \(h_{\mathrm{top}}(Y_m)=\log 2\)，可取 \(\rho(A_m)=2\)。
于是对任意长度-\(n\) 的 lift 边柱集 \([e_0\cdots e_{n-1}]\subset X_{A_m}\) 有 Parry 公式
\[
\mu_{A_m}^{\mathrm{Parry}}([e_0\cdots e_{n-1}])=\frac{l_{e_0}r_{e_{n-1}}}{2^{n-1}}.
\]
\par
构造同步积图 \(\widetilde H_m^{\mathrm{amb}}\)：其状态集为
\(
\widetilde{\mathcal{S}}_m^{\mathrm{amb}}:=\mathcal{E}_m\times\mathcal{S}_m^{\mathrm{amb}}
\)，并规定从 \((e,S)\) 到 \((e',S')\) 有边当且仅当
\[
t(e)=s(e'),\qquad S'=\delta_m(S,\lambda_m(e')).
\]
记其邻接矩阵为 \(\widetilde A_m^{\mathrm{amb}}\)。
取初始/末端向量 \(\alpha,\beta\) 为
\[
\alpha_{(e,S)}:=l_e\,\mathbf{1}\{S=\delta_m(V_m,\lambda_m(e))\},
\qquad
\beta_{(e,S)}:=r_e.
\]
由“推前 = 对 lift 求和”及上式的 Parry 公式，可将同步尾写为
\[
\nu_m(T_m>n)=2^{-(n-1)}\,\alpha^{\mathsf T}(\widetilde A_m^{\mathrm{amb}})^{n-1}\beta.
\]
\par
由于投影 \((e,S)\mapsto S\) 把 \(\widetilde H_m^{\mathrm{amb}}\) 有限对一下降到 \(H_m^{\mathrm{amb}}\)，两者的路径计数在指数标度上一致，故 \(\rho(\widetilde A_m^{\mathrm{amb}})=\rho_{\mathrm{amb}}(m)\)，且最大本质分量的周期仍为 \(p\)。
对不可约周期 \(p\) 的 Perron--Frobenius 渐近，存在常数族 \(\{d_r(m)\}_{r=0}^{p-1}\subset(0,\infty)\) 使得
\[
\alpha^{\mathsf T}(\widetilde A_m^{\mathrm{amb}})^{n-1}\beta
=d_{n\bmod p}(m)\,\rho_{\mathrm{amb}}(m)^{n-1}(1+o(1)).
\]
代回并把与 \(2^{-(n-1)}\) 的常数因子吸收进 \(c_r(m)\) 即得
\[
\nu_m(T_m>n)
=c_{n\bmod p}(m)\left(\frac{\rho_{\mathrm{amb}}(m)}{2}\right)^{\!n}(1+o(1))
=c_{n\bmod p}(m)e^{-\varepsilon_m n}(1+o(1)).
\]
最后取 \(\frac1n\log\) 并令 \(n\to\infty\) 得极限指数。证毕。
\end{proof}

\paragraph{从熵缺口到锐化曲线：未同步前缀的确定性计数律}
定理 \ref{thm:Ym-sync-tail}--\ref{thm:Ym-sync-tail-pf-sharp} 把“未同步即统计型 \(\mathsf{NULL}\)”的时间层语义转化为可审计的指数尺度与 Perron--Frobenius 主项。
然而，在同一确定化机制下，一个更为刚性的对象并非测度尾概率，而是\emph{未同步前缀集合本身}：
它构成一个有限状态正则语言，因此其块计数不仅具有熵缺口，还满足带可计算主常数与二阶谱隙的尖锐渐近。
该结果把“扫描更久 \(\Rightarrow\) 更清晰”的定性规律升级为一条有效可计算的“扫描--清晰度”锐化曲线。

\begin{definition}[未同步前缀集合与计数]\label{def:Ym-unsync-prefix-count}
沿用定理 \ref{thm:Ym-sync-tail} 的记号。
对任意 \(n\in\NN\)，定义未同步长度-\(n\) 前缀集合
\[
\mathcal{W}_{m}^{\mathrm{unsync}}(n)
:=\{\,w\in X_m^n:\ \text{从 }S_0=V_m\text{ 出发读入 }w\text{ 后满足 }|S_n|>1\,\},
\]
并记其计数为
\[
N_m^{\mathrm{unsync}}(n):=\#\mathcal{W}_{m}^{\mathrm{unsync}}(n).
\]
当 \(Y_m^{\mathrm{amb}}\neq\varnothing\) 时，\(\mathcal{W}_{m}^{\mathrm{unsync}}(n)\) 与 \(H_m^{\mathrm{amb}}\) 中从初始状态 \(V_m\) 出发、长度为 \(n\) 且始终停留在非单点态的标记路径一一对应。
\end{definition}

\leanverified{paper\_Ym\_unsync\_prefix\_pf\_sharp}
\begin{theorem}[歧义前缀计数的 Perron--Frobenius 锐化律：主项常数与二阶谱隙]\label{thm:Ym-unsync-prefix-pf-sharp}
假设 \(Y_m^{\mathrm{amb}}\neq\varnothing\)。
进一步假设非单点态诱导子图 \(H_m^{\mathrm{amb}}\) 存在唯一的最大谱半径强连通分量 \(C_m\)，且该分量是原始的（aperiodic/primitive）。
记 \(A_m^{\mathrm{amb}}\) 为定理 \ref{thm:Ym-sync-tail-pf-sharp} 中定义的按边条数计数的邻接矩阵，并令
\[
\rho_{\mathrm{amb}}(m):=\rho(A_m^{\mathrm{amb}}),\qquad
\rho_{\mathrm{amb},2}(m):=\max\bigl(\{\,|\lambda|:\ \lambda\in\sigma(A_m^{\mathrm{amb}}),\ |\lambda|<\rho_{\mathrm{amb}}(m)\,\}\cup\{0\}\bigr),
\]
则 \(\rho_{\mathrm{amb},2}(m)<\rho_{\mathrm{amb}}(m)\)，并存在显式常数 \(c_m>0\) 使得
\[
\boxed{\;
N_m^{\mathrm{unsync}}(n)
=c_m\,\rho_{\mathrm{amb}}(m)^{n}\ +\ O\!\bigl(\rho_{\mathrm{amb},2}(m)^{n}\bigr),
\qquad n\to\infty.\;}
\]
令 \(\varepsilon_m=\log 2-\log\rho_{\mathrm{amb}}(m)\)（见定理 \ref{thm:Ym-sync-tail-pf-sharp}）。
则有“锐化曲线”
\[
\boxed{\;
\frac{N_m^{\mathrm{unsync}}(n)}{2^n}
=c_m\,e^{-\varepsilon_m n}\ +\ O\!\Bigl(\bigl(\tfrac{\rho_{\mathrm{amb},2}(m)}{2}\bigr)^{n}\Bigr),
\qquad n\to\infty.\;}
\]
当 \(\rho_{\mathrm{amb},2}(m)>0\) 时，定义二阶谱隙指数
\[
\Delta_{\mathrm{amb}}(m):=\log\!\Bigl(\frac{\rho_{\mathrm{amb}}(m)}{\rho_{\mathrm{amb},2}(m)}\Bigr)>0,
\]
则上述误差项亦可写成 \(O\!\bigl(e^{-(\varepsilon_m+\Delta_{\mathrm{amb}}(m))n}\bigr)\)。
其中常数 \(c_m\) 具有可审计闭式：取 \(A_m^{\mathrm{amb}}\) 的 Perron--Frobenius 右/左特征向量对 \((v_m,w_m)\) 满足
\[
A_m^{\mathrm{amb}}v_m=\rho_{\mathrm{amb}}(m)v_m,\qquad
w_m^{\mathsf T}A_m^{\mathrm{amb}}=\rho_{\mathrm{amb}}(m)w_m^{\mathsf T},\qquad
w_m^{\mathsf T}v_m=1,
\]
并令 \(e_{V_m}\) 为初始状态 \(V_m\) 的坐标向量、\(\mathbf 1\) 为 \(\mathcal{S}_m^{\mathrm{amb}}\) 上的全 \(1\) 向量，则
\[
\boxed{\;
c_m=(e_{V_m}^{\mathsf T}v_m)\,(w_m^{\mathsf T}\mathbf 1).\;}
\]
\end{theorem}

\begin{proof}
由定义 \ref{def:Ym-unsync-prefix-count} 与确定性右解析性可知：长度-\(n\) 标记词 \(w\) 使滤波状态仍为非单点，当且仅当从 \(V_m\) 出发沿 \(H_m^{\mathrm{amb}}\) 读入 \(w\) 得到一条长度-\(n\) 的标记路径。
按定理 \ref{thm:Ym-sync-tail-pf-sharp} 的邻接矩阵定义，\(A_m^{\mathrm{amb}}\) 的条目正是一步标记选择的多重度；因此块计数满足矩阵幂公式
\[
N_m^{\mathrm{unsync}}(n)=e_{V_m}^{\mathsf T}(A_m^{\mathrm{amb}})^n\mathbf 1.
\]
\par
在“唯一最大谱半径强连通分量且原始”的假设下，\(\rho_{\mathrm{amb}}(m)\) 是 \(A_m^{\mathrm{amb}}\) 的简单特征值且为谱圆上的唯一点谱；
从而存在 \(0\le \rho_{\mathrm{amb},2}(m)<\rho_{\mathrm{amb}}(m)\) 使得矩阵幂满足 Perron--Frobenius 尖锐分解
\[
(A_m^{\mathrm{amb}})^n=\rho_{\mathrm{amb}}(m)^n\,v_m w_m^{\mathsf T}\ +\ O\!\bigl(\rho_{\mathrm{amb},2}(m)^n\bigr),
\qquad n\to\infty,
\]
其中 \((v_m,w_m)\) 为归一化的 Perron 左/右特征向量对。
左右乘以 \(e_{V_m}^{\mathsf T}\) 与 \(\mathbf 1\) 即得
\[
N_m^{\mathrm{unsync}}(n)
=(e_{V_m}^{\mathsf T}v_m)(w_m^{\mathsf T}\mathbf 1)\,\rho_{\mathrm{amb}}(m)^n\ +\ O\!\bigl(\rho_{\mathrm{amb},2}(m)^n\bigr),
\]
从而主项常数闭式与第一条渐近成立。
将 \(\rho_{\mathrm{amb}}(m)^n/2^n=e^{-(\log 2-\log\rho_{\mathrm{amb}}(m))n}=e^{-\varepsilon_m n}\) 与同样的误差缩放代入，即得“锐化曲线”形式；归一化误差由 \((\rho_{\mathrm{amb},2}(m)/2)^n\) 控制，而当 \(\rho_{\mathrm{amb},2}(m)>0\) 时其指数差为 \(\Delta_{\mathrm{amb}}(m)\)。
证毕。
\end{proof}

\leanverified{paper\_Ym\_scan\_length\_asymptotic}
\begin{corollary}[给定清晰度阈值的渐近最优扫描长度]\label{cor:Ym-scan-length-asymptotic}
在定理 \ref{thm:Ym-unsync-prefix-pf-sharp} 的假设下，给定 \(\eta\in(0,1)\)，定义最小扫描长度
\[
n_{\eta}:=\min\Bigl\{\,n\in\NN:\ \frac{N_m^{\mathrm{unsync}}(n)}{2^n}\le \eta\,\Bigr\}.
\]
则当 \(\eta\downarrow 0\) 时有渐近律
\[
n_{\eta}
=\frac{1}{\varepsilon_m}\log\!\Bigl(\frac{c_m}{\eta}\Bigr)\ +\ O(1).
\]
特别地，\(\varepsilon_m^{-1}\) 给出分辨率 \(m\) 下内禀的“扫描--清晰度”尺度。
\end{corollary}

\begin{proof}
由定理 \ref{thm:Ym-unsync-prefix-pf-sharp}，
\(
\frac{N_m^{\mathrm{unsync}}(n)}{2^n}
=c_m e^{-\varepsilon_m n}\Bigl(1+O\!\bigl((\rho_{\mathrm{amb},2}(m)/\rho_{\mathrm{amb}}(m))^{n}\bigr)\Bigr)
\)。
令 \(\eta\downarrow 0\) 并解不等式 \(c_m e^{-\varepsilon_m n}\lesssim \eta\)，即可得到
\(
n_{\eta}=\varepsilon_m^{-1}\log(c_m/\eta)+O(1)
\)。
证毕。
\end{proof}

\leanverified{paper\_ym\_finite\_delay\_spectral\_tail}
\begin{corollary}[有限延迟与谱尾的严格二分]\label{cor:Ym-finite-delay-spectral-tail}
在定理 \ref{thm:Ym-sync-tail} 的记号下，非单点态诱导子图 \(H_m^{\mathrm{amb}}\) 无有向环当且仅当
\[
\rho_{\mathrm{amb}}(m)=0.
\]
在此情形下，存在 \(n_0<\infty\) 使得对所有 \(n\ge n_0\) 有 \(N_m^{\mathrm{unsync}}(n)=0\)（有限延迟同步）。
反之，若 \(H_m^{\mathrm{amb}}\) 含有向环（从而 \(Y_m^{\mathrm{amb}}\neq\varnothing\)），则 \(1\le \rho_{\mathrm{amb}}(m)<2\) 并给出非平凡的指数型谱尾。
\end{corollary}

\begin{proof}
若 \(H_m^{\mathrm{amb}}\) 无环，则在有限图上其邻接矩阵 \(A_m^{\mathrm{amb}}\) 为幂零矩阵（某次幂为零），从而谱半径 \(\rho(A_m^{\mathrm{amb}})=0\)，即 \(\rho_{\mathrm{amb}}(m)=0\)；并且 \((A_m^{\mathrm{amb}})^n=0\) 对充分大 \(n\) 成立，于是由
\(
N_m^{\mathrm{unsync}}(n)=e_{V_m}^{\mathsf T}(A_m^{\mathrm{amb}})^n\mathbf 1
\)
得 \(N_m^{\mathrm{unsync}}(n)=0\) 终止于有限处。
\par
反向蕴含同理：若 \(\rho_{\mathrm{amb}}(m)=0\)，则 \(A_m^{\mathrm{amb}}\) 为幂零，故 \(H_m^{\mathrm{amb}}\) 不含有向环。
若 \(H_m^{\mathrm{amb}}\) 含有向环，则其最大强连通分量的谱半径至少为 \(1\)，故 \(\rho_{\mathrm{amb}}(m)\ge 1\)；另一方面由 \(h_{\mathrm{top}}(Y_m)=\log 2\) 及 \(\rho_{\mathrm{amb}}(m)=e^{h_{\mathrm{top}}(Y_m^{\mathrm{amb}})}\le e^{h_{\mathrm{top}}(Y_m)}=2\) 得 \(\rho_{\mathrm{amb}}(m)\le 2\)，并在“严格低熵壳”情形下有 \(\rho_{\mathrm{amb}}(m)<2\)。
证毕。
\end{proof}

\paragraph{分层清晰度指数：残余歧义等级的谱级分解}
上文以 \(|S_t|>1\) 作为统计型 \(\mathsf{NULL}\) 的判据，从而把“是否同步”视为最粗粒度的清晰度刻画。
由于候选集合在读入下单调不增（命题 \ref{prop:det-candidate-set-monotone}），可以进一步按残余候选规模分层，从而得到一族严格可计算的清晰度指数（而非单一的同步指数）。

\begin{definition}[残余歧义等级与诱导计数图]\label{def:Ym-ambiguity-level-k}
对整数 \(k\ge 2\)，定义阈值态集合
\[
\mathcal{S}_{m,\ge k}:=\{\,S\in\mathcal{S}_m^{\mathrm{det}}:\ |S|\ge k\,\},
\]
并令 \(H_{m,\ge k}\) 为确定化表示 \(H_m^{\mathrm{det}}\) 在 \(\mathcal{S}_{m,\ge k}\) 上的诱导子图。
记其按边条数计数的邻接矩阵为
\[
(A_{m,\ge k})_{S,S'}:=\#\{\,a\in X_m:\ \delta_m(S,a)=S'\,\},\qquad S,S'\in\mathcal{S}_{m,\ge k}.
\]
对任意 \(n\in\NN\)，定义长度-\(n\) 的 \(k\)-歧义前缀集合
\[
\mathcal{W}_{m,\ge k}(n)
:=\{\,w\in X_m^n:\ \text{从 }S_0=V_m\text{ 出发读入 }w\text{ 后满足 }|S_n|\ge k\,\},
\]
并记其计数为 \(N_{m,\ge k}(n):=\#\mathcal{W}_{m,\ge k}(n)\)。
显然 \(k=2\) 时有 \(N_{m,\ge 2}(n)=N_m^{\mathrm{unsync}}(n)\)。
\end{definition}

\begin{theorem}[歧义谱定理：残余歧义等级的 Perron--Frobenius 锐化律]\label{thm:Ym-ambiguity-spectrum-law}
固定整数 \(k\ge 2\) 并假设诱导子图 \(H_{m,\ge k}\) 存在唯一的最大谱半径强连通分量且该分量为原始的。
令
\[
\rho_{m,\ge k}:=\rho(A_{m,\ge k}),\qquad
\rho_{m,\ge k,2}:=\max\bigl(\{\,|\lambda|:\ \lambda\in\sigma(A_{m,\ge k}),\ |\lambda|<\rho_{m,\ge k}\,\}\cup\{0\}\bigr),
\]
则存在显式常数 \(c_{m,\ge k}>0\) 使得
\[
N_{m,\ge k}(n)=c_{m,\ge k}\,\rho_{m,\ge k}^{\,n}\ +\ O\!\bigl(\rho_{m,\ge k,2}^{\,n}\bigr),
\qquad n\to\infty.
\]
令 \(\varepsilon_{m,\ge k}:=\log 2-\log\rho_{m,\ge k}\)。
则有分层锐化曲线
\[
\frac{N_{m,\ge k}(n)}{2^n}
=c_{m,\ge k}\,e^{-\varepsilon_{m,\ge k} n}\ +\ O\!\Bigl(\bigl(\tfrac{\rho_{m,\ge k,2}}{2}\bigr)^{n}\Bigr),
\qquad n\to\infty.
\]
当 \(\rho_{m,\ge k,2}>0\) 时，定义二阶谱隙指数
\[
\Delta_{m,\ge k}:=\log\!\Bigl(\frac{\rho_{m,\ge k}}{\rho_{m,\ge k,2}}\Bigr)>0,
\]
则上述误差项亦可写成 \(O\!\bigl(e^{-(\varepsilon_{m,\ge k}+\Delta_{m,\ge k})n}\bigr)\)。
并且指数满足单调链
\[
\boxed{\ \varepsilon_{m,\ge 2}\ \le\ \varepsilon_{m,\ge 3}\ \le\ \varepsilon_{m,\ge 4}\ \le\ \cdots\ }.
\]
此外，若 \((v_{m,\ge k},w_{m,\ge k})\) 为 \(A_{m,\ge k}\) 的 Perron--Frobenius 右/左特征向量并归一化为 \(w_{m,\ge k}^{\mathsf T}v_{m,\ge k}=1\)，则主项常数可取闭式
\[
c_{m,\ge k}=(e_{V_m}^{\mathsf T}v_{m,\ge k})\,(w_{m,\ge k}^{\mathsf T}\mathbf 1),
\]
其中 \(e_{V_m}\) 为初始态坐标向量、\(\mathbf 1\) 为 \(\mathcal{S}_{m,\ge k}\) 上的全 \(1\) 向量。
\end{theorem}

\begin{proof}
由确定性右解析性与定义 \ref{def:Ym-ambiguity-level-k} 可知：长度-\(n\) 标记词 \(w\) 使得 \(|S_n|\ge k\)，当且仅当从 \(V_m\) 出发沿诱导子图 \(H_{m,\ge k}\) 读入 \(w\) 得到一条长度-\(n\) 的标记路径。
按邻接矩阵 \(A_{m,\ge k}\) 的多重度定义，块计数满足
\[
N_{m,\ge k}(n)=e_{V_m}^{\mathsf T}(A_{m,\ge k})^n\mathbf 1.
\]
在“唯一最大谱半径强连通分量且原始”的假设下，对 \(A_{m,\ge k}\) 应用 Perron--Frobenius 尖锐分解即可得到
\(
N_{m,\ge k}(n)=c_{m,\ge k}\rho_{m,\ge k}^{\,n}+O(\rho_{m,\ge k,2}^{\,n})
\)
及相应的归一化锐化曲线；主项常数闭式由左右乘 \(e_{V_m}^{\mathsf T}\) 与 \(\mathbf 1\) 立即得到。
\par
对单调链：由集合嵌套 \(\mathcal{S}_{m,\ge k+1}\subseteq\mathcal{S}_{m,\ge k}\)，图 \(H_{m,\ge k+1}\) 由 \(H_{m,\ge k}\) 诱导得到；
因此（将缺失顶点按零行零列补齐到同一维数后）有逐项不等式 \(0\le A_{m,\ge k+1}\le A_{m,\ge k}\)。
由非负矩阵谱半径的单调性得 \(\rho_{m,\ge k+1}\le \rho_{m,\ge k}\)，从而 \(\varepsilon_{m,\ge k}=\log 2-\log\rho_{m,\ge k}\) 随 \(k\) 单调不减，得到所示链式不等式。
证毕。
\end{proof}
\leanverified{paper\_Ym\_ambiguity\_spectrum\_law}

\begin{corollary}[时间层右可逆性的二分律：有限延迟或指数稀薄壳]\label{cor:Ym-time-layer-dichotomy}
在推论 \ref{cor:Ym-singleton-essential-finite-delay} 的“单点本质态”假设下，窗口层的多对一（即 \(\Fold_m\) 在单窗上的非单射）在时间层的右可逆性上只可能表现为以下二者之一：
\begin{enumerate}
  \item \textbf{强可逆（有限延迟同步）：}若 \(Y_m^{\mathrm{amb}}=\varnothing\)，则存在统一有限同步步数 \(L_m<\infty\)（备注 \ref{rem:Ym-amb-empty}），并在延迟 \(L_m\) 之后右解码在线成立。
  \item \textbf{弱不可逆（指数稀薄壳）：}若 \(Y_m^{\mathrm{amb}}\neq\varnothing\)，则歧义仅能长期驻留在严格低熵子系统上；对最大熵测度 \(\nu_m\)，同步时间 \(T_m\) 具有指数尾界
  \(
  \nu_m(T_m>n)\le C_m e^{-\varepsilon_m n}
  \)
（定理 \ref{thm:Ym-sync-tail}），从而“长期不可逆”在最大熵部分上以指数小概率出现。
\end{enumerate}
因此，”窗口多对一”与”时间多对一”在本模型中被严格区分：时间层的右不可逆要么根本不存在（有限延迟后必同步），要么仅存在于指数稀薄的歧义壳层中。
\leanverified{paper\_fold\_time\_layer\_dichotomy}
\end{corollary}

\begin{remark}[\texorpdfstring{$\eta_m$}{eta_m} 的可观测解释：最大熵层“可逆性”的谱级证书]\label{rem:eta_m_as_reversibility_certificate}
命题 \ref{prop:Ym-zeta-residue-ratio} 引入的 \(\eta_m\) 是 cover-\(\zeta\) 与内禀-\(\zeta\) 在共同主极点处的留数比。
它不改变指数尺度（熵/主极点模长），但以常数级方式量化“cover 上的周期结构在因子映射 \(\pi_m:X_{A_m}\to Y_m\) 下被合并”的强度。
\par
另一方面，在右解析确定化表示 \(H_m^{\mathrm{det}}\) 中，子集状态 \(S_t\subset V_m\) 记录“给定过去输出时可能的内部后缀集合”；
\(|S_t|=1\) 表示已进入可右解码（局部一一）的熵层，而 \(|S_t|>1\) 则是歧义壳的非单点记忆。
因此可把 \(\eta_m\) 视为“最大熵层可逆性”的\textbf{谱级指纹}：\(\eta_m\) 越接近 \(1\)，越接近“在熵主项上几乎不发生合并”的一一因子；
\(\eta_m<1\) 则提示仍存在不可忽略的常数级合并（尽管该合并不会造成指数级丢熵；见命题 \ref{prop:periodic-count-sandwich} 与推论 \ref{cor:Ym-periodic-compression-ratio}）。
\end{remark}

\begin{definition}[单点态占比的在线估计器（无需行列式）]\label{def:eta_m_hat_singleton_rate}
对任意输出序列 \(y=(y_t)_{t\in\ZZ}\in Y_m\)，从初始候选集 \(S_0:=V_m\) 出发递推
\[
S_{t+1}:=\delta_m(S_t,y_t)\qquad (t\ge 0),
\]
其中 \(\delta_m\) 由定义 \ref{def:Phi_m-right-resolving} 给出。
定义长度 \(N\) 的单点态占比
\[
\widehat\eta_m(N;y):=\frac{1}{N}\#\{\,0\le t<N:\ |S_t|=1\,\}.
\]
\end{definition}

\begin{proposition}[单点态占比的极限：最大熵测度下的可逆性常数]\label{prop:eta_m_hat_ergodic_limit}
\leanverified{paper\_eta\_m\_hat\_ergodic\_limit}
设 \(Y_m\) 不可约，从而其最大熵测度 \(\nu_m\) 存在且由命题 \ref{prop:Phi_m-parry-pushforward} 给出。则存在常数 \(\eta_m^{\mathrm{obs}}\in(0,1]\) 使得对 \(\nu_m\)-几乎处处的 \(y\in Y_m\)，定义 \ref{def:eta_m_hat_singleton_rate} 的在线统计量满足
\[
\widehat\eta_m(N;y)\ \xrightarrow[N\to\infty]{}\ \eta_m^{\mathrm{obs}}.
\]
在推论 \ref{cor:Ym-singleton-essential-finite-delay} 的“单点本质态”条件下，且把 \(\nu_m\) 限制在最大熵本质分量上时有 \(\eta_m^{\mathrm{obs}}=1\)；偏离 \(1\) 的程度可被视为“在最大熵层仍发生右解码歧义”的可观测强度。
\end{proposition}
\begin{proof}
将 \((y_t,S_t)\) 视为有限状态扩展移位系统，对 \(f(y,S):=\mathbf{1}\{|S|=1\}\) 应用 Birkhoff 遍历定理得到
\(
\widehat\eta_m(N;y)\to \int f\,d\bar\nu_m=:\eta_m^{\mathrm{obs}}
\)
于 \(\nu_m\)-几乎处处成立。
若满足推论 \ref{cor:Ym-singleton-essential-finite-delay} 的“单点本质态”条件，则在最大熵本质分量上几乎处处 \(|S_t|=1\)，故 \(\eta_m^{\mathrm{obs}}=1\)。证毕。
\end{proof}


% Auto-split to keep each TeX source < 600 lines.

\subsubsection{歧义壳的主奇点星形簇与分形余维：周期调制的解析等价刻画}\label{subsubsec:Ym-amb-singularity-star-codimension}
本小节在定理 \ref{thm:Ym-unsync-prefix-pf-sharp} 的“未同步前缀计数＝歧义子图路径计数”接口之上，记录两个无需额外结构假设的解析侧推论：
其一是“周期调制 \(\Longleftrightarrow\) 主奇点星形簇”的生成函数刻画；
其二是歧义壳在前缀超度量下的 Hausdorff 维数与余维闭式。

\begin{theorem}[周期调制 \texorpdfstring{$\Longleftrightarrow$}{<->} 主奇点星形簇：等模极点的不可消去性]\label{thm:Ym-amb-main-singularity-star}
假设 \(Y_m^{\mathrm{amb}}\neq\varnothing\)。
沿用定义 \ref{def:Ym-unsync-prefix-count} 与定理 \ref{thm:Ym-unsync-prefix-pf-sharp} 的记号，令
\[
N_m^{\mathrm{unsync}}(n):=\#\mathcal{W}_{m}^{\mathrm{unsync}}(n),
\qquad
U_m(z):=\sum_{n\ge 0}N_m^{\mathrm{unsync}}(n)\,z^n,
\]
并令 \(A_m^{\mathrm{amb}}\) 为按边条数计数的歧义邻接矩阵（定义于定理 \ref{thm:Ym-sync-tail-pf-sharp}，并在定理 \ref{thm:Ym-unsync-prefix-pf-sharp} 中用于未同步前缀的计数表示）。
则有严格恒等式
\[
U_m(z)=e_{V_m}^{\mathsf T}(I-zA_m^{\mathrm{amb}})^{-1}\mathbf 1,
\]
因此 \(U_m\) 为有理函数。
\par
进一步，设 \(A_m^{\mathrm{amb}}\) 的最大谱半径强连通分量唯一，且其周期为 \(p\ge 1\)。
令 \(\rho_{\mathrm{amb}}(m):=\rho(A_m^{\mathrm{amb}})\)。则 \(U_m(z)\) 的收敛半径为 \(R=\rho_{\mathrm{amb}}(m)^{-1}\)，并且 \(\det(I-zA_m^{\mathrm{amb}})=0\) 在圆周 \(|z|=R\) 上的解形成等模点集
\[
z_r=\rho_{\mathrm{amb}}(m)^{-1}e^{2\pi i r/p},\qquad r=0,1,\dots,p-1.
\]
当 \(p>1\) 时，上述点集给出“主奇点候选星形簇”：它们是 \(\det(I-zA_m^{\mathrm{amb}})=0\) 在 \(|z|=R\) 上的全部解。
由于 \(U_m(z)\) 是 resolvent 的一个矩阵元，可能发生约化消去；在不发生消去时，这 \(p\) 个点均为 \(U_m\) 的主极点，并与定理 \ref{thm:Ym-sync-tail-pf-sharp} 中的 \(n\bmod p\) 调制相一致。
\end{theorem}

\begin{proof}
由定理 \ref{thm:Ym-unsync-prefix-pf-sharp} 的证明中矩阵幂公式
\(
N_m^{\mathrm{unsync}}(n)=e_{V_m}^{\mathsf T}(A_m^{\mathrm{amb}})^n\mathbf 1
\)，对 \(|z|<\rho_{\mathrm{amb}}(m)^{-1}\) 有
\[
U_m(z)=\sum_{n\ge 0}e_{V_m}^{\mathsf T}(A_m^{\mathrm{amb}})^n\mathbf 1\,z^n
=e_{V_m}^{\mathsf T}\left(\sum_{n\ge 0}(zA_m^{\mathrm{amb}})^n\right)\mathbf 1
=e_{V_m}^{\mathsf T}(I-zA_m^{\mathrm{amb}})^{-1}\mathbf 1,
\]
从而 \(U_m\) 有理且其奇点包含 \(\det(I-zA_m^{\mathrm{amb}})=0\) 的零点。
\par
若最大谱半径强连通分量唯一且周期为 \(p\)，则 Perron--Frobenius 理论给出：谱圆上的特征值恰为
\(
\rho_{\mathrm{amb}}(m)e^{2\pi i r/p}
\)
（\(r=0,\dots,p-1\)），因而
\(
\det(I-zA_m^{\mathrm{amb}})=0
\)
在 \(|z|=R\) 上的解恰为
\(
z_r=\rho_{\mathrm{amb}}(m)^{-1}e^{2\pi i r/p}
\)。
对应地，resolvent \((I-zA_m^{\mathrm{amb}})^{-1}\) 在这些点有奇性；而 \(U_m(z)=e_{V_m}^{\mathsf T}(I-zA_m^{\mathrm{amb}})^{-1}\mathbf 1\) 的主奇点只能来自这些候选点。
若在相应的谱投影上 \(e_{V_m}^{\mathsf T}\) 与 \(\mathbf 1\) 的配对不退化，则不发生约化消去，候选点均成为 \(U_m\) 的实际主极点；此时复平面上的等模极点簇与时间域的 \(n\bmod p\) 调制互为解析等价刻画。证毕。
\end{proof}

\begin{corollary}[歧义壳的 Hausdorff 维数与余维闭式：\texorpdfstring{$\mathrm{codim}=\varepsilon_m/\log 2$}{codim=eps/log2}]\label{cor:Ym-amb-hausdorff-codim}
假设 \(Y_m^{\mathrm{amb}}\neq\varnothing\)，并以
\[
\varepsilon_m:=\log 2-\log\rho_{\mathrm{amb}}(m)>0
\]
记熵缺口（定理 \ref{thm:Ym-sync-tail-pf-sharp}）。
在 \(X_m^{\NN}\) 上取前缀超度量
\[
d(y,y'):=2^{-\ell(y,y')},
\qquad
\ell(y,y'):=\sup\{k\ge 0:\ y_0\cdots y_{k-1}=y'_0\cdots y'_{k-1}\},
\]
并将其限制到子系统 \(Y_m^{\mathrm{amb}}\subseteq X_m^{\NN}\)。
则
\[
\dim_{\mathrm H}(Y_m^{\mathrm{amb}})
=\frac{h_{\mathrm{top}}(Y_m^{\mathrm{amb}})}{\log 2}
=\frac{\log\rho_{\mathrm{amb}}(m)}{\log 2}
=1-\frac{\varepsilon_m}{\log 2}.
\]
从而歧义壳的分形余维严格为
\[
\mathrm{codim}(Y_m^{\mathrm{amb}})=\frac{\varepsilon_m}{\log 2}.
\]
\end{corollary}

\begin{proof}
在前缀超度量 \(d\) 下，长度 \(n\) 的柱集恰为半径 \(2^{-n}\) 的闭球；并且对任意子移位 \(Z\subseteq X_m^{\NN}\)，其长度-\(n\) 允许块数满足
\(
N_Z(n)=\exp(h_{\mathrm{top}}(Z)\,n+o(n))
\)。
由标准维数公式（在该超度量下，Hausdorff 维数等于熵除以 \(\log 2\)）得到
\(
\dim_{\mathrm H}(Z)=h_{\mathrm{top}}(Z)/\log 2
\)。
取 \(Z=Y_m^{\mathrm{amb}}\) 并用 \(h_{\mathrm{top}}(Y_m^{\mathrm{amb}})=\log\rho_{\mathrm{amb}}(m)\)（定理 \ref{thm:Ym-sync-tail-pf-sharp}）即可得到所示闭式；余维为 \(1-\dim_{\mathrm H}\) 的直接整理。证毕。
\end{proof}




\subsection{推广：一般 Ostrowski 计数下的稳定折叠}\label{subsec:ostrowski-generalized-fold}

本节将“Fibonacci/Zeckendorf 折叠”推广到一般无理数的 Ostrowski 计数系统，以回应“扩展到更广泛计数框架”的结构性需求。黄金分支的特殊性在于：当 $\alpha=\varphi^{-1}=[0;1,1,1,\dots]$ 时，Ostrowski 计数退化为 Fibonacci 权重与 Zeckendorf 语法，因此稳定类型空间成为最简单的二元 SFT（禁止相邻 $11$），并使 PF/$\zeta$ 接口闭式可得。

\subsubsection{一般 Ostrowski 权重与合法数字}

设 $\alpha\in(0,1)\setminus\QQ$ 的连分数展开为
\[
\alpha=[0;a_1,a_2,\dots],\qquad a_k\in\NN.
\]
令其收敛分母 $(q_n)_{n\ge 0}$ 由递推定义
\[
q_{-1}=0,\ q_0=1,\qquad q_{n+1}=a_{n+1}q_n+q_{n-1}\quad(n\ge 0).
\]

Ostrowski 表示定理保证：对任意 $N\in\NN$，存在唯一数字序列 $(d_n)_{n\ge 1}$ 满足
\(0\le d_n\le a_n\)、若 $d_n=a_n$ 则 $d_{n-1}=0$（$n\ge 2$）、且仅有限个 $d_n\neq 0$，并有
\[
N=\sum_{n\ge 1} d_n q_{n-1}.
\]
以下记该唯一序列为 $N$ 的 Ostrowski 数字，记为 $Z_\alpha(N)=(d_1,d_2,\dots)$（参见 \cite{Ostrowski1922} 及后续综述）。

\begin{definition}[一般稳定类型空间]\label{def:Xm-alpha}
对任意 $m\ge 1$，定义长度-$m$ 的 Ostrowski 合法集合
\[
X_m^{(\alpha)}:=\Bigl\{(d_1,\dots,d_m):\ 0\le d_n\le a_n,\ \ d_n=a_n\Rightarrow d_{n-1}=0\ (n\ge 2)\Bigr\}.
\]
并定义截断 $\pi_m^{(\alpha)}(d_1,d_2,\dots)=(d_1,\dots,d_m)$。
\end{definition}

\begin{proposition}[计数与递推（分母即计数）]\label{prop:Xm-alpha-cardinality}
对任意 $m\ge 1$，有
\[
\bigl|X_m^{(\alpha)}\bigr|=q_{m+1}.
\]
特别地，当 $\alpha=\varphi^{-1}$（即 $a_n\equiv 1$）时，$q_{m+1}=F_{m+2}$，而 $X_m^{(\alpha)}$ 退化为二元集合 $\{0,1\}^m$ 中禁止相邻 $11$ 的黄金均值语法 $X_m$。
\leanverified{paper\_xm\_alpha\_cardinality}
\leanverified{ostrowskiDenom}
\leanverified{ostrowskiDenom\_zero}
\leanverified{ostrowskiDenom\_one}
\leanverified{ostrowskiDenom\_succ\_succ}
\leanverified{ostrowskiDenom\_const\_one\_eq\_fib}
\leanverified{ostrowskiDenom\_pos}
\leanverified{ostrowskiDenom\_mono}
\end{proposition}

\begin{definition}[一般 Ostrowski 折叠的熵隙指数]\label{def:ostrowski-entropy-gap}
对给定 $\alpha=[0;a_1,a_2,\dots]$，长度 $m$ 的“原始数字空间”与稳定类型空间分别为
\[
\Omega_m^{(\alpha)}=\prod_{n=1}^{m}\{0,1,\dots,a_n\},
\qquad
X_m^{(\alpha)}\subseteq \Omega_m^{(\alpha)},
\]
并满足 $|X_m^{(\alpha)}|=q_{m+1}$（命题 \ref{prop:Xm-alpha-cardinality}）。
定义分辨率 $m$ 的折叠压缩量（自然对数）
\[
S_Z^{(\alpha)}(m):=\log\frac{|\Omega_m^{(\alpha)}|}{|X_m^{(\alpha)}|}
=\sum_{n=1}^{m}\log(a_n+1)-\log q_{m+1},
\]
以及对应的熵隙指数
\[
\kappa(\alpha):=\limsup_{m\to\infty}\frac{1}{m}S_Z^{(\alpha)}(m).
\]
\end{definition}

\begin{remark}[黄金分支为最简周期特例]
当 $\alpha=\varphi^{-1}=[0;1,1,1,\dots]$ 时，$\Omega_m^{(\alpha)}=\{0,1\}^m$、$|X_m^{(\alpha)}|=F_{m+2}$，从而
\[
S_Z^{(\alpha)}(m)=\log\frac{2^m}{F_{m+2}}
=m\log\!\left(\frac{2}{\varphi}\right)+O(1),
\]
因此 $\kappa(\varphi^{-1})=\log(2/\varphi)$，与本文二元折叠的熵隙常数一致（见第 \ref{subsec:pom-projection-entropy} 小节的相关公式）。
\end{remark}

\subsubsection{金属平均族：常数型 \texorpdfstring{$\alpha_A=[0;A,A,A,\dots]$}{alphaA} 的两态语法与“压缩/反共振”对偶}

\paragraph{两态 SFT 表示（合法语法的最小记忆）}
当 $a_n\equiv A$ 为常数（$A\in\NN_{\ge 1}$）时，合法约束
\(
d_n=A\Rightarrow d_{n-1}=0
\)
只依赖“上一位是否为 $0$”。因此可用两态记忆刻画：
\[
S_0:\ d_{n-1}=0,\qquad S_+:\ d_{n-1}\in\{1,\dots,A\}.
\]

\begin{proposition}[金属平均族的合法语法是两态 SFT]\label{prop:metallic-two-state-sft}
设 $\alpha_A=[0;A,A,A,\dots]$（$A\in\NN_{\ge 1}$）。则长度-$m$ 的合法集合 $X_m^{(\alpha_A)}$ 可由两态 SFT 的计数矩阵给出：
\[
M_A=
\begin{pmatrix}
1 & A\\
1 & A-1
\end{pmatrix},
\qquad
\chi_{M_A}(\lambda)=\lambda^2-A\lambda-1,
\]
其 Perron 根为
\[
\lambda_A=\frac{A+\sqrt{A^2+4}}{2}.
\]
并且拓扑熵满足
\[
h_{\mathrm{top}}\bigl(X^{(\alpha_A)}\bigr)=\log\lambda_A.
\]
\end{proposition}
\begin{proof}
转移计数为：从 $S_0$ 选 $0$（$1$ 种）留在 $S_0$，选 $1,\dots,A$（$A$ 种）到 $S_+$；从 $S_+$ 选 $0$（$1$ 种）到 $S_0$，选 $1,\dots,A-1$（$A-1$ 种）留在 $S_+$（$A$ 被禁止）。故计数矩阵为 $M_A$。
其特征多项式与 Perron 根为直接计算；而有限型移位的拓扑熵等于邻接矩阵谱半径的对数，故得结论。证毕。
\end{proof}
\leanverified{paper\_metallic\_two\_state\_sft}

\paragraph{折叠压缩的最优性：存在内部极大点}
对金属平均族，原始数字空间规模为 $|\Omega_m^{(\alpha_A)}|=(A+1)^m$，而命题 \ref{prop:Xm-alpha-cardinality} 给出
\(
|X_m^{(\alpha_A)}|=q_{m+1}\asymp \lambda_A^m
\)。
因此折叠压缩的指数斜率为
\[
h_Z(A):=\lim_{m\to\infty}\frac{1}{m}\log\frac{|\Omega_m^{(\alpha_A)}|}{|X_m^{(\alpha_A)}|}
=\log\!\left(\frac{A+1}{\lambda_A}\right).
\]

\begin{proposition}[金属平均族的投影压缩率在 \texorpdfstring{$A=\frac{3}{2}$}{A=3/2} 处取极大]\label{prop:metallic-compression-extremum}
把 $A$ 连续化为实参数（$A\ge 1$），则函数
\(
h_Z(A)=\log\!\left(\frac{A+1}{\lambda_A}\right)
\)
在区间 $[1,\infty)$ 上存在唯一极大点
\[
A_\star=\frac{3}{2}.
\]
若限制 $A\in\ZZ_{\ge 1}$，则 $h_Z(A)$ 的最大值在 $A=2$（白银分支）处取得。
\leanverified{paper\_metallic\_compression\_extremum\_seeds}
\leanverified{paper\_metallic\_compression\_extremum\_package}
\leanverified{paper\_metallic\_compression\_extremum}
\end{proposition}
\begin{proof}
由 $\lambda_A=(A+\sqrt{A^2+4})/2$，直接求导得到
\[
h_Z'(A)=\frac{1}{A+1}-\frac{1}{\sqrt{A^2+4}}.
\]
方程 $h_Z'(A)=0$ 等价于 $\sqrt{A^2+4}=A+1$，唯一解为 $A_\star=3/2$。再由
\[
h_Z''(A)=-\frac{1}{(A+1)^2}+\frac{A}{(A^2+4)^{3/2}}<0\qquad(A\ge 1)
\]
（可平方化为一条初等不等式）知该临界点为严格极大点，且全局唯一。
对整数点，因 $h_Z$ 在 $[1,A_\star]$ 上严格增、在 $[A_\star,\infty)$ 上严格减，最大只能在 $A=1,2$ 中取到。比较
\[
h_Z(2)-h_Z(1)=\log\!\left(\frac{3}{1+\sqrt2}\cdot\frac{1+\sqrt5}{4}\right)
\]
只需验证
\(
3(1+\sqrt5)>4(1+\sqrt2)
\)，即
\(
3\sqrt5-4\sqrt2>1
\)，其平方为 $77-24\sqrt{10}>1$，等价于 $(19/6)^2>10$，成立。故整数最大在 $A=2$。证毕。
\end{proof}

\paragraph{与反共振证书的不可兼得}
附录 \ref{app:kronecker-discrepancy-metallic-gap} 的 gap/差异度证书给出另一条“反共振强度”泛函，其典型形式可写为
\[
\kappa(A):=\frac{A}{\log\lambda_A},
\]
并且在 $A\ge 1$ 上严格递增（黄金分支 $A=1$ 在该证书意义下严格最优）。

\begin{remark}[两类最优性分离：反共振 vs 投影压缩]\label{rem:metallic-dual-optimality}
金属平均族的尺度更换并非“黄金最优”的单一结论：同一参数轴上存在两条竞争的极值目标。
\begin{itemize}
  \item \textbf{反共振/有限样本稳定性证书最强：}由 $\kappa(A)$ 的单调性，取 $A=1$（黄金）最优；
  \item \textbf{投影压缩率最大（每位被压扁的信息量最大）：}由命题 \ref{prop:metallic-compression-extremum}，连续极值在 $A=3/2$，整数最优在 $A=2$（白银）。
\end{itemize}
因此“反共振最优性”与“投影压缩最优性”在金属平均族上不可兼得：它们对应两条不同的可审计泛函，且极值点分离。
\end{remark}

\paragraph{金属平均族的二目标相图：帕累托前沿、标量化阈值与有限样本设计律}

\begin{proposition}[金属平均族的两目标帕累托前沿：\texorpdfstring{$\lambda$}{lambda} 一参数闭式刻画]\label{prop:metallic-pareto-frontier-lambda}
令 \(\lambda:=\lambda_A\) 为命题 \ref{prop:metallic-two-state-sft} 的 Perron 根。由特征方程
\(\lambda^2-A\lambda-1=0\) 得到参数变换
\[
A=\lambda-\lambda^{-1},\qquad \lambda\ge \varphi.
\]
于是两条可审计目标函数 \(h_Z\) 与 \(\kappa\) 可写成一维闭式曲线
\[
h_Z(\lambda)=\log\!\left(1+\lambda^{-1}-\lambda^{-2}\right),
\qquad
\kappa(\lambda)=\frac{\lambda-\lambda^{-1}}{\log\lambda}.
\]
在连续参数 \(A\ge 1\)（等价于 \(\lambda\ge\varphi\)）下，以“\textbf{最大化} \(h_Z\)”与“\textbf{最小化} \(\kappa\)”为准则的连续帕累托前沿
恰为区间 \(\lambda\in[\varphi,2]\)（等价于 \(A\in[1,3/2]\)）的像 \(\lambda\mapsto(\kappa(\lambda),h_Z(\lambda))\)；
而对 \(\lambda>2\)（等价于 \(A>3/2\)）的分支，两目标同时变差，从而被前沿严格支配。
\leanverified{paper\_metallic\_pareto\_frontier\_lambda}
\end{proposition}
\begin{proof}
由 \(\lambda^2-A\lambda-1=0\) 直接解出 \(A=\lambda-\lambda^{-1}\)。把该关系代回
\(
h_Z(A)=\log\bigl((A+1)/\lambda_A\bigr)
\)
与
\(
\kappa(A)=A/\log\lambda_A
\)
即可得到所示闭式表达。
\par
另一方面，命题 \ref{prop:metallic-compression-extremum} 表明 \(h_Z(A)\) 在 \([1,3/2]\) 上严格增、在 \([3/2,\infty)\) 上严格减，而 \(\kappa(A)\) 在 \(A\ge 1\) 上严格递增。
因此当 \(A>3/2\) 时，取 \(A_\star=3/2\) 有 \(\kappa(A_\star)<\kappa(A)\) 且 \(h_Z(A_\star)>h_Z(A)\)，故该分支被严格支配；而 \(A\in[1,3/2]\) 上任意两点一增一减，互不支配，从而构成帕累托前沿。证毕。
\end{proof}

\begin{remark}[设计空间缩减：无需讨论 \texorpdfstring{$A>3/2$}{A>3/2}]\label{rem:metallic-design-space-reduction}
命题 \ref{prop:metallic-pareto-frontier-lambda} 表明：在连续参数下，所有 \(A>3/2\) 的分支在“压缩/反共振”二目标意义下都被 \(A_\star=3/2\) 严格支配，因而不属于帕累托集合。
因此若允许 \(A\) 连续化，则设计空间可无损缩减为
\[
A\in\left[1,\frac{3}{2}\right].
\]
若强制 \(A\in\ZZ_{\ge 1}\)，则由同样的单调性可知 \(A\ge 3\) 均被 \(A=2\) 严格支配，从而整数帕累托集合缩减为 \(\{1,2\}\)。
\end{remark}

\begin{conjecture}[Pisot 规范化的末位支配：bin-\texorpdfstring{$\beta$}{beta}-fold 的一比特指数等价]\label{conj:ostrowski-pisot-last-digit-dominance-lecam}
设 \(\beta>1\) 为 Pisot 数，并考虑与 \(\beta\)-展开/相应 Ostrowski 合法语法相容的有限窗规范化折叠
\(
\mathrm{Fold}^{(\beta)}_m:\Omega_m^{(\beta)}\to X_m^{(\beta)}
\)。
令 \(U_m^{(\beta)}:=\mathrm{Unif}(X_m^{(\beta)})\)，并令 \(\mu_m^{(\beta)}\) 为 bin-\(\beta\)-fold 在 \(X_m^{(\beta)}\) 上的推前分布。
取末位统计量 \(\ell(x)\) 为最末一位数字（或等价的末位合法状态标签）。
则猜想存在常数 \(c_\beta\in(0,1)\) 与参数序列 \(\beta_m^{(\beta)}\to\log\beta\)，使得指数族近似
\[
\bar\mu_m^{(\beta)}(x)\ \propto\ U_m^{(\beta)}(x)\,\exp\!\bigl(-\beta_m^{(\beta)}\,\ell(x)\bigr)
\]
满足
\[
D_{\mathrm{TV}}\bigl(\mu_m^{(\beta)},\bar\mu_m^{(\beta)}\bigr)\ \le\ C_\beta\,c_\beta^{\,m},
\]
并且在 Le Cam 意义下，二点实验 \((\mu_m^{(\beta)},U_m^{(\beta)})\) 与其末位推前实验 \((\ell_\#\mu_m^{(\beta)},\ell_\#U_m^{(\beta)})\) 指数等价：全信息的最优检验误差在指数精度上由末位完全承载。
\end{conjecture}

\begin{corollary}[最大投影压缩的“语法熵锁定”：\texorpdfstring{$\lambda=2$}{lambda=2}]\label{cor:metallic-compression-locking-lambda2}
命题 \ref{prop:metallic-compression-extremum} 的连续极大点 \(A_\star=3/2\) 等价于
\[
\lambda_{A_\star}=2.
\]
此时最大压缩斜率为闭式常数
\[
h_Z(A_\star)=h_Z(\lambda=2)=\log\!\left(\frac{5}{4}\right).
\]
\end{corollary}
\begin{proof}
由 \(\lambda_A=(A+\sqrt{A^2+4})/2\)，代入 \(A_\star=3/2\) 得 \(\lambda_{A_\star}=2\)。
再由 \(h_Z(A)=\log((A+1)/\lambda_A)\) 得
\(
h_Z(A_\star)=\log((5/2)/2)=\log(5/4)
\)。证毕。
\end{proof}
\leanverified{paper\_metallic\_compression\_locking\_lambda2}

\begin{proposition}[线性权衡的端点相变阈值]\label{prop:metallic-linear-scalarization-threshold}
对权重 \(\beta\ge 0\)，考虑线性标量化目标
\[
J_\beta(A):=h_Z(A)-\beta\,\kappa(A),\qquad A\ge 1.
\]
定义临界阈值
\[
\beta_c:=\frac{h_Z'(1)}{\kappa'(1)}
=\frac{\left(\frac{1}{2}-\frac{1}{\sqrt5}\right)(\log\varphi)^2}{\log\varphi-\frac{1}{\sqrt5}}
\approx 0.35953.
\]
则 \(J_\beta\) 的最大化点 \(A_\beta\) 满足：
\begin{itemize}
  \item 若 \(\beta\ge \beta_c\)，则 \(A_\beta=1\)（黄金端点被强制选中）；
  \item 若 \(0\le \beta<\beta_c\)，则存在唯一 \(A_\beta\in(1,3/2]\) 使得 \(J_\beta'(A_\beta)=0\)，并且当 \(\beta\downarrow 0\) 时 \(A_\beta\to 3/2\)。
\end{itemize}
\leanverified{paper\_metallic\_linear\_scalarization\_threshold}
\end{proposition}
\begin{proof}
由命题 \ref{prop:metallic-pareto-frontier-lambda}，对任意 \(A>3/2\) 都有 \(h_Z(A)<h_Z(3/2)\) 且 \(\kappa(A)>\kappa(3/2)\)，因此对任意 \(\beta\ge 0\)，
\[
J_\beta(A)=h_Z(A)-\beta\kappa(A)\ <\ h_Z(3/2)-\beta\kappa(3/2)=J_\beta(3/2),
\]
从而最大化点必落在紧区间 \([1,3/2]\)。
\par
在该区间上，由命题 \ref{prop:metallic-compression-extremum} 的推导可得
\[
h_Z'(A)=\frac{1}{A+1}-\frac{1}{\sqrt{A^2+4}},\qquad
h_Z''(A)=-\frac{1}{(A+1)^2}+\frac{A}{(A^2+4)^{3/2}}<0.
\]
另一方面，由 \(\lambda_A^2-A\lambda_A-1=0\) 微分得到
\[
(\log\lambda_A)'=\frac{1}{\sqrt{A^2+4}},
\]
从而
\[
\kappa'(A)=\frac{\log\lambda_A-\frac{A}{\sqrt{A^2+4}}}{(\log\lambda_A)^2}.
\]
再令
\[
f(A):=\frac{2A\sqrt{A^2+4}}{A^2+8}-\log\lambda_A.
\]
直接计算得到
\[
f'(A)=\frac{A^2(8-A^2)}{\sqrt{A^2+4}\,(A^2+8)^2}>0\qquad(1\le A\le 3/2),
\]
并且 \(f(1)=2\sqrt5/9-\log\varphi\approx 1.50\times 10^{-2}>0\)，故 \(f(A)>0\) 在 \([1,3/2]\) 上成立。
注意到
\[
\kappa''(A)=\frac{-A^2-8+\frac{2A\sqrt{A^2+4}}{\log\lambda_A}}{(A^2+4)^{3/2}(\log\lambda_A)^2},
\]
其分子与 \(f(A)\) 同号，故 \(\kappa''(A)>0\) 在 \([1,3/2]\) 上成立。
于是对任意 \(\beta\ge 0\)，
\[
J_\beta''(A)=h_Z''(A)-\beta\,\kappa''(A)<0\qquad(1\le A\le 3/2),
\]
即 \(J_\beta\) 在 \([1,3/2]\) 上严格凹，从而最大化点唯一。
\par
再计算 \(J_\beta'(1)=h_Z'(1)-\beta\kappa'(1)\)。当 \(\beta\ge\beta_c\) 时，\(J_\beta'(1)\le 0\)，结合严格凹性可知最大化点为端点 \(A_\beta=1\)。
当 \(0\le\beta<\beta_c\) 时，\(J_\beta'(1)>0\)，且 \(J_\beta'(3/2)=h_Z'(3/2)-\beta\kappa'(3/2)=-\beta\kappa'(3/2)<0\)，由介值定理存在唯一 \(A_\beta\in(1,3/2]\) 使 \(J_\beta'(A_\beta)=0\)，并由严格凹性知其即为最大化点。
最后，当 \(\beta\downarrow 0\) 时，\(J_\beta'\to h_Z'\)（在紧区间上一致），而 \(h_Z'\) 在 \(A_\star=3/2\) 处唯一为零，故 \(A_\beta\to 3/2\)。证毕。
\end{proof}
\leanverified{paper\_metallic\_linear\_scalarization\_threshold}

\begin{corollary}[连续 Pareto 尺度律：权重 \texorpdfstring{$\beta$}{beta} \texorpdfstring{$\mapsto$}{->} 最优 \texorpdfstring{$A_\beta$}{A_beta}]\label{cor:metallic-pareto-scale-law}
沿命题 \ref{prop:metallic-linear-scalarization-threshold} 的记号，令 \(A_\beta\) 表示标量化目标
\(
J_\beta(A)=h_Z(A)-\beta\kappa(A)
\)
在 \(A\ge 1\) 上的唯一最大化点。
则：
\begin{enumerate}
  \item 对每个 \(\beta\ge 0\)，都存在且仅存在一个最优尺度 \(A_\beta\in[1,3/2]\)；
  \item 当 \(0\le \beta<\beta_c\) 时，\(A_\beta\in(1,3/2]\) 为唯一内点解
  \[
  h_Z'(A_\beta)=\beta\,\kappa'(A_\beta),
  \]
  并且映射 \(\beta\mapsto A_\beta\) 在 \([0,\beta_c)\) 上连续且严格递减（因此 \(\beta\) 增大把最优尺度从 \(3/2\) 单调推向 \(1\)）；
  \item 当 \(\beta\ge \beta_c\) 时，端点 \(A_\beta=1\) 被强制选中。
\end{enumerate}
\end{corollary}
\begin{proof}
命题 \ref{prop:metallic-linear-scalarization-threshold} 已给出：最大化点唯一，且必落在 \([1,3/2]\)；并且当 \(\beta<\beta_c\) 时最大化点为内点并满足一阶条件 \(J_\beta'(A_\beta)=0\)，即 \(h_Z'(A_\beta)=\beta\kappa'(A_\beta)\)。
\par
对 \(\beta\in[0,\beta_c)\) 定义
\(
F(A,\beta):=h_Z'(A)-\beta\kappa'(A)
\)。
由命题 \ref{prop:metallic-linear-scalarization-threshold} 的推导，\([1,3/2]\) 上有 \(h_Z''(A)<0\) 且 \(\kappa''(A)>0\)，从而在解点处
\(
\partial_A F(A_\beta,\beta)=h_Z''(A_\beta)-\beta\kappa''(A_\beta)<0
\)。
因此由隐函数定理，\(\beta\mapsto A_\beta\) 在 \([0,\beta_c)\) 上连续且可微，并满足
\[
\frac{dA_\beta}{d\beta}
=\frac{\kappa'(A_\beta)}{h_Z''(A_\beta)-\beta\kappa''(A_\beta)}<0,
\]
即严格递减。最后当 \(\beta\ge\beta_c\) 时的端点结论已由同一命题给出。证毕。
\end{proof}

\begin{corollary}[把权重绑定到有限样本证书的样本量阈值]\label{cor:metallic-sample-threshold}
\leanverified{paper\_metallic\_sample\_threshold}
令
\[
\beta(N):=\frac{\log N}{N}\qquad(N\ge 2),
\]
并记 \(A_N:=A_{\beta(N)}\) 为命题 \ref{prop:metallic-linear-scalarization-threshold} 的最大化点。
令 \(N_c\) 为方程 \(\beta(N)=\beta_c\) 在区间 \(N\in[e,\infty)\) 上的唯一解，则
\[
N_c\approx 3.42058.
\]
于是对所有整数 \(N\ge 4\)，均有 \(A_N\in(1,3/2]\)，并且当 \(N\to\infty\) 时 \(A_N\to 3/2\)。
\end{corollary}
\begin{proof}
函数 \(N\mapsto \beta(N)=\log N/N\) 在 \(N\ge e\) 上严格递减，且 \(\beta(e)=1/e>\beta_c\)、\(\lim_{N\to\infty}\beta(N)=0\)，故在 \([e,\infty)\) 上存在唯一解 \(N_c\)。
数值近似由该一元方程直接求得（可复现计算）。
当 \(N\ge 4>N_c\) 时有 \(\beta(N)<\beta_c\)，由命题 \ref{prop:metallic-linear-scalarization-threshold} 得 \(A_N\in(1,3/2]\)。
又因 \(\beta(N)\to 0\)，结合同一命题的极限结论得到 \(A_N\to 3/2\)。证毕。
\end{proof}

\begin{remark}[“最后 \texorpdfstring{$5\%$}{5\%} 压缩”的边际代价常数]\label{rem:metallic-marginal-price}
由推论 \ref{cor:metallic-compression-locking-lambda2} 的连续压缩最优点 \(A_\star=3/2\) 与黄金端点 \(A=1\) 的比较，
压缩斜率增量为
\[
\Delta h:=h_Z(3/2)-h_Z(1)
=\log\!\left(\frac{5}{4}\right)-\log\!\left(\frac{2}{\varphi}\right)
=\log\!\left(\frac{5\varphi}{8}\right)\approx 1.12082\times 10^{-2},
\]
相对增益 \(\Delta h/h_Z(1)\approx 5.29\%\)。
而反共振证书系数增量为
\[
\Delta \kappa:=\kappa(3/2)-\kappa(1)
=\frac{3}{2\log 2}-\frac{1}{\log\varphi}\approx 8.59556\times 10^{-2}.
\]
因此“单位压缩增益对应的证书代价”可由可复现常数指纹
\[
\boxed{\ \Pi:=\frac{\Delta \kappa}{\Delta h}\approx 7.669\ }
\]
量化。
\par
若强制整数 \(A\in\ZZ_{\ge 1}\)，则黄金 \(\to\) 白银（\(A=2\)）的对应比值为
\[
\Pi_{2:1}:=\frac{\kappa(2)-\kappa(1)}{h_Z(2)-h_Z(1)}\approx 36.03,
\]
体现出整数分支上“低回报-高代价”段更为陡峭。
\end{remark}

\begin{remark}[设计律：黄金 \texorpdfstring{$A=1$}{A=1} $\to$ 白银 \texorpdfstring{$A=2$}{A=2} 属于“低回报-高代价”段]\label{rem:metallic-design-law-low-return}
在 \(A\in[1,2]\)（等价于 \(\lambda\in[\varphi,1+\sqrt2]\)）上，压缩斜率的增益处于低量级而反共振证书单调变差：
\[
0<h_Z(2)-h_Z(1)
=\log\!\left(\frac{3}{1+\sqrt2}\cdot\frac{1+\sqrt5}{4}\right)\approx 5.2\times 10^{-3},
\]
而
\[
\kappa(2)-\kappa(1)=\frac{2}{\log(1+\sqrt2)}-\frac{1}{\log\varphi}>0.
\]
因此若主目标是“有限样本稳定性/可审计证书”，黄金端点 \(A=1\) 是性价比最优的帕累托端点；
若主目标是“最大投影压缩”，则应在连续最优 \(\lambda=2\)（或整数最优 \(A=2\)）附近选取参数。
\end{remark}

\subsubsection{金属平均族的局部规范化算法与截断误差的可审计统计}\label{subsubsec:ostrowski-metallic-local-normalization-truncation-mismatch}
本小节在常数型 \(\alpha_A=[0;A,A,A,\dots]\)（\(A\in\NN_{\ge 1}\)）的 Ostrowski 系统中，把“取规范形再截断”的折叠机制落实为一组可计算的局部重写规则，并据此给出朴素截断与折叠感知截断之间的单位长度差异率及其大字母渐近。
该结果提供了 \(\Fold_m^{(\alpha_A)}\) 的实现层可核验生成式，并把跨分辨率的规范选择差异编译为可审计的统计账本。

\paragraph{两规则规范化：终止性、数值保持与规范性}
对 \(\alpha_A\) 有 \(a_n\equiv A\)，其收敛分母递推为
\[
q_{-1}=0,\ q_0=1,\qquad q_{n+1}=Aq_n+q_{n-1}\quad(n\ge 0).
\]
为与定义 \ref{def:fold-alpha} 的取值公式 \(N_\alpha(\omega)=\sum_{n=1}^m \omega_n q_{n-1}\) 保持一致，对仅有限支撑的序列 \(\mathbf d=(d_n)_{n\ge 1}\in\NN^{(\NN)}\) 定义加权值
\[
V(\mathbf d):=\sum_{n\ge 1} d_n q_{n-1}\in\NN.
\]

\begin{theorem}[金属平均折叠的两规则规范化系统]\label{thm:ostrowski-metallic-two-rule-normalization}
设 \(A\in\NN_{\ge 1}\)。在字母表 \(\{0,1,\dots,A+1\}\) 上考虑对有限支撑序列 \(\mathbf d\) 的如下局部改写（其中 \(i\) 为触发位置）：
\begin{enumerate}
  \item 若 \(i\ge 1\) 且 \(b\in\{1,\dots,A\}\)，并且
  \(
  (d_{i+2},d_{i+1},d_i)=(x,A,b)
  \)，
  则改写
  \[
  (x,A,b)\ \Longrightarrow\ (x+1,0,b-1).
  \]
  \item 若 \(i\ge 2\)，并且
  \(
  (d_{i+1},d_i,d_{i-1},d_{i-2})=(x,A+1,0,c)
  \)，
  则改写
  \[
  (x,A+1,0,c)\ \Longrightarrow\ (x+1,0,A-1,c+1).
  \]
\end{enumerate}
对任意输入 \(\omega\in\Omega_m^{(\alpha_A)}=\{0,1,\dots,A\}^m\)（高位补零视作有限支撑序列），反复在任意可用位置应用上述改写，过程必在有限步终止于某一不可约序列 \(\mathbf d^{\mathrm{nf}}\)，并且：
\begin{enumerate}
  \item \(\mathbf d^{\mathrm{nf}}\) 仅取值于 \(\{0,1,\dots,A\}\)，且满足合法约束 \(\ d_n^{\mathrm{nf}}=A\Rightarrow d_{n-1}^{\mathrm{nf}}=0\ (n\ge 2)\)；
  \item 数值保持：\(V(\mathbf d^{\mathrm{nf}})=V(\omega)\)；
  \item \(\mathbf d^{\mathrm{nf}}=Z_{\alpha_A}\!\bigl(V(\omega)\bigr)\)，从而
  \[
  \Fold_m^{(\alpha_A)}(\omega)=\pi_m^{(\alpha_A)}\!\bigl(\mathbf d^{\mathrm{nf}}\bigr).
  \]
\end{enumerate}
\end{theorem}
\begin{proof}
\textbf{数值不变性.}
对第一条改写，利用 \(q_{i+1}=Aq_i+q_{i-1}\) 得
\[
xq_{i+1}+Aq_i+bq_{i-1}
=xq_{i+1}+(q_{i+1}-q_{i-1})+bq_{i-1}
=(x+1)q_{i+1}+(b-1)q_{i-1},
\]
与右侧权和一致。
对第二条改写，记四元组处的权重分别为 \(q_i,q_{i-1},q_{i-2},q_{i-3}\)。
该步等价于验证恒等式
\[
(A+1)q_{i-1}=q_i+(A-1)q_{i-2}+q_{i-3}.
\]
由递推 \(q_i=Aq_{i-1}+q_{i-2}\) 与 \(q_{i-1}=Aq_{i-2}+q_{i-3}\)，右端化为
\[
q_i+(A-1)q_{i-2}+q_{i-3}
=Aq_{i-1}+q_{i-2}+(A-1)q_{i-2}+q_{i-3}
=Aq_{i-1}+Aq_{i-2}+q_{i-3}
=(A+1)q_{i-1},
\]
故恒等式成立。
再加上同一高位 \(xq_i\) 与更低位 \(cq_{i-3}\) 的共同项即得该四元组改写的权和相等。
因此每一步改写均保持 \(V(\mathbf d)\) 不变。
\par
\textbf{终止性.}
固定总值 \(N:=V(\omega)\)。由于 \(q_n\to\infty\)，表示 \(N\) 的非负数字序列只能在有限多位上非零。
在该有限集合上引入从高位到低位的字典序：两表示不同则比较其最大差异指标处的数字。
在第一条改写中，最高受影响位为 \(i+2\)，其数字严格增加 \(x\mapsto x+1\)；
在第二条改写中，最高受影响位为 \(i+1\)，其数字同样严格增加。
因此每步改写使字典序严格上升；而在固定 \(N\) 的表示集合中不存在无限严格上升链，故过程必有限步终止。
\par
\textbf{不可约性推出合法性.}
若不可约序列中出现数字 \(A+1\)，则由两规则的结构可知其低一位必须为 \(0\)，从而局部形态 \((A+1,0,\cdot)\) 必出现并触发第二条改写，矛盾。
因此不可约序列各位均不超过 \(A\)。
若仍存在某处 \(d_{n}=A\) 且 \(d_{n-1}>0\)，则第一条改写可用，亦矛盾。
故不可约序列满足 \(d_n=A\Rightarrow d_{n-1}=0\)，即落入合法语法。
\par
\textbf{正确性.}
不可约序列 \(\mathbf d^{\mathrm{nf}}\) 与 \(\omega\) 表示同一整数 \(N\)，且 \(\mathbf d^{\mathrm{nf}}\) 合法。
由 Ostrowski 表示的唯一性（见 \cite{Ostrowski1922} 以及本节前述固定的规范性口径），合法表示必等于规范数字 \(Z_{\alpha_A}(N)\)。
因此 \(\mathbf d^{\mathrm{nf}}=Z_{\alpha_A}(V(\omega))\)，并得到折叠等式。证毕。
\end{proof}

\paragraph{朴素截断与折叠感知截断：单位长度差异率与大字母主项}
对 \(\omega\in\Omega_{m+k}^{(\alpha_A)}=\{0,\dots,A\}^{m+k}\) 定义朴素截断
\[
r_{m+k\to m}^{(\alpha_A)}(\omega):=\pi_{m+k\to m}^{(\alpha_A)}(\omega)=(\omega_1,\dots,\omega_m),
\]
以及折叠感知截断
\[
\widetilde r_{m+k\to m}^{(\alpha_A)}(\omega)
:=\pi_{m+k\to m}^{(\alpha_A)}\!\left(\Fold_{m+k}^{(\alpha_A)}(\omega)\right).
\]
令其 Hamming 规范差为
\[
G^{(A)}_{m,k}(\omega)
:=\#\Bigl\{1\le i\le m:\ \bigl(r_{m+k\to m}^{(\alpha_A)}(\omega)\bigr)_i\neq \bigl(\widetilde r_{m+k\to m}^{(\alpha_A)}(\omega)\bigr)_i\Bigr\}.
\]

\begin{theorem}[金属平均截断规范差的极限与步长不变性]\label{thm:ostrowski-metallic-truncation-mismatch-limit}
在 \(\Omega_{m+k}^{(\alpha_A)}\) 上取乘积均匀测度。
则对每个固定 \(A\ge 1\) 与固定 \(k\ge 1\)，极限
\[
g_A:=\lim_{m\to\infty}\frac{1}{m}\,\EE\!\left[G^{(A)}_{m,k}\right]
\]
存在，且与 \(k\) 无关。
此外，当 \(A\to\infty\) 时存在常数 \(C>0\) 使
\[
\left|g_A-\frac{3A}{(A+1)^2}\right|\ \le\ \frac{C}{(A+1)^2},
\]
从而
\(
g_A=\frac{3}{A}+O(A^{-2})
\)。
\end{theorem}
\begin{proof}
把 \(\omega\) 视为截断自一条双侧独立均匀序列，并令 \(X_i\) 表示位置 \(i\) 是否发生失配的指示变量，则
\(
G^{(A)}_{m,k}=\sum_{i=1}^m X_i
\)。
由定理 \ref{thm:ostrowski-metallic-two-rule-normalization}，高位新增信息对低位的影响只能通过连续可触发模式向下传播；而每跨越一步继续传播都需要在相邻更高位再次遇到数字 \(A\)。
在独立均匀模型下，某指定位取值为 \(A\) 的概率为 \((A+1)^{-1}\)，故传播长度具有几何型尾界。
这意味着 \(X_i\) 对远处高位的依赖以指数小概率发生，从而形成具有指数衰减依赖的平移不变过程。
于是存在常数使
\[
\EE\!\left[G^{(A)}_{m,k}\right]=\sum_{i=1}^m \EE[X_i]=m\,\EE[X_0]+O(1),
\]
其中 \(O(1)\) 为边界截断项。
因此 \(m^{-1}\EE[G^{(A)}_{m,k}]\) 收敛到 \(g_A:=\EE[X_0]\)，并且由于传播尾界对 \(k\) 的改变只产生有界的端部差异，极限与 \(k\) 无关。
\par
最后讨论 \(A\to\infty\) 的主项。
考虑局部缺陷事件
\[
D_i:=\{\omega_{i+1}=A,\ \omega_i>0\}.
\]
在独立均匀测度下
\(
\PP(D_i)=\frac{A}{(A+1)^2}.
\)
当 \(D_i\) 发生且其邻域不包含第二个高位 \(A\) 引发的连锁修正时，第一条规则在该三元组上触发一次并改变恰好三个位置 \((i,i+1,i+2)\) 的数字，因此在该“孤立缺陷”情形下对失配计数的局部贡献为 \(3\)。
使这一三位贡献偏离 \(3\) 的情形要求在长度 \(O(1)\) 的窗口内出现至少两个取值为 \(A\) 的位置，从而其概率为 \(O((A+1)^{-2})\)。
因此单位长度期望失配满足
\(
g_A=3\PP(D_0)+O((A+1)^{-2})
\)，即所述估计。证毕。
\end{proof}

\begin{proposition}[传播长度的指数尾界与方差量级]\label{prop:ostrowski-metallic-propagation-tail-variance}
在定理 \ref{thm:ostrowski-metallic-truncation-mismatch-limit} 的概率模型中，存在常数 \(\rho_A\in(0,1)\)、\(C_A<\infty\) 使得对一切 \(t\ge 0\),
\[
\PP(\mathsf L\ge t)\le C_A\,\rho_A^{t},
\qquad\text{并可取}\quad
\rho_A\le \frac{1}{A+1},
\]
其中 \(\mathsf L\) 表示高位新增信息在规范化过程中向低位影响的最大回传距离。
从而对每个固定 \(A\) 与 \(k\)，有
\[
\Var\!\left(G^{(A)}_{m,k}\right)=\Theta(m)\qquad(m\to\infty),
\]
并且当 \(A\to\infty\) 时存在绝对常数 \(C\) 使
\[
\Var\!\left(G^{(A)}_{m,k}\right)\ \le\ C\,\frac{m}{A}.
\]
\end{proposition}
\begin{proof}
传播尾界来自定理 \ref{thm:ostrowski-metallic-two-rule-normalization} 所揭示的触发链结构：每向低位继续传播一步，需要一个新的高位取值 \(A\) 以维持可触发缺陷，故可由参数 \((A+1)^{-1}\) 的几何尾随机变量支配。
\par
令 \(X_i\) 为第 \(i\) 位失配指示变量，则 \(G^{(A)}_{m,k}=\sum_{i=1}^m X_i\)。
指数尾界蕴含存在常数 \(c>0\) 使得协方差满足
\(
|\Cov(X_i,X_j)|\le C'e^{-c|i-j|}
\)。
因此
\[
\Var\!\left(G^{(A)}_{m,k}\right)
=\sum_{i=1}^m\Var(X_i)+2\sum_{1\le i<j\le m}\Cov(X_i,X_j)
=m\sigma_A^2+o(m),
\]
其中 \(\sigma_A^2:=\sum_{r\in\ZZ}\Cov(X_0,X_r)\in(0,\infty)\)。
另一方面，当 \(A\to\infty\) 时由定理 \ref{thm:ostrowski-metallic-truncation-mismatch-limit} 得 \(\EE[X_0]=g_A=\Theta(A^{-1})\)，且 \(0\le X_0\le 1\) 给出 \(\Var(X_0)\le \EE[X_0]\)。
结合指数衰减求和，得到 \(\sigma_A^2=O(g_A)\)，从而 \(\Var(G^{(A)}_{m,k})\le C m/A\)（\(A\) 足够大时）。证毕。
\end{proof}

\begin{conjecture}[大字母展开的二阶系数]\label{conj:ostrowski-metallic-ga-second-order}
在定理 \ref{thm:ostrowski-metallic-truncation-mismatch-limit} 的记号下，当 \(A\to\infty\) 时存在渐近展开
\[
g_A=\frac{3}{A}-\frac{12}{A^2}+O\!\left(\frac{1}{A^3}\right),
\]
且二阶系数由最短“非孤立缺陷簇”对三位主贡献的系统性修正给出。
\end{conjecture}

\begin{conjecture}[金属平均族的零块记忆律与 \texorpdfstring{$\beta_A^{-2n}$}{betaA^{-2n}} 衰减模量]\label{conj:ostrowski-metallic-lucas-ostrowski-memory-law}
对每个 \(A\in\NN_{\ge 1}\)，令 \(\beta_A>1\) 为方程 \(x^2-Ax-1=0\) 的 Pisot 根。
令 \((U^{(A)}_n)_{n\ge 0}\) 为 Lucas--Fibonacci 型序列
\[
U^{(A)}_0=0,\qquad U^{(A)}_1=1,\qquad U^{(A)}_{n+2}=A U^{(A)}_{n+1}+U^{(A)}_{n}\quad(n\ge 0).
\]
设 \((g^{(A)}_t)_{t\in\ZZ}\) 为 \(\alpha_A=[0;A,A,A,\dots]\) 的折叠规范化所诱导的规范差指示过程（可参照猜想 \ref{conj:fold-gauge-anomaly-metallic-continuant-law} 的结构性刻画口径）。
定义零块语境下的一步绝对误差极小化量 \(\mathcal E^{(A)}_n\) 与定理 \ref{thm:fold-gauge-anomaly-zero-block-minimax-abs-error} 同型。
则存在仅依赖于 \(A\) 的整数修正 \(c(A)\) 与阈值 \(n_0(A)\) 使得当 \(n\ge n_0(A)\) 时成立闭式
\[
\boxed{\ \mathcal E^{(A)}_n=\frac{1}{4\,U^{(A)}_{n+c(A)}\,U^{(A)}_{n+c(A)+1}}\ }.
\]
并且误差指数衰减率满足
\[
\mathcal E^{(A)}_n=\Theta\!\bigl(\beta_A^{-2n}\bigr).
\]
在 \(A=1\) 时，\(\beta_A=\varphi\)、\(U^{(1)}_n=F_n\)，并可取 \(c(1)=3\)，从而退化为定理 \ref{thm:fold-gauge-anomaly-zero-block-minimax-abs-error} 的 Fibonacci 闭式。
\end{conjecture}

\endinput

\subsubsection{金属平均族的闭式计数、有限尺度振荡与整数尺度相变}\label{subsubsec:ostrowski-metallic-binet-oscillation-integer-phase-transition}
在命题 \ref{prop:metallic-two-state-sft}、命题 \ref{prop:metallic-compression-extremum} 与命题 \ref{prop:metallic-linear-scalarization-threshold} 已经把常数型金属平均族的两态语法、连续压缩极值与连续标量化前沿固定下来之后，有限尺度计数与整数尺度选择还允许进一步闭式化。
下述结果把合法语言计数、有限分辨率压缩缺陷与离散样本阈值统一到同一组完全显式公式中。

\begin{theorem}[金属平均族稳定语言计数的闭式 Binet 公式]\label{thm:metallic-binet-closed-form}
沿命题 \ref{prop:metallic-two-state-sft} 的记号，固定整数 \(A\ge 1\)，并记
\[
\alpha_A=[0;A,A,A,\dots],
\qquad
\lambda_A=\frac{A+\sqrt{A^2+4}}{2},
\qquad
\mu_A:=\frac{A-\sqrt{A^2+4}}{2}=-\lambda_A^{-1}.
\]
令 \((q_n)_{n\ge 0}\) 为 \(\alpha_A\) 的收敛分母序列，即
\[
q_{-1}=0,\qquad q_0=1,\qquad q_{n+1}=Aq_n+q_{n-1}\quad(n\ge 0).
\]
则对一切 \(n\ge 0\) 都有严格闭式
\[
q_n=
\frac{\lambda_A^{n+2}+(-1)^n\lambda_A^{-n}}{\lambda_A^2+1}.
\]
特别地，由命题 \ref{prop:Xm-alpha-cardinality}，对一切 \(m\ge 1\),
\[
\bigl|X_m^{(\alpha_A)}\bigr|
=q_{m+1}
=
\frac{\lambda_A^{m+3}+(-1)^{m+1}\lambda_A^{-m-1}}{\lambda_A^2+1}.
\]
\leanverified{paper\_metallic\_binet\_closed\_form}
\leanverified{phi\_irrational}
\end{theorem}
\begin{proof}
递推 \(q_{n+1}=Aq_n+q_{n-1}\) 的特征方程为
\[
\lambda^2-A\lambda-1=0,
\]
其两根正是 \(\lambda_A\) 与 \(\mu_A\)。
因此存在常数 \(c_1,c_2\) 使
\[
q_n=c_1\lambda_A^n+c_2\mu_A^n.
\]
由初值 \(q_{-1}=0,\ q_0=1\) 解得
\[
q_n=\frac{\lambda_A^{n+1}-\mu_A^{n+1}}{\lambda_A-\mu_A}.
\]
再用 \(\mu_A=-\lambda_A^{-1}\) 以及
\[
\lambda_A-\mu_A
=
\lambda_A+\lambda_A^{-1}
=
\frac{\lambda_A^2+1}{\lambda_A}
\]
化简，即得
\[
q_n=
\frac{\lambda_A^{n+2}+(-1)^n\lambda_A^{-n}}{\lambda_A^2+1}.
\]
最后代入命题 \ref{prop:Xm-alpha-cardinality} 的恒等式
\(
|X_m^{(\alpha_A)}|=q_{m+1}
\)
即可。证毕。
\end{proof}

\leanverified{paper\_metallic\_finite\_scale\_compression\_oscillation}
\begin{theorem}[金属平均族投影压缩率的有限尺度奇偶振荡公式]\label{thm:metallic-finite-scale-compression-oscillation}
沿用前述记号。定义有限尺度投影压缩量
\[
\Xi_m(A):=\log\frac{(A+1)^m}{|X_m^{(\alpha_A)}|}
\qquad (m\ge 1),
\]
以及常数项
\[
\gamma_A:=\log\frac{\lambda_A^2+1}{\lambda_A^3}.
\]
则对一切 \(m\ge 1\) 都有精确恒等式
\[
\Xi_m(A)
=
m\,h_Z(A)+\gamma_A-\log\!\Bigl(1+(-1)^{m+1}\lambda_A^{-2m-4}\Bigr).
\]
从而对任意 \(n\ge 0\)，偶、奇子序列满足
\[
\Xi_{2n}(A)-\bigl(2n\,h_Z(A)+\gamma_A\bigr)
=
-\log\!\bigl(1-\lambda_A^{-4n-4}\bigr)>0,
\]
\[
\Xi_{2n+1}(A)-\bigl((2n+1)\,h_Z(A)+\gamma_A\bigr)
=
-\log\!\bigl(1+\lambda_A^{-4n-6}\bigr)<0.
\]
因此有限尺度误差并非仅具 \(O(\lambda_A^{-2m})\) 量级，而是满足严格的奇偶振荡律：偶数子序列从上方收敛，奇数子序列从下方收敛，且两者都以精确指数 \(2m\) 收敛到同一仿射主项。
\end{theorem}
\begin{proof}
由定理 \ref{thm:metallic-binet-closed-form}，
\[
\bigl|X_m^{(\alpha_A)}\bigr|
=
\frac{\lambda_A^{m+3}+(-1)^{m+1}\lambda_A^{-m-1}}{\lambda_A^2+1}
=
\frac{\lambda_A^{m+3}}{\lambda_A^2+1}
\Bigl(1+(-1)^{m+1}\lambda_A^{-2m-4}\Bigr).
\]
因此
\[
\Xi_m(A)
=
m\log(A+1)
-(m+3)\log\lambda_A
+\log(\lambda_A^2+1)
-\log\!\Bigl(1+(-1)^{m+1}\lambda_A^{-2m-4}\Bigr).
\]
把前两项合并为
\[
m\bigl(\log(A+1)-\log\lambda_A\bigr)=m\,h_Z(A),
\]
余下常数项正是
\[
\log(\lambda_A^2+1)-3\log\lambda_A
=
\log\frac{\lambda_A^2+1}{\lambda_A^3}
=
\gamma_A.
\]
于是得到主公式。
再分别代入 \(m=2n\) 与 \(m=2n+1\) 即得两条奇偶公式，其符号由 \(\lambda_A>1\) 直接给出。证毕。
\end{proof}

\begin{theorem}[整数尺度的线性标量化阈值]\label{thm:metallic-integer-scalarization-threshold}
沿用命题 \ref{prop:metallic-linear-scalarization-threshold} 的标量化目标
\[
J_\beta(A):=h_Z(A)-\beta\,\kappa(A),
\qquad
\kappa(A)=\frac{A}{\log\lambda_A}.
\]
在整数约束 \(A\in\ZZ_{\ge 1}\) 下，帕累托集合缩减为 \(\{1,2\}\)。
因此存在唯一阈值
\[
\beta_{\ZZ}
:=
\frac{h_Z(2)-h_Z(1)}{\kappa(2)-\kappa(1)}
=
\frac{\log\!\left(\dfrac{3\varphi}{2(1+\sqrt2)}\right)}
{\dfrac{2}{\log(1+\sqrt2)}-\dfrac{1}{\log\varphi}}
\approx 0.0277519138993265,
\]
其中 \(\varphi=(1+\sqrt5)/2\)。
并且最优整数尺度满足严格相变律
\[
A_\beta^{\ZZ}
=
\begin{cases}
2,& 0\le \beta<\beta_{\ZZ},\\[2mm]
\{1,2\},& \beta=\beta_{\ZZ},\\[2mm]
1,& \beta>\beta_{\ZZ}.
\end{cases}
\]
\end{theorem}
\begin{proof}
由注 \ref{rem:metallic-design-space-reduction}，所有整数 \(A\ge 3\) 在“最大化 \(h_Z\)、最小化 \(\kappa\)”的双目标意义下都被 \(A=2\) 严格支配。
故在线性标量化
\(
J_\beta(A)=h_Z(A)-\beta\kappa(A)
\)
下，只需比较 \(A=1\) 与 \(A=2\)。

由 \(\lambda_1=\varphi\) 与 \(\lambda_2=1+\sqrt2\)，有
\[
h_Z(1)=\log\frac{2}{\varphi},
\qquad
h_Z(2)=\log\frac{3}{1+\sqrt2},
\]
\[
\kappa(1)=\frac{1}{\log\varphi},
\qquad
\kappa(2)=\frac{2}{\log(1+\sqrt2)}.
\]
阈值条件 \(J_\beta(1)=J_\beta(2)\) 唯一给出
\[
\beta
=
\frac{h_Z(2)-h_Z(1)}{\kappa(2)-\kappa(1)}
=
\beta_{\ZZ}.
\]
由于 \(\kappa(2)>\kappa(1)\) 且 \(h_Z(2)>h_Z(1)\)，当 \(\beta<\beta_{\ZZ}\) 时压缩项占优，故 \(A=2\) 唯一胜出；当 \(\beta>\beta_{\ZZ}\) 时证书项占优，故 \(A=1\) 唯一胜出；等号时二者并列。证毕。
\end{proof}
\leanverified{paper_metallic_integer_scalarization_threshold}

\begin{corollary}[样本惩罚律下的整数相变与 Lambert--\texorpdfstring{$W$}{W} 闭式]\label{cor:metallic-sample-penalty-integer-threshold-lambertw}
\leanverified{paper\_metallic\_sample\_penalty\_integer\_threshold\_lambertw}
沿用定理 \ref{thm:metallic-integer-scalarization-threshold} 的阈值 \(\beta_{\ZZ}\)。
对整数样本量 \(N\ge 3\)，令
\[
\beta(N):=\frac{\log N}{N},
\]
并按同一标量化目标在整数尺度上取最优点 \(A_N^{\ZZ}\)。
则存在唯一实阈值
\[
N_{\ZZ}
=
-\frac{W_{-1}(-\beta_{\ZZ})}{\beta_{\ZZ}}
\approx 188.850182986328,
\]
其中 \(W_{-1}\) 为 Lambert--\(W\) 函数的 \((-1)\) 支。
于是对所有整数 \(N\ge 3\)，最优整数尺度满足
\[
A_N^{\ZZ}
=
\begin{cases}
1,& 3\le N\le 188,\\[2mm]
2,& N\ge 189.
\end{cases}
\]
\end{corollary}
\begin{proof}
由定理 \ref{thm:metallic-integer-scalarization-threshold}，整数最优尺度完全由 \(\beta(N)\) 与 \(\beta_{\ZZ}\) 的比较决定。
函数
\[
\beta(N)=\frac{\log N}{N}
\]
在 \(N\ge e\) 上严格递减，而 \(3>e\)，故其在全部整数 \(N\ge 3\) 上也严格递减。
阈值点满足方程
\[
\frac{\log N}{N}=\beta_{\ZZ}.
\]
令 \(u=\log N\)，则 \(N=e^u\)，上式等价于
\[
u\,e^{-u}=\beta_{\ZZ}
\qquad\Longleftrightarrow\qquad
(-u)e^{-u}=-\beta_{\ZZ}.
\]
因此
\[
u=-W(-\beta_{\ZZ}).
\]
为得到 \(N>e\) 的大根，必须选取 Lambert--\(W\) 的 \((-1)\) 支，从而
\[
N_{\ZZ}=e^u
=
e^{-W_{-1}(-\beta_{\ZZ})}
=
-\frac{W_{-1}(-\beta_{\ZZ})}{\beta_{\ZZ}}.
\]
数值上 \(N_{\ZZ}\approx 188.85018\)。
由于 \(\beta(N)\) 在 \(N\ge 3\) 上严格递减，故
\[
\beta(N)>\beta_{\ZZ}\iff N\le 188,
\qquad
\beta(N)<\beta_{\ZZ}\iff N\ge 189.
\]
再由定理 \ref{thm:metallic-integer-scalarization-threshold} 的相变律得到结论。证毕。
\end{proof}

\endinput


\subsubsection{一般折叠算子 \texorpdfstring{$\Fold_m^{(\alpha)}$}{Fold\_m(alpha)}}

为了与本文的“折叠 = 取值 + 取规范形”机制保持一致，我们把一般 Ostrowski 折叠定义为：先将“未规范数字”（允许违反相邻最大值约束）赋一个可审计整数值，再取 Ostrowski 规范形并截断。

\begin{definition}[Ostrowski 加权值与折叠]\label{def:fold-alpha}
令
\[
\Omega_m^{(\alpha)}:=\prod_{n=1}^{m}\{0,1,\dots,a_n\}
\]
为长度-$m$ 的“原始数字空间”。对 $\omega=(\omega_1,\dots,\omega_m)\in\Omega_m^{(\alpha)}$ 定义其 Ostrowski 加权值
\[
N_\alpha(\omega):=\sum_{n=1}^{m}\omega_n q_{n-1}\in\NN.
\]
定义一般折叠
\[
\Fold_m^{(\alpha)}:\Omega_m^{(\alpha)}\to X_m^{(\alpha)},\qquad
\Fold_m^{(\alpha)}(\omega):=\pi_m^{(\alpha)}\!\bigl(Z_\alpha(N_\alpha(\omega))\bigr).
\]
\end{definition}

\begin{proposition}[一般折叠的基本性质]\label{prop:fold-alpha-basic}
\leanverified{paper\_fold\_alpha\_basic}
$\Fold_m^{(\alpha)}$ 良定义且满足：
\begin{enumerate}
  \item \textbf{（规范性）} $\Fold_m^{(\alpha)}(\omega)\in X_m^{(\alpha)}$；
  \item \textbf{（幂等性）} 若 $\omega\in X_m^{(\alpha)}$，则 $\Fold_m^{(\alpha)}(\omega)=\omega$；
  \item \textbf{（满射性）} $\Fold_m^{(\alpha)}:\Omega_m^{(\alpha)}\twoheadrightarrow X_m^{(\alpha)}$ 为满射；
  \item \textbf{（多尺度相容）} 对 $m_2\ge m_1$，有
  \[
  \pi_{m_2\to m_1}^{(\alpha)}\circ \Fold_{m_2}^{(\alpha)}=\Fold_{m_1}^{(\alpha)}\circ \pi_{m_2\to m_1}^{(\alpha)},
  \]
  其中 $\pi_{m_2\to m_1}^{(\alpha)}$ 为截断到前 $m_1$ 位的投影。
\end{enumerate}
\end{proposition}

\begin{proof}
第 1 条由 Ostrowski 表示的合法性立得。第 2 条：若 $\omega$ 已合法，则 $\omega$ 即为某个整数的规范 Ostrowski 数字的前 $m$ 位，唯一性迫使规范化后截断不变。第 3 条：对任意 $x\in X_m^{(\alpha)}$，取 $\omega=x$ 即有 $\Fold_m^{(\alpha)}(\omega)=x$。第 4 条由截断与“先取值再取规范形”的函子性得到：对任意 $\omega$，两边都等于 $\pi_{m_1}^{(\alpha)}(Z_\alpha(N_\alpha(\omega)))$。
\end{proof}

\begin{theorem}[“先规范形再截断”与跨分辨率函子性]\label{thm:normalize-then-truncate-functorial}
一般 Ostrowski 折叠按定义 \ref{def:fold-alpha} 采取“取值 \(\to\) 取规范形 \(\to\) 截断”的顺序，因此跨分辨率交换图（命题 \ref{prop:fold-alpha-basic}-(4)）是\textbf{定义性}成立的：它并非额外结构，而是由规范化顺序强制给出。
\par
在本文的二元 Zeckendorf 折叠中，微态空间 \(\Omega_m=\{0,1\}^m\) 上的朴素截断 \(r_{m+1\to m}\)（小节 \ref{subsec:fold-omega-restriction}）对应的是“先截断再规范化”的顺序，因此一般不满足交换。
相反，折叠感知 restriction
\[
\widetilde r_{m+1\to m}:=\pi_{m+1\to m}\circ \Fold_{m+1}\qquad(\Omega_{m+1}\to X_m\subset\Omega_m)
\]
正是把微态层的降分辨率顺序纠正为“先规范形再截断”，从而恢复严格交换（命题 \ref{prop:fold-omega-commute}）。
\par
并且在如下“自然性”约束下该修正是唯一的：若映射 \(\widehat r_{m+1\to m}:\Omega_{m+1}\to X_m\) 满足
\[
\Fold_m\circ \widehat r_{m+1\to m}=\pi_{m+1\to m}\circ \Fold_{m+1},
\]
则必有 \(\widehat r_{m+1\to m}=\widetilde r_{m+1\to m}\)。
\leanverified{paper\_normalize\_then\_truncate\_functorial}
\end{theorem}

\begin{proof}
前两段只是对定义 \ref{def:fold-alpha} 与命题 \ref{prop:fold-alpha-basic}-(4) 的逻辑展开，以及对小节 \ref{subsec:fold-omega-restriction} 中“朴素截断不交换、折叠感知 restriction 交换”的机制解释。
\par
对唯一性：若 \(\widehat r_{m+1\to m}(\omega)\in X_m\)，则由二元折叠的幂等性（定理 \ref{thm:fold-suite}-(2)，即 \(\Fold_m|_{X_m}=\mathrm{id}\)）有
\[
(\Fold_m\circ \widehat r_{m+1\to m})(\omega)=\widehat r_{m+1\to m}(\omega).
\]
与假设 \(\Fold_m\circ \widehat r_{m+1\to m}=\pi_{m+1\to m}\circ \Fold_{m+1}\) 合并即得
\(
\widehat r_{m+1\to m}(\omega)=\pi_{m+1\to m}(\Fold_{m+1}(\omega))=\widetilde r_{m+1\to m}(\omega)
\)
对所有 \(\omega\) 成立。证毕。
\end{proof}

\begin{corollary}[跨分辨率可审计不变量的判据]\label{cor:multiscale-auditable-invariants}
\leanverified{paper\_multiscale\_auditable\_invariants}
若某一“跨分辨率规律/指纹”依赖于使用朴素截断 \(r_{m+1\to m}\) 而非折叠感知 restriction \(\widetilde r_{m+1\to m}\)，则它不具备命题 \ref{prop:fold-omega-commute} 的自然性，因而只能被视为由降分辨率规范选择诱发的 gauge 现象，而不能进入跨分辨率的可审计不变量库。
\end{corollary}

\begin{remark}[与本文二元读出的关系]
本文的主线实验固定二元读出与黄金语法（便于得到闭式 PF/Parry/$\zeta$ 基线）。上述一般化表明：若允许更一般的计数/地址字母表（或更一般的观测通道），则“折叠 = 取值 + Ostrowski 规范形”的机制可原则化推广；黄金分支只是其中最简、最可闭式分析的特例。
\end{remark}

\begin{theorem}[二次无理 \texorpdfstring{$\alpha$}{alpha} 下折叠可观测的有限状态矩阵化：压力函数的解析性与代数可计算性]\label{thm:ostrowski-quadratic-irrational-pressure-matrix-algebraicity}
设 \(\alpha\in(0,1)\) 为二次无理数，并取定义 \ref{def:fold-alpha} 的 Ostrowski 折叠 \(\Fold_m^{(\alpha)}\)。
固定一个长度 \(m\) 的可加观测量 \(G_m\)，其由规范化/折叠过程的有限窗口局部贡献逐步累加得到（例如由朴素截断与折叠感知截断之间的逐位失配指示所加和；在二元 Fold 的情形可取定义 \ref{def:fold-gauge-anomaly} 的规范差 \(G_m\)）。
\par
令 \(\mathbb P\) 为一族乘积分布，使输入数字 \(\omega_n\) 独立且其边缘分布仅依赖于相位 \(n\bmod p\)（某个固定 \(p\ge 1\)），并在每个允许取值上赋正概率。
定义对数矩母函数
\[
\Lambda_\alpha(\theta):=\lim_{m\to\infty}\frac1m\log \mathbb E_{\mathbb P}\!\big[e^{\theta G_m}\big],
\]
若极限存在。
则：
\begin{enumerate}
  \item \(\Lambda_\alpha(\theta)\) 在 \(\theta\) 的某邻域内存在且解析；
  \item 存在有限维非负矩阵 \(M_\alpha(\theta)\)，其元素属于由
  \(\{e^\theta\}\cup\{\mathbb P(\omega_n=\cdot)\}\) 生成的系数扩张，使得
  \[
  \Lambda_\alpha(\theta)=\log \rho\bigl(M_\alpha(\theta)\bigr),
  \]
  其中 \(\rho(\cdot)\) 为 Perron--Frobenius 谱半径；
  \item 密度极限
  \[
  g_\alpha:=\lim_{m\to\infty}\frac1m\,\mathbb E_{\mathbb P}[G_m]
  \]
  存在且满足 \(g_\alpha=\Lambda_\alpha'(0)\)；并且当基线概率为有理数时，\(g_\alpha\) 属于上述系数扩张的代数闭包（特别地为代数数）。
\end{enumerate}
\leanverified{paper\_ostrowski\_quadratic\_irrational\_pressure\_matrix\_algebraicity}
\end{theorem}
\begin{proof}
二次无理 \(\alpha\) 的连分数展开最终周期性意味着 Ostrowski 规范化的进位/借位规则在“相位 \(n\bmod p\)”意义下为有限型；
因此存在一个有限状态在线转导器，其状态可取为“局部寄存器（携带规范化所需的有限记忆）\(+\) 相位”，输入为 \(\omega\)，输出为 \(\Fold^{(\alpha)}(\omega)\)，并同步产生每一步对 \(G_m\) 的局部贡献。
\par
对任意 \(\theta\) 在该有限状态图上对每条转移边赋权
\[
w_\theta(\text{edge})
=\mathbb P(\text{该输入符号})\cdot e^{\theta\cdot(\text{该步贡献})},
\]
得到加权邻接矩阵 \(M_\alpha(\theta)\)。
于是存在向量 \(\mathbf u,\mathbf v\) 使
\(
\mathbb E[e^{\theta G_m}]=\mathbf u^\top M_\alpha(\theta)^m\mathbf v
\)。
由于 \(M_\alpha(\theta)\) 非负且在 \(\theta\) 的小邻域内保持原始性（由全支持与强连通性保证），Perron--Frobenius 定理给出
\[
\lim_{m\to\infty}\frac1m\log \mathbf u^\top M_\alpha(\theta)^m\mathbf v=\log\rho(M_\alpha(\theta)),
\]
并且主特征值随 \(\theta\) 解析变化，从而得到 (1)(2)。
对 \(\theta\) 在 \(0\) 处求导并用解析性确保交换极限与导数，即得 (3)。证毕。
\end{proof}

